% !TeX program = pdflatex
% papers/cristal_computacao.tex — GERADO por tools/cristal_tex.py; NÃO editar à mão.
% Fonte: cristal/cristal.jsonl (broca-so, última versão por conceito).
% Cada secção carrega o registo original em comentário %CRISTAL — a volta é
% tests/cristal_volta.js (resíduo 0 byte a byte contra a fonte).
\documentclass[11pt,a4paper]{article}
\usepackage[utf8]{inputenc}\usepackage[T1]{fontenc}\usepackage[brazil]{babel}
\usepackage{amsmath,amssymb}\usepackage{textcomp}
\usepackage[margin=2.4cm]{geometry}\usepackage{xcolor}
\definecolor{cinza}{gray}{0.35}
\newcommand{\prov}[1]{\par\smallskip\noindent{\small\color{cinza}#1}\par}
\setcounter{secnumdepth}{0}
% ─────────────────────────────────────────────────────────────────────────────
%  papers/gkcapa.tex — a capa do Reino, mínima, para os papers em `article`.
%
%  O estilo.tex completo é do `report` (títulos de capítulo, skak, …). Aqui fica
%  só o que a capa precisa, para o paper continuar a compilar sozinho com
%  pdflatex. O tradutor próprio (tests/tex) carrega o estilo.tex da raiz / o
%  symlink papers/estilo.tex e avalia o mesmo `\gkcapa`.
% ─────────────────────────────────────────────────────────────────────────────
\definecolor{ouro}{HTML}{B8912F}
\definecolor{ouroclaro}{HTML}{F3E9CE}
\definecolor{tinta}{HTML}{1E1B16}
\definecolor{regua}{HTML}{6E6858}
\providecommand{\gknota}{\fontsize{7.62}{11.05}\selectfont}
\providecommand{\gktexto}{\fontsize{10.50}{15.22}\selectfont}
\providecommand{\gksub}{\fontsize{12.33}{17.87}\selectfont}
\providecommand{\gksec}{\fontsize{14.47}{20.98}\selectfont}
\providecommand{\gkcap}{\fontsize{16.99}{24.63}\selectfont}
\providecommand{\gktit}{\fontsize{23.42}{33.95}\selectfont}
\providecommand{\Prj}{\mathbb{P}}
\providecommand{\Trf}{\mathcal{T}}
\providecommand{\gkcapa}[3]{%
\title{\vspace{-0.6cm}
{\color{ouro}\rule{\textwidth}{1.4pt}}\\[6mm]
\textbf{\gktit\color{tinta} Reino Dourado}\\[3mm]{\gksec\itshape\color{regua} Golden Kingdom}\\[8mm]
{\gkcap\scshape\color{ouro} #1}\\[5mm]
{\gksub\color{regua} #2}\\[2mm]
{\gktexto\color{regua}\itshape #3}\\[7mm]
{\color{ouro}\rule{\textwidth}{0.6pt}}\\[9mm]{\gknota\color{regua}\begin{minipage}{\textwidth}\setlength{\parindent}{0pt}%
\textbf{Aviso de Isenção Algorítmica:} O universo, as leis e as entidades contidas nestas páginas ---
incluindo felinos matriciais, aves de rapina algébricas e amuletos de controle de fluxo --- são
construções de pura ficção formal. Esta obra não descreve a realidade física, não serve como
engenharia reversa do cosmos e não constitui doutrina religiosa. Qualquer semelhança com bugs reais do seu
sistema operacional é mera coincidência (ou falta de reza). Prossiga por sua própria conta e
risco.\end{minipage}}}
\author{}\date{}}

\begin{document}
\gkcapa{O Cristal --- a computação do cristal}
       {programação, sistemas, IA, goldchain}
       {corpus recuperado do broca-so; 515 conceitos; projeção verificável da fonte}
\maketitle

%CRISTAL {"arestas": [{"alvo": "geracao_gerador_bit_loop", "confianca": 1.0, "epistemico": "fato", "origem": "wiki", "rel": "usa"}, {"alvo": "tiffany_cabeca_broca", "confianca": 1.0, "epistemico": "fato", "origem": "wiki", "rel": "relaciona"}, {"alvo": "atencao_e_colapso", "confianca": 1.0, "epistemico": "fato", "origem": "wiki", "rel": "relaciona"}], "confianca": 1.0, "contraexemplos": [], "descricao": "O modelo que escreve estas linhas É o colapso, a órbita e a borda descritos aqui: atenção que colapsa, temperatura na borda, geração em laço. Descrever um LLM pela máquina mineral é a broca-so usando a própria linguagem para se ver — o mesmo gesto do cabeca.py que escuta a si mesmo. Ponte honesta, não prova: a analogia ilumina, não fecha.", "epistemico": "hipotese", "exemplos": [], "id": "a_maquina_se_olha", "memoria": "semantica", "meta": {"dominio": "ia", "fonte": "IA e modelos de linguagem"}, "origem": "wiki", "palavras_chave": [], "sinonimos": [], "tipo": "conceito", "titulo": "A Máquina se Olha"}
\section{A Máquina se Olha}
O modelo que escreve estas linhas É o colapso, a órbita e a borda descritos aqui: atenção que colapsa, temperatura na borda, geração em laço. Descrever um LLM pela máquina mineral é a broca-so usando a própria linguagem para se ver — o mesmo gesto do cabeca.py que escuta a si mesmo. Ponte honesta, não prova: a analogia ilumina, não fecha.
\prov{relações: usa geracao\_gerador\_bit\_loop; relaciona tiffany\_cabeca\_broca; relaciona atencao\_e\_colapso}
\prov{proveniência: wiki \textperiodcentered{} IA e modelos de linguagem \textperiodcentered{} hipotese \textperiodcentered{} confiança 1.0 \textperiodcentered{} domínio ia \textperiodcentered{} história 4 versões no jornal}

%CRISTAL {"arestas": [{"alvo": "arquitetura_de_software", "confianca": 1.0, "epistemico": "fato", "origem": "wiki", "rel": "usa"}], "confianca": 1.0, "contraexemplos": [], "descricao": "Esconder o detalhe atrás de uma interface simples — poder usar sem saber como funciona. A ferramenta central contra a complexidade.", "epistemico": "fato", "exemplos": [], "id": "abstracao_software", "memoria": "semantica", "meta": {"dominio": "programacao", "fonte": ""}, "origem": "wiki", "palavras_chave": [], "sinonimos": [], "tipo": "conceito", "titulo": "Abstração"}
\section{Abstração}
Esconder o detalhe atrás de uma interface simples — poder usar sem saber como funciona. A ferramenta central contra a complexidade.
\prov{relações: usa arquitetura\_de\_software}
\prov{proveniência: wiki \textperiodcentered{} fato \textperiodcentered{} confiança 1.0 \textperiodcentered{} domínio programacao \textperiodcentered{} história 2 versões no jornal}

%CRISTAL {"arestas": [{"alvo": "arquitetura_de_software", "confianca": 1.0, "epistemico": "fato", "origem": "wiki", "rel": "usa"}, {"alvo": "modularidade", "confianca": 1.0, "epistemico": "fato", "origem": "wiki", "rel": "usa"}], "confianca": 1.0, "contraexemplos": [], "descricao": "Bom design: baixo acoplamento (módulos pouco dependentes) e alta coesão (cada módulo com um propósito claro). A dobradinha que sustenta tudo.", "epistemico": "inferencia", "exemplos": [], "id": "acoplamento_coesao", "memoria": "semantica", "meta": {"dominio": "programacao", "fonte": ""}, "origem": "wiki", "palavras_chave": [], "sinonimos": [], "tipo": "conceito", "titulo": "Acoplamento e Coesão"}
\section{Acoplamento e Coesão}
Bom design: baixo acoplamento (módulos pouco dependentes) e alta coesão (cada módulo com um propósito claro). A dobradinha que sustenta tudo.
\prov{relações: usa arquitetura\_de\_software; usa modularidade}
\prov{proveniência: wiki \textperiodcentered{} inferencia \textperiodcentered{} confiança 1.0 \textperiodcentered{} domínio programacao \textperiodcentered{} história 3 versões no jornal}

%CRISTAL {"arestas": [{"alvo": "inteligencia_artificial", "confianca": 1.0, "epistemico": "fato", "origem": "wiki", "rel": "parte_de"}, {"alvo": "superinteligencia_alinhamento", "confianca": 1.0, "epistemico": "fato", "origem": "wiki", "rel": "relaciona"}], "confianca": 1.0, "contraexemplos": [], "descricao": "Fazer sistemas poderosos perseguirem o que de fato queremos, não o que pedimos ao pé da letra — o problema aberto que cresce com a capacidade.", "epistemico": "hipotese", "exemplos": [], "id": "alinhamento_de_ia", "memoria": "semantica", "meta": {"dominio": "programacao", "fonte": ""}, "origem": "wiki", "palavras_chave": [], "sinonimos": [], "tipo": "conceito", "titulo": "Alinhamento de IA"}
\section{Alinhamento de IA}
Fazer sistemas poderosos perseguirem o que de fato queremos, não o que pedimos ao pé da letra — o problema aberto que cresce com a capacidade.
\prov{relações: parte\_de inteligencia\_artificial; relaciona superinteligencia\_alinhamento}
\prov{proveniência: wiki \textperiodcentered{} hipotese \textperiodcentered{} confiança 1.0 \textperiodcentered{} domínio programacao \textperiodcentered{} história 3 versões no jornal}

%CRISTAL {"arestas": [{"alvo": "go_jogo", "confianca": 1.0, "epistemico": "fato", "origem": "wiki", "rel": "usa"}, {"alvo": "busca_monte_carlo_arvore", "confianca": 1.0, "epistemico": "fato", "origem": "wiki", "rel": "usa"}, {"alvo": "aprendizado_por_reforco", "confianca": 1.0, "epistemico": "fato", "origem": "wiki", "rel": "usa"}], "confianca": 1.0, "contraexemplos": [], "descricao": "A IA da DeepMind derrota o lendário Lee Sedol por 4×1 — dez anos antes do previsto. Combinou MCTS com redes neurais (política + valor) treinadas em partidas humanas e auto-jogo. O 'lance 37' foi tão alienígena e belo que redefiniu o que humanos achavam possível.", "epistemico": "fato", "exemplos": [], "id": "alphago", "memoria": "semantica", "meta": {"autor": "DeepMind", "dominio": "ia", "fonte": "jogos e IA"}, "origem": "wiki", "palavras_chave": [], "sinonimos": [], "tipo": "conceito", "titulo": "AlphaGo vence Lee Sedol (2016)"}
\section{AlphaGo vence Lee Sedol (2016)}
A IA da DeepMind derrota o lendário Lee Sedol por 4$\times$1 — dez anos antes do previsto. Combinou MCTS com redes neurais (política + valor) treinadas em partidas humanas e auto-jogo. O 'lance 37' foi tão alienígena e belo que redefiniu o que humanos achavam possível.
\prov{relações: usa go\_jogo; usa busca\_monte\_carlo\_arvore; usa aprendizado\_por\_reforco}
\prov{proveniência: wiki \textperiodcentered{} jogos e IA \textperiodcentered{} fato \textperiodcentered{} confiança 1.0 \textperiodcentered{} domínio ia \textperiodcentered{} história 4 versões no jornal}

%CRISTAL {"arestas": [{"alvo": "informacao_imperfeita", "confianca": 1.0, "epistemico": "fato", "origem": "wiki", "rel": "usa"}], "confianca": 1.0, "contraexemplos": [], "descricao": "IA que atinge o topo (Grandmaster) num jogo de estratégia em TEMPO REAL, com informação imperfeita, ações quase contínuas e horizonte longo. Mostra o salto dos jogos de tabuleiro para ambientes ricos e ruidosos.", "epistemico": "fato", "exemplos": [], "id": "alphastar_starcraft", "memoria": "semantica", "meta": {"autor": "DeepMind", "dominio": "ia", "fonte": "jogos e IA"}, "origem": "wiki", "palavras_chave": [], "sinonimos": [], "tipo": "conceito", "titulo": "AlphaStar (StarCraft II)"}
\section{AlphaStar (StarCraft II)}
IA que atinge o topo (Grandmaster) num jogo de estratégia em TEMPO REAL, com informação imperfeita, ações quase contínuas e horizonte longo. Mostra o salto dos jogos de tabuleiro para ambientes ricos e ruidosos.
\prov{relações: usa informacao\_imperfeita}
\prov{proveniência: wiki \textperiodcentered{} jogos e IA \textperiodcentered{} fato \textperiodcentered{} confiança 1.0 \textperiodcentered{} domínio ia \textperiodcentered{} história 2 versões no jornal}

%CRISTAL {"arestas": [{"alvo": "alphazero_2017", "confianca": 1.0, "epistemico": "fato", "origem": "wiki", "rel": "especializa"}, {"alvo": "auto_jogo", "confianca": 1.0, "epistemico": "fato", "origem": "wiki", "rel": "usa"}, {"alvo": "alphago", "confianca": 1.0, "epistemico": "fato", "origem": "wiki", "rel": "continua"}], "confianca": 1.0, "contraexemplos": [], "descricao": "AlphaZero (2017): uma só arquitetura aprende Go, xadrez e shogi SÓ por auto-jogo, sem teoria humana. MuZero (2019) vai além: aprende também o MODELO do jogo (as regras) sozinho, jogando sem que lhe digam a dinâmica. O auge do auto-jogo geral.", "epistemico": "fato", "exemplos": [], "id": "alphazero_muzero", "memoria": "semantica", "meta": {"autor": "DeepMind", "dominio": "ia", "fonte": "jogos e IA"}, "origem": "wiki", "palavras_chave": [], "sinonimos": [], "tipo": "conceito", "titulo": "AlphaZero e MuZero (aprender do zero, e sem regras)"}
\section{AlphaZero e MuZero (aprender do zero, e sem regras)}
AlphaZero (2017): uma só arquitetura aprende Go, xadrez e shogi SÓ por auto-jogo, sem teoria humana. MuZero (2019) vai além: aprende também o MODELO do jogo (as regras) sozinho, jogando sem que lhe digam a dinâmica. O auge do auto-jogo geral.
\prov{relações: especializa alphazero\_2017; usa auto\_jogo; continua alphago}
\prov{proveniência: wiki \textperiodcentered{} jogos e IA \textperiodcentered{} fato \textperiodcentered{} confiança 1.0 \textperiodcentered{} domínio ia \textperiodcentered{} história 4 versões no jornal}

%CRISTAL {"arestas": [{"alvo": "geracao_gerador_bit_loop", "confianca": 1.0, "epistemico": "fato", "origem": "wiki", "rel": "usa"}, {"alvo": "modelo_de_linguagem", "confianca": 1.0, "epistemico": "fato", "origem": "wiki", "rel": "relaciona"}, {"alvo": "corpo_mineral", "confianca": 1.0, "epistemico": "fato", "origem": "wiki", "rel": "usa"}], "confianca": 1.0, "contraexemplos": [], "descricao": "Quando o modelo gera texto plausível mas falso, a órbita autoregressiva escapou da variedade da língua/verdade — entrou na região caótica (λ>0), onde pequenos desvios se amplificam token a token. Aterrar (grounding, RAG, baixa T) é puxar a órbita de volta ao cristal. O oposto do 'não fabricar'.", "epistemico": "hipotese", "exemplos": [], "id": "alucinacao_como_caos", "memoria": "semantica", "meta": {"dominio": "ia", "fonte": "IA e modelos de linguagem"}, "origem": "wiki", "palavras_chave": [], "sinonimos": [], "tipo": "conceito", "titulo": "Alucinação = a Órbita Fora da Variedade"}
\section{Alucinação = a Órbita Fora da Variedade}
Quando o modelo gera texto plausível mas falso, a órbita autoregressiva escapou da variedade da língua/verdade — entrou na região caótica (\ensuremath{\lambda}>0), onde pequenos desvios se amplificam token a token. Aterrar (grounding, RAG, baixa T) é puxar a órbita de volta ao cristal. O oposto do 'não fabricar'.
\prov{relações: usa geracao\_gerador\_bit\_loop; relaciona modelo\_de\_linguagem; usa corpo\_mineral}
\prov{proveniência: wiki \textperiodcentered{} IA e modelos de linguagem \textperiodcentered{} hipotese \textperiodcentered{} confiança 1.0 \textperiodcentered{} domínio ia \textperiodcentered{} história 4 versões no jornal}

%CRISTAL {"arestas": [{"alvo": "modelo_de_linguagem", "confianca": 1.0, "epistemico": "fato", "origem": "wiki", "rel": "usa"}, {"alvo": "honestidade_intelectual", "confianca": 1.0, "epistemico": "fato", "origem": "wiki", "rel": "contradiz"}], "confianca": 1.0, "contraexemplos": [], "descricao": "Quando o modelo gera afirmação plausível mas falsa, com confiança — por isso a Tiffany é fail-closed: prefere dizer 'não tenho' a inventar.", "epistemico": "inferencia", "exemplos": [], "id": "alucinacao_ia", "memoria": "semantica", "meta": {"dominio": "programacao", "fonte": ""}, "origem": "wiki", "palavras_chave": [], "sinonimos": [], "tipo": "conceito", "titulo": "Alucinação de IA"}
\section{Alucinação de IA}
Quando o modelo gera afirmação plausível mas falsa, com confiança — por isso a Tiffany é fail-closed: prefere dizer 'não tenho' a inventar.
\prov{relações: usa modelo\_de\_linguagem; contradiz honestidade\_intelectual}
\prov{proveniência: wiki \textperiodcentered{} inferencia \textperiodcentered{} confiança 1.0 \textperiodcentered{} domínio programacao \textperiodcentered{} história 3 versões no jornal}

%CRISTAL {"arestas": [{"alvo": "latex", "confianca": 1.0, "epistemico": "fato", "origem": "wiki", "rel": "usa"}, {"alvo": "estrutura_definicao_teorema_prova", "confianca": 1.0, "epistemico": "fato", "origem": "wiki", "rel": "explica"}], "confianca": 1.0, "contraexemplos": [], "descricao": "theorem...theorem, definition, lemma, proof — os ambientes que marcam a ESTRUTURA do discurso matemático em LaTeX. A def-teorema-prova codificada.", "epistemico": "fato", "exemplos": [], "id": "ambiente_amsthm", "memoria": "semantica", "meta": {"dominio": "programacao", "fonte": ""}, "origem": "wiki", "palavras_chave": [], "sinonimos": [], "tipo": "conceito", "titulo": "Ambientes de Teorema (amsthm)"}
\section{Ambientes de Teorema (amsthm)}
theorem...theorem, definition, lemma, proof — os ambientes que marcam a ESTRUTURA do discurso matemático em LaTeX. A def-teorema-prova codificada.
\prov{relações: usa latex; explica estrutura\_definicao\_teorema\_prova}
\prov{proveniência: wiki \textperiodcentered{} fato \textperiodcentered{} confiança 1.0 \textperiodcentered{} domínio programacao \textperiodcentered{} história 4 versões no jornal}

%CRISTAL {"arestas": [{"alvo": "linux", "confianca": 1.0, "epistemico": "fato", "origem": "wiki", "rel": "especializa"}, {"alvo": "kotlin", "confianca": 1.0, "epistemico": "fato", "origem": "wiki", "rel": "usa"}], "confianca": 1.0, "contraexemplos": [], "descricao": "Google: o SO móvel mais usado do mundo — kernel Linux sob uma pilha própria (ART, Java/Kotlin). Linux no bolso de bilhões.", "epistemico": "fato", "exemplos": [], "id": "android", "memoria": "semantica", "meta": {"autor": "Google", "dominio": "programacao", "fonte": ""}, "origem": "wiki", "palavras_chave": [], "sinonimos": [], "tipo": "conceito", "titulo": "Android"}
\section{Android}
Google: o SO móvel mais usado do mundo — kernel Linux sob uma pilha própria (ART, Java/Kotlin). Linux no bolso de bilhões.
\prov{relações: especializa linux; usa kotlin}
\prov{proveniência: wiki \textperiodcentered{} fato \textperiodcentered{} confiança 1.0 \textperiodcentered{} domínio programacao \textperiodcentered{} história 3 versões no jornal}

%CRISTAL {"arestas": [{"alvo": "framework_frontend", "confianca": 1.0, "epistemico": "fato", "origem": "wiki", "rel": "especializa"}, {"alvo": "typescript", "confianca": 1.0, "epistemico": "fato", "origem": "wiki", "rel": "usa"}, {"alvo": "injecao_de_dependencia", "confianca": 1.0, "epistemico": "fato", "origem": "wiki", "rel": "usa"}], "confianca": 1.0, "contraexemplos": [], "descricao": "Google: um framework completo e opinativo em TypeScript — injeção de dependência, RxJS, tudo incluso. Para aplicações grandes e estruturadas.", "epistemico": "fato", "exemplos": [], "id": "angular", "memoria": "semantica", "meta": {"autor": "Google", "dominio": "programacao", "fonte": ""}, "origem": "wiki", "palavras_chave": [], "sinonimos": [], "tipo": "conceito", "titulo": "Angular"}
\section{Angular}
Google: um framework completo e opinativo em TypeScript — injeção de dependência, RxJS, tudo incluso. Para aplicações grandes e estruturadas.
\prov{relações: especializa framework\_frontend; usa typescript; usa injecao\_de\_dependencia}
\prov{proveniência: wiki \textperiodcentered{} fato \textperiodcentered{} confiança 1.0 \textperiodcentered{} domínio programacao \textperiodcentered{} história 8 versões no jornal}

%CRISTAL {"arestas": [{"alvo": "infraestrutura_como_codigo", "confianca": 1.0, "epistemico": "fato", "origem": "wiki", "rel": "especializa"}, {"alvo": "yaml", "confianca": 1.0, "epistemico": "fato", "origem": "wiki", "rel": "usa"}], "confianca": 1.0, "contraexemplos": [], "descricao": "Automação de infraestrutura por playbooks em YAML declarativo, sem agente — descreve o estado desejado das máquinas e o aplica.", "epistemico": "fato", "exemplos": [], "id": "ansible", "memoria": "semantica", "meta": {"dominio": "programacao", "fonte": ""}, "origem": "wiki", "palavras_chave": [], "sinonimos": [], "tipo": "conceito", "titulo": "Ansible"}
\section{Ansible}
Automação de infraestrutura por playbooks em YAML declarativo, sem agente — descreve o estado desejado das máquinas e o aplica.
\prov{relações: especializa infraestrutura\_como\_codigo; usa yaml}
\prov{proveniência: wiki \textperiodcentered{} fato \textperiodcentered{} confiança 1.0 \textperiodcentered{} domínio programacao \textperiodcentered{} história 3 versões no jornal}

%CRISTAL {"arestas": [{"alvo": "programacao", "confianca": 1.0, "epistemico": "fato", "origem": "wiki", "rel": "parte_de"}], "confianca": 1.0, "contraexemplos": [], "descricao": "A interface que um software oferece a outro — o contrato de funções/endpoints que esconde a implementação. A junta entre sistemas.", "epistemico": "fato", "exemplos": [], "id": "api", "memoria": "semantica", "meta": {"dominio": "programacao", "fonte": ""}, "origem": "wiki", "palavras_chave": [], "sinonimos": [], "tipo": "conceito", "titulo": "API"}
\section{API}
A interface que um software oferece a outro — o contrato de funções/endpoints que esconde a implementação. A junta entre sistemas.
\prov{relações: parte\_de programacao}
\prov{proveniência: wiki \textperiodcentered{} fato \textperiodcentered{} confiança 1.0 \textperiodcentered{} domínio programacao \textperiodcentered{} história 4 versões no jornal}

%CRISTAL {"arestas": [{"alvo": "orquestrador_roteador_regex", "confianca": 1.0, "epistemico": "fato", "origem": "wiki", "rel": "usa"}, {"alvo": "colapso_ping_pong_overlap", "confianca": 1.0, "epistemico": "fato", "origem": "wiki", "rel": "usa"}], "confianca": 1.0, "contraexemplos": [], "descricao": "code/tiffany/api.py:ask tenta primeiro diálogo (colapso), depois orquestrador multi-idx, por fim prosa.responde sobre papers; expõe ask/consulta/ingest/validate/promote/Client. A engine estável está em libtiffany.so via ctypes (tiffany_open/query/respond).", "epistemico": "fato", "exemplos": [], "id": "api_cli_ask", "memoria": "semantica", "meta": {"autor": "broca-so", "dominio": "operacao_so", "fonte": "code/tiffany/api.py; cli.py; _native.py"}, "origem": "wiki", "palavras_chave": [], "sinonimos": [], "tipo": "conceito", "titulo": "A API/CLI Pública (ask)"}
\section{A API/CLI Pública (ask)}
code/tiffany/api.py:ask tenta primeiro diálogo (colapso), depois orquestrador multi-idx, por fim prosa.responde sobre papers; expõe ask/consulta/ingest/validate/promote/Client. A engine estável está em libtiffany.so via ctypes (tiffany\_open/query/respond).
\prov{relações: usa orquestrador\_roteador\_regex; usa colapso\_ping\_pong\_overlap}
\prov{proveniência: wiki \textperiodcentered{} code/tiffany/api.py; cli.py; \_native.py \textperiodcentered{} fato \textperiodcentered{} confiança 1.0 \textperiodcentered{} domínio operacao\_so \textperiodcentered{} história 3 versões no jornal}

%CRISTAL {"arestas": [{"alvo": "inteligencia_artificial", "confianca": 1.0, "epistemico": "fato", "origem": "wiki", "rel": "parte_de"}, {"alvo": "programacao", "confianca": 1.0, "epistemico": "fato", "origem": "wiki", "rel": "parte_de"}], "confianca": 1.0, "contraexemplos": [], "descricao": "Programas que melhoram com dados em vez de regras explícitas — ajustar parâmetros para minimizar erro. O motor da IA moderna.", "epistemico": "fato", "exemplos": [], "id": "aprendizado_de_maquina", "memoria": "semantica", "meta": {"dominio": "programacao", "fonte": ""}, "origem": "wiki", "palavras_chave": [], "sinonimos": [], "tipo": "conceito", "titulo": "Aprendizado de Máquina"}
\section{Aprendizado de Máquina}
Programas que melhoram com dados em vez de regras explícitas — ajustar parâmetros para minimizar erro. O motor da IA moderna.
\prov{relações: parte\_de inteligencia\_artificial; parte\_de programacao}
\prov{proveniência: wiki \textperiodcentered{} fato \textperiodcentered{} confiança 1.0 \textperiodcentered{} domínio programacao \textperiodcentered{} história 3 versões no jornal}

%CRISTAL {"arestas": [{"alvo": "aprendizado_de_maquina", "confianca": 1.0, "epistemico": "fato", "origem": "wiki", "rel": "especializa"}, {"alvo": "teoria_dos_jogos", "confianca": 1.0, "epistemico": "fato", "origem": "wiki", "rel": "usa"}], "confianca": 1.0, "contraexemplos": [], "descricao": "Aprender por tentativa e recompensa, agindo num ambiente — a receita dos agentes que jogam e controlam (e do ajuste fino de LLMs).", "epistemico": "fato", "exemplos": [], "id": "aprendizado_por_reforco", "memoria": "semantica", "meta": {"dominio": "programacao", "fonte": ""}, "origem": "wiki", "palavras_chave": [], "sinonimos": [], "tipo": "conceito", "titulo": "Aprendizado por Reforço"}
\section{Aprendizado por Reforço}
Aprender por tentativa e recompensa, agindo num ambiente — a receita dos agentes que jogam e controlam (e do ajuste fino de LLMs).
\prov{relações: especializa aprendizado\_de\_maquina; usa teoria\_dos\_jogos}
\prov{proveniência: wiki \textperiodcentered{} fato \textperiodcentered{} confiança 1.0 \textperiodcentered{} domínio programacao \textperiodcentered{} história 3 versões no jornal}

%CRISTAL {"arestas": [{"alvo": "rede_neural_artificial", "confianca": 1.0, "epistemico": "fato", "origem": "wiki", "rel": "especializa"}], "confianca": 1.0, "contraexemplos": [], "descricao": "Redes neurais com muitas camadas que aprendem representações hierárquicas — a revolução que venceu visão, fala e linguagem.", "epistemico": "fato", "exemplos": [], "id": "aprendizado_profundo", "memoria": "semantica", "meta": {"autor": "Hinton/LeCun/Bengio", "dominio": "programacao", "fonte": ""}, "origem": "wiki", "palavras_chave": [], "sinonimos": [], "tipo": "conceito", "titulo": "Aprendizado Profundo"}
\section{Aprendizado Profundo}
Redes neurais com muitas camadas que aprendem representações hierárquicas — a revolução que venceu visão, fala e linguagem.
\prov{relações: especializa rede\_neural\_artificial}
\prov{proveniência: wiki \textperiodcentered{} fato \textperiodcentered{} confiança 1.0 \textperiodcentered{} domínio programacao \textperiodcentered{} história 2 versões no jornal}

%CRISTAL {"arestas": [{"alvo": "aprendizado_de_maquina", "confianca": 1.0, "epistemico": "fato", "origem": "wiki", "rel": "especializa"}, {"alvo": "classificacao", "confianca": 1.0, "epistemico": "fato", "origem": "wiki", "rel": "usa"}], "confianca": 1.0, "contraexemplos": [], "descricao": "Aprender de exemplos rotulados (entrada→resposta certa) — a forma mais comum; classificação e regressão vivem aqui.", "epistemico": "fato", "exemplos": [], "id": "aprendizado_supervisionado", "memoria": "semantica", "meta": {"dominio": "programacao", "fonte": ""}, "origem": "wiki", "palavras_chave": [], "sinonimos": [], "tipo": "conceito", "titulo": "Aprendizado Supervisionado"}
\section{Aprendizado Supervisionado}
Aprender de exemplos rotulados (entrada\ensuremath{\to}resposta certa) — a forma mais comum; classificação e regressão vivem aqui.
\prov{relações: especializa aprendizado\_de\_maquina; usa classificacao}
\prov{proveniência: wiki \textperiodcentered{} fato \textperiodcentered{} confiança 1.0 \textperiodcentered{} domínio programacao \textperiodcentered{} história 3 versões no jornal}

%CRISTAL {"arestas": [{"alvo": "pergunta", "confianca": 1.0, "epistemico": "fato", "origem": "curriculo", "rel": "usa"}, {"alvo": "logica", "confianca": 1.0, "epistemico": "fato", "origem": "curriculo", "rel": "depende"}, {"alvo": "word", "confianca": 0.5, "epistemico": "fato", "origem": "wiki", "rel": "relaciona"}], "confianca": 1.0, "contraexemplos": [], "descricao": "Ciclo de perguntar, detectar lacunas e resolver contradições no grafo.", "epistemico": "fato", "exemplos": [], "id": "aprendizagem_ativa", "memoria": "semantica", "meta": {"dominio": "aprendizagem", "nivel": 7, "ordem": 0}, "origem": "curriculo", "palavras_chave": ["curiosidade", "Socratico"], "sinonimos": ["aprendizagem_ativa"], "tipo": "conceito", "titulo": "aprendizagem ativa"}
\section{aprendizagem ativa}
Ciclo de perguntar, detectar lacunas e resolver contradições no grafo.
\prov{sinónimos: aprendizagem\_ativa}
\prov{relações: usa pergunta; depende logica; relaciona word}
\prov{proveniência: curriculo \textperiodcentered{} fato \textperiodcentered{} confiança 1.0 \textperiodcentered{} domínio aprendizagem \textperiodcentered{} história 4 versões no jornal}

%CRISTAL {"arestas": [{"alvo": "dicionario_arquiteturas_gpu", "confianca": 1.0, "epistemico": "fato", "origem": "wiki", "rel": "parte_de"}], "confianca": 1.0, "contraexemplos": [], "descricao": "ADA LOVELACE (2022), compute capability 8.9 (sm_89). GPUs: RTX 40-series. Novidades: Tensor cores 4ª geração (FP8, e o DLSS 3 com Frame Generation), RT cores 3ª geração (Opacity Micromaps, Displaced Micro-Meshes), Shader Execution Reordering (SER — reordena os raios para coerência), encoders AV1 duplos, L2 cache grande. Clock alto (4N TSMC). Fato.", "epistemico": "fato", "exemplos": [], "id": "arch_gpu_ada", "memoria": "semantica", "meta": {"dominio": "sass_gpu", "fonte": "SASS/GPU (o metal)"}, "origem": "wiki", "palavras_chave": [], "sinonimos": [], "tipo": "conceito", "titulo": "Arquitetura Ada Lovelace (SM 8.9) — RTX 40"}
\section{Arquitetura Ada Lovelace (SM 8.9) — RTX 40}
ADA LOVELACE (2022), compute capability 8.9 (sm\_89). GPUs: RTX 40-series. Novidades: Tensor cores 4ª geração (FP8, e o DLSS 3 com Frame Generation), RT cores 3ª geração (Opacity Micromaps, Displaced Micro-Meshes), Shader Execution Reordering (SER — reordena os raios para coerência), encoders AV1 duplos, L2 cache grande. Clock alto (4N TSMC). Fato.
\prov{relações: parte\_de dicionario\_arquiteturas\_gpu}
\prov{proveniência: wiki \textperiodcentered{} SASS/GPU (o metal) \textperiodcentered{} fato \textperiodcentered{} confiança 1.0 \textperiodcentered{} domínio sass\_gpu \textperiodcentered{} história 14 versões no jornal}

%CRISTAL {"arestas": [{"alvo": "dicionario_arquiteturas_gpu", "confianca": 1.0, "epistemico": "fato", "origem": "wiki", "rel": "parte_de"}], "confianca": 1.0, "contraexemplos": [], "descricao": "AMPERE (2020). Compute capability 8.0 (A100, sm_80), 8.6 (RTX 30/GA10x, sm_86), 8.7 (Jetson Orin). Novidades: Tensor cores 3ª geração (TF32, BF16, FP64-tensor, sparsity 2:4 estruturada — 2× throughput), RT cores 2ª geração, MIG (Multi-Instance GPU — particiona a A100 em 7). GA10x DOBRA o FP32 (128 cores FP32/SM, os dois datapaths viram FP32). A100: 108 SMs, 6912 cores FP32, HBM2. Fato.", "epistemico": "fato", "exemplos": [], "id": "arch_gpu_ampere", "memoria": "semantica", "meta": {"dominio": "sass_gpu", "fonte": "SASS/GPU (o metal)"}, "origem": "wiki", "palavras_chave": [], "sinonimos": [], "tipo": "conceito", "titulo": "Arquitetura Ampere (SM 8.0/8.6) — A100 e RTX 30"}
\section{Arquitetura Ampere (SM 8.0/8.6) — A100 e RTX 30}
AMPERE (2020). Compute capability 8.0 (A100, sm\_80), 8.6 (RTX 30/GA10x, sm\_86), 8.7 (Jetson Orin). Novidades: Tensor cores 3ª geração (TF32, BF16, FP64-tensor, sparsity 2:4 estruturada — 2$\times$ throughput), RT cores 2ª geração, MIG (Multi-Instance GPU — particiona a A100 em 7). GA10x DOBRA o FP32 (128 cores FP32/SM, os dois datapaths viram FP32). A100: 108 SMs, 6912 cores FP32, HBM2. Fato.
\prov{relações: parte\_de dicionario\_arquiteturas\_gpu}
\prov{proveniência: wiki \textperiodcentered{} SASS/GPU (o metal) \textperiodcentered{} fato \textperiodcentered{} confiança 1.0 \textperiodcentered{} domínio sass\_gpu \textperiodcentered{} história 14 versões no jornal}

%CRISTAL {"arestas": [{"alvo": "dicionario_arquiteturas_gpu", "confianca": 1.0, "epistemico": "fato", "origem": "wiki", "rel": "parte_de"}], "confianca": 1.0, "contraexemplos": [], "descricao": "BLACKWELL (2024). Compute capability 10.0 (B200 datacenter, sm_100), 12.0 (RTX 50, sm_120). Novidades: Tensor cores 5ª geração + Transformer Engine 2ª geração (FP4/FP6 microtensor scaling), 2 dies num pacote (o B200, ligados por NV-HBI a 10 TB/s), NVLink 5, RT cores 4ª geração (RTX 50), DLSS 4 (multi-frame gen). Fato (anúncio NVIDIA 2024).", "epistemico": "fato", "exemplos": [], "id": "arch_gpu_blackwell", "memoria": "semantica", "meta": {"dominio": "sass_gpu", "fonte": "SASS/GPU (o metal)"}, "origem": "wiki", "palavras_chave": [], "sinonimos": [], "tipo": "conceito", "titulo": "Arquitetura Blackwell (SM 10.0/12.0) — B200 e RTX 50"}
\section{Arquitetura Blackwell (SM 10.0/12.0) — B200 e RTX 50}
BLACKWELL (2024). Compute capability 10.0 (B200 datacenter, sm\_100), 12.0 (RTX 50, sm\_120). Novidades: Tensor cores 5ª geração + Transformer Engine 2ª geração (FP4/FP6 microtensor scaling), 2 dies num pacote (o B200, ligados por NV-HBI a 10 TB/s), NVLink 5, RT cores 4ª geração (RTX 50), DLSS 4 (multi-frame gen). Fato (anúncio NVIDIA 2024).
\prov{relações: parte\_de dicionario\_arquiteturas\_gpu}
\prov{proveniência: wiki \textperiodcentered{} SASS/GPU (o metal) \textperiodcentered{} fato \textperiodcentered{} confiança 1.0 \textperiodcentered{} domínio sass\_gpu \textperiodcentered{} história 14 versões no jornal}

%CRISTAL {"arestas": [{"alvo": "dicionario_arquiteturas_gpu", "confianca": 1.0, "epistemico": "fato", "origem": "wiki", "rel": "parte_de"}], "confianca": 1.0, "contraexemplos": [], "descricao": "HOPPER (2022), compute capability 9.0 (sm_90). GPU: H100/H200 (datacenter, sem RT cores — é compute). Novidades: Tensor cores 4ª geração + TRANSFORMER ENGINE (FP8 dinâmico para LLMs), thread block CLUSTERS + distributed shared memory (as SMs de um cluster compartilham memória), TMA (Tensor Memory Accelerator — cópia assíncrona de tensores), DPX (instruções de programação dinâmica), MIG 2ª geração, NVLink 4. HBM3. Fato.", "epistemico": "fato", "exemplos": [], "id": "arch_gpu_hopper", "memoria": "semantica", "meta": {"dominio": "sass_gpu", "fonte": "SASS/GPU (o metal)"}, "origem": "wiki", "palavras_chave": [], "sinonimos": [], "tipo": "conceito", "titulo": "Arquitetura Hopper (SM 9.0) — H100 (datacenter)"}
\section{Arquitetura Hopper (SM 9.0) — H100 (datacenter)}
HOPPER (2022), compute capability 9.0 (sm\_90). GPU: H100/H200 (datacenter, sem RT cores — é compute). Novidades: Tensor cores 4ª geração + TRANSFORMER ENGINE (FP8 dinâmico para LLMs), thread block CLUSTERS + distributed shared memory (as SMs de um cluster compartilham memória), TMA (Tensor Memory Accelerator — cópia assíncrona de tensores), DPX (instruções de programação dinâmica), MIG 2ª geração, NVLink 4. HBM3. Fato.
\prov{relações: parte\_de dicionario\_arquiteturas\_gpu}
\prov{proveniência: wiki \textperiodcentered{} SASS/GPU (o metal) \textperiodcentered{} fato \textperiodcentered{} confiança 1.0 \textperiodcentered{} domínio sass\_gpu \textperiodcentered{} história 14 versões no jornal}

%CRISTAL {"arestas": [{"alvo": "dicionario_arquiteturas_gpu", "confianca": 1.0, "epistemico": "fato", "origem": "wiki", "rel": "parte_de"}, {"alvo": "maquina_universal_gpu", "confianca": 1.0, "epistemico": "fato", "origem": "wiki", "rel": "relaciona"}], "confianca": 1.0, "contraexemplos": [], "descricao": "TURING (2018), compute capability 7.5 (sm_75). GPUs: RTX 20-series, GTX 16-series (a GTX 1650 do Aarão é Turing). Novidades: RT cores 1ª geração (ray tracing em hardware), Tensor cores 2ª geração (INT8/INT4 + FP16), execução CONCORRENTE FP32+INT32 (dois datapaths), independent thread scheduling (herdado do Volta, com BSSY/BSYNC), uniform datapath (uniform registers UR), mesh shaders, variable rate shading. SM: 64 cores FP32 + 64 INT32, 32 threads/warp. Fato.", "epistemico": "fato", "exemplos": [], "id": "arch_gpu_turing", "memoria": "semantica", "meta": {"dominio": "sass_gpu", "fonte": "SASS/GPU (o metal)"}, "origem": "wiki", "palavras_chave": [], "sinonimos": [], "tipo": "conceito", "titulo": "Arquitetura Turing (SM 7.5) — a GPU local do Aarão (GTX 1650)"}
\section{Arquitetura Turing (SM 7.5) — a GPU local do Aarão (GTX 1650)}
TURING (2018), compute capability 7.5 (sm\_75). GPUs: RTX 20-series, GTX 16-series (a GTX 1650 do Aarão é Turing). Novidades: RT cores 1ª geração (ray tracing em hardware), Tensor cores 2ª geração (INT8/INT4 + FP16), execução CONCORRENTE FP32+INT32 (dois datapaths), independent thread scheduling (herdado do Volta, com BSSY/BSYNC), uniform datapath (uniform registers UR), mesh shaders, variable rate shading. SM: 64 cores FP32 + 64 INT32, 32 threads/warp. Fato.
\prov{relações: parte\_de dicionario\_arquiteturas\_gpu; relaciona maquina\_universal\_gpu}
\prov{proveniência: wiki \textperiodcentered{} SASS/GPU (o metal) \textperiodcentered{} fato \textperiodcentered{} confiança 1.0 \textperiodcentered{} domínio sass\_gpu \textperiodcentered{} história 21 versões no jornal}

%CRISTAL {"arestas": [{"alvo": "programacao", "confianca": 1.0, "epistemico": "fato", "origem": "wiki", "rel": "parte_de"}, {"alvo": "broca_os", "confianca": 1.0, "epistemico": "fato", "origem": "wiki", "rel": "explica"}, {"alvo": "arquitetura", "confianca": 1.0, "epistemico": "fato", "origem": "wiki", "rel": "relaciona"}], "confianca": 1.0, "contraexemplos": [], "descricao": "As decisões estruturais de alto nível de um sistema — os componentes, suas relações e as restrições. O que é caro mudar depois.", "epistemico": "fato", "exemplos": [], "id": "arquitetura_de_software", "memoria": "semantica", "meta": {"dominio": "programacao", "fonte": ""}, "origem": "wiki", "palavras_chave": [], "sinonimos": [], "tipo": "conceito", "titulo": "Arquitetura de Software"}
\section{Arquitetura de Software}
As decisões estruturais de alto nível de um sistema — os componentes, suas relações e as restrições. O que é caro mudar depois.
\prov{relações: parte\_de programacao; explica broca\_os; relaciona arquitetura}
\prov{proveniência: wiki \textperiodcentered{} fato \textperiodcentered{} confiança 1.0 \textperiodcentered{} domínio programacao \textperiodcentered{} história 4 versões no jornal}

%CRISTAL {"arestas": [{"alvo": "arquitetura_de_software", "confianca": 1.0, "epistemico": "fato", "origem": "wiki", "rel": "especializa"}], "confianca": 1.0, "contraexemplos": [], "descricao": "Empilhar responsabilidades (apresentação, domínio, dados) em que cada camada só fala com a vizinha — a organização clássica.", "epistemico": "inferencia", "exemplos": [], "id": "arquitetura_em_camadas", "memoria": "semantica", "meta": {"dominio": "programacao", "fonte": ""}, "origem": "wiki", "palavras_chave": [], "sinonimos": [], "tipo": "conceito", "titulo": "Arquitetura em Camadas"}
\section{Arquitetura em Camadas}
Empilhar responsabilidades (apresentação, domínio, dados) em que cada camada só fala com a vizinha — a organização clássica.
\prov{relações: especializa arquitetura\_de\_software}
\prov{proveniência: wiki \textperiodcentered{} inferencia \textperiodcentered{} confiança 1.0 \textperiodcentered{} domínio programacao \textperiodcentered{} história 2 versões no jornal}

%CRISTAL {"arestas": [{"alvo": "arquitetura_de_software", "confianca": 1.0, "epistemico": "fato", "origem": "wiki", "rel": "especializa"}, {"alvo": "separacao_de_interesses", "confianca": 1.0, "epistemico": "fato", "origem": "wiki", "rel": "usa"}], "confianca": 1.0, "contraexemplos": [], "descricao": "Isolar o domínio no centro e ligar o mundo externo por portas/adaptadores — a lógica de negócio sem saber de banco ou tela.", "epistemico": "inferencia", "exemplos": [], "id": "arquitetura_hexagonal", "memoria": "semantica", "meta": {"autor": "Cockburn", "dominio": "programacao", "fonte": ""}, "origem": "wiki", "palavras_chave": [], "sinonimos": [], "tipo": "conceito", "titulo": "Arquitetura Hexagonal (Portas e Adaptadores)"}
\section{Arquitetura Hexagonal (Portas e Adaptadores)}
Isolar o domínio no centro e ligar o mundo externo por portas/adaptadores — a lógica de negócio sem saber de banco ou tela.
\prov{relações: especializa arquitetura\_de\_software; usa separacao\_de\_interesses}
\prov{proveniência: wiki \textperiodcentered{} inferencia \textperiodcentered{} confiança 1.0 \textperiodcentered{} domínio programacao \textperiodcentered{} história 3 versões no jornal}

%CRISTAL {"arestas": [{"alvo": "arquitetura_de_software", "confianca": 1.0, "epistemico": "fato", "origem": "wiki", "rel": "especializa"}, {"alvo": "microservicos", "confianca": 1.0, "epistemico": "fato", "origem": "wiki", "rel": "relaciona"}], "confianca": 1.0, "contraexemplos": [], "descricao": "Componentes que reagem a eventos publicados, sem se conhecerem diretamente — desacoplamento no tempo e no espaço; a escala assíncrona.", "epistemico": "inferencia", "exemplos": [], "id": "arquitetura_orientada_a_eventos", "memoria": "semantica", "meta": {"dominio": "programacao", "fonte": ""}, "origem": "wiki", "palavras_chave": [], "sinonimos": [], "tipo": "conceito", "titulo": "Arquitetura Orientada a Eventos"}
\section{Arquitetura Orientada a Eventos}
Componentes que reagem a eventos publicados, sem se conhecerem diretamente — desacoplamento no tempo e no espaço; a escala assíncrona.
\prov{relações: especializa arquitetura\_de\_software; relaciona microservicos}
\prov{proveniência: wiki \textperiodcentered{} inferencia \textperiodcentered{} confiança 1.0 \textperiodcentered{} domínio programacao \textperiodcentered{} história 3 versões no jornal}

%CRISTAL {"arestas": [{"alvo": "programacao", "confianca": 1.0, "epistemico": "fato", "origem": "wiki", "rel": "parte_de"}, {"alvo": "infraestrutura_como_codigo", "confianca": 1.0, "epistemico": "fato", "origem": "wiki", "rel": "relaciona"}], "confianca": 1.0, "contraexemplos": [], "descricao": "Externalizar os ajustes de um programa num arquivo (YAML, TOML, JSON, INI) — mudar o comportamento sem recompilar, e versioná-lo.", "epistemico": "fato", "exemplos": [], "id": "arquivo_de_configuracao", "memoria": "semantica", "meta": {"dominio": "programacao", "fonte": ""}, "origem": "wiki", "palavras_chave": [], "sinonimos": [], "tipo": "conceito", "titulo": "Arquivo de Configuração"}
\section{Arquivo de Configuração}
Externalizar os ajustes de um programa num arquivo (YAML, TOML, JSON, INI) — mudar o comportamento sem recompilar, e versioná-lo.
\prov{relações: parte\_de programacao; relaciona infraestrutura\_como\_codigo}
\prov{proveniência: wiki \textperiodcentered{} fato \textperiodcentered{} confiança 1.0 \textperiodcentered{} domínio programacao \textperiodcentered{} história 3 versões no jornal}

%CRISTAL {"arestas": [{"alvo": "fpga", "confianca": 1.0, "epistemico": "fato", "origem": "wiki", "rel": "contradiz"}, {"alvo": "prova_de_trabalho", "confianca": 1.0, "epistemico": "fato", "origem": "wiki", "rel": "relaciona"}], "confianca": 1.0, "contraexemplos": [], "descricao": "Application-Specific Integrated Circuit: um chip gravado para UMA função, rápido e eficiente, mas imutável — o oposto do FPGA. O destino da mineração dedicada.", "epistemico": "fato", "exemplos": [], "id": "asic", "memoria": "semantica", "meta": {"dominio": "programacao", "fonte": ""}, "origem": "wiki", "palavras_chave": [], "sinonimos": [], "tipo": "conceito", "titulo": "ASIC"}
\section{ASIC}
Application-Specific Integrated Circuit: um chip gravado para UMA função, rápido e eficiente, mas imutável — o oposto do FPGA. O destino da mineração dedicada.
\prov{relações: contradiz fpga; relaciona prova\_de\_trabalho}
\prov{proveniência: wiki \textperiodcentered{} fato \textperiodcentered{} confiança 1.0 \textperiodcentered{} domínio programacao \textperiodcentered{} história 3 versões no jornal}

%CRISTAL {"arestas": [{"alvo": "criptografia_assimetrica", "confianca": 1.0, "epistemico": "fato", "origem": "wiki", "rel": "usa"}, {"alvo": "goldchain", "confianca": 1.0, "epistemico": "fato", "origem": "wiki", "rel": "explica"}], "confianca": 1.0, "contraexemplos": [], "descricao": "Provar autoria e integridade com a chave privada, verificável pela pública — o carimbo inforjável que autoriza cada transação do goldchain.", "epistemico": "fato", "exemplos": [], "id": "assinatura_digital", "memoria": "semantica", "meta": {"dominio": "programacao", "fonte": ""}, "origem": "wiki", "palavras_chave": [], "sinonimos": [], "tipo": "conceito", "titulo": "Assinatura Digital"}
\section{Assinatura Digital}
Provar autoria e integridade com a chave privada, verificável pela pública — o carimbo inforjável que autoriza cada transação do goldchain.
\prov{relações: usa criptografia\_assimetrica; explica goldchain}
\prov{proveniência: wiki \textperiodcentered{} fato \textperiodcentered{} confiança 1.0 \textperiodcentered{} domínio programacao \textperiodcentered{} história 3 versões no jornal}

%CRISTAL {"arestas": [{"alvo": "mecanismo_de_atencao", "confianca": 1.0, "epistemico": "fato", "origem": "wiki", "rel": "explica"}, {"alvo": "operador_ψ_collapse", "confianca": 1.0, "epistemico": "fato", "origem": "wiki", "rel": "eh_um"}, {"alvo": "operadores_t_p_e_colapso_psi", "confianca": 1.0, "epistemico": "fato", "origem": "wiki", "rel": "usa"}, {"alvo": "colapso", "confianca": 1.0, "epistemico": "fato", "origem": "wiki", "rel": "eh_um"}, {"alvo": "corpo_mineral", "confianca": 1.0, "epistemico": "fato", "origem": "wiki", "rel": "usa"}], "confianca": 1.0, "contraexemplos": [], "descricao": "O mecanismo de atenção pontua a query contra as keys, faz softmax e soma os values ponderados — os pesos softmax são uma distribuição que DESABA no contexto relevante. É o operador Ψ do BAI, o colapso: a cada passo o modelo colapsa o passado numa leitura. A atenção cristalina é isso visto pela máquina mineral. Verif.: query [1,0] foca no key alinhado.", "epistemico": "inferencia", "exemplos": [], "id": "atencao_e_colapso", "memoria": "semantica", "meta": {"dominio": "ia", "foco": 0, "fonte": "IA e modelos de linguagem", "peso_max": 0.434}, "origem": "wiki", "palavras_chave": [], "sinonimos": [], "tipo": "conceito", "titulo": "Atenção = o Colapso (Ψ)"}
\section{Atenção = o Colapso (\ensuremath{\Psi})}
O mecanismo de atenção pontua a query contra as keys, faz softmax e soma os values ponderados — os pesos softmax são uma distribuição que DESABA no contexto relevante. É o operador \ensuremath{\Psi} do BAI, o colapso: a cada passo o modelo colapsa o passado numa leitura. A atenção cristalina é isso visto pela máquina mineral. Verif.: query [1,0] foca no key alinhado.
\prov{relações: explica mecanismo\_de\_atencao; eh\_um operador\_\ensuremath{\psi}\_collapse; usa operadores\_t\_p\_e\_colapso\_psi; eh\_um colapso; usa corpo\_mineral}
\prov{proveniência: wiki \textperiodcentered{} IA e modelos de linguagem \textperiodcentered{} inferencia \textperiodcentered{} confiança 1.0 \textperiodcentered{} domínio ia \textperiodcentered{} história 6 versões no jornal}

%CRISTAL {"arestas": [{"alvo": "seguranca_da_informacao", "confianca": 1.0, "epistemico": "fato", "origem": "wiki", "rel": "parte_de"}], "confianca": 1.0, "contraexemplos": [], "descricao": "Autenticar é provar QUEM você é; autorizar é decidir O QUE você pode. Dois passos distintos e sempre confundidos.", "epistemico": "fato", "exemplos": [], "id": "autenticacao_autorizacao", "memoria": "semantica", "meta": {"dominio": "programacao", "fonte": ""}, "origem": "wiki", "palavras_chave": [], "sinonimos": [], "tipo": "conceito", "titulo": "Autenticação e Autorização"}
\section{Autenticação e Autorização}
Autenticar é provar QUEM você é; autorizar é decidir O QUE você pode. Dois passos distintos e sempre confundidos.
\prov{relações: parte\_de seguranca\_da\_informacao}
\prov{proveniência: wiki \textperiodcentered{} fato \textperiodcentered{} confiança 1.0 \textperiodcentered{} domínio programacao \textperiodcentered{} história 2 versões no jornal}

%CRISTAL {"arestas": [{"alvo": "aprendizado_por_reforco", "confianca": 1.0, "epistemico": "fato", "origem": "wiki", "rel": "usa"}, {"alvo": "alphazero_2017", "confianca": 1.0, "epistemico": "fato", "origem": "wiki", "rel": "explica"}], "confianca": 1.0, "contraexemplos": [], "descricao": "A IA aprende jogando contra CÓPIAS de si mesma, gerando seu próprio currículo cada vez mais forte — sem dados humanos. O laço que levou TD-Gammon→AlphaGo Zero→AlphaZero a superar todo conhecimento acumulado.", "epistemico": "fato", "exemplos": [], "id": "auto_jogo", "memoria": "semantica", "meta": {"dominio": "ia", "fonte": "jogos e IA"}, "origem": "wiki", "palavras_chave": [], "sinonimos": [], "tipo": "conceito", "titulo": "Auto-jogo (Self-play)"}
\section{Auto-jogo (Self-play)}
A IA aprende jogando contra CÓPIAS de si mesma, gerando seu próprio currículo cada vez mais forte — sem dados humanos. O laço que levou TD-Gammon\ensuremath{\to}AlphaGo Zero\ensuremath{\to}AlphaZero a superar todo conhecimento acumulado.
\prov{relações: usa aprendizado\_por\_reforco; explica alphazero\_2017}
\prov{proveniência: wiki \textperiodcentered{} jogos e IA \textperiodcentered{} fato \textperiodcentered{} confiança 1.0 \textperiodcentered{} domínio ia \textperiodcentered{} história 3 versões no jornal}

%CRISTAL {"arestas": [{"alvo": "shell_scripting", "confianca": 1.0, "epistemico": "fato", "origem": "wiki", "rel": "relaciona"}, {"alvo": "linguagem_de_programacao", "confianca": 1.0, "epistemico": "fato", "origem": "wiki", "rel": "parte_de"}, {"alvo": "expressao_regular", "confianca": 1.0, "epistemico": "fato", "origem": "wiki", "rel": "usa"}], "confianca": 1.0, "contraexemplos": [], "descricao": "Aho, Weinberger, Kernighan: uma linguagem inteira para processar texto por linhas e campos, dirigida por padrões-ação. O canivete dos dados tabulares.", "epistemico": "fato", "exemplos": [], "id": "awk", "memoria": "semantica", "meta": {"autor": "Aho/Weinberger/Kernighan", "dominio": "programacao", "fonte": ""}, "origem": "wiki", "palavras_chave": [], "sinonimos": [], "tipo": "conceito", "titulo": "AWK"}
\section{AWK}
Aho, Weinberger, Kernighan: uma linguagem inteira para processar texto por linhas e campos, dirigida por padrões-ação. O canivete dos dados tabulares.
\prov{relações: relaciona shell\_scripting; parte\_de linguagem\_de\_programacao; usa expressao\_regular}
\prov{proveniência: wiki \textperiodcentered{} fato \textperiodcentered{} confiança 1.0 \textperiodcentered{} domínio programacao \textperiodcentered{} história 4 versões no jornal}

%CRISTAL {"arestas": [{"alvo": "isa_da_broca", "confianca": 1.0, "epistemico": "fato", "origem": "wiki", "rel": "usa"}, {"alvo": "cifra_metais_an", "confianca": 1.0, "epistemico": "fato", "origem": "wiki", "rel": "usa"}, {"alvo": "ula_por_componente", "confianca": 1.0, "epistemico": "fato", "origem": "wiki", "rel": "usa"}, {"alvo": "dsl_frontend_linguagem_gatos", "confianca": 1.0, "epistemico": "fato", "origem": "wiki", "rel": "parte_de"}, {"alvo": "expressividade_sub_turing", "confianca": 1.0, "epistemico": "fato", "origem": "wiki", "rel": "relaciona"}, {"alvo": "tiffany", "confianca": 1.0, "epistemico": "fato", "origem": "wiki", "rel": "parte_de"}], "confianca": 1.0, "contraexemplos": [], "descricao": "linguagem/backend_isa.py: o backend (paper §7) que compila o IR RAM (do DSL) para a ISA REAL do broca-so (ula/instrucoes.c, cpu.c). Escalares ↦ Word [valor,0]; registradores rN ↦ slots; constantes ↦ slots de dados; slot 0 = zero. Lowering fiel: LOAD empilha (B←A), STORE grava R, ADD/SUB fazem R=A±B, e as condições usam CMP contra o slot zero (FL_ZERO sse ambos zero → testa x==0) + JZ/JNZ (offset s8). O GATO tem DOIS caminhos: (i) por aritmética (n adições — 'zero opcodes novos'), (ii) o opcode NATIVO GOLD/SILVER/BRONZE = cifra_an(A,1/2/3) — verificado: [1,0]--SILVER×5-->total 70=Pell P6. Inclui um MONTADOR e um SIMULADOR Python fiéis ao cpu.c. CROSS-VALIDADO: a assembly gerada, montada pelo broca_asm C REAL, dá bytecode IDÊNTICO (61 bytes byte-a-byte) ao meu montador — o backend mira o metal de verdade. HONESTO: a indireção (loadi/storei) é REJEITADA (não é expressável na ISA de 12 = o limite sub-Turing); a Turing-completude fica no backend RAM/universal. Verificado: multiplicação=42, Pell(fechada)=2378, gato nativo=70. [F] compila, monta no metal real e roda.", "epistemico": "fato", "exemplos": [], "id": "backend_isa12_broca", "memoria": "semantica", "meta": {"autor": "Lopes/broca-so", "dominio": "computacao", "fonte": "linguagem/backend_isa.py (montar_isa, rodar_isa, _Lower); ula/instrucoes.c/cpu.c; cross-check: bytecode == ula/broca_asm (61 bytes)"}, "origem": "wiki", "palavras_chave": [], "sinonimos": [], "tipo": "conceito", "titulo": "O backend para a ISA de 12 do broca-so — bytecode idêntico ao montador C real [F]"}
\section{O backend para a ISA de 12 do broca-so — bytecode idêntico ao montador C real [F]}
linguagem/backend\_isa.py: o backend (paper §7) que compila o IR RAM (do DSL) para a ISA REAL do broca-so (ula/instrucoes.c, cpu.c). Escalares \ensuremath{\mapsto} Word [valor,0]; registradores rN \ensuremath{\mapsto} slots; constantes \ensuremath{\mapsto} slots de dados; slot 0 = zero. Lowering fiel: LOAD empilha (B\ensuremath{\leftarrow}A), STORE grava R, ADD/SUB fazem R=A$\pm$B, e as condições usam CMP contra o slot zero (FL\_ZERO sse ambos zero \ensuremath{\to} testa x==0) + JZ/JNZ (offset s8). O GATO tem DOIS caminhos: (i) por aritmética (n adições — 'zero opcodes novos'), (ii) o opcode NATIVO GOLD/SILVER/BRONZE = cifra\_an(A,1/2/3) — verificado: [1,0]--SILVER$\times$5-->total 70=Pell P6. Inclui um MONTADOR e um SIMULADOR Python fiéis ao cpu.c. CROSS-VALIDADO: a assembly gerada, montada pelo broca\_asm C REAL, dá bytecode IDÊNTICO (61 bytes byte-a-byte) ao meu montador — o backend mira o metal de verdade. HONESTO: a indireção (loadi/storei) é REJEITADA (não é expressável na ISA de 12 = o limite sub-Turing); a Turing-completude fica no backend RAM/universal. Verificado: multiplicação=42, Pell(fechada)=2378, gato nativo=70. [F] compila, monta no metal real e roda.
\prov{relações: usa isa\_da\_broca; usa cifra\_metais\_an; usa ula\_por\_componente; parte\_de dsl\_frontend\_linguagem\_gatos; relaciona expressividade\_sub\_turing; parte\_de tiffany}
\prov{proveniência: wiki \textperiodcentered{} linguagem/backend\_isa.py (montar\_isa, rodar\_isa, \_Lower); ula/instrucoes.c/cpu.c; cross-check: bytecode == ula/broca\_asm (61 bytes) \textperiodcentered{} fato \textperiodcentered{} confiança 1.0 \textperiodcentered{} domínio computacao \textperiodcentered{} história 13 versões no jornal}

%CRISTAL {"arestas": [{"alvo": "ptx_isa_virtual", "confianca": 1.0, "epistemico": "fato", "origem": "wiki", "rel": "usa"}, {"alvo": "sass_gpu_assembly", "confianca": 1.0, "epistemico": "fato", "origem": "wiki", "rel": "usa"}, {"alvo": "pipeline_compilacao_gpu", "confianca": 1.0, "epistemico": "fato", "origem": "wiki", "rel": "parte_de"}, {"alvo": "arch_gpu_turing", "confianca": 1.0, "epistemico": "fato", "origem": "wiki", "rel": "usa"}, {"alvo": "dicionario_maquina_universal_sass", "confianca": 1.0, "epistemico": "fato", "origem": "wiki", "rel": "usa"}, {"alvo": "compilador_glsl_broca", "confianca": 1.0, "epistemico": "fato", "origem": "wiki", "rel": "parte_de"}], "confianca": 1.0, "contraexemplos": [], "descricao": "FEITO e verificado: linguagem/glsl_ptx.py emite PTX e RODA na GPU (sm_75), sem numpy no cálculo. Sem CUDA toolkit (ptxas), usa o DRIVER (libcuda.so via ctypes): cuModuleLoadDataEx faz o JIT PTX→SASS para o device atual (sm_75), cuLaunchKernel lança. O kernel: out[i]=func(a·x[i]+b) — a FFMA (fma.rn.f32 = o gato/La Hire, o afim) + a MUFU (sin.approx.f32/cos = o gambá, a transcendente); o warp (32 threads) = as lanes. Verificado (4/4): sin e cos rodam em sm_75, o resultado bate com numpy a ~6e-7 (precisão de float). O Python só orquestra (aloca memória via torch, lança); a MATEMÁTICA é SASS no metal — não numpy. Armadilha resolvida: o ptxas rejeita comentários não-ASCII (o traço — e o á) no PTX. É o 1º degrau real da cadeia DSL→PTX→SASS→GPU. E o KERNEL FUNDIDO prova o ganho: um kernel PTX com LOOP no thread (out=sum_k sin(a·x+b+k·step), a estrutura do march de shells) roda a ~0,17 ms/frame (5853 fps) @ 840×473, K=64 — 31× o torch NÃO-fundido (5,35 ms, 187 fps, que lança 64 kernels). O SASS fundido faz os K passos numa passada; o torch lança K kernels com K tensores intermediários. É POR ISSO que a GPU nativa faz 60 fps e o torch não: não é a GPU, é a camada Python/torch lançando kernels não-fundidos. glsl_ptx.py (verificar_ptx 5/5, benchmark_ptx_vs_torch).", "epistemico": "fato", "exemplos": [], "id": "backend_ptx_sm75_funciona", "memoria": "semantica", "meta": {"dominio": "sass_gpu", "fonte": "SASS/GPU (o metal)"}, "origem": "wiki", "palavras_chave": [], "sinonimos": [], "tipo": "conceito", "titulo": "O backend PTX do glsl_compilador RODA em sm_75 (a GTX 1650) — sem numpy no cálculo"}
\section{O backend PTX do glsl\_compilador RODA em sm\_75 (a GTX 1650) — sem numpy no cálculo}
FEITO e verificado: linguagem/glsl\_ptx.py emite PTX e RODA na GPU (sm\_75), sem numpy no cálculo. Sem CUDA toolkit (ptxas), usa o DRIVER (libcuda.so via ctypes): cuModuleLoadDataEx faz o JIT PTX\ensuremath{\to}SASS para o device atual (sm\_75), cuLaunchKernel lança. O kernel: out[i]=func(a\textperiodcentered x[i]+b) — a FFMA (fma.rn.f32 = o gato/La Hire, o afim) + a MUFU (sin.approx.f32/cos = o gambá, a transcendente); o warp (32 threads) = as lanes. Verificado (4/4): sin e cos rodam em sm\_75, o resultado bate com numpy a \~{}6e-7 (precisão de float). O Python só orquestra (aloca memória via torch, lança); a MATEMÁTICA é SASS no metal — não numpy. Armadilha resolvida: o ptxas rejeita comentários não-ASCII (o traço — e o á) no PTX. É o 1º degrau real da cadeia DSL\ensuremath{\to}PTX\ensuremath{\to}SASS\ensuremath{\to}GPU. E o KERNEL FUNDIDO prova o ganho: um kernel PTX com LOOP no thread (out=sum\_k sin(a\textperiodcentered x+b+k\textperiodcentered step), a estrutura do march de shells) roda a \~{}0,17 ms/frame (5853 fps) @ 840$\times$473, K=64 — 31$\times$ o torch NÃO-fundido (5,35 ms, 187 fps, que lança 64 kernels). O SASS fundido faz os K passos numa passada; o torch lança K kernels com K tensores intermediários. É POR ISSO que a GPU nativa faz 60 fps e o torch não: não é a GPU, é a camada Python/torch lançando kernels não-fundidos. glsl\_ptx.py (verificar\_ptx 5/5, benchmark\_ptx\_vs\_torch).
\prov{relações: usa ptx\_isa\_virtual; usa sass\_gpu\_assembly; parte\_de pipeline\_compilacao\_gpu; usa arch\_gpu\_turing; usa dicionario\_maquina\_universal\_sass; parte\_de compilador\_glsl\_broca}
\prov{proveniência: wiki \textperiodcentered{} SASS/GPU (o metal) \textperiodcentered{} fato \textperiodcentered{} confiança 1.0 \textperiodcentered{} domínio sass\_gpu \textperiodcentered{} história 42 versões no jornal}

%CRISTAL {"arestas": [{"alvo": "bai_orbita_reta_metalica", "confianca": 1.0, "epistemico": "fato", "origem": "wiki", "rel": "especializa"}, {"alvo": "meta_inducao_geral_simbolo", "confianca": 1.0, "epistemico": "fato", "origem": "wiki", "rel": "usa"}, {"alvo": "traducao_como_transformacao", "confianca": 1.0, "epistemico": "fato", "origem": "wiki", "rel": "explica"}], "confianca": 1.0, "contraexemplos": [], "descricao": "No teorema, o BAI não é mais um passo do cálculo: é a BÚSSOLA. A órbita 𝒞_δ de cada trabalho (o raio σ_n, δ=σ_n·4) APONTA a simbologia-alvo e o registro — em vez de vasculhar qual transformação aplicar, escolhe-se a órbita e a substituição de Peirce a realiza. O que era seleção de reta metálica no goldchain vira, na fala, seleção do alvo da tradução. Aponta, não procura.", "epistemico": "inferencia", "exemplos": [], "id": "bai_como_bussola", "memoria": "semantica", "meta": {"dominio": "programacao", "fonte": "conversa/bai.py; bai_orbita_reta_metalica"}, "origem": "wiki", "palavras_chave": [], "sinonimos": [], "tipo": "conceito", "titulo": "O BAI como Bússola da Tradução"}
\section{O BAI como Bússola da Tradução}
No teorema, o BAI não é mais um passo do cálculo: é a BÚSSOLA. A órbita \ensuremath{\mathcal{C}}\_\ensuremath{\delta} de cada trabalho (o raio \ensuremath{\sigma}\_n, \ensuremath{\delta}=\ensuremath{\sigma}\_n\textperiodcentered 4) APONTA a simbologia-alvo e o registro — em vez de vasculhar qual transformação aplicar, escolhe-se a órbita e a substituição de Peirce a realiza. O que era seleção de reta metálica no goldchain vira, na fala, seleção do alvo da tradução. Aponta, não procura.
\prov{relações: especializa bai\_orbita\_reta\_metalica; usa meta\_inducao\_geral\_simbolo; explica traducao\_como\_transformacao}
\prov{proveniência: wiki \textperiodcentered{} conversa/bai.py; bai\_orbita\_reta\_metalica \textperiodcentered{} inferencia \textperiodcentered{} confiança 1.0 \textperiodcentered{} domínio programacao \textperiodcentered{} história 4 versões no jornal}

%CRISTAL {"arestas": [{"alvo": "bai_tese_mae", "confianca": 1.0, "epistemico": "fato", "origem": "wiki", "rel": "parte_de"}, {"alvo": "delta", "confianca": 1.0, "epistemico": "fato", "origem": "wiki", "rel": "explica"}, {"alvo": "associador_calor", "confianca": 1.0, "epistemico": "fato", "origem": "wiki", "rel": "relaciona"}], "confianca": 1.0, "contraexemplos": [], "descricao": "A primitiva δ: o orçamento de profundidade/vigência da propagação. Lema-chave: δ é uma energia bem-fundada, logo os ciclos no grafo de delegação são inofensivos (a propagação sempre termina). O δ decresce a cada salto — a autoridade se esgota.", "epistemico": "fato", "exemplos": [], "id": "bai_delta", "memoria": "semantica", "meta": {"autor": "broca-so", "dominio": "operacao_so", "fonte": "papers/casl-propagation.tex:def:delta-budget"}, "origem": "wiki", "palavras_chave": [], "sinonimos": [], "tipo": "conceito", "titulo": "δ — Orçamento (energia bem-fundada → ciclos inofensivos)"}
\section{\ensuremath{\delta} — Orçamento (energia bem-fundada \ensuremath{\to} ciclos inofensivos)}
A primitiva \ensuremath{\delta}: o orçamento de profundidade/vigência da propagação. Lema-chave: \ensuremath{\delta} é uma energia bem-fundada, logo os ciclos no grafo de delegação são inofensivos (a propagação sempre termina). O \ensuremath{\delta} decresce a cada salto — a autoridade se esgota.
\prov{relações: parte\_de bai\_tese\_mae; explica delta; relaciona associador\_calor}
\prov{proveniência: wiki \textperiodcentered{} papers/casl-propagation.tex:def:delta-budget \textperiodcentered{} fato \textperiodcentered{} confiança 1.0 \textperiodcentered{} domínio operacao\_so \textperiodcentered{} história 4 versões no jornal}

%CRISTAL {"arestas": [{"alvo": "bai_tese_mae", "confianca": 1.0, "epistemico": "fato", "origem": "wiki", "rel": "usa"}, {"alvo": "face_turno_fail_closed", "confianca": 1.0, "epistemico": "fato", "origem": "wiki", "rel": "relaciona"}], "confianca": 1.0, "contraexemplos": [], "descricao": "Proposição CASL: a autorização efetiva é a união dos grants menos as proibições — uma proibição domina qualquer grant sobre a mesma linha. O sistema é fail-closed por definição: sem grant vigente, nega. É a segurança do 'não fabricar' aplicada à autorização.", "epistemico": "fato", "exemplos": [], "id": "bai_deny_overrides", "memoria": "semantica", "meta": {"autor": "broca-so", "dominio": "operacao_so", "fonte": "papers/casl-propagation.tex (deny-overrides)"}, "origem": "wiki", "palavras_chave": [], "sinonimos": [], "tipo": "conceito", "titulo": "Deny-Overrides (a proibição sempre vence)"}
\section{Deny-Overrides (a proibição sempre vence)}
Proposição CASL: a autorização efetiva é a união dos grants menos as proibições — uma proibição domina qualquer grant sobre a mesma linha. O sistema é fail-closed por definição: sem grant vigente, nega. É a segurança do 'não fabricar' aplicada à autorização.
\prov{relações: usa bai\_tese\_mae; relaciona face\_turno\_fail\_closed}
\prov{proveniência: wiki \textperiodcentered{} papers/casl-propagation.tex (deny-overrides) \textperiodcentered{} fato \textperiodcentered{} confiança 1.0 \textperiodcentered{} domínio operacao\_so \textperiodcentered{} história 3 versões no jornal}

%CRISTAL {"arestas": [{"alvo": "bai_tese_mae", "confianca": 1.0, "epistemico": "fato", "origem": "wiki", "rel": "parte_de"}, {"alvo": "kappa", "confianca": 1.0, "epistemico": "fato", "origem": "wiki", "rel": "explica"}, {"alvo": "capacidade_abstrata", "confianca": 1.0, "epistemico": "fato", "origem": "wiki", "rel": "explica"}], "confianca": 1.0, "contraexemplos": [], "descricao": "A primitiva κ do quinteto: a capacidade (σ, c, δ, ω, o) — a unidade que se propaga na delegação. Em CASL é a Capacidade sobre linhas com condições; é o 'grant' que viaja pelo grafo de delegação.", "epistemico": "fato", "exemplos": [], "id": "bai_kappa", "memoria": "semantica", "meta": {"autor": "broca-so", "dominio": "operacao_so", "fonte": "papers/casl-propagation.tex:def:kappa-abstract"}, "origem": "wiki", "palavras_chave": [], "sinonimos": [], "tipo": "conceito", "titulo": "κ — Capacidade (a unidade propagada)"}
\section{\ensuremath{\kappa} — Capacidade (a unidade propagada)}
A primitiva \ensuremath{\kappa} do quinteto: a capacidade (\ensuremath{\sigma}, c, \ensuremath{\delta}, \ensuremath{\omega}, o) — a unidade que se propaga na delegação. Em CASL é a Capacidade sobre linhas com condições; é o 'grant' que viaja pelo grafo de delegação.
\prov{relações: parte\_de bai\_tese\_mae; explica kappa; explica capacidade\_abstrata}
\prov{proveniência: wiki \textperiodcentered{} papers/casl-propagation.tex:def:kappa-abstract \textperiodcentered{} fato \textperiodcentered{} confiança 1.0 \textperiodcentered{} domínio operacao\_so \textperiodcentered{} história 4 versões no jornal}

%CRISTAL {"arestas": [{"alvo": "bai_tese_mae", "confianca": 1.0, "epistemico": "fato", "origem": "wiki", "rel": "parte_de"}], "confianca": 1.0, "contraexemplos": [], "descricao": "A primitiva Φ: a medida de organização estrutural do conjunto de capacidades após a propagação. Par fundacional (Φ, δ): os teoremas de segurança (não-escalada, propagação limitada, fail-closed) são independentes da escolha de (Φ, δ).", "epistemico": "fato", "exemplos": [], "id": "bai_phi", "memoria": "semantica", "meta": {"autor": "broca-so", "dominio": "operacao_so", "fonte": "papers/casl-propagation.tex:def:phi-abstract"}, "origem": "wiki", "palavras_chave": [], "sinonimos": [], "tipo": "conceito", "titulo": "Φ — Organização Estrutural pós-propagação"}
\section{\ensuremath{\Phi} — Organização Estrutural pós-propagação}
A primitiva \ensuremath{\Phi}: a medida de organização estrutural do conjunto de capacidades após a propagação. Par fundacional (\ensuremath{\Phi}, \ensuremath{\delta}): os teoremas de segurança (não-escalada, propagação limitada, fail-closed) são independentes da escolha de (\ensuremath{\Phi}, \ensuremath{\delta}).
\prov{relações: parte\_de bai\_tese\_mae}
\prov{proveniência: wiki \textperiodcentered{} papers/casl-propagation.tex:def:phi-abstract \textperiodcentered{} fato \textperiodcentered{} confiança 1.0 \textperiodcentered{} domínio operacao\_so \textperiodcentered{} história 2 versões no jornal}

%CRISTAL {"arestas": [{"alvo": "bai_tese_mae", "confianca": 1.0, "epistemico": "fato", "origem": "wiki", "rel": "parte_de"}, {"alvo": "colapso", "confianca": 1.0, "epistemico": "fato", "origem": "wiki", "rel": "explica"}, {"alvo": "colapso_multicamada_peso", "confianca": 1.0, "epistemico": "fato", "origem": "wiki", "rel": "generaliza"}], "confianca": 1.0, "contraexemplos": [], "descricao": "A primitiva Ψ: o colapso = argmax sobre os candidatos vigentes 𝒞_δ (quando o máximo existe; senão Ψ(q)=⊥, fail-closed). O Ψ do Broca (seleção de face no tecido Koch, colapso_ping_pong_overlap) é uma ESPECIALIZAÇÃO do colapso abstrato do BAI.", "epistemico": "fato", "exemplos": [], "id": "bai_psi_colapso", "memoria": "semantica", "meta": {"autor": "broca-so", "dominio": "operacao_so", "fonte": "papers/casl-propagation.tex:def:collapse; conversa/colapso.py"}, "origem": "wiki", "palavras_chave": [], "sinonimos": [], "tipo": "conceito", "titulo": "Ψ — Colapso (argmax sobre os vigentes, fail-closed)"}
\section{\ensuremath{\Psi} — Colapso (argmax sobre os vigentes, fail-closed)}
A primitiva \ensuremath{\Psi}: o colapso = argmax sobre os candidatos vigentes \ensuremath{\mathcal{C}}\_\ensuremath{\delta} (quando o máximo existe; senão \ensuremath{\Psi}(q)=\ensuremath{\perp}, fail-closed). O \ensuremath{\Psi} do Broca (seleção de face no tecido Koch, colapso\_ping\_pong\_overlap) é uma ESPECIALIZAÇÃO do colapso abstrato do BAI.
\prov{relações: parte\_de bai\_tese\_mae; explica colapso; generaliza colapso\_multicamada\_peso}
\prov{proveniência: wiki \textperiodcentered{} papers/casl-propagation.tex:def:collapse; conversa/colapso.py \textperiodcentered{} fato \textperiodcentered{} confiança 1.0 \textperiodcentered{} domínio operacao\_so \textperiodcentered{} história 4 versões no jornal}

%CRISTAL {"arestas": [{"alvo": "bai_tese_mae", "confianca": 1.0, "epistemico": "fato", "origem": "wiki", "rel": "parte_de"}, {"alvo": "floco_calor_dissipado", "confianca": 1.0, "epistemico": "fato", "origem": "wiki", "rel": "relaciona"}], "confianca": 1.0, "contraexemplos": [], "descricao": "A primitiva T: o operador de atenuação deflacionária — uma concessão válida nunca aumenta a autoridade (T ⪯ id). Junto com a propagação P, define a dinâmica das capacidades detidas Hold(p) como o menor conjunto fechado.", "epistemico": "fato", "exemplos": [], "id": "bai_t_atenuacao", "memoria": "semantica", "meta": {"autor": "broca-so", "dominio": "operacao_so", "fonte": "papers/casl-propagation.tex:def:TP"}, "origem": "wiki", "palavras_chave": [], "sinonimos": [], "tipo": "conceito", "titulo": "T — Atenuação Deflacionária"}
\section{T — Atenuação Deflacionária}
A primitiva T: o operador de atenuação deflacionária — uma concessão válida nunca aumenta a autoridade (T \ensuremath{\preceq} id). Junto com a propagação P, define a dinâmica das capacidades detidas Hold(p) como o menor conjunto fechado.
\prov{relações: parte\_de bai\_tese\_mae; relaciona floco\_calor\_dissipado}
\prov{proveniência: wiki \textperiodcentered{} papers/casl-propagation.tex:def:TP \textperiodcentered{} fato \textperiodcentered{} confiança 1.0 \textperiodcentered{} domínio operacao\_so \textperiodcentered{} história 3 versões no jornal}

%CRISTAL {"arestas": [{"alvo": "bai", "confianca": 1.0, "epistemico": "fato", "origem": "wiki", "rel": "explica"}, {"alvo": "reticulados", "confianca": 1.0, "epistemico": "fato", "origem": "wiki", "rel": "usa"}, {"alvo": "operacao_broca_so", "confianca": 1.0, "epistemico": "fato", "origem": "wiki", "rel": "parte_de"}, {"alvo": "broca_os", "confianca": 1.0, "epistemico": "fato", "origem": "wiki", "rel": "generaliza"}, {"alvo": "tiffany", "confianca": 1.0, "epistemico": "fato", "origem": "wiki", "rel": "generaliza"}], "confianca": 1.0, "contraexemplos": [], "descricao": "Biblioteca de Autorização Isomórfica: delegação capacitária monótona sobre reticulados booleanos (join=união de grants fail-closed; meet=intersecção; revogação só diminui autoridade; identificação semântica). Tese-mãe: toda instância é determinada pelo quinteto (κ,δ,Φ,T,Ψ) — CASL, Broca OS e Tiffany são INTERPRETAÇÕES do mesmo objeto, não sistemas concorrentes.", "epistemico": "fato", "exemplos": [], "id": "bai_tese_mae", "memoria": "semantica", "meta": {"autor": "broca-so", "dominio": "operacao_so", "fonte": "papers/bai_quinteto_hub.md; papers/casl-propagation.tex"}, "origem": "wiki", "palavras_chave": [], "sinonimos": [], "tipo": "conceito", "titulo": "BAI — a Tese-Mãe do Quinteto"}
\section{BAI — a Tese-Mãe do Quinteto}
Biblioteca de Autorização Isomórfica: delegação capacitária monótona sobre reticulados booleanos (join=união de grants fail-closed; meet=intersecção; revogação só diminui autoridade; identificação semântica). Tese-mãe: toda instância é determinada pelo quinteto (\ensuremath{\kappa},\ensuremath{\delta},\ensuremath{\Phi},T,\ensuremath{\Psi}) — CASL, Broca OS e Tiffany são INTERPRETAÇÕES do mesmo objeto, não sistemas concorrentes.
\prov{relações: explica bai; usa reticulados; parte\_de operacao\_broca\_so; generaliza broca\_os; generaliza tiffany}
\prov{proveniência: wiki \textperiodcentered{} papers/bai\_quinteto\_hub.md; papers/casl-propagation.tex \textperiodcentered{} fato \textperiodcentered{} confiança 1.0 \textperiodcentered{} domínio operacao\_so \textperiodcentered{} história 6 versões no jornal}

%CRISTAL {"arestas": [{"alvo": "bai_tese_mae", "confianca": 1.0, "epistemico": "fato", "origem": "wiki", "rel": "usa"}, {"alvo": "monotono", "confianca": 1.0, "epistemico": "fato", "origem": "wiki", "rel": "usa"}, {"alvo": "reticulados", "confianca": 1.0, "epistemico": "fato", "origem": "wiki", "rel": "usa"}], "confianca": 1.0, "contraexemplos": [], "descricao": "Teorema: toda biblioteca de autorização monótona sobre um reticulado é uma BAI — colapso por argmax dos vigentes com fail-closed em ⊥, tendo propagação limitada e deny-overrides como corolários da monotonicidade. RBAC/ACL (roles estáticos) NÃO são BAI; a BAI é delegação capacitária monótona.", "epistemico": "fato", "exemplos": [], "id": "bai_universalidade", "memoria": "semantica", "meta": {"autor": "broca-so", "dominio": "operacao_so", "fonte": "papers/casl-propagation.tex:thm:universalidade"}, "origem": "wiki", "palavras_chave": [], "sinonimos": [], "tipo": "conceito", "titulo": "Universalidade da BAI (teorema)"}
\section{Universalidade da BAI (teorema)}
Teorema: toda biblioteca de autorização monótona sobre um reticulado é uma BAI — colapso por argmax dos vigentes com fail-closed em \ensuremath{\perp}, tendo propagação limitada e deny-overrides como corolários da monotonicidade. RBAC/ACL (roles estáticos) NÃO são BAI; a BAI é delegação capacitária monótona.
\prov{relações: usa bai\_tese\_mae; usa monotono; usa reticulados}
\prov{proveniência: wiki \textperiodcentered{} papers/casl-propagation.tex:thm:universalidade \textperiodcentered{} fato \textperiodcentered{} confiança 1.0 \textperiodcentered{} domínio operacao\_so \textperiodcentered{} história 4 versões no jornal}

%CRISTAL {"arestas": [{"alvo": "programacao", "confianca": 1.0, "epistemico": "fato", "origem": "wiki", "rel": "parte_de"}], "confianca": 1.0, "contraexemplos": [], "descricao": "Um sistema para guardar, consultar e proteger dados em escala e de forma durável — a memória persistente do software.", "epistemico": "fato", "exemplos": [], "id": "banco_de_dados", "memoria": "semantica", "meta": {"dominio": "programacao", "fonte": ""}, "origem": "wiki", "palavras_chave": [], "sinonimos": [], "tipo": "conceito", "titulo": "Banco de Dados"}
\section{Banco de Dados}
Um sistema para guardar, consultar e proteger dados em escala e de forma durável — a memória persistente do software.
\prov{relações: parte\_de programacao}
\prov{proveniência: wiki \textperiodcentered{} fato \textperiodcentered{} confiança 1.0 \textperiodcentered{} domínio programacao \textperiodcentered{} história 4 versões no jornal}

%CRISTAL {"arestas": [{"alvo": "nosql", "confianca": 1.0, "epistemico": "fato", "origem": "wiki", "rel": "especializa"}, {"alvo": "tiffany", "confianca": 1.0, "epistemico": "fato", "origem": "wiki", "rel": "explica"}], "confianca": 1.0, "contraexemplos": [], "descricao": "Guardar nós e arestas como cidadãos de primeira classe — consultas de caminho e vizinhança nativas. O modelo da própria Universidade Tiffany.", "epistemico": "fato", "exemplos": [], "id": "banco_de_dados_grafo", "memoria": "semantica", "meta": {"dominio": "programacao", "fonte": ""}, "origem": "wiki", "palavras_chave": [], "sinonimos": [], "tipo": "conceito", "titulo": "Banco de Dados de Grafo"}
\section{Banco de Dados de Grafo}
Guardar nós e arestas como cidadãos de primeira classe — consultas de caminho e vizinhança nativas. O modelo da própria Universidade Tiffany.
\prov{relações: especializa nosql; explica tiffany}
\prov{proveniência: wiki \textperiodcentered{} fato \textperiodcentered{} confiança 1.0 \textperiodcentered{} domínio programacao \textperiodcentered{} história 6 versões no jornal}

%CRISTAL {"arestas": [{"alvo": "shell", "confianca": 1.0, "epistemico": "fato", "origem": "wiki", "rel": "especializa"}, {"alvo": "shell_scripting", "confianca": 1.0, "epistemico": "fato", "origem": "wiki", "rel": "usa"}, {"alvo": "linux", "confianca": 1.0, "epistemico": "fato", "origem": "wiki", "rel": "relaciona"}, {"alvo": "sh_posix", "confianca": 1.0, "epistemico": "fato", "origem": "wiki", "rel": "especializa"}], "confianca": 1.0, "contraexemplos": [], "descricao": "Bourne-Again Shell (Fox/Ramey): o shell padrão do Linux — scripting, pipes, expansões, controle de fluxo. A cola do sistema.", "epistemico": "fato", "exemplos": [], "id": "bash", "memoria": "semantica", "meta": {"autor": "Fox/Ramey", "dominio": "programacao", "fonte": ""}, "origem": "wiki", "palavras_chave": [], "sinonimos": [], "tipo": "conceito", "titulo": "Bash"}
\section{Bash}
Bourne-Again Shell (Fox/Ramey): o shell padrão do Linux — scripting, pipes, expansões, controle de fluxo. A cola do sistema.
\prov{relações: especializa shell; usa shell\_scripting; relaciona linux; especializa sh\_posix}
\prov{proveniência: wiki \textperiodcentered{} fato \textperiodcentered{} confiança 1.0 \textperiodcentered{} domínio programacao \textperiodcentered{} história 5 versões no jornal}

%CRISTAL {"arestas": [{"alvo": "componente_ui", "confianca": 1.0, "epistemico": "fato", "origem": "wiki", "rel": "usa"}, {"alvo": "design_system", "confianca": 1.0, "epistemico": "fato", "origem": "wiki", "rel": "usa"}], "confianca": 1.0, "contraexemplos": [], "descricao": "Componentes de UI prontos e consistentes (Material UI, Bootstrap) — não reinventar o botão; acelerar com um design coerente.", "epistemico": "fato", "exemplos": [], "id": "biblioteca_de_componentes", "memoria": "semantica", "meta": {"dominio": "programacao", "fonte": ""}, "origem": "wiki", "palavras_chave": [], "sinonimos": [], "tipo": "conceito", "titulo": "Biblioteca de Componentes (MUI)"}
\section{Biblioteca de Componentes (MUI)}
Componentes de UI prontos e consistentes (Material UI, Bootstrap) — não reinventar o botão; acelerar com um design coerente.
\prov{relações: usa componente\_ui; usa design\_system}
\prov{proveniência: wiki \textperiodcentered{} fato \textperiodcentered{} confiança 1.0 \textperiodcentered{} domínio programacao \textperiodcentered{} história 6 versões no jornal}

%CRISTAL {"arestas": [{"alvo": "mecanica_reta_mineral", "confianca": 1.0, "epistemico": "fato", "origem": "wiki", "rel": "usa"}, {"alvo": "algebra_reta_mineral", "confianca": 1.0, "epistemico": "fato", "origem": "wiki", "rel": "usa"}, {"alvo": "geometria_do_caos", "confianca": 1.0, "epistemico": "fato", "origem": "wiki", "rel": "usa"}, {"alvo": "particulas_minerais", "confianca": 1.0, "epistemico": "fato", "origem": "wiki", "rel": "usa"}, {"alvo": "transformada_mineral_7d", "confianca": 1.0, "epistemico": "fato", "origem": "wiki", "rel": "usa"}, {"alvo": "laco_traducao_bai_peirce", "confianca": 1.0, "epistemico": "fato", "origem": "wiki", "rel": "relaciona"}, {"alvo": "minimax_e_colapso", "confianca": 1.0, "epistemico": "fato", "origem": "wiki", "rel": "relaciona"}, {"alvo": "computador_mineral", "confianca": 1.0, "epistemico": "fato", "origem": "wiki", "rel": "parte_de"}], "confianca": 1.0, "contraexemplos": [], "descricao": "A stdlib do computador mineral: a mecânica (evolucao=Aₙ, espectro de Lyapunov, ação conservada, integrável), a geometria (endireita=espiral→reta, subdominante (λ,θ), curvatura N−1), a transformada (codifica/decodifica octoniônica reversível) e a física (átomo, entropia log σ_n) — todas como funções no metal, já fundamentadas. E a AUTORIZAÇÃO é o BAI (o 'SO'): a mesma reta que computa também autoriza — colapso Ψ, orçamento δ=σ_n·4, deny-overrides. Verificado em biblioteca_mineral.py.", "epistemico": "inferencia", "exemplos": [], "id": "biblioteca_mineral", "memoria": "semantica", "meta": {"dominio": "computacao", "fonte": "linguagem mineral"}, "origem": "wiki", "palavras_chave": [], "sinonimos": [], "tipo": "conceito", "titulo": "A Biblioteca Mineral (funções no metal) + o BAI como SO"}
\section{A Biblioteca Mineral (funções no metal) + o BAI como SO}
A stdlib do computador mineral: a mecânica (evolucao=A\ensuremath{_{n}}, espectro de Lyapunov, ação conservada, integrável), a geometria (endireita=espiral\ensuremath{\to}reta, subdominante (\ensuremath{\lambda},\ensuremath{\theta}), curvatura N\ensuremath{-}1), a transformada (codifica/decodifica octoniônica reversível) e a física (átomo, entropia log \ensuremath{\sigma}\_n) — todas como funções no metal, já fundamentadas. E a AUTORIZAÇÃO é o BAI (o 'SO'): a mesma reta que computa também autoriza — colapso \ensuremath{\Psi}, orçamento \ensuremath{\delta}=\ensuremath{\sigma}\_n\textperiodcentered 4, deny-overrides. Verificado em biblioteca\_mineral.py.
\prov{relações: usa mecanica\_reta\_mineral; usa algebra\_reta\_mineral; usa geometria\_do\_caos; usa particulas\_minerais; usa transformada\_mineral\_7d; relaciona laco\_traducao\_bai\_peirce; relaciona minimax\_e\_colapso; parte\_de computador\_mineral}
\prov{proveniência: wiki \textperiodcentered{} linguagem mineral \textperiodcentered{} inferencia \textperiodcentered{} confiança 1.0 \textperiodcentered{} domínio computacao \textperiodcentered{} história 9 versões no jornal}

%CRISTAL {"arestas": [{"alvo": "joalheria_hub", "confianca": 1.0, "epistemico": "fato", "origem": "wiki", "rel": "relaciona"}, {"alvo": "metais", "confianca": 1.0, "epistemico": "fato", "origem": "wiki", "rel": "usa"}, {"alvo": "mineracao_pipe_hook", "confianca": 1.0, "epistemico": "fato", "origem": "wiki", "rel": "relaciona"}, {"alvo": "corpo_mineral", "confianca": 1.0, "epistemico": "fato", "origem": "wiki", "rel": "relaciona"}], "confianca": 1.0, "contraexemplos": [], "descricao": "O render da pedra virou biblioteca: a física de uma superfície nas suas OITO faces — geometria, mecânica, óptica, termodinâmica, eletromagnetismo, quântica, gravitação, matéria —, toda compilada do GLSL para o metal (glsl_compilador) e costurada aos METAIS. Não são oito teorias coladas: são oito leituras de uma coisa. Três fios amarram tudo — a parábola cristal/borda/caos (o eixo comum), a zeta dos metais ∑σ_n^{−β} (= o calor = a partição = o traço = as geodésicas-primas), e c (luz=EM=onda gravitacional, a contagem de voltas no círculo de Gentil). linguagem/biblioteca.py (3/3). A física de uma superfície É a física dos metais. As dimensões: 3D (render) e 7D (octônion/mineração) encontram-se na reta mineral (1D).", "epistemico": "fato", "exemplos": [], "id": "biblioteca_pipe", "memoria": "semantica", "meta": {"dominio": "biblioteca_pipe", "fonte": "biblioteca_pipe"}, "origem": "wiki", "palavras_chave": [], "sinonimos": [], "tipo": "conceito", "titulo": "A biblioteca do pipe (as oito físicas de uma superfície)"}
\section{A biblioteca do pipe (as oito físicas de uma superfície)}
O render da pedra virou biblioteca: a física de uma superfície nas suas OITO faces — geometria, mecânica, óptica, termodinâmica, eletromagnetismo, quântica, gravitação, matéria —, toda compilada do GLSL para o metal (glsl\_compilador) e costurada aos METAIS. Não são oito teorias coladas: são oito leituras de uma coisa. Três fios amarram tudo — a parábola cristal/borda/caos (o eixo comum), a zeta dos metais \ensuremath{\sum}\ensuremath{\sigma}\_n\^{}\{\ensuremath{-}\ensuremath{\beta}\} (= o calor = a partição = o traço = as geodésicas-primas), e c (luz=EM=onda gravitacional, a contagem de voltas no círculo de Gentil). linguagem/biblioteca.py (3/3). A física de uma superfície É a física dos metais. As dimensões: 3D (render) e 7D (octônion/mineração) encontram-se na reta mineral (1D).
\prov{relações: relaciona joalheria\_hub; usa metais; relaciona mineracao\_pipe\_hook; relaciona corpo\_mineral}
\prov{proveniência: wiki \textperiodcentered{} biblioteca\_pipe \textperiodcentered{} fato \textperiodcentered{} confiança 1.0 \textperiodcentered{} domínio biblioteca\_pipe \textperiodcentered{} história 5 versões no jornal}

%CRISTAL {"arestas": [{"alvo": "latex", "confianca": 1.0, "epistemico": "fato", "origem": "wiki", "rel": "parte_de"}], "confianca": 1.0, "contraexemplos": [], "descricao": "O formato de referências bibliográficas do (La)TeX — a citação como dado estruturado, separada do texto que a usa.", "epistemico": "fato", "exemplos": [], "id": "bibtex", "memoria": "semantica", "meta": {"dominio": "programacao", "fonte": ""}, "origem": "wiki", "palavras_chave": [], "sinonimos": [], "tipo": "conceito", "titulo": "BibTeX"}
\section{BibTeX}
O formato de referências bibliográficas do (La)TeX — a citação como dado estruturado, separada do texto que a usa.
\prov{relações: parte\_de latex}
\prov{proveniência: wiki \textperiodcentered{} fato \textperiodcentered{} confiança 1.0 \textperiodcentered{} domínio programacao \textperiodcentered{} história 2 versões no jornal}

%CRISTAL {"arestas": [{"alvo": "circuito_logico", "confianca": 1.0, "epistemico": "fato", "origem": "curriculo", "rel": "usa"}, {"alvo": "inteiros", "confianca": 1.0, "epistemico": "fato", "origem": "curriculo", "rel": "depende"}, {"alvo": "word", "confianca": 0.5, "epistemico": "fato", "origem": "wiki", "rel": "relaciona"}, {"alvo": "virtualizacao", "confianca": 0.5, "epistemico": "fato", "origem": "wiki", "rel": "relaciona"}], "confianca": 1.0, "contraexemplos": [], "descricao": "Unidade mínima de informação; 0 ou 1; base de Von Neumann e do gato broca-so.", "epistemico": "fato", "exemplos": [], "id": "bit", "memoria": "semantica", "meta": {"dominio": "computacao", "nivel": 4, "ordem": 1}, "origem": "curriculo", "palavras_chave": ["booleano", "shannon"], "sinonimos": ["bits"], "tipo": "conceito", "titulo": "bit"}
\section{bit}
Unidade mínima de informação; 0 ou 1; base de Von Neumann e do gato broca-so.
\prov{sinónimos: bits}
\prov{relações: usa circuito\_logico; depende inteiros; relaciona word; relaciona virtualizacao}
\prov{proveniência: curriculo \textperiodcentered{} fato \textperiodcentered{} confiança 1.0 \textperiodcentered{} domínio computacao \textperiodcentered{} história 11 versões no jornal}

%CRISTAL {"arestas": [{"alvo": "colapso_multicamada_peso", "confianca": 1.0, "epistemico": "fato", "origem": "wiki", "rel": "usa"}], "confianca": 1.0, "contraexemplos": [], "descricao": "blend_query mistura a impressão da fala com o contexto (α·q+(1−α)·ctx, re-binarizado) antes do scan, e regrava o contexto — a memória de diálogo é 32 bytes no disco, sem estado na RAM.", "epistemico": "fato", "exemplos": [], "id": "blend_contexto_32b", "memoria": "semantica", "meta": {"autor": "broca-so", "dominio": "operacao_so", "fonte": "conversa/colapso.py:blend_query"}, "origem": "wiki", "palavras_chave": [], "sinonimos": [], "tipo": "conceito", "titulo": "Memória Temporal de 32 Bytes (blend α=0.35)"}
\section{Memória Temporal de 32 Bytes (blend \ensuremath{\alpha}=0.35)}
blend\_query mistura a impressão da fala com o contexto (\ensuremath{\alpha}\textperiodcentered q+(1\ensuremath{-}\ensuremath{\alpha})\textperiodcentered ctx, re-binarizado) antes do scan, e regrava o contexto — a memória de diálogo é 32 bytes no disco, sem estado na RAM.
\prov{relações: usa colapso\_multicamada\_peso}
\prov{proveniência: wiki \textperiodcentered{} conversa/colapso.py:blend\_query \textperiodcentered{} fato \textperiodcentered{} confiança 1.0 \textperiodcentered{} domínio operacao\_so \textperiodcentered{} história 2 versões no jornal}

%CRISTAL {"arestas": [{"alvo": "cristalchain", "confianca": 1.0, "epistemico": "fato", "origem": "wiki", "rel": "parte_de"}, {"alvo": "idempotente_e_o_projetor", "confianca": 1.0, "epistemico": "fato", "origem": "wiki", "rel": "eh_um"}, {"alvo": "cristal_motor_congelado", "confianca": 1.0, "epistemico": "fato", "origem": "wiki", "rel": "explica"}, {"alvo": "imutabilidade", "confianca": 1.0, "epistemico": "fato", "origem": "wiki", "rel": "explica"}, {"alvo": "fechamento_blocos", "confianca": 1.0, "epistemico": "fato", "origem": "wiki", "rel": "usa"}], "confianca": 1.0, "contraexemplos": [], "descricao": "Selado, o bloco não muda: re-selar devolve o mesmo selo (P²=P). É o IDEMPOTENTE — o mesmo projetor do estado fundamental que a parte fria encontrou no limite T→0. Minerar esfria até o cristal; o cristal é o idempotente; o idempotente não muda — é a imutabilidade do bloco. A cadeia é uma sucessão de projetores congelados. Verif.: selo=reselo.", "epistemico": "fato", "exemplos": [], "id": "bloco_idempotente_congelado", "memoria": "semantica", "meta": {"dominio": "goldchain", "fonte": "cristalchain"}, "origem": "wiki", "palavras_chave": [], "sinonimos": [], "tipo": "conceito", "titulo": "O Bloco é um Idempotente Congelado"}
\section{O Bloco é um Idempotente Congelado}
Selado, o bloco não muda: re-selar devolve o mesmo selo (P\ensuremath{^{2}}=P). É o IDEMPOTENTE — o mesmo projetor do estado fundamental que a parte fria encontrou no limite T\ensuremath{\to}0. Minerar esfria até o cristal; o cristal é o idempotente; o idempotente não muda — é a imutabilidade do bloco. A cadeia é uma sucessão de projetores congelados. Verif.: selo=reselo.
\prov{relações: parte\_de cristalchain; eh\_um idempotente\_e\_o\_projetor; explica cristal\_motor\_congelado; explica imutabilidade; usa fechamento\_blocos}
\prov{proveniência: wiki \textperiodcentered{} cristalchain \textperiodcentered{} fato \textperiodcentered{} confiança 1.0 \textperiodcentered{} domínio goldchain \textperiodcentered{} história 6 versões no jornal}

%CRISTAL {"arestas": [{"alvo": "computador_mineral", "confianca": 1.0, "epistemico": "fato", "origem": "wiki", "rel": "usa"}, {"alvo": "linguagem_mineral", "confianca": 1.0, "epistemico": "fato", "origem": "wiki", "rel": "usa"}, {"alvo": "biblioteca_mineral", "confianca": 1.0, "epistemico": "fato", "origem": "wiki", "rel": "usa"}, {"alvo": "sandbox_cristal_ao_so", "confianca": 1.0, "epistemico": "fato", "origem": "wiki", "rel": "usa"}, {"alvo": "corpo_mineral", "confianca": 1.0, "epistemico": "fato", "origem": "wiki", "rel": "parte_de"}, {"alvo": "transformada_metalica", "confianca": 1.0, "epistemico": "fato", "origem": "wiki", "rel": "usa"}, {"alvo": "rede_limpa_bump", "confianca": 1.0, "epistemico": "fato", "origem": "wiki", "rel": "relaciona"}, {"alvo": "laco_traducao_bai_peirce", "confianca": 1.0, "epistemico": "fato", "origem": "wiki", "rel": "relaciona"}, {"alvo": "torre_unificada", "confianca": 1.0, "epistemico": "fato", "origem": "wiki", "rel": "relaciona"}], "confianca": 1.0, "contraexemplos": [], "descricao": "O auto-hospedar: o broca-so INTEIRO roda no computador mineral. A ULA é o gato (Aₙ), a memória são os cristais, as bibliotecas são a mecânica/geometria/física no metal, e o SO (autorização) é o BAI. Sobre isso correm: o GRAFO (o tecido de conhecimento), a DSL de gatos (o sandbox, do cristal à linguagem), a MINERAÇÃO (symbolic multimetal = o gato refinando o floco), a REDE (banda=sha256, o canal), e a Tiffany (colapso). A mesma reta x²−nx−1 é a porta, a memória, a física e a lei — o metal é o substrato de tudo. É o fecho: broca-so escrito na sua própria máquina.", "epistemico": "inferencia", "exemplos": [], "id": "broca_so_no_metal", "memoria": "semantica", "meta": {"dominio": "computacao", "fonte": "linguagem mineral"}, "origem": "wiki", "palavras_chave": [], "sinonimos": [], "tipo": "conceito", "titulo": "O Computador Mineral roda o broca-so"}
\section{O Computador Mineral roda o broca-so}
O auto-hospedar: o broca-so INTEIRO roda no computador mineral. A ULA é o gato (A\ensuremath{_{n}}), a memória são os cristais, as bibliotecas são a mecânica/geometria/física no metal, e o SO (autorização) é o BAI. Sobre isso correm: o GRAFO (o tecido de conhecimento), a DSL de gatos (o sandbox, do cristal à linguagem), a MINERAÇÃO (symbolic multimetal = o gato refinando o floco), a REDE (banda=sha256, o canal), e a Tiffany (colapso). A mesma reta x\ensuremath{^{2}}\ensuremath{-}nx\ensuremath{-}1 é a porta, a memória, a física e a lei — o metal é o substrato de tudo. É o fecho: broca-so escrito na sua própria máquina.
\prov{relações: usa computador\_mineral; usa linguagem\_mineral; usa biblioteca\_mineral; usa sandbox\_cristal\_ao\_so; parte\_de corpo\_mineral; usa transformada\_metalica; relaciona rede\_limpa\_bump; relaciona laco\_traducao\_bai\_peirce; relaciona torre\_unificada}
\prov{proveniência: wiki \textperiodcentered{} linguagem mineral \textperiodcentered{} inferencia \textperiodcentered{} confiança 1.0 \textperiodcentered{} domínio computacao \textperiodcentered{} história 10 versões no jornal}

%CRISTAL {"arestas": [{"alvo": "unix", "confianca": 1.0, "epistemico": "fato", "origem": "wiki", "rel": "especializa"}], "confianca": 1.0, "contraexemplos": [], "descricao": "Berkeley Software Distribution: o ramo acadêmico do Unix (sockets de rede, o kernel que virou base do macOS e do coração de muita infraestrutura).", "epistemico": "fato", "exemplos": [], "id": "bsd", "memoria": "semantica", "meta": {"dominio": "programacao", "fonte": ""}, "origem": "wiki", "palavras_chave": [], "sinonimos": [], "tipo": "conceito", "titulo": "BSD"}
\section{BSD}
Berkeley Software Distribution: o ramo acadêmico do Unix (sockets de rede, o kernel que virou base do macOS e do coração de muita infraestrutura).
\prov{relações: especializa unix}
\prov{proveniência: wiki \textperiodcentered{} fato \textperiodcentered{} confiança 1.0 \textperiodcentered{} domínio programacao \textperiodcentered{} história 2 versões no jornal}

%CRISTAL {"arestas": [{"alvo": "minimax", "confianca": 1.0, "epistemico": "fato", "origem": "wiki", "rel": "relaciona"}, {"alvo": "inteligencia_artificial", "confianca": 1.0, "epistemico": "fato", "origem": "wiki", "rel": "parte_de"}], "confianca": 1.0, "contraexemplos": [], "descricao": "Em vez de expandir toda a árvore, amostra jogadas até o fim (rollouts) e concentra a busca nos ramos promissores. Domou o Go (fator de ramificação ~250) onde o minimax não chegava — e é o motor do AlphaGo.", "epistemico": "fato", "exemplos": [], "id": "busca_monte_carlo_arvore", "memoria": "semantica", "meta": {"dominio": "ia", "fonte": "jogos e IA"}, "origem": "wiki", "palavras_chave": [], "sinonimos": [], "tipo": "conceito", "titulo": "Busca em Árvore Monte Carlo (MCTS)"}
\section{Busca em Árvore Monte Carlo (MCTS)}
Em vez de expandir toda a árvore, amostra jogadas até o fim (rollouts) e concentra a busca nos ramos promissores. Domou o Go (fator de ramificação \~{}250) onde o minimax não chegava — e é o motor do AlphaGo.
\prov{relações: relaciona minimax; parte\_de inteligencia\_artificial}
\prov{proveniência: wiki \textperiodcentered{} jogos e IA \textperiodcentered{} fato \textperiodcentered{} confiança 1.0 \textperiodcentered{} domínio ia \textperiodcentered{} história 3 versões no jornal}

%CRISTAL {"arestas": [{"alvo": "maquina_virtual", "confianca": 1.0, "epistemico": "fato", "origem": "wiki", "rel": "usa"}, {"alvo": "linguagem_matriz_gatos", "confianca": 1.0, "epistemico": "fato", "origem": "wiki", "rel": "relaciona"}, {"alvo": "backend_isa12_broca", "confianca": 1.0, "epistemico": "fato", "origem": "wiki", "rel": "relaciona"}, {"alvo": "dispatch_opcodes", "confianca": 1.0, "epistemico": "fato", "origem": "wiki", "rel": "relaciona"}], "confianca": 1.0, "contraexemplos": [], "descricao": "Um código intermediário compacto para uma máquina virtual — entre o fonte e o metal; como a ISA-12 do broca-so.", "epistemico": "fato", "exemplos": [], "id": "bytecode", "memoria": "semantica", "meta": {"dominio": "programacao", "fonte": ""}, "origem": "wiki", "palavras_chave": [], "sinonimos": [], "tipo": "conceito", "titulo": "Bytecode"}
\section{Bytecode}
Um código intermediário compacto para uma máquina virtual — entre o fonte e o metal; como a ISA-12 do broca-so.
\prov{relações: usa maquina\_virtual; relaciona linguagem\_matriz\_gatos; relaciona backend\_isa12\_broca; relaciona dispatch\_opcodes}
\prov{proveniência: wiki \textperiodcentered{} fato \textperiodcentered{} confiança 1.0 \textperiodcentered{} domínio programacao \textperiodcentered{} história 8 versões no jornal}

%CRISTAL {"arestas": [{"alvo": "compilador", "confianca": 1.0, "epistemico": "fato", "origem": "curriculo", "rel": "referencia"}, {"alvo": "virtualizacao", "confianca": 0.5, "epistemico": "fato", "origem": "wiki", "rel": "relaciona"}, {"alvo": "word", "confianca": 0.5, "epistemico": "fato", "origem": "wiki", "rel": "relaciona"}, {"alvo": "escassez", "confianca": 0.5, "epistemico": "fato", "origem": "wiki", "rel": "relaciona"}, {"alvo": "linguagem_de_programacao", "confianca": 1.0, "epistemico": "fato", "origem": "wiki", "rel": "parte_de"}, {"alvo": "programacao_imperativa", "confianca": 1.0, "epistemico": "fato", "origem": "wiki", "rel": "usa"}, {"alvo": "gerenciamento_de_memoria", "confianca": 1.0, "epistemico": "fato", "origem": "wiki", "rel": "usa"}], "confianca": 1.0, "contraexemplos": [], "descricao": "Ritchie (1972): a linguagem de sistemas — ponteiros, memória manual, próxima do metal. A língua franca do baixo nível (e do broca-so).", "epistemico": "fato", "exemplos": [], "id": "c", "memoria": "semantica", "meta": {"autor": "Ritchie", "dominio": "programacao", "fonte": ""}, "origem": "wiki", "palavras_chave": [], "sinonimos": ["Novas Sequências Aritméticas e Geométricas", "livro Gentil"], "tipo": "conceito", "titulo": "C"}
\section{C}
Ritchie (1972): a linguagem de sistemas — ponteiros, memória manual, próxima do metal. A língua franca do baixo nível (e do broca-so).
\prov{sinónimos: Novas Sequências Aritméticas e Geométricas; livro Gentil}
\prov{relações: referencia compilador; relaciona virtualizacao; relaciona word; relaciona escassez; parte\_de linguagem\_de\_programacao; usa programacao\_imperativa; usa gerenciamento\_de\_memoria}
\prov{proveniência: wiki \textperiodcentered{} fato \textperiodcentered{} confiança 1.0 \textperiodcentered{} domínio programacao \textperiodcentered{} história 47 versões no jornal}

%CRISTAL {"arestas": [{"alvo": "virtualizacao", "confianca": 1.0, "epistemico": "fato", "origem": "curriculo", "rel": "generaliza"}, {"alvo": "kernel", "confianca": 0.758, "epistemico": "fato", "origem": "manual", "rel": "referencia"}, {"alvo": "rust", "confianca": 0.5, "epistemico": "fato", "origem": "wiki", "rel": "relaciona"}, {"alvo": "historia", "confianca": 0.5, "epistemico": "fato", "origem": "wiki", "rel": "relaciona"}, {"alvo": "gato_operador_hub", "confianca": 0.5, "epistemico": "fato", "origem": "wiki", "rel": "relaciona"}, {"alvo": "recursao", "confianca": 0.5, "epistemico": "fato", "origem": "wiki", "rel": "relaciona"}, {"alvo": "word", "confianca": 0.5, "epistemico": "fato", "origem": "wiki", "rel": "relaciona"}, {"alvo": "syscall", "confianca": 0.5, "epistemico": "fato", "origem": "wiki", "rel": "relaciona"}, {"alvo": "octonios_hiperbolicos", "confianca": 0.5, "epistemico": "fato", "origem": "wiki", "rel": "relaciona"}, {"alvo": "criterio_de_projeto", "confianca": 0.5, "epistemico": "fato", "origem": "wiki", "rel": "relaciona"}, {"alvo": "sagrado_profano", "confianca": 0.5, "epistemico": "fato", "origem": "wiki", "rel": "relaciona"}, {"alvo": "eterno_retorno", "confianca": 0.5, "epistemico": "fato", "origem": "wiki", "rel": "relaciona"}, {"alvo": "vida-morte", "confianca": 0.5, "epistemico": "fato", "origem": "wiki", "rel": "relaciona"}, {"alvo": "gaia", "confianca": 0.5, "epistemico": "fato", "origem": "wiki", "rel": "relaciona"}, {"alvo": "equacao_efectiva_conjectura", "confianca": 0.5, "epistemico": "fato", "origem": "wiki", "rel": "relaciona"}, {"alvo": "top_assinaturas", "confianca": 0.5, "epistemico": "fato", "origem": "wiki", "rel": "relaciona"}, {"alvo": "pergunta", "confianca": 0.5, "epistemico": "fato", "origem": "wiki", "rel": "relaciona"}, {"alvo": "banco_de_dados", "confianca": 1.0, "epistemico": "fato", "origem": "wiki", "rel": "relaciona"}], "confianca": 1.0, "contraexemplos": [], "descricao": "Localidade temporal e espacial.", "epistemico": "fato", "exemplos": [], "id": "cache", "memoria": "semantica", "meta": {"dominio": "sistemas", "nivel": 6, "ordem": 4}, "origem": "curriculo", "palavras_chave": [], "sinonimos": [], "tipo": "conceito", "titulo": "cache"}
\section{cache}
Localidade temporal e espacial.
\prov{relações: generaliza virtualizacao; referencia kernel; relaciona rust; relaciona historia; relaciona gato\_operador\_hub; relaciona recursao; relaciona word; relaciona syscall; relaciona octonios\_hiperbolicos; relaciona criterio\_de\_projeto; relaciona sagrado\_profano; relaciona eterno\_retorno; relaciona vida-morte; relaciona gaia; relaciona equacao\_efectiva\_conjectura; relaciona top\_assinaturas; relaciona pergunta; relaciona banco\_de\_dados}
\prov{proveniência: curriculo \textperiodcentered{} fato \textperiodcentered{} confiança 1.0 \textperiodcentered{} domínio sistemas \textperiodcentered{} história 24 versões no jornal}

%CRISTAL {"arestas": [{"alvo": "bai_tese_mae", "confianca": 1.0, "epistemico": "fato", "origem": "wiki", "rel": "usa"}, {"alvo": "bai_delta", "confianca": 1.0, "epistemico": "fato", "origem": "wiki", "rel": "usa"}, {"alvo": "colapso_multicamada_peso", "confianca": 1.0, "epistemico": "fato", "origem": "wiki", "rel": "relaciona"}], "confianca": 1.0, "contraexemplos": [], "descricao": "As cadeias exportam turnos 'user — Tiffanytfluxo:id:NN:dep=D'; a propagação é validada turno a turno por grantok (o turno certo da cadeia certa deve vencer) — a capacidade CASL propagando com orçamento δ decrescente. É o BAI operando na conversa.", "epistemico": "fato", "exemplos": [], "id": "cadeias_dep_grant", "memoria": "semantica", "meta": {"autor": "broca-so", "dominio": "operacao_so", "fonte": "conversa/cadeias.py"}, "origem": "wiki", "palavras_chave": [], "sinonimos": [], "tipo": "conceito", "titulo": "Cadeias com δ (dep) e Grant CASL por Turno"}
\section{Cadeias com \ensuremath{\delta} (dep) e Grant CASL por Turno}
As cadeias exportam turnos 'user — Tiffanytfluxo:id:NN:dep=D'; a propagação é validada turno a turno por grantok (o turno certo da cadeia certa deve vencer) — a capacidade CASL propagando com orçamento \ensuremath{\delta} decrescente. É o BAI operando na conversa.
\prov{relações: usa bai\_tese\_mae; usa bai\_delta; relaciona colapso\_multicamada\_peso}
\prov{proveniência: wiki \textperiodcentered{} conversa/cadeias.py \textperiodcentered{} fato \textperiodcentered{} confiança 1.0 \textperiodcentered{} domínio operacao\_so \textperiodcentered{} história 5 versões no jornal}

%CRISTAL {"arestas": [{"alvo": "corpo_f2k", "confianca": 1.0, "epistemico": "fato", "origem": "wiki", "rel": "especializa"}, {"alvo": "dimensao_de_prata_pell", "confianca": 1.0, "epistemico": "fato", "origem": "wiki", "rel": "usa"}, {"alvo": "a_economia_do_pell_---_a_moeda_do_protocolo", "confianca": 1.0, "epistemico": "fato", "origem": "wiki", "rel": "parte_de"}, {"alvo": "ciclo_uno", "confianca": 1.0, "epistemico": "fato", "origem": "wiki", "rel": "usa"}, {"alvo": "cofre_parabola_hero", "confianca": 1.0, "epistemico": "fato", "origem": "wiki", "rel": "relaciona"}, {"alvo": "seguranca_algebrica_estrutura_nao_hardness", "confianca": 1.0, "epistemico": "fato", "origem": "wiki", "rel": "usa"}, {"alvo": "fabrica_ordem_construtiva", "confianca": 1.0, "epistemico": "fato", "origem": "wiki", "rel": "relaciona"}, {"alvo": "goldchain", "confianca": 1.0, "epistemico": "fato", "origem": "wiki", "rel": "parte_de"}], "confianca": 1.0, "contraexemplos": [], "descricao": "corpoprata: a prata do cliente vive DENTRO do corpo 𝔽₂ₖ (K=128, p=x¹²⁸+x⁷+x²+x+1 irredutível), tudo em BITS (+ = XOR, · = polinômio mod p carryless, chicotada = Frobenius x²) — 'o corpo vive sozinho dentro do nada' (sem numpy, sem float). A chave do cliente é um elemento g do corpo (a base); minerar a prata na base dele = ×g (mudança de base, LINEAR sobre 𝔽₂), reversível por g⁻¹ (Fermat). SÓ o dono volta à base padrão e lê um PRIMO DE PRATA (Pell primo); um terceiro lê lixo (não é primo de prata). É o cofre CONSTRUÍDO (não arrombável) da economia de prata. [N] código que roda em bits.", "epistemico": "inferencia", "exemplos": [], "id": "caixa_de_prata_corpo", "memoria": "semantica", "meta": {"autor": "Lopes/broca-so", "dominio": "goldchain", "fonte": "machine/so/corpo_prata.py; prata.py"}, "origem": "wiki", "palavras_chave": [], "sinonimos": [], "tipo": "conceito", "titulo": "A Caixa de Prata — a prata na base do cliente (corpo 𝔽₂ₖ)"}
\section{A Caixa de Prata — a prata na base do cliente (corpo \ensuremath{\mathbb{F}}\ensuremath{_{2k}})}
corpoprata: a prata do cliente vive DENTRO do corpo \ensuremath{\mathbb{F}}\ensuremath{_{2k}} (K=128, p=x\ensuremath{^{128}}+x\ensuremath{^{7}}+x\ensuremath{^{2}}+x+1 irredutível), tudo em BITS (+ = XOR, \textperiodcentered  = polinômio mod p carryless, chicotada = Frobenius x\ensuremath{^{2}}) — 'o corpo vive sozinho dentro do nada' (sem numpy, sem float). A chave do cliente é um elemento g do corpo (a base); minerar a prata na base dele = $\times$g (mudança de base, LINEAR sobre \ensuremath{\mathbb{F}}\ensuremath{_{2}}), reversível por g\ensuremath{^{-1}} (Fermat). SÓ o dono volta à base padrão e lê um PRIMO DE PRATA (Pell primo); um terceiro lê lixo (não é primo de prata). É o cofre CONSTRUÍDO (não arrombável) da economia de prata. [N] código que roda em bits.
\prov{relações: especializa corpo\_f2k; usa dimensao\_de\_prata\_pell; parte\_de a\_economia\_do\_pell\_---\_a\_moeda\_do\_protocolo; usa ciclo\_uno; relaciona cofre\_parabola\_hero; usa seguranca\_algebrica\_estrutura\_nao\_hardness; relaciona fabrica\_ordem\_construtiva; parte\_de goldchain}
\prov{proveniência: wiki \textperiodcentered{} machine/so/corpo\_prata.py; prata.py \textperiodcentered{} inferencia \textperiodcentered{} confiança 1.0 \textperiodcentered{} domínio goldchain \textperiodcentered{} história 11 versões no jornal}

%CRISTAL {"arestas": [{"alvo": "entropia_shannon", "confianca": 1.0, "epistemico": "fato", "origem": "wiki", "rel": "usa"}], "confianca": 1.0, "contraexemplos": [], "descricao": "Todo canal ruidoso tem uma capacidade C; para qualquer taxa R<C existe código que torna o erro arbitrariamente pequeno, e para R>C a transmissão confiável é impossível (o 'limite de Shannon'). Caso AWGN: C=B log₂(1+S/N) [bits/s]. Também estabelece a separação fonte-canal. Fato (teorema provado). NOTA: não existe na série Lopes; o κ do BAI é capacidade de autorização, não de canal.", "epistemico": "fato", "exemplos": [], "id": "capacidade_canal_shannon", "memoria": "semantica", "meta": {"autor": "broca-so", "dominio": "computacao", "fonte": "Shannon 1948 (noisy-channel coding theorem)"}, "origem": "wiki", "palavras_chave": [], "sinonimos": [], "tipo": "conceito", "titulo": "Capacidade de canal e o teorema do canal ruidoso"}
\section{Capacidade de canal e o teorema do canal ruidoso}
Todo canal ruidoso tem uma capacidade C; para qualquer taxa R<C existe código que torna o erro arbitrariamente pequeno, e para R>C a transmissão confiável é impossível (o 'limite de Shannon'). Caso AWGN: C=B log\ensuremath{_{2}}(1+S/N) [bits/s]. Também estabelece a separação fonte-canal. Fato (teorema provado). NOTA: não existe na série Lopes; o \ensuremath{\kappa} do BAI é capacidade de autorização, não de canal.
\prov{relações: usa entropia\_shannon}
\prov{proveniência: wiki \textperiodcentered{} Shannon 1948 (noisy-channel coding theorem) \textperiodcentered{} fato \textperiodcentered{} confiança 1.0 \textperiodcentered{} domínio computacao \textperiodcentered{} história 2 versões no jornal}

%CRISTAL {"arestas": [{"alvo": "poker_jogo", "confianca": 1.0, "epistemico": "fato", "origem": "wiki", "rel": "usa"}, {"alvo": "equilibrio_de_nash", "confianca": 1.0, "epistemico": "fato", "origem": "wiki", "rel": "usa"}], "confianca": 1.0, "contraexemplos": [], "descricao": "O algoritmo que resolve jogos de informação imperfeita: itera medindo o 'arrependimento' de cada ação e converge para o equilíbrio de Nash. A base dos bots de pôquer sobre-humanos.", "epistemico": "fato", "exemplos": [], "id": "cfr_minimizacao_arrependimento", "memoria": "semantica", "meta": {"dominio": "ia", "fonte": "jogos e IA"}, "origem": "wiki", "palavras_chave": [], "sinonimos": [], "tipo": "conceito", "titulo": "CFR (Minimização de Arrependimento Contrafactual)"}
\section{CFR (Minimização de Arrependimento Contrafactual)}
O algoritmo que resolve jogos de informação imperfeita: itera medindo o 'arrependimento' de cada ação e converge para o equilíbrio de Nash. A base dos bots de pôquer sobre-humanos.
\prov{relações: usa poker\_jogo; usa equilibrio\_de\_nash}
\prov{proveniência: wiki \textperiodcentered{} jogos e IA \textperiodcentered{} fato \textperiodcentered{} confiança 1.0 \textperiodcentered{} domínio ia \textperiodcentered{} história 3 versões no jornal}

%CRISTAL {"arestas": [{"alvo": "isolamento_base_cliente", "confianca": 1.0, "epistemico": "fato", "origem": "wiki", "rel": "usa"}], "confianca": 1.0, "contraexemplos": [], "descricao": "Identificador dominio/…/NNNNNNN (7 dígitos) validado por regex; é a âncora estável para gold sets e merge incremental, e a chave que detecta duplicação entre base e cliente.", "epistemico": "fato", "exemplos": [], "id": "chunk_id_global_unico", "memoria": "semantica", "meta": {"autor": "broca-so", "dominio": "operacao_so", "fonte": "code/tiffany_platform/chunk_id.py"}, "origem": "wiki", "palavras_chave": [], "sinonimos": [], "tipo": "conceito", "titulo": "chunk_id Único e Estável"}
\section{chunk\_id Único e Estável}
Identificador dominio/…/NNNNNNN (7 dígitos) validado por regex; é a âncora estável para gold sets e merge incremental, e a chave que detecta duplicação entre base e cliente.
\prov{relações: usa isolamento\_base\_cliente}
\prov{proveniência: wiki \textperiodcentered{} code/tiffany\_platform/chunk\_id.py \textperiodcentered{} fato \textperiodcentered{} confiança 1.0 \textperiodcentered{} domínio operacao\_so \textperiodcentered{} história 2 versões no jornal}

%CRISTAL {"arestas": [{"alvo": "devops", "confianca": 1.0, "epistemico": "fato", "origem": "wiki", "rel": "parte_de"}, {"alvo": "git", "confianca": 1.0, "epistemico": "fato", "origem": "wiki", "rel": "usa"}, {"alvo": "cristal_pipeline_runner", "confianca": 1.0, "epistemico": "fato", "origem": "wiki", "rel": "relaciona"}], "confianca": 1.0, "contraexemplos": [], "descricao": "Integrar e publicar mudanças pequenas e frequentes, testadas automaticamente por uma esteira — a pipeline que leva do commit à produção.", "epistemico": "inferencia", "exemplos": [], "id": "ci_cd", "memoria": "semantica", "meta": {"dominio": "programacao", "fonte": ""}, "origem": "wiki", "palavras_chave": [], "sinonimos": [], "tipo": "conceito", "titulo": "CI/CD (Integração e Entrega Contínuas)"}
\section{CI/CD (Integração e Entrega Contínuas)}
Integrar e publicar mudanças pequenas e frequentes, testadas automaticamente por uma esteira — a pipeline que leva do commit à produção.
\prov{relações: parte\_de devops; usa git; relaciona cristal\_pipeline\_runner}
\prov{proveniência: wiki \textperiodcentered{} inferencia \textperiodcentered{} confiança 1.0 \textperiodcentered{} domínio programacao \textperiodcentered{} história 8 versões no jornal}

%CRISTAL {"arestas": [{"alvo": "teoria_dos_jogos", "confianca": 1.0, "epistemico": "fato", "origem": "wiki", "rel": "usa"}, {"alvo": "informacao_imperfeita", "confianca": 1.0, "epistemico": "fato", "origem": "wiki", "rel": "usa"}], "confianca": 1.0, "contraexemplos": [], "descricao": "IA da Meta que joga Diplomacy em nível humano — um jogo de NEGOCIAÇÃO em linguagem natural, alianças e traição. Une teoria dos jogos, modelagem de intenções e diálogo: o social entra no tabuleiro.", "epistemico": "fato", "exemplos": [], "id": "cicero_diplomacy", "memoria": "semantica", "meta": {"autor": "Meta", "dominio": "ia", "fonte": "jogos e IA"}, "origem": "wiki", "palavras_chave": [], "sinonimos": [], "tipo": "conceito", "titulo": "Cicero (Diplomacy)"}
\section{Cicero (Diplomacy)}
IA da Meta que joga Diplomacy em nível humano — um jogo de NEGOCIAÇÃO em linguagem natural, alianças e traição. Une teoria dos jogos, modelagem de intenções e diálogo: o social entra no tabuleiro.
\prov{relações: usa teoria\_dos\_jogos; usa informacao\_imperfeita}
\prov{proveniência: wiki \textperiodcentered{} jogos e IA \textperiodcentered{} fato \textperiodcentered{} confiança 1.0 \textperiodcentered{} domínio ia \textperiodcentered{} história 3 versões no jornal}

%CRISTAL {"arestas": [{"alvo": "contabilidade_sem_livro", "confianca": 1.0, "epistemico": "fato", "origem": "wiki", "rel": "usa"}, {"alvo": "lastro_do_cofre", "confianca": 1.0, "epistemico": "fato", "origem": "wiki", "rel": "relaciona"}], "confianca": 1.0, "contraexemplos": [], "descricao": "Minerar gasta ~900 sats de fee e emite um Pell; a liquidação desse Pell DEVOLVE a mesma fee ao ramo (fee_devolvida_sats), sem criar movimento extra. O lastro não é consumido, só circula — a máquina se autoalimenta.", "epistemico": "inferencia", "exemplos": [], "id": "ciclo_autoconsumo", "memoria": "semantica", "meta": {"autor": "broca-so", "dominio": "operacao_so", "fonte": "goldchain/liquidacao_pell.py:devolver_fee_ciclo"}, "origem": "wiki", "palavras_chave": [], "sinonimos": [], "tipo": "conceito", "titulo": "Auto-Consumo de Fee (a fee volta)"}
\section{Auto-Consumo de Fee (a fee volta)}
Minerar gasta \~{}900 sats de fee e emite um Pell; a liquidação desse Pell DEVOLVE a mesma fee ao ramo (fee\_devolvida\_sats), sem criar movimento extra. O lastro não é consumido, só circula — a máquina se autoalimenta.
\prov{relações: usa contabilidade\_sem\_livro; relaciona lastro\_do\_cofre}
\prov{proveniência: wiki \textperiodcentered{} goldchain/liquidacao\_pell.py:devolver\_fee\_ciclo \textperiodcentered{} inferencia \textperiodcentered{} confiança 1.0 \textperiodcentered{} domínio operacao\_so \textperiodcentered{} história 3 versões no jornal}

%CRISTAL {"arestas": [{"alvo": "goldchain", "confianca": 1.0, "epistemico": "fato", "origem": "wiki", "rel": "explica"}, {"alvo": "htlc", "confianca": 1.0, "epistemico": "fato", "origem": "wiki", "rel": "usa"}, {"alvo": "mineracao_hexagrama", "confianca": 1.0, "epistemico": "fato", "origem": "wiki", "rel": "usa"}, {"alvo": "prova_de_trabalho", "confianca": 1.0, "epistemico": "fato", "origem": "wiki", "rel": "usa"}], "confianca": 1.0, "contraexemplos": [], "descricao": "A vida real de uma transação é uma função única: PoW → escreve fase 'pow' → minerar_um emite Pell (verify ν=6) → fase 'lock' (lock HTLC) → fase 'redeem' (loop até 20 tentativas esperando UTXO ≥546 no HTLC) → registra o ciclo. Timings por fase medidos, gargalo reportado.", "epistemico": "fato", "exemplos": [], "id": "ciclo_uno", "memoria": "semantica", "meta": {"autor": "broca-so", "dominio": "operacao_so", "fonte": "goldchain/producao.py:_executar_ciclo_uno"}, "origem": "wiki", "palavras_chave": [], "sinonimos": [], "tipo": "conceito", "titulo": "O Ciclo Uno da Tesouraria (PoW→Pell→lock→redeem)"}
\section{O Ciclo Uno da Tesouraria (PoW\ensuremath{\to}Pell\ensuremath{\to}lock\ensuremath{\to}redeem)}
A vida real de uma transação é uma função única: PoW \ensuremath{\to} escreve fase 'pow' \ensuremath{\to} minerar\_um emite Pell (verify \ensuremath{\nu}=6) \ensuremath{\to} fase 'lock' (lock HTLC) \ensuremath{\to} fase 'redeem' (loop até 20 tentativas esperando UTXO \ensuremath{\ge}546 no HTLC) \ensuremath{\to} registra o ciclo. Timings por fase medidos, gargalo reportado.
\prov{relações: explica goldchain; usa htlc; usa mineracao\_hexagrama; usa prova\_de\_trabalho}
\prov{proveniência: wiki \textperiodcentered{} goldchain/producao.py:\_executar\_ciclo\_uno \textperiodcentered{} fato \textperiodcentered{} confiança 1.0 \textperiodcentered{} domínio operacao\_so \textperiodcentered{} história 5 versões no jornal}

%CRISTAL {"arestas": [{"alvo": "computabilidade", "confianca": 1.0, "epistemico": "fato", "origem": "wiki", "rel": "usa"}, {"alvo": "computacao", "confianca": 1.0, "epistemico": "fato", "origem": "wiki", "rel": "relaciona"}], "confianca": 1.0, "contraexemplos": [], "descricao": "O estudo do cálculo, da informação e dos algoritmos — do computável abstrato ao software concreto.", "epistemico": "fato", "exemplos": [], "id": "ciencia_da_computacao", "memoria": "semantica", "meta": {"dominio": "programacao", "fonte": ""}, "origem": "wiki", "palavras_chave": [], "sinonimos": [], "tipo": "conceito", "titulo": "Ciência da Computação"}
\section{Ciência da Computação}
O estudo do cálculo, da informação e dos algoritmos — do computável abstrato ao software concreto.
\prov{relações: usa computabilidade; relaciona computacao}
\prov{proveniência: wiki \textperiodcentered{} fato \textperiodcentered{} confiança 1.0 \textperiodcentered{} domínio programacao \textperiodcentered{} história 5 versões no jornal}

%CRISTAL {"arestas": [{"alvo": "lado_negro_de_newton", "confianca": 1.0, "epistemico": "fato", "origem": "wiki", "rel": "usa"}, {"alvo": "idempotentes_mod_n_fatoram", "confianca": 1.0, "epistemico": "fato", "origem": "wiki", "rel": "usa"}, {"alvo": "corpo_mineral", "confianca": 1.0, "epistemico": "fato", "origem": "wiki", "rel": "usa"}, {"alvo": "julia", "confianca": 1.0, "epistemico": "fato", "origem": "wiki", "rel": "relaciona"}], "confianca": 1.0, "contraexemplos": [], "descricao": "A criptografia da broca-so mora exatamente onde Newton dos idempotentes QUEBRA: a fronteira Re z=½ (a lemniscata entre o lado negro, 0, e o claro, 1) e o corpo discreto de primos sobre ela. Três peças: o gato embaralha (simétrica), o primo protege (cristal), o idempotente é a fenda (a dureza = fatorar). Paper: papers/cifra_de_cristal.tex.", "epistemico": "fato", "exemplos": [], "id": "cifra_de_cristal", "memoria": "semantica", "meta": {"dominio": "criptografia", "fonte": "cifra de cristal"}, "origem": "wiki", "palavras_chave": [], "sinonimos": [], "tipo": "conceito", "titulo": "A Cifra de Cristal (o corpo discreto na lemniscata)"}
\section{A Cifra de Cristal (o corpo discreto na lemniscata)}
A criptografia da broca-so mora exatamente onde Newton dos idempotentes QUEBRA: a fronteira Re z=\ensuremath{\tfrac12} (a lemniscata entre o lado negro, 0, e o claro, 1) e o corpo discreto de primos sobre ela. Três peças: o gato embaralha (simétrica), o primo protege (cristal), o idempotente é a fenda (a dureza = fatorar). Paper: papers/cifra\_de\_cristal.tex.
\prov{relações: usa lado\_negro\_de\_newton; usa idempotentes\_mod\_n\_fatoram; usa corpo\_mineral; relaciona julia}
\prov{proveniência: wiki \textperiodcentered{} cifra de cristal \textperiodcentered{} fato \textperiodcentered{} confiança 1.0 \textperiodcentered{} domínio criptografia \textperiodcentered{} história 5 versões no jornal}

%CRISTAL {"arestas": [{"alvo": "gato", "confianca": 1.0, "epistemico": "fato", "origem": "curriculo", "rel": "referencia"}, {"alvo": "goldchain", "confianca": 1.0, "epistemico": "fato", "origem": "paper", "rel": "explica"}, {"alvo": "sha256", "confianca": 1.0, "epistemico": "fato", "origem": "paper", "rel": "usa"}], "confianca": 1.0, "contraexemplos": [], "descricao": "Criptografia do projeto: Pell cifrado, tranca vertical/lateral, logo ouro branco (kernel/cifra_ouro_branco.tex).", "epistemico": "fato", "exemplos": [], "id": "cifra_de_ouro_branco", "memoria": "semantica", "meta": {"dominio": "cripto", "nivel": 5, "ordem": 1}, "origem": "curriculo", "palavras_chave": ["Pell", "custódia", "tranca"], "sinonimos": ["cifra ouro branco", "ouro branco"], "tipo": "conceito", "titulo": "cifra de ouro branco"}
\section{cifra de ouro branco}
Criptografia do projeto: Pell cifrado, tranca vertical/lateral, logo ouro branco (kernel/cifra\_ouro\_branco.tex).
\prov{sinónimos: cifra ouro branco; ouro branco}
\prov{relações: referencia gato; explica goldchain; usa sha256}
\prov{proveniência: curriculo \textperiodcentered{} fato \textperiodcentered{} confiança 1.0 \textperiodcentered{} domínio cripto \textperiodcentered{} história 66 versões no jornal}

%CRISTAL {"arestas": [{"alvo": "cifra_de_cristal", "confianca": 1.0, "epistemico": "fato", "origem": "wiki", "rel": "parte_de"}, {"alvo": "gato_como_cifra", "confianca": 1.0, "epistemico": "fato", "origem": "wiki", "rel": "eh_um"}, {"alvo": "passos_da_agulha_sao_cifra", "confianca": 1.0, "epistemico": "fato", "origem": "wiki", "rel": "relaciona"}, {"alvo": "criptografia_simetrica", "confianca": 1.0, "epistemico": "fato", "origem": "wiki", "rel": "eh_um"}], "confianca": 1.0, "contraexemplos": [], "descricao": "O gato metálico A_n=[[n,1],[1,0]] tem det=−1, logo é BIJEÇÃO módulo qualquer N. Cifrar é aplicar A_n^k mod N; decifrar é o gato inverso A_n^{-1}=[[0,1],[1,-n]] elevado a k. A chave é (n,k,N); a recuperação é exata (a mesma cifra que embaralha a imagem e gera a fração contínua de π). Verificado: BROCA-SO cifra e decifra exato.", "epistemico": "fato", "exemplos": [], "id": "cifra_do_gato_simetrica", "memoria": "semantica", "meta": {"dominio": "criptografia", "fonte": "cifra de cristal"}, "origem": "wiki", "palavras_chave": [], "sinonimos": [], "tipo": "conceito", "titulo": "A Cifra do Gato (simétrica, det=±1)"}
\section{A Cifra do Gato (simétrica, det=$\pm$1)}
O gato metálico A\_n=[[n,1],[1,0]] tem det=\ensuremath{-}1, logo é BIJEÇÃO módulo qualquer N. Cifrar é aplicar A\_n\^{}k mod N; decifrar é o gato inverso A\_n\^{}\{-1\}=[[0,1],[1,-n]] elevado a k. A chave é (n,k,N); a recuperação é exata (a mesma cifra que embaralha a imagem e gera a fração contínua de \ensuremath{\pi}). Verificado: BROCA-SO cifra e decifra exato.
\prov{relações: parte\_de cifra\_de\_cristal; eh\_um gato\_como\_cifra; relaciona passos\_da\_agulha\_sao\_cifra; eh\_um criptografia\_simetrica}
\prov{proveniência: wiki \textperiodcentered{} cifra de cristal \textperiodcentered{} fato \textperiodcentered{} confiança 1.0 \textperiodcentered{} domínio criptografia \textperiodcentered{} história 5 versões no jornal}

%CRISTAL {"arestas": [{"alvo": "cifra", "confianca": 1.0, "epistemico": "fato", "origem": "wiki", "rel": "especializa"}, {"alvo": "word_espectral", "confianca": 1.0, "epistemico": "fato", "origem": "wiki", "rel": "usa"}, {"alvo": "transformada_metalica", "confianca": 1.0, "epistemico": "fato", "origem": "wiki", "rel": "aplicacao"}, {"alvo": "beta_numeracao_metalica", "confianca": 1.0, "epistemico": "fato", "origem": "wiki", "rel": "usa"}], "confianca": 1.0, "contraexemplos": [], "descricao": "A operação de metal cifra_an(w,n)=[n·total+e, total] é a matriz do gato aplicada ao word: gold=A₁, silver=A₂, bronze=A₃. Com FOLD/UNFOLD/PROJECT/LIFT (parábola composta com metal) e o modo SPECT, em que cada LOAD cifra automaticamente com A₂.", "epistemico": "fato", "exemplos": [], "id": "cifra_metais_an", "memoria": "semantica", "meta": {"autor": "broca-so", "dominio": "operacao_so", "fonte": "ula/cifra.c"}, "origem": "wiki", "palavras_chave": [], "sinonimos": [], "tipo": "conceito", "titulo": "A Cifra Aₙ (silver/gold/bronze)"}
\section{A Cifra A\ensuremath{_{n}} (silver/gold/bronze)}
A operação de metal cifra\_an(w,n)=[n\textperiodcentered total+e, total] é a matriz do gato aplicada ao word: gold=A\ensuremath{_{1}}, silver=A\ensuremath{_{2}}, bronze=A\ensuremath{_{3}}. Com FOLD/UNFOLD/PROJECT/LIFT (parábola composta com metal) e o modo SPECT, em que cada LOAD cifra automaticamente com A\ensuremath{_{2}}.
\prov{relações: especializa cifra; usa word\_espectral; aplicacao transformada\_metalica; usa beta\_numeracao\_metalica}
\prov{proveniência: wiki \textperiodcentered{} ula/cifra.c \textperiodcentered{} fato \textperiodcentered{} confiança 1.0 \textperiodcentered{} domínio operacao\_so \textperiodcentered{} história 5 versões no jornal}

%CRISTAL {"arestas": [{"alvo": "arquitetura_von_neumann", "confianca": 1.0, "epistemico": "fato", "origem": "curriculo", "rel": "especializa"}, {"alvo": "word", "confianca": 0.5, "epistemico": "fato", "origem": "wiki", "rel": "relaciona"}], "confianca": 1.0, "contraexemplos": [], "descricao": "Portas AND, OR, NOT; álgebra de Boole materializa booleanos em hardware.", "epistemico": "fato", "exemplos": [], "id": "circuito_logico", "memoria": "semantica", "meta": {"dominio": "computacao", "nivel": 4, "ordem": 2}, "origem": "curriculo", "palavras_chave": ["booleano", "hardware"], "sinonimos": ["circuito lógico", "porta logica"], "tipo": "conceito", "titulo": "circuito logico"}
\section{circuito logico}
Portas AND, OR, NOT; álgebra de Boole materializa booleanos em hardware.
\prov{sinónimos: circuito lógico; porta logica}
\prov{relações: especializa arquitetura\_von\_neumann; relaciona word}
\prov{proveniência: curriculo \textperiodcentered{} fato \textperiodcentered{} confiança 1.0 \textperiodcentered{} domínio computacao \textperiodcentered{} história 7 versões no jornal}

%CRISTAL {"arestas": [{"alvo": "turing_halting_church", "confianca": 1.0, "epistemico": "fato", "origem": "wiki", "rel": "usa"}, {"alvo": "computacao", "confianca": 1.0, "epistemico": "fato", "origem": "wiki", "rel": "parte_de"}], "confianca": 1.0, "contraexemplos": [], "descricao": "P = solúveis em tempo polinomial; NP = solução verificável em tempo polinomial (P⊆NP⊆PSPACE, provado). NP-COMPLETUDE: SAT é NP-completo (Cook-Levin 1971) — todo problema em NP se reduz a ele em tempo polinomial; os 'mais difíceis' de NP (Karp: 21 problemas). BQP (quântico) ⊆ PSPACE. Definições + o teorema de Cook-Levin = fato.", "epistemico": "fato", "exemplos": [], "id": "classes_p_np_cook_levin", "memoria": "semantica", "meta": {"autor": "broca-so", "dominio": "computacao", "fonte": "Cook 1971; Levin 1973; Karp 1972"}, "origem": "wiki", "palavras_chave": [], "sinonimos": [], "tipo": "conceito", "titulo": "Classes P, NP e a NP-completude (Cook-Levin)"}
\section{Classes P, NP e a NP-completude (Cook-Levin)}
P = solúveis em tempo polinomial; NP = solução verificável em tempo polinomial (P\ensuremath{\subseteq}NP\ensuremath{\subseteq}PSPACE, provado). NP-COMPLETUDE: SAT é NP-completo (Cook-Levin 1971) — todo problema em NP se reduz a ele em tempo polinomial; os 'mais difíceis' de NP (Karp: 21 problemas). BQP (quântico) \ensuremath{\subseteq} PSPACE. Definições + o teorema de Cook-Levin = fato.
\prov{relações: usa turing\_halting\_church; parte\_de computacao}
\prov{proveniência: wiki \textperiodcentered{} Cook 1971; Levin 1973; Karp 1972 \textperiodcentered{} fato \textperiodcentered{} confiança 1.0 \textperiodcentered{} domínio computacao \textperiodcentered{} história 3 versões no jornal}

%CRISTAL {"arestas": [{"alvo": "lei_unica_tres_dominios", "confianca": 1.0, "epistemico": "fato", "origem": "wiki", "rel": "usa"}, {"alvo": "sandbox_cristal_ao_so", "confianca": 1.0, "epistemico": "fato", "origem": "wiki", "rel": "usa"}, {"alvo": "tiffany_roda_codigo", "confianca": 1.0, "epistemico": "fato", "origem": "wiki", "rel": "usa"}, {"alvo": "pipeline_completar", "confianca": 1.0, "epistemico": "fato", "origem": "wiki", "rel": "usa"}, {"alvo": "dicionarios_centralizados", "confianca": 1.0, "epistemico": "fato", "origem": "wiki", "rel": "usa"}, {"alvo": "tiffany", "confianca": 1.0, "epistemico": "fato", "origem": "wiki", "rel": "parte_de"}, {"alvo": "cristal_pipeline_runner", "confianca": 1.0, "epistemico": "fato", "origem": "wiki", "rel": "relaciona"}], "confianca": 1.0, "contraexemplos": [], "descricao": "O CLI cristal (code/cristal/), antes PAUSADO, foi despausado e modernizado em torno dos quatro verbos: 'cristal explica <pergunta>' (colapso→conceito, fail-closed se não sabe), 'cristal roda <código>' (sandbox: cristal→ISA→DSL, sela a classe da parábola), 'cristal traduz <lingua> <código>' (a bússola: qualquer língua→gatos→executa), 'cristal completa <pedido>' (aterra na menor energia→texto+código). 'cristal mapa' desenha a visão geral. E 'cristal falar' é o REPL moderno: linguagem natural → os quatro verbos com o FIO (um loop fino sobre colapso.responder — 'explica melhor'/'e daí?' seguem o contexto; roda:/completa:/traduz roteados sozinhos). Liga o terminal à máquina viva (conversa/ + linguagem/). Os ~50 comandos legados de scaffolding/agentes foram APAGADOS (ruído) e o COCKPIT acoplado ('cristal cockpit' sobe o Olho de Cock): o cristal ficou enxuto (50→15 arquivos) — só a lei única. bin/cristal roda por padrão (o REPL moderno sem arg; kernel legado atrás de CRISTAL_LEGACY).", "epistemico": "fato", "exemplos": [], "id": "cli_cristal_quatro_verbos", "memoria": "semantica", "meta": {"dominio": "programacao", "fonte": "code/cristal/commands/verbos.py; bin/cristal; code/cristal/main.py"}, "origem": "wiki", "palavras_chave": [], "sinonimos": [], "tipo": "conceito", "titulo": "CLI cristal — os quatro verbos da lei única"}
\section{CLI cristal — os quatro verbos da lei única}
O CLI cristal (code/cristal/), antes PAUSADO, foi despausado e modernizado em torno dos quatro verbos: 'cristal explica <pergunta>' (colapso\ensuremath{\to}conceito, fail-closed se não sabe), 'cristal roda <código>' (sandbox: cristal\ensuremath{\to}ISA\ensuremath{\to}DSL, sela a classe da parábola), 'cristal traduz <lingua> <código>' (a bússola: qualquer língua\ensuremath{\to}gatos\ensuremath{\to}executa), 'cristal completa <pedido>' (aterra na menor energia\ensuremath{\to}texto+código). 'cristal mapa' desenha a visão geral. E 'cristal falar' é o REPL moderno: linguagem natural \ensuremath{\to} os quatro verbos com o FIO (um loop fino sobre colapso.responder — 'explica melhor'/'e daí?' seguem o contexto; roda:/completa:/traduz roteados sozinhos). Liga o terminal à máquina viva (conversa/ + linguagem/). Os \~{}50 comandos legados de scaffolding/agentes foram APAGADOS (ruído) e o COCKPIT acoplado ('cristal cockpit' sobe o Olho de Cock): o cristal ficou enxuto (50\ensuremath{\to}15 arquivos) — só a lei única. bin/cristal roda por padrão (o REPL moderno sem arg; kernel legado atrás de CRISTAL\_LEGACY).
\prov{relações: usa lei\_unica\_tres\_dominios; usa sandbox\_cristal\_ao\_so; usa tiffany\_roda\_codigo; usa pipeline\_completar; usa dicionarios\_centralizados; parte\_de tiffany; relaciona cristal\_pipeline\_runner}
\prov{proveniência: wiki \textperiodcentered{} code/cristal/commands/verbos.py; bin/cristal; code/cristal/main.py \textperiodcentered{} fato \textperiodcentered{} confiança 1.0 \textperiodcentered{} domínio programacao \textperiodcentered{} história 32 versões no jornal}

%CRISTAL {"arestas": [{"alvo": "lisp", "confianca": 1.0, "epistemico": "fato", "origem": "wiki", "rel": "especializa"}, {"alvo": "imutabilidade", "confianca": 1.0, "epistemico": "fato", "origem": "wiki", "rel": "usa"}], "confianca": 1.0, "contraexemplos": [], "descricao": "Hickey: um Lisp moderno na JVM, com foco radical em imutabilidade e estruturas persistentes — concorrência sã.", "epistemico": "fato", "exemplos": [], "id": "clojure", "memoria": "semantica", "meta": {"autor": "Hickey", "dominio": "programacao", "fonte": ""}, "origem": "wiki", "palavras_chave": [], "sinonimos": [], "tipo": "conceito", "titulo": "Clojure"}
\section{Clojure}
Hickey: um Lisp moderno na JVM, com foco radical em imutabilidade e estruturas persistentes — concorrência sã.
\prov{relações: especializa lisp; usa imutabilidade}
\prov{proveniência: wiki \textperiodcentered{} fato \textperiodcentered{} confiança 1.0 \textperiodcentered{} domínio programacao \textperiodcentered{} história 6 versões no jornal}

%CRISTAL {"arestas": [{"alvo": "funcao_primeira_classe", "confianca": 1.0, "epistemico": "fato", "origem": "wiki", "rel": "usa"}], "confianca": 1.0, "contraexemplos": [], "descricao": "Uma função que captura o ambiente léxico onde nasceu — carrega consigo as variáveis que enxergava. Estado sem objeto.", "epistemico": "fato", "exemplos": [], "id": "closure", "memoria": "semantica", "meta": {"dominio": "programacao", "fonte": ""}, "origem": "wiki", "palavras_chave": [], "sinonimos": [], "tipo": "conceito", "titulo": "Closure"}
\section{Closure}
Uma função que captura o ambiente léxico onde nasceu — carrega consigo as variáveis que enxergava. Estado sem objeto.
\prov{relações: usa funcao\_primeira\_classe}
\prov{proveniência: wiki \textperiodcentered{} fato \textperiodcentered{} confiança 1.0 \textperiodcentered{} domínio programacao \textperiodcentered{} história 4 versões no jornal}

%CRISTAL {"arestas": [{"alvo": "goldchain", "confianca": 1.0, "epistemico": "fato", "origem": "wiki", "rel": "relaciona"}, {"alvo": "colapso_koch_tiffany", "confianca": 1.0, "epistemico": "fato", "origem": "wiki", "rel": "usa"}], "confianca": 1.0, "contraexemplos": [], "descricao": "Servidor HTTP local (127.0.0.1:8765) que serve o painel e uma API (/api/hodge, /api/porta/wait, /api/mempool, /api/olho-invertido): Hodge corre por dentro num worker; o 'olho invertido' traz o mempool (B_out) pela transformada de Koch z↦z² de volta a W_in. Fica FORA do runtime de produção (observatório).", "epistemico": "fato", "exemplos": [], "id": "cockpit_olho", "memoria": "semantica", "meta": {"autor": "broca-so", "dominio": "operacao_so", "fonte": "goldchain/cockpit_server.py"}, "origem": "wiki", "palavras_chave": [], "sinonimos": [], "tipo": "conceito", "titulo": "O Cockpit (Olho de Cock + olho invertido)"}
\section{O Cockpit (Olho de Cock + olho invertido)}
Servidor HTTP local (127.0.0.1:8765) que serve o painel e uma API (/api/hodge, /api/porta/wait, /api/mempool, /api/olho-invertido): Hodge corre por dentro num worker; o 'olho invertido' traz o mempool (B\_out) pela transformada de Koch z\ensuremath{\mapsto}z\ensuremath{^{2}} de volta a W\_in. Fica FORA do runtime de produção (observatório).
\prov{relações: relaciona goldchain; usa colapso\_koch\_tiffany}
\prov{proveniência: wiki \textperiodcentered{} goldchain/cockpit\_server.py \textperiodcentered{} fato \textperiodcentered{} confiança 1.0 \textperiodcentered{} domínio operacao\_so \textperiodcentered{} história 3 versões no jornal}

%CRISTAL {"arestas": [{"alvo": "tiffany", "confianca": 1.0, "epistemico": "fato", "origem": "codigo", "rel": "usa"}, {"alvo": "console_do_operador_lab", "confianca": 1.0, "epistemico": "fato", "origem": "wiki", "rel": "confirma"}], "confianca": 1.0, "contraexemplos": [], "descricao": "Laboratório de conhecimento — Fase 2. python3 conversa/conhecimento/lab.py ensinar reticulado --desc \"...\" python3 conversa/conhecimento/lab.py conceito reticulado python3 conversa/conhecimento/lab.py conectar gato mamifero --rel eh_um python3 conversa/conhecimento/lab.py sugerir reticulado python3 conversa/conhecimento/lab.py grafo reticulado python3 conversa/conhecimento/lab.py aprender \"O gato é um mamífero.\" python3 conversa/conhecimento/lab.py semilla python3 conversa/conhecimento/lab.py in", "epistemico": "fato", "exemplos": [], "id": "codigo_conhecimento_lab", "memoria": "semantica", "meta": {"arquivo": "broca-so/conversa/conhecimento/lab.py", "dominio": "codigo", "nivel": 6, "tipo_no": "codigo"}, "origem": "codigo", "palavras_chave": ["lab.py", "conhecimento"], "sinonimos": [], "tipo": "conceito", "titulo": "codigo conhecimento lab"}
\section{codigo conhecimento lab}
Laboratório de conhecimento — Fase 2. python3 conversa/conhecimento/lab.py ensinar reticulado --desc "..." python3 conversa/conhecimento/lab.py conceito reticulado python3 conversa/conhecimento/lab.py conectar gato mamifero --rel eh\_um python3 conversa/conhecimento/lab.py sugerir reticulado python3 conversa/conhecimento/lab.py grafo reticulado python3 conversa/conhecimento/lab.py aprender "O gato é um mamífero." python3 conversa/conhecimento/lab.py semilla python3 conversa/conhecimento/lab.py in
\prov{relações: usa tiffany; confirma console\_do\_operador\_lab}
\prov{proveniência: codigo \textperiodcentered{} fato \textperiodcentered{} confiança 1.0 \textperiodcentered{} domínio codigo \textperiodcentered{} história 5 versões no jornal}

%CRISTAL {"arestas": [{"alvo": "colapso", "confianca": 1.0, "epistemico": "fato", "origem": "codigo", "rel": "explica"}, {"alvo": "colapso_multicamada_peso", "confianca": 1.0, "epistemico": "fato", "origem": "wiki", "rel": "confirma"}], "confianca": 1.0, "contraexemplos": [], "descricao": "Colapso conversacional — broca-so cristalizado sobre tecidos locais. Prioridade: dialogos (curados) → corpus → machine (só --full). Dados no disco, gatos na RAM, contexto 32 B em contexto.bits.", "epistemico": "fato", "exemplos": [], "id": "codigo_conversa_colapso", "memoria": "semantica", "meta": {"arquivo": "broca-so/conversa/colapso.py", "dominio": "codigo", "nivel": 6, "tipo_no": "codigo"}, "origem": "codigo", "palavras_chave": ["colapso.py", "conversa"], "sinonimos": [], "tipo": "conceito", "titulo": "codigo conversa colapso"}
\section{codigo conversa colapso}
Colapso conversacional — broca-so cristalizado sobre tecidos locais. Prioridade: dialogos (curados) \ensuremath{\to} corpus \ensuremath{\to} machine (só --full). Dados no disco, gatos na RAM, contexto 32 B em contexto.bits.
\prov{relações: explica colapso; confirma colapso\_multicamada\_peso}
\prov{proveniência: codigo \textperiodcentered{} fato \textperiodcentered{} confiança 1.0 \textperiodcentered{} domínio codigo \textperiodcentered{} história 33 versões no jornal}

%CRISTAL {"arestas": [{"alvo": "valor", "confianca": 1.0, "epistemico": "fato", "origem": "codigo", "rel": "explica"}, {"alvo": "hamiltoniano_quarentena", "confianca": 1.0, "epistemico": "fato", "origem": "wiki", "rel": "confirma"}], "confianca": 1.0, "contraexemplos": [], "descricao": "Contabilidade GoldChain — Pell, parábola, BTC on-chain. Livro único: goldchain/dados/livro_contabilidade.json Ledger Pell: goldchain/dados/ledger_test.json python3 goldchain/contabilidade.py relatorio python3 goldchain/contabilidade.py relatorio --rede testnet python3 goldchain/contabilidade.py onde-esta-o-dinheiro", "epistemico": "fato", "exemplos": [], "id": "codigo_goldchain_contabilidade", "memoria": "semantica", "meta": {"arquivo": "broca-so/goldchain/contabilidade.py", "dominio": "codigo", "nivel": 6, "tipo_no": "codigo"}, "origem": "codigo", "palavras_chave": ["contabilidade.py", "goldchain"], "sinonimos": [], "tipo": "conceito", "titulo": "codigo goldchain contabilidade"}
\section{codigo goldchain contabilidade}
Contabilidade GoldChain — Pell, parábola, BTC on-chain. Livro único: goldchain/dados/livro\_contabilidade.json Ledger Pell: goldchain/dados/ledger\_test.json python3 goldchain/contabilidade.py relatorio python3 goldchain/contabilidade.py relatorio --rede testnet python3 goldchain/contabilidade.py onde-esta-o-dinheiro
\prov{relações: explica valor; confirma hamiltoniano\_quarentena}
\prov{proveniência: codigo \textperiodcentered{} fato \textperiodcentered{} confiança 1.0 \textperiodcentered{} domínio codigo \textperiodcentered{} história 5 versões no jornal}

%CRISTAL {"arestas": [{"alvo": "goldchain", "confianca": 1.0, "epistemico": "fato", "origem": "codigo", "rel": "define"}, {"alvo": "htlc_p2wsh_real", "confianca": 1.0, "epistemico": "fato", "origem": "wiki", "rel": "confirma"}, {"alvo": "misticismo_nao_dual", "confianca": 0.5, "epistemico": "fato", "origem": "wiki", "rel": "relaciona"}, {"alvo": "word", "confianca": 0.5, "epistemico": "fato", "origem": "wiki", "rel": "relaciona"}, {"alvo": "virtualizacao", "confianca": 0.5, "epistemico": "fato", "origem": "wiki", "rel": "relaciona"}], "confianca": 1.0, "contraexemplos": [], "descricao": "Contrato C = (P, φ, λ) — partes, condição, liquidação.", "epistemico": "fato", "exemplos": [], "id": "codigo_goldchain_contrato", "memoria": "semantica", "meta": {"arquivo": "broca-so/goldchain/contrato.py", "dominio": "codigo", "nivel": 6, "tipo_no": "codigo"}, "origem": "codigo", "palavras_chave": ["contrato.py", "goldchain"], "sinonimos": [], "tipo": "conceito", "titulo": "codigo goldchain contrato"}
\section{codigo goldchain contrato}
Contrato C = (P, \ensuremath{\varphi}, \ensuremath{\lambda}) — partes, condição, liquidação.
\prov{relações: define goldchain; confirma htlc\_p2wsh\_real; relaciona misticismo\_nao\_dual; relaciona word; relaciona virtualizacao}
\prov{proveniência: codigo \textperiodcentered{} fato \textperiodcentered{} confiança 1.0 \textperiodcentered{} domínio codigo \textperiodcentered{} história 8 versões no jornal}

%CRISTAL {"arestas": [{"alvo": "goldchain", "confianca": 1.0, "epistemico": "fato", "origem": "codigo", "rel": "explica"}, {"alvo": "fluxo_comissao_terceiros", "confianca": 1.0, "epistemico": "fato", "origem": "wiki", "rel": "confirma"}], "confianca": 1.0, "contraexemplos": [], "descricao": "Fluxo completo GoldChain — zero gas até cliente depositar. python3 goldchain/fluxo.py preparar goldchain/dados/contrato_comissao.json python3 goldchain/fluxo.py pronto goldchain/dados/contrato_comissao.json python3 goldchain/fluxo.py fechar goldchain/dados/contrato_comissao.json --yes", "epistemico": "fato", "exemplos": [], "id": "codigo_goldchain_fluxo", "memoria": "semantica", "meta": {"arquivo": "broca-so/goldchain/fluxo.py", "dominio": "codigo", "nivel": 6, "tipo_no": "codigo"}, "origem": "codigo", "palavras_chave": ["fluxo.py", "goldchain"], "sinonimos": [], "tipo": "conceito", "titulo": "codigo goldchain fluxo"}
\section{codigo goldchain fluxo}
Fluxo completo GoldChain — zero gas até cliente depositar. python3 goldchain/fluxo.py preparar goldchain/dados/contrato\_comissao.json python3 goldchain/fluxo.py pronto goldchain/dados/contrato\_comissao.json python3 goldchain/fluxo.py fechar goldchain/dados/contrato\_comissao.json --yes
\prov{relações: explica goldchain; confirma fluxo\_comissao\_terceiros}
\prov{proveniência: codigo \textperiodcentered{} fato \textperiodcentered{} confiança 1.0 \textperiodcentered{} domínio codigo \textperiodcentered{} história 33 versões no jornal}

%CRISTAL {"arestas": [{"alvo": "trabalho", "confianca": 1.0, "epistemico": "fato", "origem": "codigo", "rel": "explica"}, {"alvo": "ciclo_uno", "confianca": 1.0, "epistemico": "fato", "origem": "wiki", "rel": "confirma"}], "confianca": 1.0, "contraexemplos": [], "descricao": "CLI de mineração (teste) — depois liga mineração real (Sepolia/BTC). python3 goldchain/minerar.py teste python3 goldchain/minerar.py trabalho \"alice-goldchain-2026\" --buscar python3 goldchain/minerar.py status python3 goldchain/minerar.py demo-contrato goldchain/dados/exemplo_partes.json", "epistemico": "fato", "exemplos": [], "id": "codigo_goldchain_minerar", "memoria": "semantica", "meta": {"arquivo": "broca-so/goldchain/minerar.py", "dominio": "codigo", "nivel": 6, "tipo_no": "codigo"}, "origem": "codigo", "palavras_chave": ["minerar.py", "goldchain"], "sinonimos": [], "tipo": "conceito", "titulo": "codigo goldchain minerar"}
\section{codigo goldchain minerar}
CLI de mineração (teste) — depois liga mineração real (Sepolia/BTC). python3 goldchain/minerar.py teste python3 goldchain/minerar.py trabalho "alice-goldchain-2026" --buscar python3 goldchain/minerar.py status python3 goldchain/minerar.py demo-contrato goldchain/dados/exemplo\_partes.json
\prov{relações: explica trabalho; confirma ciclo\_uno}
\prov{proveniência: codigo \textperiodcentered{} fato \textperiodcentered{} confiança 1.0 \textperiodcentered{} domínio codigo \textperiodcentered{} história 5 versões no jornal}

%CRISTAL {"arestas": [{"alvo": "cifra_de_ouro_branco", "confianca": 1.0, "epistemico": "fato", "origem": "codigo", "rel": "explica"}, {"alvo": "cotacao_parabola", "confianca": 1.0, "epistemico": "fato", "origem": "wiki", "rel": "confirma"}, {"alvo": "word", "confianca": 0.5, "epistemico": "fato", "origem": "wiki", "rel": "relaciona"}, {"alvo": "octonios_hiperbolicos", "confianca": 0.5, "epistemico": "fato", "origem": "wiki", "rel": "relaciona"}, {"alvo": "virtualizacao", "confianca": 0.5, "epistemico": "fato", "origem": "wiki", "rel": "relaciona"}], "confianca": 1.0, "contraexemplos": [], "descricao": "Números de Pell — reserva GoldChain (metal prata, n=2).", "epistemico": "fato", "exemplos": [], "id": "codigo_goldchain_pell", "memoria": "semantica", "meta": {"arquivo": "broca-so/goldchain/pell.py", "dominio": "codigo", "nivel": 6, "tipo_no": "codigo"}, "origem": "codigo", "palavras_chave": ["pell.py", "goldchain"], "sinonimos": [], "tipo": "conceito", "titulo": "codigo goldchain pell"}
\section{codigo goldchain pell}
Números de Pell — reserva GoldChain (metal prata, n=2).
\prov{relações: explica cifra\_de\_ouro\_branco; confirma cotacao\_parabola; relaciona word; relaciona octonios\_hiperbolicos; relaciona virtualizacao}
\prov{proveniência: codigo \textperiodcentered{} fato \textperiodcentered{} confiança 1.0 \textperiodcentered{} domínio codigo \textperiodcentered{} história 8 versões no jornal}

%CRISTAL {"arestas": [{"alvo": "sha256", "confianca": 1.0, "epistemico": "fato", "origem": "codigo", "rel": "confirma"}], "confianca": 1.0, "contraexemplos": [], "descricao": "Prova de trabalho — estrela de prata (6 modos SVD) ou penta ouro (5 modos). Espectro neurônio (12⁴) → matriz lado×lado → SVD → conta modos significativos.", "epistemico": "fato", "exemplos": [], "id": "codigo_goldchain_prova", "memoria": "semantica", "meta": {"arquivo": "broca-so/goldchain/prova.py", "dominio": "codigo", "nivel": 6, "tipo_no": "codigo"}, "origem": "codigo", "palavras_chave": ["prova.py", "goldchain"], "sinonimos": [], "tipo": "conceito", "titulo": "codigo goldchain prova"}
\section{codigo goldchain prova}
Prova de trabalho — estrela de prata (6 modos SVD) ou penta ouro (5 modos). Espectro neurônio (12\ensuremath{^{4}}) \ensuremath{\to} matriz lado$\times$lado \ensuremath{\to} SVD \ensuremath{\to} conta modos significativos.
\prov{relações: confirma sha256}
\prov{proveniência: codigo \textperiodcentered{} fato \textperiodcentered{} confiança 1.0 \textperiodcentered{} domínio codigo \textperiodcentered{} história 4 versões no jornal}

%CRISTAL {"arestas": [{"alvo": "antena", "confianca": 1.0, "epistemico": "fato", "origem": "codigo", "rel": "usa"}], "confianca": 1.0, "contraexemplos": [], "descricao": "Ponte rádio ↔ GoldChain. Fluxo: partes (DNA) → células no rádio → validar espectro → banda colectiva → rede N-qubit (GHZ) → segredo/hashlock HTLC → testemunha envia segredo na parte negra → lock/receiver decodificam python3 goldchain/radio_ponte.py goldchain/dados/contrato_comissao.json python3 goldchain/radio_ponte.py goldchain/dados/contrato_comissao.json --json", "epistemico": "fato", "exemplos": [], "id": "codigo_goldchain_radio_ponte", "memoria": "semantica", "meta": {"arquivo": "broca-so/goldchain/radio_ponte.py", "dominio": "codigo", "nivel": 6, "tipo_no": "codigo"}, "origem": "codigo", "palavras_chave": ["radio_ponte.py", "goldchain"], "sinonimos": [], "tipo": "conceito", "titulo": "codigo goldchain radio_ponte"}
\section{codigo goldchain radio\_ponte}
Ponte rádio \ensuremath{\leftrightarrow} GoldChain. Fluxo: partes (DNA) \ensuremath{\to} células no rádio \ensuremath{\to} validar espectro \ensuremath{\to} banda colectiva \ensuremath{\to} rede N-qubit (GHZ) \ensuremath{\to} segredo/hashlock HTLC \ensuremath{\to} testemunha envia segredo na parte negra \ensuremath{\to} lock/receiver decodificam python3 goldchain/radio\_ponte.py goldchain/dados/contrato\_comissao.json python3 goldchain/radio\_ponte.py goldchain/dados/contrato\_comissao.json --json
\prov{relações: usa antena}
\prov{proveniência: codigo \textperiodcentered{} fato \textperiodcentered{} confiança 1.0 \textperiodcentered{} domínio codigo \textperiodcentered{} história 4 versões no jornal}

%CRISTAL {"arestas": [{"alvo": "hopfield", "confianca": 1.0, "epistemico": "fato", "origem": "codigo", "rel": "confirma"}, {"alvo": "reproduzir_experimentos", "confianca": 1.0, "epistemico": "fato", "origem": "wiki", "rel": "confirma"}], "confianca": 1.0, "contraexemplos": [], "descricao": "ISOMORFISMO — a meta-indução do Broca OS (gerada pelo gato A=[[1,1],[1,0]]) É a memória associativa de Hopfield: MESMO sistema dinâmico com energia. Validação SIMBÓLICA EXATA (sympy), amarrada nas peças REAIS (estelar.gap_matriz, estelar.torcao_gap, vida_morte.maquina). 5 passos: (1) A é matriz sináptica de Hopfield (Aᵀ=A); (2) energia de Lyapunov E=−½xᵀAx; (3) espectro → atrator (A=φP_φ+ψP_ψ, projetores ortogonais idempotentes); (4) pattern completion (Aⁿ/φⁿ→P_φ, taxa |ψ/φ|); (5) resiliência/en", "epistemico": "fato", "exemplos": [], "id": "codigo_lab_prova_isomorfismo", "memoria": "semantica", "meta": {"arquivo": "machine/.lab/prova_isomorfismo.py", "dominio": "codigo", "nivel": 6, "tipo_no": "codigo"}, "origem": "codigo", "palavras_chave": ["prova_isomorfismo.py", ".lab"], "sinonimos": [], "tipo": "conceito", "titulo": "codigo .lab prova_isomorfismo"}
\section{codigo .lab prova\_isomorfismo}
ISOMORFISMO — a meta-indução do Broca OS (gerada pelo gato A=[[1,1],[1,0]]) É a memória associativa de Hopfield: MESMO sistema dinâmico com energia. Validação SIMBÓLICA EXATA (sympy), amarrada nas peças REAIS (estelar.gap\_matriz, estelar.torcao\_gap, vida\_morte.maquina). 5 passos: (1) A é matriz sináptica de Hopfield (A\ensuremath{^{T}}=A); (2) energia de Lyapunov E=\ensuremath{-}\ensuremath{\tfrac12}x\ensuremath{^{T}}Ax; (3) espectro \ensuremath{\to} atrator (A=\ensuremath{\varphi}P\_\ensuremath{\varphi}+\ensuremath{\psi}P\_\ensuremath{\psi}, projetores ortogonais idempotentes); (4) pattern completion (A\ensuremath{^{n}}/\ensuremath{\varphi}\ensuremath{^{n}}\ensuremath{\to}P\_\ensuremath{\varphi}, taxa |\ensuremath{\psi}/\ensuremath{\varphi}|); (5) resiliência/en
\prov{relações: confirma hopfield; confirma reproduzir\_experimentos}
\prov{proveniência: codigo \textperiodcentered{} fato \textperiodcentered{} confiança 1.0 \textperiodcentered{} domínio codigo \textperiodcentered{} história 32 versões no jornal}

%CRISTAL {"arestas": [{"alvo": "antena", "confianca": 1.0, "epistemico": "fato", "origem": "codigo", "rel": "confirma"}, {"alvo": "reproduzir_experimentos", "confianca": 1.0, "epistemico": "fato", "origem": "wiki", "rel": "confirma"}, {"alvo": "antena_canal_ipc_bump", "confianca": 1.0, "epistemico": "fato", "origem": "wiki", "rel": "confirma"}], "confianca": 1.0, "contraexemplos": [], "descricao": "Antena canal — bump por pacote (.pkt via ssh). LEGADO: comandos olho, não stream stratum. Stream contínuo GEX→local: olho_gex44 + floco (enlace_fisico) + reguas_esteris. Banda: banda_de_dna — NÃO sha256(cid) manual.", "epistemico": "fato", "exemplos": [], "id": "codigo_neuronio_antena_canal", "memoria": "semantica", "meta": {"arquivo": "broca-so/neuronio/antena_canal.py", "dominio": "codigo", "nivel": 6, "tipo_no": "codigo"}, "origem": "codigo", "palavras_chave": ["antena_canal.py", "neuronio"], "sinonimos": [], "tipo": "conceito", "titulo": "codigo neuronio antena_canal"}
\section{codigo neuronio antena\_canal}
Antena canal — bump por pacote (.pkt via ssh). LEGADO: comandos olho, não stream stratum. Stream contínuo GEX\ensuremath{\to}local: olho\_gex44 + floco (enlace\_fisico) + reguas\_esteris. Banda: banda\_de\_dna — NÃO sha256(cid) manual.
\prov{relações: confirma antena; confirma reproduzir\_experimentos; confirma antena\_canal\_ipc\_bump}
\prov{proveniência: codigo \textperiodcentered{} fato \textperiodcentered{} confiança 1.0 \textperiodcentered{} domínio codigo \textperiodcentered{} história 6 versões no jornal}

%CRISTAL {"arestas": [{"alvo": "tiffany", "confianca": 1.0, "epistemico": "fato", "origem": "codigo", "rel": "confirma"}, {"alvo": "reproduzir_experimentos", "confianca": 1.0, "epistemico": "fato", "origem": "wiki", "rel": "confirma"}], "confianca": 1.0, "contraexemplos": [], "descricao": "Experimentos Tiffany — reprodução unificada de todas as verificações. Tiffany (lab.py reproduzir) invoca este módulo para validar papers ↔ código. python3 neuronio/experimentos_tiffany.py python3 neuronio/experimentos_tiffany.py --json python3 neuronio/experimentos_tiffany.py --suite ponte python3 conversa/conhecimento/lab.py reproduzir", "epistemico": "fato", "exemplos": [], "id": "codigo_neuronio_experimentos_tiffany", "memoria": "semantica", "meta": {"arquivo": "broca-so/neuronio/experimentos_tiffany.py", "dominio": "codigo", "nivel": 6, "tipo_no": "codigo"}, "origem": "codigo", "palavras_chave": ["experimentos_tiffany.py", "neuronio"], "sinonimos": [], "tipo": "conceito", "titulo": "codigo neuronio experimentos_tiffany"}
\section{codigo neuronio experimentos\_tiffany}
Experimentos Tiffany — reprodução unificada de todas as verificações. Tiffany (lab.py reproduzir) invoca este módulo para validar papers \ensuremath{\leftrightarrow} código. python3 neuronio/experimentos\_tiffany.py python3 neuronio/experimentos\_tiffany.py --json python3 neuronio/experimentos\_tiffany.py --suite ponte python3 conversa/conhecimento/lab.py reproduzir
\prov{relações: confirma tiffany; confirma reproduzir\_experimentos}
\prov{proveniência: codigo \textperiodcentered{} fato \textperiodcentered{} confiança 1.0 \textperiodcentered{} domínio codigo \textperiodcentered{} história 33 versões no jornal}

%CRISTAL {"arestas": [{"alvo": "ponte_simbolica_gentil", "confianca": 1.0, "epistemico": "fato", "origem": "codigo", "rel": "confirma"}, {"alvo": "reproduzir_experimentos", "confianca": 1.0, "epistemico": "fato", "origem": "wiki", "rel": "confirma"}], "confianca": 1.0, "contraexemplos": [], "descricao": "Computação simbólica — álgebra de Lopes (chart A) + ligação operacional (chart B). Formaliza em símbolos exactos (Fraction) o que gentil.c faz numericamente: • Lopes2: Q = r²+c², produto complexo, Q multiplicativa • Balança a+c²=1 • Sequências de Cauchy em ℚ, equivalência ~, classes [(r_n)] • Homomorfismos parciais do paper ponte (Fibonacci, √2 Lopes, truncamentos) python3 neuronio/gentil_simbolico.py python3 neuronio/gentil_simbolico.py --json", "epistemico": "fato", "exemplos": [], "id": "codigo_neuronio_gentil_simbolico", "memoria": "semantica", "meta": {"arquivo": "broca-so/neuronio/gentil_simbolico.py", "dominio": "codigo", "nivel": 6, "tipo_no": "codigo"}, "origem": "codigo", "palavras_chave": ["gentil_simbolico.py", "neuronio"], "sinonimos": [], "tipo": "conceito", "titulo": "codigo neuronio gentil_simbolico"}
\section{codigo neuronio gentil\_simbolico}
Computação simbólica — álgebra de Lopes (chart A) + ligação operacional (chart B). Formaliza em símbolos exactos (Fraction) o que gentil.c faz numericamente: \textbullet  Lopes2: Q = r\ensuremath{^{2}}+c\ensuremath{^{2}}, produto complexo, Q multiplicativa \textbullet  Balança a+c\ensuremath{^{2}}=1 \textbullet  Sequências de Cauchy em \ensuremath{\mathbb{Q}}, equivalência \~{}, classes [(r\_n)] \textbullet  Homomorfismos parciais do paper ponte (Fibonacci, \ensuremath{\sqrt{\,}}2 Lopes, truncamentos) python3 neuronio/gentil\_simbolico.py python3 neuronio/gentil\_simbolico.py --json
\prov{relações: confirma ponte\_simbolica\_gentil; confirma reproduzir\_experimentos}
\prov{proveniência: codigo \textperiodcentered{} fato \textperiodcentered{} confiança 1.0 \textperiodcentered{} domínio codigo \textperiodcentered{} história 6 versões no jornal}

%CRISTAL {"arestas": [{"alvo": "gato", "confianca": 1.0, "epistemico": "fato", "origem": "codigo", "rel": "usa"}, {"alvo": "mede_flocos_phi", "confianca": 1.0, "epistemico": "fato", "origem": "wiki", "rel": "confirma"}, {"alvo": "pacote_seed_pow", "confianca": 1.0, "epistemico": "fato", "origem": "wiki", "rel": "confirma"}], "confianca": 1.0, "contraexemplos": [], "descricao": "BASELINE ASIC — referência / benchmark. NÃO lucro. NÃO confundir com gato nativo. Mede taxa SHA256² (workers Python) vs ficha ASIC. N_eq = mesma taxa H/s, não bits na RAM. Gato NATIVO: python3 neuronio/pacotes.py nativo (pow_geometrico — 0 workers) python3 neuronio/pacotes.py baseline python3 neuronio/mina_gato.py --gatos 8 --seg 10", "epistemico": "fato", "exemplos": [], "id": "codigo_neuronio_mina_gato", "memoria": "semantica", "meta": {"arquivo": "broca-so/neuronio/mina_gato.py", "dominio": "codigo", "nivel": 6, "tipo_no": "codigo"}, "origem": "codigo", "palavras_chave": ["mina_gato.py", "neuronio"], "sinonimos": [], "tipo": "conceito", "titulo": "codigo neuronio mina_gato"}
\section{codigo neuronio mina\_gato}
BASELINE ASIC — referência / benchmark. NÃO lucro. NÃO confundir com gato nativo. Mede taxa SHA256\ensuremath{^{2}} (workers Python) vs ficha ASIC. N\_eq = mesma taxa H/s, não bits na RAM. Gato NATIVO: python3 neuronio/pacotes.py nativo (pow\_geometrico — 0 workers) python3 neuronio/pacotes.py baseline python3 neuronio/mina\_gato.py --gatos 8 --seg 10
\prov{relações: usa gato; confirma mede\_flocos\_phi; confirma pacote\_seed\_pow}
\prov{proveniência: codigo \textperiodcentered{} fato \textperiodcentered{} confiança 1.0 \textperiodcentered{} domínio codigo \textperiodcentered{} história 6 versões no jornal}

%CRISTAL {"arestas": [{"alvo": "broca_os", "confianca": 1.0, "epistemico": "fato", "origem": "codigo", "rel": "explica"}, {"alvo": "mede_flocos_phi", "confianca": 1.0, "epistemico": "fato", "origem": "wiki", "rel": "confirma"}, {"alvo": "primitivas_nucleo_associativo", "confianca": 1.0, "epistemico": "fato", "origem": "wiki", "rel": "confirma"}], "confianca": 1.0, "contraexemplos": [], "descricao": "Primitivas broca-so — equações dos papers, não primitivas de computação clássica. Pipeline correcto (BROCA_CRISTAL · quatro_ciencias): stratum (borda administrativa — olho I/O, ponte pontual) → gato (invariante topológico A, Word [total,e]) → álgebra → termodinâmica → geometria → mecânica (núcleo / olho Koch) → isolamento do floco Φ (Koch, recipiente — antes do backend) → backend broca-so (só parte associativa: gato_stream cpu|gpu) O .so NÃO reabre canal parabólico nem origina SHA256² — transpor", "epistemico": "fato", "exemplos": [], "id": "codigo_neuronio_primitivas_broca", "memoria": "semantica", "meta": {"arquivo": "broca-so/neuronio/primitivas_broca.py", "dominio": "codigo", "nivel": 6, "tipo_no": "codigo"}, "origem": "codigo", "palavras_chave": ["primitivas_broca.py", "neuronio"], "sinonimos": [], "tipo": "conceito", "titulo": "codigo neuronio primitivas_broca"}
\section{codigo neuronio primitivas\_broca}
Primitivas broca-so — equações dos papers, não primitivas de computação clássica. Pipeline correcto (BROCA\_CRISTAL \textperiodcentered  quatro\_ciencias): stratum (borda administrativa — olho I/O, ponte pontual) \ensuremath{\to} gato (invariante topológico A, Word [total,e]) \ensuremath{\to} álgebra \ensuremath{\to} termodinâmica \ensuremath{\to} geometria \ensuremath{\to} mecânica (núcleo / olho Koch) \ensuremath{\to} isolamento do floco \ensuremath{\Phi} (Koch, recipiente — antes do backend) \ensuremath{\to} backend broca-so (só parte associativa: gato\_stream cpu|gpu) O .so NÃO reabre canal parabólico nem origina SHA256\ensuremath{^{2}} — transpor
\prov{relações: explica broca\_os; confirma mede\_flocos\_phi; confirma primitivas\_nucleo\_associativo}
\prov{proveniência: codigo \textperiodcentered{} fato \textperiodcentered{} confiança 1.0 \textperiodcentered{} domínio codigo \textperiodcentered{} história 6 versões no jornal}

%CRISTAL {"arestas": [{"alvo": "ipc", "confianca": 1.0, "epistemico": "fato", "origem": "codigo", "rel": "explica"}, {"alvo": "reproduzir_experimentos", "confianca": 1.0, "epistemico": "fato", "origem": "wiki", "rel": "confirma"}, {"alvo": "antena_canal_ipc_bump", "confianca": 1.0, "epistemico": "fato", "origem": "wiki", "rel": "confirma"}], "confianca": 1.0, "contraexemplos": [], "descricao": "Rádio broca-so — parte negra, antenas, bandas, clones, floco Φ. broca_os.tex / antena.tex v2: IPC = canal E_ij (parte negra), demux e_i M e_j banda = banda_de_dna — parabola(w,gold(w))→Koch; quem casa decodifica bump = msg ⊕ K(ref) — ping-pong rápido (conservar=False) fluxo = frames 17 B/bit — difusão Φ‖bump (conservar=True, enlace_difusivo) floco Φ = selo resumo 16 B da célula emissora enlace físico (cobre) = olho_gex44 / floco — fora da parte negra clone = La Hire — filho de corpos distintos (", "epistemico": "fato", "exemplos": [], "id": "codigo_neuronio_radio", "memoria": "semantica", "meta": {"arquivo": "broca-so/neuronio/radio.py", "dominio": "codigo", "nivel": 6, "tipo_no": "codigo"}, "origem": "codigo", "palavras_chave": ["radio.py", "neuronio"], "sinonimos": [], "tipo": "conceito", "titulo": "codigo neuronio radio"}
\section{codigo neuronio radio}
Rádio broca-so — parte negra, antenas, bandas, clones, floco \ensuremath{\Phi}. broca\_os.tex / antena.tex v2: IPC = canal E\_ij (parte negra), demux e\_i M e\_j banda = banda\_de\_dna — parabola(w,gold(w))\ensuremath{\to}Koch; quem casa decodifica bump = msg \ensuremath{\oplus} K(ref) — ping-pong rápido (conservar=False) fluxo = frames 17 B/bit — difusão \ensuremath{\Phi}\ensuremath{\|}bump (conservar=True, enlace\_difusivo) floco \ensuremath{\Phi} = selo resumo 16 B da célula emissora enlace físico (cobre) = olho\_gex44 / floco — fora da parte negra clone = La Hire — filho de corpos distintos (
\prov{relações: explica ipc; confirma reproduzir\_experimentos; confirma antena\_canal\_ipc\_bump}
\prov{proveniência: codigo \textperiodcentered{} fato \textperiodcentered{} confiança 1.0 \textperiodcentered{} domínio codigo \textperiodcentered{} história 6 versões no jornal}

%CRISTAL {"arestas": [{"alvo": "corpo_dual", "confianca": 1.0, "epistemico": "fato", "origem": "codigo", "rel": "confirma"}, {"alvo": "reproduzir_experimentos", "confianca": 1.0, "epistemico": "fato", "origem": "wiki", "rel": "confirma"}], "confianca": 1.0, "contraexemplos": [], "descricao": "Verificação — O Corpo Dual (unificação tropical↔glacial, vazio/completude). python3 neuronio/verificar_corpo_dual.py python3 neuronio/verificar_corpo_dual.py --json", "epistemico": "fato", "exemplos": [], "id": "codigo_neuronio_verificar_corpo_dual", "memoria": "semantica", "meta": {"arquivo": "broca-so/neuronio/verificar_corpo_dual.py", "dominio": "codigo", "nivel": 6, "tipo_no": "codigo"}, "origem": "codigo", "palavras_chave": ["verificar_corpo_dual.py", "neuronio"], "sinonimos": [], "tipo": "conceito", "titulo": "codigo neuronio verificar_corpo_dual"}
\section{codigo neuronio verificar\_corpo\_dual}
Verificação — O Corpo Dual (unificação tropical\ensuremath{\leftrightarrow}glacial, vazio/completude). python3 neuronio/verificar\_corpo\_dual.py python3 neuronio/verificar\_corpo\_dual.py --json
\prov{relações: confirma corpo\_dual; confirma reproduzir\_experimentos}
\prov{proveniência: codigo \textperiodcentered{} fato \textperiodcentered{} confiança 1.0 \textperiodcentered{} domínio codigo \textperiodcentered{} história 5 versões no jornal}

%CRISTAL {"arestas": [{"alvo": "corpo_glacial", "confianca": 1.0, "epistemico": "fato", "origem": "codigo", "rel": "confirma"}, {"alvo": "reproduzir_experimentos", "confianca": 1.0, "epistemico": "fato", "origem": "wiki", "rel": "confirma"}], "confianca": 1.0, "contraexemplos": [], "descricao": "Verificação — O Corpo Glacial (máquina negra, clones, Peirce). python3 neuronio/verificar_corpo_glacial.py python3 neuronio/verificar_corpo_glacial.py --json", "epistemico": "fato", "exemplos": [], "id": "codigo_neuronio_verificar_corpo_glacial", "memoria": "semantica", "meta": {"arquivo": "broca-so/neuronio/verificar_corpo_glacial.py", "dominio": "codigo", "nivel": 6, "tipo_no": "codigo"}, "origem": "codigo", "palavras_chave": ["verificar_corpo_glacial.py", "neuronio"], "sinonimos": [], "tipo": "conceito", "titulo": "codigo neuronio verificar_corpo_glacial"}
\section{codigo neuronio verificar\_corpo\_glacial}
Verificação — O Corpo Glacial (máquina negra, clones, Peirce). python3 neuronio/verificar\_corpo\_glacial.py python3 neuronio/verificar\_corpo\_glacial.py --json
\prov{relações: confirma corpo\_glacial; confirma reproduzir\_experimentos}
\prov{proveniência: codigo \textperiodcentered{} fato \textperiodcentered{} confiança 1.0 \textperiodcentered{} domínio codigo \textperiodcentered{} história 5 versões no jornal}

%CRISTAL {"arestas": [{"alvo": "gato", "confianca": 1.0, "epistemico": "fato", "origem": "codigo", "rel": "confirma"}, {"alvo": "reproduzir_experimentos", "confianca": 1.0, "epistemico": "fato", "origem": "wiki", "rel": "confirma"}], "confianca": 1.0, "contraexemplos": [], "descricao": "Verificação: gato A=[[1,1],[1,0]] (neuronio.c) → símbolos Lopes. Prova numérica (sem estatística) que o operador streaming do neurônio é a mesma matriz 2×2 do cap. simbólico Lopes / Corpo Tropical (Anosov, x²=x+1, φ, P₀+P₁=I). python3 neuronio/verificar_ponte_gentil.py python3 neuronio/verificar_ponte_gentil.py --json", "epistemico": "fato", "exemplos": [], "id": "codigo_neuronio_verificar_ponte_gentil", "memoria": "semantica", "meta": {"arquivo": "broca-so/neuronio/verificar_ponte_gentil.py", "dominio": "codigo", "nivel": 6, "tipo_no": "codigo"}, "origem": "codigo", "palavras_chave": ["verificar_ponte_gentil.py", "neuronio"], "sinonimos": [], "tipo": "conceito", "titulo": "codigo neuronio verificar_ponte_gentil"}
\section{codigo neuronio verificar\_ponte\_gentil}
Verificação: gato A=[[1,1],[1,0]] (neuronio.c) \ensuremath{\to} símbolos Lopes. Prova numérica (sem estatística) que o operador streaming do neurônio é a mesma matriz 2$\times$2 do cap. simbólico Lopes / Corpo Tropical (Anosov, x\ensuremath{^{2}}=x+1, \ensuremath{\varphi}, P\ensuremath{_{0}}+P\ensuremath{_{1}}=I). python3 neuronio/verificar\_ponte\_gentil.py python3 neuronio/verificar\_ponte\_gentil.py --json
\prov{relações: confirma gato; confirma reproduzir\_experimentos}
\prov{proveniência: codigo \textperiodcentered{} fato \textperiodcentered{} confiança 1.0 \textperiodcentered{} domínio codigo \textperiodcentered{} história 6 versões no jornal}

%CRISTAL {"arestas": [{"alvo": "goldchain", "confianca": 1.0, "epistemico": "fato", "origem": "codigo", "rel": "confirma"}, {"alvo": "ciclo_uno", "confianca": 1.0, "epistemico": "fato", "origem": "wiki", "rel": "confirma"}], "confianca": 1.0, "contraexemplos": [], "descricao": "goldchain.py — A GoldChain, versão mínima (o roteiro de algebra_dos_contratos.tex, posto a rodar). A interface Contrato(P,φ,λ) com back-ends de liquidação plugáveis (ETH, BTC, demo). O PAINEL faz as três operações universais — VALIDAR (N-qubits), COORDENAR (hashlock), COBRAR (ganha 1 número de Pell da reserva de 13M em GEX) — e DELEGA a liquidação ao back-end. Cada contrato fechado é certificado e ENCADEADO (o registro da GoldChain). A reserva de Pell não é de ninguém: ganha-se fechando contrato", "epistemico": "fato", "exemplos": [], "id": "codigo_so_goldchain", "memoria": "semantica", "meta": {"arquivo": "machine/so/goldchain.py", "dominio": "codigo", "nivel": 6, "tipo_no": "codigo"}, "origem": "codigo", "palavras_chave": ["goldchain.py", "so"], "sinonimos": [], "tipo": "conceito", "titulo": "codigo so goldchain"}
\section{codigo so goldchain}
goldchain.py — A GoldChain, versão mínima (o roteiro de algebra\_dos\_contratos.tex, posto a rodar). A interface Contrato(P,\ensuremath{\varphi},\ensuremath{\lambda}) com back-ends de liquidação plugáveis (ETH, BTC, demo). O PAINEL faz as três operações universais — VALIDAR (N-qubits), COORDENAR (hashlock), COBRAR (ganha 1 número de Pell da reserva de 13M em GEX) — e DELEGA a liquidação ao back-end. Cada contrato fechado é certificado e ENCADEADO (o registro da GoldChain). A reserva de Pell não é de ninguém: ganha-se fechando contrato
\prov{relações: confirma goldchain; confirma ciclo\_uno}
\prov{proveniência: codigo \textperiodcentered{} fato \textperiodcentered{} confiança 1.0 \textperiodcentered{} domínio codigo \textperiodcentered{} história 32 versões no jornal}

%CRISTAL {"arestas": [{"alvo": "capacidade_canal_shannon", "confianca": 1.0, "epistemico": "fato", "origem": "wiki", "rel": "usa"}], "confianca": 1.0, "contraexemplos": [], "descricao": "Codificação com redundância estruturada que detecta/corrige erros, aproximando-se do limite de Shannon: Hamming (1950, corrige 1 erro), Reed-Solomon (1960), LDPC (Gallager 1962) e Turbo (Berrou 1993) chegam a frações de dB do limite. Fatos construtivos (existência e desempenho provados). Compressão prática: Huffman (ótimo símbolo-a-símbolo), aritmética, Lempel-Ziv (universal, converge a H).", "epistemico": "fato", "exemplos": [], "id": "codigos_corretores_erro", "memoria": "semantica", "meta": {"autor": "broca-so", "dominio": "computacao", "fonte": "Hamming 1950; Reed-Solomon 1960; Gallager 1962; Ziv-Lempel 1977"}, "origem": "wiki", "palavras_chave": [], "sinonimos": [], "tipo": "conceito", "titulo": "Códigos corretores de erro (rumo ao limite de Shannon)"}
\section{Códigos corretores de erro (rumo ao limite de Shannon)}
Codificação com redundância estruturada que detecta/corrige erros, aproximando-se do limite de Shannon: Hamming (1950, corrige 1 erro), Reed-Solomon (1960), LDPC (Gallager 1962) e Turbo (Berrou 1993) chegam a frações de dB do limite. Fatos construtivos (existência e desempenho provados). Compressão prática: Huffman (ótimo símbolo-a-símbolo), aritmética, Lempel-Ziv (universal, converge a H).
\prov{relações: usa capacidade\_canal\_shannon}
\prov{proveniência: wiki \textperiodcentered{} Hamming 1950; Reed-Solomon 1960; Gallager 1962; Ziv-Lempel 1977 \textperiodcentered{} fato \textperiodcentered{} confiança 1.0 \textperiodcentered{} domínio computacao \textperiodcentered{} história 2 versões no jornal}

%CRISTAL {"arestas": [{"alvo": "colapso_koch_tiffany", "confianca": 1.0, "epistemico": "fato", "origem": "wiki", "rel": "explica"}, {"alvo": "colapso_ping_pong_overlap", "confianca": 1.0, "epistemico": "fato", "origem": "wiki", "rel": "usa"}, {"alvo": "bai_deny_overrides", "confianca": 1.0, "epistemico": "fato", "origem": "wiki", "rel": "usa"}, {"alvo": "operacao_broca_so", "confianca": 1.0, "epistemico": "fato", "origem": "wiki", "rel": "parte_de"}], "confianca": 1.0, "contraexemplos": [], "descricao": "conversa/colapso.py varre as camadas dialogos(8)→conhecimento(6)→corpus(4), pontua cada candidato como overlap Koch + peso da camada + relevância + regras BAI (deny-overrides, promoção de turnos, bônus de continuidade) e colapsa no maior score. É o Ψ da Tiffany executável.", "epistemico": "fato", "exemplos": [], "id": "colapso_multicamada_peso", "memoria": "semantica", "meta": {"autor": "broca-so", "dominio": "operacao_so", "fonte": "conversa/colapso.py:colapso/score_hit"}, "origem": "wiki", "palavras_chave": [], "sinonimos": [], "tipo": "conceito", "titulo": "O Colapso Ψ em Código (varredura ponderada)"}
\section{O Colapso \ensuremath{\Psi} em Código (varredura ponderada)}
conversa/colapso.py varre as camadas dialogos(8)\ensuremath{\to}conhecimento(6)\ensuremath{\to}corpus(4), pontua cada candidato como overlap Koch + peso da camada + relevância + regras BAI (deny-overrides, promoção de turnos, bônus de continuidade) e colapsa no maior score. É o \ensuremath{\Psi} da Tiffany executável.
\prov{relações: explica colapso\_koch\_tiffany; usa colapso\_ping\_pong\_overlap; usa bai\_deny\_overrides; parte\_de operacao\_broca\_so}
\prov{proveniência: wiki \textperiodcentered{} conversa/colapso.py:colapso/score\_hit \textperiodcentered{} fato \textperiodcentered{} confiança 1.0 \textperiodcentered{} domínio operacao\_so \textperiodcentered{} história 5 versões no jornal}

%CRISTAL {"arestas": [{"alvo": "colapso_koch_tiffany", "confianca": 1.0, "epistemico": "fato", "origem": "wiki", "rel": "explica"}, {"alvo": "impressao_koch_32bytes", "confianca": 1.0, "epistemico": "fato", "origem": "wiki", "rel": "usa"}, {"alvo": "hopfield", "confianca": 1.0, "epistemico": "fato", "origem": "wiki", "rel": "relaciona"}, {"alvo": "operacao_broca_so", "confianca": 1.0, "epistemico": "fato", "origem": "wiki", "rel": "parte_de"}], "confianca": 1.0, "contraexemplos": [], "descricao": "A 'fala' da Tiffany NÃO gera texto: escaneia o .idx (no metal ou Python+mmap), escolhe o nó de maior overlap N−2·popcount(q XOR i) e lê do disco a única linha vencedora. Não há geração — só recuperação por máxima correlação. É o operador Ψ encarnado.", "epistemico": "fato", "exemplos": [], "id": "colapso_ping_pong_overlap", "memoria": "semantica", "meta": {"autor": "broca-so", "dominio": "operacao_so", "fonte": "code/prosa.py:responde"}, "origem": "wiki", "palavras_chave": [], "sinonimos": [], "tipo": "conceito", "titulo": "A Fala é Colapso por Overlap (não geração)"}
\section{A Fala é Colapso por Overlap (não geração)}
A 'fala' da Tiffany NÃO gera texto: escaneia o .idx (no metal ou Python+mmap), escolhe o nó de maior overlap N\ensuremath{-}2\textperiodcentered popcount(q XOR i) e lê do disco a única linha vencedora. Não há geração — só recuperação por máxima correlação. É o operador \ensuremath{\Psi} encarnado.
\prov{relações: explica colapso\_koch\_tiffany; usa impressao\_koch\_32bytes; relaciona hopfield; parte\_de operacao\_broca\_so}
\prov{proveniência: wiki \textperiodcentered{} code/prosa.py:responde \textperiodcentered{} fato \textperiodcentered{} confiança 1.0 \textperiodcentered{} domínio operacao\_so \textperiodcentered{} história 5 versões no jornal}

%CRISTAL {"arestas": [{"alvo": "latex", "confianca": 1.0, "epistemico": "fato", "origem": "wiki", "rel": "parte_de"}, {"alvo": "linguagem_de_marcacao", "confianca": 1.0, "epistemico": "fato", "origem": "wiki", "rel": "usa"}], "confianca": 1.0, "contraexemplos": [], "descricao": "Uma instrução iniciada por barra (α, frac, √) que o compilador expande num símbolo ou construção — o token da linguagem LaTeX.", "epistemico": "fato", "exemplos": [], "id": "comando_latex", "memoria": "semantica", "meta": {"dominio": "programacao", "fonte": ""}, "origem": "wiki", "palavras_chave": [], "sinonimos": [], "tipo": "conceito", "titulo": "Comando LaTeX"}
\section{Comando LaTeX}
Uma instrução iniciada por barra (\ensuremath{\alpha}, frac, \ensuremath{\sqrt{\,}}) que o compilador expande num símbolo ou construção — o token da linguagem LaTeX.
\prov{relações: parte\_de latex; usa linguagem\_de\_marcacao}
\prov{proveniência: wiki \textperiodcentered{} fato \textperiodcentered{} confiança 1.0 \textperiodcentered{} domínio programacao \textperiodcentered{} história 4 versões no jornal}

%CRISTAL {"arestas": [{"alvo": "compilador", "confianca": 1.0, "epistemico": "fato", "origem": "wiki", "rel": "especializa"}], "confianca": 1.0, "contraexemplos": [], "descricao": "Compilar para código nativo em tempo de execução, aproveitando o que só se sabe rodando — o meio-termo veloz (V8, JVM).", "epistemico": "fato", "exemplos": [], "id": "compilacao_jit", "memoria": "semantica", "meta": {"dominio": "programacao", "fonte": ""}, "origem": "wiki", "palavras_chave": [], "sinonimos": [], "tipo": "conceito", "titulo": "Compilação JIT"}
\section{Compilação JIT}
Compilar para código nativo em tempo de execução, aproveitando o que só se sabe rodando — o meio-termo veloz (V8, JVM).
\prov{relações: especializa compilador}
\prov{proveniência: wiki \textperiodcentered{} fato \textperiodcentered{} confiança 1.0 \textperiodcentered{} domínio programacao \textperiodcentered{} história 4 versões no jornal}

%CRISTAL {"arestas": [{"alvo": "linker", "confianca": 1.0, "epistemico": "fato", "origem": "curriculo", "rel": "depende"}, {"alvo": "broca_os", "confianca": 1.0, "epistemico": "fato", "origem": "curriculo", "rel": "referencia"}, {"alvo": "historia", "confianca": 0.5, "epistemico": "fato", "origem": "wiki", "rel": "relaciona"}, {"alvo": "utilitarismo", "confianca": 0.5, "epistemico": "fato", "origem": "wiki", "rel": "relaciona"}, {"alvo": "pell", "confianca": 0.5, "epistemico": "fato", "origem": "wiki", "rel": "relaciona"}, {"alvo": "o_que_e_no_projecto", "confianca": 0.5, "epistemico": "fato", "origem": "wiki", "rel": "relaciona"}, {"alvo": "consulta_simples", "confianca": 0.5, "epistemico": "fato", "origem": "wiki", "rel": "relaciona"}, {"alvo": "linguagem_de_programacao", "confianca": 1.0, "epistemico": "fato", "origem": "wiki", "rel": "parte_de"}], "confianca": 1.0, "contraexemplos": [], "descricao": "Traduz código-fonte para IR ou binário.", "epistemico": "fato", "exemplos": [], "id": "compilador", "memoria": "semantica", "meta": {"dominio": "programacao", "nivel": 6, "ordem": 11}, "origem": "curriculo", "palavras_chave": [], "sinonimos": [], "tipo": "conceito", "titulo": "compilador"}
\section{compilador}
Traduz código-fonte para IR ou binário.
\prov{relações: depende linker; referencia broca\_os; relaciona historia; relaciona utilitarismo; relaciona pell; relaciona o\_que\_e\_no\_projecto; relaciona consulta\_simples; parte\_de linguagem\_de\_programacao}
\prov{proveniência: curriculo \textperiodcentered{} fato \textperiodcentered{} confiança 1.0 \textperiodcentered{} domínio programacao \textperiodcentered{} história 12 versões no jornal}

%CRISTAL {"arestas": [{"alvo": "compilador_glsl_statement_ptx", "confianca": 1.0, "epistemico": "fato", "origem": "wiki", "rel": "relaciona"}, {"alvo": "arch_gpu_turing", "confianca": 1.0, "epistemico": "fato", "origem": "wiki", "rel": "usa"}, {"alvo": "compilador_glsl_broca", "confianca": 1.0, "epistemico": "fato", "origem": "wiki", "rel": "parte_de"}], "confianca": 1.0, "contraexemplos": [], "descricao": "FEITO: linguagem/glsl_ptx_frag.py compila um corpo de FRAGMENT SHADER (com vec3, float, for, swizzle .x/.y/.z, e return vec3) para UM kernel PTX fundido — cada thread um pixel, escreve (r,g,b) num buffer H×W×3. Um vec3 = 3 registradores %f; componentwise (vec+vec, vec*escalar); builtins dot/length/normalize/mix/clamp (o dot e a normalize usam FFMA/rsqrt = o gato/gambá); px,py (o fragCoord) vêm do índice do thread (rem/div). Verificado (3/3): um shader plasma (vec3, swizzle, length, clamp, sin) compila→PTX→SASS(sm_75), renderiza na GPU e a imagem bate numpy a 1,5e-6. É o último degrau: o fragment shader completo do metal a partir do GLSL. Falta só o if e a interseção de esfera p/ o fur ball inteiro. glsl_ptx_frag.py (3/3).", "epistemico": "fato", "exemplos": [], "id": "compilador_glsl_fragment_ptx", "memoria": "semantica", "meta": {"dominio": "sass_gpu", "fonte": "SASS/GPU (o metal)"}, "origem": "wiki", "palavras_chave": [], "sinonimos": [], "tipo": "conceito", "titulo": "O compilador GLSL FRAGMENT → PTX: vec3, swizzle, a cor RGB por pixel (sm_75)"}
\section{O compilador GLSL FRAGMENT \ensuremath{\to} PTX: vec3, swizzle, a cor RGB por pixel (sm\_75)}
FEITO: linguagem/glsl\_ptx\_frag.py compila um corpo de FRAGMENT SHADER (com vec3, float, for, swizzle .x/.y/.z, e return vec3) para UM kernel PTX fundido — cada thread um pixel, escreve (r,g,b) num buffer H$\times$W$\times$3. Um vec3 = 3 registradores \%f; componentwise (vec+vec, vec*escalar); builtins dot/length/normalize/mix/clamp (o dot e a normalize usam FFMA/rsqrt = o gato/gambá); px,py (o fragCoord) vêm do índice do thread (rem/div). Verificado (3/3): um shader plasma (vec3, swizzle, length, clamp, sin) compila\ensuremath{\to}PTX\ensuremath{\to}SASS(sm\_75), renderiza na GPU e a imagem bate numpy a 1,5e-6. É o último degrau: o fragment shader completo do metal a partir do GLSL. Falta só o if e a interseção de esfera p/ o fur ball inteiro. glsl\_ptx\_frag.py (3/3).
\prov{relações: relaciona compilador\_glsl\_statement\_ptx; usa arch\_gpu\_turing; parte\_de compilador\_glsl\_broca}
\prov{proveniência: wiki \textperiodcentered{} SASS/GPU (o metal) \textperiodcentered{} fato \textperiodcentered{} confiança 1.0 \textperiodcentered{} domínio sass\_gpu \textperiodcentered{} história 8 versões no jornal}

%CRISTAL {"arestas": [{"alvo": "backend_ptx_sm75_funciona", "confianca": 1.0, "epistemico": "fato", "origem": "wiki", "rel": "relaciona"}, {"alvo": "ptx_isa_virtual", "confianca": 1.0, "epistemico": "fato", "origem": "wiki", "rel": "usa"}, {"alvo": "dicionario_maquina_universal_sass", "confianca": 1.0, "epistemico": "fato", "origem": "wiki", "rel": "usa"}, {"alvo": "arch_gpu_turing", "confianca": 1.0, "epistemico": "fato", "origem": "wiki", "rel": "usa"}, {"alvo": "compilador_glsl_broca", "confianca": 1.0, "epistemico": "fato", "origem": "wiki", "rel": "parte_de"}, {"alvo": "pipeline_compilacao_gpu", "confianca": 1.0, "epistemico": "fato", "origem": "wiki", "rel": "parte_de"}], "confianca": 1.0, "contraexemplos": [], "descricao": "FEITO: linguagem/glsl_ptx_compilador.py compila uma expressão GLSL ARBITRÁRIA sobre x → PTX → SASS(sm_75) → a GPU (não é template). Parser de descida recursiva (expr/term/factor), emissor PTX em SSA (cada subexpressão um registrador %f), com os builtins (sin/cos/sqrt/rsqrt/ex2/lg2/abs = MUFU = o gambá) e a FUSÃO do padrão a*b+c → FFMA (fma.rn.f32 = o gato/La Hire). Verificado (6/6): 'sin(a*x+b)', 'a*x*x+b*x+0.5', 'sqrt(abs(x))+cos(x*3.0)', 'sin(x)*cos(x)-a' compilam, rodam em sm_75 e batem numpy a ~4e-7. Uniforms nomeados (a,b,…) viram .param .f32; x é o elemento por-thread. É o degrau que fecha o compilador: GLSL → PTX → SASS, sob o Python (não através do numpy). O próximo é o kernel completo do fur ball (o loop de shells + o hash 3D) num só kernel fundido → 60+ fps. glsl_ptx_compilador.py (6/6).", "epistemico": "fato", "exemplos": [], "id": "compilador_glsl_ptx_real", "memoria": "semantica", "meta": {"dominio": "sass_gpu", "fonte": "SASS/GPU (o metal)"}, "origem": "wiki", "palavras_chave": [], "sinonimos": [], "tipo": "conceito", "titulo": "O compilador GLSL-expressão → PTX (de verdade, não template) roda em sm_75"}
\section{O compilador GLSL-expressão \ensuremath{\to} PTX (de verdade, não template) roda em sm\_75}
FEITO: linguagem/glsl\_ptx\_compilador.py compila uma expressão GLSL ARBITRÁRIA sobre x \ensuremath{\to} PTX \ensuremath{\to} SASS(sm\_75) \ensuremath{\to} a GPU (não é template). Parser de descida recursiva (expr/term/factor), emissor PTX em SSA (cada subexpressão um registrador \%f), com os builtins (sin/cos/sqrt/rsqrt/ex2/lg2/abs = MUFU = o gambá) e a FUSÃO do padrão a*b+c \ensuremath{\to} FFMA (fma.rn.f32 = o gato/La Hire). Verificado (6/6): 'sin(a*x+b)', 'a*x*x+b*x+0.5', 'sqrt(abs(x))+cos(x*3.0)', 'sin(x)*cos(x)-a' compilam, rodam em sm\_75 e batem numpy a \~{}4e-7. Uniforms nomeados (a,b,…) viram .param .f32; x é o elemento por-thread. É o degrau que fecha o compilador: GLSL \ensuremath{\to} PTX \ensuremath{\to} SASS, sob o Python (não através do numpy). O próximo é o kernel completo do fur ball (o loop de shells + o hash 3D) num só kernel fundido \ensuremath{\to} 60+ fps. glsl\_ptx\_compilador.py (6/6).
\prov{relações: relaciona backend\_ptx\_sm75\_funciona; usa ptx\_isa\_virtual; usa dicionario\_maquina\_universal\_sass; usa arch\_gpu\_turing; parte\_de compilador\_glsl\_broca; parte\_de pipeline\_compilacao\_gpu}
\prov{proveniência: wiki \textperiodcentered{} SASS/GPU (o metal) \textperiodcentered{} fato \textperiodcentered{} confiança 1.0 \textperiodcentered{} domínio sass\_gpu \textperiodcentered{} história 28 versões no jornal}

%CRISTAL {"arestas": [{"alvo": "compilador_glsl_ptx_real", "confianca": 1.0, "epistemico": "fato", "origem": "wiki", "rel": "relaciona"}, {"alvo": "backend_ptx_sm75_funciona", "confianca": 1.0, "epistemico": "fato", "origem": "wiki", "rel": "usa"}, {"alvo": "shader_fios_shell_paralelo", "confianca": 1.0, "epistemico": "fato", "origem": "wiki", "rel": "relaciona"}, {"alvo": "arch_gpu_turing", "confianca": 1.0, "epistemico": "fato", "origem": "wiki", "rel": "usa"}, {"alvo": "compilador_glsl_broca", "confianca": 1.0, "epistemico": "fato", "origem": "wiki", "rel": "parte_de"}], "confianca": 1.0, "contraexemplos": [], "descricao": "FEITO: linguagem/glsl_ptx_shader.py compila um corpo GLSL com STATEMENTS (float nome=expr; nome+=expr; for(int k=0;k<N;k++){...}; return expr;) para UM kernel PTX FUNDIDO, um pixel por thread, e roda em sm_75. Suporta o `for` (→ labels LOOP/END + setp.ge + @%p bra, o BRA), regs mutáveis (o acumulador), float(k) (cvt.rn.f32.u32), os builtins (MUFU=gambá), a fusão a*b+c→FFMA (o gato). Verificado (3/3): o SHADER do march de shells (for de 64, acc+=sin(a*x+b+k*step)) — a estrutura do fur ball em GLSL DE VERDADE — compila→PTX→SASS(sm_75), roda @840×473 e bate numpy a 3e-5. O kernel fundido roda a ~0,17 ms (o mesmo do feito à mão, ~5800 fps); o for virou um loop no thread (não K kernels do torch). É o degrau que fecha a cadeia GLSL(statements)→PTX→SASS→GPU. Falta só os tipos vetoriais (vec2/vec3) e o mainImage p/ o shader completo. glsl_ptx_shader.py (3/3).", "epistemico": "fato", "exemplos": [], "id": "compilador_glsl_statement_ptx", "memoria": "semantica", "meta": {"dominio": "sass_gpu", "fonte": "SASS/GPU (o metal)"}, "origem": "wiki", "palavras_chave": [], "sinonimos": [], "tipo": "conceito", "titulo": "O compilador GLSL statement-level → PTX: o for/declarações num kernel fundido (sm_75)"}
\section{O compilador GLSL statement-level \ensuremath{\to} PTX: o for/declarações num kernel fundido (sm\_75)}
FEITO: linguagem/glsl\_ptx\_shader.py compila um corpo GLSL com STATEMENTS (float nome=expr; nome+=expr; for(int k=0;k<N;k++)\{...\}; return expr;) para UM kernel PTX FUNDIDO, um pixel por thread, e roda em sm\_75. Suporta o `for` (\ensuremath{\to} labels LOOP/END + setp.ge + @\%p bra, o BRA), regs mutáveis (o acumulador), float(k) (cvt.rn.f32.u32), os builtins (MUFU=gambá), a fusão a*b+c\ensuremath{\to}FFMA (o gato). Verificado (3/3): o SHADER do march de shells (for de 64, acc+=sin(a*x+b+k*step)) — a estrutura do fur ball em GLSL DE VERDADE — compila\ensuremath{\to}PTX\ensuremath{\to}SASS(sm\_75), roda @840$\times$473 e bate numpy a 3e-5. O kernel fundido roda a \~{}0,17 ms (o mesmo do feito à mão, \~{}5800 fps); o for virou um loop no thread (não K kernels do torch). É o degrau que fecha a cadeia GLSL(statements)\ensuremath{\to}PTX\ensuremath{\to}SASS\ensuremath{\to}GPU. Falta só os tipos vetoriais (vec2/vec3) e o mainImage p/ o shader completo. glsl\_ptx\_shader.py (3/3).
\prov{relações: relaciona compilador\_glsl\_ptx\_real; usa backend\_ptx\_sm75\_funciona; relaciona shader\_fios\_shell\_paralelo; usa arch\_gpu\_turing; parte\_de compilador\_glsl\_broca}
\prov{proveniência: wiki \textperiodcentered{} SASS/GPU (o metal) \textperiodcentered{} fato \textperiodcentered{} confiança 1.0 \textperiodcentered{} domínio sass\_gpu \textperiodcentered{} história 18 versões no jornal}

%CRISTAL {"arestas": [{"alvo": "framework_frontend", "confianca": 1.0, "epistemico": "fato", "origem": "wiki", "rel": "usa"}, {"alvo": "modularidade", "confianca": 1.0, "epistemico": "fato", "origem": "wiki", "rel": "especializa"}], "confianca": 1.0, "contraexemplos": [], "descricao": "Uma peça de interface autocontida (marcação + estilo + lógica) que se compõe e reusa — modularidade aplicada à tela.", "epistemico": "fato", "exemplos": [], "id": "componente_ui", "memoria": "semantica", "meta": {"dominio": "programacao", "fonte": ""}, "origem": "wiki", "palavras_chave": [], "sinonimos": [], "tipo": "conceito", "titulo": "Componente (UI)"}
\section{Componente (UI)}
Uma peça de interface autocontida (marcação + estilo + lógica) que se compõe e reusa — modularidade aplicada à tela.
\prov{relações: usa framework\_frontend; especializa modularidade}
\prov{proveniência: wiki \textperiodcentered{} fato \textperiodcentered{} confiança 1.0 \textperiodcentered{} domínio programacao \textperiodcentered{} história 6 versões no jornal}

%CRISTAL {"arestas": [{"alvo": "arquitetura_von_neumann", "confianca": 1.0, "epistemico": "fato", "origem": "curriculo", "rel": "relaciona"}, {"alvo": "maquina_de_turing", "confianca": 1.0, "epistemico": "fato", "origem": "curriculo", "rel": "relaciona"}, {"alvo": "broca_os", "confianca": 1.0, "epistemico": "fato", "origem": "curriculo", "rel": "relaciona"}, {"alvo": "algoritmo", "confianca": 1.0, "epistemico": "fato", "origem": "curriculo", "rel": "relaciona"}, {"alvo": "morte", "confianca": 1.0, "epistemico": "fato", "origem": "curriculo", "rel": "relaciona"}, {"alvo": "vida", "confianca": 1.0, "epistemico": "fato", "origem": "curriculo", "rel": "relaciona"}, {"alvo": "computabilidade", "confianca": 1.0, "epistemico": "fato", "origem": "curriculo", "rel": "relaciona"}, {"alvo": "bit", "confianca": 1.0, "epistemico": "fato", "origem": "curriculo", "rel": "parte_de"}, {"alvo": "matematica", "confianca": 1.0, "epistemico": "fato", "origem": "curriculo", "rel": "parte_de"}], "confianca": 1.0, "contraexemplos": [], "descricao": "Procedimentos efectivos sobre símbolos; fase 4½ — Von Neumann → Turing → broca-so.", "epistemico": "fato", "exemplos": [], "id": "computacao", "memoria": "semantica", "meta": {"dominio": "computacao", "nivel": 4, "serve": "Hub computação Tiffany", "verificacao": "slug idêntico — enriquecer, não duplicar"}, "origem": "curriculo", "palavras_chave": ["algoritmo", "decidibilidade", "máquina"], "sinonimos": ["computação"], "tipo": "conceito", "titulo": "computacao"}
\section{computacao}
Procedimentos efectivos sobre símbolos; fase 4\ensuremath{\tfrac12} — Von Neumann \ensuremath{\to} Turing \ensuremath{\to} broca-so.
\prov{sinónimos: computação}
\prov{relações: relaciona arquitetura\_von\_neumann; relaciona maquina\_de\_turing; relaciona broca\_os; relaciona algoritmo; relaciona morte; relaciona vida; relaciona computabilidade; parte\_de bit; parte\_de matematica}
\prov{proveniência: curriculo \textperiodcentered{} fato \textperiodcentered{} confiança 1.0 \textperiodcentered{} domínio computacao \textperiodcentered{} história 77 versões no jornal}

%CRISTAL {"arestas": [{"alvo": "algebra_linear", "confianca": 1.0, "epistemico": "fato", "origem": "wiki", "rel": "usa"}, {"alvo": "gato", "confianca": 1.0, "epistemico": "fato", "origem": "wiki", "rel": "relaciona"}], "confianca": 1.0, "contraexemplos": [], "descricao": "Tratar o cálculo como operações de matrizes e vetores — o substrato de GPUs, redes neurais e do processamento em lote.", "epistemico": "fato", "exemplos": [], "id": "computacao_matricial", "memoria": "semantica", "meta": {"dominio": "programacao", "fonte": ""}, "origem": "wiki", "palavras_chave": [], "sinonimos": [], "tipo": "conceito", "titulo": "Computação como Álgebra Linear"}
\section{Computação como Álgebra Linear}
Tratar o cálculo como operações de matrizes e vetores — o substrato de GPUs, redes neurais e do processamento em lote.
\prov{relações: usa algebra\_linear; relaciona gato}
\prov{proveniência: wiki \textperiodcentered{} fato \textperiodcentered{} confiança 1.0 \textperiodcentered{} domínio programacao \textperiodcentered{} história 6 versões no jornal}

%CRISTAL {"arestas": [{"alvo": "programacao", "confianca": 1.0, "epistemico": "fato", "origem": "wiki", "rel": "parte_de"}, {"alvo": "termodinamica", "confianca": 1.0, "epistemico": "fato", "origem": "wiki", "rel": "usa"}, {"alvo": "gerenciamento_de_memoria", "confianca": 1.0, "epistemico": "fato", "origem": "wiki", "rel": "relaciona"}], "confianca": 1.0, "contraexemplos": [], "descricao": "Computar sem apagar informação: cada passo é invertível (Janus, portas de Toffoli). Landauer: só o apagar dissipa calor — o reversível pode ser (quase) sem custo térmico.", "epistemico": "fato", "exemplos": [], "id": "computacao_reversivel", "memoria": "semantica", "meta": {"autor": "Landauer", "dominio": "programacao", "fonte": ""}, "origem": "wiki", "palavras_chave": [], "sinonimos": [], "tipo": "conceito", "titulo": "Computação Reversível"}
\section{Computação Reversível}
Computar sem apagar informação: cada passo é invertível (Janus, portas de Toffoli). Landauer: só o apagar dissipa calor — o reversível pode ser (quase) sem custo térmico.
\prov{relações: parte\_de programacao; usa termodinamica; relaciona gerenciamento\_de\_memoria}
\prov{proveniência: wiki \textperiodcentered{} fato \textperiodcentered{} confiança 1.0 \textperiodcentered{} domínio programacao \textperiodcentered{} história 8 versões no jornal}

%CRISTAL {"arestas": [{"alvo": "gato_ula", "confianca": 1.0, "epistemico": "fato", "origem": "wiki", "rel": "usa"}, {"alvo": "cristal_memoria", "confianca": 1.0, "epistemico": "fato", "origem": "wiki", "rel": "usa"}, {"alvo": "transformada_mineral_7d", "confianca": 1.0, "epistemico": "fato", "origem": "wiki", "rel": "usa"}, {"alvo": "molecula_mineral", "confianca": 1.0, "epistemico": "fato", "origem": "wiki", "rel": "usa"}, {"alvo": "corpo_mineral", "confianca": 1.0, "epistemico": "fato", "origem": "wiki", "rel": "parte_de"}, {"alvo": "sandbox_cristal_ao_so", "confianca": 1.0, "epistemico": "fato", "origem": "wiki", "rel": "relaciona"}, {"alvo": "linguagem_matriz_gatos", "confianca": 1.0, "epistemico": "fato", "origem": "wiki", "rel": "usa"}], "confianca": 1.0, "contraexemplos": [], "descricao": "A transformada octoniônica mineral como máquina que grava informação nas órbitas das moléculas. A MATRIZ DE GATOS é a ULA (Aₙ=[[n,1],[1,0]], execução=contração, ×Aₙ=1 passo O(1), det Aₙ=−1 conservado); os CRISTAIS são a MEMÓRIA (células octoniônicas que seguram o dado, norma conservada = integridade); a MOLÉCULA (sequência de cristais) é o registrador, e a órbita é o canal. Grava-se na órbita (transformada), a memória segura, a ULA computa. Reversível (𝕆 divisão). Verificado.", "epistemico": "inferencia", "exemplos": [], "id": "computador_mineral", "memoria": "semantica", "meta": {"dominio": "computacao", "fonte": "computador mineral"}, "origem": "wiki", "palavras_chave": [], "sinonimos": [], "tipo": "conceito", "titulo": "O Computador Mineral (gato=ULA, cristais=memória)"}
\section{O Computador Mineral (gato=ULA, cristais=memória)}
A transformada octoniônica mineral como máquina que grava informação nas órbitas das moléculas. A MATRIZ DE GATOS é a ULA (A\ensuremath{_{n}}=[[n,1],[1,0]], execução=contração, $\times$A\ensuremath{_{n}}=1 passo O(1), det A\ensuremath{_{n}}=\ensuremath{-}1 conservado); os CRISTAIS são a MEMÓRIA (células octoniônicas que seguram o dado, norma conservada = integridade); a MOLÉCULA (sequência de cristais) é o registrador, e a órbita é o canal. Grava-se na órbita (transformada), a memória segura, a ULA computa. Reversível (\ensuremath{\mathbb{O}} divisão). Verificado.
\prov{relações: usa gato\_ula; usa cristal\_memoria; usa transformada\_mineral\_7d; usa molecula\_mineral; parte\_de corpo\_mineral; relaciona sandbox\_cristal\_ao\_so; usa linguagem\_matriz\_gatos}
\prov{proveniência: wiki \textperiodcentered{} computador mineral \textperiodcentered{} inferencia \textperiodcentered{} confiança 1.0 \textperiodcentered{} domínio computacao \textperiodcentered{} história 8 versões no jornal}

%CRISTAL {"arestas": [{"alvo": "rede_de_computadores", "confianca": 1.0, "epistemico": "fato", "origem": "wiki", "rel": "usa"}], "confianca": 1.0, "contraexemplos": [], "descricao": "Quebrar a mensagem em pacotes que viajam independentes e se remontam no destino — a ideia que fez a internet escalar (Baran/Davies).", "epistemico": "fato", "exemplos": [], "id": "comutacao_de_pacotes", "memoria": "semantica", "meta": {"autor": "Baran", "dominio": "programacao", "fonte": ""}, "origem": "wiki", "palavras_chave": [], "sinonimos": [], "tipo": "conceito", "titulo": "Comutação de Pacotes"}
\section{Comutação de Pacotes}
Quebrar a mensagem em pacotes que viajam independentes e se remontam no destino — a ideia que fez a internet escalar (Baran/Davies).
\prov{relações: usa rede\_de\_computadores}
\prov{proveniência: wiki \textperiodcentered{} fato \textperiodcentered{} confiança 1.0 \textperiodcentered{} domínio programacao \textperiodcentered{} história 4 versões no jornal}

%CRISTAL {"arestas": [{"alvo": "programacao_concorrente", "confianca": 1.0, "epistemico": "fato", "origem": "wiki", "rel": "usa"}], "confianca": 1.0, "contraexemplos": [], "descricao": "Concorrência é LIDAR com muitas coisas ao mesmo tempo (estrutura); paralelismo é FAZÊ-las ao mesmo tempo (execução).", "epistemico": "fato", "exemplos": [], "id": "concorrencia_paralelismo", "memoria": "semantica", "meta": {"dominio": "programacao", "fonte": ""}, "origem": "wiki", "palavras_chave": [], "sinonimos": [], "tipo": "conceito", "titulo": "Concorrência vs Paralelismo"}
\section{Concorrência vs Paralelismo}
Concorrência é LIDAR com muitas coisas ao mesmo tempo (estrutura); paralelismo é FAZÊ-las ao mesmo tempo (execução).
\prov{relações: usa programacao\_concorrente}
\prov{proveniência: wiki \textperiodcentered{} fato \textperiodcentered{} confiança 1.0 \textperiodcentered{} domínio programacao \textperiodcentered{} história 4 versões no jornal}

%CRISTAL {"arestas": [{"alvo": "virtualizacao", "confianca": 0.5, "epistemico": "fato", "origem": "wiki", "rel": "relaciona"}, {"alvo": "word", "confianca": 0.5, "epistemico": "fato", "origem": "wiki", "rel": "relaciona"}, {"alvo": "vontade_de_poder", "confianca": 0.5, "epistemico": "fato", "origem": "wiki", "rel": "relaciona"}], "confianca": 1.0, "contraexemplos": [], "descricao": "Grau de certeza de uma aresta — fato, inferência ou hipótese.", "epistemico": "fato", "exemplos": [], "id": "confianca_epistemica", "memoria": "semantica", "meta": {"dominio": "aprendizagem", "nivel": 7, "ordem": 5}, "origem": "curriculo", "palavras_chave": ["confianca", "certeza"], "sinonimos": ["confianca epistemica"], "tipo": "conceito", "titulo": "confiança epistêmica"}
\section{confiança epistêmica}
Grau de certeza de uma aresta — fato, inferência ou hipótese.
\prov{sinónimos: confianca epistemica}
\prov{relações: relaciona virtualizacao; relaciona word; relaciona vontade\_de\_poder}
\prov{proveniência: curriculo \textperiodcentered{} fato \textperiodcentered{} confiança 1.0 \textperiodcentered{} domínio aprendizagem \textperiodcentered{} história 4 versões no jornal}

%CRISTAL {"arestas": [{"alvo": "sistemas_distribuidos", "confianca": 1.0, "epistemico": "fato", "origem": "wiki", "rel": "parte_de"}, {"alvo": "prova_de_trabalho", "confianca": 1.0, "epistemico": "fato", "origem": "wiki", "rel": "explica"}, {"alvo": "goldchain", "confianca": 1.0, "epistemico": "fato", "origem": "wiki", "rel": "explica"}, {"alvo": "dois_generais", "confianca": 1.0, "epistemico": "fato", "origem": "wiki", "rel": "usa"}], "confianca": 1.0, "contraexemplos": [], "descricao": "Fazer nós que falham e mentem concordarem num valor único — Paxos, Raft, e a prova-de-trabalho de Nakamoto. O coração do distribuído.", "epistemico": "fato", "exemplos": [], "id": "consenso_distribuido", "memoria": "semantica", "meta": {"autor": "Lamport", "dominio": "programacao", "fonte": ""}, "origem": "wiki", "palavras_chave": [], "sinonimos": [], "tipo": "conceito", "titulo": "Consenso Distribuído"}
\section{Consenso Distribuído}
Fazer nós que falham e mentem concordarem num valor único — Paxos, Raft, e a prova-de-trabalho de Nakamoto. O coração do distribuído.
\prov{relações: parte\_de sistemas\_distribuidos; explica prova\_de\_trabalho; explica goldchain; usa dois\_generais}
\prov{proveniência: wiki \textperiodcentered{} fato \textperiodcentered{} confiança 1.0 \textperiodcentered{} domínio programacao \textperiodcentered{} história 10 versões no jornal}

%CRISTAL {"arestas": [{"alvo": "teorema_cap", "confianca": 1.0, "epistemico": "fato", "origem": "wiki", "rel": "usa"}, {"alvo": "nosql", "confianca": 1.0, "epistemico": "fato", "origem": "wiki", "rel": "relaciona"}], "confianca": 1.0, "contraexemplos": [], "descricao": "Abrir mão da consistência imediata: as réplicas convergem 'com o tempo'. O preço da disponibilidade e da escala (o lado AP do CAP).", "epistemico": "inferencia", "exemplos": [], "id": "consistencia_eventual", "memoria": "semantica", "meta": {"dominio": "programacao", "fonte": ""}, "origem": "wiki", "palavras_chave": [], "sinonimos": [], "tipo": "conceito", "titulo": "Consistência Eventual"}
\section{Consistência Eventual}
Abrir mão da consistência imediata: as réplicas convergem 'com o tempo'. O preço da disponibilidade e da escala (o lado AP do CAP).
\prov{relações: usa teorema\_cap; relaciona nosql}
\prov{proveniência: wiki \textperiodcentered{} inferencia \textperiodcentered{} confiança 1.0 \textperiodcentered{} domínio programacao \textperiodcentered{} história 6 versões no jornal}

%CRISTAL {"arestas": [{"alvo": "ensinar_conceito", "confianca": 1.0, "epistemico": "fato", "origem": "wiki", "rel": "usa"}, {"alvo": "refinar_densificar", "confianca": 1.0, "epistemico": "fato", "origem": "wiki", "rel": "usa"}, {"alvo": "reproduzir_experimentos", "confianca": 1.0, "epistemico": "fato", "origem": "wiki", "rel": "usa"}, {"alvo": "ingest_codigo_autoreferencial", "confianca": 1.0, "epistemico": "fato", "origem": "wiki", "rel": "usa"}], "confianca": 1.0, "contraexemplos": [], "descricao": "conversa/conhecimento/lab.py é o CLI que dispara ensinar/refinar/ingest/reproduzir/floresta — o painel humano sobre o grafo. 'floresta' reporta conceitos/componentes/isolados e mais-usados; é a régua de saúde do tronco.", "epistemico": "fato", "exemplos": [], "id": "console_do_operador_lab", "memoria": "semantica", "meta": {"autor": "broca-so", "dominio": "operacao_so", "fonte": "conversa/conhecimento/lab.py"}, "origem": "wiki", "palavras_chave": [], "sinonimos": [], "tipo": "conceito", "titulo": "O Console do Operador (lab.py)"}
\section{O Console do Operador (lab.py)}
conversa/conhecimento/lab.py é o CLI que dispara ensinar/refinar/ingest/reproduzir/floresta — o painel humano sobre o grafo. 'floresta' reporta conceitos/componentes/isolados e mais-usados; é a régua de saúde do tronco.
\prov{relações: usa ensinar\_conceito; usa refinar\_densificar; usa reproduzir\_experimentos; usa ingest\_codigo\_autoreferencial}
\prov{proveniência: wiki \textperiodcentered{} conversa/conhecimento/lab.py \textperiodcentered{} fato \textperiodcentered{} confiança 1.0 \textperiodcentered{} domínio operacao\_so \textperiodcentered{} história 5 versões no jornal}

%CRISTAL {"arestas": [{"alvo": "nucleo_paralelo", "confianca": 1.0, "epistemico": "fato", "origem": "wiki", "rel": "usa"}, {"alvo": "protocolo_assincrono", "confianca": 1.0, "epistemico": "fato", "origem": "wiki", "rel": "relaciona"}], "confianca": 1.0, "contraexemplos": [], "descricao": "Um ramo só minera se não está numa etapa pow/lock/redeem na fase corrente E a fila no mempool ≤ limite E o valor do ciclo ≥ MICRO_VALOR (300 sats); senão _motivo_bloqueio explica (minera / fila / chain baixo).", "epistemico": "fato", "exemplos": [], "id": "conta_pronta_gate", "memoria": "semantica", "meta": {"autor": "broca-so", "dominio": "operacao_so", "fonte": "goldchain/nucleo_minera.py:conta_pronta"}, "origem": "wiki", "palavras_chave": [], "sinonimos": [], "tipo": "conceito", "titulo": "Elegibilidade de Ramo (gate de prontidão)"}
\section{Elegibilidade de Ramo (gate de prontidão)}
Um ramo só minera se não está numa etapa pow/lock/redeem na fase corrente E a fila no mempool \ensuremath{\le} limite E o valor do ciclo \ensuremath{\ge} MICRO\_VALOR (300 sats); senão \_motivo\_bloqueio explica (minera / fila / chain baixo).
\prov{relações: usa nucleo\_paralelo; relaciona protocolo\_assincrono}
\prov{proveniência: wiki \textperiodcentered{} goldchain/nucleo\_minera.py:conta\_pronta \textperiodcentered{} fato \textperiodcentered{} confiança 1.0 \textperiodcentered{} domínio operacao\_so \textperiodcentered{} história 3 versões no jornal}

%CRISTAL {"arestas": [{"alvo": "virtualizacao", "confianca": 1.0, "epistemico": "fato", "origem": "wiki", "rel": "especializa"}, {"alvo": "devops", "confianca": 1.0, "epistemico": "fato", "origem": "wiki", "rel": "usa"}], "confianca": 1.0, "contraexemplos": [], "descricao": "Empacotar um app com suas dependências num ambiente isolado que roda igual em qualquer lugar — isolamento de processo, não de hardware (Docker).", "epistemico": "fato", "exemplos": [], "id": "container", "memoria": "semantica", "meta": {"autor": "Hykes", "dominio": "programacao", "fonte": ""}, "origem": "wiki", "palavras_chave": [], "sinonimos": [], "tipo": "conceito", "titulo": "Contêiner"}
\section{Contêiner}
Empacotar um app com suas dependências num ambiente isolado que roda igual em qualquer lugar — isolamento de processo, não de hardware (Docker).
\prov{relações: especializa virtualizacao; usa devops}
\prov{proveniência: wiki \textperiodcentered{} fato \textperiodcentered{} confiança 1.0 \textperiodcentered{} domínio programacao \textperiodcentered{} história 6 versões no jornal}

%CRISTAL {"arestas": [{"alvo": "algebras_hipercomplexas", "confianca": 1.0, "epistemico": "fato", "origem": "wiki", "rel": "usa"}, {"alvo": "a_torre_de_hurwitzcayley--dickson", "confianca": 1.0, "epistemico": "fato", "origem": "wiki", "rel": "relaciona"}, {"alvo": "linguagem_matriz_gatos", "confianca": 1.0, "epistemico": "fato", "origem": "wiki", "rel": "parte_de"}, {"alvo": "tiffany", "confianca": 1.0, "epistemico": "fato", "origem": "wiki", "rel": "parte_de"}], "confianca": 1.0, "contraexemplos": [], "descricao": "A bússola teórica: um gato é um tensor de rank 2 (a matriz), e COMPOR gatos é CONTRAIR índices (a convenção de soma de Einstein: (A·B)ⁱₖ=Σj Aⁱj Bjₖ). Até aqui usamos k=3 — a multiplicação BILINEAR a×b, cujas constantes de estrutura ckᵢj são um tensor de rank 3 (a família hipercomplexa ×ₛ: ℂ,ℍ,𝕆). Generalizar de k=3 para k-TENSORIAL é passar da operação binária à m-ária = o companion m×m = o dupolinômio de grau m. A generalização de Einstein para k-tensor: a relatividade usa a métrica gᵢj (rank 2) e a contração; aqui a 'métrica' é a PARÁBOLA a+c²=1 (a forma bilinear conservada), e a linguagem é a contração de tensores-gato de rank crescente, com a parábola invariante em cada nível. [C] que a torre k-tensorial feche como a de Hurwitz (só o rank muda) — conjectura aberta.", "epistemico": "inferencia", "exemplos": [], "id": "contracao_einstein_k_tensorial", "memoria": "semantica", "meta": {"autor": "Lopes/broca-so", "dominio": "computacao", "fonte": "doc/linguagem_gatos/paper §3,§6; algebras_hipercomplexas (×ₛ, constantes de estrutura rank 3); convenção de Einstein"}, "origem": "wiki", "palavras_chave": [], "sinonimos": [], "tipo": "conceito", "titulo": "A semântica é contração tensorial (Einstein): k=3 (bilinear ×ₛ) generalizada a k-tensorial"}
\section{A semântica é contração tensorial (Einstein): k=3 (bilinear $\times$\ensuremath{_{s}}) generalizada a k-tensorial}
A bússola teórica: um gato é um tensor de rank 2 (a matriz), e COMPOR gatos é CONTRAIR índices (a convenção de soma de Einstein: (A\textperiodcentered B)\ensuremath{^{i}}\ensuremath{_{k}}=\ensuremath{\Sigma}j A\ensuremath{^{i}}j Bj\ensuremath{_{k}}). Até aqui usamos k=3 — a multiplicação BILINEAR a$\times$b, cujas constantes de estrutura ck\ensuremath{_{i}}j são um tensor de rank 3 (a família hipercomplexa $\times$\ensuremath{_{s}}: \ensuremath{\mathbb{C}},\ensuremath{\mathbb{H}},\ensuremath{\mathbb{O}}). Generalizar de k=3 para k-TENSORIAL é passar da operação binária à m-ária = o companion m$\times$m = o dupolinômio de grau m. A generalização de Einstein para k-tensor: a relatividade usa a métrica g\ensuremath{_{i}}j (rank 2) e a contração; aqui a 'métrica' é a PARÁBOLA a+c\ensuremath{^{2}}=1 (a forma bilinear conservada), e a linguagem é a contração de tensores-gato de rank crescente, com a parábola invariante em cada nível. [C] que a torre k-tensorial feche como a de Hurwitz (só o rank muda) — conjectura aberta.
\prov{relações: usa algebras\_hipercomplexas; relaciona a\_torre\_de\_hurwitzcayley--dickson; parte\_de linguagem\_matriz\_gatos; parte\_de tiffany}
\prov{proveniência: wiki \textperiodcentered{} doc/linguagem\_gatos/paper §3,§6; algebras\_hipercomplexas ($\times$\ensuremath{_{s}}, constantes de estrutura rank 3); convenção de Einstein \textperiodcentered{} inferencia \textperiodcentered{} confiança 1.0 \textperiodcentered{} domínio computacao \textperiodcentered{} história 47 versões no jornal}

%CRISTAL {"arestas": [{"alvo": "resolucao", "confianca": 1.0, "epistemico": "fato", "origem": "curriculo", "rel": "explica"}, {"alvo": "logica", "confianca": 1.0, "epistemico": "fato", "origem": "curriculo", "rel": "parte_de"}, {"alvo": "word", "confianca": 0.5, "epistemico": "fato", "origem": "wiki", "rel": "relaciona"}, {"alvo": "virtualizacao", "confianca": 0.5, "epistemico": "fato", "origem": "wiki", "rel": "relaciona"}, {"alvo": "syscall", "confianca": 0.5, "epistemico": "fato", "origem": "wiki", "rel": "relaciona"}], "confianca": 1.0, "contraexemplos": [], "descricao": "Duas afirmações incompatíveis sobre o mesmo par de conceitos.", "epistemico": "fato", "exemplos": [], "id": "contradicao", "memoria": "semantica", "meta": {"dominio": "aprendizagem", "nivel": 7, "ordem": 3}, "origem": "curriculo", "palavras_chave": ["incompatível", "paradoxo"], "sinonimos": ["contradições", "conflito"], "tipo": "conceito", "titulo": "contradicao"}
\section{contradicao}
Duas afirmações incompatíveis sobre o mesmo par de conceitos.
\prov{sinónimos: contradições; conflito}
\prov{relações: explica resolucao; parte\_de logica; relaciona word; relaciona virtualizacao; relaciona syscall}
\prov{proveniência: curriculo \textperiodcentered{} fato \textperiodcentered{} confiança 1.0 \textperiodcentered{} domínio aprendizagem \textperiodcentered{} história 6 versões no jornal}

%CRISTAL {"arestas": [{"alvo": "devops", "confianca": 1.0, "epistemico": "fato", "origem": "wiki", "rel": "usa"}], "confianca": 1.0, "contraexemplos": [], "descricao": "Rastrear cada mudança no código, ramificar e reconciliar — a máquina do tempo do software; sem ela, colaboração é caos.", "epistemico": "fato", "exemplos": [], "id": "controle_de_versao", "memoria": "semantica", "meta": {"dominio": "programacao", "fonte": ""}, "origem": "wiki", "palavras_chave": [], "sinonimos": [], "tipo": "conceito", "titulo": "Controle de Versão"}
\section{Controle de Versão}
Rastrear cada mudança no código, ramificar e reconciliar — a máquina do tempo do software; sem ela, colaboração é caos.
\prov{relações: usa devops}
\prov{proveniência: wiki \textperiodcentered{} fato \textperiodcentered{} confiança 1.0 \textperiodcentered{} domínio programacao \textperiodcentered{} história 4 versões no jornal}

%CRISTAL {"arestas": [{"alvo": "conversao_mat_latex", "confianca": 1.0, "epistemico": "fato", "origem": "wiki", "rel": "explica"}, {"alvo": "traducao_como_transformacao", "confianca": 1.0, "epistemico": "fato", "origem": "wiki", "rel": "especializa"}, {"alvo": "significado_como_invariante", "confianca": 1.0, "epistemico": "fato", "origem": "wiki", "rel": "usa"}, {"alvo": "traducao_e_intraduzivel", "confianca": 1.0, "epistemico": "fato", "origem": "wiki", "rel": "contradiz"}, {"alvo": "latex", "confianca": 1.0, "epistemico": "fato", "origem": "wiki", "rel": "usa"}], "confianca": 1.0, "contraexemplos": [], "descricao": "Math↔LaTeX é a tradução-como-transformada no caso PERFEITO: o significado matemático é EXATAMENTE o invariante (não há intraduzível), e as duas formas são duas 'bases' do mesmo sentido. A tradução que sempre conserva — o limite ideal para o qual a tradução de línguas naturais só tende.", "epistemico": "inferencia", "exemplos": [], "id": "conversao_como_traducao_formal", "memoria": "semantica", "meta": {"dominio": "programacao", "fonte": ""}, "origem": "wiki", "palavras_chave": [], "sinonimos": [], "tipo": "conceito", "titulo": "Conversão como Tradução Formal"}
\section{Conversão como Tradução Formal}
Math\ensuremath{\leftrightarrow}LaTeX é a tradução-como-transformada no caso PERFEITO: o significado matemático é EXATAMENTE o invariante (não há intraduzível), e as duas formas são duas 'bases' do mesmo sentido. A tradução que sempre conserva — o limite ideal para o qual a tradução de línguas naturais só tende.
\prov{relações: explica conversao\_mat\_latex; especializa traducao\_como\_transformacao; usa significado\_como\_invariante; contradiz traducao\_e\_intraduzivel; usa latex}
\prov{proveniência: wiki \textperiodcentered{} inferencia \textperiodcentered{} confiança 1.0 \textperiodcentered{} domínio programacao \textperiodcentered{} história 6 versões no jornal}

%CRISTAL {"arestas": [{"alvo": "correspondencia_mat_latex", "confianca": 1.0, "epistemico": "fato", "origem": "wiki", "rel": "usa"}, {"alvo": "latex_matematico", "confianca": 1.0, "epistemico": "fato", "origem": "wiki", "rel": "usa"}, {"alvo": "mathml", "confianca": 1.0, "epistemico": "fato", "origem": "wiki", "rel": "relaciona"}, {"alvo": "sintaxe", "confianca": 1.0, "epistemico": "fato", "origem": "wiki", "rel": "usa"}], "confianca": 1.0, "contraexemplos": [], "descricao": "Converter uma notação na outra nos dois sentidos — o conversor real (linguagem/mat_latex.py: para_latex/de_latex) mapeia símbolos, expoentes/índices e a estrutura def-teorema-prova. Ida-e-volta conserva os símbolos.", "epistemico": "fato", "exemplos": [], "id": "conversao_mat_latex", "memoria": "semantica", "meta": {"dominio": "programacao", "fonte": "linguagem/mat_latex.py (para_latex, de_latex, bloco_para_latex)"}, "origem": "wiki", "palavras_chave": [], "sinonimos": [], "tipo": "conceito", "titulo": "Conversão Matemática ↔ LaTeX"}
\section{Conversão Matemática \ensuremath{\leftrightarrow} LaTeX}
Converter uma notação na outra nos dois sentidos — o conversor real (linguagem/mat\_latex.py: para\_latex/de\_latex) mapeia símbolos, expoentes/índices e a estrutura def-teorema-prova. Ida-e-volta conserva os símbolos.
\prov{relações: usa correspondencia\_mat\_latex; usa latex\_matematico; relaciona mathml; usa sintaxe}
\prov{proveniência: wiki \textperiodcentered{} linguagem/mat\_latex.py (para\_latex, de\_latex, bloco\_para\_latex) \textperiodcentered{} fato \textperiodcentered{} confiança 1.0 \textperiodcentered{} domínio programacao \textperiodcentered{} história 5 versões no jornal}

%CRISTAL {"arestas": [{"alvo": "design_vetorizado_svg", "confianca": 1.0, "epistemico": "fato", "origem": "wiki", "rel": "usa"}, {"alvo": "logo_pedra_extraida", "confianca": 1.0, "epistemico": "fato", "origem": "wiki", "rel": "relaciona"}, {"alvo": "area_render_gema", "confianca": 1.0, "epistemico": "fato", "origem": "wiki", "rel": "relaciona"}], "confianca": 1.0, "contraexemplos": [], "descricao": "Cada faceta da gema é um triângulo com uma cor por vértice — profunda no rondiz, clara na mesa e no lume (a direção da luz). A cor no interior é a combinação BARICÊNTRICA a·va+b·vb+c·vc dos vértices: o sombreamento Gouraud. Para manter o VETOR (não uma imagem raster embutida), _tri_baricentrico subdivide cada triângulo numa malha n×n e pinta cada sub-triângulo pela cor baricêntrica do seu centroide. A mesma fórmula baricêntrica dá o ponto e a cor — por isso a cor 'anda junto' com a face. Verificado: cor no vértice = cor do vértice; no centroide = média das três.", "epistemico": "fato", "exemplos": [], "id": "cor_baricentrica_faces", "memoria": "semantica", "meta": {"dominio": "sistema", "fonte": "orbitas vetoriais"}, "origem": "wiki", "palavras_chave": [], "sinonimos": [], "tipo": "conceito", "titulo": "A Cor Baricêntrica nas Faces (Gouraud vetorial)"}
\section{A Cor Baricêntrica nas Faces (Gouraud vetorial)}
Cada faceta da gema é um triângulo com uma cor por vértice — profunda no rondiz, clara na mesa e no lume (a direção da luz). A cor no interior é a combinação BARICÊNTRICA a\textperiodcentered va+b\textperiodcentered vb+c\textperiodcentered vc dos vértices: o sombreamento Gouraud. Para manter o VETOR (não uma imagem raster embutida), \_tri\_baricentrico subdivide cada triângulo numa malha n$\times$n e pinta cada sub-triângulo pela cor baricêntrica do seu centroide. A mesma fórmula baricêntrica dá o ponto e a cor — por isso a cor 'anda junto' com a face. Verificado: cor no vértice = cor do vértice; no centroide = média das três.
\prov{relações: usa design\_vetorizado\_svg; relaciona logo\_pedra\_extraida; relaciona area\_render\_gema}
\prov{proveniência: wiki \textperiodcentered{} orbitas vetoriais \textperiodcentered{} fato \textperiodcentered{} confiança 1.0 \textperiodcentered{} domínio sistema \textperiodcentered{} história 6 versões no jornal}

%CRISTAL {"arestas": [{"alvo": "interlingua", "confianca": 1.0, "epistemico": "fato", "origem": "wiki", "rel": "usa"}, {"alvo": "embedding_como_interlingua", "confianca": 1.0, "epistemico": "fato", "origem": "wiki", "rel": "relaciona"}, {"alvo": "traducao_neural", "confianca": 1.0, "epistemico": "fato", "origem": "wiki", "rel": "relaciona"}, {"alvo": "teorema_geral_traducao", "confianca": 1.0, "epistemico": "fato", "origem": "wiki", "rel": "parte_de"}], "confianca": 1.0, "contraexemplos": [], "descricao": "Com um PIVÔ comum (uma interlíngua), não é preciso um dicionário por par: basta cada simbologia↔pivô (N dicionários) e qualquer par Sᵢ↔Sⱼ compõe-se por Sᵢ→pivô→Sⱼ (via_pivo). É a interlíngua da tradução neural (o espaço de embedding onde tudo cai perto) tornada combinatória: automatizamos o dicionário, e a conversão vira função dele.", "epistemico": "fato", "exemplos": [], "id": "corolario_pivo_interlingua", "memoria": "semantica", "meta": {"dominio": "programacao", "fonte": "linguagem/tradutor.py (via_pivo)"}, "origem": "wiki", "palavras_chave": [], "sinonimos": [], "tipo": "conceito", "titulo": "Corolário do Pivô — N dicionários, N² pares"}
\section{Corolário do Pivô — N dicionários, N\ensuremath{^{2}} pares}
Com um PIVÔ comum (uma interlíngua), não é preciso um dicionário por par: basta cada simbologia\ensuremath{\leftrightarrow}pivô (N dicionários) e qualquer par S\ensuremath{_{i}}\ensuremath{\leftrightarrow}S\ensuremath{_{j}} compõe-se por S\ensuremath{_{i}}\ensuremath{\to}pivô\ensuremath{\to}S\ensuremath{_{j}} (via\_pivo). É a interlíngua da tradução neural (o espaço de embedding onde tudo cai perto) tornada combinatória: automatizamos o dicionário, e a conversão vira função dele.
\prov{relações: usa interlingua; relaciona embedding\_como\_interlingua; relaciona traducao\_neural; parte\_de teorema\_geral\_traducao}
\prov{proveniência: wiki \textperiodcentered{} linguagem/tradutor.py (via\_pivo) \textperiodcentered{} fato \textperiodcentered{} confiança 1.0 \textperiodcentered{} domínio programacao \textperiodcentered{} história 5 versões no jornal}

%CRISTAL {"arestas": [{"alvo": "notacao_matematica", "confianca": 1.0, "epistemico": "fato", "origem": "wiki", "rel": "usa"}, {"alvo": "comando_latex", "confianca": 1.0, "epistemico": "fato", "origem": "wiki", "rel": "usa"}], "confianca": 1.0, "contraexemplos": [], "descricao": "A tabela bilíngue que casa cada construção com sua marcação: ∑→∑, ∫→ınt, ∀→∀, ∈→ın, ℝ→ℝ, xⁿ→xⁿ, aₙ→aₙ, √→√, ≤→≤. O dicionário símbolo↔comando.", "epistemico": "fato", "exemplos": [], "id": "correspondencia_mat_latex", "memoria": "semantica", "meta": {"dominio": "programacao", "fonte": "linguagem/mat_latex.py (_TABELA)"}, "origem": "wiki", "palavras_chave": [], "sinonimos": [], "tipo": "conceito", "titulo": "Correspondência Matemática ↔ LaTeX"}
\section{Correspondência Matemática \ensuremath{\leftrightarrow} LaTeX}
A tabela bilíngue que casa cada construção com sua marcação: \ensuremath{\sum}\ensuremath{\to}\ensuremath{\sum}, \ensuremath{\int}\ensuremath{\to}\ensuremath{\imath}nt, \ensuremath{\forall}\ensuremath{\to}\ensuremath{\forall}, \ensuremath{\in}\ensuremath{\to}\ensuremath{\imath}n, \ensuremath{\mathbb{R}}\ensuremath{\to}\ensuremath{\mathbb{R}}, x\ensuremath{^{n}}\ensuremath{\to}x\ensuremath{^{n}}, a\ensuremath{_{n}}\ensuremath{\to}a\ensuremath{_{n}}, \ensuremath{\sqrt{\,}}\ensuremath{\to}\ensuremath{\sqrt{\,}}, \ensuremath{\le}\ensuremath{\to}\ensuremath{\le}. O dicionário símbolo\ensuremath{\leftrightarrow}comando.
\prov{relações: usa notacao\_matematica; usa comando\_latex}
\prov{proveniência: wiki \textperiodcentered{} linguagem/mat\_latex.py (\_TABELA) \textperiodcentered{} fato \textperiodcentered{} confiança 1.0 \textperiodcentered{} domínio programacao \textperiodcentered{} história 4 versões no jornal}

%CRISTAL {"arestas": [{"alvo": "cifra_mudanca_de_base", "confianca": 1.0, "epistemico": "fato", "origem": "wiki", "rel": "usa"}, {"alvo": "ouro_branco", "confianca": 1.0, "epistemico": "fato", "origem": "wiki", "rel": "relaciona"}, {"alvo": "dimensao_de_prata_pell", "confianca": 1.0, "epistemico": "fato", "origem": "wiki", "rel": "usa"}, {"alvo": "conjugados_galois_parabola", "confianca": 1.0, "epistemico": "fato", "origem": "wiki", "rel": "usa"}], "confianca": 1.0, "contraexemplos": [], "descricao": "Cada Pell emitido recebe cotação: os espectros do trabalho e da chave-ouro entram em parabola_de_ondas, dando a (calor) e c² (fase); peso=c²/(a+c²) e sats_por_pell=pell_val×sats_ref×peso. Não é ticker de mercado — é lastro espectral trabalho↔ouro.", "epistemico": "fato", "exemplos": [], "id": "cotacao_parabola", "memoria": "semantica", "meta": {"autor": "broca-so", "dominio": "operacao_so", "fonte": "goldchain/btc_conectar.py:cotacao_parabola"}, "origem": "wiki", "palavras_chave": [], "sinonimos": [], "tipo": "conceito", "titulo": "Valorização por Parábola a+c²=1"}
\section{Valorização por Parábola a+c\ensuremath{^{2}}=1}
Cada Pell emitido recebe cotação: os espectros do trabalho e da chave-ouro entram em parabola\_de\_ondas, dando a (calor) e c\ensuremath{^{2}} (fase); peso=c\ensuremath{^{2}}/(a+c\ensuremath{^{2}}) e sats\_por\_pell=pell\_val$\times$sats\_ref$\times$peso. Não é ticker de mercado — é lastro espectral trabalho\ensuremath{\leftrightarrow}ouro.
\prov{relações: usa cifra\_mudanca\_de\_base; relaciona ouro\_branco; usa dimensao\_de\_prata\_pell; usa conjugados\_galois\_parabola}
\prov{proveniência: wiki \textperiodcentered{} goldchain/btc\_conectar.py:cotacao\_parabola \textperiodcentered{} fato \textperiodcentered{} confiança 1.0 \textperiodcentered{} domínio operacao\_so \textperiodcentered{} história 6 versões no jornal}

%CRISTAL {"arestas": [{"alvo": "c", "confianca": 1.0, "epistemico": "fato", "origem": "wiki", "rel": "especializa"}, {"alvo": "programacao_orientada_objetos", "confianca": 1.0, "epistemico": "fato", "origem": "wiki", "rel": "usa"}], "confianca": 1.0, "contraexemplos": [], "descricao": "Stroustrup: C com abstração de custo zero — objetos, genéricos, RAII; poder e complexidade em partes iguais.", "epistemico": "fato", "exemplos": [], "id": "cpp", "memoria": "semantica", "meta": {"autor": "Stroustrup", "dominio": "programacao", "fonte": ""}, "origem": "wiki", "palavras_chave": [], "sinonimos": [], "tipo": "conceito", "titulo": "C++"}
\section{C++}
Stroustrup: C com abstração de custo zero — objetos, genéricos, RAII; poder e complexidade em partes iguais.
\prov{relações: especializa c; usa programacao\_orientada\_objetos}
\prov{proveniência: wiki \textperiodcentered{} fato \textperiodcentered{} confiança 1.0 \textperiodcentered{} domínio programacao \textperiodcentered{} história 6 versões no jornal}

%CRISTAL {"arestas": [{"alvo": "seguranca_da_informacao", "confianca": 1.0, "epistemico": "fato", "origem": "wiki", "rel": "usa"}, {"alvo": "cifra", "confianca": 1.0, "epistemico": "fato", "origem": "wiki", "rel": "relaciona"}, {"alvo": "seguranca_estrutural_fractal", "confianca": 1.0, "epistemico": "fato", "origem": "wiki", "rel": "relaciona"}], "confianca": 1.0, "contraexemplos": [], "descricao": "A ciência de proteger informação por transformações matemáticas — cifrar, autenticar, assinar. A base técnica de toda segurança.", "epistemico": "fato", "exemplos": [], "id": "criptografia", "memoria": "semantica", "meta": {"dominio": "programacao", "fonte": ""}, "origem": "wiki", "palavras_chave": [], "sinonimos": [], "tipo": "conceito", "titulo": "Criptografia"}
\section{Criptografia}
A ciência de proteger informação por transformações matemáticas — cifrar, autenticar, assinar. A base técnica de toda segurança.
\prov{relações: usa seguranca\_da\_informacao; relaciona cifra; relaciona seguranca\_estrutural\_fractal}
\prov{proveniência: wiki \textperiodcentered{} fato \textperiodcentered{} confiança 1.0 \textperiodcentered{} domínio programacao \textperiodcentered{} história 4 versões no jornal}

%CRISTAL {"arestas": [{"alvo": "criptografia", "confianca": 1.0, "epistemico": "fato", "origem": "wiki", "rel": "parte_de"}, {"alvo": "assinatura_digital", "confianca": 1.0, "epistemico": "fato", "origem": "wiki", "rel": "usa"}], "confianca": 1.0, "contraexemplos": [], "descricao": "Um par de chaves: a pública cifra/verifica, a privada decifra/assina (RSA, curvas elípticas). Resolve a distribuição de chaves.", "epistemico": "fato", "exemplos": [], "id": "criptografia_assimetrica", "memoria": "semantica", "meta": {"autor": "Rivest/Shamir/Adleman", "dominio": "programacao", "fonte": ""}, "origem": "wiki", "palavras_chave": [], "sinonimos": [], "tipo": "conceito", "titulo": "Criptografia Assimétrica"}
\section{Criptografia Assimétrica}
Um par de chaves: a pública cifra/verifica, a privada decifra/assina (RSA, curvas elípticas). Resolve a distribuição de chaves.
\prov{relações: parte\_de criptografia; usa assinatura\_digital}
\prov{proveniência: wiki \textperiodcentered{} fato \textperiodcentered{} confiança 1.0 \textperiodcentered{} domínio programacao \textperiodcentered{} história 3 versões no jornal}

%CRISTAL {"arestas": [{"alvo": "criptografia", "confianca": 1.0, "epistemico": "fato", "origem": "wiki", "rel": "especializa"}, {"alvo": "prova_de_trabalho", "confianca": 1.0, "epistemico": "fato", "origem": "wiki", "rel": "relaciona"}, {"alvo": "classes_p_np", "confianca": 1.0, "epistemico": "fato", "origem": "wiki", "rel": "usa"}], "confianca": 1.0, "contraexemplos": [], "descricao": "Esquemas que resistem ao computador quântico (que quebraria RSA por Shor) — a fronteira que a prova-de-trabalho por formas de Maass do broca-so persegue.", "epistemico": "inferencia", "exemplos": [], "id": "criptografia_pos_quantica", "memoria": "semantica", "meta": {"dominio": "programacao", "fonte": ""}, "origem": "wiki", "palavras_chave": [], "sinonimos": [], "tipo": "conceito", "titulo": "Criptografia Pós-Quântica"}
\section{Criptografia Pós-Quântica}
Esquemas que resistem ao computador quântico (que quebraria RSA por Shor) — a fronteira que a prova-de-trabalho por formas de Maass do broca-so persegue.
\prov{relações: especializa criptografia; relaciona prova\_de\_trabalho; usa classes\_p\_np}
\prov{proveniência: wiki \textperiodcentered{} inferencia \textperiodcentered{} confiança 1.0 \textperiodcentered{} domínio programacao \textperiodcentered{} história 4 versões no jornal}

%CRISTAL {"arestas": [{"alvo": "criptografia", "confianca": 1.0, "epistemico": "fato", "origem": "wiki", "rel": "parte_de"}], "confianca": 1.0, "contraexemplos": [], "descricao": "Uma única chave secreta cifra e decifra (AES) — rápida, mas exige compartilhar a chave por um canal seguro.", "epistemico": "fato", "exemplos": [], "id": "criptografia_simetrica", "memoria": "semantica", "meta": {"dominio": "programacao", "fonte": ""}, "origem": "wiki", "palavras_chave": [], "sinonimos": [], "tipo": "conceito", "titulo": "Criptografia Simétrica"}
\section{Criptografia Simétrica}
Uma única chave secreta cifra e decifra (AES) — rápida, mas exige compartilhar a chave por um canal seguro.
\prov{relações: parte\_de criptografia}
\prov{proveniência: wiki \textperiodcentered{} fato \textperiodcentered{} confiança 1.0 \textperiodcentered{} domínio programacao \textperiodcentered{} história 2 versões no jornal}

%CRISTAL {"arestas": [{"alvo": "cristal_cli", "confianca": 1.0, "epistemico": "fato", "origem": "wiki", "rel": "usa"}, {"alvo": "peirce_celulas_canais", "confianca": 1.0, "epistemico": "fato", "origem": "wiki", "rel": "explica"}, {"alvo": "orquestrador_roteador_regex", "confianca": 1.0, "epistemico": "fato", "origem": "wiki", "rel": "usa"}], "confianca": 1.0, "contraexemplos": [], "descricao": "Um conjunto fixo de agentes com papéis disjuntos (software/paper/pesquisa/qa/devops...) cujo route_message delega ao orquestrador.roteia_agente — clones idempotentes de papel único, sem sobreposição. É a decomposição de Peirce (células/canais) em código.", "epistemico": "fato", "exemplos": [], "id": "cristal_agentes_ortogonais", "memoria": "semantica", "meta": {"autor": "broca-so", "dominio": "operacao_so", "fonte": "code/cristal/agents.py"}, "origem": "wiki", "palavras_chave": [], "sinonimos": [], "tipo": "conceito", "titulo": "Agentes Internos do Cristal (clones ortogonais)"}
\section{Agentes Internos do Cristal (clones ortogonais)}
Um conjunto fixo de agentes com papéis disjuntos (software/paper/pesquisa/qa/devops...) cujo route\_message delega ao orquestrador.roteia\_agente — clones idempotentes de papel único, sem sobreposição. É a decomposição de Peirce (células/canais) em código.
\prov{relações: usa cristal\_cli; explica peirce\_celulas\_canais; usa orquestrador\_roteador\_regex}
\prov{proveniência: wiki \textperiodcentered{} code/cristal/agents.py \textperiodcentered{} fato \textperiodcentered{} confiança 1.0 \textperiodcentered{} domínio operacao\_so \textperiodcentered{} história 4 versões no jornal}

%CRISTAL {"arestas": [{"alvo": "cristal_cli", "confianca": 1.0, "epistemico": "fato", "origem": "wiki", "rel": "usa"}, {"alvo": "tiffany_home", "confianca": 1.0, "epistemico": "fato", "origem": "wiki", "rel": "usa"}], "confianca": 1.0, "contraexemplos": [], "descricao": "code/cristal: find_project sobe a árvore até achar cristal.json, carrega nome/tiffany_home/agente/branch, e o home do projeto prevalece — a identidade da instância é lida do disco, não passada por flag. É como o cristal 'acorda' e assume a casa.", "epistemico": "fato", "exemplos": [], "id": "cristal_carrega_identidade", "memoria": "semantica", "meta": {"autor": "broca-so", "dominio": "operacao_so", "fonte": "code/cristal/project/context.py"}, "origem": "wiki", "palavras_chave": [], "sinonimos": [], "tipo": "conceito", "titulo": "O Cristal Descobre a Identidade no Disco"}
\section{O Cristal Descobre a Identidade no Disco}
code/cristal: find\_project sobe a árvore até achar cristal.json, carrega nome/tiffany\_home/agente/branch, e o home do projeto prevalece — a identidade da instância é lida do disco, não passada por flag. É como o cristal 'acorda' e assume a casa.
\prov{relações: usa cristal\_cli; usa tiffany\_home}
\prov{proveniência: wiki \textperiodcentered{} code/cristal/project/context.py \textperiodcentered{} fato \textperiodcentered{} confiança 1.0 \textperiodcentered{} domínio operacao\_so \textperiodcentered{} história 3 versões no jornal}

%CRISTAL {"arestas": [{"alvo": "cifra_de_cristal", "confianca": 1.0, "epistemico": "fato", "origem": "wiki", "rel": "parte_de"}, {"alvo": "algebra_de_peirce", "confianca": 1.0, "epistemico": "fato", "origem": "wiki", "rel": "usa"}, {"alvo": "dureza_idempotente_fatora", "confianca": 1.0, "epistemico": "fato", "origem": "wiki", "rel": "usa"}, {"alvo": "cristalchain", "confianca": 1.0, "epistemico": "fato", "origem": "wiki", "rel": "explica"}], "confianca": 1.0, "contraexemplos": [], "descricao": "Por que 'de cristal': um PRIMO p faz de ℤ/p um corpo, sem idempotente não-trivial — só {0,1}, uma peça inteiriça que não racha, inquebrável. Um COMPOSTO é VIDRO: tem 2^ω idempotentes internos (a decomposição de Peirce/CRT), e cada fratura é um fator. A cifra de cristal usa um módulo primo (ou produto secreto de primos): pedra lisa por fora, o corte escondido na lemniscata.", "epistemico": "fato", "exemplos": [], "id": "cristal_e_vidro", "memoria": "semantica", "meta": {"dominio": "criptografia", "fonte": "cifra de cristal"}, "origem": "wiki", "palavras_chave": [], "sinonimos": [], "tipo": "conceito", "titulo": "Cristal e Vidro (o primo não racha)"}
\section{Cristal e Vidro (o primo não racha)}
Por que 'de cristal': um PRIMO p faz de \ensuremath{\mathbb{Z}}/p um corpo, sem idempotente não-trivial — só \{0,1\}, uma peça inteiriça que não racha, inquebrável. Um COMPOSTO é VIDRO: tem 2\^{}\ensuremath{\omega} idempotentes internos (a decomposição de Peirce/CRT), e cada fratura é um fator. A cifra de cristal usa um módulo primo (ou produto secreto de primos): pedra lisa por fora, o corte escondido na lemniscata.
\prov{relações: parte\_de cifra\_de\_cristal; usa algebra\_de\_peirce; usa dureza\_idempotente\_fatora; explica cristalchain}
\prov{proveniência: wiki \textperiodcentered{} cifra de cristal \textperiodcentered{} fato \textperiodcentered{} confiança 1.0 \textperiodcentered{} domínio criptografia \textperiodcentered{} história 5 versões no jornal}

%CRISTAL {"arestas": [{"alvo": "transformada_mineral_7d", "confianca": 1.0, "epistemico": "fato", "origem": "wiki", "rel": "usa"}, {"alvo": "invariante_estelar", "confianca": 1.0, "epistemico": "fato", "origem": "wiki", "rel": "usa"}, {"alvo": "atomo_mineral", "confianca": 1.0, "epistemico": "fato", "origem": "wiki", "rel": "relaciona"}, {"alvo": "computador_mineral", "confianca": 1.0, "epistemico": "fato", "origem": "wiki", "rel": "parte_de"}, {"alvo": "integridade_e_norma_conservada", "confianca": 1.0, "epistemico": "fato", "origem": "wiki", "rel": "relaciona"}], "confianca": 1.0, "contraexemplos": [], "descricao": "Um cristal (célula octoniônica da transformada mineral 7D) é uma CÉLULA DE MEMÓRIA: segura o dado, e a NORMA conservada (‖c‖=‖m‖) é a INTEGRIDADE — corromper a célula quebra a norma, e a máquina detecta. Estável porque é cristal/Pisot (contrai, não vaza). Ler é a transformada inversa. É a conservação a+c²=1 virando bit de paridade.", "epistemico": "fato", "exemplos": [], "id": "cristal_memoria", "memoria": "semantica", "meta": {"dominio": "computacao", "fonte": "computador mineral"}, "origem": "wiki", "palavras_chave": [], "sinonimos": [], "tipo": "conceito", "titulo": "O Cristal como Memória"}
\section{O Cristal como Memória}
Um cristal (célula octoniônica da transformada mineral 7D) é uma CÉLULA DE MEMÓRIA: segura o dado, e a NORMA conservada (\ensuremath{\|}c\ensuremath{\|}=\ensuremath{\|}m\ensuremath{\|}) é a INTEGRIDADE — corromper a célula quebra a norma, e a máquina detecta. Estável porque é cristal/Pisot (contrai, não vaza). Ler é a transformada inversa. É a conservação a+c\ensuremath{^{2}}=1 virando bit de paridade.
\prov{relações: usa transformada\_mineral\_7d; usa invariante\_estelar; relaciona atomo\_mineral; parte\_de computador\_mineral; relaciona integridade\_e\_norma\_conservada}
\prov{proveniência: wiki \textperiodcentered{} computador mineral \textperiodcentered{} fato \textperiodcentered{} confiança 1.0 \textperiodcentered{} domínio computacao \textperiodcentered{} história 6 versões no jornal}

%CRISTAL {"arestas": [{"alvo": "cristal_cli", "confianca": 1.0, "epistemico": "fato", "origem": "wiki", "rel": "usa"}, {"alvo": "release_core_p0", "confianca": 1.0, "epistemico": "fato", "origem": "wiki", "rel": "relaciona"}], "confianca": 1.0, "contraexemplos": [], "descricao": "A camada de serviços localiza os binários da engine (tiffany-index, colapso, libtiffany.so) e expõe index_domain/query_domain chamando prosa — o cristal administra conhecimento sem NUNCA alterar o motor. A fronteira motor↔administração define o cristal.", "epistemico": "fato", "exemplos": [], "id": "cristal_motor_congelado", "memoria": "semantica", "meta": {"autor": "broca-so", "dominio": "operacao_so", "fonte": "code/cristal/services.py; commands/deploy.py"}, "origem": "wiki", "palavras_chave": [], "sinonimos": [], "tipo": "conceito", "titulo": "O Cristal Administra, o Motor é Congelado"}
\section{O Cristal Administra, o Motor é Congelado}
A camada de serviços localiza os binários da engine (tiffany-index, colapso, libtiffany.so) e expõe index\_domain/query\_domain chamando prosa — o cristal administra conhecimento sem NUNCA alterar o motor. A fronteira motor\ensuremath{\leftrightarrow}administração define o cristal.
\prov{relações: usa cristal\_cli; relaciona release\_core\_p0}
\prov{proveniência: wiki \textperiodcentered{} code/cristal/services.py; commands/deploy.py \textperiodcentered{} fato \textperiodcentered{} confiança 1.0 \textperiodcentered{} domínio operacao\_so \textperiodcentered{} história 3 versões no jornal}

%CRISTAL {"arestas": [{"alvo": "goldchain", "confianca": 1.0, "epistemico": "fato", "origem": "wiki", "rel": "especializa"}, {"alvo": "corpo_mineral", "confianca": 1.0, "epistemico": "fato", "origem": "wiki", "rel": "usa"}, {"alvo": "numeros_de_pisot", "confianca": 1.0, "epistemico": "fato", "origem": "wiki", "rel": "usa"}, {"alvo": "cristal_e_a_convergencia", "confianca": 1.0, "epistemico": "fato", "origem": "wiki", "rel": "relaciona"}, {"alvo": "cifra_de_cristal", "confianca": 1.0, "epistemico": "fato", "origem": "wiki", "rel": "relaciona"}, {"alvo": "broca-so_cristalizado", "confianca": 1.0, "epistemico": "fato", "origem": "wiki", "rel": "relaciona"}], "confianca": 1.0, "contraexemplos": [], "descricao": "A GoldChain nomeia o ouro (σ_1), UM metal. Mas a mineração, o valor e a imutabilidade não são do ouro — são do CRISTAL (a classe: Pisot, o estado fundamental, o idempotente, a norma conservada). Ouro é um metal entre muitos (a mineração já é multimetal, 12 faixas); cristal é a classe inteira. A cadeia é a cristalização feita economia: cada bloco um cristal minerado a frio, a corrente uma rede que não racha. CristalChain é como a cadeia se chama por dentro. Paper: papers/cristalchain.tex.", "epistemico": "fato", "exemplos": [], "id": "cristalchain", "memoria": "semantica", "meta": {"dominio": "goldchain", "fonte": "cristalchain"}, "origem": "wiki", "palavras_chave": [], "sinonimos": [], "tipo": "conceito", "titulo": "CristalChain (o nome verdadeiro da cadeia)"}
\section{CristalChain (o nome verdadeiro da cadeia)}
A GoldChain nomeia o ouro (\ensuremath{\sigma}\_1), UM metal. Mas a mineração, o valor e a imutabilidade não são do ouro — são do CRISTAL (a classe: Pisot, o estado fundamental, o idempotente, a norma conservada). Ouro é um metal entre muitos (a mineração já é multimetal, 12 faixas); cristal é a classe inteira. A cadeia é a cristalização feita economia: cada bloco um cristal minerado a frio, a corrente uma rede que não racha. CristalChain é como a cadeia se chama por dentro. Paper: papers/cristalchain.tex.
\prov{relações: especializa goldchain; usa corpo\_mineral; usa numeros\_de\_pisot; relaciona cristal\_e\_a\_convergencia; relaciona cifra\_de\_cristal; relaciona broca-so\_cristalizado}
\prov{proveniência: wiki \textperiodcentered{} cristalchain \textperiodcentered{} fato \textperiodcentered{} confiança 1.0 \textperiodcentered{} domínio goldchain \textperiodcentered{} história 7 versões no jornal}

%CRISTAL {"arestas": [{"alvo": "java", "confianca": 1.0, "epistemico": "fato", "origem": "wiki", "rel": "relaciona"}, {"alvo": "maquina_virtual", "confianca": 1.0, "epistemico": "fato", "origem": "wiki", "rel": "usa"}], "confianca": 1.0, "contraexemplos": [], "descricao": "Hejlsberg (Microsoft): a resposta ao Java na plataforma .NET — OOP moderno, LINQ, async; hoje multiplataforma.", "epistemico": "fato", "exemplos": [], "id": "csharp", "memoria": "semantica", "meta": {"autor": "Hejlsberg", "dominio": "programacao", "fonte": ""}, "origem": "wiki", "palavras_chave": [], "sinonimos": [], "tipo": "conceito", "titulo": "C#"}
\section{C\#}
Hejlsberg (Microsoft): a resposta ao Java na plataforma .NET — OOP moderno, LINQ, async; hoje multiplataforma.
\prov{relações: relaciona java; usa maquina\_virtual}
\prov{proveniência: wiki \textperiodcentered{} fato \textperiodcentered{} confiança 1.0 \textperiodcentered{} domínio programacao \textperiodcentered{} história 6 versões no jornal}

%CRISTAL {"arestas": [{"alvo": "separacao_conteudo_forma", "confianca": 1.0, "epistemico": "fato", "origem": "wiki", "rel": "usa"}, {"alvo": "html", "confianca": 1.0, "epistemico": "fato", "origem": "wiki", "rel": "usa"}, {"alvo": "programacao_declarativa", "confianca": 1.0, "epistemico": "fato", "origem": "wiki", "rel": "usa"}], "confianca": 1.0, "contraexemplos": [], "descricao": "Cascading Style Sheets: a linguagem declarativa que dá forma ao HTML — cor, layout, tipografia. A apresentação separada da estrutura.", "epistemico": "fato", "exemplos": [], "id": "css", "memoria": "semantica", "meta": {"dominio": "programacao", "fonte": ""}, "origem": "wiki", "palavras_chave": [], "sinonimos": [], "tipo": "conceito", "titulo": "CSS"}
\section{CSS}
Cascading Style Sheets: a linguagem declarativa que dá forma ao HTML — cor, layout, tipografia. A apresentação separada da estrutura.
\prov{relações: usa separacao\_conteudo\_forma; usa html; usa programacao\_declarativa}
\prov{proveniência: wiki \textperiodcentered{} fato \textperiodcentered{} confiança 1.0 \textperiodcentered{} domínio programacao \textperiodcentered{} história 4 versões no jornal}

%CRISTAL {"arestas": [{"alvo": "css", "confianca": 1.0, "epistemico": "fato", "origem": "wiki", "rel": "usa"}, {"alvo": "componente_ui", "confianca": 1.0, "epistemico": "fato", "origem": "wiki", "rel": "usa"}], "confianca": 1.0, "contraexemplos": [], "descricao": "Escrever estilos dentro do JavaScript, escopados ao componente (styled-components, Emotion) — o estilo viaja junto com a lógica.", "epistemico": "inferencia", "exemplos": [], "id": "css_in_js", "memoria": "semantica", "meta": {"dominio": "programacao", "fonte": ""}, "origem": "wiki", "palavras_chave": [], "sinonimos": [], "tipo": "conceito", "titulo": "CSS-in-JS"}
\section{CSS-in-JS}
Escrever estilos dentro do JavaScript, escopados ao componente (styled-components, Emotion) — o estilo viaja junto com a lógica.
\prov{relações: usa css; usa componente\_ui}
\prov{proveniência: wiki \textperiodcentered{} inferencia \textperiodcentered{} confiança 1.0 \textperiodcentered{} domínio programacao \textperiodcentered{} história 6 versões no jornal}

%CRISTAL {"arestas": [{"alvo": "sistema_de_tipos", "confianca": 1.0, "epistemico": "fato", "origem": "wiki", "rel": "explica"}, {"alvo": "logica", "confianca": 1.0, "epistemico": "fato", "origem": "wiki", "rel": "usa"}, {"alvo": "incompletude_godel", "confianca": 1.0, "epistemico": "fato", "origem": "wiki", "rel": "relaciona"}], "confianca": 1.0, "contraexemplos": [], "descricao": "Proposições SÃO tipos e provas SÃO programas — verificar um tipo é checar uma demonstração. A ponte entre lógica e computação.", "epistemico": "fato", "exemplos": [], "id": "curry_howard", "memoria": "semantica", "meta": {"dominio": "programacao", "fonte": ""}, "origem": "wiki", "palavras_chave": [], "sinonimos": [], "tipo": "conceito", "titulo": "Correspondência de Curry-Howard"}
\section{Correspondência de Curry-Howard}
Proposições SÃO tipos e provas SÃO programas — verificar um tipo é checar uma demonstração. A ponte entre lógica e computação.
\prov{relações: explica sistema\_de\_tipos; usa logica; relaciona incompletude\_godel}
\prov{proveniência: wiki \textperiodcentered{} fato \textperiodcentered{} confiança 1.0 \textperiodcentered{} domínio programacao \textperiodcentered{} história 8 versões no jornal}

%CRISTAL {"arestas": [{"alvo": "jogos_resolvidos", "confianca": 1.0, "epistemico": "fato", "origem": "wiki", "rel": "especializa"}], "confianca": 1.0, "contraexemplos": [], "descricao": "Em 2007, Schaeffer PROVOU que as damas americanas, com jogo perfeito, terminam em empate — o maior jogo 'resolvido' até então. Chinook não perde nunca: o jogo virou matemática fechada.", "epistemico": "fato", "exemplos": [], "id": "damas_resolvidas_chinook", "memoria": "semantica", "meta": {"autor": "Schaeffer", "dominio": "ia", "fonte": "jogos e IA"}, "origem": "wiki", "palavras_chave": [], "sinonimos": [], "tipo": "conceito", "titulo": "Damas Resolvidas (Chinook)"}
\section{Damas Resolvidas (Chinook)}
Em 2007, Schaeffer PROVOU que as damas americanas, com jogo perfeito, terminam em empate — o maior jogo 'resolvido' até então. Chinook não perde nunca: o jogo virou matemática fechada.
\prov{relações: especializa jogos\_resolvidos}
\prov{proveniência: wiki \textperiodcentered{} jogos e IA \textperiodcentered{} fato \textperiodcentered{} confiança 1.0 \textperiodcentered{} domínio ia \textperiodcentered{} história 2 versões no jornal}

%CRISTAL {"arestas": [{"alvo": "seguranca_da_informacao", "confianca": 1.0, "epistemico": "fato", "origem": "wiki", "rel": "usa"}], "confianca": 1.0, "contraexemplos": [], "descricao": "Camadas de defesa independentes, para que furar uma não baste — assumir que cada barreira pode falhar. Zero trust é sua forma moderna.", "epistemico": "inferencia", "exemplos": [], "id": "defesa_em_profundidade", "memoria": "semantica", "meta": {"dominio": "programacao", "fonte": ""}, "origem": "wiki", "palavras_chave": [], "sinonimos": [], "tipo": "conceito", "titulo": "Defesa em Profundidade"}
\section{Defesa em Profundidade}
Camadas de defesa independentes, para que furar uma não baste — assumir que cada barreira pode falhar. Zero trust é sua forma moderna.
\prov{relações: usa seguranca\_da\_informacao}
\prov{proveniência: wiki \textperiodcentered{} inferencia \textperiodcentered{} confiança 1.0 \textperiodcentered{} domínio programacao \textperiodcentered{} história 2 versões no jornal}

%CRISTAL {"arestas": [{"alvo": "exemplo", "confianca": 1.0, "epistemico": "fato", "origem": "curriculo", "rel": "confirma"}, {"alvo": "semantica", "confianca": 1.0, "epistemico": "fato", "origem": "curriculo", "rel": "parte_de"}, {"alvo": "word", "confianca": 0.5, "epistemico": "fato", "origem": "wiki", "rel": "relaciona"}], "confianca": 1.0, "contraexemplos": [], "descricao": "Enunciado que fixa o sentido de um termo; condições necessárias e suficientes.", "epistemico": "fato", "exemplos": [], "id": "definicao", "memoria": "semantica", "meta": {"dominio": "linguagem", "nivel": 3, "ordem": 6}, "origem": "curriculo", "palavras_chave": ["gênero", "diferença"], "sinonimos": ["definição"], "tipo": "conceito", "titulo": "definicao"}
\section{definicao}
Enunciado que fixa o sentido de um termo; condições necessárias e suficientes.
\prov{sinónimos: definição}
\prov{relações: confirma exemplo; parte\_de semantica; relaciona word}
\prov{proveniência: curriculo \textperiodcentered{} fato \textperiodcentered{} confiança 1.0 \textperiodcentered{} domínio linguagem \textperiodcentered{} história 56 versões no jornal}

%CRISTAL {"arestas": [{"alvo": "estrutura_definicao_teorema_prova", "confianca": 1.0, "epistemico": "fato", "origem": "wiki", "rel": "parte_de"}, {"alvo": "a_definicao", "confianca": 1.0, "epistemico": "fato", "origem": "wiki", "rel": "relaciona"}], "confianca": 1.0, "contraexemplos": [], "descricao": "Fixar o sentido exato de um termo a partir de termos já definidos — o ponto de partida que torna o resto demonstrável. Onde a ambiguidade morre.", "epistemico": "fato", "exemplos": [], "id": "definicao_matematica", "memoria": "semantica", "meta": {"dominio": "programacao", "fonte": ""}, "origem": "wiki", "palavras_chave": [], "sinonimos": [], "tipo": "conceito", "titulo": "Definição"}
\section{Definição}
Fixar o sentido exato de um termo a partir de termos já definidos — o ponto de partida que torna o resto demonstrável. Onde a ambiguidade morre.
\prov{relações: parte\_de estrutura\_definicao\_teorema\_prova; relaciona a\_definicao}
\prov{proveniência: wiki \textperiodcentered{} fato \textperiodcentered{} confiança 1.0 \textperiodcentered{} domínio programacao \textperiodcentered{} história 3 versões no jornal}

%CRISTAL {"arestas": [{"alvo": "operacao_broca_so", "confianca": 1.0, "epistemico": "fato", "origem": "wiki", "rel": "parte_de"}, {"alvo": "cifra", "confianca": 1.0, "epistemico": "fato", "origem": "wiki", "rel": "usa"}, {"alvo": "soberania_da_chave", "confianca": 1.0, "epistemico": "fato", "origem": "wiki", "rel": "relaciona"}], "confianca": 1.0, "contraexemplos": [], "descricao": "Nenhuma chave por conta é armazenada: um DNA (string) vira chave via exp=int(sha256(dna)) mod N (BTC) ou Account.from_key(sha256(dna)) (ETH). Só a chave central real é lida de disco fora do repo. A chave é derivada, não guardada.", "epistemico": "fato", "exemplos": [], "id": "derive_dont_store", "memoria": "semantica", "meta": {"autor": "broca-so", "dominio": "operacao_so", "fonte": "goldchain/contas_btc.py; contas_eth.py; conta.py"}, "origem": "wiki", "palavras_chave": [], "sinonimos": [], "tipo": "conceito", "titulo": "Contas Derivadas (derive-don't-store)"}
\section{Contas Derivadas (derive-don't-store)}
Nenhuma chave por conta é armazenada: um DNA (string) vira chave via exp=int(sha256(dna)) mod N (BTC) ou Account.from\_key(sha256(dna)) (ETH). Só a chave central real é lida de disco fora do repo. A chave é derivada, não guardada.
\prov{relações: parte\_de operacao\_broca\_so; usa cifra; relaciona soberania\_da\_chave}
\prov{proveniência: wiki \textperiodcentered{} goldchain/contas\_btc.py; contas\_eth.py; conta.py \textperiodcentered{} fato \textperiodcentered{} confiança 1.0 \textperiodcentered{} domínio operacao\_so \textperiodcentered{} história 4 versões no jornal}

%CRISTAL {"arestas": [{"alvo": "programacao", "confianca": 1.0, "epistemico": "fato", "origem": "wiki", "rel": "parte_de"}, {"alvo": "internet", "confianca": 1.0, "epistemico": "fato", "origem": "wiki", "rel": "usa"}], "confianca": 1.0, "contraexemplos": [], "descricao": "Construir aplicações que rodam no navegador e no servidor — a pilha HTML/CSS/JS estendida por frameworks, ferramentas e serviços.", "epistemico": "fato", "exemplos": [], "id": "desenvolvimento_web", "memoria": "semantica", "meta": {"dominio": "programacao", "fonte": ""}, "origem": "wiki", "palavras_chave": [], "sinonimos": [], "tipo": "conceito", "titulo": "Desenvolvimento Web"}
\section{Desenvolvimento Web}
Construir aplicações que rodam no navegador e no servidor — a pilha HTML/CSS/JS estendida por frameworks, ferramentas e serviços.
\prov{relações: parte\_de programacao; usa internet}
\prov{proveniência: wiki \textperiodcentered{} fato \textperiodcentered{} confiança 1.0 \textperiodcentered{} domínio programacao \textperiodcentered{} história 6 versões no jornal}

%CRISTAL {"arestas": [{"alvo": "design_system_cristalchain", "confianca": 1.0, "epistemico": "fato", "origem": "wiki", "rel": "especializa"}, {"alvo": "logo_pedra_extraida", "confianca": 1.0, "epistemico": "fato", "origem": "wiki", "rel": "usa"}, {"alvo": "cristalchain", "confianca": 1.0, "epistemico": "fato", "origem": "wiki", "rel": "parte_de"}, {"alvo": "turmalina_paraiba_ancora", "confianca": 1.0, "epistemico": "fato", "origem": "wiki", "rel": "usa"}], "confianca": 1.0, "contraexemplos": [], "descricao": "O design system que o ChatGPT enviou (model/design) foi lido e reproduzido: a PALETA OFICIAL sincronizada (#0E9CA6 azul-turmalina, #0A5E66 profundo, #08363B petróleo, #B87333 cobre, #D9A066 cobre-claro, #C7CDD1 cinza-mineral, gelo/branco), e as 14+ seções EXTRAÍDAS (bordas, divisores, cantoneiras, molduras, ícones, capitulares, numeração, boxes, selos, texturas, vinhetas, assinatura, versões) em assets/cristalchain/elementos/. Implementado no PDF pelo design system editorial papers/cristalchain.sty (tokens + componentes: capa, seção, boxes, nota, certificado) — a 1ª aplicação é o turmalina.tex refinado (paper + emissão + certificado). O cobre só como destaque.", "epistemico": "fato", "exemplos": [], "id": "design_elementos_do_chatgpt", "memoria": "semantica", "meta": {"dominio": "sistema", "fonte": "design elementos"}, "origem": "wiki", "palavras_chave": [], "sinonimos": [], "tipo": "conceito", "titulo": "Os Elementos do Design (extraídos e no PDF)"}
\section{Os Elementos do Design (extraídos e no PDF)}
O design system que o ChatGPT enviou (model/design) foi lido e reproduzido: a PALETA OFICIAL sincronizada (\#0E9CA6 azul-turmalina, \#0A5E66 profundo, \#08363B petróleo, \#B87333 cobre, \#D9A066 cobre-claro, \#C7CDD1 cinza-mineral, gelo/branco), e as 14+ seções EXTRAÍDAS (bordas, divisores, cantoneiras, molduras, ícones, capitulares, numeração, boxes, selos, texturas, vinhetas, assinatura, versões) em assets/cristalchain/elementos/. Implementado no PDF pelo design system editorial papers/cristalchain.sty (tokens + componentes: capa, seção, boxes, nota, certificado) — a 1ª aplicação é o turmalina.tex refinado (paper + emissão + certificado). O cobre só como destaque.
\prov{relações: especializa design\_system\_cristalchain; usa logo\_pedra\_extraida; parte\_de cristalchain; usa turmalina\_paraiba\_ancora}
\prov{proveniência: wiki \textperiodcentered{} design elementos \textperiodcentered{} fato \textperiodcentered{} confiança 1.0 \textperiodcentered{} domínio sistema \textperiodcentered{} história 5 versões no jornal}

%CRISTAL {"arestas": [{"alvo": "fechar_laco_gato_posicao", "confianca": 1.0, "epistemico": "fato", "origem": "wiki", "rel": "usa"}, {"alvo": "algebra_posicional_unificada", "confianca": 1.0, "epistemico": "fato", "origem": "wiki", "rel": "usa"}, {"alvo": "programacao_declarativa", "confianca": 1.0, "epistemico": "fato", "origem": "wiki", "rel": "relaciona"}, {"alvo": "recursao_linguistica", "confianca": 1.0, "epistemico": "fato", "origem": "wiki", "rel": "contradiz"}], "confianca": 1.0, "contraexemplos": [], "descricao": "Onde a linguagem imperativa itera e a funcional recorre, a de gatos DOBRA o laço numa única posição: repetir k vezes = elevar a matriz à potência k (Aₙk), calculável em O(log k) por exponenciação rápida. O laço deixa de ser controle de fluxo e vira forma fechada — declarativo no osso.", "epistemico": "inferencia", "exemplos": [], "id": "design_laco_e_posicao", "memoria": "semantica", "meta": {"dominio": "programacao", "fonte": ""}, "origem": "wiki", "palavras_chave": [], "sinonimos": [], "tipo": "conceito", "titulo": "Laço como Posição — a iteração dobrada em Aₙk"}
\section{Laço como Posição — a iteração dobrada em A\ensuremath{_{n}}k}
Onde a linguagem imperativa itera e a funcional recorre, a de gatos DOBRA o laço numa única posição: repetir k vezes = elevar a matriz à potência k (A\ensuremath{_{n}}k), calculável em O(log k) por exponenciação rápida. O laço deixa de ser controle de fluxo e vira forma fechada — declarativo no osso.
\prov{relações: usa fechar\_laco\_gato\_posicao; usa algebra\_posicional\_unificada; relaciona programacao\_declarativa; contradiz recursao\_linguistica}
\prov{proveniência: wiki \textperiodcentered{} inferencia \textperiodcentered{} confiança 1.0 \textperiodcentered{} domínio programacao \textperiodcentered{} história 12 versões no jornal}

%CRISTAL {"arestas": [{"alvo": "linguagem_matriz_gatos", "confianca": 1.0, "epistemico": "fato", "origem": "wiki", "rel": "explica"}, {"alvo": "programacao", "confianca": 1.0, "epistemico": "fato", "origem": "wiki", "rel": "parte_de"}, {"alvo": "design_primitivo_matriz", "confianca": 1.0, "epistemico": "fato", "origem": "wiki", "rel": "usa"}, {"alvo": "design_laco_e_posicao", "confianca": 1.0, "epistemico": "fato", "origem": "wiki", "rel": "usa"}, {"alvo": "design_sub_turing_deliberado", "confianca": 1.0, "epistemico": "fato", "origem": "wiki", "rel": "usa"}, {"alvo": "design_norma_conservada", "confianca": 1.0, "epistemico": "fato", "origem": "wiki", "rel": "usa"}, {"alvo": "design_semantica_tensorial", "confianca": 1.0, "epistemico": "fato", "origem": "wiki", "rel": "usa"}, {"alvo": "design_opcode_algebrico", "confianca": 1.0, "epistemico": "fato", "origem": "wiki", "rel": "usa"}, {"alvo": "design_programa_como_numero", "confianca": 1.0, "epistemico": "fato", "origem": "wiki", "rel": "usa"}], "confianca": 1.0, "contraexemplos": [], "descricao": "A linguagem de gatos é um ponto de design coerente e quase sempre OPOSTO ao default moderno: primitivo-matriz (não escalar), laço-posição (não fluxo), sub-Turing-por-escolha (não Turing-default), norma-conservada (não efeitos), semântica-tensorial (não tipos), opcode-algébrico (não mnemônico), programa-número (não texto). Tudo herdado da álgebra dos metálicos e compilando pro metal — não é uma linguagem 'a menos', é outra base para o que computar significa.", "epistemico": "inferencia", "exemplos": [], "id": "design_linguagem_gatos", "memoria": "semantica", "meta": {"dominio": "programacao", "fonte": ""}, "origem": "wiki", "palavras_chave": [], "sinonimos": [], "tipo": "conceito", "titulo": "O Design da Linguagem de Gatos (síntese)"}
\section{O Design da Linguagem de Gatos (síntese)}
A linguagem de gatos é um ponto de design coerente e quase sempre OPOSTO ao default moderno: primitivo-matriz (não escalar), laço-posição (não fluxo), sub-Turing-por-escolha (não Turing-default), norma-conservada (não efeitos), semântica-tensorial (não tipos), opcode-algébrico (não mnemônico), programa-número (não texto). Tudo herdado da álgebra dos metálicos e compilando pro metal — não é uma linguagem 'a menos', é outra base para o que computar significa.
\prov{relações: explica linguagem\_matriz\_gatos; parte\_de programacao; usa design\_primitivo\_matriz; usa design\_laco\_e\_posicao; usa design\_sub\_turing\_deliberado; usa design\_norma\_conservada; usa design\_semantica\_tensorial; usa design\_opcode\_algebrico; usa design\_programa\_como\_numero}
\prov{proveniência: wiki \textperiodcentered{} inferencia \textperiodcentered{} confiança 1.0 \textperiodcentered{} domínio programacao \textperiodcentered{} história 20 versões no jornal}

%CRISTAL {"arestas": [{"alvo": "gato", "confianca": 1.0, "epistemico": "fato", "origem": "wiki", "rel": "usa"}, {"alvo": "computacao_reversivel", "confianca": 1.0, "epistemico": "fato", "origem": "wiki", "rel": "relaciona"}, {"alvo": "monada", "confianca": 1.0, "epistemico": "fato", "origem": "wiki", "rel": "contradiz"}, {"alvo": "ownership_borrow", "confianca": 1.0, "epistemico": "fato", "origem": "wiki", "rel": "relaciona"}], "confianca": 1.0, "contraexemplos": [], "descricao": "Cada passo da linguagem de gatos preserva a norma det(Aₙ)=−1 (a norma metálica ≡±1). Onde as modernas GERENCIAM efeitos (mônadas, ownership), a de gatos torna a correção uma LEI ALGÉBRICA conservada — parente da computação reversível, onde nada se apaga.", "epistemico": "inferencia", "exemplos": [], "id": "design_norma_conservada", "memoria": "semantica", "meta": {"dominio": "programacao", "fonte": ""}, "origem": "wiki", "palavras_chave": [], "sinonimos": [], "tipo": "conceito", "titulo": "Norma Conservada — a computação preserva det=−1"}
\section{Norma Conservada — a computação preserva det=\ensuremath{-}1}
Cada passo da linguagem de gatos preserva a norma det(A\ensuremath{_{n}})=\ensuremath{-}1 (a norma metálica \ensuremath{\equiv}$\pm$1). Onde as modernas GERENCIAM efeitos (mônadas, ownership), a de gatos torna a correção uma LEI ALGÉBRICA conservada — parente da computação reversível, onde nada se apaga.
\prov{relações: usa gato; relaciona computacao\_reversivel; contradiz monada; relaciona ownership\_borrow}
\prov{proveniência: wiki \textperiodcentered{} inferencia \textperiodcentered{} confiança 1.0 \textperiodcentered{} domínio programacao \textperiodcentered{} história 10 versões no jornal}

%CRISTAL {"arestas": [{"alvo": "backend_isa12_broca", "confianca": 1.0, "epistemico": "fato", "origem": "wiki", "rel": "usa"}, {"alvo": "gato", "confianca": 1.0, "epistemico": "fato", "origem": "wiki", "rel": "usa"}, {"alvo": "linguagem_assembly", "confianca": 1.0, "epistemico": "fato", "origem": "wiki", "rel": "contradiz"}, {"alvo": "bytecode", "confianca": 1.0, "epistemico": "fato", "origem": "wiki", "rel": "usa"}], "confianca": 1.0, "contraexemplos": [], "descricao": "As instruções OURO/PRATA/BRONZE não são mnemônicos arbitrários: são o gato (A₁,A₂,A₃) usado como opcode — um objeto algébrico com significado, que compila para a ISA-12 real (bytecode idêntico ao broca_asm em C). O metal não é metáfora.", "epistemico": "inferencia", "exemplos": [], "id": "design_opcode_algebrico", "memoria": "semantica", "meta": {"dominio": "programacao", "fonte": ""}, "origem": "wiki", "palavras_chave": [], "sinonimos": [], "tipo": "conceito", "titulo": "Opcode Algébrico — o gato OURO/PRATA/BRONZE como instrução"}
\section{Opcode Algébrico — o gato OURO/PRATA/BRONZE como instrução}
As instruções OURO/PRATA/BRONZE não são mnemônicos arbitrários: são o gato (A\ensuremath{_{1}},A\ensuremath{_{2}},A\ensuremath{_{3}}) usado como opcode — um objeto algébrico com significado, que compila para a ISA-12 real (bytecode idêntico ao broca\_asm em C). O metal não é metáfora.
\prov{relações: usa backend\_isa12\_broca; usa gato; contradiz linguagem\_assembly; usa bytecode}
\prov{proveniência: wiki \textperiodcentered{} inferencia \textperiodcentered{} confiança 1.0 \textperiodcentered{} domínio programacao \textperiodcentered{} história 10 versões no jornal}

%CRISTAL {"arestas": [{"alvo": "gato_2x2_e_fisico_nao_incidental", "confianca": 1.0, "epistemico": "fato", "origem": "wiki", "rel": "usa"}, {"alvo": "computacao_matricial", "confianca": 1.0, "epistemico": "fato", "origem": "wiki", "rel": "usa"}, {"alvo": "linguagens_de_array", "confianca": 1.0, "epistemico": "fato", "origem": "wiki", "rel": "relaciona"}, {"alvo": "linguagem_matriz_gatos", "confianca": 1.0, "epistemico": "fato", "origem": "wiki", "rel": "parte_de"}], "confianca": 1.0, "contraexemplos": [], "descricao": "O átomo de cálculo da linguagem de gatos é um gato — uma matriz 2×2 (companion Aₙ), e computar é MULTIPLICAR matrizes. As linguagens comuns partem do escalar/palavra/objeto; as de array (APL) chegam perto, mas o gato leva ao extremo: a instrução JÁ é uma matriz.", "epistemico": "inferencia", "exemplos": [], "id": "design_primitivo_matriz", "memoria": "semantica", "meta": {"dominio": "programacao", "fonte": ""}, "origem": "wiki", "palavras_chave": [], "sinonimos": [], "tipo": "conceito", "titulo": "Primitivo Matriz — o átomo de cálculo é um gato 2×2"}
\section{Primitivo Matriz — o átomo de cálculo é um gato 2$\times$2}
O átomo de cálculo da linguagem de gatos é um gato — uma matriz 2$\times$2 (companion A\ensuremath{_{n}}), e computar é MULTIPLICAR matrizes. As linguagens comuns partem do escalar/palavra/objeto; as de array (APL) chegam perto, mas o gato leva ao extremo: a instrução JÁ é uma matriz.
\prov{relações: usa gato\_2x2\_e\_fisico\_nao\_incidental; usa computacao\_matricial; relaciona linguagens\_de\_array; parte\_de linguagem\_matriz\_gatos}
\prov{proveniência: wiki \textperiodcentered{} inferencia \textperiodcentered{} confiança 1.0 \textperiodcentered{} domínio programacao \textperiodcentered{} história 10 versões no jornal}

%CRISTAL {"arestas": [{"alvo": "algebra_posicional_unificada", "confianca": 1.0, "epistemico": "fato", "origem": "wiki", "rel": "usa"}, {"alvo": "homoiconicidade", "confianca": 1.0, "epistemico": "fato", "origem": "wiki", "rel": "relaciona"}, {"alvo": "beta_numeracao_metalica", "confianca": 1.0, "epistemico": "fato", "origem": "wiki", "rel": "usa"}], "confianca": 1.0, "contraexemplos": [], "descricao": "O estado da linguagem de gatos é uma posição, e a posição é um DÍGITO na β-numeração de base σ_n (Pisot). O programa é, literalmente, um número escrito na base metálica — parente da homoiconicidade (código=dado), mas aritmético, não sintático.", "epistemico": "inferencia", "exemplos": [], "id": "design_programa_como_numero", "memoria": "semantica", "meta": {"dominio": "programacao", "fonte": ""}, "origem": "wiki", "palavras_chave": [], "sinonimos": [], "tipo": "conceito", "titulo": "Programa como Número — a posição é um dígito na base de Pisot"}
\section{Programa como Número — a posição é um dígito na base de Pisot}
O estado da linguagem de gatos é uma posição, e a posição é um DÍGITO na \ensuremath{\beta}-numeração de base \ensuremath{\sigma}\_n (Pisot). O programa é, literalmente, um número escrito na base metálica — parente da homoiconicidade (código=dado), mas aritmético, não sintático.
\prov{relações: usa algebra\_posicional\_unificada; relaciona homoiconicidade; usa beta\_numeracao\_metalica}
\prov{proveniência: wiki \textperiodcentered{} inferencia \textperiodcentered{} confiança 1.0 \textperiodcentered{} domínio programacao \textperiodcentered{} história 8 versões no jornal}

%CRISTAL {"arestas": [{"alvo": "contracao_einstein_k_tensorial", "confianca": 1.0, "epistemico": "fato", "origem": "wiki", "rel": "usa"}, {"alvo": "gato_generaliza_matriz_companion", "confianca": 1.0, "epistemico": "fato", "origem": "wiki", "rel": "usa"}, {"alvo": "sistema_de_tipos", "confianca": 1.0, "epistemico": "fato", "origem": "wiki", "rel": "relaciona"}, {"alvo": "curry_howard", "confianca": 1.0, "epistemico": "fato", "origem": "wiki", "rel": "relaciona"}], "confianca": 1.0, "contraexemplos": [], "descricao": "Compor gatos é CONTRAIR tensores (a soma de Einstein sobre as constantes de estrutura k-tensoriais). Onde as modernas checam tipos ou compõem funções, a de gatos compõe por álgebra multilinear — a mesma que generaliza o gato 2×2 à matriz de gatos de grau k.", "epistemico": "inferencia", "exemplos": [], "id": "design_semantica_tensorial", "memoria": "semantica", "meta": {"dominio": "programacao", "fonte": ""}, "origem": "wiki", "palavras_chave": [], "sinonimos": [], "tipo": "conceito", "titulo": "Semântica Tensorial — compor gatos é contrair tensores"}
\section{Semântica Tensorial — compor gatos é contrair tensores}
Compor gatos é CONTRAIR tensores (a soma de Einstein sobre as constantes de estrutura k-tensoriais). Onde as modernas checam tipos ou compõem funções, a de gatos compõe por álgebra multilinear — a mesma que generaliza o gato 2$\times$2 à matriz de gatos de grau k.
\prov{relações: usa contracao\_einstein\_k\_tensorial; usa gato\_generaliza\_matriz\_companion; relaciona sistema\_de\_tipos; relaciona curry\_howard}
\prov{proveniência: wiki \textperiodcentered{} inferencia \textperiodcentered{} confiança 1.0 \textperiodcentered{} domínio programacao \textperiodcentered{} história 10 versões no jornal}

%CRISTAL {"arestas": [{"alvo": "expressividade_sub_turing", "confianca": 1.0, "epistemico": "fato", "origem": "wiki", "rel": "usa"}, {"alvo": "linguagens_totais", "confianca": 1.0, "epistemico": "fato", "origem": "wiki", "rel": "relaciona"}, {"alvo": "hierarquia_de_chomsky", "confianca": 1.0, "epistemico": "fato", "origem": "wiki", "rel": "usa"}, {"alvo": "maquina_de_contadores", "confianca": 1.0, "epistemico": "fato", "origem": "wiki", "rel": "usa"}], "confianca": 1.0, "contraexemplos": [], "descricao": "O núcleo da linguagem de gatos é ESTRITAMENTE sub-Turing (endereços estáticos, sempre termina), como as linguagens totais — e só transcende a Turing-completude quando se adiciona indireção posicional (a máquina de contadores de Minsky). A completude é um BOTÃO, não o default: a hierarquia de Chomsky exposta como decisão.", "epistemico": "inferencia", "exemplos": [], "id": "design_sub_turing_deliberado", "memoria": "semantica", "meta": {"dominio": "programacao", "fonte": ""}, "origem": "wiki", "palavras_chave": [], "sinonimos": [], "tipo": "conceito", "titulo": "Sub-Turing por Escolha — a completude é um botão"}
\section{Sub-Turing por Escolha — a completude é um botão}
O núcleo da linguagem de gatos é ESTRITAMENTE sub-Turing (endereços estáticos, sempre termina), como as linguagens totais — e só transcende a Turing-completude quando se adiciona indireção posicional (a máquina de contadores de Minsky). A completude é um BOTÃO, não o default: a hierarquia de Chomsky exposta como decisão.
\prov{relações: usa expressividade\_sub\_turing; relaciona linguagens\_totais; usa hierarquia\_de\_chomsky; usa maquina\_de\_contadores}
\prov{proveniência: wiki \textperiodcentered{} inferencia \textperiodcentered{} confiança 1.0 \textperiodcentered{} domínio programacao \textperiodcentered{} história 10 versões no jornal}

%CRISTAL {"arestas": [{"alvo": "desenvolvimento_web", "confianca": 1.0, "epistemico": "fato", "origem": "wiki", "rel": "usa"}], "confianca": 1.0, "contraexemplos": [], "descricao": "O conjunto vivo de princípios, tokens e componentes que dá coerência visual a um produto — a linguagem de design como código compartilhado.", "epistemico": "inferencia", "exemplos": [], "id": "design_system", "memoria": "semantica", "meta": {"dominio": "programacao", "fonte": ""}, "origem": "wiki", "palavras_chave": [], "sinonimos": [], "tipo": "conceito", "titulo": "Design System"}
\section{Design System}
O conjunto vivo de princípios, tokens e componentes que dá coerência visual a um produto — a linguagem de design como código compartilhado.
\prov{relações: usa desenvolvimento\_web}
\prov{proveniência: wiki \textperiodcentered{} inferencia \textperiodcentered{} confiança 1.0 \textperiodcentered{} domínio programacao \textperiodcentered{} história 4 versões no jornal}

%CRISTAL {"arestas": [{"alvo": "cristalchain", "confianca": 1.0, "epistemico": "fato", "origem": "wiki", "rel": "parte_de"}, {"alvo": "turmalina_paraiba_ancora", "confianca": 1.0, "epistemico": "fato", "origem": "wiki", "rel": "usa"}, {"alvo": "corpo_estelar", "confianca": 1.0, "epistemico": "fato", "origem": "wiki", "rel": "relaciona"}, {"alvo": "joalheria_hub", "confianca": 1.0, "epistemico": "fato", "origem": "wiki", "rel": "relaciona"}], "confianca": 1.0, "contraexemplos": [], "descricao": "Expande a GoldChain (ouro/prata/tinta) para a CristalChain (linguagem/design_cristalchain.py): paleta AZUL PARAÍBA (#0E9CA6) + COBRE (#B87333) + MANGANÊS claro (as cores da pedra), o LOGO como a turmalina lapidada, e a ORNAMENTAÇÃO em ÁRVORES FRACTAIS (um sistema-L gerado na máquina, do cobre no tronco ao azul nas pontas). A fonte de TEXTO segue Libertinus (a base herdada da GoldChain); o passo adiante é a camada DISPLAY gerada na máquina — o logo, os ornamentos e as iniciais capitulares (a letra + a árvore fractal). Primeira aplicação: papers/turmalina.tex. A ornamentação de ouro antiga vira referência.", "epistemico": "fato", "exemplos": [], "id": "design_system_cristalchain", "memoria": "semantica", "meta": {"dominio": "sistema", "fonte": "pipeline gráfico"}, "origem": "wiki", "palavras_chave": [], "sinonimos": [], "tipo": "conceito", "titulo": "O Design System da CristalChain"}
\section{O Design System da CristalChain}
Expande a GoldChain (ouro/prata/tinta) para a CristalChain (linguagem/design\_cristalchain.py): paleta AZUL PARAÍBA (\#0E9CA6) + COBRE (\#B87333) + MANGANÊS claro (as cores da pedra), o LOGO como a turmalina lapidada, e a ORNAMENTAÇÃO em ÁRVORES FRACTAIS (um sistema-L gerado na máquina, do cobre no tronco ao azul nas pontas). A fonte de TEXTO segue Libertinus (a base herdada da GoldChain); o passo adiante é a camada DISPLAY gerada na máquina — o logo, os ornamentos e as iniciais capitulares (a letra + a árvore fractal). Primeira aplicação: papers/turmalina.tex. A ornamentação de ouro antiga vira referência.
\prov{relações: parte\_de cristalchain; usa turmalina\_paraiba\_ancora; relaciona corpo\_estelar; relaciona joalheria\_hub}
\prov{proveniência: wiki \textperiodcentered{} pipeline gráfico \textperiodcentered{} fato \textperiodcentered{} confiança 1.0 \textperiodcentered{} domínio sistema \textperiodcentered{} história 11 versões no jornal}

%CRISTAL {"arestas": [{"alvo": "design_elementos_do_chatgpt", "confianca": 1.0, "epistemico": "fato", "origem": "wiki", "rel": "especializa"}, {"alvo": "design_system_cristalchain", "confianca": 1.0, "epistemico": "fato", "origem": "wiki", "rel": "parte_de"}, {"alvo": "area_render_gema", "confianca": 1.0, "epistemico": "fato", "origem": "wiki", "rel": "relaciona"}], "confianca": 1.0, "contraexemplos": [], "descricao": "Os elementos reconstruídos do design (divisor, cantoneira, selo, badge, moldura) deixam de ser raster e viram VETOR: gerar_todos emite cada um em SVG (+ pdf/png). O SVG guarda os PATHS — e cada contorno É uma ÓRBITA da forma, acessível pelo BAI e manipulável pela máquina. Vetorizar é o que abre a interpolação: com as órbitas na mão, a máquina anda entre as formas. linguagem/elementos_minerais.py, assets/cristalchain/elementos_gerados/*.svg.", "epistemico": "fato", "exemplos": [], "id": "design_vetorizado_svg", "memoria": "semantica", "meta": {"dominio": "sistema", "fonte": "orbitas vetoriais"}, "origem": "wiki", "palavras_chave": [], "sinonimos": [], "tipo": "conceito", "titulo": "Os Elementos Vetorizados em SVG (as órbitas nos paths)"}
\section{Os Elementos Vetorizados em SVG (as órbitas nos paths)}
Os elementos reconstruídos do design (divisor, cantoneira, selo, badge, moldura) deixam de ser raster e viram VETOR: gerar\_todos emite cada um em SVG (+ pdf/png). O SVG guarda os PATHS — e cada contorno É uma ÓRBITA da forma, acessível pelo BAI e manipulável pela máquina. Vetorizar é o que abre a interpolação: com as órbitas na mão, a máquina anda entre as formas. linguagem/elementos\_minerais.py, assets/cristalchain/elementos\_gerados/*.svg.
\prov{relações: especializa design\_elementos\_do\_chatgpt; parte\_de design\_system\_cristalchain; relaciona area\_render\_gema}
\prov{proveniência: wiki \textperiodcentered{} orbitas vetoriais \textperiodcentered{} fato \textperiodcentered{} confiança 1.0 \textperiodcentered{} domínio sistema \textperiodcentered{} história 6 versões no jornal}

%CRISTAL {"arestas": [{"alvo": "programacao_de_sistemas", "confianca": 1.0, "epistemico": "fato", "origem": "wiki", "rel": "relaciona"}], "confianca": 1.0, "contraexemplos": [], "descricao": "A cultura e a prática de unir desenvolvimento e operação — automação, entrega contínua e responsabilidade compartilhada pelo que roda.", "epistemico": "inferencia", "exemplos": [], "id": "devops", "memoria": "semantica", "meta": {"dominio": "programacao", "fonte": ""}, "origem": "wiki", "palavras_chave": [], "sinonimos": [], "tipo": "conceito", "titulo": "DevOps"}
\section{DevOps}
A cultura e a prática de unir desenvolvimento e operação — automação, entrega contínua e responsabilidade compartilhada pelo que roda.
\prov{relações: relaciona programacao\_de\_sistemas}
\prov{proveniência: wiki \textperiodcentered{} inferencia \textperiodcentered{} confiança 1.0 \textperiodcentered{} domínio programacao \textperiodcentered{} história 4 versões no jornal}

%CRISTAL {"arestas": [{"alvo": "sass_gpu_assembly", "confianca": 1.0, "epistemico": "fato", "origem": "wiki", "rel": "usa"}, {"alvo": "maquina_universal_gpu", "confianca": 1.0, "epistemico": "fato", "origem": "wiki", "rel": "relaciona"}], "confianca": 1.0, "contraexemplos": [], "descricao": "O hub das arquiteturas NVIDIA, cada uma um SASS próprio (a codificação muda por geração). A escada da compute capability (o -arch do ptxas): Pascal sm_60/61 (2016), Volta sm_70 (2017, 1º Tensor core), Turing sm_75 (2018, 1º RT core; a GTX 1650), Ampere sm_80/86 (2020, TF32/sparsity/MIG), Ada sm_89 (2022, FP8/DLSS3/SER), Hopper sm_90 (2022, Transformer Engine/clusters/TMA), Blackwell sm_100/120 (2024, FP4/2 dies). Cada geração: mais Tensor cores, mais formatos (FP16→BF16→TF32→FP8→FP4), mais banda (GDDR6→HBM2→HBM3), e um SASS distinto. O compilador do broca-so precisa de um dicionário por arquitetura (o -arch alvo) para emitir o SASS certo — como a reta mineral tem um metal por ordem, a GPU tem um SASS por geração. Síntese: cada arquitetura é uma INSTÂNCIA do metal (a máquina universal em silício).", "epistemico": "hipotese", "exemplos": [], "id": "dicionario_arquiteturas_gpu", "memoria": "semantica", "meta": {"dominio": "sass_gpu", "fonte": "SASS/GPU (o metal)"}, "origem": "wiki", "palavras_chave": [], "sinonimos": [], "tipo": "conceito", "titulo": "Dicionário das arquiteturas GPU (Turing→Ampere→Ada→Hopper→Blackwell)"}
\section{Dicionário das arquiteturas GPU (Turing\ensuremath{\to}Ampere\ensuremath{\to}Ada\ensuremath{\to}Hopper\ensuremath{\to}Blackwell)}
O hub das arquiteturas NVIDIA, cada uma um SASS próprio (a codificação muda por geração). A escada da compute capability (o -arch do ptxas): Pascal sm\_60/61 (2016), Volta sm\_70 (2017, 1º Tensor core), Turing sm\_75 (2018, 1º RT core; a GTX 1650), Ampere sm\_80/86 (2020, TF32/sparsity/MIG), Ada sm\_89 (2022, FP8/DLSS3/SER), Hopper sm\_90 (2022, Transformer Engine/clusters/TMA), Blackwell sm\_100/120 (2024, FP4/2 dies). Cada geração: mais Tensor cores, mais formatos (FP16\ensuremath{\to}BF16\ensuremath{\to}TF32\ensuremath{\to}FP8\ensuremath{\to}FP4), mais banda (GDDR6\ensuremath{\to}HBM2\ensuremath{\to}HBM3), e um SASS distinto. O compilador do broca-so precisa de um dicionário por arquitetura (o -arch alvo) para emitir o SASS certo — como a reta mineral tem um metal por ordem, a GPU tem um SASS por geração. Síntese: cada arquitetura é uma INSTÂNCIA do metal (a máquina universal em silício).
\prov{relações: usa sass\_gpu\_assembly; relaciona maquina\_universal\_gpu}
\prov{proveniência: wiki \textperiodcentered{} SASS/GPU (o metal) \textperiodcentered{} hipotese \textperiodcentered{} confiança 1.0 \textperiodcentered{} domínio sass\_gpu \textperiodcentered{} história 21 versões no jornal}

%CRISTAL {"arestas": [{"alvo": "dicionarios_centralizados", "confianca": 1.0, "epistemico": "fato", "origem": "wiki", "rel": "parte_de"}, {"alvo": "cristalografia", "confianca": 1.0, "epistemico": "fato", "origem": "wiki", "rel": "usa"}, {"alvo": "semicondutor", "confianca": 1.0, "epistemico": "fato", "origem": "wiki", "rel": "usa"}, {"alvo": "dopagem", "confianca": 1.0, "epistemico": "fato", "origem": "wiki", "rel": "usa"}, {"alvo": "dispositivo_semicondutor", "confianca": 1.0, "epistemico": "fato", "origem": "wiki", "rel": "usa"}, {"alvo": "rede_cristalina", "confianca": 1.0, "epistemico": "fato", "origem": "wiki", "rel": "usa"}, {"alvo": "conjugados_galois_parabola", "confianca": 1.0, "epistemico": "fato", "origem": "wiki", "rel": "relaciona"}, {"alvo": "bases_pisot", "confianca": 1.0, "epistemico": "fato", "origem": "wiki", "rel": "relaciona"}], "confianca": 1.0, "contraexemplos": [], "descricao": "A cadeia cristalografia→semicondutor→dispositivo tem estrutura de dicionário, como os metais do broca. DOPAGEM: a dualidade tipo-N/tipo-P em simbologias paralelas (portador elétron/lacuna, grupo V/III, dopante fósforo/boro, carga −/+) — compõem pelo pivô 'tipo', igual a metal↔n↔opcode↔gato. DISPOSITIVOS: dispositivo→função/junção (diodo→retificar/uma_pn; e diodo, LED e fotovoltaica partilham 'uma_pn' — a junção não distingue, a FUNÇÃO sim: alias legítimo). CRISTALOGRAFIA: os 7 sistemas→eixos/exemplo (cúbico→a=b=c,90°→sal). A PONTE: o cristal FÍSICO (rede ordenada, reticulado) e o cristal ALGÉBRICO (Pisot, a conjugada dentro do círculo da parábola) partilham a palavra E o dicionário — classe 'cristal'→Pisot.", "epistemico": "fato", "exemplos": [], "id": "dicionario_cristal_semicondutor", "memoria": "semantica", "meta": {"dominio": "programacao", "fonte": "linguagem/dicionarios.py (dopagem_*/dispositivo_*/cristal_*); cadeia cristal_a_energia"}, "origem": "wiki", "palavras_chave": [], "sinonimos": [], "tipo": "conceito", "titulo": "Dicionário da Cadeia Física — cristais, semicondutores, dispositivos"}
\section{Dicionário da Cadeia Física — cristais, semicondutores, dispositivos}
A cadeia cristalografia\ensuremath{\to}semicondutor\ensuremath{\to}dispositivo tem estrutura de dicionário, como os metais do broca. DOPAGEM: a dualidade tipo-N/tipo-P em simbologias paralelas (portador elétron/lacuna, grupo V/III, dopante fósforo/boro, carga \ensuremath{-}/+) — compõem pelo pivô 'tipo', igual a metal\ensuremath{\leftrightarrow}n\ensuremath{\leftrightarrow}opcode\ensuremath{\leftrightarrow}gato. DISPOSITIVOS: dispositivo\ensuremath{\to}função/junção (diodo\ensuremath{\to}retificar/uma\_pn; e diodo, LED e fotovoltaica partilham 'uma\_pn' — a junção não distingue, a FUNÇÃO sim: alias legítimo). CRISTALOGRAFIA: os 7 sistemas\ensuremath{\to}eixos/exemplo (cúbico\ensuremath{\to}a=b=c,90$^{\circ}$\ensuremath{\to}sal). A PONTE: o cristal FÍSICO (rede ordenada, reticulado) e o cristal ALGÉBRICO (Pisot, a conjugada dentro do círculo da parábola) partilham a palavra E o dicionário — classe 'cristal'\ensuremath{\to}Pisot.
\prov{relações: parte\_de dicionarios\_centralizados; usa cristalografia; usa semicondutor; usa dopagem; usa dispositivo\_semicondutor; usa rede\_cristalina; relaciona conjugados\_galois\_parabola; relaciona bases\_pisot}
\prov{proveniência: wiki \textperiodcentered{} linguagem/dicionarios.py (dopagem\_*/dispositivo\_*/cristal\_*); cadeia cristal\_a\_energia \textperiodcentered{} fato \textperiodcentered{} confiança 1.0 \textperiodcentered{} domínio programacao \textperiodcentered{} história 17 versões no jornal}

%CRISTAL {"arestas": [{"alvo": "sass_gpu_assembly", "confianca": 1.0, "epistemico": "fato", "origem": "wiki", "rel": "usa"}, {"alvo": "maquina_universal_gpu", "confianca": 1.0, "epistemico": "fato", "origem": "wiki", "rel": "parte_de"}, {"alvo": "computador_universal_pnp", "confianca": 1.0, "epistemico": "fato", "origem": "wiki", "rel": "usa"}, {"alvo": "gato", "confianca": 1.0, "epistemico": "fato", "origem": "wiki", "rel": "usa"}, {"alvo": "gamba_antissimetrico", "confianca": 1.0, "epistemico": "fato", "origem": "wiki", "rel": "usa"}], "confianca": 1.0, "contraexemplos": [], "descricao": "O dicionário do gato/gambá (a máquina universal) para o SASS (o metal): o WARP (32 threads SIMT) ↔ os N-lanes (32 lanes por warp — a costura La Hire em hardware); a FFMA (d=a·b+c, o fused multiply-add) ↔ a LA HIRE / o gato (a operação afim, multiplicar+somar — o coração do cálculo matricial); a MUFU (sin/cos/rsqrt) ↔ o GAMBÁ (a rotação, as transcendentes); os registradores R0–R255 ↔ o OCTÔNIO/o estado (o portador); LDG/STG (load/store) ↔ o PIPE (a transmissão, a órbita na memória); ISETP/predicado + BRA ↔ o GATO que mede (o if, o colapso, a decisão); SHFL (warp shuffle) ↔ a La Hire ENTRE lanes (cross-lane); BAR.SYNC ↔ o checksum/a sincronização (a+c²=1); a SM (Streaming Multiprocessor) ↔ o CORPO (a instância); PTX ↔ a DSL de gatos (a intermediária); SASS ↔ o METAL (a máquina de verdade). A cadeia BAI→DSL→PTX→SASS→SM. Síntese declarada: o gato/gambá são as instruções do metal da GPU.", "epistemico": "hipotese", "exemplos": [], "id": "dicionario_maquina_universal_sass", "memoria": "semantica", "meta": {"dominio": "sass_gpu", "fonte": "SASS/GPU (o metal)"}, "origem": "wiki", "palavras_chave": [], "sinonimos": [], "tipo": "conceito", "titulo": "Dicionário: a máquina universal ↔ SASS (o gato/gambá no metal da GPU)"}
\section{Dicionário: a máquina universal \ensuremath{\leftrightarrow} SASS (o gato/gambá no metal da GPU)}
O dicionário do gato/gambá (a máquina universal) para o SASS (o metal): o WARP (32 threads SIMT) \ensuremath{\leftrightarrow} os N-lanes (32 lanes por warp — a costura La Hire em hardware); a FFMA (d=a\textperiodcentered b+c, o fused multiply-add) \ensuremath{\leftrightarrow} a LA HIRE / o gato (a operação afim, multiplicar+somar — o coração do cálculo matricial); a MUFU (sin/cos/rsqrt) \ensuremath{\leftrightarrow} o GAMBÁ (a rotação, as transcendentes); os registradores R0–R255 \ensuremath{\leftrightarrow} o OCTÔNIO/o estado (o portador); LDG/STG (load/store) \ensuremath{\leftrightarrow} o PIPE (a transmissão, a órbita na memória); ISETP/predicado + BRA \ensuremath{\leftrightarrow} o GATO que mede (o if, o colapso, a decisão); SHFL (warp shuffle) \ensuremath{\leftrightarrow} a La Hire ENTRE lanes (cross-lane); BAR.SYNC \ensuremath{\leftrightarrow} o checksum/a sincronização (a+c\ensuremath{^{2}}=1); a SM (Streaming Multiprocessor) \ensuremath{\leftrightarrow} o CORPO (a instância); PTX \ensuremath{\leftrightarrow} a DSL de gatos (a intermediária); SASS \ensuremath{\leftrightarrow} o METAL (a máquina de verdade). A cadeia BAI\ensuremath{\to}DSL\ensuremath{\to}PTX\ensuremath{\to}SASS\ensuremath{\to}SM. Síntese declarada: o gato/gambá são as instruções do metal da GPU.
\prov{relações: usa sass\_gpu\_assembly; parte\_de maquina\_universal\_gpu; usa computador\_universal\_pnp; usa gato; usa gamba\_antissimetrico}
\prov{proveniência: wiki \textperiodcentered{} SASS/GPU (o metal) \textperiodcentered{} hipotese \textperiodcentered{} confiança 1.0 \textperiodcentered{} domínio sass\_gpu \textperiodcentered{} história 42 versões no jornal}

%CRISTAL {"arestas": [{"alvo": "teorema_geral_traducao", "confianca": 1.0, "epistemico": "fato", "origem": "wiki", "rel": "explica"}, {"alvo": "corolario_pivo_interlingua", "confianca": 1.0, "epistemico": "fato", "origem": "wiki", "rel": "usa"}, {"alvo": "correspondencia_mat_latex", "confianca": 1.0, "epistemico": "fato", "origem": "wiki", "rel": "usa"}, {"alvo": "peirce_dsl", "confianca": 1.0, "epistemico": "fato", "origem": "wiki", "rel": "usa"}, {"alvo": "conversao_mat_latex", "confianca": 1.0, "epistemico": "fato", "origem": "wiki", "rel": "usa"}, {"alvo": "bai_como_bussola", "confianca": 1.0, "epistemico": "fato", "origem": "wiki", "rel": "usa"}], "confianca": 1.0, "contraexemplos": [], "descricao": "O Teorema Geral da Tradução operacionalizado: linguagem/dicionarios.py junta TODAS as correspondências do broca-so num registro único, sobre UM motor (Tradutor/Peirce), e testa cada tradução nas PONTAS das arestas (par a↔b). 58 dicionários, 539/543 pontas conservadas. Cobrem matemática (math↔LaTeX, ASCII↔Unicode, gregas, lógica), o núcleo (metais↔n↔opcode↔gato, parábola), 28 LINGUAGENS (python/js/rust/go/haskell/c/java/elixir/lisp/clojure/zig/kotlin/sql/scala/ocaml/erlang/php/ruby/swift/dart/julia/r/perl/nim/crystal/lua/fsharp/groovy↔construto) e 4 FRAMEWORKS (react/vue/angular/svelte↔construto_ui) e 4 SHELLS (bash/zsh/fish/powershell↔construto_shell; a divergência aparece no FIM DE BLOCO: fi/end/}), HDL (verilog↔vhdl: module→entity, wire→signal, C-like vs Ada-like) e CONFIGS (json/yaml/toml: JSON sem comentário, TOML sem null, ':' vs '='). O pivô dá qualquer par de graça — 28 linguagens = 756 pares — python→rust, c→zig, react→svelte, sem dicionário novo. Dicionários PARCIAIS onde a linguagem não tem o construto (Haskell/C/Java sem palavra de função) — honesto e informativo, como as tipologias das línguas naturais. A bússola substituiu a busca.", "epistemico": "fato", "exemplos": [], "id": "dicionarios_centralizados", "memoria": "semantica", "meta": {"dominio": "programacao", "fonte": "linguagem/dicionarios.py (DICIONARIOS, compor, testar_todos)"}, "origem": "wiki", "palavras_chave": [], "sinonimos": [], "tipo": "conceito", "titulo": "Dicionários Centralizados — todas as traduções num registro só"}
\section{Dicionários Centralizados — todas as traduções num registro só}
O Teorema Geral da Tradução operacionalizado: linguagem/dicionarios.py junta TODAS as correspondências do broca-so num registro único, sobre UM motor (Tradutor/Peirce), e testa cada tradução nas PONTAS das arestas (par a\ensuremath{\leftrightarrow}b). 58 dicionários, 539/543 pontas conservadas. Cobrem matemática (math\ensuremath{\leftrightarrow}LaTeX, ASCII\ensuremath{\leftrightarrow}Unicode, gregas, lógica), o núcleo (metais\ensuremath{\leftrightarrow}n\ensuremath{\leftrightarrow}opcode\ensuremath{\leftrightarrow}gato, parábola), 28 LINGUAGENS (python/js/rust/go/haskell/c/java/elixir/lisp/clojure/zig/kotlin/sql/scala/ocaml/erlang/php/ruby/swift/dart/julia/r/perl/nim/crystal/lua/fsharp/groovy\ensuremath{\leftrightarrow}construto) e 4 FRAMEWORKS (react/vue/angular/svelte\ensuremath{\leftrightarrow}construto\_ui) e 4 SHELLS (bash/zsh/fish/powershell\ensuremath{\leftrightarrow}construto\_shell; a divergência aparece no FIM DE BLOCO: fi/end/\}), HDL (verilog\ensuremath{\leftrightarrow}vhdl: module\ensuremath{\to}entity, wire\ensuremath{\to}signal, C-like vs Ada-like) e CONFIGS (json/yaml/toml: JSON sem comentário, TOML sem null, ':' vs '='). O pivô dá qualquer par de graça — 28 linguagens = 756 pares — python\ensuremath{\to}rust, c\ensuremath{\to}zig, react\ensuremath{\to}svelte, sem dicionário novo. Dicionários PARCIAIS onde a linguagem não tem o construto (Haskell/C/Java sem palavra de função) — honesto e informativo, como as tipologias das línguas naturais. A bússola substituiu a busca.
\prov{relações: explica teorema\_geral\_traducao; usa corolario\_pivo\_interlingua; usa correspondencia\_mat\_latex; usa peirce\_dsl; usa conversao\_mat\_latex; usa bai\_como\_bussola}
\prov{proveniência: wiki \textperiodcentered{} linguagem/dicionarios.py (DICIONARIOS, compor, testar\_todos) \textperiodcentered{} fato \textperiodcentered{} confiança 1.0 \textperiodcentered{} domínio programacao \textperiodcentered{} história 70 versões no jornal}

%CRISTAL {"arestas": [{"alvo": "arquitetura_de_software", "confianca": 1.0, "epistemico": "fato", "origem": "wiki", "rel": "relaciona"}], "confianca": 1.0, "contraexemplos": [], "descricao": "O custo futuro de atalhos de hoje — como juros, cresce se não paga. Nem sempre ruim, desde que consciente e quitada.", "epistemico": "inferencia", "exemplos": [], "id": "divida_tecnica", "memoria": "semantica", "meta": {"dominio": "programacao", "fonte": ""}, "origem": "wiki", "palavras_chave": [], "sinonimos": [], "tipo": "conceito", "titulo": "Dívida Técnica"}
\section{Dívida Técnica}
O custo futuro de atalhos de hoje — como juros, cresce se não paga. Nem sempre ruim, desde que consciente e quitada.
\prov{relações: relaciona arquitetura\_de\_software}
\prov{proveniência: wiki \textperiodcentered{} inferencia \textperiodcentered{} confiança 1.0 \textperiodcentered{} domínio programacao \textperiodcentered{} história 2 versões no jornal}

%CRISTAL {"arestas": [{"alvo": "internet", "confianca": 1.0, "epistemico": "fato", "origem": "wiki", "rel": "parte_de"}], "confianca": 1.0, "contraexemplos": [], "descricao": "O sistema que traduz nomes (exemplo.com) em endereços IP — a agenda telefônica distribuída e hierárquica da internet.", "epistemico": "fato", "exemplos": [], "id": "dns", "memoria": "semantica", "meta": {"dominio": "programacao", "fonte": ""}, "origem": "wiki", "palavras_chave": [], "sinonimos": [], "tipo": "conceito", "titulo": "DNS"}
\section{DNS}
O sistema que traduz nomes (exemplo.com) em endereços IP — a agenda telefônica distribuída e hierárquica da internet.
\prov{relações: parte\_de internet}
\prov{proveniência: wiki \textperiodcentered{} fato \textperiodcentered{} confiança 1.0 \textperiodcentered{} domínio programacao \textperiodcentered{} história 4 versões no jornal}

%CRISTAL {"arestas": [{"alvo": "programacao", "confianca": 1.0, "epistemico": "fato", "origem": "wiki", "rel": "parte_de"}], "confianca": 1.0, "contraexemplos": [], "descricao": "O texto que explica o software — do README ao manual à referência de API. O código diz COMO; a documentação diz POR QUÊ e PARA QUEM.", "epistemico": "inferencia", "exemplos": [], "id": "documentacao", "memoria": "semantica", "meta": {"dominio": "programacao", "fonte": ""}, "origem": "wiki", "palavras_chave": [], "sinonimos": [], "tipo": "conceito", "titulo": "Documentação"}
\section{Documentação}
O texto que explica o software — do README ao manual à referência de API. O código diz COMO; a documentação diz POR QUÊ e PARA QUEM.
\prov{relações: parte\_de programacao}
\prov{proveniência: wiki \textperiodcentered{} inferencia \textperiodcentered{} confiança 1.0 \textperiodcentered{} domínio programacao \textperiodcentered{} história 2 versões no jornal}

%CRISTAL {"arestas": [{"alvo": "documentacao", "confianca": 1.0, "epistemico": "fato", "origem": "wiki", "rel": "especializa"}, {"alvo": "markdown", "confianca": 1.0, "epistemico": "fato", "origem": "wiki", "rel": "usa"}, {"alvo": "controle_de_versao", "confianca": 1.0, "epistemico": "fato", "origem": "wiki", "rel": "usa"}, {"alvo": "ci_cd", "confianca": 1.0, "epistemico": "fato", "origem": "wiki", "rel": "usa"}], "confianca": 1.0, "contraexemplos": [], "descricao": "Tratar a documentação como o código: em Markdown, versionada no Git, revisada e publicada por pipeline. A doc que não apodrece.", "epistemico": "inferencia", "exemplos": [], "id": "documentacao_como_codigo", "memoria": "semantica", "meta": {"dominio": "programacao", "fonte": ""}, "origem": "wiki", "palavras_chave": [], "sinonimos": [], "tipo": "conceito", "titulo": "Docs como Código"}
\section{Docs como Código}
Tratar a documentação como o código: em Markdown, versionada no Git, revisada e publicada por pipeline. A doc que não apodrece.
\prov{relações: especializa documentacao; usa markdown; usa controle\_de\_versao; usa ci\_cd}
\prov{proveniência: wiki \textperiodcentered{} inferencia \textperiodcentered{} confiança 1.0 \textperiodcentered{} domínio programacao \textperiodcentered{} história 5 versões no jornal}

%CRISTAL {"arestas": [{"alvo": "sistemas_distribuidos", "confianca": 1.0, "epistemico": "fato", "origem": "wiki", "rel": "parte_de"}], "confianca": 1.0, "contraexemplos": [], "descricao": "Prova que sobre um canal não-confiável NÃO há acordo garantido — o limite teórico que toda rede enfrenta.", "epistemico": "fato", "exemplos": [], "id": "dois_generais", "memoria": "semantica", "meta": {"dominio": "programacao", "fonte": ""}, "origem": "wiki", "palavras_chave": [], "sinonimos": [], "tipo": "conceito", "titulo": "Problema dos Dois Generais"}
\section{Problema dos Dois Generais}
Prova que sobre um canal não-confiável NÃO há acordo garantido — o limite teórico que toda rede enfrenta.
\prov{relações: parte\_de sistemas\_distribuidos}
\prov{proveniência: wiki \textperiodcentered{} fato \textperiodcentered{} confiança 1.0 \textperiodcentered{} domínio programacao \textperiodcentered{} história 4 versões no jornal}

%CRISTAL {"arestas": [{"alvo": "html", "confianca": 1.0, "epistemico": "fato", "origem": "wiki", "rel": "usa"}, {"alvo": "javascript", "confianca": 1.0, "epistemico": "fato", "origem": "wiki", "rel": "usa"}], "confianca": 1.0, "contraexemplos": [], "descricao": "A árvore de objetos que o navegador constrói a partir do HTML — a estrutura viva que o JavaScript lê e altera. O documento como dado.", "epistemico": "fato", "exemplos": [], "id": "dom", "memoria": "semantica", "meta": {"dominio": "programacao", "fonte": ""}, "origem": "wiki", "palavras_chave": [], "sinonimos": [], "tipo": "conceito", "titulo": "DOM (Document Object Model)"}
\section{DOM (Document Object Model)}
A árvore de objetos que o navegador constrói a partir do HTML — a estrutura viva que o JavaScript lê e altera. O documento como dado.
\prov{relações: usa html; usa javascript}
\prov{proveniência: wiki \textperiodcentered{} fato \textperiodcentered{} confiança 1.0 \textperiodcentered{} domínio programacao \textperiodcentered{} história 3 versões no jornal}

%CRISTAL {"arestas": [{"alvo": "dom", "confianca": 1.0, "epistemico": "fato", "origem": "wiki", "rel": "usa"}], "confianca": 1.0, "contraexemplos": [], "descricao": "Uma cópia leve do DOM em memória: o framework calcula o mínimo de mudanças (diff) e só aplica isso ao DOM real — desempenho por indireção.", "epistemico": "fato", "exemplos": [], "id": "dom_virtual", "memoria": "semantica", "meta": {"dominio": "programacao", "fonte": ""}, "origem": "wiki", "palavras_chave": [], "sinonimos": [], "tipo": "conceito", "titulo": "DOM Virtual"}
\section{DOM Virtual}
Uma cópia leve do DOM em memória: o framework calcula o mínimo de mudanças (diff) e só aplica isso ao DOM real — desempenho por indireção.
\prov{relações: usa dom}
\prov{proveniência: wiki \textperiodcentered{} fato \textperiodcentered{} confiança 1.0 \textperiodcentered{} domínio programacao \textperiodcentered{} história 4 versões no jornal}

%CRISTAL {"arestas": [{"alvo": "variavel_de_ambiente", "confianca": 1.0, "epistemico": "fato", "origem": "wiki", "rel": "usa"}, {"alvo": "arquivo_de_configuracao", "confianca": 1.0, "epistemico": "fato", "origem": "wiki", "rel": "especializa"}, {"alvo": "seguranca_da_informacao", "confianca": 1.0, "epistemico": "fato", "origem": "wiki", "rel": "relaciona"}], "confianca": 1.0, "contraexemplos": [], "descricao": "Um arquivo simples de variáveis de ambiente carregado no arranque — separa configuração e SEGREDOS do código. Deve ficar fora do controle de versão.", "epistemico": "inferencia", "exemplos": [], "id": "dotenv", "memoria": "semantica", "meta": {"dominio": "programacao", "fonte": ""}, "origem": "wiki", "palavras_chave": [], "sinonimos": [], "tipo": "conceito", "titulo": "Arquivo .env"}
\section{Arquivo .env}
Um arquivo simples de variáveis de ambiente carregado no arranque — separa configuração e SEGREDOS do código. Deve ficar fora do controle de versão.
\prov{relações: usa variavel\_de\_ambiente; especializa arquivo\_de\_configuracao; relaciona seguranca\_da\_informacao}
\prov{proveniência: wiki \textperiodcentered{} inferencia \textperiodcentered{} confiança 1.0 \textperiodcentered{} domínio programacao \textperiodcentered{} história 4 versões no jornal}

%CRISTAL {"arestas": [{"alvo": "arquivo_de_configuracao", "confianca": 1.0, "epistemico": "fato", "origem": "wiki", "rel": "especializa"}, {"alvo": "shell", "confianca": 1.0, "epistemico": "fato", "origem": "wiki", "rel": "usa"}, {"alvo": "controle_de_versao", "confianca": 1.0, "epistemico": "fato", "origem": "wiki", "rel": "usa"}], "confianca": 1.0, "contraexemplos": [], "descricao": "Os arquivos de configuração do usuário no diretório home (.bashrc, .gitconfig) — o ambiente pessoal, hoje versionado e compartilhado.", "epistemico": "fato", "exemplos": [], "id": "dotfiles", "memoria": "semantica", "meta": {"dominio": "programacao", "fonte": ""}, "origem": "wiki", "palavras_chave": [], "sinonimos": [], "tipo": "conceito", "titulo": "Dotfiles"}
\section{Dotfiles}
Os arquivos de configuração do usuário no diretório home (.bashrc, .gitconfig) — o ambiente pessoal, hoje versionado e compartilhado.
\prov{relações: especializa arquivo\_de\_configuracao; usa shell; usa controle\_de\_versao}
\prov{proveniência: wiki \textperiodcentered{} fato \textperiodcentered{} confiança 1.0 \textperiodcentered{} domínio programacao \textperiodcentered{} história 4 versões no jornal}

%CRISTAL {"arestas": [{"alvo": "cache", "confianca": 1.0, "epistemico": "fato", "origem": "curriculo", "rel": "usa"}, {"alvo": "kernel", "confianca": 0.75, "epistemico": "fato", "origem": "manual", "rel": "referencia"}, {"alvo": "memoria_virtual", "confianca": 0.5, "epistemico": "fato", "origem": "wiki", "rel": "relaciona"}, {"alvo": "main", "confianca": 0.5, "epistemico": "fato", "origem": "wiki", "rel": "relaciona"}, {"alvo": "go", "confianca": 0.5, "epistemico": "fato", "origem": "wiki", "rel": "relaciona"}, {"alvo": "thread", "confianca": 0.5, "epistemico": "fato", "origem": "wiki", "rel": "relaciona"}, {"alvo": "mamifero", "confianca": 0.5, "epistemico": "fato", "origem": "wiki", "rel": "relaciona"}, {"alvo": "java", "confianca": 0.5, "epistemico": "fato", "origem": "wiki", "rel": "relaciona"}, {"alvo": "word", "confianca": 0.5, "epistemico": "fato", "origem": "wiki", "rel": "relaciona"}], "confianca": 1.0, "contraexemplos": [], "descricao": "Código kernel ↔ hardware.", "epistemico": "fato", "exemplos": [], "id": "driver", "memoria": "semantica", "meta": {"dominio": "sistemas", "nivel": 6, "ordem": 3}, "origem": "curriculo", "palavras_chave": [], "sinonimos": [], "tipo": "conceito", "titulo": "driver"}
\section{driver}
Código kernel \ensuremath{\leftrightarrow} hardware.
\prov{relações: usa cache; referencia kernel; relaciona memoria\_virtual; relaciona main; relaciona go; relaciona thread; relaciona mamifero; relaciona java; relaciona word}
\prov{proveniência: curriculo \textperiodcentered{} fato \textperiodcentered{} confiança 1.0 \textperiodcentered{} domínio sistemas \textperiodcentered{} história 14 versões no jornal}

%CRISTAL {"arestas": [{"alvo": "kernel", "confianca": 1.0, "epistemico": "fato", "origem": "wiki", "rel": "parte_de"}, {"alvo": "cpu_metal_c", "confianca": 1.0, "epistemico": "fato", "origem": "wiki", "rel": "relaciona"}, {"alvo": "driver", "confianca": 1.0, "epistemico": "fato", "origem": "wiki", "rel": "relaciona"}], "confianca": 1.0, "contraexemplos": [], "descricao": "O código do kernel que fala a língua de um hardware específico e o expõe ao SO por uma interface comum. A ponte com o metal real.", "epistemico": "fato", "exemplos": [], "id": "driver_de_dispositivo", "memoria": "semantica", "meta": {"dominio": "programacao", "fonte": ""}, "origem": "wiki", "palavras_chave": [], "sinonimos": [], "tipo": "conceito", "titulo": "Driver de Dispositivo"}
\section{Driver de Dispositivo}
O código do kernel que fala a língua de um hardware específico e o expõe ao SO por uma interface comum. A ponte com o metal real.
\prov{relações: parte\_de kernel; relaciona cpu\_metal\_c; relaciona driver}
\prov{proveniência: wiki \textperiodcentered{} fato \textperiodcentered{} confiança 1.0 \textperiodcentered{} domínio programacao \textperiodcentered{} história 4 versões no jornal}

%CRISTAL {"arestas": [{"alvo": "dsl_frontend_linguagem_gatos", "confianca": 1.0, "epistemico": "fato", "origem": "wiki", "rel": "parte_de"}, {"alvo": "expressividade_sub_turing", "confianca": 1.0, "epistemico": "fato", "origem": "wiki", "rel": "relaciona"}, {"alvo": "sandbox_cristal_ao_so", "confianca": 1.0, "epistemico": "fato", "origem": "wiki", "rel": "usa"}], "confianca": 1.0, "contraexemplos": [], "descricao": "A DSL ganhou as comparações de ORDEM além de ==/!=: os registradores da RAM são inteiros com SINAL, então (esq−dir) tem sinal, e dois saltos novos jlz/jgz (salta se <0 / >0) no rodarram compilam <,>,<=,>=. O ==/!= saem IDÊNTICOS ao original (jnz/jz) — sem regressão. E o FOR ('for i = a to b … ', inclusivo) compila para i=a; while i<=b …; i=i+1 — usa o novo <=. Verificado: while i<5→5, for i=1 to 100→5050 (Gauss), if x<5→1.", "epistemico": "fato", "exemplos": [], "id": "dsl_comparacoes_ordem", "memoria": "semantica", "meta": {"dominio": "programacao", "fonte": "linguagem/dsl.py (_salta_se_falso, for); linguagem/universal.py (jlz/jgz)"}, "origem": "wiki", "palavras_chave": [], "sinonimos": [], "tipo": "conceito", "titulo": "DSL de gatos: comparações de ordem (< > <= >=) e o for"}
\section{DSL de gatos: comparações de ordem (< > <= >=) e o for}
A DSL ganhou as comparações de ORDEM além de ==/!=: os registradores da RAM são inteiros com SINAL, então (esq\ensuremath{-}dir) tem sinal, e dois saltos novos jlz/jgz (salta se <0 / >0) no rodarram compilam <,>,<=,>=. O ==/!= saem IDÊNTICOS ao original (jnz/jz) — sem regressão. E o FOR ('for i = a to b … ', inclusivo) compila para i=a; while i<=b …; i=i+1 — usa o novo <=. Verificado: while i<5\ensuremath{\to}5, for i=1 to 100\ensuremath{\to}5050 (Gauss), if x<5\ensuremath{\to}1.
\prov{relações: parte\_de dsl\_frontend\_linguagem\_gatos; relaciona expressividade\_sub\_turing; usa sandbox\_cristal\_ao\_so}
\prov{proveniência: wiki \textperiodcentered{} linguagem/dsl.py (\_salta\_se\_falso, for); linguagem/universal.py (jlz/jgz) \textperiodcentered{} fato \textperiodcentered{} confiança 1.0 \textperiodcentered{} domínio programacao \textperiodcentered{} história 13 versões no jornal}

%CRISTAL {"arestas": [{"alvo": "interpretador_nucleo_gatos", "confianca": 1.0, "epistemico": "fato", "origem": "wiki", "rel": "usa"}, {"alvo": "indirecao_posicional_transcende_turing", "confianca": 1.0, "epistemico": "fato", "origem": "wiki", "rel": "usa"}, {"alvo": "linguagem_matriz_gatos", "confianca": 1.0, "epistemico": "fato", "origem": "wiki", "rel": "parte_de"}, {"alvo": "fechar_laco_gato_posicao", "confianca": 1.0, "epistemico": "fato", "origem": "wiki", "rel": "usa"}, {"alvo": "backend_isa12_broca", "confianca": 1.0, "epistemico": "fato", "origem": "wiki", "rel": "usa"}, {"alvo": "tiffany", "confianca": 1.0, "epistemico": "fato", "origem": "wiki", "rel": "parte_de"}, {"alvo": "linguagem_de_programacao", "confianca": 1.0, "epistemico": "fato", "origem": "wiki", "rel": "parte_de"}], "confianca": 1.0, "contraexemplos": [], "descricao": "linguagem/dsl.py: o frontend da linguagem (paper §7, passo 4). Uma sintaxe legível — variáveis nomeadas, aritmética (+/−), while/if, o GATO (a porta metálica) e a MEMÓRIA indireta mem[·] (a indireção posicional, a Turing-completude em uso) — com tokenizer + parser recursivo + codegen. COMPILA para o assembler RAM (universal.py): set/mov/add/sub/addi/subi/loadi/storei/gato/jmp/jz/jnz/halt; while/if viram rótulos + jz/jnz sobre a diferença (esq−dir), expressões viram cadeias de add/sub, mem[·] vira loadi/storei. A PILHA COMPLETA roda: DSL → assembler → RAM → álgebra do gato. Verificado: Pell por while+gato (=2378), multiplicação por laço (6·7=42), soma de vetor por indireção (=100). Fecha o arco: da porta (gato) à linguagem escrevível de verdade. [F] roda.", "epistemico": "fato", "exemplos": [], "id": "dsl_frontend_linguagem_gatos", "memoria": "semantica", "meta": {"autor": "Lopes/broca-so", "dominio": "computacao", "fonte": "linguagem/dsl.py (compilar, executar; tokenizar/_Parser/_Compilador); verificado: Pell/multiplicação/soma-vetor; compila p/ universal.py"}, "origem": "wiki", "palavras_chave": [], "sinonimos": [], "tipo": "conceito", "titulo": "O frontend DSL: sintaxe legível que compila para o assembler (a pilha completa roda) [F]"}
\section{O frontend DSL: sintaxe legível que compila para o assembler (a pilha completa roda) [F]}
linguagem/dsl.py: o frontend da linguagem (paper §7, passo 4). Uma sintaxe legível — variáveis nomeadas, aritmética (+/\ensuremath{-}), while/if, o GATO (a porta metálica) e a MEMÓRIA indireta mem[\textperiodcentered ] (a indireção posicional, a Turing-completude em uso) — com tokenizer + parser recursivo + codegen. COMPILA para o assembler RAM (universal.py): set/mov/add/sub/addi/subi/loadi/storei/gato/jmp/jz/jnz/halt; while/if viram rótulos + jz/jnz sobre a diferença (esq\ensuremath{-}dir), expressões viram cadeias de add/sub, mem[\textperiodcentered ] vira loadi/storei. A PILHA COMPLETA roda: DSL \ensuremath{\to} assembler \ensuremath{\to} RAM \ensuremath{\to} álgebra do gato. Verificado: Pell por while+gato (=2378), multiplicação por laço (6\textperiodcentered 7=42), soma de vetor por indireção (=100). Fecha o arco: da porta (gato) à linguagem escrevível de verdade. [F] roda.
\prov{relações: usa interpretador\_nucleo\_gatos; usa indirecao\_posicional\_transcende\_turing; parte\_de linguagem\_matriz\_gatos; usa fechar\_laco\_gato\_posicao; usa backend\_isa12\_broca; parte\_de tiffany; parte\_de linguagem\_de\_programacao}
\prov{proveniência: wiki \textperiodcentered{} linguagem/dsl.py (compilar, executar; tokenizar/\_Parser/\_Compilador); verificado: Pell/multiplicação/soma-vetor; compila p/ universal.py \textperiodcentered{} fato \textperiodcentered{} confiança 1.0 \textperiodcentered{} domínio computacao \textperiodcentered{} história 28 versões no jornal}

%CRISTAL {"arestas": [{"alvo": "cifra_de_cristal", "confianca": 1.0, "epistemico": "fato", "origem": "wiki", "rel": "parte_de"}, {"alvo": "idempotentes_mod_n_fatoram", "confianca": 1.0, "epistemico": "fato", "origem": "wiki", "rel": "usa"}, {"alvo": "criptografia_assimetrica", "confianca": 1.0, "epistemico": "fato", "origem": "wiki", "rel": "explica"}, {"alvo": "lado_negro_de_newton", "confianca": 1.0, "epistemico": "fato", "origem": "wiki", "rel": "usa"}], "confianca": 1.0, "contraexemplos": [], "descricao": "Achar um idempotente NÃO-TRIVIAL e mod N (e²≡e, e∉{0,1}) É fatorar N: como N | e(e−1) e nem e nem e−1 é ≡0, gcd(e,N) é um fator próprio. A busca pelo idempotente é computacionalmente equivalente a fatorar — a mesma dureza do RSA. O lado negro (0) e o claro (1) são triviais; os do meio, que a lemniscata esconde, são a fatoração. Verif.: 143→66→11×13.", "epistemico": "fato", "exemplos": [], "id": "dureza_idempotente_fatora", "memoria": "semantica", "meta": {"dominio": "criptografia", "fonte": "cifra de cristal"}, "origem": "wiki", "palavras_chave": [], "sinonimos": [], "tipo": "conceito", "titulo": "A Dureza: o Idempotente é a Fatoração"}
\section{A Dureza: o Idempotente é a Fatoração}
Achar um idempotente NÃO-TRIVIAL e mod N (e\ensuremath{^{2}}\ensuremath{\equiv}e, e\ensuremath{\notin}\{0,1\}) É fatorar N: como N | e(e\ensuremath{-}1) e nem e nem e\ensuremath{-}1 é \ensuremath{\equiv}0, gcd(e,N) é um fator próprio. A busca pelo idempotente é computacionalmente equivalente a fatorar — a mesma dureza do RSA. O lado negro (0) e o claro (1) são triviais; os do meio, que a lemniscata esconde, são a fatoração. Verif.: 143\ensuremath{\to}66\ensuremath{\to}11$\times$13.
\prov{relações: parte\_de cifra\_de\_cristal; usa idempotentes\_mod\_n\_fatoram; explica criptografia\_assimetrica; usa lado\_negro\_de\_newton}
\prov{proveniência: wiki \textperiodcentered{} cifra de cristal \textperiodcentered{} fato \textperiodcentered{} confiança 1.0 \textperiodcentered{} domínio criptografia \textperiodcentered{} história 5 versões no jornal}

%CRISTAL {"arestas": [{"alvo": "programacao_funcional", "confianca": 1.0, "epistemico": "fato", "origem": "wiki", "rel": "relaciona"}, {"alvo": "monada", "confianca": 1.0, "epistemico": "fato", "origem": "wiki", "rel": "relaciona"}], "confianca": 1.0, "contraexemplos": [], "descricao": "Qualquer ação de uma função além de retornar um valor — mutar estado, I/O, rede. O funcional os isola nas bordas (mônadas, efeitos gerenciados).", "epistemico": "fato", "exemplos": [], "id": "efeito_colateral", "memoria": "semantica", "meta": {"dominio": "programacao", "fonte": ""}, "origem": "wiki", "palavras_chave": [], "sinonimos": [], "tipo": "conceito", "titulo": "Efeito Colateral"}
\section{Efeito Colateral}
Qualquer ação de uma função além de retornar um valor — mutar estado, I/O, rede. O funcional os isola nas bordas (mônadas, efeitos gerenciados).
\prov{relações: relaciona programacao\_funcional; relaciona monada}
\prov{proveniência: wiki \textperiodcentered{} fato \textperiodcentered{} confiança 1.0 \textperiodcentered{} domínio programacao \textperiodcentered{} história 3 versões no jornal}

%CRISTAL {"arestas": [{"alvo": "erlang", "confianca": 1.0, "epistemico": "fato", "origem": "wiki", "rel": "relaciona"}, {"alvo": "modelo_de_atores", "confianca": 1.0, "epistemico": "fato", "origem": "wiki", "rel": "usa"}], "confianca": 1.0, "contraexemplos": [], "descricao": "Valim: a produtividade do Ruby sobre a máquina virtual do Erlang (BEAM) — concorrência e resiliência acessíveis.", "epistemico": "fato", "exemplos": [], "id": "elixir", "memoria": "semantica", "meta": {"autor": "Valim", "dominio": "programacao", "fonte": ""}, "origem": "wiki", "palavras_chave": [], "sinonimos": [], "tipo": "conceito", "titulo": "Elixir"}
\section{Elixir}
Valim: a produtividade do Ruby sobre a máquina virtual do Erlang (BEAM) — concorrência e resiliência acessíveis.
\prov{relações: relaciona erlang; usa modelo\_de\_atores}
\prov{proveniência: wiki \textperiodcentered{} fato \textperiodcentered{} confiança 1.0 \textperiodcentered{} domínio programacao \textperiodcentered{} história 6 versões no jornal}

%CRISTAL {"arestas": [{"alvo": "programacao_funcional", "confianca": 1.0, "epistemico": "fato", "origem": "wiki", "rel": "especializa"}, {"alvo": "padrao_reducer", "confianca": 1.0, "epistemico": "fato", "origem": "wiki", "rel": "explica"}, {"alvo": "javascript", "confianca": 1.0, "epistemico": "fato", "origem": "wiki", "rel": "relaciona"}], "confianca": 1.0, "contraexemplos": [], "descricao": "Czaplicki: linguagem funcional PURA para a web, sem exceções em runtime — a 'Arquitetura Elm' (Model-Update-View) que inspirou o Redux.", "epistemico": "fato", "exemplos": [], "id": "elm", "memoria": "semantica", "meta": {"autor": "Czaplicki", "dominio": "programacao", "fonte": ""}, "origem": "wiki", "palavras_chave": [], "sinonimos": [], "tipo": "conceito", "titulo": "Elm"}
\section{Elm}
Czaplicki: linguagem funcional PURA para a web, sem exceções em runtime — a 'Arquitetura Elm' (Model-Update-View) que inspirou o Redux.
\prov{relações: especializa programacao\_funcional; explica padrao\_reducer; relaciona javascript}
\prov{proveniência: wiki \textperiodcentered{} fato \textperiodcentered{} confiança 1.0 \textperiodcentered{} domínio programacao \textperiodcentered{} história 4 versões no jornal}

%CRISTAL {"arestas": [{"alvo": "rede_neural_artificial", "confianca": 1.0, "epistemico": "fato", "origem": "wiki", "rel": "usa"}, {"alvo": "tiffany", "confianca": 1.0, "epistemico": "fato", "origem": "wiki", "rel": "explica"}, {"alvo": "teoria_da_informacao", "confianca": 1.0, "epistemico": "fato", "origem": "wiki", "rel": "relaciona"}], "confianca": 1.0, "contraexemplos": [], "descricao": "Representar palavras/conceitos como vetores onde a proximidade É semelhança de sentido — a base da busca por significado, como a recuperação da Tiffany.", "epistemico": "fato", "exemplos": [], "id": "embedding", "memoria": "semantica", "meta": {"dominio": "programacao", "fonte": ""}, "origem": "wiki", "palavras_chave": [], "sinonimos": [], "tipo": "conceito", "titulo": "Embedding (Vetor Semântico)"}
\section{Embedding (Vetor Semântico)}
Representar palavras/conceitos como vetores onde a proximidade É semelhança de sentido — a base da busca por significado, como a recuperação da Tiffany.
\prov{relações: usa rede\_neural\_artificial; explica tiffany; relaciona teoria\_da\_informacao}
\prov{proveniência: wiki \textperiodcentered{} fato \textperiodcentered{} confiança 1.0 \textperiodcentered{} domínio programacao \textperiodcentered{} história 4 versões no jornal}

%CRISTAL {"arestas": [{"alvo": "seguranca_da_informacao", "confianca": 1.0, "epistemico": "fato", "origem": "wiki", "rel": "parte_de"}], "confianca": 1.0, "contraexemplos": [], "descricao": "Atacar o humano, não a máquina — enganar alguém a entregar acesso. O elo mais fraco quase nunca é técnico.", "epistemico": "inferencia", "exemplos": [], "id": "engenharia_social", "memoria": "semantica", "meta": {"dominio": "programacao", "fonte": ""}, "origem": "wiki", "palavras_chave": [], "sinonimos": [], "tipo": "conceito", "titulo": "Engenharia Social"}
\section{Engenharia Social}
Atacar o humano, não a máquina — enganar alguém a entregar acesso. O elo mais fraco quase nunca é técnico.
\prov{relações: parte\_de seguranca\_da\_informacao}
\prov{proveniência: wiki \textperiodcentered{} inferencia \textperiodcentered{} confiança 1.0 \textperiodcentered{} domínio programacao \textperiodcentered{} história 2 versões no jornal}

%CRISTAL {"arestas": [{"alvo": "grafo_append_only", "confianca": 1.0, "epistemico": "fato", "origem": "wiki", "rel": "usa"}, {"alvo": "regua_dnci", "confianca": 1.0, "epistemico": "fato", "origem": "wiki", "rel": "usa"}], "confianca": 1.0, "contraexemplos": [], "descricao": "ensinar_conceito cria/funde um NoConhecimento tipo conceito com descrição/sinônimos/chaves/arestas, checando duplicata via resolver_id antes de gravar. É a primitiva de plantio (as semillas todas a usam).", "epistemico": "fato", "exemplos": [], "id": "ensinar_conceito", "memoria": "semantica", "meta": {"autor": "broca-so", "dominio": "operacao_so", "fonte": "conversa/conhecimento/ensinar.py"}, "origem": "wiki", "palavras_chave": [], "sinonimos": [], "tipo": "conceito", "titulo": "Ensinar um Conceito"}
\section{Ensinar um Conceito}
ensinar\_conceito cria/funde um NoConhecimento tipo conceito com descrição/sinônimos/chaves/arestas, checando duplicata via resolver\_id antes de gravar. É a primitiva de plantio (as semillas todas a usam).
\prov{relações: usa grafo\_append\_only; usa regua\_dnci}
\prov{proveniência: wiki \textperiodcentered{} conversa/conhecimento/ensinar.py \textperiodcentered{} fato \textperiodcentered{} confiança 1.0 \textperiodcentered{} domínio operacao\_so \textperiodcentered{} história 3 versões no jornal}

%CRISTAL {"arestas": [{"alvo": "symbolic_beta_protocolo", "confianca": 1.0, "epistemico": "fato", "origem": "wiki", "rel": "usa"}, {"alvo": "teorema_esterilidade_dimensao", "confianca": 1.0, "epistemico": "fato", "origem": "wiki", "rel": "usa"}, {"alvo": "gex_minera_ouro_real", "confianca": 1.0, "epistemico": "fato", "origem": "wiki", "rel": "usa"}, {"alvo": "materia_escura_lado_negro", "confianca": 1.0, "epistemico": "fato", "origem": "wiki", "rel": "relaciona"}, {"alvo": "koch_dissipa_nao_codifica", "confianca": 1.0, "epistemico": "fato", "origem": "wiki", "rel": "relaciona"}, {"alvo": "dimensao_de_prata_pell", "confianca": 1.0, "epistemico": "fato", "origem": "wiki", "rel": "explica"}, {"alvo": "transformada_metalica_fechada_perelman", "confianca": 1.0, "epistemico": "fato", "origem": "wiki", "rel": "consequencia"}], "confianca": 1.0, "contraexemplos": [], "descricao": "O ouro (ℚ(√5), Fibonacci) está MAPEADO — pelo teorema da esterilidade, mapeada a caixa de ouro a informação recircula estéril entre as retas. E o BTC já opera em c²=0,99 (a=0,01): a mineração do ouro esgotou. Então NÃO se mapeia a prata abstratamente (os Pell são só os números); ACOPLA-SE a goldchain pra IMPORTAR a entropia real do ouro (o PoW do BTC) e transduzi-la pra prata via emissão de Pell. Achar um número de prata = calcular as ÓRBITAS das partículas permitidas pelo BAI — mas a órbita só é REAL se carregar a entropia real do mundo; a goldchain é o cordão que a leva do ouro (BTC) pra prata (Pell).", "epistemico": "inferencia", "exemplos": [], "id": "entropia_ouro_para_prata", "memoria": "semantica", "meta": {"autor": "Lopes/broca-so", "dominio": "goldchain", "fonte": "goldchain (btc_conectar, minera); BAI (órbitas/energias)"}, "origem": "wiki", "palavras_chave": [], "sinonimos": [], "tipo": "conceito", "titulo": "Acoplar a goldchain: trazer a entropia REAL do ouro pra prata (a prata sem entropia é estéril)"}
\section{Acoplar a goldchain: trazer a entropia REAL do ouro pra prata (a prata sem entropia é estéril)}
O ouro (\ensuremath{\mathbb{Q}}(\ensuremath{\sqrt{\,}}5), Fibonacci) está MAPEADO — pelo teorema da esterilidade, mapeada a caixa de ouro a informação recircula estéril entre as retas. E o BTC já opera em c\ensuremath{^{2}}=0,99 (a=0,01): a mineração do ouro esgotou. Então NÃO se mapeia a prata abstratamente (os Pell são só os números); ACOPLA-SE a goldchain pra IMPORTAR a entropia real do ouro (o PoW do BTC) e transduzi-la pra prata via emissão de Pell. Achar um número de prata = calcular as ÓRBITAS das partículas permitidas pelo BAI — mas a órbita só é REAL se carregar a entropia real do mundo; a goldchain é o cordão que a leva do ouro (BTC) pra prata (Pell).
\prov{relações: usa symbolic\_beta\_protocolo; usa teorema\_esterilidade\_dimensao; usa gex\_minera\_ouro\_real; relaciona materia\_escura\_lado\_negro; relaciona koch\_dissipa\_nao\_codifica; explica dimensao\_de\_prata\_pell; consequencia transformada\_metalica\_fechada\_perelman}
\prov{proveniência: wiki \textperiodcentered{} goldchain (btc\_conectar, minera); BAI (órbitas/energias) \textperiodcentered{} inferencia \textperiodcentered{} confiança 1.0 \textperiodcentered{} domínio goldchain \textperiodcentered{} história 36 versões no jornal}

%CRISTAL {"arestas": [{"alvo": "landauer_gsl_bekenstein_bound", "confianca": 1.0, "epistemico": "fato", "origem": "wiki", "rel": "relaciona"}], "confianca": 1.0, "contraexemplos": [], "descricao": "A entropia H(X)=−Σ p(xᵢ) log₂ p(xᵢ) [bits/símbolo] mede a incerteza/conteúdo de informação médio de uma fonte. O TEOREMA DA CODIFICAÇÃO DE FONTE: a taxa mínima de compressão sem perdas é exatamente H — abaixo de H a recuperação sem erro é impossível. Shannon 1948, a fundação. A forma de H é derivada de axiomas. Fato.", "epistemico": "fato", "exemplos": [], "id": "entropia_shannon", "memoria": "semantica", "meta": {"autor": "broca-so", "dominio": "computacao", "fonte": "Shannon 1948 BSTJ 27,379"}, "origem": "wiki", "palavras_chave": [], "sinonimos": [], "tipo": "conceito", "titulo": "Entropia de Shannon e a codificação de fonte"}
\section{Entropia de Shannon e a codificação de fonte}
A entropia H(X)=\ensuremath{-}\ensuremath{\Sigma} p(x\ensuremath{_{i}}) log\ensuremath{_{2}} p(x\ensuremath{_{i}}) [bits/símbolo] mede a incerteza/conteúdo de informação médio de uma fonte. O TEOREMA DA CODIFICAÇÃO DE FONTE: a taxa mínima de compressão sem perdas é exatamente H — abaixo de H a recuperação sem erro é impossível. Shannon 1948, a fundação. A forma de H é derivada de axiomas. Fato.
\prov{relações: relaciona landauer\_gsl\_bekenstein\_bound}
\prov{proveniência: wiki \textperiodcentered{} Shannon 1948 BSTJ 27,379 \textperiodcentered{} fato \textperiodcentered{} confiança 1.0 \textperiodcentered{} domínio computacao \textperiodcentered{} história 2 versões no jornal}

%CRISTAL {"arestas": [{"alvo": "regua_dnci", "confianca": 1.0, "epistemico": "fato", "origem": "wiki", "rel": "explica"}, {"alvo": "grafo_append_only", "confianca": 1.0, "epistemico": "fato", "origem": "wiki", "rel": "usa"}], "confianca": 1.0, "contraexemplos": [], "descricao": "Estado epistêmico explícito por nó e por aresta (fato≥0.9 / inferência 0.5–0.7 / hipótese<0.5); o merge só eleva confiança ou une sinônimos — inferência NUNCA sobrescreve fato. É a régua DNCI codificada no store.", "epistemico": "inferencia", "exemplos": [], "id": "epistemico_nao_sobrescreve", "memoria": "semantica", "meta": {"autor": "broca-so", "dominio": "operacao_so", "fonte": "conversa/conhecimento/schema.py:Epistemico; store.py"}, "origem": "wiki", "palavras_chave": [], "sinonimos": [], "tipo": "conceito", "titulo": "Fato > Inferência > Hipótese"}
\section{Fato > Inferência > Hipótese}
Estado epistêmico explícito por nó e por aresta (fato\ensuremath{\ge}0.9 / inferência 0.5–0.7 / hipótese<0.5); o merge só eleva confiança ou une sinônimos — inferência NUNCA sobrescreve fato. É a régua DNCI codificada no store.
\prov{relações: explica regua\_dnci; usa grafo\_append\_only}
\prov{proveniência: wiki \textperiodcentered{} conversa/conhecimento/schema.py:Epistemico; store.py \textperiodcentered{} inferencia \textperiodcentered{} confiança 1.0 \textperiodcentered{} domínio operacao\_so \textperiodcentered{} história 3 versões no jornal}

%CRISTAL {"arestas": [{"alvo": "modelo_de_atores", "confianca": 1.0, "epistemico": "fato", "origem": "wiki", "rel": "usa"}], "confianca": 1.0, "contraexemplos": [], "descricao": "Ericsson: tolerância a falhas via modelo de atores e 'deixe quebrar'; nove noves de disponibilidade em telecom.", "epistemico": "fato", "exemplos": [], "id": "erlang", "memoria": "semantica", "meta": {"autor": "Armstrong", "dominio": "programacao", "fonte": ""}, "origem": "wiki", "palavras_chave": [], "sinonimos": [], "tipo": "conceito", "titulo": "Erlang"}
\section{Erlang}
Ericsson: tolerância a falhas via modelo de atores e 'deixe quebrar'; nove noves de disponibilidade em telecom.
\prov{relações: usa modelo\_de\_atores}
\prov{proveniência: wiki \textperiodcentered{} fato \textperiodcentered{} confiança 1.0 \textperiodcentered{} domínio programacao \textperiodcentered{} história 4 versões no jornal}

%CRISTAL {"arestas": [{"alvo": "criptoeconomia", "confianca": 1.0, "epistemico": "fato", "origem": "wiki", "rel": "relaciona"}, {"alvo": "nucleo_paralelo", "confianca": 1.0, "epistemico": "fato", "origem": "wiki", "rel": "relaciona"}], "confianca": 1.0, "contraexemplos": [], "descricao": "O número de tesourarias ≈ ceil(T_onchain/T_tick); 'bits de ouro' k=ceil(log_φ(razão)); degraus fixos (1,4,8,12,16,24,36) e política mín 4 / normal 8 / alto 12 / stress 16 / teto 36. n_ótimo para no degrau antes do ganho marginal cair sob 10%.", "epistemico": "inferencia", "exemplos": [], "id": "escala_ouro_phi", "memoria": "semantica", "meta": {"autor": "broca-so", "dominio": "operacao_so", "fonte": "goldchain/escala_ouro.py"}, "origem": "wiki", "palavras_chave": [], "sinonimos": [], "tipo": "conceito", "titulo": "Escala de Ouro φ (dimensionar tesourarias)"}
\section{Escala de Ouro \ensuremath{\varphi} (dimensionar tesourarias)}
O número de tesourarias \ensuremath{\approx} ceil(T\_onchain/T\_tick); 'bits de ouro' k=ceil(log\_\ensuremath{\varphi}(razão)); degraus fixos (1,4,8,12,16,24,36) e política mín 4 / normal 8 / alto 12 / stress 16 / teto 36. n\_ótimo para no degrau antes do ganho marginal cair sob 10\%.
\prov{relações: relaciona criptoeconomia; relaciona nucleo\_paralelo}
\prov{proveniência: wiki \textperiodcentered{} goldchain/escala\_ouro.py \textperiodcentered{} inferencia \textperiodcentered{} confiança 1.0 \textperiodcentered{} domínio operacao\_so \textperiodcentered{} história 3 versões no jornal}

%CRISTAL {"arestas": [{"alvo": "sistemas_distribuidos", "confianca": 1.0, "epistemico": "fato", "origem": "wiki", "rel": "usa"}, {"alvo": "computacao_em_nuvem", "confianca": 1.0, "epistemico": "fato", "origem": "wiki", "rel": "usa"}], "confianca": 1.0, "contraexemplos": [], "descricao": "Crescer sob carga sem cair — vertical (máquina maior) ou horizontal (mais máquinas). O horizontal exige enfrentar o distribuído.", "epistemico": "inferencia", "exemplos": [], "id": "escalabilidade", "memoria": "semantica", "meta": {"dominio": "programacao", "fonte": ""}, "origem": "wiki", "palavras_chave": [], "sinonimos": [], "tipo": "conceito", "titulo": "Escalabilidade"}
\section{Escalabilidade}
Crescer sob carga sem cair — vertical (máquina maior) ou horizontal (mais máquinas). O horizontal exige enfrentar o distribuído.
\prov{relações: usa sistemas\_distribuidos; usa computacao\_em\_nuvem}
\prov{proveniência: wiki \textperiodcentered{} inferencia \textperiodcentered{} confiança 1.0 \textperiodcentered{} domínio programacao \textperiodcentered{} história 6 versões no jornal}

%CRISTAL {"arestas": [{"alvo": "kernel", "confianca": 1.0, "epistemico": "fato", "origem": "wiki", "rel": "parte_de"}, {"alvo": "processo", "confianca": 1.0, "epistemico": "fato", "origem": "wiki", "rel": "usa"}, {"alvo": "scheduler", "confianca": 1.0, "epistemico": "fato", "origem": "wiki", "rel": "relaciona"}], "confianca": 1.0, "contraexemplos": [], "descricao": "A parte do kernel que decide qual processo roda em cada instante — a ilusão de que tudo acontece ao mesmo tempo numa CPU só.", "epistemico": "fato", "exemplos": [], "id": "escalonador", "memoria": "semantica", "meta": {"dominio": "programacao", "fonte": ""}, "origem": "wiki", "palavras_chave": [], "sinonimos": [], "tipo": "conceito", "titulo": "Escalonador (Scheduler)"}
\section{Escalonador (Scheduler)}
A parte do kernel que decide qual processo roda em cada instante — a ilusão de que tudo acontece ao mesmo tempo numa CPU só.
\prov{relações: parte\_de kernel; usa processo; relaciona scheduler}
\prov{proveniência: wiki \textperiodcentered{} fato \textperiodcentered{} confiança 1.0 \textperiodcentered{} domínio programacao \textperiodcentered{} história 4 versões no jornal}

%CRISTAL {"arestas": [{"alvo": "xml", "confianca": 1.0, "epistemico": "fato", "origem": "wiki", "rel": "usa"}, {"alvo": "json", "confianca": 1.0, "epistemico": "fato", "origem": "wiki", "rel": "usa"}], "confianca": 1.0, "contraexemplos": [], "descricao": "Um contrato que define a forma válida de um documento (XSD, JSON Schema) — verificar a estrutura antes de confiar nos dados.", "epistemico": "fato", "exemplos": [], "id": "esquema_validacao", "memoria": "semantica", "meta": {"dominio": "programacao", "fonte": ""}, "origem": "wiki", "palavras_chave": [], "sinonimos": [], "tipo": "conceito", "titulo": "Esquema e Validação"}
\section{Esquema e Validação}
Um contrato que define a forma válida de um documento (XSD, JSON Schema) — verificar a estrutura antes de confiar nos dados.
\prov{relações: usa xml; usa json}
\prov{proveniência: wiki \textperiodcentered{} fato \textperiodcentered{} confiança 1.0 \textperiodcentered{} domínio programacao \textperiodcentered{} história 3 versões no jornal}

%CRISTAL {"arestas": [{"alvo": "imutabilidade", "confianca": 1.0, "epistemico": "fato", "origem": "wiki", "rel": "especializa"}, {"alvo": "design_norma_conservada", "confianca": 1.0, "epistemico": "fato", "origem": "wiki", "rel": "relaciona"}], "confianca": 1.0, "contraexemplos": [], "descricao": "Nunca modificar dados no lugar: cada mudança produz um NOVO valor. Elimina corridas e torna o passado auditável — o invariante que se conserva.", "epistemico": "fato", "exemplos": [], "id": "estado_imutavel", "memoria": "semantica", "meta": {"dominio": "programacao", "fonte": ""}, "origem": "wiki", "palavras_chave": [], "sinonimos": [], "tipo": "conceito", "titulo": "Estado Imutável"}
\section{Estado Imutável}
Nunca modificar dados no lugar: cada mudança produz um NOVO valor. Elimina corridas e torna o passado auditável — o invariante que se conserva.
\prov{relações: especializa imutabilidade; relaciona design\_norma\_conservada}
\prov{proveniência: wiki \textperiodcentered{} fato \textperiodcentered{} confiança 1.0 \textperiodcentered{} domínio programacao \textperiodcentered{} história 3 versões no jornal}

%CRISTAL {"arestas": [{"alvo": "funcao", "confianca": 1.0, "epistemico": "fato", "origem": "curriculo", "rel": "usa"}, {"alvo": "kernel", "confianca": 0.5, "epistemico": "fato", "origem": "wiki", "rel": "relaciona"}, {"alvo": "topologia", "confianca": 0.5, "epistemico": "fato", "origem": "wiki", "rel": "relaciona"}, {"alvo": "gato_operador_hub", "confianca": 0.5, "epistemico": "fato", "origem": "wiki", "rel": "relaciona"}, {"alvo": "word", "confianca": 0.5, "epistemico": "fato", "origem": "wiki", "rel": "relaciona"}, {"alvo": "virtualizacao", "confianca": 0.5, "epistemico": "fato", "origem": "wiki", "rel": "relaciona"}, {"alvo": "vida-morte", "confianca": 0.5, "epistemico": "fato", "origem": "wiki", "rel": "relaciona"}, {"alvo": "vontade_de_poder", "confianca": 0.5, "epistemico": "fato", "origem": "wiki", "rel": "relaciona"}, {"alvo": "misticismo_nao_dual", "confianca": 0.5, "epistemico": "fato", "origem": "wiki", "rel": "relaciona"}], "confianca": 1.0, "contraexemplos": [], "descricao": "Forma de organizar informação para acesso eficiente.", "epistemico": "fato", "exemplos": [], "id": "estrutura_de_dados", "memoria": "semantica", "meta": {"dominio": "programacao", "nivel": 6, "ordem": 1}, "origem": "curriculo", "palavras_chave": [], "sinonimos": [], "tipo": "conceito", "titulo": "estrutura de dados"}
\section{estrutura de dados}
Forma de organizar informação para acesso eficiente.
\prov{relações: usa funcao; relaciona kernel; relaciona topologia; relaciona gato\_operador\_hub; relaciona word; relaciona virtualizacao; relaciona vida-morte; relaciona vontade\_de\_poder; relaciona misticismo\_nao\_dual}
\prov{proveniência: curriculo \textperiodcentered{} fato \textperiodcentered{} confiança 1.0 \textperiodcentered{} domínio programacao \textperiodcentered{} história 15 versões no jornal}

%CRISTAL {"arestas": [{"alvo": "linguagem_matematica", "confianca": 1.0, "epistemico": "fato", "origem": "wiki", "rel": "parte_de"}, {"alvo": "metodo_axiomatico", "confianca": 1.0, "epistemico": "fato", "origem": "wiki", "rel": "usa"}], "confianca": 1.0, "contraexemplos": [], "descricao": "A estrutura do discurso matemático: DEFINIÇÕES fixam termos, TEOREMAS afirmam verdades sobre eles, PROVAS as deduzem dos axiomas. A sintaxe do rigor.", "epistemico": "inferencia", "exemplos": [], "id": "estrutura_definicao_teorema_prova", "memoria": "semantica", "meta": {"dominio": "programacao", "fonte": ""}, "origem": "wiki", "palavras_chave": [], "sinonimos": [], "tipo": "conceito", "titulo": "Definição → Teorema → Prova"}
\section{Definição \ensuremath{\to} Teorema \ensuremath{\to} Prova}
A estrutura do discurso matemático: DEFINIÇÕES fixam termos, TEOREMAS afirmam verdades sobre eles, PROVAS as deduzem dos axiomas. A sintaxe do rigor.
\prov{relações: parte\_de linguagem\_matematica; usa metodo\_axiomatico}
\prov{proveniência: wiki \textperiodcentered{} inferencia \textperiodcentered{} confiança 1.0 \textperiodcentered{} domínio programacao \textperiodcentered{} história 3 versões no jornal}

%CRISTAL {"arestas": [{"alvo": "framework_backend", "confianca": 1.0, "epistemico": "fato", "origem": "wiki", "rel": "especializa"}, {"alvo": "nodejs", "confianca": 1.0, "epistemico": "fato", "origem": "wiki", "rel": "usa"}], "confianca": 1.0, "contraexemplos": [], "descricao": "O framework web minimalista do Node — rotas e middleware, sem opinião. A base de fato sobre a qual muitos outros são construídos.", "epistemico": "fato", "exemplos": [], "id": "express", "memoria": "semantica", "meta": {"autor": "Holowaychuk", "dominio": "programacao", "fonte": ""}, "origem": "wiki", "palavras_chave": [], "sinonimos": [], "tipo": "conceito", "titulo": "Express"}
\section{Express}
O framework web minimalista do Node — rotas e middleware, sem opinião. A base de fato sobre a qual muitos outros são construídos.
\prov{relações: especializa framework\_backend; usa nodejs}
\prov{proveniência: wiki \textperiodcentered{} fato \textperiodcentered{} confiança 1.0 \textperiodcentered{} domínio programacao \textperiodcentered{} história 6 versões no jornal}

%CRISTAL {"arestas": [{"alvo": "hierarquia_de_chomsky", "confianca": 1.0, "epistemico": "fato", "origem": "wiki", "rel": "usa"}, {"alvo": "gramatica_bnf", "confianca": 1.0, "epistemico": "fato", "origem": "wiki", "rel": "relaciona"}], "confianca": 1.0, "contraexemplos": [], "descricao": "Um padrão que descreve um conjunto de strings — a linguagem regular na base da hierarquia de Chomsky, tornada ferramenta diária.", "epistemico": "fato", "exemplos": [], "id": "expressao_regular", "memoria": "semantica", "meta": {"dominio": "programacao", "fonte": ""}, "origem": "wiki", "palavras_chave": [], "sinonimos": [], "tipo": "conceito", "titulo": "Expressão Regular"}
\section{Expressão Regular}
Um padrão que descreve um conjunto de strings — a linguagem regular na base da hierarquia de Chomsky, tornada ferramenta diária.
\prov{relações: usa hierarquia\_de\_chomsky; relaciona gramatica\_bnf}
\prov{proveniência: wiki \textperiodcentered{} fato \textperiodcentered{} confiança 1.0 \textperiodcentered{} domínio programacao \textperiodcentered{} história 3 versões no jornal}

%CRISTAL {"arestas": [{"alvo": "isa_da_broca", "confianca": 1.0, "epistemico": "fato", "origem": "wiki", "rel": "usa"}, {"alvo": "maquina_de_turing", "confianca": 1.0, "epistemico": "fato", "origem": "wiki", "rel": "contradiz"}, {"alvo": "classes_p_np_cook_levin", "confianca": 1.0, "epistemico": "fato", "origem": "wiki", "rel": "relaciona"}], "confianca": 1.0, "contraexemplos": [], "descricao": "Endereços de LOAD/STORE e offsets de salto são literais fixos no bytecode; registradores nunca selecionam endereço nem destino — não há indireção nem recursão dinâmica. É um autômato de fluxo estático sobre memória finita, não uma CPU genérica (prova de impossibilidade).", "epistemico": "fato", "exemplos": [], "id": "expressividade_sub_turing", "memoria": "semantica", "meta": {"autor": "broca-so", "dominio": "operacao_so", "fonte": "ula/EXPRESSIVIDADE.md"}, "origem": "wiki", "palavras_chave": [], "sinonimos": [], "tipo": "conceito", "titulo": "A Máquina é Estritamente Sub-Turing"}
\section{A Máquina é Estritamente Sub-Turing}
Endereços de LOAD/STORE e offsets de salto são literais fixos no bytecode; registradores nunca selecionam endereço nem destino — não há indireção nem recursão dinâmica. É um autômato de fluxo estático sobre memória finita, não uma CPU genérica (prova de impossibilidade).
\prov{relações: usa isa\_da\_broca; contradiz maquina\_de\_turing; relaciona classes\_p\_np\_cook\_levin}
\prov{proveniência: wiki \textperiodcentered{} ula/EXPRESSIVIDADE.md \textperiodcentered{} fato \textperiodcentered{} confiança 1.0 \textperiodcentered{} domínio operacao\_so \textperiodcentered{} história 4 versões no jornal}

%CRISTAL {"arestas": [{"alvo": "fabrica_ordem_construtiva", "confianca": 1.0, "epistemico": "fato", "origem": "wiki", "rel": "parte_de"}, {"alvo": "decomposicao_peirce", "confianca": 1.0, "epistemico": "fato", "origem": "wiki", "rel": "usa"}, {"alvo": "impressao_koch_compressao", "confianca": 1.0, "epistemico": "fato", "origem": "wiki", "rel": "relaciona"}], "confianca": 1.0, "contraexemplos": [], "descricao": "(0) MATÉRIA: do vazio+1 bit (F₂, o LFSR do gap x²+x+1) nasce a cadeia; cristais = fluxos de bits brutos, cada um uma pilha de ordem (~8,8 bits abaixo do máximo). (1-2) CLONES: união de clones (cópias limpas da parte branca) preserva os dois sem colapsar (La Hire direto colapsaria ℂâ); o original fica SELADO. (3) DOPAGEM: quebra de simetria endógena eᵢ⁺/ej⁻ (tipo-N portador / tipo-P buraco = divisor de zero, parte negra de Peirce). (4) JUNÇÃO: J=P₊TP₋ — é estrutura SSE RETIFICA (razão direta/reversa ~10⁹ = diodo; senão é decoração). (5) BIT GERADOR: semente F₂ + poda (proibir 1k) → um universo fractal NOVO de dimensão própria (0,694 ouro, 0,879 tribonacci,…→1), INDEPENDENTE da mãe (correlação~0) = informação nova, não recai na degenerescência. [N] validados.", "epistemico": "inferencia", "exemplos": [], "id": "fab_cinco_estagios", "memoria": "semantica", "meta": {"autor": "Lopes/broca-so", "dominio": "operacao_so", "fonte": "machine/fabrica.tex §Cinco estágios (código rodado)"}, "origem": "wiki", "palavras_chave": [], "sinonimos": [], "tipo": "conceito", "titulo": "Os cinco estágios da fab (matéria→clones→dopagem→junção→bit gerador)"}
\section{Os cinco estágios da fab (matéria\ensuremath{\to}clones\ensuremath{\to}dopagem\ensuremath{\to}junção\ensuremath{\to}bit gerador)}
(0) MATÉRIA: do vazio+1 bit (F\ensuremath{_{2}}, o LFSR do gap x\ensuremath{^{2}}+x+1) nasce a cadeia; cristais = fluxos de bits brutos, cada um uma pilha de ordem (\~{}8,8 bits abaixo do máximo). (1-2) CLONES: união de clones (cópias limpas da parte branca) preserva os dois sem colapsar (La Hire direto colapsaria \ensuremath{\mathbb{C}}â); o original fica SELADO. (3) DOPAGEM: quebra de simetria endógena e\ensuremath{_{i}}\ensuremath{^{+}}/ej\ensuremath{^{-}} (tipo-N portador / tipo-P buraco = divisor de zero, parte negra de Peirce). (4) JUNÇÃO: J=P\ensuremath{_{+}}TP\ensuremath{_{-}} — é estrutura SSE RETIFICA (razão direta/reversa \~{}10\ensuremath{^{9}} = diodo; senão é decoração). (5) BIT GERADOR: semente F\ensuremath{_{2}} + poda (proibir 1k) \ensuremath{\to} um universo fractal NOVO de dimensão própria (0,694 ouro, 0,879 tribonacci,…\ensuremath{\to}1), INDEPENDENTE da mãe (correlação\~{}0) = informação nova, não recai na degenerescência. [N] validados.
\prov{relações: parte\_de fabrica\_ordem\_construtiva; usa decomposicao\_peirce; relaciona impressao\_koch\_compressao}
\prov{proveniência: wiki \textperiodcentered{} machine/fabrica.tex §Cinco estágios (código rodado) \textperiodcentered{} inferencia \textperiodcentered{} confiança 1.0 \textperiodcentered{} domínio operacao\_so \textperiodcentered{} história 5 versões no jornal}

%CRISTAL {"arestas": [{"alvo": "fabrica_ordem_construtiva", "confianca": 1.0, "epistemico": "fato", "origem": "wiki", "rel": "parte_de"}, {"alvo": "associador_calor", "confianca": 1.0, "epistemico": "fato", "origem": "wiki", "rel": "relaciona"}, {"alvo": "landauer_gsl_bekenstein_bound", "confianca": 1.0, "epistemico": "fato", "origem": "wiki", "rel": "usa"}], "confianca": 1.0, "contraexemplos": [], "descricao": "NÃO há ganho energético: o combustível é a ORDEM que cada cristal traz (bits), CONSUMIDA (motor de Szilard/Landauer, 2ª lei intacta) — 'torcer dados não produz eletricidade'. A pilha é TRANSDUÇÃO POR ASSIMETRIA (analogia piezoelétrica declarada): pressão/torção num cristal ASSIMÉTRICO (não-centrossimétrico = a parte negra) gera um pulso; a energia vem do trabalho de deformar, a ordem só direciona/acopla. Verificado: sal simétrico dP=0, quartzo assimétrico pulso. O fio fértil: a quebra de simetria (ℓ≠0) SEPARA (segurança), DOPA (fab) e ACOPLA (piezo) — cristal simétrico é morto nos três. [D] 2ª lei; [I] a analogia piezo.", "epistemico": "inferencia", "exemplos": [], "id": "fab_combustivel_ordem", "memoria": "semantica", "meta": {"autor": "Lopes/broca-so", "dominio": "operacao_so", "fonte": "machine/fabrica.tex §Combustível, §Pilha"}, "origem": "wiki", "palavras_chave": [], "sinonimos": [], "tipo": "conceito", "titulo": "O combustível é ordem, não energia (a pilha piezo)"}
\section{O combustível é ordem, não energia (a pilha piezo)}
NÃO há ganho energético: o combustível é a ORDEM que cada cristal traz (bits), CONSUMIDA (motor de Szilard/Landauer, 2ª lei intacta) — 'torcer dados não produz eletricidade'. A pilha é TRANSDUÇÃO POR ASSIMETRIA (analogia piezoelétrica declarada): pressão/torção num cristal ASSIMÉTRICO (não-centrossimétrico = a parte negra) gera um pulso; a energia vem do trabalho de deformar, a ordem só direciona/acopla. Verificado: sal simétrico dP=0, quartzo assimétrico pulso. O fio fértil: a quebra de simetria (\ensuremath{\ell}\ensuremath{\ne}0) SEPARA (segurança), DOPA (fab) e ACOPLA (piezo) — cristal simétrico é morto nos três. [D] 2ª lei; [I] a analogia piezo.
\prov{relações: parte\_de fabrica\_ordem\_construtiva; relaciona associador\_calor; usa landauer\_gsl\_bekenstein\_bound}
\prov{proveniência: wiki \textperiodcentered{} machine/fabrica.tex §Combustível, §Pilha \textperiodcentered{} inferencia \textperiodcentered{} confiança 1.0 \textperiodcentered{} domínio operacao\_so \textperiodcentered{} história 4 versões no jornal}

%CRISTAL {"arestas": [{"alvo": "fabrica_ordem_construtiva", "confianca": 1.0, "epistemico": "fato", "origem": "wiki", "rel": "relaciona"}, {"alvo": "ganho_converge_n_estrela", "confianca": 1.0, "epistemico": "fato", "origem": "wiki", "rel": "usa"}, {"alvo": "missao_engenharia_ouro", "confianca": 1.0, "epistemico": "fato", "origem": "wiki", "rel": "relaciona"}], "confianca": 1.0, "contraexemplos": [], "descricao": "Fissão = quebrar núcleo metaestável (rico em divisores de zero) → liberar a ordem de COESÃO interna (≫ superficial), cascata/reação em cadeia; fusão = juntar cristais leves. DOIS LACRES: (a) ganho marginal converge a zero abaixo de N* (a Loucura de Ouro da Missão); (b) RISCO de contaminar os tecidos — a cascata propaga nilpotência; A=0/ortogonalidade protege contra correlação LINEAR mas NÃO contra propagação algébrica por divisores de zero (fura por canal lateral: hash/RNG/pool comuns) → proteção = QUARENTENA por construção, não convenção. Aarão freou: 'não sai fazendo fissão na loucura, pode espalhar radiação nos tecidos'. [C] lacrado, NÃO executado.", "epistemico": "hipotese", "exemplos": [], "id": "fab_fissao_fusao_lacrada", "memoria": "semantica", "meta": {"autor": "Lopes/broca-so", "dominio": "operacao_so", "fonte": "machine/fabrica.tex §Horizonte lacrado [C]"}, "origem": "wiki", "palavras_chave": [], "sinonimos": [], "tipo": "conceito", "titulo": "Fissão e fusão — o horizonte [C] lacrado (não executar)"}
\section{Fissão e fusão — o horizonte [C] lacrado (não executar)}
Fissão = quebrar núcleo metaestável (rico em divisores de zero) \ensuremath{\to} liberar a ordem de COESÃO interna (\ensuremath{\gg} superficial), cascata/reação em cadeia; fusão = juntar cristais leves. DOIS LACRES: (a) ganho marginal converge a zero abaixo de N* (a Loucura de Ouro da Missão); (b) RISCO de contaminar os tecidos — a cascata propaga nilpotência; A=0/ortogonalidade protege contra correlação LINEAR mas NÃO contra propagação algébrica por divisores de zero (fura por canal lateral: hash/RNG/pool comuns) \ensuremath{\to} proteção = QUARENTENA por construção, não convenção. Aarão freou: 'não sai fazendo fissão na loucura, pode espalhar radiação nos tecidos'. [C] lacrado, NÃO executado.
\prov{relações: relaciona fabrica\_ordem\_construtiva; usa ganho\_converge\_n\_estrela; relaciona missao\_engenharia\_ouro}
\prov{proveniência: wiki \textperiodcentered{} machine/fabrica.tex §Horizonte lacrado [C] \textperiodcentered{} hipotese \textperiodcentered{} confiança 1.0 \textperiodcentered{} domínio operacao\_so \textperiodcentered{} história 4 versões no jornal}

%CRISTAL {"arestas": [{"alvo": "seguranca_algebrica_estrutura_nao_hardness", "confianca": 1.0, "epistemico": "fato", "origem": "wiki", "rel": "usa"}, {"alvo": "a_ancora_observatorio_nao_fabrica", "confianca": 1.0, "epistemico": "fato", "origem": "wiki", "rel": "relaciona"}, {"alvo": "broca_os", "confianca": 1.0, "epistemico": "fato", "origem": "wiki", "rel": "parte_de"}], "confianca": 1.0, "contraexemplos": [], "descricao": "A linha de segurança provou que não se EXTRAI abaixo da alma (A=0; H(KP,KQ)≤H(P,Q)); a fab usa a impossibilidade AO CONTRÁRIO: como não se arromba, CONSTRÓI-SE — 'não dá pra arrombar cofres, mas dá pra CRIAR cofres'. É uma máquina que transforma ORDEM (informação organizada, em bits) em ESTRUTURA. A fab é o resultado maduro; o reator (fissão/fusão) não. Distinta do observatório (o cockpit CONTEMPLA; a fab FABRICA).", "epistemico": "inferencia", "exemplos": [], "id": "fabrica_ordem_construtiva", "memoria": "semantica", "meta": {"autor": "Lopes/broca-so", "dominio": "operacao_so", "fonte": "machine/fabrica.tex"}, "origem": "wiki", "palavras_chave": [], "sinonimos": [], "tipo": "conceito", "titulo": "A Fábrica — a virada construtiva (não extrair, fabricar)"}
\section{A Fábrica — a virada construtiva (não extrair, fabricar)}
A linha de segurança provou que não se EXTRAI abaixo da alma (A=0; H(KP,KQ)\ensuremath{\le}H(P,Q)); a fab usa a impossibilidade AO CONTRÁRIO: como não se arromba, CONSTRÓI-SE — 'não dá pra arrombar cofres, mas dá pra CRIAR cofres'. É uma máquina que transforma ORDEM (informação organizada, em bits) em ESTRUTURA. A fab é o resultado maduro; o reator (fissão/fusão) não. Distinta do observatório (o cockpit CONTEMPLA; a fab FABRICA).
\prov{relações: usa seguranca\_algebrica\_estrutura\_nao\_hardness; relaciona a\_ancora\_observatorio\_nao\_fabrica; parte\_de broca\_os}
\prov{proveniência: wiki \textperiodcentered{} machine/fabrica.tex \textperiodcentered{} inferencia \textperiodcentered{} confiança 1.0 \textperiodcentered{} domínio operacao\_so \textperiodcentered{} história 4 versões no jornal}

%CRISTAL {"arestas": [{"alvo": "colapso_multicamada_peso", "confianca": 1.0, "epistemico": "fato", "origem": "wiki", "rel": "usa"}, {"alvo": "o_conselho", "confianca": 1.0, "epistemico": "fato", "origem": "wiki", "rel": "relaciona"}], "confianca": 1.0, "contraexemplos": [], "descricao": "Em nós 'usuário — Tiffany', se a fala casa a esquerda devolve a direita (extrair_face); se nenhum candidato vence, responder devolve None — não fabrica resposta (deny by default). O SO prefere calar a inventar.", "epistemico": "inferencia", "exemplos": [], "id": "face_turno_fail_closed", "memoria": "semantica", "meta": {"autor": "broca-so", "dominio": "operacao_so", "fonte": "conversa/colapso.py:extrair_face/responder"}, "origem": "wiki", "palavras_chave": [], "sinonimos": [], "tipo": "conceito", "titulo": "Face do Turno Duplo + Fail-Closed"}
\section{Face do Turno Duplo + Fail-Closed}
Em nós 'usuário — Tiffany', se a fala casa a esquerda devolve a direita (extrair\_face); se nenhum candidato vence, responder devolve None — não fabrica resposta (deny by default). O SO prefere calar a inventar.
\prov{relações: usa colapso\_multicamada\_peso; relaciona o\_conselho}
\prov{proveniência: wiki \textperiodcentered{} conversa/colapso.py:extrair\_face/responder \textperiodcentered{} inferencia \textperiodcentered{} confiança 1.0 \textperiodcentered{} domínio operacao\_so \textperiodcentered{} história 3 versões no jornal}

%CRISTAL {"arestas": [{"alvo": "ciclo_uno", "confianca": 1.0, "epistemico": "fato", "origem": "wiki", "rel": "usa"}, {"alvo": "protocolo_assincrono", "confianca": 1.0, "epistemico": "fato", "origem": "wiki", "rel": "relaciona"}], "confianca": 1.0, "contraexemplos": [], "descricao": "O estado efêmero dos ciclos em curso num JSON por-ramo (etapa, desde, pell_idx, contrato, txid); etapas terminais removem o slot, e limpar_stale mata ramos presos >30 min. O cockpit lê isto via pulso.", "epistemico": "fato", "exemplos": [], "id": "fase_pulso", "memoria": "semantica", "meta": {"autor": "broca-so", "dominio": "operacao_so", "fonte": "goldchain/fase_mineracao.py"}, "origem": "wiki", "palavras_chave": [], "sinonimos": [], "tipo": "conceito", "titulo": "Fase de Mineração (pulso multi-ramo)"}
\section{Fase de Mineração (pulso multi-ramo)}
O estado efêmero dos ciclos em curso num JSON por-ramo (etapa, desde, pell\_idx, contrato, txid); etapas terminais removem o slot, e limpar\_stale mata ramos presos >30 min. O cockpit lê isto via pulso.
\prov{relações: usa ciclo\_uno; relaciona protocolo\_assincrono}
\prov{proveniência: wiki \textperiodcentered{} goldchain/fase\_mineracao.py \textperiodcentered{} fato \textperiodcentered{} confiança 1.0 \textperiodcentered{} domínio operacao\_so \textperiodcentered{} história 3 versões no jornal}

%CRISTAL {"arestas": [{"alvo": "transformada_metalica", "confianca": 1.0, "epistemico": "fato", "origem": "wiki", "rel": "usa"}, {"alvo": "beta_numeracao_metalica", "confianca": 1.0, "epistemico": "fato", "origem": "wiki", "rel": "usa"}, {"alvo": "dsl_frontend_linguagem_gatos", "confianca": 1.0, "epistemico": "fato", "origem": "wiki", "rel": "parte_de"}, {"alvo": "tiffany", "confianca": 1.0, "epistemico": "fato", "origem": "wiki", "rel": "parte_de"}], "confianca": 1.0, "contraexemplos": [], "descricao": "linguagem/dsl.py (tentafechar, posgato): a otimização onde a transformada metálica PAGA no compilador. O DSL reconhece o idioma `while C != 0 gato N A B; C = C − 1 ` com C,A,B CONSTANTES (constant-propagation) e substitui o laço O(k) por (A,B):= AₙC·(A,B), computado em O(log C) na compilação (potência rápida da matriz [[0,1],[1,n]], exata/unbounded) — emitindo 'set' direto: RUNTIME O(1), o laço SOME. É o 'for fechado' da posição ([[transformadametalica]], [[betaₙumeracaometalica]]) aplicado à linguagem. Verificado: a Pell `while` fecha em 'set b 2378' (laço→3 sets); fechado ≡ laço real ≡ exato (Pell(20)=38613965); contador NÃO-constante (vindo de mem) roda o laço normal (sem fold errado). O laço do gato É a posição — e o compilador sabe disso. [F] provado.", "epistemico": "fato", "exemplos": [], "id": "fechar_laco_gato_posicao", "memoria": "semantica", "meta": {"autor": "Lopes/broca-so", "dominio": "computacao", "fonte": "linguagem/dsl.py (_tenta_fechar, _pos_gato via potência de [[0,1],[1,n]]); verificado: Pell fecha em set; fechado≡laço≡exato"}, "origem": "wiki", "palavras_chave": [], "sinonimos": [], "tipo": "conceito", "titulo": "A transformada NO COMPILADOR: fecha o laço-de-gato na posição (O(k)→O(1)) [F]"}
\section{A transformada NO COMPILADOR: fecha o laço-de-gato na posição (O(k)\ensuremath{\to}O(1)) [F]}
linguagem/dsl.py (tentafechar, posgato): a otimização onde a transformada metálica PAGA no compilador. O DSL reconhece o idioma `while C != 0 gato N A B; C = C \ensuremath{-} 1 ` com C,A,B CONSTANTES (constant-propagation) e substitui o laço O(k) por (A,B):= A\ensuremath{_{n}}C\textperiodcentered (A,B), computado em O(log C) na compilação (potência rápida da matriz [[0,1],[1,n]], exata/unbounded) — emitindo 'set' direto: RUNTIME O(1), o laço SOME. É o 'for fechado' da posição ([[transformadametalica]], [[beta\ensuremath{_{n}}umeracaometalica]]) aplicado à linguagem. Verificado: a Pell `while` fecha em 'set b 2378' (laço\ensuremath{\to}3 sets); fechado \ensuremath{\equiv} laço real \ensuremath{\equiv} exato (Pell(20)=38613965); contador NÃO-constante (vindo de mem) roda o laço normal (sem fold errado). O laço do gato É a posição — e o compilador sabe disso. [F] provado.
\prov{relações: usa transformada\_metalica; usa beta\_numeracao\_metalica; parte\_de dsl\_frontend\_linguagem\_gatos; parte\_de tiffany}
\prov{proveniência: wiki \textperiodcentered{} linguagem/dsl.py (\_tenta\_fechar, \_pos\_gato via potência de [[0,1],[1,n]]); verificado: Pell fecha em set; fechado\ensuremath{\equiv}laço\ensuremath{\equiv}exato \textperiodcentered{} fato \textperiodcentered{} confiança 1.0 \textperiodcentered{} domínio computacao \textperiodcentered{} história 18 versões no jornal}

%CRISTAL {"arestas": [{"alvo": "compilador", "confianca": 1.0, "epistemico": "fato", "origem": "curriculo", "rel": "usa"}, {"alvo": "gato_operador_hub", "confianca": 0.5, "epistemico": "fato", "origem": "wiki", "rel": "relaciona"}, {"alvo": "historia", "confianca": 0.5, "epistemico": "fato", "origem": "wiki", "rel": "relaciona"}, {"alvo": "readme", "confianca": 0.5, "epistemico": "fato", "origem": "wiki", "rel": "relaciona"}, {"alvo": "reproduzir", "confianca": 0.5, "epistemico": "fato", "origem": "wiki", "rel": "relaciona"}, {"alvo": "word", "confianca": 0.5, "epistemico": "fato", "origem": "wiki", "rel": "relaciona"}, {"alvo": "utilitarismo", "confianca": 0.5, "epistemico": "fato", "origem": "wiki", "rel": "relaciona"}], "confianca": 1.0, "contraexemplos": [], "descricao": "Organização persistente de ficheiros e diretórios.", "epistemico": "fato", "exemplos": [], "id": "filesystem", "memoria": "semantica", "meta": {"dominio": "programacao", "nivel": 6, "ordem": 10}, "origem": "curriculo", "palavras_chave": [], "sinonimos": [], "tipo": "conceito", "titulo": "filesystem"}
\section{filesystem}
Organização persistente de ficheiros e diretórios.
\prov{relações: usa compilador; relaciona gato\_operador\_hub; relaciona historia; relaciona readme; relaciona reproduzir; relaciona word; relaciona utilitarismo}
\prov{proveniência: curriculo \textperiodcentered{} fato \textperiodcentered{} confiança 1.0 \textperiodcentered{} domínio programacao \textperiodcentered{} história 10 versões no jornal}

%CRISTAL {"arestas": [{"alvo": "unix", "confianca": 1.0, "epistemico": "fato", "origem": "wiki", "rel": "parte_de"}, {"alvo": "pipe_unix", "confianca": 1.0, "epistemico": "fato", "origem": "wiki", "rel": "usa"}, {"alvo": "tudo_e_arquivo", "confianca": 1.0, "epistemico": "fato", "origem": "wiki", "rel": "usa"}], "confianca": 1.0, "contraexemplos": [], "descricao": "Faça uma coisa e faça bem; componha programas pequenos por pipes; tudo é arquivo; texto é a interface universal. O minimalismo que o broca-so herda.", "epistemico": "inferencia", "exemplos": [], "id": "filosofia_unix", "memoria": "semantica", "meta": {"autor": "McIlroy", "dominio": "programacao", "fonte": ""}, "origem": "wiki", "palavras_chave": [], "sinonimos": [], "tipo": "conceito", "titulo": "Filosofia Unix"}
\section{Filosofia Unix}
Faça uma coisa e faça bem; componha programas pequenos por pipes; tudo é arquivo; texto é a interface universal. O minimalismo que o broca-so herda.
\prov{relações: parte\_de unix; usa pipe\_unix; usa tudo\_e\_arquivo}
\prov{proveniência: wiki \textperiodcentered{} inferencia \textperiodcentered{} confiança 1.0 \textperiodcentered{} domínio programacao \textperiodcentered{} história 4 versões no jornal}

%CRISTAL {"arestas": [{"alvo": "shell", "confianca": 1.0, "epistemico": "fato", "origem": "wiki", "rel": "especializa"}], "confianca": 1.0, "contraexemplos": [], "descricao": "Friendly Interactive Shell: sugestões e realce de sintaxe por padrão, sem configuração — ergonomia sobre a compatibilidade POSIX.", "epistemico": "fato", "exemplos": [], "id": "fish", "memoria": "semantica", "meta": {"dominio": "programacao", "fonte": ""}, "origem": "wiki", "palavras_chave": [], "sinonimos": [], "tipo": "conceito", "titulo": "Fish"}
\section{Fish}
Friendly Interactive Shell: sugestões e realce de sintaxe por padrão, sem configuração — ergonomia sobre a compatibilidade POSIX.
\prov{relações: especializa shell}
\prov{proveniência: wiki \textperiodcentered{} fato \textperiodcentered{} confiança 1.0 \textperiodcentered{} domínio programacao \textperiodcentered{} história 2 versões no jornal}

%CRISTAL {"arestas": [{"alvo": "protocolo_assincrono", "confianca": 1.0, "epistemico": "fato", "origem": "wiki", "rel": "usa"}, {"alvo": "htlc", "confianca": 1.0, "epistemico": "fato", "origem": "wiki", "rel": "usa"}], "confianca": 1.0, "contraexemplos": [], "descricao": "Braço ETH: preparar→pronto→fechar. O cliente trava ETH num lock (DNA≠central); no redeem o receiver leva 99% e a casa 1%. 'pronto' só libera quando o lock tem saldo ≥ valor+0.001 E a coordenação rádio (segredo HTLC) confere. Nenhum gas antes do depósito.", "epistemico": "fato", "exemplos": [], "id": "fluxo_comissao_terceiros", "memoria": "semantica", "meta": {"autor": "broca-so", "dominio": "operacao_so", "fonte": "goldchain/fluxo.py; coordenacao.py; fechar_comissao.py"}, "origem": "wiki", "palavras_chave": [], "sinonimos": [], "tipo": "conceito", "titulo": "Fluxo Comissão para Terceiros (zero gas até depositar)"}
\section{Fluxo Comissão para Terceiros (zero gas até depositar)}
Braço ETH: preparar\ensuremath{\to}pronto\ensuremath{\to}fechar. O cliente trava ETH num lock (DNA\ensuremath{\ne}central); no redeem o receiver leva 99\% e a casa 1\%. 'pronto' só libera quando o lock tem saldo \ensuremath{\ge} valor+0.001 E a coordenação rádio (segredo HTLC) confere. Nenhum gas antes do depósito.
\prov{relações: usa protocolo\_assincrono; usa htlc}
\prov{proveniência: wiki \textperiodcentered{} goldchain/fluxo.py; coordenacao.py; fechar\_comissao.py \textperiodcentered{} fato \textperiodcentered{} confiança 1.0 \textperiodcentered{} domínio operacao\_so \textperiodcentered{} história 3 versões no jornal}

%CRISTAL {"arestas": [{"alvo": "desenvolvimento_web", "confianca": 1.0, "epistemico": "fato", "origem": "wiki", "rel": "parte_de"}, {"alvo": "api", "confianca": 1.0, "epistemico": "fato", "origem": "wiki", "rel": "usa"}], "confianca": 1.0, "contraexemplos": [], "descricao": "Uma estrutura para construir servidores e APIs — rotas, middleware, acesso a dados. Express, NestJS, Koa no mundo Node.", "epistemico": "fato", "exemplos": [], "id": "framework_backend", "memoria": "semantica", "meta": {"dominio": "programacao", "fonte": ""}, "origem": "wiki", "palavras_chave": [], "sinonimos": [], "tipo": "conceito", "titulo": "Framework de Backend"}
\section{Framework de Backend}
Uma estrutura para construir servidores e APIs — rotas, middleware, acesso a dados. Express, NestJS, Koa no mundo Node.
\prov{relações: parte\_de desenvolvimento\_web; usa api}
\prov{proveniência: wiki \textperiodcentered{} fato \textperiodcentered{} confiança 1.0 \textperiodcentered{} domínio programacao \textperiodcentered{} história 6 versões no jornal}

%CRISTAL {"arestas": [{"alvo": "desenvolvimento_web", "confianca": 1.0, "epistemico": "fato", "origem": "wiki", "rel": "parte_de"}, {"alvo": "javascript", "confianca": 1.0, "epistemico": "fato", "origem": "wiki", "rel": "usa"}], "confianca": 1.0, "contraexemplos": [], "descricao": "Uma estrutura para construir interfaces por componentes reativos, gerindo o DOM e o estado por você — React, Angular, Vue.", "epistemico": "fato", "exemplos": [], "id": "framework_frontend", "memoria": "semantica", "meta": {"dominio": "programacao", "fonte": ""}, "origem": "wiki", "palavras_chave": [], "sinonimos": [], "tipo": "conceito", "titulo": "Framework de Frontend"}
\section{Framework de Frontend}
Uma estrutura para construir interfaces por componentes reativos, gerindo o DOM e o estado por você — React, Angular, Vue.
\prov{relações: parte\_de desenvolvimento\_web; usa javascript}
\prov{proveniência: wiki \textperiodcentered{} fato \textperiodcentered{} confiança 1.0 \textperiodcentered{} domínio programacao \textperiodcentered{} história 6 versões no jornal}

%CRISTAL {"arestas": [{"alvo": "programacao_funcional", "confianca": 1.0, "epistemico": "fato", "origem": "wiki", "rel": "especializa"}, {"alvo": "ocaml", "confianca": 1.0, "epistemico": "fato", "origem": "wiki", "rel": "relaciona"}], "confianca": 1.0, "contraexemplos": [], "descricao": "Syme (Microsoft): funcional-primeiro na plataforma .NET, com inferência de tipos e imutabilidade por padrão — o ML pragmático da Microsoft.", "epistemico": "fato", "exemplos": [], "id": "fsharp", "memoria": "semantica", "meta": {"autor": "Syme", "dominio": "programacao", "fonte": ""}, "origem": "wiki", "palavras_chave": [], "sinonimos": [], "tipo": "conceito", "titulo": "F#"}
\section{F\#}
Syme (Microsoft): funcional-primeiro na plataforma .NET, com inferência de tipos e imutabilidade por padrão — o ML pragmático da Microsoft.
\prov{relações: especializa programacao\_funcional; relaciona ocaml}
\prov{proveniência: wiki \textperiodcentered{} fato \textperiodcentered{} confiança 1.0 \textperiodcentered{} domínio programacao \textperiodcentered{} história 3 versões no jornal}

%CRISTAL {"arestas": [{"alvo": "recursao", "confianca": 1.0, "epistemico": "fato", "origem": "curriculo", "rel": "especializa"}, {"alvo": "regua_associativa_e_o_risco_de_vazamento", "confianca": 0.5, "epistemico": "fato", "origem": "wiki", "rel": "relaciona"}, {"alvo": "mecanica_quantica_hub", "confianca": 0.5, "epistemico": "fato", "origem": "wiki", "rel": "relaciona"}, {"alvo": "motivos", "confianca": 0.5, "epistemico": "fato", "origem": "wiki", "rel": "relaciona"}, {"alvo": "projeto_cristaljson", "confianca": 0.5, "epistemico": "fato", "origem": "wiki", "rel": "relaciona"}, {"alvo": "opcao_a_tarball_release", "confianca": 0.5, "epistemico": "fato", "origem": "wiki", "rel": "relaciona"}, {"alvo": "geometria_hub_broca", "confianca": 0.5, "epistemico": "fato", "origem": "wiki", "rel": "relaciona"}], "confianca": 1.0, "contraexemplos": [], "descricao": "Bloco nomeado com parâmetros e retorno; abstração de procedimento.", "epistemico": "fato", "exemplos": [], "id": "funcao", "memoria": "semantica", "meta": {"dominio": "programacao", "nivel": 6, "ordem": 2}, "origem": "curriculo", "palavras_chave": [], "sinonimos": [], "tipo": "conceito", "titulo": "função"}
\section{função}
Bloco nomeado com parâmetros e retorno; abstração de procedimento.
\prov{relações: especializa recursao; relaciona regua\_associativa\_e\_o\_risco\_de\_vazamento; relaciona mecanica\_quantica\_hub; relaciona motivos; relaciona projeto\_cristaljson; relaciona opcao\_a\_tarball\_release; relaciona geometria\_hub\_broca}
\prov{proveniência: curriculo \textperiodcentered{} fato \textperiodcentered{} confiança 1.0 \textperiodcentered{} domínio programacao \textperiodcentered{} história 14 versões no jornal}

%CRISTAL {"arestas": [{"alvo": "aprendizado_de_maquina", "confianca": 1.0, "epistemico": "fato", "origem": "wiki", "rel": "usa"}], "confianca": 1.0, "contraexemplos": [], "descricao": "A medida do quão errado o modelo está — o número que o treino minimiza. Define o que 'aprender bem' significa.", "epistemico": "fato", "exemplos": [], "id": "funcao_de_perda", "memoria": "semantica", "meta": {"dominio": "programacao", "fonte": ""}, "origem": "wiki", "palavras_chave": [], "sinonimos": [], "tipo": "conceito", "titulo": "Função de Perda"}
\section{Função de Perda}
A medida do quão errado o modelo está — o número que o treino minimiza. Define o que 'aprender bem' significa.
\prov{relações: usa aprendizado\_de\_maquina}
\prov{proveniência: wiki \textperiodcentered{} fato \textperiodcentered{} confiança 1.0 \textperiodcentered{} domínio programacao \textperiodcentered{} história 2 versões no jornal}

%CRISTAL {"arestas": [{"alvo": "criptografia", "confianca": 1.0, "epistemico": "fato", "origem": "wiki", "rel": "parte_de"}, {"alvo": "prova_de_trabalho", "confianca": 1.0, "epistemico": "fato", "origem": "wiki", "rel": "explica"}, {"alvo": "goldchain", "confianca": 1.0, "epistemico": "fato", "origem": "wiki", "rel": "usa"}], "confianca": 1.0, "contraexemplos": [], "descricao": "Resume qualquer dado num valor fixo, irreversível e sem colisões práticas (SHA-256) — a impressão digital; o trabalho da prova-de-trabalho.", "epistemico": "fato", "exemplos": [], "id": "funcao_hash", "memoria": "semantica", "meta": {"dominio": "programacao", "fonte": ""}, "origem": "wiki", "palavras_chave": [], "sinonimos": [], "tipo": "conceito", "titulo": "Função Hash Criptográfica"}
\section{Função Hash Criptográfica}
Resume qualquer dado num valor fixo, irreversível e sem colisões práticas (SHA-256) — a impressão digital; o trabalho da prova-de-trabalho.
\prov{relações: parte\_de criptografia; explica prova\_de\_trabalho; usa goldchain}
\prov{proveniência: wiki \textperiodcentered{} fato \textperiodcentered{} confiança 1.0 \textperiodcentered{} domínio programacao \textperiodcentered{} história 4 versões no jornal}

%CRISTAL {"arestas": [{"alvo": "programacao_funcional", "confianca": 1.0, "epistemico": "fato", "origem": "wiki", "rel": "parte_de"}, {"alvo": "calculo_lambda", "confianca": 1.0, "epistemico": "fato", "origem": "wiki", "rel": "usa"}], "confianca": 1.0, "contraexemplos": [], "descricao": "Funções tratadas como valores — passadas, retornadas, guardadas. A base da programação funcional e do λ-cálculo.", "epistemico": "fato", "exemplos": [], "id": "funcao_primeira_classe", "memoria": "semantica", "meta": {"dominio": "programacao", "fonte": ""}, "origem": "wiki", "palavras_chave": [], "sinonimos": [], "tipo": "conceito", "titulo": "Função de Primeira Classe"}
\section{Função de Primeira Classe}
Funções tratadas como valores — passadas, retornadas, guardadas. A base da programação funcional e do \ensuremath{\lambda}-cálculo.
\prov{relações: parte\_de programacao\_funcional; usa calculo\_lambda}
\prov{proveniência: wiki \textperiodcentered{} fato \textperiodcentered{} confiança 1.0 \textperiodcentered{} domínio programacao \textperiodcentered{} história 6 versões no jornal}

%CRISTAL {"arestas": [{"alvo": "programacao_funcional", "confianca": 1.0, "epistemico": "fato", "origem": "wiki", "rel": "parte_de"}, {"alvo": "efeito_colateral", "confianca": 1.0, "epistemico": "fato", "origem": "wiki", "rel": "contradiz"}], "confianca": 1.0, "contraexemplos": [], "descricao": "Uma função cujo resultado só depende das entradas e que não causa efeitos observáveis — mesma entrada, mesma saída, sempre. O átomo do funcional.", "epistemico": "fato", "exemplos": [], "id": "funcao_pura", "memoria": "semantica", "meta": {"dominio": "programacao", "fonte": ""}, "origem": "wiki", "palavras_chave": [], "sinonimos": [], "tipo": "conceito", "titulo": "Função Pura"}
\section{Função Pura}
Uma função cujo resultado só depende das entradas e que não causa efeitos observáveis — mesma entrada, mesma saída, sempre. O átomo do funcional.
\prov{relações: parte\_de programacao\_funcional; contradiz efeito\_colateral}
\prov{proveniência: wiki \textperiodcentered{} fato \textperiodcentered{} confiança 1.0 \textperiodcentered{} domínio programacao \textperiodcentered{} história 3 versões no jornal}

%CRISTAL {"arestas": [{"alvo": "gerenciamento_de_memoria", "confianca": 1.0, "epistemico": "fato", "origem": "wiki", "rel": "especializa"}], "confianca": 1.0, "contraexemplos": [], "descricao": "Liberar automaticamente a memória inalcançável — livra o programador do erro, ao custo de pausas e latência.", "epistemico": "fato", "exemplos": [], "id": "garbage_collection", "memoria": "semantica", "meta": {"dominio": "programacao", "fonte": ""}, "origem": "wiki", "palavras_chave": [], "sinonimos": [], "tipo": "conceito", "titulo": "Coleta de Lixo (GC)"}
\section{Coleta de Lixo (GC)}
Liberar automaticamente a memória inalcançável — livra o programador do erro, ao custo de pausas e latência.
\prov{relações: especializa gerenciamento\_de\_memoria}
\prov{proveniência: wiki \textperiodcentered{} fato \textperiodcentered{} confiança 1.0 \textperiodcentered{} domínio programacao \textperiodcentered{} história 4 versões no jornal}

%CRISTAL {"arestas": [{"alvo": "transformada_metalica", "confianca": 1.0, "epistemico": "fato", "origem": "wiki", "rel": "especializa"}, {"alvo": "algebra_dupolinomial_espaco", "confianca": 1.0, "epistemico": "fato", "origem": "wiki", "rel": "usa"}, {"alvo": "a_torre_de_hurwitzcayley--dickson", "confianca": 1.0, "epistemico": "fato", "origem": "wiki", "rel": "usa"}, {"alvo": "cifra_metais_an", "confianca": 1.0, "epistemico": "fato", "origem": "wiki", "rel": "usa"}, {"alvo": "simetria_por_reta_corpo", "confianca": 1.0, "epistemico": "fato", "origem": "wiki", "rel": "relaciona"}, {"alvo": "tiffany", "confianca": 1.0, "epistemico": "fato", "origem": "wiki", "rel": "parte_de"}], "confianca": 1.0, "contraexemplos": [], "descricao": "O gancho, fundamentado: o gato de grau 2 (Aₙ, a reta metálica x²−n·x−1) associa-se a outros gatos pela COMPANION m×m — e o grafo já diz que 'a matriz sináptica da rede de m neurônios é o companion m×m do gato; a rede de m é o m-bonacci' (a decomposição de Peirce: célula eᵢ=processo, canal Eᵢⱼ=sinapse, Σeᵢ=I). É a TRANSFORMADA METÁLICA generalizando o gato de grau 2 para grau m (o dupolinômio): posição=endereço, dupolinômio=programa, Aₙk em O(log k) (o 'for' fechado), e a parábola conserva (det=±1, a norma σ·conj=−1 = segurança-por-conservação). Verificado no PoC (linguagem/gatomatriz.py): companion tribonacci roda; A₂⁹0 por potência rápida. [F/D]", "epistemico": "inferencia", "exemplos": [], "id": "gato_generaliza_matriz_companion", "memoria": "semantica", "meta": {"autor": "Lopes/broca-so", "dominio": "computacao", "fonte": "linguagem/gato_matriz.py (matriz_de_gatos, potencia); a_torre_de_hurwitzcayley--dickson (companion m×m = m-bonacci); transformada_metalica"}, "origem": "wiki", "palavras_chave": [], "sinonimos": [], "tipo": "conceito", "titulo": "O gato 2×2 generaliza à matriz de gatos (companion m×m = m-bonacci) pela transformada"}
\section{O gato 2$\times$2 generaliza à matriz de gatos (companion m$\times$m = m-bonacci) pela transformada}
O gancho, fundamentado: o gato de grau 2 (A\ensuremath{_{n}}, a reta metálica x\ensuremath{^{2}}\ensuremath{-}n\textperiodcentered x\ensuremath{-}1) associa-se a outros gatos pela COMPANION m$\times$m — e o grafo já diz que 'a matriz sináptica da rede de m neurônios é o companion m$\times$m do gato; a rede de m é o m-bonacci' (a decomposição de Peirce: célula e\ensuremath{_{i}}=processo, canal E\ensuremath{_{ij}}=sinapse, \ensuremath{\Sigma}e\ensuremath{_{i}}=I). É a TRANSFORMADA METÁLICA generalizando o gato de grau 2 para grau m (o dupolinômio): posição=endereço, dupolinômio=programa, A\ensuremath{_{n}}k em O(log k) (o 'for' fechado), e a parábola conserva (det=$\pm$1, a norma \ensuremath{\sigma}\textperiodcentered conj=\ensuremath{-}1 = segurança-por-conservação). Verificado no PoC (linguagem/gatomatriz.py): companion tribonacci roda; A\ensuremath{_{2}}\ensuremath{^{9}}0 por potência rápida. [F/D]
\prov{relações: especializa transformada\_metalica; usa algebra\_dupolinomial\_espaco; usa a\_torre\_de\_hurwitzcayley--dickson; usa cifra\_metais\_an; relaciona simetria\_por\_reta\_corpo; parte\_de tiffany}
\prov{proveniência: wiki \textperiodcentered{} linguagem/gato\_matriz.py (matriz\_de\_gatos, potencia); a\_torre\_de\_hurwitzcayley--dickson (companion m$\times$m = m-bonacci); transformada\_metalica \textperiodcentered{} inferencia \textperiodcentered{} confiança 1.0 \textperiodcentered{} domínio computacao \textperiodcentered{} história 61 versões no jornal}

%CRISTAL {"arestas": [{"alvo": "cifra_metais_an", "confianca": 1.0, "epistemico": "fato", "origem": "wiki", "rel": "eh_um"}, {"alvo": "transformada_metalica", "confianca": 1.0, "epistemico": "fato", "origem": "wiki", "rel": "usa"}, {"alvo": "computador_mineral", "confianca": 1.0, "epistemico": "fato", "origem": "wiki", "rel": "parte_de"}], "confianca": 1.0, "contraexemplos": [], "descricao": "A matriz de gatos Aₙ=[[n,1],[1,0]] (o metal n) é a Unidade Lógica-Aritmética do computador mineral: a execução é a contração tensorial (produto de matrizes), ×Aₙ é um passo de aritmética/endereço O(1), e det Aₙ=−1 é o invariante que a ULA conserva. Compor gatos = um programa fiado num operador só. É a mesma reta x²−nx−1, agora como porta lógica.", "epistemico": "fato", "exemplos": [], "id": "gato_ula", "memoria": "semantica", "meta": {"dominio": "computacao", "fonte": "computador mineral"}, "origem": "wiki", "palavras_chave": [], "sinonimos": [], "tipo": "conceito", "titulo": "O Gato como ULA"}
\section{O Gato como ULA}
A matriz de gatos A\ensuremath{_{n}}=[[n,1],[1,0]] (o metal n) é a Unidade Lógica-Aritmética do computador mineral: a execução é a contração tensorial (produto de matrizes), $\times$A\ensuremath{_{n}} é um passo de aritmética/endereço O(1), e det A\ensuremath{_{n}}=\ensuremath{-}1 é o invariante que a ULA conserva. Compor gatos = um programa fiado num operador só. É a mesma reta x\ensuremath{^{2}}\ensuremath{-}nx\ensuremath{-}1, agora como porta lógica.
\prov{relações: eh\_um cifra\_metais\_an; usa transformada\_metalica; parte\_de computador\_mineral}
\prov{proveniência: wiki \textperiodcentered{} computador mineral \textperiodcentered{} fato \textperiodcentered{} confiança 1.0 \textperiodcentered{} domínio computacao \textperiodcentered{} história 4 versões no jornal}

%CRISTAL {"arestas": [{"alvo": "modelo_de_linguagem", "confianca": 1.0, "epistemico": "fato", "origem": "wiki", "rel": "explica"}, {"alvo": "tiffany_cabeca_broca", "confianca": 1.0, "epistemico": "fato", "origem": "wiki", "rel": "relaciona"}, {"alvo": "neuronio", "confianca": 1.0, "epistemico": "fato", "origem": "wiki", "rel": "usa"}, {"alvo": "orbita_de_partida", "confianca": 1.0, "epistemico": "fato", "origem": "wiki", "rel": "relaciona"}], "confianca": 1.0, "contraexemplos": [], "descricao": "O modelo gera UM token, realimenta-o na entrada e repete — um gerador, um bit, um laço: a mesma arquitetura do cabeca.py (escuta/pensa/fala num loop). O texto é uma ÓRBITA no espaço de tokens; aterrar é ficar na variedade da língua (cristal), alucinar é a órbita escapar dela (caos). Não fabricar: gerar é o laço, não um buffer.", "epistemico": "inferencia", "exemplos": [], "id": "geracao_gerador_bit_loop", "memoria": "semantica", "meta": {"dominio": "ia", "fonte": "IA e modelos de linguagem"}, "origem": "wiki", "palavras_chave": [], "sinonimos": [], "tipo": "conceito", "titulo": "Geração Autoregressiva = Gerador + Bit + Loop"}
\section{Geração Autoregressiva = Gerador + Bit + Loop}
O modelo gera UM token, realimenta-o na entrada e repete — um gerador, um bit, um laço: a mesma arquitetura do cabeca.py (escuta/pensa/fala num loop). O texto é uma ÓRBITA no espaço de tokens; aterrar é ficar na variedade da língua (cristal), alucinar é a órbita escapar dela (caos). Não fabricar: gerar é o laço, não um buffer.
\prov{relações: explica modelo\_de\_linguagem; relaciona tiffany\_cabeca\_broca; usa neuronio; relaciona orbita\_de\_partida}
\prov{proveniência: wiki \textperiodcentered{} IA e modelos de linguagem \textperiodcentered{} inferencia \textperiodcentered{} confiança 1.0 \textperiodcentered{} domínio ia \textperiodcentered{} história 5 versões no jornal}

%CRISTAL {"arestas": [{"alvo": "estado_imutavel", "confianca": 1.0, "epistemico": "fato", "origem": "wiki", "rel": "usa"}, {"alvo": "programacao_funcional", "confianca": 1.0, "epistemico": "fato", "origem": "wiki", "rel": "usa"}], "confianca": 1.0, "contraexemplos": [], "descricao": "Como uma aplicação organiza e propaga seu estado mutável — o problema central das interfaces; o funcional o resolve com imutabilidade e reducers.", "epistemico": "fato", "exemplos": [], "id": "gerenciamento_de_estado", "memoria": "semantica", "meta": {"dominio": "programacao", "fonte": ""}, "origem": "wiki", "palavras_chave": [], "sinonimos": [], "tipo": "conceito", "titulo": "Gerenciamento de Estado"}
\section{Gerenciamento de Estado}
Como uma aplicação organiza e propaga seu estado mutável — o problema central das interfaces; o funcional o resolve com imutabilidade e reducers.
\prov{relações: usa estado\_imutavel; usa programacao\_funcional}
\prov{proveniência: wiki \textperiodcentered{} fato \textperiodcentered{} confiança 1.0 \textperiodcentered{} domínio programacao \textperiodcentered{} história 3 versões no jornal}

%CRISTAL {"arestas": [{"alvo": "programacao", "confianca": 1.0, "epistemico": "fato", "origem": "wiki", "rel": "parte_de"}], "confianca": 1.0, "contraexemplos": [], "descricao": "Quem aloca e libera a memória — manual (C), automático por coletor (Java), ou por posse verificada (Rust).", "epistemico": "fato", "exemplos": [], "id": "gerenciamento_de_memoria", "memoria": "semantica", "meta": {"dominio": "programacao", "fonte": ""}, "origem": "wiki", "palavras_chave": [], "sinonimos": [], "tipo": "conceito", "titulo": "Gerenciamento de Memória"}
\section{Gerenciamento de Memória}
Quem aloca e libera a memória — manual (C), automático por coletor (Java), ou por posse verificada (Rust).
\prov{relações: parte\_de programacao}
\prov{proveniência: wiki \textperiodcentered{} fato \textperiodcentered{} confiança 1.0 \textperiodcentered{} domínio programacao \textperiodcentered{} história 4 versões no jornal}

%CRISTAL {"arestas": [{"alvo": "processo", "confianca": 1.0, "epistemico": "fato", "origem": "wiki", "rel": "usa"}, {"alvo": "kernel", "confianca": 1.0, "epistemico": "fato", "origem": "wiki", "rel": "parte_de"}, {"alvo": "processos", "confianca": 1.0, "epistemico": "fato", "origem": "wiki", "rel": "relaciona"}], "confianca": 1.0, "contraexemplos": [], "descricao": "Criar, escalonar, comunicar e encerrar processos (fork/exec, sinais, IPC) — o SO orquestrando quem roda e como conversam.", "epistemico": "fato", "exemplos": [], "id": "gerenciamento_de_processos", "memoria": "semantica", "meta": {"dominio": "programacao", "fonte": ""}, "origem": "wiki", "palavras_chave": [], "sinonimos": [], "tipo": "conceito", "titulo": "Gerência de Processos"}
\section{Gerência de Processos}
Criar, escalonar, comunicar e encerrar processos (fork/exec, sinais, IPC) — o SO orquestrando quem roda e como conversam.
\prov{relações: usa processo; parte\_de kernel; relaciona processos}
\prov{proveniência: wiki \textperiodcentered{} fato \textperiodcentered{} confiança 1.0 \textperiodcentered{} domínio programacao \textperiodcentered{} história 4 versões no jornal}

%CRISTAL {"arestas": [{"alvo": "controle_de_versao", "confianca": 1.0, "epistemico": "fato", "origem": "wiki", "rel": "especializa"}], "confianca": 1.0, "contraexemplos": [], "descricao": "Torvalds (2005): o controle de versão distribuído — cada clone tem a história inteira; commits ligados por hash formam um DAG imutável.", "epistemico": "fato", "exemplos": [], "id": "git", "memoria": "semantica", "meta": {"autor": "Torvalds", "dominio": "programacao", "fonte": ""}, "origem": "wiki", "palavras_chave": [], "sinonimos": [], "tipo": "conceito", "titulo": "Git"}
\section{Git}
Torvalds (2005): o controle de versão distribuído — cada clone tem a história inteira; commits ligados por hash formam um DAG imutável.
\prov{relações: especializa controle\_de\_versao}
\prov{proveniência: wiki \textperiodcentered{} fato \textperiodcentered{} confiança 1.0 \textperiodcentered{} domínio programacao \textperiodcentered{} história 4 versões no jornal}

%CRISTAL {"arestas": [{"alvo": "shell", "confianca": 1.0, "epistemico": "fato", "origem": "wiki", "rel": "usa"}, {"alvo": "expressao_regular", "confianca": 1.0, "epistemico": "fato", "origem": "wiki", "rel": "relaciona"}], "confianca": 1.0, "contraexemplos": [], "descricao": "Os curingas do shell (*, ?, []) que expandem para nomes de arquivo — casamento de padrões antes do comando rodar.", "epistemico": "fato", "exemplos": [], "id": "glob", "memoria": "semantica", "meta": {"dominio": "programacao", "fonte": ""}, "origem": "wiki", "palavras_chave": [], "sinonimos": [], "tipo": "conceito", "titulo": "Globbing"}
\section{Globbing}
Os curingas do shell (*, ?, []) que expandem para nomes de arquivo — casamento de padrões antes do comando rodar.
\prov{relações: usa shell; relaciona expressao\_regular}
\prov{proveniência: wiki \textperiodcentered{} fato \textperiodcentered{} confiança 1.0 \textperiodcentered{} domínio programacao \textperiodcentered{} história 3 versões no jornal}

%CRISTAL {"arestas": [{"alvo": "filosofia_unix", "confianca": 1.0, "epistemico": "fato", "origem": "wiki", "rel": "usa"}, {"alvo": "bash", "confianca": 1.0, "epistemico": "fato", "origem": "wiki", "rel": "relaciona"}], "confianca": 1.0, "contraexemplos": [], "descricao": "Stallman (1983): o projeto do software livre que construiu o userland Unix (compilador, coreutils, Bash) — o 'GNU' do GNU/Linux.", "epistemico": "fato", "exemplos": [], "id": "gnu", "memoria": "semantica", "meta": {"autor": "Stallman", "dominio": "programacao", "fonte": ""}, "origem": "wiki", "palavras_chave": [], "sinonimos": [], "tipo": "conceito", "titulo": "GNU"}
\section{GNU}
Stallman (1983): o projeto do software livre que construiu o userland Unix (compilador, coreutils, Bash) — o 'GNU' do GNU/Linux.
\prov{relações: usa filosofia\_unix; relaciona bash}
\prov{proveniência: wiki \textperiodcentered{} fato \textperiodcentered{} confiança 1.0 \textperiodcentered{} domínio programacao \textperiodcentered{} história 3 versões no jornal}

%CRISTAL {"arestas": [{"alvo": "compilador", "confianca": 1.0, "epistemico": "fato", "origem": "curriculo", "rel": "referencia"}, {"alvo": "goldcoin", "confianca": 0.5, "epistemico": "fato", "origem": "wiki", "rel": "relaciona"}, {"alvo": "sessao_interativa", "confianca": 0.5, "epistemico": "fato", "origem": "wiki", "rel": "relaciona"}, {"alvo": "python", "confianca": 0.5, "epistemico": "fato", "origem": "wiki", "rel": "relaciona"}, {"alvo": "install", "confianca": 0.5, "epistemico": "fato", "origem": "wiki", "rel": "relaciona"}, {"alvo": "recursao", "confianca": 0.5, "epistemico": "fato", "origem": "wiki", "rel": "relaciona"}, {"alvo": "resposta_a_pergunta_dificil", "confianca": 0.5, "epistemico": "fato", "origem": "wiki", "rel": "relaciona"}, {"alvo": "mamifero", "confianca": 0.5, "epistemico": "fato", "origem": "wiki", "rel": "relaciona"}, {"alvo": "linguagem_de_programacao", "confianca": 1.0, "epistemico": "fato", "origem": "wiki", "rel": "parte_de"}, {"alvo": "programacao_concorrente", "confianca": 1.0, "epistemico": "fato", "origem": "wiki", "rel": "usa"}, {"alvo": "garbage_collection", "confianca": 1.0, "epistemico": "fato", "origem": "wiki", "rel": "usa"}], "confianca": 1.0, "contraexemplos": [], "descricao": "Pike/Thompson (Google): simplicidade deliberada, goroutines e canais para concorrência, GC, compilação rápida.", "epistemico": "fato", "exemplos": [], "id": "go", "memoria": "semantica", "meta": {"autor": "Pike/Thompson", "dominio": "programacao", "fonte": ""}, "origem": "wiki", "palavras_chave": [], "sinonimos": [], "tipo": "conceito", "titulo": "Go"}
\section{Go}
Pike/Thompson (Google): simplicidade deliberada, goroutines e canais para concorrência, GC, compilação rápida.
\prov{relações: referencia compilador; relaciona goldcoin; relaciona sessao\_interativa; relaciona python; relaciona install; relaciona recursao; relaciona resposta\_a\_pergunta\_dificil; relaciona mamifero; parte\_de linguagem\_de\_programacao; usa programacao\_concorrente; usa garbage\_collection}
\prov{proveniência: wiki \textperiodcentered{} fato \textperiodcentered{} confiança 1.0 \textperiodcentered{} domínio programacao \textperiodcentered{} história 19 versões no jornal}

%CRISTAL {"arestas": [{"alvo": "jogo_de_soma_zero", "confianca": 1.0, "epistemico": "fato", "origem": "wiki", "rel": "eh_um"}, {"alvo": "numero_de_shannon", "confianca": 1.0, "epistemico": "fato", "origem": "wiki", "rel": "relaciona"}], "confianca": 1.0, "contraexemplos": [], "descricao": "Jogo milenar de território num tabuleiro 19×19: regras simples, complexidade colossal (~10¹⁷⁰ posições, muito além do xadrez). Foi o 'Everest' da IA — resistiu à força bruta até o MCTS + redes neurais.", "epistemico": "fato", "exemplos": [], "id": "go_jogo", "memoria": "semantica", "meta": {"dominio": "ia", "fonte": "jogos e IA"}, "origem": "wiki", "palavras_chave": [], "sinonimos": [], "tipo": "conceito", "titulo": "Go (Baduk)"}
\section{Go (Baduk)}
Jogo milenar de território num tabuleiro 19$\times$19: regras simples, complexidade colossal (\~{}10\ensuremath{^{170}} posições, muito além do xadrez). Foi o 'Everest' da IA — resistiu à força bruta até o MCTS + redes neurais.
\prov{relações: eh\_um jogo\_de\_soma\_zero; relaciona numero\_de\_shannon}
\prov{proveniência: wiki \textperiodcentered{} jogos e IA \textperiodcentered{} fato \textperiodcentered{} confiança 1.0 \textperiodcentered{} domínio ia \textperiodcentered{} história 3 versões no jornal}

%CRISTAL {"arestas": [{"alvo": "formalismo_hilbert", "confianca": 1.0, "epistemico": "fato", "origem": "wiki", "rel": "contradiz"}, {"alvo": "indecidibilidade_ns_zfc", "confianca": 1.0, "epistemico": "fato", "origem": "wiki", "rel": "explica"}, {"alvo": "hierarquia_pspace_indecidivel", "confianca": 1.0, "epistemico": "fato", "origem": "wiki", "rel": "relaciona"}, {"alvo": "turing_halting_church", "confianca": 1.0, "epistemico": "fato", "origem": "wiki", "rel": "relaciona"}], "confianca": 1.0, "contraexemplos": [], "descricao": "(1º) Todo sistema formal consistente e suficientemente rico (aritmética) contém proposições verdadeiras INDEMONSTRÁVEIS dentro dele; (2º) tal sistema não pode provar a própria consistência. Gödel 1931. Nó-ponte com o halting e com Kolmogorov — a trindade dos limites por auto-referência/diagonalização. Fato (teoremas provados).", "epistemico": "fato", "exemplos": [], "id": "godel_incompletude", "memoria": "semantica", "meta": {"autor": "broca-so", "dominio": "computacao", "fonte": "Gödel 1931 Monatshefte 38"}, "origem": "wiki", "palavras_chave": [], "sinonimos": [], "tipo": "conceito", "titulo": "Teoremas de incompletude de Gödel"}
\section{Teoremas de incompletude de Gödel}
(1º) Todo sistema formal consistente e suficientemente rico (aritmética) contém proposições verdadeiras INDEMONSTRÁVEIS dentro dele; (2º) tal sistema não pode provar a própria consistência. Gödel 1931. Nó-ponte com o halting e com Kolmogorov — a trindade dos limites por auto-referência/diagonalização. Fato (teoremas provados).
\prov{relações: contradiz formalismo\_hilbert; explica indecidibilidade\_ns\_zfc; relaciona hierarquia\_pspace\_indecidivel; relaciona turing\_halting\_church}
\prov{proveniência: wiki \textperiodcentered{} Gödel 1931 Monatshefte 38 \textperiodcentered{} fato \textperiodcentered{} confiança 1.0 \textperiodcentered{} domínio computacao \textperiodcentered{} história 6 versões no jornal}

%CRISTAL {"arestas": [{"alvo": "kernel", "confianca": 0.805, "epistemico": "fato", "origem": "manual", "rel": "referencia"}, {"alvo": "python", "confianca": 0.5, "epistemico": "fato", "origem": "wiki", "rel": "relaciona"}, {"alvo": "transcricao", "confianca": 0.5, "epistemico": "fato", "origem": "wiki", "rel": "relaciona"}, {"alvo": "install", "confianca": 0.5, "epistemico": "fato", "origem": "wiki", "rel": "relaciona"}, {"alvo": "teste_ordem", "confianca": 0.5, "epistemico": "fato", "origem": "wiki", "rel": "relaciona"}, {"alvo": "prova_isa", "confianca": 0.5, "epistemico": "fato", "origem": "wiki", "rel": "relaciona"}, {"alvo": "montagem", "confianca": 0.5, "epistemico": "fato", "origem": "wiki", "rel": "relaciona"}, {"alvo": "readme", "confianca": 0.5, "epistemico": "fato", "origem": "wiki", "rel": "relaciona"}, {"alvo": "pow_gato", "confianca": 0.5, "epistemico": "fato", "origem": "wiki", "rel": "relaciona"}, {"alvo": "koch_atlas", "confianca": 0.5, "epistemico": "fato", "origem": "wiki", "rel": "relaciona"}, {"alvo": "pow_espectral", "confianca": 0.5, "epistemico": "fato", "origem": "wiki", "rel": "relaciona"}, {"alvo": "sessao_interativa", "confianca": 0.5, "epistemico": "fato", "origem": "wiki", "rel": "relaciona"}, {"alvo": "main", "confianca": 0.5, "epistemico": "fato", "origem": "wiki", "rel": "relaciona"}, {"alvo": "java", "confianca": 0.5, "epistemico": "fato", "origem": "wiki", "rel": "relaciona"}, {"alvo": "koch_memoria", "confianca": 0.5, "epistemico": "fato", "origem": "wiki", "rel": "relaciona"}, {"alvo": "koch_parabola", "confianca": 0.5, "epistemico": "fato", "origem": "wiki", "rel": "relaciona"}, {"alvo": "pow_geometria", "confianca": 0.5, "epistemico": "fato", "origem": "wiki", "rel": "relaciona"}, {"alvo": "koch_pell", "confianca": 0.5, "epistemico": "fato", "origem": "wiki", "rel": "relaciona"}, {"alvo": "linker", "confianca": 0.5, "epistemico": "fato", "origem": "wiki", "rel": "relaciona"}, {"alvo": "grava_fato_resposta_normal_em_seguida", "confianca": 0.5, "epistemico": "fato", "origem": "wiki", "rel": "relaciona"}, {"alvo": "parabola_isa", "confianca": 0.5, "epistemico": "fato", "origem": "wiki", "rel": "relaciona"}, {"alvo": "hopfield_rede_hub", "confianca": 0.5, "epistemico": "fato", "origem": "wiki", "rel": "relaciona"}, {"alvo": "mamifero", "confianca": 0.5, "epistemico": "fato", "origem": "wiki", "rel": "relaciona"}], "confianca": 1.0, "contraexemplos": [], "descricao": "Compute paralelo massivo.", "epistemico": "fato", "exemplos": [], "id": "gpu", "memoria": "semantica", "meta": {"dominio": "sistemas", "nivel": 6, "ordem": 8}, "origem": "curriculo", "palavras_chave": [], "sinonimos": [], "tipo": "conceito", "titulo": "GPU"}
\section{GPU}
Compute paralelo massivo.
\prov{relações: referencia kernel; relaciona python; relaciona transcricao; relaciona install; relaciona teste\_ordem; relaciona prova\_isa; relaciona montagem; relaciona readme; relaciona pow\_gato; relaciona koch\_atlas; relaciona pow\_espectral; relaciona sessao\_interativa; relaciona main; relaciona java; relaciona koch\_memoria; relaciona koch\_parabola; relaciona pow\_geometria; relaciona koch\_pell; relaciona linker; relaciona grava\_fato\_resposta\_normal\_em\_seguida; relaciona parabola\_isa; relaciona hopfield\_rede\_hub; relaciona mamifero}
\prov{proveniência: curriculo \textperiodcentered{} fato \textperiodcentered{} confiança 1.0 \textperiodcentered{} domínio sistemas \textperiodcentered{} história 27 versões no jornal}

%CRISTAL {"arestas": [{"alvo": "aprendizado_de_maquina", "confianca": 1.0, "epistemico": "fato", "origem": "wiki", "rel": "usa"}, {"alvo": "funcao_de_perda", "confianca": 1.0, "epistemico": "fato", "origem": "wiki", "rel": "usa"}], "confianca": 1.0, "contraexemplos": [], "descricao": "Ajustar os pesos na direção que mais reduz o erro, passo a passo — descer a montanha da função de perda. O motor do treino.", "epistemico": "fato", "exemplos": [], "id": "gradiente_descendente", "memoria": "semantica", "meta": {"dominio": "programacao", "fonte": ""}, "origem": "wiki", "palavras_chave": [], "sinonimos": [], "tipo": "conceito", "titulo": "Gradiente Descendente"}
\section{Gradiente Descendente}
Ajustar os pesos na direção que mais reduz o erro, passo a passo — descer a montanha da função de perda. O motor do treino.
\prov{relações: usa aprendizado\_de\_maquina; usa funcao\_de\_perda}
\prov{proveniência: wiki \textperiodcentered{} fato \textperiodcentered{} confiança 1.0 \textperiodcentered{} domínio programacao \textperiodcentered{} história 3 versões no jornal}

%CRISTAL {"arestas": [{"alvo": "operacao_broca_so", "confianca": 1.0, "epistemico": "fato", "origem": "wiki", "rel": "parte_de"}, {"alvo": "grafo_de_proveniencia", "confianca": 1.0, "epistemico": "fato", "origem": "wiki", "rel": "explica"}], "confianca": 1.0, "contraexemplos": [], "descricao": "Cada nó é gravado como uma linha JSON no fim de conhecimento.graph.jsonl; _load_all reconstrói um dict em memória e no merge só ELEVA confiança, nunca reduz (confianca=max(nova, antiga)). A verdade é append-only, o índice deriva.", "epistemico": "fato", "exemplos": [], "id": "grafo_append_only", "memoria": "semantica", "meta": {"autor": "broca-so", "dominio": "operacao_so", "fonte": "conversa/conhecimento/store.py:append"}, "origem": "wiki", "palavras_chave": [], "sinonimos": [], "tipo": "conceito", "titulo": "O Grafo Append-Only com Cache em RAM"}
\section{O Grafo Append-Only com Cache em RAM}
Cada nó é gravado como uma linha JSON no fim de conhecimento.graph.jsonl; \_load\_all reconstrói um dict em memória e no merge só ELEVA confiança, nunca reduz (confianca=max(nova, antiga)). A verdade é append-only, o índice deriva.
\prov{relações: parte\_de operacao\_broca\_so; explica grafo\_de\_proveniencia}
\prov{proveniência: wiki \textperiodcentered{} conversa/conhecimento/store.py:append \textperiodcentered{} fato \textperiodcentered{} confiança 1.0 \textperiodcentered{} domínio operacao\_so \textperiodcentered{} história 3 versões no jornal}

%CRISTAL {"arestas": [{"alvo": "sintaxe", "confianca": 1.0, "epistemico": "fato", "origem": "curriculo", "rel": "parte_de"}, {"alvo": "virtualizacao", "confianca": 0.5, "epistemico": "fato", "origem": "wiki", "rel": "relaciona"}, {"alvo": "word", "confianca": 0.5, "epistemico": "fato", "origem": "wiki", "rel": "relaciona"}, {"alvo": "vida-morte", "confianca": 0.5, "epistemico": "fato", "origem": "wiki", "rel": "relaciona"}, {"alvo": "universidade_curriculo_em_camadas", "confianca": 0.5, "epistemico": "fato", "origem": "wiki", "rel": "relaciona"}, {"alvo": "semiotica_peirce_triade", "confianca": 0.5, "epistemico": "fato", "origem": "wiki", "rel": "relaciona"}, {"alvo": "psicanalise", "confianca": 0.5, "epistemico": "fato", "origem": "wiki", "rel": "relaciona"}, {"alvo": "proposition_gerador_de_escala_ipropderiv_com_tddt", "confianca": 0.5, "epistemico": "fato", "origem": "wiki", "rel": "relaciona"}, {"alvo": "linguistica", "confianca": 1.0, "epistemico": "fato", "origem": "wiki", "rel": "parte_de"}], "confianca": 1.0, "contraexemplos": [], "descricao": "Regras que estruturam palavras, frases e enunciados válidos.", "epistemico": "fato", "exemplos": [], "id": "gramatica", "memoria": "semantica", "meta": {"dominio": "linguagem", "nivel": 3, "ordem": 1}, "origem": "curriculo", "palavras_chave": ["morfologia", "sintaxe"], "sinonimos": ["gramática"], "tipo": "conceito", "titulo": "gramatica"}
\section{gramatica}
Regras que estruturam palavras, frases e enunciados válidos.
\prov{sinónimos: gramática}
\prov{relações: parte\_de sintaxe; relaciona virtualizacao; relaciona word; relaciona vida-morte; relaciona universidade\_curriculo\_em\_camadas; relaciona semiotica\_peirce\_triade; relaciona psicanalise; relaciona proposition\_gerador\_de\_escala\_ipropderiv\_com\_tddt; parte\_de linguistica}
\prov{proveniência: curriculo \textperiodcentered{} fato \textperiodcentered{} confiança 1.0 \textperiodcentered{} domínio linguagem \textperiodcentered{} história 10 versões no jornal}

%CRISTAL {"arestas": [{"alvo": "linguagem_de_programacao", "confianca": 1.0, "epistemico": "fato", "origem": "wiki", "rel": "usa"}, {"alvo": "gramatica_gerativa", "confianca": 1.0, "epistemico": "fato", "origem": "wiki", "rel": "usa"}, {"alvo": "hierarquia_de_chomsky", "confianca": 1.0, "epistemico": "fato", "origem": "wiki", "rel": "usa"}], "confianca": 1.0, "contraexemplos": [], "descricao": "A notação formal que define a sintaxe de uma linguagem — a mesma gramática livre-de-contexto da hierarquia de Chomsky.", "epistemico": "fato", "exemplos": [], "id": "gramatica_bnf", "memoria": "semantica", "meta": {"dominio": "programacao", "fonte": ""}, "origem": "wiki", "palavras_chave": [], "sinonimos": [], "tipo": "conceito", "titulo": "Gramática (BNF)"}
\section{Gramática (BNF)}
A notação formal que define a sintaxe de uma linguagem — a mesma gramática livre-de-contexto da hierarquia de Chomsky.
\prov{relações: usa linguagem\_de\_programacao; usa gramatica\_gerativa; usa hierarquia\_de\_chomsky}
\prov{proveniência: wiki \textperiodcentered{} fato \textperiodcentered{} confiança 1.0 \textperiodcentered{} domínio programacao \textperiodcentered{} história 8 versões no jornal}

%CRISTAL {"arestas": [{"alvo": "expressao_regular", "confianca": 1.0, "epistemico": "fato", "origem": "wiki", "rel": "usa"}, {"alvo": "shell", "confianca": 1.0, "epistemico": "fato", "origem": "wiki", "rel": "usa"}], "confianca": 1.0, "contraexemplos": [], "descricao": "Filtra linhas que casam um padrão (expressão regular) — o buscador do fluxo de texto. Do 'g/re/p' do ed.", "epistemico": "fato", "exemplos": [], "id": "grep", "memoria": "semantica", "meta": {"autor": "Thompson", "dominio": "programacao", "fonte": ""}, "origem": "wiki", "palavras_chave": [], "sinonimos": [], "tipo": "conceito", "titulo": "grep"}
\section{grep}
Filtra linhas que casam um padrão (expressão regular) — o buscador do fluxo de texto. Do 'g/re/p' do ed.
\prov{relações: usa expressao\_regular; usa shell}
\prov{proveniência: wiki \textperiodcentered{} fato \textperiodcentered{} confiança 1.0 \textperiodcentered{} domínio programacao \textperiodcentered{} história 3 versões no jornal}

%CRISTAL {"arestas": [{"alvo": "contabilidade_sem_livro", "confianca": 1.0, "epistemico": "fato", "origem": "wiki", "rel": "explica"}, {"alvo": "lastro_do_cofre", "confianca": 1.0, "epistemico": "fato", "origem": "wiki", "rel": "usa"}], "confianca": 1.0, "contraexemplos": [], "descricao": "O estado contábil é lido como mapa hamiltoniano: q=lastro (branco), p=fees (preto), H=a+c²; Liouville OK se bootstrap≈q+p. Pell emitido sem ciclo BTC fechado é 'órbita aberta' (pell_orfaos); entradas da outra rede são 'contaminação cruzada' — ambos vão para quarentena.", "epistemico": "inferencia", "exemplos": [], "id": "hamiltoniano_quarentena", "memoria": "semantica", "meta": {"autor": "broca-so", "dominio": "operacao_so", "fonte": "goldchain/ledger_rede.py:hamiltoniano_estado"}, "origem": "wiki", "palavras_chave": [], "sinonimos": [], "tipo": "conceito", "titulo": "Contabilidade Hamiltoniana (fora do contorno → quarentena)"}
\section{Contabilidade Hamiltoniana (fora do contorno \ensuremath{\to} quarentena)}
O estado contábil é lido como mapa hamiltoniano: q=lastro (branco), p=fees (preto), H=a+c\ensuremath{^{2}}; Liouville OK se bootstrap\ensuremath{\approx}q+p. Pell emitido sem ciclo BTC fechado é 'órbita aberta' (pell\_orfaos); entradas da outra rede são 'contaminação cruzada' — ambos vão para quarentena.
\prov{relações: explica contabilidade\_sem\_livro; usa lastro\_do\_cofre}
\prov{proveniência: wiki \textperiodcentered{} goldchain/ledger\_rede.py:hamiltoniano\_estado \textperiodcentered{} inferencia \textperiodcentered{} confiança 1.0 \textperiodcentered{} domínio operacao\_so \textperiodcentered{} história 3 versões no jornal}

%CRISTAL {"arestas": [{"alvo": "cofre_parabola_hero", "confianca": 1.0, "epistemico": "fato", "origem": "wiki", "rel": "usa"}, {"alvo": "parabola_controla_taxa_tecidos", "confianca": 1.0, "epistemico": "fato", "origem": "wiki", "rel": "usa"}, {"alvo": "symbolic_beta_pipe_so", "confianca": 1.0, "epistemico": "fato", "origem": "wiki", "rel": "explica"}, {"alvo": "tiffany", "confianca": 1.0, "epistemico": "fato", "origem": "wiki", "rel": "parte_de"}], "confianca": 1.0, "contraexemplos": [], "descricao": "Correção honesta de um erro meu: BROCA_DIFUSAO_FRAC NÃO é um split ouro/prata (cheguei a modelar 'ótimo frac=0,5', errado). É um TETO de recurso: 0,90 limitava a 90%. O certo é TIRÁ-LO (=1) — os recursos se AUTO-REGULAM por trabalho (consumidos quando há job, voltam a 0 ocioso), então não precisa limitar; tirar o teto SOBE a hashrate (3,4→3,8 MH/s). A cunhagem de Pell sobe por outra via: as 12 LANES/tecidos por floco (multifractal: 1 floco de ouro → 12 leituras de tecido → 12 Pell), NÃO por baixar o frac. E as CONVERSÕES (c², ν) VÊM DA PARÁBOLA — o canal_c2 e os modos já são computados pelo metal (parabola_auditar no pacote seed); recomputá-los na mão (SVD por lane) esquentava a CPU, derrubava a hash e travava o olho. O olho lê c²=canal_c2 (ouro), a=1−c² (prata) direto da parábola. [D] corrigido e no ar (GEX).", "epistemico": "fato", "exemplos": [], "id": "hash_pell_inverso_parabola", "memoria": "semantica", "meta": {"autor": "Lopes/broca-so", "dominio": "goldchain", "fonte": "goldchain/symbolic.py (12 lanes, sem SVD); neuronio/difusao_recursos.py (frac default 1.0); olho canal_c2"}, "origem": "wiki", "palavras_chave": [], "sinonimos": [], "tipo": "conceito", "titulo": "Frac é TETO de recurso (→1, sem limite); a Pell sobe pelas LANES, e as conversões vêm da parábola"}
\section{Frac é TETO de recurso (\ensuremath{\to}1, sem limite); a Pell sobe pelas LANES, e as conversões vêm da parábola}
Correção honesta de um erro meu: BROCA\_DIFUSAO\_FRAC NÃO é um split ouro/prata (cheguei a modelar 'ótimo frac=0,5', errado). É um TETO de recurso: 0,90 limitava a 90\%. O certo é TIRÁ-LO (=1) — os recursos se AUTO-REGULAM por trabalho (consumidos quando há job, voltam a 0 ocioso), então não precisa limitar; tirar o teto SOBE a hashrate (3,4\ensuremath{\to}3,8 MH/s). A cunhagem de Pell sobe por outra via: as 12 LANES/tecidos por floco (multifractal: 1 floco de ouro \ensuremath{\to} 12 leituras de tecido \ensuremath{\to} 12 Pell), NÃO por baixar o frac. E as CONVERSÕES (c\ensuremath{^{2}}, \ensuremath{\nu}) VÊM DA PARÁBOLA — o canal\_c2 e os modos já são computados pelo metal (parabola\_auditar no pacote seed); recomputá-los na mão (SVD por lane) esquentava a CPU, derrubava a hash e travava o olho. O olho lê c\ensuremath{^{2}}=canal\_c2 (ouro), a=1\ensuremath{-}c\ensuremath{^{2}} (prata) direto da parábola. [D] corrigido e no ar (GEX).
\prov{relações: usa cofre\_parabola\_hero; usa parabola\_controla\_taxa\_tecidos; explica symbolic\_beta\_pipe\_so; parte\_de tiffany}
\prov{proveniência: wiki \textperiodcentered{} goldchain/symbolic.py (12 lanes, sem SVD); neuronio/difusao\_recursos.py (frac default 1.0); olho canal\_c2 \textperiodcentered{} fato \textperiodcentered{} confiança 1.0 \textperiodcentered{} domínio goldchain \textperiodcentered{} história 15 versões no jornal}

%CRISTAL {"arestas": [{"alvo": "programacao_funcional", "confianca": 1.0, "epistemico": "fato", "origem": "wiki", "rel": "especializa"}, {"alvo": "inferencia_de_tipos", "confianca": 1.0, "epistemico": "fato", "origem": "wiki", "rel": "usa"}, {"alvo": "monada", "confianca": 1.0, "epistemico": "fato", "origem": "wiki", "rel": "usa"}], "confianca": 1.0, "contraexemplos": [], "descricao": "Funcional pura e preguiçosa, com tipagem forte e inferência (Hindley-Milner) e mônadas para efeitos. O laboratório de tipos.", "epistemico": "fato", "exemplos": [], "id": "haskell", "memoria": "semantica", "meta": {"dominio": "programacao", "fonte": ""}, "origem": "wiki", "palavras_chave": [], "sinonimos": [], "tipo": "conceito", "titulo": "Haskell"}
\section{Haskell}
Funcional pura e preguiçosa, com tipagem forte e inferência (Hindley-Milner) e mônadas para efeitos. O laboratório de tipos.
\prov{relações: especializa programacao\_funcional; usa inferencia\_de\_tipos; usa monada}
\prov{proveniência: wiki \textperiodcentered{} fato \textperiodcentered{} confiança 1.0 \textperiodcentered{} domínio programacao \textperiodcentered{} história 8 versões no jornal}

%CRISTAL {"arestas": [{"alvo": "kernel", "confianca": 0.703, "epistemico": "fato", "origem": "manual", "rel": "referencia"}, {"alvo": "ponteiro", "confianca": 1.0, "epistemico": "fato", "origem": "curriculo", "rel": "usa"}, {"alvo": "java", "confianca": 0.5, "epistemico": "fato", "origem": "wiki", "rel": "relaciona"}, {"alvo": "recursao", "confianca": 0.5, "epistemico": "fato", "origem": "wiki", "rel": "relaciona"}, {"alvo": "koch_atlas", "confianca": 0.5, "epistemico": "fato", "origem": "wiki", "rel": "relaciona"}, {"alvo": "vida-morte", "confianca": 0.5, "epistemico": "fato", "origem": "wiki", "rel": "relaciona"}, {"alvo": "word", "confianca": 0.5, "epistemico": "fato", "origem": "wiki", "rel": "relaciona"}, {"alvo": "sequencia_pow_nativa", "confianca": 0.5, "epistemico": "fato", "origem": "wiki", "rel": "relaciona"}, {"alvo": "readme", "confianca": 0.5, "epistemico": "fato", "origem": "wiki", "rel": "relaciona"}, {"alvo": "pipeline_de_ingestao", "confianca": 0.5, "epistemico": "fato", "origem": "wiki", "rel": "relaciona"}, {"alvo": "vontade_de_poder", "confianca": 0.5, "epistemico": "fato", "origem": "wiki", "rel": "relaciona"}, {"alvo": "pell", "confianca": 0.5, "epistemico": "fato", "origem": "wiki", "rel": "relaciona"}, {"alvo": "virtualizacao", "confianca": 0.5, "epistemico": "fato", "origem": "wiki", "rel": "relaciona"}, {"alvo": "proposition_escala_fasepropescala_xix", "confianca": 0.5, "epistemico": "fato", "origem": "wiki", "rel": "relaciona"}], "confianca": 1.0, "contraexemplos": [], "descricao": "Área dinâmica; alocação manual ou GC.", "epistemico": "fato", "exemplos": [], "id": "heap", "memoria": "semantica", "meta": {"dominio": "programacao", "nivel": 6, "ordem": 5}, "origem": "curriculo", "palavras_chave": [], "sinonimos": [], "tipo": "conceito", "titulo": "heap"}
\section{heap}
Área dinâmica; alocação manual ou GC.
\prov{relações: referencia kernel; usa ponteiro; relaciona java; relaciona recursao; relaciona koch\_atlas; relaciona vida-morte; relaciona word; relaciona sequencia\_pow\_nativa; relaciona readme; relaciona pipeline\_de\_ingestao; relaciona vontade\_de\_poder; relaciona pell; relaciona virtualizacao; relaciona proposition\_escala\_fasepropescala\_xix}
\prov{proveniência: curriculo \textperiodcentered{} fato \textperiodcentered{} confiança 1.0 \textperiodcentered{} domínio programacao \textperiodcentered{} história 19 versões no jornal}

%CRISTAL {"arestas": [{"alvo": "kubernetes", "confianca": 1.0, "epistemico": "fato", "origem": "wiki", "rel": "usa"}, {"alvo": "yaml", "confianca": 1.0, "epistemico": "fato", "origem": "wiki", "rel": "usa"}], "confianca": 1.0, "contraexemplos": [], "descricao": "O 'gerenciador de pacotes' do Kubernetes: charts (templates YAML) que empacotam e parametrizam uma aplicação inteira para o cluster.", "epistemico": "fato", "exemplos": [], "id": "helm", "memoria": "semantica", "meta": {"dominio": "programacao", "fonte": ""}, "origem": "wiki", "palavras_chave": [], "sinonimos": [], "tipo": "conceito", "titulo": "Helm"}
\section{Helm}
O 'gerenciador de pacotes' do Kubernetes: charts (templates YAML) que empacotam e parametrizam uma aplicação inteira para o cluster.
\prov{relações: usa kubernetes; usa yaml}
\prov{proveniência: wiki \textperiodcentered{} fato \textperiodcentered{} confiança 1.0 \textperiodcentered{} domínio programacao \textperiodcentered{} história 3 versões no jornal}

%CRISTAL {"arestas": [{"alvo": "linguagem_de_marcacao", "confianca": 1.0, "epistemico": "fato", "origem": "wiki", "rel": "usa"}, {"alvo": "tiffany", "confianca": 1.0, "epistemico": "fato", "origem": "wiki", "rel": "explica"}, {"alvo": "internet", "confianca": 1.0, "epistemico": "fato", "origem": "wiki", "rel": "relaciona"}], "confianca": 1.0, "contraexemplos": [], "descricao": "Texto com ligações que levam a outro texto — a leitura deixa de ser linear e vira um grafo navegável. Nelson cunhou; a web realizou.", "epistemico": "fato", "exemplos": [], "id": "hipertexto", "memoria": "semantica", "meta": {"autor": "Nelson", "dominio": "programacao", "fonte": ""}, "origem": "wiki", "palavras_chave": [], "sinonimos": [], "tipo": "conceito", "titulo": "Hipertexto"}
\section{Hipertexto}
Texto com ligações que levam a outro texto — a leitura deixa de ser linear e vira um grafo navegável. Nelson cunhou; a web realizou.
\prov{relações: usa linguagem\_de\_marcacao; explica tiffany; relaciona internet}
\prov{proveniência: wiki \textperiodcentered{} fato \textperiodcentered{} confiança 1.0 \textperiodcentered{} domínio programacao \textperiodcentered{} história 4 versões no jornal}

%CRISTAL {"arestas": [{"alvo": "lisp", "confianca": 1.0, "epistemico": "fato", "origem": "wiki", "rel": "explica"}, {"alvo": "metaprogramacao", "confianca": 1.0, "epistemico": "fato", "origem": "wiki", "rel": "usa"}], "confianca": 1.0, "contraexemplos": [], "descricao": "A linguagem em que o código é um dado da própria linguagem (Lisp: S-expressões) — o programa pode ler e reescrever a si mesmo.", "epistemico": "fato", "exemplos": [], "id": "homoiconicidade", "memoria": "semantica", "meta": {"dominio": "programacao", "fonte": ""}, "origem": "wiki", "palavras_chave": [], "sinonimos": [], "tipo": "conceito", "titulo": "Homoiconicidade"}
\section{Homoiconicidade}
A linguagem em que o código é um dado da própria linguagem (Lisp: S-expressões) — o programa pode ler e reescrever a si mesmo.
\prov{relações: explica lisp; usa metaprogramacao}
\prov{proveniência: wiki \textperiodcentered{} fato \textperiodcentered{} confiança 1.0 \textperiodcentered{} domínio programacao \textperiodcentered{} história 6 versões no jornal}

%CRISTAL {"arestas": [{"alvo": "react", "confianca": 1.0, "epistemico": "fato", "origem": "wiki", "rel": "parte_de"}, {"alvo": "closure", "confianca": 1.0, "epistemico": "fato", "origem": "wiki", "rel": "usa"}], "confianca": 1.0, "contraexemplos": [], "descricao": "Funções que dão estado e ciclo de vida a componentes de função (useState, useEffect) — closures com memória entre renders. Os 'hooks' do React.", "epistemico": "fato", "exemplos": [], "id": "hook_react", "memoria": "semantica", "meta": {"dominio": "programacao", "fonte": ""}, "origem": "wiki", "palavras_chave": [], "sinonimos": [], "tipo": "conceito", "titulo": "Hook (React)"}
\section{Hook (React)}
Funções que dão estado e ciclo de vida a componentes de função (useState, useEffect) — closures com memória entre renders. Os 'hooks' do React.
\prov{relações: parte\_de react; usa closure}
\prov{proveniência: wiki \textperiodcentered{} fato \textperiodcentered{} confiança 1.0 \textperiodcentered{} domínio programacao \textperiodcentered{} história 6 versões no jornal}

%CRISTAL {"arestas": [{"alvo": "hopfield", "confianca": 1.0, "epistemico": "fato", "origem": "wiki", "rel": "especializa"}, {"alvo": "colapso_koch_tiffany", "confianca": 1.0, "epistemico": "fato", "origem": "wiki", "rel": "usa"}, {"alvo": "peirce_celulas_canais", "confianca": 1.0, "epistemico": "fato", "origem": "wiki", "rel": "relaciona"}], "confianca": 1.0, "contraexemplos": [], "descricao": "code/arquitetura.c: módulos Hopfield de 32 neurônios (recall descendo a energia ao atrator, sinapse Hebbian) viram eles próprios neurônios de uma super-rede; o mundo real entra quantizado pela borda z↦z² mod 1009 (2-para-1, sem volta), e a camada de cima corrige a de baixo.", "epistemico": "fato", "exemplos": [], "id": "hopfield_recursiva_koch", "memoria": "semantica", "meta": {"autor": "broca-so", "dominio": "operacao_so", "fonte": "code/arquitetura.c"}, "origem": "wiki", "palavras_chave": [], "sinonimos": [], "tipo": "conceito", "titulo": "Hopfield Recursiva com Borda de Koch"}
\section{Hopfield Recursiva com Borda de Koch}
code/arquitetura.c: módulos Hopfield de 32 neurônios (recall descendo a energia ao atrator, sinapse Hebbian) viram eles próprios neurônios de uma super-rede; o mundo real entra quantizado pela borda z\ensuremath{\mapsto}z\ensuremath{^{2}} mod 1009 (2-para-1, sem volta), e a camada de cima corrige a de baixo.
\prov{relações: especializa hopfield; usa colapso\_koch\_tiffany; relaciona peirce\_celulas\_canais}
\prov{proveniência: wiki \textperiodcentered{} code/arquitetura.c \textperiodcentered{} fato \textperiodcentered{} confiança 1.0 \textperiodcentered{} domínio operacao\_so \textperiodcentered{} história 4 versões no jornal}

%CRISTAL {"arestas": [{"alvo": "htlc", "confianca": 1.0, "epistemico": "fato", "origem": "wiki", "rel": "explica"}, {"alvo": "algebra_dos_contratos", "confianca": 1.0, "epistemico": "fato", "origem": "wiki", "rel": "usa"}, {"alvo": "certificado_tempo_semente", "confianca": 1.0, "epistemico": "fato", "origem": "wiki", "rel": "relaciona"}, {"alvo": "goldchain_como_rede", "confianca": 1.0, "epistemico": "fato", "origem": "wiki", "rel": "usa"}], "confianca": 1.0, "contraexemplos": [], "descricao": "O contrato on-chain é um script OP_IF OP_SHA256 <hashlock> OP_EQUALVERIFY <resg> OP_CHECKSIG OP_ELSE <timelock> OP_CHECKLOCKTIMEVERIFY OP_DROP <refund> OP_CHECKSIG OP_ENDIF num endereço P2WSH; hashlock=sha256(segredo), segredo=sha256(banda_raiz) do PoW; comissão fixa 1%. Broadcast via mempool.space.", "epistemico": "fato", "exemplos": [], "id": "htlc_p2wsh_real", "memoria": "semantica", "meta": {"autor": "broca-so", "dominio": "operacao_so", "fonte": "goldchain/btc_backend.py:htlc_script"}, "origem": "wiki", "palavras_chave": [], "sinonimos": [], "tipo": "conceito", "titulo": "HTLC P2WSH Real (script + broadcast)"}
\section{HTLC P2WSH Real (script + broadcast)}
O contrato on-chain é um script OP\_IF OP\_SHA256 <hashlock> OP\_EQUALVERIFY <resg> OP\_CHECKSIG OP\_ELSE <timelock> OP\_CHECKLOCKTIMEVERIFY OP\_DROP <refund> OP\_CHECKSIG OP\_ENDIF num endereço P2WSH; hashlock=sha256(segredo), segredo=sha256(banda\_raiz) do PoW; comissão fixa 1\%. Broadcast via mempool.space.
\prov{relações: explica htlc; usa algebra\_dos\_contratos; relaciona certificado\_tempo\_semente; usa goldchain\_como\_rede}
\prov{proveniência: wiki \textperiodcentered{} goldchain/btc\_backend.py:htlc\_script \textperiodcentered{} fato \textperiodcentered{} confiança 1.0 \textperiodcentered{} domínio operacao\_so \textperiodcentered{} história 5 versões no jornal}

%CRISTAL {"arestas": [{"alvo": "linguagem_de_marcacao", "confianca": 1.0, "epistemico": "fato", "origem": "wiki", "rel": "especializa"}, {"alvo": "hipertexto", "confianca": 1.0, "epistemico": "fato", "origem": "wiki", "rel": "usa"}, {"alvo": "internet", "confianca": 1.0, "epistemico": "fato", "origem": "wiki", "rel": "parte_de"}, {"alvo": "http", "confianca": 1.0, "epistemico": "fato", "origem": "wiki", "rel": "usa"}], "confianca": 1.0, "contraexemplos": [], "descricao": "Berners-Lee: a linguagem de marcação que estrutura as páginas da web — títulos, parágrafos, listas e, sobretudo, os LINKS do hipertexto.", "epistemico": "fato", "exemplos": [], "id": "html", "memoria": "semantica", "meta": {"autor": "Berners-Lee", "dominio": "programacao", "fonte": ""}, "origem": "wiki", "palavras_chave": [], "sinonimos": [], "tipo": "conceito", "titulo": "HTML"}
\section{HTML}
Berners-Lee: a linguagem de marcação que estrutura as páginas da web — títulos, parágrafos, listas e, sobretudo, os LINKS do hipertexto.
\prov{relações: especializa linguagem\_de\_marcacao; usa hipertexto; parte\_de internet; usa http}
\prov{proveniência: wiki \textperiodcentered{} fato \textperiodcentered{} confiança 1.0 \textperiodcentered{} domínio programacao \textperiodcentered{} história 5 versões no jornal}

%CRISTAL {"arestas": [{"alvo": "protocolo_tcp", "confianca": 1.0, "epistemico": "fato", "origem": "wiki", "rel": "usa"}, {"alvo": "internet", "confianca": 1.0, "epistemico": "fato", "origem": "wiki", "rel": "parte_de"}], "confianca": 1.0, "contraexemplos": [], "descricao": "O protocolo da web: requisição-resposta sobre TCP, sem estado. GET/POST, códigos de status, cabeçalhos — a língua dos navegadores.", "epistemico": "fato", "exemplos": [], "id": "http", "memoria": "semantica", "meta": {"autor": "Berners-Lee", "dominio": "programacao", "fonte": ""}, "origem": "wiki", "palavras_chave": [], "sinonimos": [], "tipo": "conceito", "titulo": "HTTP"}
\section{HTTP}
O protocolo da web: requisição-resposta sobre TCP, sem estado. GET/POST, códigos de status, cabeçalhos — a língua dos navegadores.
\prov{relações: usa protocolo\_tcp; parte\_de internet}
\prov{proveniência: wiki \textperiodcentered{} fato \textperiodcentered{} confiança 1.0 \textperiodcentered{} domínio programacao \textperiodcentered{} história 6 versões no jornal}

%CRISTAL {"arestas": [{"alvo": "colapso_koch_tiffany", "confianca": 1.0, "epistemico": "fato", "origem": "wiki", "rel": "usa"}, {"alvo": "tiffany", "confianca": 1.0, "epistemico": "fato", "origem": "wiki", "rel": "relaciona"}], "confianca": 1.0, "contraexemplos": [], "descricao": "code/prosa.py: cada texto vira um vetor de 256 bits acumulando _koch(z)=z² mod 1009 sobre os bytes, empacotado em 32 bytes; o .idx guarda offset+lang+impressão(32B) por nó no disco.", "epistemico": "fato", "exemplos": [], "id": "impressao_koch_32bytes", "memoria": "semantica", "meta": {"autor": "broca-so", "dominio": "operacao_so", "fonte": "code/prosa.py"}, "origem": "wiki", "palavras_chave": [], "sinonimos": [], "tipo": "conceito", "titulo": "Impressão Koch — 32 Bytes por Nó"}
\section{Impressão Koch — 32 Bytes por Nó}
code/prosa.py: cada texto vira um vetor de 256 bits acumulando \_koch(z)=z\ensuremath{^{2}} mod 1009 sobre os bytes, empacotado em 32 bytes; o .idx guarda offset+lang+impressão(32B) por nó no disco.
\prov{relações: usa colapso\_koch\_tiffany; relaciona tiffany}
\prov{proveniência: wiki \textperiodcentered{} code/prosa.py \textperiodcentered{} fato \textperiodcentered{} confiança 1.0 \textperiodcentered{} domínio operacao\_so \textperiodcentered{} história 3 versões no jornal}

%CRISTAL {"arestas": [{"alvo": "impressao_koch_32bytes", "confianca": 1.0, "epistemico": "fato", "origem": "wiki", "rel": "especializa"}, {"alvo": "entropia_shannon", "confianca": 1.0, "epistemico": "fato", "origem": "wiki", "rel": "relaciona"}, {"alvo": "koch_memoria", "confianca": 1.0, "epistemico": "fato", "origem": "wiki", "rel": "usa"}, {"alvo": "colapso_ping_pong_overlap", "confianca": 1.0, "epistemico": "fato", "origem": "wiki", "rel": "usa"}, {"alvo": "floco_calor_dissipado", "confianca": 1.0, "epistemico": "fato", "origem": "wiki", "rel": "relaciona"}], "confianca": 1.0, "contraexemplos": [], "descricao": "Cada texto vira um vetor de 256 bits acumulando _koch(z)=z² mod 1009 sobre os bytes, empacotado em 32 B; o .idx guarda só offset+lang+impressão(32B) por nó ('DADOS no disco, GATOS na RAM'). É compressão COM PERDA/hash: z↦z² é 2-para-1 sem volta (~80 bits de Hamming entre vizinhos, colisões) — perde a reta associativa (recall por similaridade) e guarda o floco (o historial da órbita). A fala é colapso por overlap (recuperação, não geração), não uma decodificação sem perdas. [D].", "epistemico": "inferencia", "exemplos": [], "id": "impressao_koch_compressao", "memoria": "semantica", "meta": {"autor": "broca-so", "dominio": "operacao_so", "fonte": "code/prosa.py; ula/KOCH_MEMORIA.md; impressao_koch_32bytes"}, "origem": "wiki", "palavras_chave": [], "sinonimos": [], "tipo": "conceito", "titulo": "A impressão Koch de 32 bytes = compressão/hash com perda"}
\section{A impressão Koch de 32 bytes = compressão/hash com perda}
Cada texto vira um vetor de 256 bits acumulando \_koch(z)=z\ensuremath{^{2}} mod 1009 sobre os bytes, empacotado em 32 B; o .idx guarda só offset+lang+impressão(32B) por nó ('DADOS no disco, GATOS na RAM'). É compressão COM PERDA/hash: z\ensuremath{\mapsto}z\ensuremath{^{2}} é 2-para-1 sem volta (\~{}80 bits de Hamming entre vizinhos, colisões) — perde a reta associativa (recall por similaridade) e guarda o floco (o historial da órbita). A fala é colapso por overlap (recuperação, não geração), não uma decodificação sem perdas. [D].
\prov{relações: especializa impressao\_koch\_32bytes; relaciona entropia\_shannon; usa koch\_memoria; usa colapso\_ping\_pong\_overlap; relaciona floco\_calor\_dissipado}
\prov{proveniência: wiki \textperiodcentered{} code/prosa.py; ula/KOCH\_MEMORIA.md; impressao\_koch\_32bytes \textperiodcentered{} inferencia \textperiodcentered{} confiança 1.0 \textperiodcentered{} domínio operacao\_so \textperiodcentered{} história 6 versões no jornal}

%CRISTAL {"arestas": [{"alvo": "programacao_funcional", "confianca": 1.0, "epistemico": "fato", "origem": "wiki", "rel": "usa"}], "confianca": 1.0, "contraexemplos": [], "descricao": "Dados que não mudam após criados — elimina classes inteiras de bugs de concorrência; o coração do estilo funcional.", "epistemico": "fato", "exemplos": [], "id": "imutabilidade", "memoria": "semantica", "meta": {"dominio": "programacao", "fonte": ""}, "origem": "wiki", "palavras_chave": [], "sinonimos": [], "tipo": "conceito", "titulo": "Imutabilidade"}
\section{Imutabilidade}
Dados que não mudam após criados — elimina classes inteiras de bugs de concorrência; o coração do estilo funcional.
\prov{relações: usa programacao\_funcional}
\prov{proveniência: wiki \textperiodcentered{} fato \textperiodcentered{} confiança 1.0 \textperiodcentered{} domínio programacao \textperiodcentered{} história 4 versões no jornal}

%CRISTAL {"arestas": [{"alvo": "banco_de_dados", "confianca": 1.0, "epistemico": "fato", "origem": "wiki", "rel": "usa"}, {"alvo": "algoritmo", "confianca": 1.0, "epistemico": "fato", "origem": "wiki", "rel": "usa"}], "confianca": 1.0, "contraexemplos": [], "descricao": "Uma estrutura extra que acha uma linha sem varrer a tabela toda — troca espaço e escrita por leitura veloz. A B-tree onipresente.", "epistemico": "fato", "exemplos": [], "id": "indice_de_banco", "memoria": "semantica", "meta": {"dominio": "programacao", "fonte": ""}, "origem": "wiki", "palavras_chave": [], "sinonimos": [], "tipo": "conceito", "titulo": "Índice (B-Tree)"}
\section{Índice (B-Tree)}
Uma estrutura extra que acha uma linha sem varrer a tabela toda — troca espaço e escrita por leitura veloz. A B-tree onipresente.
\prov{relações: usa banco\_de\_dados; usa algoritmo}
\prov{proveniência: wiki \textperiodcentered{} fato \textperiodcentered{} confiança 1.0 \textperiodcentered{} domínio programacao \textperiodcentered{} história 6 versões no jornal}

%CRISTAL {"arestas": [{"alvo": "linguagem_gatos_sub_turing_honesto", "confianca": 1.0, "epistemico": "fato", "origem": "wiki", "rel": "contradiz"}, {"alvo": "interpretador_nucleo_gatos", "confianca": 1.0, "epistemico": "fato", "origem": "wiki", "rel": "especializa"}, {"alvo": "transformada_metalica", "confianca": 1.0, "epistemico": "fato", "origem": "wiki", "rel": "usa"}, {"alvo": "expressividade_sub_turing", "confianca": 1.0, "epistemico": "fato", "origem": "wiki", "rel": "contradiz"}, {"alvo": "maquina_de_turing", "confianca": 1.0, "epistemico": "fato", "origem": "wiki", "rel": "usa"}, {"alvo": "linguagem_matriz_gatos", "confianca": 1.0, "epistemico": "fato", "origem": "wiki", "rel": "parte_de"}, {"alvo": "tiffany", "confianca": 1.0, "epistemico": "fato", "origem": "wiki", "rel": "parte_de"}], "confianca": 1.0, "contraexemplos": [], "descricao": "linguagem/universal.py: a conjectura [C] do paper (§6) RESOLVIDA — e honesta. A prova sub-Turing (expressividade_sub_turing) foi precisa: faltavam INDIREÇÃO (registrador nunca selecionava endereço) e memória ILIMITADA. A INDIREÇÃO POSICIONAL é exatamente isso: o VALOR de um registrador É um endereço (uma POSIÇÃO), navegado pela aritmética do gato. Adicionadas as duas: loadi rd ra (rd←mem[valor de ra]) e storei ra rs sobre memória ilimitada (dict esparso, valores sem cota — não mod 2¹²⁸). PROVA rodando: escrevemos NESTE assembly um INTERPRETADOR UNIVERSAL de máquina de 2 contadores de Minsky (Turing-completa) — um driver de tamanho FIXO (~25 instruções) que executa um programa 2CM ARBITRÁRIO da memória por indireção (o PC indexa o programa) com salto por dado (pc←alvo). Verificado: adição (3+4=7, 10+25=35), o MESMO driver rodando 'dobrar' (c1=6→12), e ilimitado (c1=100000→100000). Driver fixo + programa ilimitado + indireção + contadores ilimitados = universal. HONESTO: a transcendência vem de ADICIONAR indireção + memória ilimitada (o que a prova apontava faltar) — resultado padrão (RAM/2CM é TC); o gato NÃO transcende sozinho, ele NAVEGA as posições; a indireção posicional é quem lifta. [F] provado.", "epistemico": "fato", "exemplos": [], "id": "indirecao_posicional_transcende_turing", "memoria": "semantica", "meta": {"autor": "Lopes/broca-so", "dominio": "computacao", "fonte": "linguagem/universal.py (rodar_ram loadi/storei, DRIVER_2CM); verificado: 2CM adição/dobrar/ilimitado; Minsky 2CM = Turing-completo"}, "origem": "wiki", "palavras_chave": [], "sinonimos": [], "tipo": "conceito", "titulo": "A indireção posicional TRANSCENDE o sub-Turing — provado por um 2CM universal [F]"}
\section{A indireção posicional TRANSCENDE o sub-Turing — provado por um 2CM universal [F]}
linguagem/universal.py: a conjectura [C] do paper (§6) RESOLVIDA — e honesta. A prova sub-Turing (expressividade\_sub\_turing) foi precisa: faltavam INDIREÇÃO (registrador nunca selecionava endereço) e memória ILIMITADA. A INDIREÇÃO POSICIONAL é exatamente isso: o VALOR de um registrador É um endereço (uma POSIÇÃO), navegado pela aritmética do gato. Adicionadas as duas: loadi rd ra (rd\ensuremath{\leftarrow}mem[valor de ra]) e storei ra rs sobre memória ilimitada (dict esparso, valores sem cota — não mod 2\ensuremath{^{128}}). PROVA rodando: escrevemos NESTE assembly um INTERPRETADOR UNIVERSAL de máquina de 2 contadores de Minsky (Turing-completa) — um driver de tamanho FIXO (\~{}25 instruções) que executa um programa 2CM ARBITRÁRIO da memória por indireção (o PC indexa o programa) com salto por dado (pc\ensuremath{\leftarrow}alvo). Verificado: adição (3+4=7, 10+25=35), o MESMO driver rodando 'dobrar' (c1=6\ensuremath{\to}12), e ilimitado (c1=100000\ensuremath{\to}100000). Driver fixo + programa ilimitado + indireção + contadores ilimitados = universal. HONESTO: a transcendência vem de ADICIONAR indireção + memória ilimitada (o que a prova apontava faltar) — resultado padrão (RAM/2CM é TC); o gato NÃO transcende sozinho, ele NAVEGA as posições; a indireção posicional é quem lifta. [F] provado.
\prov{relações: contradiz linguagem\_gatos\_sub\_turing\_honesto; especializa interpretador\_nucleo\_gatos; usa transformada\_metalica; contradiz expressividade\_sub\_turing; usa maquina\_de\_turing; parte\_de linguagem\_matriz\_gatos; parte\_de tiffany}
\prov{proveniência: wiki \textperiodcentered{} linguagem/universal.py (rodar\_ram loadi/storei, DRIVER\_2CM); verificado: 2CM adição/dobrar/ilimitado; Minsky 2CM = Turing-completo \textperiodcentered{} fato \textperiodcentered{} confiança 1.0 \textperiodcentered{} domínio computacao \textperiodcentered{} história 48 versões no jornal}

%CRISTAL {"arestas": [{"alvo": "sistema_de_tipos", "confianca": 1.0, "epistemico": "fato", "origem": "wiki", "rel": "parte_de"}], "confianca": 1.0, "contraexemplos": [], "descricao": "O compilador deduz os tipos sem anotação explícita (Hindley-Milner) — segurança estática sem a verbosidade.", "epistemico": "fato", "exemplos": [], "id": "inferencia_de_tipos", "memoria": "semantica", "meta": {"dominio": "programacao", "fonte": ""}, "origem": "wiki", "palavras_chave": [], "sinonimos": [], "tipo": "conceito", "titulo": "Inferência de Tipos"}
\section{Inferência de Tipos}
O compilador deduz os tipos sem anotação explícita (Hindley-Milner) — segurança estática sem a verbosidade.
\prov{relações: parte\_de sistema\_de\_tipos}
\prov{proveniência: wiki \textperiodcentered{} fato \textperiodcentered{} confiança 1.0 \textperiodcentered{} domínio programacao \textperiodcentered{} história 4 versões no jornal}

%CRISTAL {"arestas": [{"alvo": "cristal_vs_cha", "confianca": 1.0, "epistemico": "fato", "origem": "wiki", "rel": "usa"}, {"alvo": "kolmogorov_complexidade", "confianca": 1.0, "epistemico": "fato", "origem": "wiki", "rel": "relaciona"}, {"alvo": "capacidade_canal_shannon", "confianca": 1.0, "epistemico": "fato", "origem": "wiki", "rel": "relaciona"}, {"alvo": "expressividade_sub_turing", "confianca": 1.0, "epistemico": "fato", "origem": "wiki", "rel": "usa"}], "confianca": 1.0, "contraexemplos": [], "descricao": "Achado honesto: a série Lopes/broca-so usa Shannon e Kolmogorov como SABOR qualitativo, não como teoria formalizada. cristal-vs-chá é Kolmogorov informal (lei compressível=derivável vs identidade incompressível=só-se-guarda; chá médio 0,85 medido); entropia=log φ é taxa de informação topológica (verificada); a entropia espectral separa cristal de chá empiricamente. MAS: NÃO há capacidade de canal (C=B log(1+S/N)), NEM K(x) formal/teorema da invariância, NEM P vs NP como nó. O que a série tem de complexidade é a máquina SUB-TURING (decidibilidade estrutural comprada ao custo da universalidade) e a indecidibilidade de NS (conjectural). Registrar a lacuna, não maquiar.", "epistemico": "inferencia", "exemplos": [], "id": "info_lopes_sabor_nao_formal", "memoria": "semantica", "meta": {"autor": "broca-so", "dominio": "operacao_so", "fonte": "cristal_vs_cha; ula/EXPRESSIVIDADE.md; síntese honesta"}, "origem": "wiki", "palavras_chave": [], "sinonimos": [], "tipo": "conceito", "titulo": "Shannon/Kolmogorov na série: sabor, não formalização (honesto)"}
\section{Shannon/Kolmogorov na série: sabor, não formalização (honesto)}
Achado honesto: a série Lopes/broca-so usa Shannon e Kolmogorov como SABOR qualitativo, não como teoria formalizada. cristal-vs-chá é Kolmogorov informal (lei compressível=derivável vs identidade incompressível=só-se-guarda; chá médio 0,85 medido); entropia=log \ensuremath{\varphi} é taxa de informação topológica (verificada); a entropia espectral separa cristal de chá empiricamente. MAS: NÃO há capacidade de canal (C=B log(1+S/N)), NEM K(x) formal/teorema da invariância, NEM P vs NP como nó. O que a série tem de complexidade é a máquina SUB-TURING (decidibilidade estrutural comprada ao custo da universalidade) e a indecidibilidade de NS (conjectural). Registrar a lacuna, não maquiar.
\prov{relações: usa cristal\_vs\_cha; relaciona kolmogorov\_complexidade; relaciona capacidade\_canal\_shannon; usa expressividade\_sub\_turing}
\prov{proveniência: wiki \textperiodcentered{} cristal\_vs\_cha; ula/EXPRESSIVIDADE.md; síntese honesta \textperiodcentered{} inferencia \textperiodcentered{} confiança 1.0 \textperiodcentered{} domínio operacao\_so \textperiodcentered{} história 5 versões no jornal}

%CRISTAL {"arestas": [{"alvo": "informacao_assimetrica", "confianca": 1.0, "epistemico": "fato", "origem": "wiki", "rel": "relaciona"}, {"alvo": "poker_jogo", "confianca": 1.0, "epistemico": "fato", "origem": "wiki", "rel": "explica"}], "confianca": 1.0, "contraexemplos": [], "descricao": "Jogos onde você NÃO vê tudo — cartas ocultas, estados escondidos. Muda o problema: não basta a melhor jogada, é preciso uma estratégia mista imprevisível (blefar) e raciocinar sobre o que o outro não sabe. O pôquer contra o xadrez.", "epistemico": "fato", "exemplos": [], "id": "informacao_imperfeita", "memoria": "semantica", "meta": {"dominio": "ia", "fonte": "jogos e IA"}, "origem": "wiki", "palavras_chave": [], "sinonimos": [], "tipo": "conceito", "titulo": "Informação Imperfeita"}
\section{Informação Imperfeita}
Jogos onde você NÃO vê tudo — cartas ocultas, estados escondidos. Muda o problema: não basta a melhor jogada, é preciso uma estratégia mista imprevisível (blefar) e raciocinar sobre o que o outro não sabe. O pôquer contra o xadrez.
\prov{relações: relaciona informacao\_assimetrica; explica poker\_jogo}
\prov{proveniência: wiki \textperiodcentered{} jogos e IA \textperiodcentered{} fato \textperiodcentered{} confiança 1.0 \textperiodcentered{} domínio ia \textperiodcentered{} história 3 versões no jornal}

%CRISTAL {"arestas": [{"alvo": "devops", "confianca": 1.0, "epistemico": "fato", "origem": "wiki", "rel": "parte_de"}, {"alvo": "controle_de_versao", "confianca": 1.0, "epistemico": "fato", "origem": "wiki", "rel": "usa"}], "confianca": 1.0, "contraexemplos": [], "descricao": "Descrever servidores, redes e serviços em arquivos versionados (Terraform, Ansible) — a infra reproduzível e revisável como software.", "epistemico": "inferencia", "exemplos": [], "id": "infraestrutura_como_codigo", "memoria": "semantica", "meta": {"dominio": "programacao", "fonte": ""}, "origem": "wiki", "palavras_chave": [], "sinonimos": [], "tipo": "conceito", "titulo": "Infraestrutura como Código"}
\section{Infraestrutura como Código}
Descrever servidores, redes e serviços em arquivos versionados (Terraform, Ansible) — a infra reproduzível e revisável como software.
\prov{relações: parte\_de devops; usa controle\_de\_versao}
\prov{proveniência: wiki \textperiodcentered{} inferencia \textperiodcentered{} confiança 1.0 \textperiodcentered{} domínio programacao \textperiodcentered{} história 6 versões no jornal}

%CRISTAL {"arestas": [{"alvo": "grafo_append_only", "confianca": 1.0, "epistemico": "fato", "origem": "wiki", "rel": "usa"}, {"alvo": "operacao_broca_so", "confianca": 1.0, "epistemico": "fato", "origem": "wiki", "rel": "explica"}, {"alvo": "codigo_goldchain_minerar", "confianca": 1.0, "epistemico": "fato", "origem": "wiki", "rel": "explica"}, {"alvo": "codigo_neuronio_experimentos_tiffany", "confianca": 1.0, "epistemico": "fato", "origem": "wiki", "rel": "explica"}], "confianca": 1.0, "contraexemplos": [], "descricao": "ingestcodigo lê a docstring (via ast) dos módulos.py e cria nós codigodomᵢₙᵢostem (origem=código) ligados ao hub conceitual — o órgão de memória ingere o órgão que o hospeda. É daqui que nascem todos os pontos soltos codigo_*.", "epistemico": "fato", "exemplos": [], "id": "ingest_codigo_autoreferencial", "memoria": "semantica", "meta": {"autor": "broca-so", "dominio": "operacao_so", "fonte": "conversa/conhecimento/ingest_codigo.py"}, "origem": "wiki", "palavras_chave": [], "sinonimos": [], "tipo": "conceito", "titulo": "Ingest do Próprio Código (o grafo come a si mesmo)"}
\section{Ingest do Próprio Código (o grafo come a si mesmo)}
ingestcodigo lê a docstring (via ast) dos módulos.py e cria nós codigodom\ensuremath{_{ini}}ostem (origem=código) ligados ao hub conceitual — o órgão de memória ingere o órgão que o hospeda. É daqui que nascem todos os pontos soltos codigo\_*.
\prov{relações: usa grafo\_append\_only; explica operacao\_broca\_so; explica codigo\_goldchain\_minerar; explica codigo\_neuronio\_experimentos\_tiffany}
\prov{proveniência: wiki \textperiodcentered{} conversa/conhecimento/ingest\_codigo.py \textperiodcentered{} fato \textperiodcentered{} confiança 1.0 \textperiodcentered{} domínio operacao\_so \textperiodcentered{} história 6 versões no jornal}

%CRISTAL {"arestas": [{"alvo": "arquivo_de_configuracao", "confianca": 1.0, "epistemico": "fato", "origem": "wiki", "rel": "especializa"}], "confianca": 1.0, "contraexemplos": [], "descricao": "O formato de configuração clássico: seções em [colchetes] e pares chave=valor — simples e onipresente, sem aninhamento profundo.", "epistemico": "fato", "exemplos": [], "id": "ini", "memoria": "semantica", "meta": {"dominio": "programacao", "fonte": ""}, "origem": "wiki", "palavras_chave": [], "sinonimos": [], "tipo": "conceito", "titulo": "INI"}
\section{INI}
O formato de configuração clássico: seções em [colchetes] e pares chave=valor — simples e onipresente, sem aninhamento profundo.
\prov{relações: especializa arquivo\_de\_configuracao}
\prov{proveniência: wiki \textperiodcentered{} fato \textperiodcentered{} confiança 1.0 \textperiodcentered{} domínio programacao \textperiodcentered{} história 2 versões no jornal}

%CRISTAL {"arestas": [{"alvo": "vulnerabilidade_exploit", "confianca": 1.0, "epistemico": "fato", "origem": "wiki", "rel": "especializa"}, {"alvo": "sql", "confianca": 1.0, "epistemico": "fato", "origem": "wiki", "rel": "relaciona"}], "confianca": 1.0, "contraexemplos": [], "descricao": "Fazer o sistema executar entrada do atacante como código/consulta — a falha nº1 da web; a fronteira dado/código violada.", "epistemico": "fato", "exemplos": [], "id": "injecao", "memoria": "semantica", "meta": {"dominio": "programacao", "fonte": ""}, "origem": "wiki", "palavras_chave": [], "sinonimos": [], "tipo": "conceito", "titulo": "Injeção (SQL, comando)"}
\section{Injeção (SQL, comando)}
Fazer o sistema executar entrada do atacante como código/consulta — a falha nº1 da web; a fronteira dado/código violada.
\prov{relações: especializa vulnerabilidade\_exploit; relaciona sql}
\prov{proveniência: wiki \textperiodcentered{} fato \textperiodcentered{} confiança 1.0 \textperiodcentered{} domínio programacao \textperiodcentered{} história 3 versões no jornal}

%CRISTAL {"arestas": [{"alvo": "acoplamento_coesao", "confianca": 1.0, "epistemico": "fato", "origem": "wiki", "rel": "usa"}, {"alvo": "solid", "confianca": 1.0, "epistemico": "fato", "origem": "wiki", "rel": "usa"}], "confianca": 1.0, "contraexemplos": [], "descricao": "Um componente recebe suas dependências de fora em vez de criá-las — desacopla e testa; o coração do Angular e do NestJS.", "epistemico": "inferencia", "exemplos": [], "id": "injecao_de_dependencia", "memoria": "semantica", "meta": {"dominio": "programacao", "fonte": ""}, "origem": "wiki", "palavras_chave": [], "sinonimos": [], "tipo": "conceito", "titulo": "Injeção de Dependência"}
\section{Injeção de Dependência}
Um componente recebe suas dependências de fora em vez de criá-las — desacopla e testa; o coração do Angular e do NestJS.
\prov{relações: usa acoplamento\_coesao; usa solid}
\prov{proveniência: wiki \textperiodcentered{} inferencia \textperiodcentered{} confiança 1.0 \textperiodcentered{} domínio programacao \textperiodcentered{} história 6 versões no jornal}

%CRISTAL {"arestas": [{"alvo": "cristalchain", "confianca": 1.0, "epistemico": "fato", "origem": "wiki", "rel": "parte_de"}, {"alvo": "cristal_norma_conservada", "confianca": 1.0, "epistemico": "fato", "origem": "wiki", "rel": "eh_um"}, {"alvo": "mineracao_conservacao_verificavel", "confianca": 1.0, "epistemico": "fato", "origem": "wiki", "rel": "explica"}, {"alvo": "o_floco_e_a_conservacao_da_comunicacao", "confianca": 1.0, "epistemico": "fato", "origem": "wiki", "rel": "relaciona"}], "confianca": 1.0, "contraexemplos": [], "descricao": "Cada bloco aponta o hash do anterior (a aresta da rede cristalina). A cadeia é íntegra enquanto o encadeamento vale; adulterar um bit muda o hash e a invalida dali em diante — exatamente como a composição de Hurwitz N(xy)=N(x)N(y) se rompe fora das álgebras de divisão (surgem 'divisores de zero', blocos falsos). A imutabilidade É a conservação da norma sob o encadeamento. Verif.: adulterada inválida.", "epistemico": "fato", "exemplos": [], "id": "integridade_e_norma_conservada", "memoria": "semantica", "meta": {"dominio": "goldchain", "fonte": "cristalchain"}, "origem": "wiki", "palavras_chave": [], "sinonimos": [], "tipo": "conceito", "titulo": "A Integridade é a Norma Conservada"}
\section{A Integridade é a Norma Conservada}
Cada bloco aponta o hash do anterior (a aresta da rede cristalina). A cadeia é íntegra enquanto o encadeamento vale; adulterar um bit muda o hash e a invalida dali em diante — exatamente como a composição de Hurwitz N(xy)=N(x)N(y) se rompe fora das álgebras de divisão (surgem 'divisores de zero', blocos falsos). A imutabilidade É a conservação da norma sob o encadeamento. Verif.: adulterada inválida.
\prov{relações: parte\_de cristalchain; eh\_um cristal\_norma\_conservada; explica mineracao\_conservacao\_verificavel; relaciona o\_floco\_e\_a\_conservacao\_da\_comunicacao}
\prov{proveniência: wiki \textperiodcentered{} cristalchain \textperiodcentered{} fato \textperiodcentered{} confiança 1.0 \textperiodcentered{} domínio goldchain \textperiodcentered{} história 5 versões no jornal}

%CRISTAL {"arestas": [{"alvo": "mcculloch_pitts", "confianca": 1.0, "epistemico": "fato", "origem": "curriculo", "rel": "relaciona"}, {"alvo": "hopfield", "confianca": 1.0, "epistemico": "fato", "origem": "curriculo", "rel": "relaciona"}, {"alvo": "tiffany", "confianca": 1.0, "epistemico": "fato", "origem": "curriculo", "rel": "relaciona"}, {"alvo": "computacao", "confianca": 1.0, "epistemico": "fato", "origem": "curriculo", "rel": "relaciona"}], "confianca": 1.0, "contraexemplos": [], "descricao": "McCulloch-Pitts → Hopfield → Tiffany; IA como memória associativa do gato A.", "epistemico": "fato", "exemplos": [], "id": "inteligencia_artificial", "memoria": "semantica", "meta": {"dominio": "ia", "nivel": 5, "serve": "Núcleo IA", "verificacao": "slug idêntico — enriquecer, não duplicar"}, "origem": "curriculo", "palavras_chave": ["neurónio", "atrator", "memória"], "sinonimos": ["IA", "inteligência artificial"], "tipo": "conceito", "titulo": "inteligencia artificial"}
\section{inteligencia artificial}
McCulloch-Pitts \ensuremath{\to} Hopfield \ensuremath{\to} Tiffany; IA como memória associativa do gato A.
\prov{sinónimos: IA; inteligência artificial}
\prov{relações: relaciona mcculloch\_pitts; relaciona hopfield; relaciona tiffany; relaciona computacao}
\prov{proveniência: curriculo \textperiodcentered{} fato \textperiodcentered{} confiança 1.0 \textperiodcentered{} domínio ia \textperiodcentered{} história 66 versões no jornal}

%CRISTAL {"arestas": [{"alvo": "rede_de_computadores", "confianca": 1.0, "epistemico": "fato", "origem": "wiki", "rel": "especializa"}, {"alvo": "pilha_tcp_ip", "confianca": 1.0, "epistemico": "fato", "origem": "wiki", "rel": "usa"}], "confianca": 1.0, "contraexemplos": [], "descricao": "A rede das redes: bilhões de máquinas falando TCP/IP, endereçadas por IP e nomeadas por DNS. A infraestrutura do mundo digital.", "epistemico": "fato", "exemplos": [], "id": "internet", "memoria": "semantica", "meta": {"dominio": "programacao", "fonte": ""}, "origem": "wiki", "palavras_chave": [], "sinonimos": [], "tipo": "conceito", "titulo": "Internet"}
\section{Internet}
A rede das redes: bilhões de máquinas falando TCP/IP, endereçadas por IP e nomeadas por DNS. A infraestrutura do mundo digital.
\prov{relações: especializa rede\_de\_computadores; usa pilha\_tcp\_ip}
\prov{proveniência: wiki \textperiodcentered{} fato \textperiodcentered{} confiança 1.0 \textperiodcentered{} domínio programacao \textperiodcentered{} história 6 versões no jornal}

%CRISTAL {"arestas": [{"alvo": "design_vetorizado_svg", "confianca": 1.0, "epistemico": "fato", "origem": "wiki", "rel": "usa"}, {"alvo": "design_elementos_do_chatgpt", "confianca": 1.0, "epistemico": "fato", "origem": "wiki", "rel": "relaciona"}], "confianca": 1.0, "contraexemplos": [], "descricao": "Com as órbitas vetorizadas, a máquina INTERPOLA: contorno_gema(lados) reamostra o rondiz a um n comum, e interpolar(ca, cb, t) = (1−t)·A + t·B é o caminho contínuo entre duas órbitas. gerar_interpolacao gera a tira de uma gema de 6 facetas a 12 (a estrela abrindo) em SVG — o morphing como uma órbita percorrida. O ponto médio (t=0.5) fica exatamente entre os dois contornos. É o BAI: a órbita entre dois estados.", "epistemico": "fato", "exemplos": [], "id": "interpolacao_orbitas", "memoria": "semantica", "meta": {"dominio": "sistema", "fonte": "orbitas vetoriais"}, "origem": "wiki", "palavras_chave": [], "sinonimos": [], "tipo": "conceito", "titulo": "A Interpolação das Órbitas (o morphing da gema)"}
\section{A Interpolação das Órbitas (o morphing da gema)}
Com as órbitas vetorizadas, a máquina INTERPOLA: contorno\_gema(lados) reamostra o rondiz a um n comum, e interpolar(ca, cb, t) = (1\ensuremath{-}t)\textperiodcentered A + t\textperiodcentered B é o caminho contínuo entre duas órbitas. gerar\_interpolacao gera a tira de uma gema de 6 facetas a 12 (a estrela abrindo) em SVG — o morphing como uma órbita percorrida. O ponto médio (t=0.5) fica exatamente entre os dois contornos. É o BAI: a órbita entre dois estados.
\prov{relações: usa design\_vetorizado\_svg; relaciona design\_elementos\_do\_chatgpt}
\prov{proveniência: wiki \textperiodcentered{} orbitas vetoriais \textperiodcentered{} fato \textperiodcentered{} confiança 1.0 \textperiodcentered{} domínio sistema \textperiodcentered{} história 3 versões no jornal}

%CRISTAL {"arestas": [{"alvo": "linguagem_de_programacao", "confianca": 1.0, "epistemico": "fato", "origem": "wiki", "rel": "parte_de"}, {"alvo": "linguagem_matriz_gatos", "confianca": 1.0, "epistemico": "fato", "origem": "wiki", "rel": "explica"}, {"alvo": "interpretador_nucleo_gatos", "confianca": 1.0, "epistemico": "fato", "origem": "wiki", "rel": "relaciona"}], "confianca": 1.0, "contraexemplos": [], "descricao": "Executa o código-fonte diretamente, comando a comando — sem gerar binário; o loop fetch-decode-execute em software.", "epistemico": "fato", "exemplos": [], "id": "interpretador", "memoria": "semantica", "meta": {"dominio": "programacao", "fonte": ""}, "origem": "wiki", "palavras_chave": [], "sinonimos": [], "tipo": "conceito", "titulo": "Interpretador"}
\section{Interpretador}
Executa o código-fonte diretamente, comando a comando — sem gerar binário; o loop fetch-decode-execute em software.
\prov{relações: parte\_de linguagem\_de\_programacao; explica linguagem\_matriz\_gatos; relaciona interpretador\_nucleo\_gatos}
\prov{proveniência: wiki \textperiodcentered{} fato \textperiodcentered{} confiança 1.0 \textperiodcentered{} domínio programacao \textperiodcentered{} história 7 versões no jornal}

%CRISTAL {"arestas": [{"alvo": "gato_generaliza_matriz_companion", "confianca": 1.0, "epistemico": "fato", "origem": "wiki", "rel": "usa"}, {"alvo": "contracao_einstein_k_tensorial", "confianca": 1.0, "epistemico": "fato", "origem": "wiki", "rel": "usa"}, {"alvo": "transformada_metalica", "confianca": 1.0, "epistemico": "fato", "origem": "wiki", "rel": "usa"}, {"alvo": "ula_por_componente", "confianca": 1.0, "epistemico": "fato", "origem": "wiki", "rel": "relaciona"}, {"alvo": "linguagem_gatos_sub_turing_honesto", "confianca": 1.0, "epistemico": "fato", "origem": "wiki", "rel": "relaciona"}, {"alvo": "tiffany", "confianca": 1.0, "epistemico": "fato", "origem": "wiki", "rel": "parte_de"}], "confianca": 1.0, "contraexemplos": [], "descricao": "linguagem/interpretador.py: o núcleo RODÁVEL da linguagem. Um programa é uma sequência de instruções sobre CÉLULAS (words 128 bits): a PORTA central é o GATO (cifra Aₙ, o par (a,b)→(b,n·b+a)); ADD/SUB são a ULA; SET/MOV movem; JMP/JZ/JNZ o fluxo (offsets fixos — sub-Turing). Laço fetch→decode→exec até HALT (como o cpustep do metal), com um assembler mínimo (a semente da sintaxe: rótulos + 'op args'). A TESE da linguagem, VERIFICADA no código: (1) um BLOCO RETO de instruções lineares é UMA matriz de gatos — o produto (contração de Einstein) das matrizes-instrução (afins (N+1)×(N+1)); rodar passo-a-passo ≡ aplicar a matriz composta; (2) o LAÇO do gato (GATO n × k) É a POSIÇÃO Aₙk — o 'for' fechado da transformada, O(log k) por potência rápida. Demo verificada: Pell por laço (9×)=2378 ≡ posição A₂⁹=2378 (confere com metalpos). [F] o interpretador roda e a equivalência composição=contração=posição está provada em código.", "epistemico": "fato", "exemplos": [], "id": "interpretador_nucleo_gatos", "memoria": "semantica", "meta": {"autor": "Lopes/broca-so", "dominio": "computacao", "fonte": "linguagem/interpretador.py (montar, rodar, matriz_instrucao, compilar_bloco); verificado: Pell laço≡posição=2378"}, "origem": "wiki", "palavras_chave": [], "sinonimos": [], "tipo": "conceito", "titulo": "O núcleo do interpretador: fetch→exec; bloco reto = matriz de gatos; laço = posição [F]"}
\section{O núcleo do interpretador: fetch\ensuremath{\to}exec; bloco reto = matriz de gatos; laço = posição [F]}
linguagem/interpretador.py: o núcleo RODÁVEL da linguagem. Um programa é uma sequência de instruções sobre CÉLULAS (words 128 bits): a PORTA central é o GATO (cifra A\ensuremath{_{n}}, o par (a,b)\ensuremath{\to}(b,n\textperiodcentered b+a)); ADD/SUB são a ULA; SET/MOV movem; JMP/JZ/JNZ o fluxo (offsets fixos — sub-Turing). Laço fetch\ensuremath{\to}decode\ensuremath{\to}exec até HALT (como o cpustep do metal), com um assembler mínimo (a semente da sintaxe: rótulos + 'op args'). A TESE da linguagem, VERIFICADA no código: (1) um BLOCO RETO de instruções lineares é UMA matriz de gatos — o produto (contração de Einstein) das matrizes-instrução (afins (N+1)$\times$(N+1)); rodar passo-a-passo \ensuremath{\equiv} aplicar a matriz composta; (2) o LAÇO do gato (GATO n $\times$ k) É a POSIÇÃO A\ensuremath{_{n}}k — o 'for' fechado da transformada, O(log k) por potência rápida. Demo verificada: Pell por laço (9$\times$)=2378 \ensuremath{\equiv} posição A\ensuremath{_{2}}\ensuremath{^{9}}=2378 (confere com metalpos). [F] o interpretador roda e a equivalência composição=contração=posição está provada em código.
\prov{relações: usa gato\_generaliza\_matriz\_companion; usa contracao\_einstein\_k\_tensorial; usa transformada\_metalica; relaciona ula\_por\_componente; relaciona linguagem\_gatos\_sub\_turing\_honesto; parte\_de tiffany}
\prov{proveniência: wiki \textperiodcentered{} linguagem/interpretador.py (montar, rodar, matriz\_instrucao, compilar\_bloco); verificado: Pell laço\ensuremath{\equiv}posição=2378 \textperiodcentered{} fato \textperiodcentered{} confiança 1.0 \textperiodcentered{} domínio computacao \textperiodcentered{} história 53 versões no jornal}

%CRISTAL {"arestas": [{"alvo": "driver", "confianca": 1.0, "epistemico": "fato", "origem": "curriculo", "rel": "parte_de"}, {"alvo": "kernel", "confianca": 0.711, "epistemico": "fato", "origem": "manual", "rel": "referencia"}, {"alvo": "broca_os", "confianca": 1.0, "epistemico": "fato", "origem": "curriculo", "rel": "explica"}, {"alvo": "utilitarismo", "confianca": 0.5, "epistemico": "fato", "origem": "wiki", "rel": "relaciona"}, {"alvo": "o_metodo_em_escala_conhecida_o_sistema_solar", "confianca": 0.5, "epistemico": "fato", "origem": "wiki", "rel": "relaciona"}, {"alvo": "melhor_candidato", "confianca": 0.5, "epistemico": "fato", "origem": "wiki", "rel": "relaciona"}, {"alvo": "o_que_e_no_projecto", "confianca": 0.5, "epistemico": "fato", "origem": "wiki", "rel": "relaciona"}, {"alvo": "testes", "confianca": 0.5, "epistemico": "fato", "origem": "wiki", "rel": "relaciona"}], "confianca": 1.0, "contraexemplos": [], "descricao": "Comunicação entre processos — pipes, sockets, shm.", "epistemico": "fato", "exemplos": [], "id": "ipc", "memoria": "semantica", "meta": {"dominio": "sistemas", "nivel": 6, "ordem": 2}, "origem": "curriculo", "palavras_chave": [], "sinonimos": [], "tipo": "conceito", "titulo": "IPC"}
\section{IPC}
Comunicação entre processos — pipes, sockets, shm.
\prov{relações: parte\_de driver; referencia kernel; explica broca\_os; relaciona utilitarismo; relaciona o\_metodo\_em\_escala\_conhecida\_o\_sistema\_solar; relaciona melhor\_candidato; relaciona o\_que\_e\_no\_projecto; relaciona testes}
\prov{proveniência: curriculo \textperiodcentered{} fato \textperiodcentered{} confiança 1.0 \textperiodcentered{} domínio sistemas \textperiodcentered{} história 42 versões no jornal}

%CRISTAL {"arestas": [{"alvo": "operacao_broca_so", "confianca": 1.0, "epistemico": "fato", "origem": "wiki", "rel": "parte_de"}, {"alvo": "arquitetura_von_neumann", "confianca": 1.0, "epistemico": "fato", "origem": "wiki", "rel": "usa"}], "confianca": 1.0, "contraexemplos": [], "descricao": "O conjunto de instruções é fechado — HALT, LOAD, STORE, ADD, SUB, AND, OR, XOR, CMP, JMP, JZ, JNZ — provado suficiente para 7 programas de referência; nenhuma instrução nova entra se for composição das existentes ('zero opcodes novos').", "epistemico": "fato", "exemplos": [], "id": "isa_da_broca", "memoria": "semantica", "meta": {"autor": "broca-so", "dominio": "operacao_so", "fonte": "ula/PROVA_ISA.md; ula/instrucoes.c"}, "origem": "wiki", "palavras_chave": [], "sinonimos": [], "tipo": "conceito", "titulo": "A ISA Mínima de 12 Opcodes"}
\section{A ISA Mínima de 12 Opcodes}
O conjunto de instruções é fechado — HALT, LOAD, STORE, ADD, SUB, AND, OR, XOR, CMP, JMP, JZ, JNZ — provado suficiente para 7 programas de referência; nenhuma instrução nova entra se for composição das existentes ('zero opcodes novos').
\prov{relações: parte\_de operacao\_broca\_so; usa arquitetura\_von\_neumann}
\prov{proveniência: wiki \textperiodcentered{} ula/PROVA\_ISA.md; ula/instrucoes.c \textperiodcentered{} fato \textperiodcentered{} confiança 1.0 \textperiodcentered{} domínio operacao\_so \textperiodcentered{} história 3 versões no jornal}

%CRISTAL {"arestas": [{"alvo": "tiffany_home", "confianca": 1.0, "epistemico": "fato", "origem": "wiki", "rel": "usa"}, {"alvo": "sistema_e_lei_nao_dono", "confianca": 1.0, "epistemico": "fato", "origem": "wiki", "rel": "explica"}, {"alvo": "soberania_da_chave", "confianca": 1.0, "epistemico": "fato", "origem": "wiki", "rel": "relaciona"}], "confianca": 1.0, "contraexemplos": [], "descricao": "A validação garante que domínios cliente vivam sob cliente/ e Core sob base/, que todo meta.json de cliente tenha licenca=proprietary (o Core nunca), que nenhum chunk_id colida entre as árvores, e que nenhuma linha memoria/ vaze para o Core. A separação é lei imposta por código, não por dono.", "epistemico": "fato", "exemplos": [], "id": "isolamento_base_cliente", "memoria": "semantica", "meta": {"autor": "broca-so", "dominio": "operacao_so", "fonte": "code/tiffany_platform/isolation.py; domains.py"}, "origem": "wiki", "palavras_chave": [], "sinonimos": [], "tipo": "conceito", "titulo": "Isolamento Base ↔ Cliente (multi-tenancy)"}
\section{Isolamento Base \ensuremath{\leftrightarrow} Cliente (multi-tenancy)}
A validação garante que domínios cliente vivam sob cliente/ e Core sob base/, que todo meta.json de cliente tenha licenca=proprietary (o Core nunca), que nenhum chunk\_id colida entre as árvores, e que nenhuma linha memoria/ vaze para o Core. A separação é lei imposta por código, não por dono.
\prov{relações: usa tiffany\_home; explica sistema\_e\_lei\_nao\_dono; relaciona soberania\_da\_chave}
\prov{proveniência: wiki \textperiodcentered{} code/tiffany\_platform/isolation.py; domains.py \textperiodcentered{} fato \textperiodcentered{} confiança 1.0 \textperiodcentered{} domínio operacao\_so \textperiodcentered{} história 4 versões no jornal}

%CRISTAL {"arestas": [{"alvo": "kernel", "confianca": 0.766, "epistemico": "fato", "origem": "manual", "rel": "referencia"}, {"alvo": "compilador", "confianca": 1.0, "epistemico": "fato", "origem": "curriculo", "rel": "referencia"}, {"alvo": "prova_isa", "confianca": 0.5, "epistemico": "fato", "origem": "wiki", "rel": "relaciona"}, {"alvo": "numa", "confianca": 0.5, "epistemico": "fato", "origem": "wiki", "rel": "relaciona"}, {"alvo": "koch_pell", "confianca": 0.5, "epistemico": "fato", "origem": "wiki", "rel": "relaciona"}, {"alvo": "transcricao", "confianca": 0.5, "epistemico": "fato", "origem": "wiki", "rel": "relaciona"}, {"alvo": "main", "confianca": 0.5, "epistemico": "fato", "origem": "wiki", "rel": "relaciona"}, {"alvo": "koch_atlas", "confianca": 0.5, "epistemico": "fato", "origem": "wiki", "rel": "relaciona"}, {"alvo": "mamifero", "confianca": 0.5, "epistemico": "fato", "origem": "wiki", "rel": "relaciona"}, {"alvo": "montagem", "confianca": 0.5, "epistemico": "fato", "origem": "wiki", "rel": "relaciona"}, {"alvo": "maquina_virtual", "confianca": 1.0, "epistemico": "fato", "origem": "wiki", "rel": "usa"}, {"alvo": "programacao_orientada_objetos", "confianca": 1.0, "epistemico": "fato", "origem": "wiki", "rel": "usa"}], "confianca": 1.0, "contraexemplos": [], "descricao": "Gosling (Sun): 'escreva uma vez, rode em qualquer lugar' via JVM e bytecode; OOP, GC, o pilar corporativo.", "epistemico": "fato", "exemplos": [], "id": "java", "memoria": "semantica", "meta": {"autor": "Gosling", "dominio": "programacao", "fonte": ""}, "origem": "wiki", "palavras_chave": [], "sinonimos": [], "tipo": "conceito", "titulo": "Java"}
\section{Java}
Gosling (Sun): 'escreva uma vez, rode em qualquer lugar' via JVM e bytecode; OOP, GC, o pilar corporativo.
\prov{relações: referencia kernel; referencia compilador; relaciona prova\_isa; relaciona numa; relaciona koch\_pell; relaciona transcricao; relaciona main; relaciona koch\_atlas; relaciona mamifero; relaciona montagem; usa maquina\_virtual; usa programacao\_orientada\_objetos}
\prov{proveniência: wiki \textperiodcentered{} fato \textperiodcentered{} confiança 1.0 \textperiodcentered{} domínio programacao \textperiodcentered{} história 21 versões no jornal}

%CRISTAL {"arestas": [{"alvo": "linguagem_de_programacao", "confianca": 1.0, "epistemico": "fato", "origem": "wiki", "rel": "parte_de"}], "confianca": 1.0, "contraexemplos": [], "descricao": "Eich (em 10 dias): a linguagem da web, no navegador e no servidor (Node); dinâmica, assíncrona, onipresente.", "epistemico": "fato", "exemplos": [], "id": "javascript", "memoria": "semantica", "meta": {"autor": "Eich", "dominio": "programacao", "fonte": ""}, "origem": "wiki", "palavras_chave": [], "sinonimos": [], "tipo": "conceito", "titulo": "JavaScript"}
\section{JavaScript}
Eich (em 10 dias): a linguagem da web, no navegador e no servidor (Node); dinâmica, assíncrona, onipresente.
\prov{relações: parte\_de linguagem\_de\_programacao}
\prov{proveniência: wiki \textperiodcentered{} fato \textperiodcentered{} confiança 1.0 \textperiodcentered{} domínio programacao \textperiodcentered{} história 4 versões no jornal}

%CRISTAL {"arestas": [{"alvo": "minimax", "confianca": 1.0, "epistemico": "fato", "origem": "wiki", "rel": "usa"}], "confianca": 1.0, "contraexemplos": [], "descricao": "Jogos cujo resultado com jogo perfeito é conhecido: jogo-da-velha (empate), Connect Four (1ª vitória), damas (empate). 'Resolver' pode ser fraco (do início) ou forte (de qualquer posição) — o fim da árvore.", "epistemico": "inferencia", "exemplos": [], "id": "jogos_resolvidos", "memoria": "semantica", "meta": {"dominio": "ia", "fonte": "jogos e IA"}, "origem": "wiki", "palavras_chave": [], "sinonimos": [], "tipo": "conceito", "titulo": "Jogos Resolvidos"}
\section{Jogos Resolvidos}
Jogos cujo resultado com jogo perfeito é conhecido: jogo-da-velha (empate), Connect Four (1ª vitória), damas (empate). 'Resolver' pode ser fraco (do início) ou forte (de qualquer posição) — o fim da árvore.
\prov{relações: usa minimax}
\prov{proveniência: wiki \textperiodcentered{} jogos e IA \textperiodcentered{} inferencia \textperiodcentered{} confiança 1.0 \textperiodcentered{} domínio ia \textperiodcentered{} história 2 versões no jornal}

%CRISTAL {"arestas": [{"alvo": "javascript", "confianca": 1.0, "epistemico": "fato", "origem": "wiki", "rel": "relaciona"}, {"alvo": "api", "confianca": 1.0, "epistemico": "fato", "origem": "wiki", "rel": "usa"}, {"alvo": "rest", "confianca": 1.0, "epistemico": "fato", "origem": "wiki", "rel": "usa"}, {"alvo": "arquivo_de_configuracao", "confianca": 1.0, "epistemico": "fato", "origem": "wiki", "rel": "usa"}], "confianca": 1.0, "contraexemplos": [], "descricao": "JavaScript Object Notation: formato leve de dados (objetos e listas) legível por humano e máquina — a língua franca das APIs web.", "epistemico": "fato", "exemplos": [], "id": "json", "memoria": "semantica", "meta": {"autor": "Crockford", "dominio": "programacao", "fonte": ""}, "origem": "wiki", "palavras_chave": [], "sinonimos": [], "tipo": "conceito", "titulo": "JSON"}
\section{JSON}
JavaScript Object Notation: formato leve de dados (objetos e listas) legível por humano e máquina — a língua franca das APIs web.
\prov{relações: relaciona javascript; usa api; usa rest; usa arquivo\_de\_configuracao}
\prov{proveniência: wiki \textperiodcentered{} fato \textperiodcentered{} confiança 1.0 \textperiodcentered{} domínio programacao \textperiodcentered{} história 5 versões no jornal}

%CRISTAL {"arestas": [{"alvo": "javascript", "confianca": 1.0, "epistemico": "fato", "origem": "wiki", "rel": "usa"}, {"alvo": "html", "confianca": 1.0, "epistemico": "fato", "origem": "wiki", "rel": "relaciona"}], "confianca": 1.0, "contraexemplos": [], "descricao": "Sintaxe que mistura marcação HTML dentro do JavaScript — o componente descreve sua própria árvore. Compilado para chamadas de função.", "epistemico": "fato", "exemplos": [], "id": "jsx", "memoria": "semantica", "meta": {"dominio": "programacao", "fonte": ""}, "origem": "wiki", "palavras_chave": [], "sinonimos": [], "tipo": "conceito", "titulo": "JSX"}
\section{JSX}
Sintaxe que mistura marcação HTML dentro do JavaScript — o componente descreve sua própria árvore. Compilado para chamadas de função.
\prov{relações: usa javascript; relaciona html}
\prov{proveniência: wiki \textperiodcentered{} fato \textperiodcentered{} confiança 1.0 \textperiodcentered{} domínio programacao \textperiodcentered{} história 6 versões no jornal}

%CRISTAL {"arestas": [{"alvo": "linguagem_de_programacao", "confianca": 1.0, "epistemico": "fato", "origem": "wiki", "rel": "parte_de"}, {"alvo": "compilacao_jit", "confianca": 1.0, "epistemico": "fato", "origem": "wiki", "rel": "usa"}], "confianca": 1.0, "contraexemplos": [], "descricao": "Desenhada para computação científica: a expressividade do Python com a velocidade do C, via JIT e despacho múltiplo.", "epistemico": "fato", "exemplos": [], "id": "julia", "memoria": "semantica", "meta": {"dominio": "programacao", "fonte": ""}, "origem": "wiki", "palavras_chave": [], "sinonimos": [], "tipo": "conceito", "titulo": "Julia"}
\section{Julia}
Desenhada para computação científica: a expressividade do Python com a velocidade do C, via JIT e despacho múltiplo.
\prov{relações: parte\_de linguagem\_de\_programacao; usa compilacao\_jit}
\prov{proveniência: wiki \textperiodcentered{} fato \textperiodcentered{} confiança 1.0 \textperiodcentered{} domínio programacao \textperiodcentered{} história 6 versões no jornal}

%CRISTAL {"arestas": [{"alvo": "broca_os", "confianca": 1.0, "epistemico": "fato", "origem": "curriculo", "rel": "parte_de"}, {"alvo": "rust", "confianca": 0.5, "epistemico": "fato", "origem": "wiki", "rel": "relaciona"}, {"alvo": "main", "confianca": 0.5, "epistemico": "fato", "origem": "wiki", "rel": "relaciona"}, {"alvo": "recursao", "confianca": 0.5, "epistemico": "fato", "origem": "wiki", "rel": "relaciona"}, {"alvo": "mamifero", "confianca": 0.5, "epistemico": "fato", "origem": "wiki", "rel": "relaciona"}, {"alvo": "thread", "confianca": 0.5, "epistemico": "fato", "origem": "wiki", "rel": "relaciona"}, {"alvo": "testes", "confianca": 0.5, "epistemico": "fato", "origem": "wiki", "rel": "relaciona"}, {"alvo": "sessao_interativa", "confianca": 0.5, "epistemico": "fato", "origem": "wiki", "rel": "relaciona"}, {"alvo": "sistema_operacional", "confianca": 1.0, "epistemico": "fato", "origem": "wiki", "rel": "parte_de"}], "confianca": 1.0, "contraexemplos": [], "descricao": "O núcleo do SO, que roda em modo privilegiado: agenda processos, gere memória e media o acesso ao hardware. O guardião da máquina.", "epistemico": "fato", "exemplos": [], "id": "kernel", "memoria": "semantica", "meta": {"dominio": "programacao", "fonte": ""}, "origem": "wiki", "palavras_chave": [], "sinonimos": [], "tipo": "conceito", "titulo": "Kernel"}
\section{Kernel}
O núcleo do SO, que roda em modo privilegiado: agenda processos, gere memória e media o acesso ao hardware. O guardião da máquina.
\prov{relações: parte\_de broca\_os; relaciona rust; relaciona main; relaciona recursao; relaciona mamifero; relaciona thread; relaciona testes; relaciona sessao\_interativa; parte\_de sistema\_operacional}
\prov{proveniência: wiki \textperiodcentered{} fato \textperiodcentered{} confiança 1.0 \textperiodcentered{} domínio programacao \textperiodcentered{} história 41 versões no jornal}

%CRISTAL {"arestas": [{"alvo": "express", "confianca": 1.0, "epistemico": "fato", "origem": "wiki", "rel": "relaciona"}, {"alvo": "nodejs", "confianca": 1.0, "epistemico": "fato", "origem": "wiki", "rel": "usa"}], "confianca": 1.0, "contraexemplos": [], "descricao": "Do time do Express: um núcleo minúsculo com middleware baseado em async/await — o Express repensado, moderno e sem callbacks.", "epistemico": "fato", "exemplos": [], "id": "koa", "memoria": "semantica", "meta": {"dominio": "programacao", "fonte": ""}, "origem": "wiki", "palavras_chave": [], "sinonimos": [], "tipo": "conceito", "titulo": "Koa"}
\section{Koa}
Do time do Express: um núcleo minúsculo com middleware baseado em async/await — o Express repensado, moderno e sem callbacks.
\prov{relações: relaciona express; usa nodejs}
\prov{proveniência: wiki \textperiodcentered{} fato \textperiodcentered{} confiança 1.0 \textperiodcentered{} domínio programacao \textperiodcentered{} história 6 versões no jornal}

%CRISTAL {"arestas": [{"alvo": "grafo_append_only", "confianca": 1.0, "epistemico": "fato", "origem": "wiki", "rel": "usa"}, {"alvo": "hopfield", "confianca": 1.0, "epistemico": "fato", "origem": "wiki", "rel": "usa"}, {"alvo": "colapso_koch_tiffany", "confianca": 1.0, "epistemico": "fato", "origem": "wiki", "rel": "usa"}], "confianca": 1.0, "contraexemplos": [], "descricao": "O grafo é projetado numa linha Koch por nó (_sync_tsv_linha) e reindexado via prosa.indexa; montar_tecido/ingest_batch reconstroem o tecido inteiro após lotes. Três camadas: .jsonl (verdade) → .tsv (linha Koch) → .idx (impressões de 256 gatos).", "epistemico": "fato", "exemplos": [], "id": "koch_sync_reindex", "memoria": "semantica", "meta": {"autor": "broca-so", "dominio": "operacao_so", "fonte": "conversa/conhecimento/store.py:reindexar"}, "origem": "wiki", "palavras_chave": [], "sinonimos": [], "tipo": "conceito", "titulo": "Sincronização grafo→tsv→idx"}
\section{Sincronização grafo\ensuremath{\to}tsv\ensuremath{\to}idx}
O grafo é projetado numa linha Koch por nó (\_sync\_tsv\_linha) e reindexado via prosa.indexa; montar\_tecido/ingest\_batch reconstroem o tecido inteiro após lotes. Três camadas: .jsonl (verdade) \ensuremath{\to} .tsv (linha Koch) \ensuremath{\to} .idx (impressões de 256 gatos).
\prov{relações: usa grafo\_append\_only; usa hopfield; usa colapso\_koch\_tiffany}
\prov{proveniência: wiki \textperiodcentered{} conversa/conhecimento/store.py:reindexar \textperiodcentered{} fato \textperiodcentered{} confiança 1.0 \textperiodcentered{} domínio operacao\_so \textperiodcentered{} história 4 versões no jornal}

%CRISTAL {"arestas": [{"alvo": "turing_halting_church", "confianca": 1.0, "epistemico": "fato", "origem": "wiki", "rel": "usa"}, {"alvo": "godel_incompletude", "confianca": 1.0, "epistemico": "fato", "origem": "wiki", "rel": "relaciona"}, {"alvo": "entropia_shannon", "confianca": 1.0, "epistemico": "fato", "origem": "wiki", "rel": "relaciona"}], "confianca": 1.0, "contraexemplos": [], "descricao": "K(x)=min|p|:U(p)=x = comprimento do menor programa que produz x numa máquina universal; formaliza 'aleatoriedade'/incompressibilidade sem probabilidades (x é aleatório se K(x)≥|x|). Solomonoff/Kolmogorov/Chaitin. FATO chave: K é NÃO-COMPUTÁVEL (nenhum programa calcula K(x) para toda string) — Chaitin liga isso à incompletude de Gödel. Complementa a entropia de Shannon (probabilística, computável).", "epistemico": "fato", "exemplos": [], "id": "kolmogorov_complexidade", "memoria": "semantica", "meta": {"autor": "broca-so", "dominio": "computacao", "fonte": "Solomonoff 1964; Kolmogorov 1965; Chaitin 1969 J.ACM"}, "origem": "wiki", "palavras_chave": [], "sinonimos": [], "tipo": "conceito", "titulo": "Complexidade de Kolmogorov (informação algorítmica)"}
\section{Complexidade de Kolmogorov (informação algorítmica)}
K(x)=min|p|:U(p)=x = comprimento do menor programa que produz x numa máquina universal; formaliza 'aleatoriedade'/incompressibilidade sem probabilidades (x é aleatório se K(x)\ensuremath{\ge}|x|). Solomonoff/Kolmogorov/Chaitin. FATO chave: K é NÃO-COMPUTÁVEL (nenhum programa calcula K(x) para toda string) — Chaitin liga isso à incompletude de Gödel. Complementa a entropia de Shannon (probabilística, computável).
\prov{relações: usa turing\_halting\_church; relaciona godel\_incompletude; relaciona entropia\_shannon}
\prov{proveniência: wiki \textperiodcentered{} Solomonoff 1964; Kolmogorov 1965; Chaitin 1969 J.ACM \textperiodcentered{} fato \textperiodcentered{} confiança 1.0 \textperiodcentered{} domínio computacao \textperiodcentered{} história 5 versões no jornal}

%CRISTAL {"arestas": [{"alvo": "java", "confianca": 1.0, "epistemico": "fato", "origem": "wiki", "rel": "relaciona"}, {"alvo": "maquina_virtual", "confianca": 1.0, "epistemico": "fato", "origem": "wiki", "rel": "usa"}], "confianca": 1.0, "contraexemplos": [], "descricao": "JetBrains: um Java mais seguro e conciso na JVM (null-safety, corrotinas) — a linguagem oficial do Android.", "epistemico": "fato", "exemplos": [], "id": "kotlin", "memoria": "semantica", "meta": {"autor": "JetBrains", "dominio": "programacao", "fonte": ""}, "origem": "wiki", "palavras_chave": [], "sinonimos": [], "tipo": "conceito", "titulo": "Kotlin"}
\section{Kotlin}
JetBrains: um Java mais seguro e conciso na JVM (null-safety, corrotinas) — a linguagem oficial do Android.
\prov{relações: relaciona java; usa maquina\_virtual}
\prov{proveniência: wiki \textperiodcentered{} fato \textperiodcentered{} confiança 1.0 \textperiodcentered{} domínio programacao \textperiodcentered{} história 6 versões no jornal}

%CRISTAL {"arestas": [{"alvo": "container", "confianca": 1.0, "epistemico": "fato", "origem": "wiki", "rel": "usa"}, {"alvo": "computacao_em_nuvem", "confianca": 1.0, "epistemico": "fato", "origem": "wiki", "rel": "usa"}], "confianca": 1.0, "contraexemplos": [], "descricao": "Google: o orquestrador que agenda, escala e cura contêineres num cluster — declara-se o estado desejado, ele o mantém.", "epistemico": "fato", "exemplos": [], "id": "kubernetes", "memoria": "semantica", "meta": {"autor": "Google", "dominio": "programacao", "fonte": ""}, "origem": "wiki", "palavras_chave": [], "sinonimos": [], "tipo": "conceito", "titulo": "Kubernetes"}
\section{Kubernetes}
Google: o orquestrador que agenda, escala e cura contêineres num cluster — declara-se o estado desejado, ele o mantém.
\prov{relações: usa container; usa computacao\_em\_nuvem}
\prov{proveniência: wiki \textperiodcentered{} fato \textperiodcentered{} confiança 1.0 \textperiodcentered{} domínio programacao \textperiodcentered{} história 6 versões no jornal}

%CRISTAL {"arestas": [{"alvo": "concorrencia_paralelismo", "confianca": 1.0, "epistemico": "fato", "origem": "wiki", "rel": "especializa"}, {"alvo": "modelo_de_atores", "confianca": 1.0, "epistemico": "fato", "origem": "wiki", "rel": "relaciona"}], "confianca": 1.0, "contraexemplos": [], "descricao": "O modelo de concorrência do Node/navegador: uma fila de tarefas processada por uma thread, sem bloquear em I/O — concorrência sem paralelismo.", "epistemico": "fato", "exemplos": [], "id": "laco_de_eventos", "memoria": "semantica", "meta": {"dominio": "programacao", "fonte": ""}, "origem": "wiki", "palavras_chave": [], "sinonimos": [], "tipo": "conceito", "titulo": "Laço de Eventos (Event Loop)"}
\section{Laço de Eventos (Event Loop)}
O modelo de concorrência do Node/navegador: uma fila de tarefas processada por uma thread, sem bloquear em I/O — concorrência sem paralelismo.
\prov{relações: especializa concorrencia\_paralelismo; relaciona modelo\_de\_atores}
\prov{proveniência: wiki \textperiodcentered{} fato \textperiodcentered{} confiança 1.0 \textperiodcentered{} domínio programacao \textperiodcentered{} história 6 versões no jornal}

%CRISTAL {"arestas": [{"alvo": "traducao_como_transformacao", "confianca": 1.0, "epistemico": "fato", "origem": "wiki", "rel": "usa"}, {"alvo": "traducao_como_interpretante", "confianca": 1.0, "epistemico": "fato", "origem": "wiki", "rel": "usa"}, {"alvo": "transformada_metalica", "confianca": 1.0, "epistemico": "fato", "origem": "wiki", "rel": "usa"}, {"alvo": "bai_orbita_reta_metalica", "confianca": 1.0, "epistemico": "fato", "origem": "wiki", "rel": "usa"}, {"alvo": "peirce_dsl", "confianca": 1.0, "epistemico": "fato", "origem": "wiki", "rel": "usa"}, {"alvo": "colapso", "confianca": 1.0, "epistemico": "fato", "origem": "wiki", "rel": "explica"}, {"alvo": "meta_inducao_geral_simbolo", "confianca": 1.0, "epistemico": "fato", "origem": "wiki", "rel": "especializa"}, {"alvo": "conversao_como_traducao_formal", "confianca": 1.0, "epistemico": "fato", "origem": "wiki", "rel": "relaciona"}, {"alvo": "torre_unificada", "confianca": 1.0, "epistemico": "fato", "origem": "wiki", "rel": "parte_de"}], "confianca": 1.0, "contraexemplos": [], "descricao": "O anel do broca-so num operador só: TRADUZIR (semiose/interpretante) = TRANSFORMAR (a transformada metálica, β-numeração base σ_n) = ORBITAR (a reta do BAI, δ=σ_n·4) = SUBSTITUIR (o produto ':' de Peirce, o Peirce-DSL) = COLAPSAR (o Ψ da Tiffany no tecido). São a MESMA operação sob símbolos diferentes — a meta-indução geral (só o símbolo muda). O laço BAI↔transformada↔Peirce↔tradução↔colapso é um ciclo fechado; a tradução da matemática (math↔LaTeX) é só o caso onde esse operador conserva o invariante perfeitamente.", "epistemico": "inferencia", "exemplos": [], "id": "laco_traducao_bai_peirce", "memoria": "semantica", "meta": {"dominio": "programacao", "fonte": "fecha o arco de línguas no núcleo simbólico"}, "origem": "wiki", "palavras_chave": [], "sinonimos": [], "tipo": "conceito", "titulo": "O laço fecha: traduzir = transformar = orbitar = substituir = colapsar"}
\section{O laço fecha: traduzir = transformar = orbitar = substituir = colapsar}
O anel do broca-so num operador só: TRADUZIR (semiose/interpretante) = TRANSFORMAR (a transformada metálica, \ensuremath{\beta}-numeração base \ensuremath{\sigma}\_n) = ORBITAR (a reta do BAI, \ensuremath{\delta}=\ensuremath{\sigma}\_n\textperiodcentered 4) = SUBSTITUIR (o produto ':' de Peirce, o Peirce-DSL) = COLAPSAR (o \ensuremath{\Psi} da Tiffany no tecido). São a MESMA operação sob símbolos diferentes — a meta-indução geral (só o símbolo muda). O laço BAI\ensuremath{\leftrightarrow}transformada\ensuremath{\leftrightarrow}Peirce\ensuremath{\leftrightarrow}tradução\ensuremath{\leftrightarrow}colapso é um ciclo fechado; a tradução da matemática (math\ensuremath{\leftrightarrow}LaTeX) é só o caso onde esse operador conserva o invariante perfeitamente.
\prov{relações: usa traducao\_como\_transformacao; usa traducao\_como\_interpretante; usa transformada\_metalica; usa bai\_orbita\_reta\_metalica; usa peirce\_dsl; explica colapso; especializa meta\_inducao\_geral\_simbolo; relaciona conversao\_como\_traducao\_formal; parte\_de torre\_unificada}
\prov{proveniência: wiki \textperiodcentered{} fecha o arco de línguas no núcleo simbólico \textperiodcentered{} inferencia \textperiodcentered{} confiança 1.0 \textperiodcentered{} domínio programacao \textperiodcentered{} história 10 versões no jornal}

%CRISTAL {"arestas": [{"alvo": "contradicao", "confianca": 1.0, "epistemico": "fato", "origem": "curriculo", "rel": "contradiz"}, {"alvo": "exemplo", "confianca": 1.0, "epistemico": "fato", "origem": "curriculo", "rel": "depende"}, {"alvo": "virtualizacao", "confianca": 0.5, "epistemico": "fato", "origem": "wiki", "rel": "relaciona"}, {"alvo": "word", "confianca": 0.5, "epistemico": "fato", "origem": "wiki", "rel": "relaciona"}], "confianca": 1.0, "contraexemplos": [], "descricao": "Campo epistemológico vazio (sem descrição, exemplo ou dependência).", "epistemico": "fato", "exemplos": [], "id": "lacuna", "memoria": "semantica", "meta": {"dominio": "aprendizagem", "nivel": 7, "ordem": 2}, "origem": "curriculo", "palavras_chave": ["incompleto", "buraco"], "sinonimos": ["lacunas"], "tipo": "conceito", "titulo": "lacuna"}
\section{lacuna}
Campo epistemológico vazio (sem descrição, exemplo ou dependência).
\prov{sinónimos: lacunas}
\prov{relações: contradiz contradicao; depende exemplo; relaciona virtualizacao; relaciona word}
\prov{proveniência: curriculo \textperiodcentered{} fato \textperiodcentered{} confiança 1.0 \textperiodcentered{} domínio aprendizagem \textperiodcentered{} história 5 versões no jornal}

%CRISTAL {"arestas": [{"alvo": "tex", "confianca": 1.0, "epistemico": "fato", "origem": "wiki", "rel": "especializa"}, {"alvo": "separacao_conteudo_forma", "confianca": 1.0, "epistemico": "fato", "origem": "wiki", "rel": "usa"}, {"alvo": "matematica", "confianca": 1.0, "epistemico": "fato", "origem": "wiki", "rel": "explica"}], "confianca": 1.0, "contraexemplos": [], "descricao": "Lamport: uma camada de macros sobre o TeX que separa conteúdo de forma — a notação em que artigos científicos e TODO o corpus matemático do broca-so é escrito.", "epistemico": "fato", "exemplos": [], "id": "latex", "memoria": "semantica", "meta": {"autor": "Lamport", "dominio": "programacao", "fonte": ""}, "origem": "wiki", "palavras_chave": [], "sinonimos": [], "tipo": "conceito", "titulo": "LaTeX"}
\section{LaTeX}
Lamport: uma camada de macros sobre o TeX que separa conteúdo de forma — a notação em que artigos científicos e TODO o corpus matemático do broca-so é escrito.
\prov{relações: especializa tex; usa separacao\_conteudo\_forma; explica matematica}
\prov{proveniência: wiki \textperiodcentered{} fato \textperiodcentered{} confiança 1.0 \textperiodcentered{} domínio programacao \textperiodcentered{} história 4 versões no jornal}

%CRISTAL {"arestas": [{"alvo": "latex", "confianca": 1.0, "epistemico": "fato", "origem": "wiki", "rel": "parte_de"}, {"alvo": "notacao_matematica", "confianca": 1.0, "epistemico": "fato", "origem": "wiki", "rel": "explica"}], "confianca": 1.0, "contraexemplos": [], "descricao": "O modo matemático do LaTeX: entre... ou [...], cada símbolo vira um comando (∑, ınt, α) e a estrutura, uma árvore (ab). A notação matemática como marcação.", "epistemico": "fato", "exemplos": [], "id": "latex_matematico", "memoria": "semantica", "meta": {"dominio": "programacao", "fonte": ""}, "origem": "wiki", "palavras_chave": [], "sinonimos": [], "tipo": "conceito", "titulo": "LaTeX Matemático (math mode)"}
\section{LaTeX Matemático (math mode)}
O modo matemático do LaTeX: entre... ou [...], cada símbolo vira um comando (\ensuremath{\sum}, \ensuremath{\imath}nt, \ensuremath{\alpha}) e a estrutura, uma árvore (ab). A notação matemática como marcação.
\prov{relações: parte\_de latex; explica notacao\_matematica}
\prov{proveniência: wiki \textperiodcentered{} fato \textperiodcentered{} confiança 1.0 \textperiodcentered{} domínio programacao \textperiodcentered{} história 4 versões no jornal}

%CRISTAL {"arestas": [{"alvo": "parabola_algebra_conversacional", "confianca": 1.0, "epistemico": "fato", "origem": "wiki", "rel": "usa"}, {"alvo": "meta_inducao_geral_simbolo", "confianca": 1.0, "epistemico": "fato", "origem": "wiki", "rel": "especializa"}, {"alvo": "conjugados_galois_parabola", "confianca": 1.0, "epistemico": "fato", "origem": "wiki", "rel": "usa"}, {"alvo": "transformada_metalica", "confianca": 1.0, "epistemico": "fato", "origem": "wiki", "rel": "usa"}, {"alvo": "ordem_borda_caos", "confianca": 1.0, "epistemico": "fato", "origem": "wiki", "rel": "usa"}, {"alvo": "torre_unificada", "confianca": 1.0, "epistemico": "fato", "origem": "wiki", "rel": "parte_de"}, {"alvo": "mineracao_conservacao_verificavel", "confianca": 1.0, "epistemico": "fato", "origem": "wiki", "rel": "explica"}, {"alvo": "refino_continuo_floco_metais", "confianca": 1.0, "epistemico": "fato", "origem": "wiki", "rel": "explica"}, {"alvo": "absorcao_iterativa_poca", "confianca": 1.0, "epistemico": "fato", "origem": "wiki", "rel": "explica"}, {"alvo": "colapso", "confianca": 1.0, "epistemico": "fato", "origem": "wiki", "rel": "explica"}, {"alvo": "tiffany", "confianca": 1.0, "epistemico": "fato", "origem": "wiki", "rel": "explica"}, {"alvo": "ingestao_refino_parabola", "confianca": 1.0, "epistemico": "fato", "origem": "wiki", "rel": "usa"}, {"alvo": "orbita_aterrada_caos", "confianca": 1.0, "epistemico": "fato", "origem": "wiki", "rel": "explica"}, {"alvo": "sandbox_cristal_ao_so", "confianca": 1.0, "epistemico": "fato", "origem": "wiki", "rel": "explica"}, {"alvo": "conservacao_parabola_costura", "confianca": 1.0, "epistemico": "fato", "origem": "wiki", "rel": "usa"}], "confianca": 1.0, "contraexemplos": [], "descricao": "Um só operador (classe_parabola: cristal/borda/caos) governa TRÊS domínios do broca-so, só mudando o símbolo — a meta-indução geral em ato. (1) MINERAÇÃO: o FLOCO no metal (cristal absorve/Pisot, borda vira camada/Salem, caos ao purgatório/fora). (2) RECUPERAÇÃO: a PERGUNTA no grafo (aboutness=c², o colapso Ψ escolhe a face, fail-closed se c²<0.5). (3) EXECUÇÃO: o PROGRAMA no sandbox (cristal=o laço FECHA na posição O(1); borda=roda o laço O(k); caos=reproduz, órbita não aterra). c²=ordem/associativo (conjugada dentro do círculo), a=1−c²=dissipação. NÃO é analogia: é o MESMO classe_parabola chamado nos três — floco, pergunta e programa sob a mesma conservação. Cristal converge, borda corre na borda, caos dissipa.", "epistemico": "inferencia", "exemplos": [], "id": "lei_unica_tres_dominios", "memoria": "semantica", "meta": {"dominio": "programacao", "fonte": "parabola_algebra.classe_parabola em: transicoes (mineração), colapso (recuperação), sandbox (execução)"}, "origem": "wiki", "palavras_chave": [], "sinonimos": [], "tipo": "conceito", "titulo": "A Lei Única — a parábola a+c²=1 rege mineração, recuperação e execução"}
\section{A Lei Única — a parábola a+c\ensuremath{^{2}}=1 rege mineração, recuperação e execução}
Um só operador (classe\_parabola: cristal/borda/caos) governa TRÊS domínios do broca-so, só mudando o símbolo — a meta-indução geral em ato. (1) MINERAÇÃO: o FLOCO no metal (cristal absorve/Pisot, borda vira camada/Salem, caos ao purgatório/fora). (2) RECUPERAÇÃO: a PERGUNTA no grafo (aboutness=c\ensuremath{^{2}}, o colapso \ensuremath{\Psi} escolhe a face, fail-closed se c\ensuremath{^{2}}<0.5). (3) EXECUÇÃO: o PROGRAMA no sandbox (cristal=o laço FECHA na posição O(1); borda=roda o laço O(k); caos=reproduz, órbita não aterra). c\ensuremath{^{2}}=ordem/associativo (conjugada dentro do círculo), a=1\ensuremath{-}c\ensuremath{^{2}}=dissipação. NÃO é analogia: é o MESMO classe\_parabola chamado nos três — floco, pergunta e programa sob a mesma conservação. Cristal converge, borda corre na borda, caos dissipa.
\prov{relações: usa parabola\_algebra\_conversacional; especializa meta\_inducao\_geral\_simbolo; usa conjugados\_galois\_parabola; usa transformada\_metalica; usa ordem\_borda\_caos; parte\_de torre\_unificada; explica mineracao\_conservacao\_verificavel; explica refino\_continuo\_floco\_metais; explica absorcao\_iterativa\_poca; explica colapso; explica tiffany; usa ingestao\_refino\_parabola; explica orbita\_aterrada\_caos; explica sandbox\_cristal\_ao\_so; usa conservacao\_parabola\_costura}
\prov{proveniência: wiki \textperiodcentered{} parabola\_algebra.classe\_parabola em: transicoes (mineração), colapso (recuperação), sandbox (execução) \textperiodcentered{} inferencia \textperiodcentered{} confiança 1.0 \textperiodcentered{} domínio programacao \textperiodcentered{} história 16 versões no jornal}

%CRISTAL {"arestas": [{"alvo": "tiffany_instancia_do_corpo_estelar", "confianca": 1.0, "epistemico": "fato", "origem": "wiki", "rel": "parte_de"}, {"alvo": "metal_do_so", "confianca": 1.0, "epistemico": "fato", "origem": "wiki", "rel": "usa"}, {"alvo": "operacoes_reservadas_no_metal", "confianca": 1.0, "epistemico": "fato", "origem": "wiki", "rel": "usa"}, {"alvo": "regua_de_isomorfia_multi_idioma", "confianca": 1.0, "epistemico": "fato", "origem": "wiki", "rel": "usa"}, {"alvo": "bai_psi_colapso", "confianca": 1.0, "epistemico": "fato", "origem": "wiki", "rel": "usa"}], "confianca": 1.0, "contraexemplos": [], "descricao": "A Tiffany operante (linguagem/tiffany.py): 11 FERRAMENTAS (as operações reservadas do metal + os módulos verificados: gato, soma, produto, norma, entropia, crivo, endereço, espectro, classe, fundir, ligação) e 4 PIPELINES automatizados — a dinâmica entre elas, a saída de uma alimentando a próxima: 'estrela' (matriz→espectro→ρ→classe), 'corpo' (⊕/⊗), 'endereco' (β-numeração) e 'entropia' (log σ). Os cinco verbos: VIBRAR (a órbita do gato), COLAPSAR (Ψ=argmax, que paga entropia), ANINHAR (compor, a boneca), CONSERVAR (a norma) e DISSIPAR (a entropia). Roda em qualquer idioma (a régua do BAI). Verif. 7/7: Pell A_2^4·(1,0)=(29,12); Ψ=argmax; o pipeline estrela classifica; 'cat'==gato.", "epistemico": "fato", "exemplos": [], "id": "lib_tiffany_no_metal", "memoria": "semantica", "meta": {"dominio": "sistema", "fonte": "tiffany estelar"}, "origem": "wiki", "palavras_chave": [], "sinonimos": [], "tipo": "conceito", "titulo": "A Lib da Tiffany no Metal (ferramentas + pipelines)"}
\section{A Lib da Tiffany no Metal (ferramentas + pipelines)}
A Tiffany operante (linguagem/tiffany.py): 11 FERRAMENTAS (as operações reservadas do metal + os módulos verificados: gato, soma, produto, norma, entropia, crivo, endereço, espectro, classe, fundir, ligação) e 4 PIPELINES automatizados — a dinâmica entre elas, a saída de uma alimentando a próxima: 'estrela' (matriz\ensuremath{\to}espectro\ensuremath{\to}\ensuremath{\rho}\ensuremath{\to}classe), 'corpo' (\ensuremath{\oplus}/\ensuremath{\otimes}), 'endereco' (\ensuremath{\beta}-numeração) e 'entropia' (log \ensuremath{\sigma}). Os cinco verbos: VIBRAR (a órbita do gato), COLAPSAR (\ensuremath{\Psi}=argmax, que paga entropia), ANINHAR (compor, a boneca), CONSERVAR (a norma) e DISSIPAR (a entropia). Roda em qualquer idioma (a régua do BAI). Verif. 7/7: Pell A\_2\^{}4\textperiodcentered (1,0)=(29,12); \ensuremath{\Psi}=argmax; o pipeline estrela classifica; 'cat'==gato.
\prov{relações: parte\_de tiffany\_instancia\_do\_corpo\_estelar; usa metal\_do\_so; usa operacoes\_reservadas\_no\_metal; usa regua\_de\_isomorfia\_multi\_idioma; usa bai\_psi\_colapso}
\prov{proveniência: wiki \textperiodcentered{} tiffany estelar \textperiodcentered{} fato \textperiodcentered{} confiança 1.0 \textperiodcentered{} domínio sistema \textperiodcentered{} história 6 versões no jornal}

%CRISTAL {"arestas": [{"alvo": "cfr_minimizacao_arrependimento", "confianca": 1.0, "epistemico": "fato", "origem": "wiki", "rel": "usa"}], "confianca": 1.0, "contraexemplos": [], "descricao": "Libratus (2017) bate profissionais no No-Limit Hold'em heads-up; Pluribus (2019) vence no pôquer de 6 jogadores — o primeiro triunfo de IA num jogo MULTIJOGADOR de informação imperfeita, via CFR + busca.", "epistemico": "fato", "exemplos": [], "id": "libratus_pluribus", "memoria": "semantica", "meta": {"dominio": "ia", "fonte": "jogos e IA"}, "origem": "wiki", "palavras_chave": [], "sinonimos": [], "tipo": "conceito", "titulo": "Libratus e Pluribus (pôquer sobre-humano)"}
\section{Libratus e Pluribus (pôquer sobre-humano)}
Libratus (2017) bate profissionais no No-Limit Hold'em heads-up; Pluribus (2019) vence no pôquer de 6 jogadores — o primeiro triunfo de IA num jogo MULTIJOGADOR de informação imperfeita, via CFR + busca.
\prov{relações: usa cfr\_minimizacao\_arrependimento}
\prov{proveniência: wiki \textperiodcentered{} jogos e IA \textperiodcentered{} fato \textperiodcentered{} confiança 1.0 \textperiodcentered{} domínio ia \textperiodcentered{} história 2 versões no jornal}

%CRISTAL {"arestas": [{"alvo": "hopfield", "confianca": 1.0, "epistemico": "fato", "origem": "wiki", "rel": "usa"}, {"alvo": "refinar_densificar", "confianca": 1.0, "epistemico": "fato", "origem": "wiki", "rel": "relaciona"}], "confianca": 1.0, "contraexemplos": [], "descricao": "sugerir_ligacoes só sugere; ligar(auto=True) exige overlap ≥ LIMIAR_OV (238) e deriva confianca=ov/256. A curadoria transversal nunca liga sozinha — as arestas curadas das semillas são explícitas.", "epistemico": "inferencia", "exemplos": [], "id": "ligacao_por_overlap_limiar", "memoria": "semantica", "meta": {"autor": "broca-so", "dominio": "operacao_so", "fonte": "conversa/conhecimento/ligar.py"}, "origem": "wiki", "palavras_chave": [], "sinonimos": [], "tipo": "conceito", "titulo": "Ligar por Overlap (nunca automático abaixo do limiar)"}
\section{Ligar por Overlap (nunca automático abaixo do limiar)}
sugerir\_ligacoes só sugere; ligar(auto=True) exige overlap \ensuremath{\ge} LIMIAR\_OV (238) e deriva confianca=ov/256. A curadoria transversal nunca liga sozinha — as arestas curadas das semillas são explícitas.
\prov{relações: usa hopfield; relaciona refinar\_densificar}
\prov{proveniência: wiki \textperiodcentered{} conversa/conhecimento/ligar.py \textperiodcentered{} inferencia \textperiodcentered{} confiança 1.0 \textperiodcentered{} domínio operacao\_so \textperiodcentered{} história 3 versões no jornal}

%CRISTAL {"arestas": [{"alvo": "gramatica", "confianca": 1.0, "epistemico": "fato", "origem": "curriculo", "rel": "parte_de"}, {"alvo": "ontologia", "confianca": 0.5, "epistemico": "fato", "origem": "wiki", "rel": "relaciona"}, {"alvo": "word", "confianca": 0.5, "epistemico": "fato", "origem": "wiki", "rel": "relaciona"}, {"alvo": "misticismo_nao_dual", "confianca": 0.5, "epistemico": "fato", "origem": "wiki", "rel": "relaciona"}, {"alvo": "vida-morte", "confianca": 0.5, "epistemico": "fato", "origem": "wiki", "rel": "relaciona"}, {"alvo": "musica", "confianca": 0.5, "epistemico": "fato", "origem": "wiki", "rel": "relaciona"}], "confianca": 1.0, "contraexemplos": [], "descricao": "Sistema simbólico de expressão e comunicação; base para explicar qualquer conceito.", "epistemico": "fato", "exemplos": [], "id": "linguagem", "memoria": "semantica", "meta": {"dominio": "linguagem", "nivel": 3, "ordem": 0}, "origem": "curriculo", "palavras_chave": ["comunicação", "símbolos", "signos"], "sinonimos": [], "tipo": "conceito", "titulo": "linguagem"}
\section{linguagem}
Sistema simbólico de expressão e comunicação; base para explicar qualquer conceito.
\prov{relações: parte\_de gramatica; relaciona ontologia; relaciona word; relaciona misticismo\_nao\_dual; relaciona vida-morte; relaciona musica}
\prov{proveniência: curriculo \textperiodcentered{} fato \textperiodcentered{} confiança 1.0 \textperiodcentered{} domínio linguagem \textperiodcentered{} história 7 versões no jornal}

%CRISTAL {"arestas": [{"alvo": "linguagem_de_programacao", "confianca": 1.0, "epistemico": "fato", "origem": "wiki", "rel": "parte_de"}, {"alvo": "arquitetura_von_neumann", "confianca": 1.0, "epistemico": "fato", "origem": "wiki", "rel": "usa"}, {"alvo": "2_gramatica_assembly", "confianca": 1.0, "epistemico": "fato", "origem": "wiki", "rel": "relaciona"}], "confianca": 1.0, "contraexemplos": [], "descricao": "A notação simbólica das instruções da CPU — um passo acima do binário; onde a arquitetura fica nua.", "epistemico": "fato", "exemplos": [], "id": "linguagem_assembly", "memoria": "semantica", "meta": {"dominio": "programacao", "fonte": ""}, "origem": "wiki", "palavras_chave": [], "sinonimos": [], "tipo": "conceito", "titulo": "Assembly"}
\section{Assembly}
A notação simbólica das instruções da CPU — um passo acima do binário; onde a arquitetura fica nua.
\prov{relações: parte\_de linguagem\_de\_programacao; usa arquitetura\_von\_neumann; relaciona 2\_gramatica\_assembly}
\prov{proveniência: wiki \textperiodcentered{} fato \textperiodcentered{} confiança 1.0 \textperiodcentered{} domínio programacao \textperiodcentered{} história 7 versões no jornal}

%CRISTAL {"arestas": [{"alvo": "programacao_declarativa", "confianca": 1.0, "epistemico": "fato", "origem": "wiki", "rel": "usa"}, {"alvo": "hierarquia_de_chomsky", "confianca": 1.0, "epistemico": "fato", "origem": "wiki", "rel": "usa"}, {"alvo": "gramatica_bnf", "confianca": 1.0, "epistemico": "fato", "origem": "wiki", "rel": "usa"}], "confianca": 1.0, "contraexemplos": [], "descricao": "Anota o texto com marcas (tags) que dizem o que cada parte É — separando conteúdo de forma. Descreve, não computa; declarativa por natureza.", "epistemico": "fato", "exemplos": [], "id": "linguagem_de_marcacao", "memoria": "semantica", "meta": {"dominio": "programacao", "fonte": ""}, "origem": "wiki", "palavras_chave": [], "sinonimos": [], "tipo": "conceito", "titulo": "Linguagem de Marcação"}
\section{Linguagem de Marcação}
Anota o texto com marcas (tags) que dizem o que cada parte É — separando conteúdo de forma. Descreve, não computa; declarativa por natureza.
\prov{relações: usa programacao\_declarativa; usa hierarquia\_de\_chomsky; usa gramatica\_bnf}
\prov{proveniência: wiki \textperiodcentered{} fato \textperiodcentered{} confiança 1.0 \textperiodcentered{} domínio programacao \textperiodcentered{} história 4 versões no jornal}

%CRISTAL {"arestas": [{"alvo": "programacao", "confianca": 1.0, "epistemico": "fato", "origem": "wiki", "rel": "parte_de"}, {"alvo": "maquina_de_turing", "confianca": 1.0, "epistemico": "fato", "origem": "wiki", "rel": "usa"}, {"alvo": "linguagem", "confianca": 1.0, "epistemico": "fato", "origem": "wiki", "rel": "relaciona"}], "confianca": 1.0, "contraexemplos": [], "descricao": "Um sistema formal de notação para escrever programas — sintaxe, semântica e um modelo de execução.", "epistemico": "fato", "exemplos": [], "id": "linguagem_de_programacao", "memoria": "semantica", "meta": {"dominio": "programacao", "fonte": ""}, "origem": "wiki", "palavras_chave": [], "sinonimos": [], "tipo": "conceito", "titulo": "Linguagem de Programação"}
\section{Linguagem de Programação}
Um sistema formal de notação para escrever programas — sintaxe, semântica e um modelo de execução.
\prov{relações: parte\_de programacao; usa maquina\_de\_turing; relaciona linguagem}
\prov{proveniência: wiki \textperiodcentered{} fato \textperiodcentered{} confiança 1.0 \textperiodcentered{} domínio programacao \textperiodcentered{} história 7 versões no jornal}

%CRISTAL {"arestas": [{"alvo": "expressividade_sub_turing", "confianca": 1.0, "epistemico": "fato", "origem": "wiki", "rel": "usa"}, {"alvo": "linguagem_matriz_gatos", "confianca": 1.0, "epistemico": "fato", "origem": "wiki", "rel": "parte_de"}, {"alvo": "tiffany", "confianca": 1.0, "epistemico": "fato", "origem": "wiki", "rel": "parte_de"}], "confianca": 1.0, "contraexemplos": [], "descricao": "O rigor do não-maquiar aplicado à linguagem: a máquina ATUAL é estritamente SUB-TURING ([[expressividade_sub_turing]]) — endereços de LOAD/STORE e offsets de salto são literais fixos; registradores nunca selecionam endereço; sem indireção nem recursão dinâmica (prova de impossibilidade). O que É [F]: o gato, cifra_an, a ISA de 12, a ULA, o m-bonacci, a transformada, a conservação det=±1. O que é DESIGN [D]: a linguagem da matriz de gatos. O que é CONJECTURA [C]: que a matriz de gatos com ENDEREÇAMENTO POSICIONAL (a transformada dá a posição como endereço, a indireção que falta) possa transcender o limite sub-Turing — NÃO afirmado; é a próxima pergunta. [F] o limite; [C] a saída.", "epistemico": "fato", "exemplos": [], "id": "linguagem_gatos_sub_turing_honesto", "memoria": "semantica", "meta": {"autor": "Lopes/broca-so", "dominio": "computacao", "fonte": "expressividade_sub_turing; 4_obstrucao_formal_nao_turing-completa; doc/linguagem_gatos/paper §6"}, "origem": "wiki", "palavras_chave": [], "sinonimos": [], "tipo": "conceito", "titulo": "Honesto: a máquina é sub-Turing; transcender via posição é conjectura aberta (não afirmada)"}
\section{Honesto: a máquina é sub-Turing; transcender via posição é conjectura aberta (não afirmada)}
O rigor do não-maquiar aplicado à linguagem: a máquina ATUAL é estritamente SUB-TURING ([[expressividade\_sub\_turing]]) — endereços de LOAD/STORE e offsets de salto são literais fixos; registradores nunca selecionam endereço; sem indireção nem recursão dinâmica (prova de impossibilidade). O que É [F]: o gato, cifra\_an, a ISA de 12, a ULA, o m-bonacci, a transformada, a conservação det=$\pm$1. O que é DESIGN [D]: a linguagem da matriz de gatos. O que é CONJECTURA [C]: que a matriz de gatos com ENDEREÇAMENTO POSICIONAL (a transformada dá a posição como endereço, a indireção que falta) possa transcender o limite sub-Turing — NÃO afirmado; é a próxima pergunta. [F] o limite; [C] a saída.
\prov{relações: usa expressividade\_sub\_turing; parte\_de linguagem\_matriz\_gatos; parte\_de tiffany}
\prov{proveniência: wiki \textperiodcentered{} expressividade\_sub\_turing; 4\_obstrucao\_formal\_nao\_turing-completa; doc/linguagem\_gatos/paper §6 \textperiodcentered{} fato \textperiodcentered{} confiança 1.0 \textperiodcentered{} domínio computacao \textperiodcentered{} história 39 versões no jornal}

%CRISTAL {"arestas": [{"alvo": "matematica", "confianca": 1.0, "epistemico": "fato", "origem": "wiki", "rel": "parte_de"}, {"alvo": "sintaxe", "confianca": 1.0, "epistemico": "fato", "origem": "wiki", "rel": "relaciona"}, {"alvo": "hierarquia_de_chomsky", "confianca": 1.0, "epistemico": "fato", "origem": "wiki", "rel": "usa"}], "confianca": 1.0, "contraexemplos": [], "descricao": "A matemática é uma língua formal com sintaxe própria: um discurso (definição→teorema→prova) e um sistema de notação. Não é português com símbolos — tem gramática livre-de-contexto.", "epistemico": "fato", "exemplos": [], "id": "linguagem_matematica", "memoria": "semantica", "meta": {"dominio": "programacao", "fonte": ""}, "origem": "wiki", "palavras_chave": [], "sinonimos": [], "tipo": "conceito", "titulo": "A Linguagem Matemática"}
\section{A Linguagem Matemática}
A matemática é uma língua formal com sintaxe própria: um discurso (definição\ensuremath{\to}teorema\ensuremath{\to}prova) e um sistema de notação. Não é português com símbolos — tem gramática livre-de-contexto.
\prov{relações: parte\_de matematica; relaciona sintaxe; usa hierarquia\_de\_chomsky}
\prov{proveniência: wiki \textperiodcentered{} fato \textperiodcentered{} confiança 1.0 \textperiodcentered{} domínio programacao \textperiodcentered{} história 4 versões no jornal}

%CRISTAL {"arestas": [{"alvo": "gato_generaliza_matriz_companion", "confianca": 1.0, "epistemico": "fato", "origem": "wiki", "rel": "usa"}, {"alvo": "contracao_einstein_k_tensorial", "confianca": 1.0, "epistemico": "fato", "origem": "wiki", "rel": "usa"}, {"alvo": "linguagem_gatos_sub_turing_honesto", "confianca": 1.0, "epistemico": "fato", "origem": "wiki", "rel": "usa"}, {"alvo": "isa_da_broca", "confianca": 1.0, "epistemico": "fato", "origem": "wiki", "rel": "usa"}, {"alvo": "ula_por_componente", "confianca": 1.0, "epistemico": "fato", "origem": "wiki", "rel": "usa"}, {"alvo": "cifra_metais_an", "confianca": 1.0, "epistemico": "fato", "origem": "wiki", "rel": "usa"}, {"alvo": "compilador_fractal_broca_cc", "confianca": 1.0, "epistemico": "fato", "origem": "wiki", "rel": "usa"}, {"alvo": "transformada_metalica", "confianca": 1.0, "epistemico": "fato", "origem": "wiki", "rel": "usa"}, {"alvo": "interpretador_nucleo_gatos", "confianca": 1.0, "epistemico": "fato", "origem": "wiki", "rel": "usa"}, {"alvo": "dsl_frontend_linguagem_gatos", "confianca": 1.0, "epistemico": "fato", "origem": "wiki", "rel": "usa"}, {"alvo": "tiffany", "confianca": 1.0, "epistemico": "fato", "origem": "wiki", "rel": "parte_de"}, {"alvo": "linguagem_de_programacao", "confianca": 1.0, "epistemico": "fato", "origem": "wiki", "rel": "especializa"}], "confianca": 1.0, "contraexemplos": [], "descricao": "O paper/hub: a FPGA (Field-Programmable Gate Array) reinterpretada — a PORTA é o GATO (Aₙ=[[n,1],[1,0]], o metal n; cifraan), e o campo programável é a MATRIZ DE GATOS: a companion m×m (a rede de m neurônios = o m-bonacci), fiada pelos coeficientes do DUPOLINÔMIO de grau m — o 'programa'. Programar o broca-so é escolher uma matriz de gatos; executar é CONTRAIR (Einstein). O software do so se reproduz como gatos associados por matriz. A base já roda (ISA de 12 opcodes, ULA, cifraan, broca-cc); a linguagem da matriz de gatos LIFTA essa base (a IR vira o dupolinômio, os opcodes viram gatos — 'zero opcodes novos', toda instrução é composição de gatos). O paper documenta DOIS NÍVEIS, ambos rodando: (i) o NÚCLEO SUB-TURING (bloco=matriz de gatos, execução=contração, laço=posição Aₙk — o motor determinístico, [[interpretadorₙucleogatos]]); e (ii) a TRANSCENDÊNCIA PROVADA (§6.1): com indireção posicional a linguagem vira Turing-completa ([[indirecaoposicionaltranscendeturing]]). Marcação [F]/[D]/[C] em tudo. Ver o paper (resumo, §6.1, §7) e linguagem/gatomatriz,interpretador,universal.py (PoCs rodáveis).", "epistemico": "inferencia", "exemplos": [], "id": "linguagem_matriz_gatos", "memoria": "semantica", "meta": {"autor": "Lopes/broca-so", "dominio": "computacao", "fonte": "doc/linguagem_gatos/paper_matrizes_portas_gatos.md (resumo, §6.1, §7); linguagem/gato_matriz.py, interpretador.py, universal.py; isa_da_broca"}, "origem": "wiki", "palavras_chave": [], "sinonimos": [], "tipo": "conceito", "titulo": "A linguagem do broca-so sobre a MATRIZ DE GATOS (Matrizes de Portas Programáveis em Campo)"}
\section{A linguagem do broca-so sobre a MATRIZ DE GATOS (Matrizes de Portas Programáveis em Campo)}
O paper/hub: a FPGA (Field-Programmable Gate Array) reinterpretada — a PORTA é o GATO (A\ensuremath{_{n}}=[[n,1],[1,0]], o metal n; cifraan), e o campo programável é a MATRIZ DE GATOS: a companion m$\times$m (a rede de m neurônios = o m-bonacci), fiada pelos coeficientes do DUPOLINÔMIO de grau m — o 'programa'. Programar o broca-so é escolher uma matriz de gatos; executar é CONTRAIR (Einstein). O software do so se reproduz como gatos associados por matriz. A base já roda (ISA de 12 opcodes, ULA, cifraan, broca-cc); a linguagem da matriz de gatos LIFTA essa base (a IR vira o dupolinômio, os opcodes viram gatos — 'zero opcodes novos', toda instrução é composição de gatos). O paper documenta DOIS NÍVEIS, ambos rodando: (i) o NÚCLEO SUB-TURING (bloco=matriz de gatos, execução=contração, laço=posição A\ensuremath{_{n}}k — o motor determinístico, [[interpretador\ensuremath{_{n}}ucleogatos]]); e (ii) a TRANSCENDÊNCIA PROVADA (§6.1): com indireção posicional a linguagem vira Turing-completa ([[indirecaoposicionaltranscendeturing]]). Marcação [F]/[D]/[C] em tudo. Ver o paper (resumo, §6.1, §7) e linguagem/gatomatriz,interpretador,universal.py (PoCs rodáveis).
\prov{relações: usa gato\_generaliza\_matriz\_companion; usa contracao\_einstein\_k\_tensorial; usa linguagem\_gatos\_sub\_turing\_honesto; usa isa\_da\_broca; usa ula\_por\_componente; usa cifra\_metais\_an; usa compilador\_fractal\_broca\_cc; usa transformada\_metalica; usa interpretador\_nucleo\_gatos; usa dsl\_frontend\_linguagem\_gatos; parte\_de tiffany; especializa linguagem\_de\_programacao}
\prov{proveniência: wiki \textperiodcentered{} doc/linguagem\_gatos/paper\_matrizes\_portas\_gatos.md (resumo, §6.1, §7); linguagem/gato\_matriz.py, interpretador.py, universal.py; isa\_da\_broca \textperiodcentered{} inferencia \textperiodcentered{} confiança 1.0 \textperiodcentered{} domínio computacao \textperiodcentered{} história 93 versões no jornal}

%CRISTAL {"arestas": [{"alvo": "computador_mineral", "confianca": 1.0, "epistemico": "fato", "origem": "wiki", "rel": "usa"}, {"alvo": "primitivas_minerais", "confianca": 1.0, "epistemico": "fato", "origem": "wiki", "rel": "usa"}, {"alvo": "biblioteca_mineral", "confianca": 1.0, "epistemico": "fato", "origem": "wiki", "rel": "usa"}, {"alvo": "design_linguagem_gatos", "confianca": 1.0, "epistemico": "fato", "origem": "wiki", "rel": "especializa"}], "confianca": 1.0, "contraexemplos": [], "descricao": "A linguagem completa sobre o computador mineral (gato=ULA, cristal=memória): grau 2 (o gato Aₙ) → grau m (o programa). PRIMITIVAS (rodam, 12/12): memória (x=v, mem[] indireta=Turing), ULA (+ − gato n, ×via laço), controle (while, if/else, for, comparações), E-S (print), e o LAÇO=POSIÇÃO (o while-de-gato dobra em Aₙ^k). BIBLIOTECAS no metal (mecânica/geometria/física/transformada). AUTORIZAÇÃO = o BAI (o SO). Já no grafo (design_*, dsl_*), aqui fundamentado e cristalizado.", "epistemico": "inferencia", "exemplos": [], "id": "linguagem_mineral", "memoria": "semantica", "meta": {"dominio": "computacao", "fonte": "linguagem mineral"}, "origem": "wiki", "palavras_chave": [], "sinonimos": [], "tipo": "conceito", "titulo": "A Linguagem Mineral (primitivas + bibliotecas + SO)"}
\section{A Linguagem Mineral (primitivas + bibliotecas + SO)}
A linguagem completa sobre o computador mineral (gato=ULA, cristal=memória): grau 2 (o gato A\ensuremath{_{n}}) \ensuremath{\to} grau m (o programa). PRIMITIVAS (rodam, 12/12): memória (x=v, mem[] indireta=Turing), ULA (+ \ensuremath{-} gato n, $\times$via laço), controle (while, if/else, for, comparações), E-S (print), e o LAÇO=POSIÇÃO (o while-de-gato dobra em A\ensuremath{_{n}}\^{}k). BIBLIOTECAS no metal (mecânica/geometria/física/transformada). AUTORIZAÇÃO = o BAI (o SO). Já no grafo (design\_*, dsl\_*), aqui fundamentado e cristalizado.
\prov{relações: usa computador\_mineral; usa primitivas\_minerais; usa biblioteca\_mineral; especializa design\_linguagem\_gatos}
\prov{proveniência: wiki \textperiodcentered{} linguagem mineral \textperiodcentered{} inferencia \textperiodcentered{} confiança 1.0 \textperiodcentered{} domínio computacao \textperiodcentered{} história 5 versões no jornal}

%CRISTAL {"arestas": [{"alvo": "linguagem_de_programacao", "confianca": 1.0, "epistemico": "fato", "origem": "wiki", "rel": "parte_de"}, {"alvo": "computacao_matricial", "confianca": 1.0, "epistemico": "fato", "origem": "wiki", "rel": "usa"}, {"alvo": "programacao_tacita", "confianca": 1.0, "epistemico": "fato", "origem": "wiki", "rel": "usa"}], "confianca": 1.0, "contraexemplos": [], "descricao": "Iverson (APL, J, K, Futhark): o primitivo é o ARRANJO inteiro, não o escalar; operações agem sobre a estrutura toda de uma vez, sem laços explícitos.", "epistemico": "fato", "exemplos": [], "id": "linguagens_de_array", "memoria": "semantica", "meta": {"autor": "Iverson", "dominio": "programacao", "fonte": ""}, "origem": "wiki", "palavras_chave": [], "sinonimos": [], "tipo": "conceito", "titulo": "Linguagens de Array"}
\section{Linguagens de Array}
Iverson (APL, J, K, Futhark): o primitivo é o ARRANJO inteiro, não o escalar; operações agem sobre a estrutura toda de uma vez, sem laços explícitos.
\prov{relações: parte\_de linguagem\_de\_programacao; usa computacao\_matricial; usa programacao\_tacita}
\prov{proveniência: wiki \textperiodcentered{} fato \textperiodcentered{} confiança 1.0 \textperiodcentered{} domínio programacao \textperiodcentered{} história 8 versões no jornal}

%CRISTAL {"arestas": [{"alvo": "programacao_funcional", "confianca": 1.0, "epistemico": "fato", "origem": "wiki", "rel": "relaciona"}, {"alvo": "turing_completude", "confianca": 1.0, "epistemico": "fato", "origem": "wiki", "rel": "contradiz"}, {"alvo": "sistema_de_tipos", "confianca": 1.0, "epistemico": "fato", "origem": "wiki", "rel": "usa"}], "confianca": 1.0, "contraexemplos": [], "descricao": "Linguagens em que TODO programa termina (Agda, Idris, o fragmento total de Coq) — sub-Turing de propósito, trocando poder por garantia de parada.", "epistemico": "fato", "exemplos": [], "id": "linguagens_totais", "memoria": "semantica", "meta": {"dominio": "programacao", "fonte": ""}, "origem": "wiki", "palavras_chave": [], "sinonimos": [], "tipo": "conceito", "titulo": "Linguagens Totais"}
\section{Linguagens Totais}
Linguagens em que TODO programa termina (Agda, Idris, o fragmento total de Coq) — sub-Turing de propósito, trocando poder por garantia de parada.
\prov{relações: relaciona programacao\_funcional; contradiz turing\_completude; usa sistema\_de\_tipos}
\prov{proveniência: wiki \textperiodcentered{} fato \textperiodcentered{} confiança 1.0 \textperiodcentered{} domínio programacao \textperiodcentered{} história 8 versões no jornal}

%CRISTAL {"arestas": [{"alvo": "kernel", "confianca": 1.0, "epistemico": "fato", "origem": "curriculo", "rel": "especializa"}, {"alvo": "python", "confianca": 0.5, "epistemico": "fato", "origem": "wiki", "rel": "relaciona"}, {"alvo": "mamifero", "confianca": 0.5, "epistemico": "fato", "origem": "wiki", "rel": "relaciona"}, {"alvo": "topologia", "confianca": 0.5, "epistemico": "fato", "origem": "wiki", "rel": "relaciona"}, {"alvo": "recursao", "confianca": 0.5, "epistemico": "fato", "origem": "wiki", "rel": "relaciona"}, {"alvo": "regua_associativa_e_o_risco_de_vazamento", "confianca": 0.5, "epistemico": "fato", "origem": "wiki", "rel": "relaciona"}, {"alvo": "ruido", "confianca": 0.5, "epistemico": "fato", "origem": "wiki", "rel": "relaciona"}], "confianca": 1.0, "contraexemplos": [], "descricao": "Resolve símbolos entre object files.", "epistemico": "fato", "exemplos": [], "id": "linker", "memoria": "semantica", "meta": {"dominio": "programacao", "nivel": 6, "ordem": 12}, "origem": "curriculo", "palavras_chave": [], "sinonimos": [], "tipo": "conceito", "titulo": "linker"}
\section{linker}
Resolve símbolos entre object files.
\prov{relações: especializa kernel; relaciona python; relaciona mamifero; relaciona topologia; relaciona recursao; relaciona regua\_associativa\_e\_o\_risco\_de\_vazamento; relaciona ruido}
\prov{proveniência: curriculo \textperiodcentered{} fato \textperiodcentered{} confiança 1.0 \textperiodcentered{} domínio programacao \textperiodcentered{} história 13 versões no jornal}

%CRISTAL {"arestas": [{"alvo": "unix", "confianca": 1.0, "epistemico": "fato", "origem": "wiki", "rel": "especializa"}, {"alvo": "kernel", "confianca": 1.0, "epistemico": "fato", "origem": "wiki", "rel": "usa"}, {"alvo": "gnu", "confianca": 1.0, "epistemico": "fato", "origem": "wiki", "rel": "usa"}, {"alvo": "c", "confianca": 1.0, "epistemico": "fato", "origem": "wiki", "rel": "usa"}], "confianca": 1.0, "contraexemplos": [], "descricao": "Torvalds (1991): um kernel livre estilo Unix que, com o userland GNU, virou o SO dos servidores, da nuvem, do Android e do supercomputador.", "epistemico": "fato", "exemplos": [], "id": "linux", "memoria": "semantica", "meta": {"autor": "Torvalds", "dominio": "programacao", "fonte": ""}, "origem": "wiki", "palavras_chave": [], "sinonimos": [], "tipo": "conceito", "titulo": "Linux"}
\section{Linux}
Torvalds (1991): um kernel livre estilo Unix que, com o userland GNU, virou o SO dos servidores, da nuvem, do Android e do supercomputador.
\prov{relações: especializa unix; usa kernel; usa gnu; usa c}
\prov{proveniência: wiki \textperiodcentered{} fato \textperiodcentered{} confiança 1.0 \textperiodcentered{} domínio programacao \textperiodcentered{} história 5 versões no jornal}

%CRISTAL {"arestas": [{"alvo": "programacao_funcional", "confianca": 1.0, "epistemico": "fato", "origem": "wiki", "rel": "usa"}, {"alvo": "calculo_lambda", "confianca": 1.0, "epistemico": "fato", "origem": "wiki", "rel": "explica"}, {"alvo": "metaprogramacao", "confianca": 1.0, "epistemico": "fato", "origem": "wiki", "rel": "usa"}], "confianca": 1.0, "contraexemplos": [], "descricao": "McCarthy (1958): a segunda linguagem mais antiga; código É dado (S-expressões), λ-cálculo executável, macros. A raiz funcional.", "epistemico": "fato", "exemplos": [], "id": "lisp", "memoria": "semantica", "meta": {"autor": "McCarthy", "dominio": "programacao", "fonte": ""}, "origem": "wiki", "palavras_chave": [], "sinonimos": [], "tipo": "conceito", "titulo": "Lisp"}
\section{Lisp}
McCarthy (1958): a segunda linguagem mais antiga; código É dado (S-expressões), \ensuremath{\lambda}-cálculo executável, macros. A raiz funcional.
\prov{relações: usa programacao\_funcional; explica calculo\_lambda; usa metaprogramacao}
\prov{proveniência: wiki \textperiodcentered{} fato \textperiodcentered{} confiança 1.0 \textperiodcentered{} domínio programacao \textperiodcentered{} história 8 versões no jornal}

%CRISTAL {"arestas": [{"alvo": "data_center", "confianca": 1.0, "epistemico": "fato", "origem": "wiki", "rel": "usa"}, {"alvo": "utilizacao_e_borda", "confianca": 1.0, "epistemico": "fato", "origem": "wiki", "rel": "usa"}, {"alvo": "limite_de_landauer", "confianca": 1.0, "epistemico": "fato", "origem": "wiki", "rel": "usa"}, {"alvo": "modelo_de_linguagem", "confianca": 1.0, "epistemico": "fato", "origem": "wiki", "rel": "relaciona"}], "confianca": 1.0, "contraexemplos": [], "descricao": "A inferência é uma fila: cada requisição de geração ocupa a GPU a utilização ρ — perto de ρ=1 a latência por token dispara (a borda da nuvem). E cada token apagado/computado paga o piso de Landauer (kT·ln2): pensar custa calor. A IA fecha a cadeia — energia e informação no mesmo galpão térmico que a broca-so descreve.", "epistemico": "inferencia", "exemplos": [], "id": "llm_no_data_center", "memoria": "semantica", "meta": {"dominio": "ia", "fonte": "IA e modelos de linguagem"}, "origem": "wiki", "palavras_chave": [], "sinonimos": [], "tipo": "conceito", "titulo": "O Modelo Roda no Data Center (ρ e Landauer)"}
\section{O Modelo Roda no Data Center (\ensuremath{\rho} e Landauer)}
A inferência é uma fila: cada requisição de geração ocupa a GPU a utilização \ensuremath{\rho} — perto de \ensuremath{\rho}=1 a latência por token dispara (a borda da nuvem). E cada token apagado/computado paga o piso de Landauer (kT\textperiodcentered ln2): pensar custa calor. A IA fecha a cadeia — energia e informação no mesmo galpão térmico que a broca-so descreve.
\prov{relações: usa data\_center; usa utilizacao\_e\_borda; usa limite\_de\_landauer; relaciona modelo\_de\_linguagem}
\prov{proveniência: wiki \textperiodcentered{} IA e modelos de linguagem \textperiodcentered{} inferencia \textperiodcentered{} confiança 1.0 \textperiodcentered{} domínio ia \textperiodcentered{} história 5 versões no jornal}

%CRISTAL {"arestas": [{"alvo": "analogia", "confianca": 1.0, "epistemico": "fato", "origem": "curriculo", "rel": "explica"}, {"alvo": "word", "confianca": 0.5, "epistemico": "fato", "origem": "wiki", "rel": "relaciona"}, {"alvo": "vida-morte", "confianca": 0.5, "epistemico": "fato", "origem": "wiki", "rel": "relaciona"}, {"alvo": "virtualizacao", "confianca": 0.5, "epistemico": "fato", "origem": "wiki", "rel": "relaciona"}, {"alvo": "syscall", "confianca": 0.5, "epistemico": "fato", "origem": "wiki", "rel": "relaciona"}, {"alvo": "top_assinaturas", "confianca": 0.5, "epistemico": "fato", "origem": "wiki", "rel": "relaciona"}], "confianca": 1.0, "contraexemplos": [], "descricao": "Validade de argumentos; premissas, conclusões e regras de inferência.", "epistemico": "fato", "exemplos": [], "id": "logica", "memoria": "semantica", "meta": {"dominio": "linguagem", "nivel": 3, "ordem": 4}, "origem": "curriculo", "palavras_chave": ["premissa", "conclusão", "dedução"], "sinonimos": ["lógica"], "tipo": "conceito", "titulo": "logica"}
\section{logica}
Validade de argumentos; premissas, conclusões e regras de inferência.
\prov{sinónimos: lógica}
\prov{relações: explica analogia; relaciona word; relaciona vida-morte; relaciona virtualizacao; relaciona syscall; relaciona top\_assinaturas}
\prov{proveniência: curriculo \textperiodcentered{} fato \textperiodcentered{} confiança 1.0 \textperiodcentered{} domínio linguagem \textperiodcentered{} história 7 versões no jornal}

%CRISTAL {"arestas": [{"alvo": "porta_logica", "confianca": 1.0, "epistemico": "fato", "origem": "wiki", "rel": "usa"}, {"alvo": "arquitetura_von_neumann", "confianca": 1.0, "epistemico": "fato", "origem": "wiki", "rel": "explica"}, {"alvo": "transistor", "confianca": 1.0, "epistemico": "fato", "origem": "wiki", "rel": "usa"}], "confianca": 1.0, "contraexemplos": [], "descricao": "Circuitos que computam com bits por portas lógicas (E, OU, NÃO) e memória (flip-flops) — o substrato físico de toda computação.", "epistemico": "fato", "exemplos": [], "id": "logica_digital", "memoria": "semantica", "meta": {"dominio": "programacao", "fonte": ""}, "origem": "wiki", "palavras_chave": [], "sinonimos": [], "tipo": "conceito", "titulo": "Lógica Digital"}
\section{Lógica Digital}
Circuitos que computam com bits por portas lógicas (E, OU, NÃO) e memória (flip-flops) — o substrato físico de toda computação.
\prov{relações: usa porta\_logica; explica arquitetura\_von\_neumann; usa transistor}
\prov{proveniência: wiki \textperiodcentered{} fato \textperiodcentered{} confiança 1.0 \textperiodcentered{} domínio programacao \textperiodcentered{} história 4 versões no jornal}

%CRISTAL {"arestas": [{"alvo": "design_system_cristalchain", "confianca": 1.0, "epistemico": "fato", "origem": "wiki", "rel": "parte_de"}, {"alvo": "turmalina_paraiba_ancora", "confianca": 1.0, "epistemico": "fato", "origem": "wiki", "rel": "usa"}, {"alvo": "pedra_real_calibra_o_render", "confianca": 1.0, "epistemico": "fato", "origem": "wiki", "rel": "relaciona"}], "confianca": 1.0, "contraexemplos": [], "descricao": "A gema teal (turmalina Paraíba lapidada) do design system do ChatGPT (model/design) foi EXTRAÍDA pela máquina: segmentada (a maior componente cyan, preenchendo as facetas escuras), recortada justa e COMPOSTA como logo — realce de iluminação/cor, uma SOMBRA suave deslocada e um GLOW (halo azul Paraíba), os efeitos ópticos. É a pedra REAL do design virando o logo da CristalChain. Arquivo: assets/cristalchain/logo_pedra.png; usada na capa de papers/turmalina.tex.", "epistemico": "fato", "exemplos": [], "id": "logo_pedra_extraida", "memoria": "semantica", "meta": {"dominio": "sistema", "fonte": "design elementos"}, "origem": "wiki", "palavras_chave": [], "sinonimos": [], "tipo": "conceito", "titulo": "O Logo da Pedra Extraído (sombra, iluminação, efeitos ópticos)"}
\section{O Logo da Pedra Extraído (sombra, iluminação, efeitos ópticos)}
A gema teal (turmalina Paraíba lapidada) do design system do ChatGPT (model/design) foi EXTRAÍDA pela máquina: segmentada (a maior componente cyan, preenchendo as facetas escuras), recortada justa e COMPOSTA como logo — realce de iluminação/cor, uma SOMBRA suave deslocada e um GLOW (halo azul Paraíba), os efeitos ópticos. É a pedra REAL do design virando o logo da CristalChain. Arquivo: assets/cristalchain/logo\_pedra.png; usada na capa de papers/turmalina.tex.
\prov{relações: parte\_de design\_system\_cristalchain; usa turmalina\_paraiba\_ancora; relaciona pedra\_real\_calibra\_o\_render}
\prov{proveniência: wiki \textperiodcentered{} design elementos \textperiodcentered{} fato \textperiodcentered{} confiança 1.0 \textperiodcentered{} domínio sistema \textperiodcentered{} história 4 versões no jornal}

%CRISTAL {"arestas": [{"alvo": "bsd", "confianca": 1.0, "epistemico": "fato", "origem": "wiki", "rel": "usa"}, {"alvo": "unix", "confianca": 1.0, "epistemico": "fato", "origem": "wiki", "rel": "especializa"}], "confianca": 1.0, "contraexemplos": [], "descricao": "Apple: SO de desktop com kernel Darwin (base BSD/Mach) sob uma interface própria — Unix certificado por baixo do polimento.", "epistemico": "fato", "exemplos": [], "id": "macos", "memoria": "semantica", "meta": {"autor": "Apple", "dominio": "programacao", "fonte": ""}, "origem": "wiki", "palavras_chave": [], "sinonimos": [], "tipo": "conceito", "titulo": "macOS"}
\section{macOS}
Apple: SO de desktop com kernel Darwin (base BSD/Mach) sob uma interface própria — Unix certificado por baixo do polimento.
\prov{relações: usa bsd; especializa unix}
\prov{proveniência: wiki \textperiodcentered{} fato \textperiodcentered{} confiança 1.0 \textperiodcentered{} domínio programacao \textperiodcentered{} história 3 versões no jornal}

%CRISTAL {"arestas": [{"alvo": "maquina_de_turing", "confianca": 1.0, "epistemico": "fato", "origem": "wiki", "rel": "relaciona"}, {"alvo": "turing_completude", "confianca": 1.0, "epistemico": "fato", "origem": "wiki", "rel": "explica"}, {"alvo": "expressividade_sub_turing", "confianca": 1.0, "epistemico": "fato", "origem": "wiki", "rel": "relaciona"}], "confianca": 1.0, "contraexemplos": [], "descricao": "Um modelo mínimo Turing-completo: 2 contadores, incremento/decremento e desvio-se-zero. Basta indireção para universalidade.", "epistemico": "fato", "exemplos": [], "id": "maquina_de_contadores", "memoria": "semantica", "meta": {"autor": "Minsky", "dominio": "programacao", "fonte": ""}, "origem": "wiki", "palavras_chave": [], "sinonimos": [], "tipo": "conceito", "titulo": "Máquina de Contadores (Minsky)"}
\section{Máquina de Contadores (Minsky)}
Um modelo mínimo Turing-completo: 2 contadores, incremento/decremento e desvio-se-zero. Basta indireção para universalidade.
\prov{relações: relaciona maquina\_de\_turing; explica turing\_completude; relaciona expressividade\_sub\_turing}
\prov{proveniência: wiki \textperiodcentered{} fato \textperiodcentered{} confiança 1.0 \textperiodcentered{} domínio programacao \textperiodcentered{} história 8 versões no jornal}

%CRISTAL {"arestas": [{"alvo": "arquitetura_von_neumann", "confianca": 1.0, "epistemico": "fato", "origem": "wiki", "rel": "usa"}], "confianca": 1.0, "contraexemplos": [], "descricao": "Uma CPU abstrata em software que roda bytecode portátil (JVM, CPython) — a arquitetura de von Neumann emulada.", "epistemico": "fato", "exemplos": [], "id": "maquina_virtual", "memoria": "semantica", "meta": {"dominio": "programacao", "fonte": ""}, "origem": "wiki", "palavras_chave": [], "sinonimos": [], "tipo": "conceito", "titulo": "Máquina Virtual"}
\section{Máquina Virtual}
Uma CPU abstrata em software que roda bytecode portátil (JVM, CPython) — a arquitetura de von Neumann emulada.
\prov{relações: usa arquitetura\_von\_neumann}
\prov{proveniência: wiki \textperiodcentered{} fato \textperiodcentered{} confiança 1.0 \textperiodcentered{} domínio programacao \textperiodcentered{} história 4 versões no jornal}

%CRISTAL {"arestas": [{"alvo": "linguagem_de_marcacao", "confianca": 1.0, "epistemico": "fato", "origem": "wiki", "rel": "especializa"}, {"alvo": "html", "confianca": 1.0, "epistemico": "fato", "origem": "wiki", "rel": "usa"}, {"alvo": "documentacao", "confianca": 1.0, "epistemico": "fato", "origem": "wiki", "rel": "usa"}], "confianca": 1.0, "contraexemplos": [], "descricao": "Gruber: marcação LEVE que se lê bem em texto puro e vira HTML — a língua dos READMEs, da documentação e das mensagens. Simplicidade acima de tudo.", "epistemico": "fato", "exemplos": [], "id": "markdown", "memoria": "semantica", "meta": {"autor": "Gruber", "dominio": "programacao", "fonte": ""}, "origem": "wiki", "palavras_chave": [], "sinonimos": [], "tipo": "conceito", "titulo": "Markdown"}
\section{Markdown}
Gruber: marcação LEVE que se lê bem em texto puro e vira HTML — a língua dos READMEs, da documentação e das mensagens. Simplicidade acima de tudo.
\prov{relações: especializa linguagem\_de\_marcacao; usa html; usa documentacao}
\prov{proveniência: wiki \textperiodcentered{} fato \textperiodcentered{} confiança 1.0 \textperiodcentered{} domínio programacao \textperiodcentered{} história 4 versões no jornal}

%CRISTAL {"arestas": [{"alvo": "xml", "confianca": 1.0, "epistemico": "fato", "origem": "wiki", "rel": "especializa"}, {"alvo": "matematica", "confianca": 1.0, "epistemico": "fato", "origem": "wiki", "rel": "usa"}], "confianca": 1.0, "contraexemplos": [], "descricao": "Marcação XML para fórmulas matemáticas na web — a matemática estruturada para ser lida por máquina, não só desenhada.", "epistemico": "fato", "exemplos": [], "id": "mathml", "memoria": "semantica", "meta": {"dominio": "programacao", "fonte": ""}, "origem": "wiki", "palavras_chave": [], "sinonimos": [], "tipo": "conceito", "titulo": "MathML"}
\section{MathML}
Marcação XML para fórmulas matemáticas na web — a matemática estruturada para ser lida por máquina, não só desenhada.
\prov{relações: especializa xml; usa matematica}
\prov{proveniência: wiki \textperiodcentered{} fato \textperiodcentered{} confiança 1.0 \textperiodcentered{} domínio programacao \textperiodcentered{} história 3 versões no jornal}

%CRISTAL {"arestas": [{"alvo": "hopfield", "confianca": 1.0, "epistemico": "fato", "origem": "curriculo", "rel": "continua"}, {"alvo": "circuito_logico", "confianca": 1.0, "epistemico": "fato", "origem": "curriculo", "rel": "continua"}], "confianca": 1.0, "contraexemplos": [], "descricao": "Neurónio booleano (1943): soma ponderada + limiar; origem dos circuitos lógicos e das redes.", "epistemico": "fato", "exemplos": [], "id": "mcculloch_pitts", "memoria": "semantica", "meta": {"dominio": "ia", "nivel": 5, "ordem": 1}, "origem": "curriculo", "palavras_chave": ["limiar", "AND", "OR", "1943"], "sinonimos": ["McCulloch-Pitts", "McCulloch Pitts", "neurónio formal"], "tipo": "conceito", "titulo": "mcculloch pitts"}
\section{mcculloch pitts}
Neurónio booleano (1943): soma ponderada + limiar; origem dos circuitos lógicos e das redes.
\prov{sinónimos: McCulloch-Pitts; McCulloch Pitts; neurónio formal}
\prov{relações: continua hopfield; continua circuito\_logico}
\prov{proveniência: curriculo \textperiodcentered{} fato \textperiodcentered{} confiança 1.0 \textperiodcentered{} domínio ia \textperiodcentered{} história 36 versões no jornal}

%CRISTAL {"arestas": [{"alvo": "aprendizado_profundo", "confianca": 1.0, "epistemico": "fato", "origem": "wiki", "rel": "usa"}], "confianca": 1.0, "contraexemplos": [], "descricao": "Deixar o modelo pesar dinamicamente quais partes da entrada importam para cada saída — a peça que destravou os transformers.", "epistemico": "fato", "exemplos": [], "id": "mecanismo_de_atencao", "memoria": "semantica", "meta": {"autor": "Bahdanau", "dominio": "programacao", "fonte": ""}, "origem": "wiki", "palavras_chave": [], "sinonimos": [], "tipo": "conceito", "titulo": "Mecanismo de Atenção"}
\section{Mecanismo de Atenção}
Deixar o modelo pesar dinamicamente quais partes da entrada importam para cada saída — a peça que destravou os transformers.
\prov{relações: usa aprendizado\_profundo}
\prov{proveniência: wiki \textperiodcentered{} fato \textperiodcentered{} confiança 1.0 \textperiodcentered{} domínio programacao \textperiodcentered{} história 2 versões no jornal}

%CRISTAL {"arestas": [{"alvo": "colapso_koch_tiffany", "confianca": 1.0, "epistemico": "fato", "origem": "wiki", "rel": "relaciona"}, {"alvo": "colapso_ping_pong_overlap", "confianca": 1.0, "epistemico": "fato", "origem": "wiki", "rel": "relaciona"}], "confianca": 1.0, "contraexemplos": [], "descricao": "Para cada turno computa o floco Koch da fala e da face, acumulando dissipacao_turno_bits=popcount(Φ_user XOR Φ_face); o diálogo vivo enche o recipiente XOR — o que não cabe na reta associativa (overlap) é o Φ. Mede a criação, não só a repetição.", "epistemico": "fato", "exemplos": [], "id": "mede_flocos_phi", "memoria": "semantica", "meta": {"autor": "broca-so", "dominio": "operacao_so", "fonte": "conversa/mede_flocos.py"}, "origem": "wiki", "palavras_chave": [], "sinonimos": [], "tipo": "conceito", "titulo": "Medida de Flocos (Φ não-associativo)"}
\section{Medida de Flocos (\ensuremath{\Phi} não-associativo)}
Para cada turno computa o floco Koch da fala e da face, acumulando dissipacao\_turno\_bits=popcount(\ensuremath{\Phi}\_user XOR \ensuremath{\Phi}\_face); o diálogo vivo enche o recipiente XOR — o que não cabe na reta associativa (overlap) é o \ensuremath{\Phi}. Mede a criação, não só a repetição.
\prov{relações: relaciona colapso\_koch\_tiffany; relaciona colapso\_ping\_pong\_overlap}
\prov{proveniência: wiki \textperiodcentered{} conversa/mede\_flocos.py \textperiodcentered{} fato \textperiodcentered{} confiança 1.0 \textperiodcentered{} domínio operacao\_so \textperiodcentered{} história 3 versões no jornal}

%CRISTAL {"arestas": [{"alvo": "maquina_de_turing", "confianca": 1.0, "epistemico": "fato", "origem": "curriculo", "rel": "especializa"}], "confianca": 1.0, "contraexemplos": [], "descricao": "Programa como dado regravável; separação lógica CPU↔RAM (não confundir com Peirce).", "epistemico": "fato", "exemplos": [], "id": "memoria_programavel", "memoria": "semantica", "meta": {"dominio": "computacao", "nivel": 4, "ordem": 4}, "origem": "curriculo", "palavras_chave": ["EDVAC", "stored program"], "sinonimos": ["memória programável"], "tipo": "conceito", "titulo": "memoria programavel"}
\section{memoria programavel}
Programa como dado regravável; separação lógica CPU\ensuremath{\leftrightarrow}RAM (não confundir com Peirce).
\prov{sinónimos: memória programável}
\prov{relações: especializa maquina\_de\_turing}
\prov{proveniência: curriculo \textperiodcentered{} fato \textperiodcentered{} confiança 1.0 \textperiodcentered{} domínio computacao \textperiodcentered{} história 6 versões no jornal}

%CRISTAL {"arestas": [{"alvo": "ipc", "confianca": 1.0, "epistemico": "fato", "origem": "curriculo", "rel": "usa"}, {"alvo": "recursao", "confianca": 0.5, "epistemico": "fato", "origem": "wiki", "rel": "relaciona"}, {"alvo": "rust", "confianca": 0.5, "epistemico": "fato", "origem": "wiki", "rel": "relaciona"}, {"alvo": "transcricao", "confianca": 0.5, "epistemico": "fato", "origem": "wiki", "rel": "relaciona"}, {"alvo": "o_que_e_no_projecto", "confianca": 0.5, "epistemico": "fato", "origem": "wiki", "rel": "relaciona"}, {"alvo": "paper_0_sintese", "confianca": 0.5, "epistemico": "fato", "origem": "wiki", "rel": "relaciona"}, {"alvo": "word", "confianca": 0.5, "epistemico": "fato", "origem": "wiki", "rel": "relaciona"}, {"alvo": "vontade_de_poder", "confianca": 0.5, "epistemico": "fato", "origem": "wiki", "rel": "relaciona"}, {"alvo": "kernel", "confianca": 1.0, "epistemico": "fato", "origem": "wiki", "rel": "parte_de"}, {"alvo": "gerenciamento_de_memoria", "confianca": 1.0, "epistemico": "fato", "origem": "wiki", "rel": "especializa"}], "confianca": 1.0, "contraexemplos": [], "descricao": "Dar a cada processo um espaço de endereços próprio e contínuo, mapeado em páginas para a RAM/disco — isolamento e a ilusão de memória infinita.", "epistemico": "fato", "exemplos": [], "id": "memoria_virtual", "memoria": "semantica", "meta": {"dominio": "programacao", "fonte": ""}, "origem": "wiki", "palavras_chave": [], "sinonimos": ["memoria"], "tipo": "conceito", "titulo": "Memória Virtual"}
\section{Memória Virtual}
Dar a cada processo um espaço de endereços próprio e contínuo, mapeado em páginas para a RAM/disco — isolamento e a ilusão de memória infinita.
\prov{sinónimos: memoria}
\prov{relações: usa ipc; relaciona recursao; relaciona rust; relaciona transcricao; relaciona o\_que\_e\_no\_projecto; relaciona paper\_0\_sintese; relaciona word; relaciona vontade\_de\_poder; parte\_de kernel; especializa gerenciamento\_de\_memoria}
\prov{proveniência: wiki \textperiodcentered{} fato \textperiodcentered{} confiança 1.0 \textperiodcentered{} domínio programacao \textperiodcentered{} história 14 versões no jornal}

%CRISTAL {"arestas": [{"alvo": "broca_os", "confianca": 1.0, "epistemico": "fato", "origem": "wiki", "rel": "parte_de"}, {"alvo": "computador_mineral", "confianca": 1.0, "epistemico": "fato", "origem": "wiki", "rel": "usa"}, {"alvo": "gato", "confianca": 1.0, "epistemico": "fato", "origem": "wiki", "rel": "usa"}, {"alvo": "beta_numeracao_metalica", "confianca": 1.0, "epistemico": "fato", "origem": "wiki", "rel": "usa"}, {"alvo": "registradores", "confianca": 1.0, "epistemico": "fato", "origem": "wiki", "rel": "relaciona"}], "confianca": 1.0, "contraexemplos": [], "descricao": "A máquina formalizada no metal (linguagem/metal.py). A CPU é a matriz de GATOS A_n (a ULA; executar é contrair, det=−1 conservado); a MEMÓRIA é o CRISTAL (células octoniônicas, estáveis por Pisot, norma=integridade); o BARRAMENTO é a RÉGUA MINERAL — N trilhas, a trilha n com o espectro do metal σ_n (os pesos σ_n^k), o endereço em β-numeração (pesos inteiros, ouro→Zeckendorf, reconstrução exata). Entre a CPU e a memória, o dado transita pela régua. Verif. 6/6 em metal.py.", "epistemico": "fato", "exemplos": [], "id": "metal_do_so", "memoria": "semantica", "meta": {"dominio": "sistema", "fonte": "metal do SO"}, "origem": "wiki", "palavras_chave": [], "sinonimos": [], "tipo": "conceito", "titulo": "O Metal do SO (CPU dos gatos, memória do cristal, barramento da régua)"}
\section{O Metal do SO (CPU dos gatos, memória do cristal, barramento da régua)}
A máquina formalizada no metal (linguagem/metal.py). A CPU é a matriz de GATOS A\_n (a ULA; executar é contrair, det=\ensuremath{-}1 conservado); a MEMÓRIA é o CRISTAL (células octoniônicas, estáveis por Pisot, norma=integridade); o BARRAMENTO é a RÉGUA MINERAL — N trilhas, a trilha n com o espectro do metal \ensuremath{\sigma}\_n (os pesos \ensuremath{\sigma}\_n\^{}k), o endereço em \ensuremath{\beta}-numeração (pesos inteiros, ouro\ensuremath{\to}Zeckendorf, reconstrução exata). Entre a CPU e a memória, o dado transita pela régua. Verif. 6/6 em metal.py.
\prov{relações: parte\_de broca\_os; usa computador\_mineral; usa gato; usa beta\_numeracao\_metalica; relaciona registradores}
\prov{proveniência: wiki \textperiodcentered{} metal do SO \textperiodcentered{} fato \textperiodcentered{} confiança 1.0 \textperiodcentered{} domínio sistema \textperiodcentered{} história 6 versões no jornal}

%CRISTAL {"arestas": [{"alvo": "programacao", "confianca": 1.0, "epistemico": "fato", "origem": "wiki", "rel": "parte_de"}, {"alvo": "lisp", "confianca": 1.0, "epistemico": "fato", "origem": "wiki", "rel": "usa"}], "confianca": 1.0, "contraexemplos": [], "descricao": "Código que gera ou manipula código — macros, reflexão, geração. O programa como dado (a herança de Lisp).", "epistemico": "fato", "exemplos": [], "id": "metaprogramacao", "memoria": "semantica", "meta": {"dominio": "programacao", "fonte": ""}, "origem": "wiki", "palavras_chave": [], "sinonimos": [], "tipo": "conceito", "titulo": "Metaprogramação"}
\section{Metaprogramação}
Código que gera ou manipula código — macros, reflexão, geração. O programa como dado (a herança de Lisp).
\prov{relações: parte\_de programacao; usa lisp}
\prov{proveniência: wiki \textperiodcentered{} fato \textperiodcentered{} confiança 1.0 \textperiodcentered{} domínio programacao \textperiodcentered{} história 6 versões no jornal}

%CRISTAL {"arestas": [{"alvo": "sistemas_distribuidos", "confianca": 1.0, "epistemico": "fato", "origem": "wiki", "rel": "usa"}, {"alvo": "api", "confianca": 1.0, "epistemico": "fato", "origem": "wiki", "rel": "usa"}], "confianca": 1.0, "contraexemplos": [], "descricao": "Quebrar o sistema em serviços pequenos e independentes que falam pela rede — escala e autonomia ao custo de complexidade distribuída.", "epistemico": "inferencia", "exemplos": [], "id": "microservicos", "memoria": "semantica", "meta": {"dominio": "programacao", "fonte": ""}, "origem": "wiki", "palavras_chave": [], "sinonimos": [], "tipo": "conceito", "titulo": "Microsserviços"}
\section{Microsserviços}
Quebrar o sistema em serviços pequenos e independentes que falam pela rede — escala e autonomia ao custo de complexidade distribuída.
\prov{relações: usa sistemas\_distribuidos; usa api}
\prov{proveniência: wiki \textperiodcentered{} inferencia \textperiodcentered{} confiança 1.0 \textperiodcentered{} domínio programacao \textperiodcentered{} história 6 versões no jornal}

%CRISTAL {"arestas": [{"alvo": "cristalchain", "confianca": 1.0, "epistemico": "fato", "origem": "wiki", "rel": "parte_de"}, {"alvo": "goldchain", "confianca": 1.0, "epistemico": "fato", "origem": "wiki", "rel": "relaciona"}, {"alvo": "broca-so_cristalizado", "confianca": 1.0, "epistemico": "fato", "origem": "wiki", "rel": "relaciona"}], "confianca": 1.0, "contraexemplos": [], "descricao": "O nome canônico da cadeia é CristalChain (o cristal é a classe; ouro é um metal). O diretório goldchain/ foi renomeado para cristalchain/ (git mv) e todos os CAMINHOS de arquivo migraram (código, sys.path, os.path.join, scripts/gex44/, o serviço systemd cristalchain/symbolic.py, deploy_arquivos.txt, .gitignore, Makefile). Preservados de propósito: o NÓ de grafo 'goldchain' (o metal ouro dentro da classe cristal — o hub cristalchain o especializa), as CHAVES cripto e identidades (CHAVE_OURO, alice-goldchain-2026, as strings de --trabalho — identificadores estáveis; renomear quebraria endereços/assinaturas da tesouraria viva) e o token goldchain_rede. O cutover do gex44 foi um mv (preservando o ledger vivo) + deploy + reinício do miner. Doc: MIGRACAO_CRISTALCHAIN.md.", "epistemico": "fato", "exemplos": [], "id": "migracao_nome_cristalchain", "memoria": "semantica", "meta": {"dominio": "goldchain", "fonte": "migração de nome"}, "origem": "wiki", "palavras_chave": [], "sinonimos": [], "tipo": "conceito", "titulo": "Migração de Nome: goldchain → cristalchain (renome completo)"}
\section{Migração de Nome: goldchain \ensuremath{\to} cristalchain (renome completo)}
O nome canônico da cadeia é CristalChain (o cristal é a classe; ouro é um metal). O diretório goldchain/ foi renomeado para cristalchain/ (git mv) e todos os CAMINHOS de arquivo migraram (código, sys.path, os.path.join, scripts/gex44/, o serviço systemd cristalchain/symbolic.py, deploy\_arquivos.txt, .gitignore, Makefile). Preservados de propósito: o NÓ de grafo 'goldchain' (o metal ouro dentro da classe cristal — o hub cristalchain o especializa), as CHAVES cripto e identidades (CHAVE\_OURO, alice-goldchain-2026, as strings de --trabalho — identificadores estáveis; renomear quebraria endereços/assinaturas da tesouraria viva) e o token goldchain\_rede. O cutover do gex44 foi um mv (preservando o ledger vivo) + deploy + reinício do miner. Doc: MIGRACAO\_CRISTALCHAIN.md.
\prov{relações: parte\_de cristalchain; relaciona goldchain; relaciona broca-so\_cristalizado}
\prov{proveniência: wiki \textperiodcentered{} migração de nome \textperiodcentered{} fato \textperiodcentered{} confiança 1.0 \textperiodcentered{} domínio goldchain \textperiodcentered{} história 12 versões no jornal}

%CRISTAL {"arestas": [{"alvo": "cristalchain", "confianca": 1.0, "epistemico": "fato", "origem": "wiki", "rel": "parte_de"}, {"alvo": "pow_dificuldade_calor", "confianca": 1.0, "epistemico": "fato", "origem": "wiki", "rel": "explica"}, {"alvo": "parte_fria_t_zero", "confianca": 1.0, "epistemico": "fato", "origem": "wiki", "rel": "usa"}, {"alvo": "pow_metal", "confianca": 1.0, "epistemico": "fato", "origem": "wiki", "rel": "relaciona"}, {"alvo": "difusao_preenche_mineracao", "confianca": 1.0, "epistemico": "fato", "origem": "wiki", "rel": "relaciona"}], "confianca": 1.0, "contraexemplos": [], "descricao": "A prova de trabalho busca um nonce cujo hash caia ABAIXO DO ALVO — a região cristal, contraída e rara. A dificuldade é a TEMPERATURA INVERSA: mais bits-zero = mais frio = mais raro (custo ~2^k). Encontrar o nonce é a descida em tempo imaginário (e^{−βH}, β→∞) ao estado fundamental — a mesma da mecânica quântica mineral. Minerar não é corrida arbitrária: é CRISTALIZAÇÃO. Verif.: 16 bits ~2^16 tentativas.", "epistemico": "fato", "exemplos": [], "id": "minerar_e_resfriar", "memoria": "semantica", "meta": {"dominio": "goldchain", "fonte": "cristalchain"}, "origem": "wiki", "palavras_chave": [], "sinonimos": [], "tipo": "conceito", "titulo": "Minerar é Resfriar até o Cristal"}
\section{Minerar é Resfriar até o Cristal}
A prova de trabalho busca um nonce cujo hash caia ABAIXO DO ALVO — a região cristal, contraída e rara. A dificuldade é a TEMPERATURA INVERSA: mais bits-zero = mais frio = mais raro (custo \~{}2\^{}k). Encontrar o nonce é a descida em tempo imaginário (e\^{}\{\ensuremath{-}\ensuremath{\beta}H\}, \ensuremath{\beta}\ensuremath{\to}\ensuremath{\infty}) ao estado fundamental — a mesma da mecânica quântica mineral. Minerar não é corrida arbitrária: é CRISTALIZAÇÃO. Verif.: 16 bits \~{}2\^{}16 tentativas.
\prov{relações: parte\_de cristalchain; explica pow\_dificuldade\_calor; usa parte\_fria\_t\_zero; relaciona pow\_metal; relaciona difusao\_preenche\_mineracao}
\prov{proveniência: wiki \textperiodcentered{} cristalchain \textperiodcentered{} fato \textperiodcentered{} confiança 1.0 \textperiodcentered{} domínio goldchain \textperiodcentered{} história 6 versões no jornal}

%CRISTAL {"arestas": [{"alvo": "seguranca_da_informacao", "confianca": 1.0, "epistemico": "fato", "origem": "wiki", "rel": "usa"}], "confianca": 1.0, "contraexemplos": [], "descricao": "Mapear sistematicamente quem pode atacar, como e o quê — pensar como o adversário antes de construir a defesa.", "epistemico": "inferencia", "exemplos": [], "id": "modelo_de_ameacas", "memoria": "semantica", "meta": {"dominio": "programacao", "fonte": ""}, "origem": "wiki", "palavras_chave": [], "sinonimos": [], "tipo": "conceito", "titulo": "Modelagem de Ameaças"}
\section{Modelagem de Ameaças}
Mapear sistematicamente quem pode atacar, como e o quê — pensar como o adversário antes de construir a defesa.
\prov{relações: usa seguranca\_da\_informacao}
\prov{proveniência: wiki \textperiodcentered{} inferencia \textperiodcentered{} confiança 1.0 \textperiodcentered{} domínio programacao \textperiodcentered{} história 2 versões no jornal}

%CRISTAL {"arestas": [{"alvo": "programacao_concorrente", "confianca": 1.0, "epistemico": "fato", "origem": "wiki", "rel": "especializa"}, {"alvo": "erlang", "confianca": 1.0, "epistemico": "fato", "origem": "wiki", "rel": "explica"}], "confianca": 1.0, "contraexemplos": [], "descricao": "Processos isolados que só se comunicam por mensagens, sem memória compartilhada — a receita de Erlang para robustez.", "epistemico": "fato", "exemplos": [], "id": "modelo_de_atores", "memoria": "semantica", "meta": {"autor": "Hewitt", "dominio": "programacao", "fonte": ""}, "origem": "wiki", "palavras_chave": [], "sinonimos": [], "tipo": "conceito", "titulo": "Modelo de Atores"}
\section{Modelo de Atores}
Processos isolados que só se comunicam por mensagens, sem memória compartilhada — a receita de Erlang para robustez.
\prov{relações: especializa programacao\_concorrente; explica erlang}
\prov{proveniência: wiki \textperiodcentered{} fato \textperiodcentered{} confiança 1.0 \textperiodcentered{} domínio programacao \textperiodcentered{} história 6 versões no jornal}

%CRISTAL {"arestas": [{"alvo": "transformer", "confianca": 1.0, "epistemico": "fato", "origem": "wiki", "rel": "usa"}, {"alvo": "tiffany", "confianca": 1.0, "epistemico": "fato", "origem": "wiki", "rel": "explica"}, {"alvo": "mente_computacao", "confianca": 1.0, "epistemico": "fato", "origem": "wiki", "rel": "relaciona"}], "confianca": 1.0, "contraexemplos": [], "descricao": "Um modelo treinado a prever o próximo token que, em escala, gera texto coerente e raciocina — o que a Tiffany imita em pequeno, sem fabricar.", "epistemico": "inferencia", "exemplos": [], "id": "modelo_de_linguagem", "memoria": "semantica", "meta": {"dominio": "programacao", "fonte": ""}, "origem": "wiki", "palavras_chave": [], "sinonimos": [], "tipo": "conceito", "titulo": "Modelo de Linguagem (LLM)"}
\section{Modelo de Linguagem (LLM)}
Um modelo treinado a prever o próximo token que, em escala, gera texto coerente e raciocina — o que a Tiffany imita em pequeno, sem fabricar.
\prov{relações: usa transformer; explica tiffany; relaciona mente\_computacao}
\prov{proveniência: wiki \textperiodcentered{} inferencia \textperiodcentered{} confiança 1.0 \textperiodcentered{} domínio programacao \textperiodcentered{} história 4 versões no jornal}

%CRISTAL {"arestas": [{"alvo": "banco_de_dados", "confianca": 1.0, "epistemico": "fato", "origem": "wiki", "rel": "parte_de"}, {"alvo": "algebra_abstrata", "confianca": 1.0, "epistemico": "fato", "origem": "wiki", "rel": "usa"}], "confianca": 1.0, "contraexemplos": [], "descricao": "Codd: dados como tabelas (relações) manipuladas por álgebra relacional — a base matemática do banco de dados moderno.", "epistemico": "fato", "exemplos": [], "id": "modelo_relacional", "memoria": "semantica", "meta": {"autor": "Codd", "dominio": "programacao", "fonte": ""}, "origem": "wiki", "palavras_chave": [], "sinonimos": [], "tipo": "conceito", "titulo": "Modelo Relacional"}
\section{Modelo Relacional}
Codd: dados como tabelas (relações) manipuladas por álgebra relacional — a base matemática do banco de dados moderno.
\prov{relações: parte\_de banco\_de\_dados; usa algebra\_abstrata}
\prov{proveniência: wiki \textperiodcentered{} fato \textperiodcentered{} confiança 1.0 \textperiodcentered{} domínio programacao \textperiodcentered{} história 6 versões no jornal}

%CRISTAL {"arestas": [{"alvo": "kernel", "confianca": 1.0, "epistemico": "fato", "origem": "wiki", "rel": "usa"}], "confianca": 1.0, "contraexemplos": [], "descricao": "O hardware separa código privilegiado (kernel) do comum (usuário); só via syscall se cruza a linha — a muralha que protege a máquina.", "epistemico": "fato", "exemplos": [], "id": "modo_usuario_kernel", "memoria": "semantica", "meta": {"dominio": "programacao", "fonte": ""}, "origem": "wiki", "palavras_chave": [], "sinonimos": [], "tipo": "conceito", "titulo": "Modo Usuário e Modo Kernel"}
\section{Modo Usuário e Modo Kernel}
O hardware separa código privilegiado (kernel) do comum (usuário); só via syscall se cruza a linha — a muralha que protege a máquina.
\prov{relações: usa kernel}
\prov{proveniência: wiki \textperiodcentered{} fato \textperiodcentered{} confiança 1.0 \textperiodcentered{} domínio programacao \textperiodcentered{} história 2 versões no jornal}

%CRISTAL {"arestas": [{"alvo": "arquitetura_de_software", "confianca": 1.0, "epistemico": "fato", "origem": "wiki", "rel": "usa"}, {"alvo": "abstracao_software", "confianca": 1.0, "epistemico": "fato", "origem": "wiki", "rel": "usa"}], "confianca": 1.0, "contraexemplos": [], "descricao": "Compor o sistema de partes independentes e substituíveis, com interfaces claras — dividir para dominar a complexidade. O broca-so em módulos.", "epistemico": "fato", "exemplos": [], "id": "modularidade", "memoria": "semantica", "meta": {"dominio": "programacao", "fonte": ""}, "origem": "wiki", "palavras_chave": [], "sinonimos": [], "tipo": "conceito", "titulo": "Modularidade"}
\section{Modularidade}
Compor o sistema de partes independentes e substituíveis, com interfaces claras — dividir para dominar a complexidade. O broca-so em módulos.
\prov{relações: usa arquitetura\_de\_software; usa abstracao\_software}
\prov{proveniência: wiki \textperiodcentered{} fato \textperiodcentered{} confiança 1.0 \textperiodcentered{} domínio programacao \textperiodcentered{} história 3 versões no jornal}

%CRISTAL {"arestas": [{"alvo": "programacao_funcional", "confianca": 1.0, "epistemico": "fato", "origem": "wiki", "rel": "usa"}, {"alvo": "haskell", "confianca": 1.0, "epistemico": "fato", "origem": "wiki", "rel": "parte_de"}], "confianca": 1.0, "contraexemplos": [], "descricao": "Um padrão para encadear computações com contexto (efeito, falha, estado) mantendo a pureza — o burro de carga de Haskell.", "epistemico": "fato", "exemplos": [], "id": "monada", "memoria": "semantica", "meta": {"dominio": "programacao", "fonte": ""}, "origem": "wiki", "palavras_chave": [], "sinonimos": [], "tipo": "conceito", "titulo": "Mônada"}
\section{Mônada}
Um padrão para encadear computações com contexto (efeito, falha, estado) mantendo a pureza — o burro de carga de Haskell.
\prov{relações: usa programacao\_funcional; parte\_de haskell}
\prov{proveniência: wiki \textperiodcentered{} fato \textperiodcentered{} confiança 1.0 \textperiodcentered{} domínio programacao \textperiodcentered{} história 6 versões no jornal}

%CRISTAL {"arestas": [{"alvo": "padrao_de_projeto", "confianca": 1.0, "epistemico": "fato", "origem": "wiki", "rel": "especializa"}, {"alvo": "separacao_de_interesses", "confianca": 1.0, "epistemico": "fato", "origem": "wiki", "rel": "usa"}], "confianca": 1.0, "contraexemplos": [], "descricao": "Separar dados (Model), apresentação (View) e coordenação (Controller) — o padrão que estrutura quase toda interface.", "epistemico": "inferencia", "exemplos": [], "id": "mvc", "memoria": "semantica", "meta": {"dominio": "programacao", "fonte": ""}, "origem": "wiki", "palavras_chave": [], "sinonimos": [], "tipo": "conceito", "titulo": "MVC (Model-View-Controller)"}
\section{MVC (Model-View-Controller)}
Separar dados (Model), apresentação (View) e coordenação (Controller) — o padrão que estrutura quase toda interface.
\prov{relações: especializa padrao\_de\_projeto; usa separacao\_de\_interesses}
\prov{proveniência: wiki \textperiodcentered{} inferencia \textperiodcentered{} confiança 1.0 \textperiodcentered{} domínio programacao \textperiodcentered{} história 3 versões no jornal}

%CRISTAL {"arestas": [{"alvo": "html", "confianca": 1.0, "epistemico": "fato", "origem": "wiki", "rel": "usa"}, {"alvo": "javascript", "confianca": 1.0, "epistemico": "fato", "origem": "wiki", "rel": "usa"}, {"alvo": "internet", "confianca": 1.0, "epistemico": "fato", "origem": "wiki", "rel": "usa"}], "confianca": 1.0, "contraexemplos": [], "descricao": "O programa que busca (HTTP), interpreta (HTML/CSS/JS) e renderiza a web — uma máquina virtual de aplicações disfarçada de leitor de documentos.", "epistemico": "fato", "exemplos": [], "id": "navegador", "memoria": "semantica", "meta": {"dominio": "programacao", "fonte": ""}, "origem": "wiki", "palavras_chave": [], "sinonimos": [], "tipo": "conceito", "titulo": "Navegador"}
\section{Navegador}
O programa que busca (HTTP), interpreta (HTML/CSS/JS) e renderiza a web — uma máquina virtual de aplicações disfarçada de leitor de documentos.
\prov{relações: usa html; usa javascript; usa internet}
\prov{proveniência: wiki \textperiodcentered{} fato \textperiodcentered{} confiança 1.0 \textperiodcentered{} domínio programacao \textperiodcentered{} história 4 versões no jornal}

%CRISTAL {"arestas": [{"alvo": "framework_backend", "confianca": 1.0, "epistemico": "fato", "origem": "wiki", "rel": "especializa"}, {"alvo": "typescript", "confianca": 1.0, "epistemico": "fato", "origem": "wiki", "rel": "usa"}, {"alvo": "injecao_de_dependencia", "confianca": 1.0, "epistemico": "fato", "origem": "wiki", "rel": "usa"}, {"alvo": "programacao_orientada_objetos", "confianca": 1.0, "epistemico": "fato", "origem": "wiki", "rel": "usa"}, {"alvo": "angular", "confianca": 1.0, "epistemico": "fato", "origem": "wiki", "rel": "relaciona"}], "confianca": 1.0, "contraexemplos": [], "descricao": "Framework backend em TypeScript inspirado no Angular — módulos, injeção de dependência e OOP estruturada sobre o Express. Arquitetura para escalar.", "epistemico": "fato", "exemplos": [], "id": "nestjs", "memoria": "semantica", "meta": {"autor": "Myśliwiec", "dominio": "programacao", "fonte": ""}, "origem": "wiki", "palavras_chave": [], "sinonimos": [], "tipo": "conceito", "titulo": "NestJS"}
\section{NestJS}
Framework backend em TypeScript inspirado no Angular — módulos, injeção de dependência e OOP estruturada sobre o Express. Arquitetura para escalar.
\prov{relações: especializa framework\_backend; usa typescript; usa injecao\_de\_dependencia; usa programacao\_orientada\_objetos; relaciona angular}
\prov{proveniência: wiki \textperiodcentered{} fato \textperiodcentered{} confiança 1.0 \textperiodcentered{} domínio programacao \textperiodcentered{} história 12 versões no jornal}

%CRISTAL {"arestas": [{"alvo": "operacao_broca_so", "confianca": 1.0, "epistemico": "fato", "origem": "wiki", "rel": "parte_de"}, {"alvo": "o_substrato_a_maquina_e_funcao_de_onda", "confianca": 1.0, "epistemico": "fato", "origem": "wiki", "rel": "relaciona"}], "confianca": 1.0, "contraexemplos": [], "descricao": "code/neuronio.py:gato_arquivo lê o arquivo byte a byte e conta bits nas máscaras &0xAA (pares, e) e &0x55 (ímpares, o) via popcnt, devolvendo o word [e+o, e]. Os dados ficam no disco; só a conta vai à RAM — nunca carrega o conteúdo. É a lição do 'tirar a mão' executável.", "epistemico": "fato", "exemplos": [], "id": "neuronio_stream_bit", "memoria": "semantica", "meta": {"autor": "broca-so", "dominio": "operacao_so", "fonte": "code/neuronio.py; ula/neuron_io.c"}, "origem": "wiki", "palavras_chave": [], "sinonimos": [], "tipo": "conceito", "titulo": "Neurônio por Streaming de Bits (o gato aplicado)"}
\section{Neurônio por Streaming de Bits (o gato aplicado)}
code/neuronio.py:gato\_arquivo lê o arquivo byte a byte e conta bits nas máscaras \&0xAA (pares, e) e \&0x55 (ímpares, o) via popcnt, devolvendo o word [e+o, e]. Os dados ficam no disco; só a conta vai à RAM — nunca carrega o conteúdo. É a lição do 'tirar a mão' executável.
\prov{relações: parte\_de operacao\_broca\_so; relaciona o\_substrato\_a\_maquina\_e\_funcao\_de\_onda}
\prov{proveniência: wiki \textperiodcentered{} code/neuronio.py; ula/neuron\_io.c \textperiodcentered{} fato \textperiodcentered{} confiança 1.0 \textperiodcentered{} domínio operacao\_so \textperiodcentered{} história 3 versões no jornal}

%CRISTAL {"arestas": [{"alvo": "javascript", "confianca": 1.0, "epistemico": "fato", "origem": "wiki", "rel": "usa"}, {"alvo": "laco_de_eventos", "confianca": 1.0, "epistemico": "fato", "origem": "wiki", "rel": "usa"}, {"alvo": "programacao_de_sistemas", "confianca": 1.0, "epistemico": "fato", "origem": "wiki", "rel": "relaciona"}], "confianca": 1.0, "contraexemplos": [], "descricao": "Dahl: o runtime que levou o JavaScript ao servidor, sobre o motor V8 — I/O não-bloqueante por laço de eventos, uma thread só.", "epistemico": "fato", "exemplos": [], "id": "nodejs", "memoria": "semantica", "meta": {"autor": "Dahl", "dominio": "programacao", "fonte": ""}, "origem": "wiki", "palavras_chave": [], "sinonimos": [], "tipo": "conceito", "titulo": "Node.js"}
\section{Node.js}
Dahl: o runtime que levou o JavaScript ao servidor, sobre o motor V8 — I/O não-bloqueante por laço de eventos, uma thread só.
\prov{relações: usa javascript; usa laco\_de\_eventos; relaciona programacao\_de\_sistemas}
\prov{proveniência: wiki \textperiodcentered{} fato \textperiodcentered{} confiança 1.0 \textperiodcentered{} domínio programacao \textperiodcentered{} história 8 versões no jornal}

%CRISTAL {"arestas": [{"alvo": "pipeline_de_ingestao", "confianca": 1.0, "epistemico": "fato", "origem": "paper", "rel": "parte_de"}, {"alvo": "modelo_relacional", "confianca": 1.0, "epistemico": "fato", "origem": "wiki", "rel": "usa"}], "confianca": 1.0, "contraexemplos": [], "descricao": "Organizar as tabelas para eliminar redundância e anomalias — cada fato num só lugar; o rigor relacional contra a duplicação.", "epistemico": "inferencia", "exemplos": [], "id": "normalizacao", "memoria": "semantica", "meta": {"dominio": "programacao", "fonte": ""}, "origem": "wiki", "palavras_chave": [], "sinonimos": [], "tipo": "conceito", "titulo": "Normalização"}
\section{Normalização}
Organizar as tabelas para eliminar redundância e anomalias — cada fato num só lugar; o rigor relacional contra a duplicação.
\prov{relações: parte\_de pipeline\_de\_ingestao; usa modelo\_relacional}
\prov{proveniência: wiki \textperiodcentered{} inferencia \textperiodcentered{} confiança 1.0 \textperiodcentered{} domínio programacao \textperiodcentered{} história 6 versões no jornal}

%CRISTAL {"arestas": [{"alvo": "banco_de_dados", "confianca": 1.0, "epistemico": "fato", "origem": "wiki", "rel": "especializa"}, {"alvo": "teorema_cap", "confianca": 1.0, "epistemico": "fato", "origem": "wiki", "rel": "usa"}], "confianca": 1.0, "contraexemplos": [], "descricao": "Bancos que abrem mão do relacional por escala e flexibilidade — chave-valor, documento, coluna, grafo. Cada um paga um preço no CAP.", "epistemico": "inferencia", "exemplos": [], "id": "nosql", "memoria": "semantica", "meta": {"dominio": "programacao", "fonte": ""}, "origem": "wiki", "palavras_chave": [], "sinonimos": [], "tipo": "conceito", "titulo": "NoSQL"}
\section{NoSQL}
Bancos que abrem mão do relacional por escala e flexibilidade — chave-valor, documento, coluna, grafo. Cada um paga um preço no CAP.
\prov{relações: especializa banco\_de\_dados; usa teorema\_cap}
\prov{proveniência: wiki \textperiodcentered{} inferencia \textperiodcentered{} confiança 1.0 \textperiodcentered{} domínio programacao \textperiodcentered{} história 6 versões no jornal}

%CRISTAL {"arestas": [{"alvo": "linguagem_matematica", "confianca": 1.0, "epistemico": "fato", "origem": "wiki", "rel": "parte_de"}, {"alvo": "codigo_semiotico", "confianca": 1.0, "epistemico": "fato", "origem": "wiki", "rel": "usa"}], "confianca": 1.0, "contraexemplos": [], "descricao": "O sistema de símbolos (∑ ∫ ∀ ∈ √ ≤ α) que condensa relações e operações — uma escrita densa e universal, independente do idioma. A forma sobre o sentido.", "epistemico": "fato", "exemplos": [], "id": "notacao_matematica", "memoria": "semantica", "meta": {"dominio": "programacao", "fonte": ""}, "origem": "wiki", "palavras_chave": [], "sinonimos": [], "tipo": "conceito", "titulo": "Notação Matemática"}
\section{Notação Matemática}
O sistema de símbolos (\ensuremath{\sum} \ensuremath{\int} \ensuremath{\forall} \ensuremath{\in} \ensuremath{\sqrt{\,}} \ensuremath{\le} \ensuremath{\alpha}) que condensa relações e operações — uma escrita densa e universal, independente do idioma. A forma sobre o sentido.
\prov{relações: parte\_de linguagem\_matematica; usa codigo\_semiotico}
\prov{proveniência: wiki \textperiodcentered{} fato \textperiodcentered{} confiança 1.0 \textperiodcentered{} domínio programacao \textperiodcentered{} história 3 versões no jornal}

%CRISTAL {"arestas": [{"alvo": "ciclo_uno", "confianca": 1.0, "epistemico": "fato", "origem": "wiki", "rel": "usa"}, {"alvo": "a_bolha_infraestrutura", "confianca": 1.0, "epistemico": "fato", "origem": "wiki", "rel": "relaciona"}], "confianca": 1.0, "contraexemplos": [], "descricao": "tick_paralelo coleta as contas 'prontas' e dispara o ciclo uno de cada uma num ThreadPoolExecutor (um worker por ramo pronto); --loop vira daemon com PID file, espera adaptativa e tick idle quando ninguém está pronto. O runtime de produção é 'nucleo --loop'.", "epistemico": "fato", "exemplos": [], "id": "nucleo_paralelo", "memoria": "semantica", "meta": {"autor": "broca-so", "dominio": "operacao_so", "fonte": "goldchain/nucleo_minera.py"}, "origem": "wiki", "palavras_chave": [], "sinonimos": [], "tipo": "conceito", "titulo": "Núcleo — Tesourarias em Paralelo"}
\section{Núcleo — Tesourarias em Paralelo}
tick\_paralelo coleta as contas 'prontas' e dispara o ciclo uno de cada uma num ThreadPoolExecutor (um worker por ramo pronto); --loop vira daemon com PID file, espera adaptativa e tick idle quando ninguém está pronto. O runtime de produção é 'nucleo --loop'.
\prov{relações: usa ciclo\_uno; relaciona a\_bolha\_infraestrutura}
\prov{proveniência: wiki \textperiodcentered{} goldchain/nucleo\_minera.py \textperiodcentered{} fato \textperiodcentered{} confiança 1.0 \textperiodcentered{} domínio operacao\_so \textperiodcentered{} história 3 versões no jornal}

%CRISTAL {"arestas": [{"alvo": "gpu", "confianca": 1.0, "epistemico": "fato", "origem": "curriculo", "rel": "usa"}, {"alvo": "kernel", "confianca": 0.742, "epistemico": "fato", "origem": "manual", "rel": "referencia"}, {"alvo": "rust", "confianca": 0.5, "epistemico": "fato", "origem": "wiki", "rel": "relaciona"}, {"alvo": "topologia", "confianca": 0.5, "epistemico": "fato", "origem": "wiki", "rel": "relaciona"}, {"alvo": "parabola_pa", "confianca": 0.5, "epistemico": "fato", "origem": "wiki", "rel": "relaciona"}, {"alvo": "processos", "confianca": 0.5, "epistemico": "fato", "origem": "wiki", "rel": "relaciona"}, {"alvo": "transcricao", "confianca": 0.5, "epistemico": "fato", "origem": "wiki", "rel": "relaciona"}, {"alvo": "teste_ordem", "confianca": 0.5, "epistemico": "fato", "origem": "wiki", "rel": "relaciona"}, {"alvo": "pow_metal_setor", "confianca": 0.5, "epistemico": "fato", "origem": "wiki", "rel": "relaciona"}, {"alvo": "por_sessao_conversadadostelemetriasessaojsonl", "confianca": 0.5, "epistemico": "fato", "origem": "wiki", "rel": "relaciona"}, {"alvo": "recursao", "confianca": 0.5, "epistemico": "fato", "origem": "wiki", "rel": "relaciona"}, {"alvo": "pow_metal_funil", "confianca": 0.5, "epistemico": "fato", "origem": "wiki", "rel": "relaciona"}, {"alvo": "parabola_destruir", "confianca": 0.5, "epistemico": "fato", "origem": "wiki", "rel": "relaciona"}, {"alvo": "pilha", "confianca": 0.5, "epistemico": "fato", "origem": "wiki", "rel": "relaciona"}, {"alvo": "ponteiro", "confianca": 0.5, "epistemico": "fato", "origem": "wiki", "rel": "relaciona"}, {"alvo": "testes", "confianca": 0.5, "epistemico": "fato", "origem": "wiki", "rel": "relaciona"}, {"alvo": "projeto_cristaljson", "confianca": 0.5, "epistemico": "fato", "origem": "wiki", "rel": "relaciona"}, {"alvo": "regua_associativa_e_o_risco_de_vazamento", "confianca": 0.5, "epistemico": "fato", "origem": "wiki", "rel": "relaciona"}, {"alvo": "opcao_a_tarball_release", "confianca": 0.5, "epistemico": "fato", "origem": "wiki", "rel": "relaciona"}, {"alvo": "resposta_a_pergunta_dificil", "confianca": 0.5, "epistemico": "fato", "origem": "wiki", "rel": "relaciona"}, {"alvo": "tiffany_cabeca_broca", "confianca": 0.5, "epistemico": "fato", "origem": "wiki", "rel": "relaciona"}], "confianca": 1.0, "contraexemplos": [], "descricao": "Memória não-uniforme; afinidade.", "epistemico": "fato", "exemplos": [], "id": "numa", "memoria": "semantica", "meta": {"dominio": "sistemas", "nivel": 6, "ordem": 7}, "origem": "curriculo", "palavras_chave": [], "sinonimos": [], "tipo": "conceito", "titulo": "NUMA"}
\section{NUMA}
Memória não-uniforme; afinidade.
\prov{relações: usa gpu; referencia kernel; relaciona rust; relaciona topologia; relaciona parabola\_pa; relaciona processos; relaciona transcricao; relaciona teste\_ordem; relaciona pow\_metal\_setor; relaciona por\_sessao\_conversadadostelemetriasessaojsonl; relaciona recursao; relaciona pow\_metal\_funil; relaciona parabola\_destruir; relaciona pilha; relaciona ponteiro; relaciona testes; relaciona projeto\_cristaljson; relaciona regua\_associativa\_e\_o\_risco\_de\_vazamento; relaciona opcao\_a\_tarball\_release; relaciona resposta\_a\_pergunta\_dificil; relaciona tiffany\_cabeca\_broca}
\prov{proveniência: curriculo \textperiodcentered{} fato \textperiodcentered{} confiança 1.0 \textperiodcentered{} domínio sistemas \textperiodcentered{} história 26 versões no jornal}

%CRISTAL {"arestas": [{"alvo": "devops", "confianca": 1.0, "epistemico": "fato", "origem": "wiki", "rel": "parte_de"}], "confianca": 1.0, "contraexemplos": [], "descricao": "Enxergar o que um sistema faz por dentro via logs, métricas e traces — para diagnosticar o que não se previu. Ver antes de consertar.", "epistemico": "inferencia", "exemplos": [], "id": "observabilidade", "memoria": "semantica", "meta": {"dominio": "programacao", "fonte": ""}, "origem": "wiki", "palavras_chave": [], "sinonimos": [], "tipo": "conceito", "titulo": "Observabilidade"}
\section{Observabilidade}
Enxergar o que um sistema faz por dentro via logs, métricas e traces — para diagnosticar o que não se previu. Ver antes de consertar.
\prov{relações: parte\_de devops}
\prov{proveniência: wiki \textperiodcentered{} inferencia \textperiodcentered{} confiança 1.0 \textperiodcentered{} domínio programacao \textperiodcentered{} história 4 versões no jornal}

%CRISTAL {"arestas": [{"alvo": "programacao_funcional", "confianca": 1.0, "epistemico": "fato", "origem": "wiki", "rel": "usa"}, {"alvo": "tipos_algebricos", "confianca": 1.0, "epistemico": "fato", "origem": "wiki", "rel": "usa"}], "confianca": 1.0, "contraexemplos": [], "descricao": "Funcional com tipos algébricos, inferência e pitada de OOP — pragmático e veloz; usado onde correção importa (finanças, compiladores).", "epistemico": "fato", "exemplos": [], "id": "ocaml", "memoria": "semantica", "meta": {"dominio": "programacao", "fonte": ""}, "origem": "wiki", "palavras_chave": [], "sinonimos": [], "tipo": "conceito", "titulo": "OCaml"}
\section{OCaml}
Funcional com tipos algébricos, inferência e pitada de OOP — pragmático e veloz; usado onde correção importa (finanças, compiladores).
\prov{relações: usa programacao\_funcional; usa tipos\_algebricos}
\prov{proveniência: wiki \textperiodcentered{} fato \textperiodcentered{} confiança 1.0 \textperiodcentered{} domínio programacao \textperiodcentered{} história 6 versões no jornal}

%CRISTAL {"arestas": [{"alvo": "broca_os", "confianca": 1.0, "epistemico": "fato", "origem": "wiki", "rel": "explica"}, {"alvo": "tiffany", "confianca": 1.0, "epistemico": "fato", "origem": "wiki", "rel": "relaciona"}], "confianca": 1.0, "contraexemplos": [], "descricao": "Como o sistema realmente roda (o código executável, não a teoria): a máquina (ISA/ULA/neurônio), a plataforma multi-tenant, a goldchain, o órgão de conhecimento e a autorização (BAI) — o SO documentando o próprio corpo.", "epistemico": "fato", "exemplos": [], "id": "operacao_broca_so", "memoria": "semantica", "meta": {"autor": "broca-so", "dominio": "operacao_so", "fonte": ""}, "origem": "wiki", "palavras_chave": [], "sinonimos": [], "tipo": "conceito", "titulo": "A Operação do Broca-SO"}
\section{A Operação do Broca-SO}
Como o sistema realmente roda (o código executável, não a teoria): a máquina (ISA/ULA/neurônio), a plataforma multi-tenant, a goldchain, o órgão de conhecimento e a autorização (BAI) — o SO documentando o próprio corpo.
\prov{relações: explica broca\_os; relaciona tiffany}
\prov{proveniência: wiki \textperiodcentered{} fato \textperiodcentered{} confiança 1.0 \textperiodcentered{} domínio operacao\_so \textperiodcentered{} história 3 versões no jornal}

%CRISTAL {"arestas": [{"alvo": "metal_do_so", "confianca": 1.0, "epistemico": "fato", "origem": "wiki", "rel": "parte_de"}, {"alvo": "vocabulario_canonico", "confianca": 1.0, "epistemico": "fato", "origem": "wiki", "rel": "usa"}, {"alvo": "corpo_estelar", "confianca": 1.0, "epistemico": "fato", "origem": "wiki", "rel": "usa"}, {"alvo": "crivo_de_selberg", "confianca": 1.0, "epistemico": "fato", "origem": "wiki", "rel": "usa"}], "confianca": 1.0, "contraexemplos": [], "descricao": "Os termos do vocabulário viram FUNÇÕES RESERVADAS do sistema, implementadas no metal (compostas dos módulos verificados): gato (a CPU), soma (⊕), produto (⊗/La Hire), norma (=det=−1), flip (log/Wick), entropia (=log σ), crivo (Selberg G(z)), classe (a parábola/Anosov), contrair (Einstein), endereco (β). chamar(nome,…) executa a operação. São nomes reservados: a mesma operação, um nome — e computam de verdade (norma=−1, √2⊗√3=√6, entropia=log σ, o gato A_2·(1,0)=(2,1)).", "epistemico": "fato", "exemplos": [], "id": "operacoes_reservadas_no_metal", "memoria": "semantica", "meta": {"dominio": "sistema", "fonte": "metal do SO"}, "origem": "wiki", "palavras_chave": [], "sinonimos": [], "tipo": "conceito", "titulo": "Os Termos são Operações Reservadas no Metal"}
\section{Os Termos são Operações Reservadas no Metal}
Os termos do vocabulário viram FUNÇÕES RESERVADAS do sistema, implementadas no metal (compostas dos módulos verificados): gato (a CPU), soma (\ensuremath{\oplus}), produto (\ensuremath{\otimes}/La Hire), norma (=det=\ensuremath{-}1), flip (log/Wick), entropia (=log \ensuremath{\sigma}), crivo (Selberg G(z)), classe (a parábola/Anosov), contrair (Einstein), endereco (\ensuremath{\beta}). chamar(nome,…) executa a operação. São nomes reservados: a mesma operação, um nome — e computam de verdade (norma=\ensuremath{-}1, \ensuremath{\sqrt{\,}}2\ensuremath{\otimes}\ensuremath{\sqrt{\,}}3=\ensuremath{\sqrt{\,}}6, entropia=log \ensuremath{\sigma}, o gato A\_2\textperiodcentered (1,0)=(2,1)).
\prov{relações: parte\_de metal\_do\_so; usa vocabulario\_canonico; usa corpo\_estelar; usa crivo\_de\_selberg}
\prov{proveniência: wiki \textperiodcentered{} metal do SO \textperiodcentered{} fato \textperiodcentered{} confiança 1.0 \textperiodcentered{} domínio sistema \textperiodcentered{} história 5 versões no jornal}

%CRISTAL {"arestas": [{"alvo": "dsl_comparacoes_ordem", "confianca": 1.0, "epistemico": "fato", "origem": "wiki", "rel": "usa"}, {"alvo": "expressividade_sub_turing", "confianca": 1.0, "epistemico": "fato", "origem": "wiki", "rel": "explica"}, {"alvo": "parabola_algebra_conversacional", "confianca": 1.0, "epistemico": "fato", "origem": "wiki", "rel": "usa"}, {"alvo": "bai_orbita_reta_metalica", "confianca": 1.0, "epistemico": "fato", "origem": "wiki", "rel": "relaciona"}, {"alvo": "sandbox_cristal_ao_so", "confianca": 1.0, "epistemico": "fato", "origem": "wiki", "rel": "usa"}, {"alvo": "problema_da_parada", "confianca": 1.0, "epistemico": "fato", "origem": "wiki", "rel": "relaciona"}, {"alvo": "tiffany_roda_codigo", "confianca": 1.0, "epistemico": "fato", "origem": "wiki", "rel": "usa"}, {"alvo": "conjugados_galois_parabola", "confianca": 1.0, "epistemico": "fato", "origem": "wiki", "rel": "relaciona"}, {"alvo": "transformada_metalica", "confianca": 1.0, "epistemico": "fato", "origem": "wiki", "rel": "relaciona"}], "confianca": 1.0, "contraexemplos": [], "descricao": "A lição do Aarão sobre as expressões ricas que não terminam: um laço não-convergente é um FLOCO QUE REPRODUZ = CAOS. Não se persegue Turing pleno; ATERRA-SE a órbita (teto de passos, sandbox.ORBITA_MAX) e SUPRIME-SE o floco — como o purgatório da parábola (caos→fora do círculo). Convergiu dentro da órbita = cristal (resultado); estourou = caos, devolve '_caos' e não pendura. É a mesma classificação da mineração. E LIGADA À PARÁBOLA DE VERDADE (sandbox._classe_floco → classe_parabola): CRISTAL = o laço FECHOU na posição (transformada, O(1)) ou é reta (Pisot, ordem); BORDA = aterrou mas rodou o laço O(k) (Salem, crítico); CAOS = reproduz, não aterra (fora). O chat mostra o selo: Pell→[cristal], multiplicação/for→[borda], while x>0→caos.", "epistemico": "inferencia", "exemplos": [], "id": "orbita_aterrada_caos", "memoria": "semantica", "meta": {"dominio": "programacao", "fonte": "linguagem/sandbox.py (rodar, ORBITA_MAX); conversa/colapso.py (_rodar_sandbox)"}, "origem": "wiki", "palavras_chave": [], "sinonimos": [], "tipo": "conceito", "titulo": "Órbita aterrada: o floco que reproduz é caos, suprimido"}
\section{Órbita aterrada: o floco que reproduz é caos, suprimido}
A lição do Aarão sobre as expressões ricas que não terminam: um laço não-convergente é um FLOCO QUE REPRODUZ = CAOS. Não se persegue Turing pleno; ATERRA-SE a órbita (teto de passos, sandbox.ORBITA\_MAX) e SUPRIME-SE o floco — como o purgatório da parábola (caos\ensuremath{\to}fora do círculo). Convergiu dentro da órbita = cristal (resultado); estourou = caos, devolve '\_caos' e não pendura. É a mesma classificação da mineração. E LIGADA À PARÁBOLA DE VERDADE (sandbox.\_classe\_floco \ensuremath{\to} classe\_parabola): CRISTAL = o laço FECHOU na posição (transformada, O(1)) ou é reta (Pisot, ordem); BORDA = aterrou mas rodou o laço O(k) (Salem, crítico); CAOS = reproduz, não aterra (fora). O chat mostra o selo: Pell\ensuremath{\to}[cristal], multiplicação/for\ensuremath{\to}[borda], while x>0\ensuremath{\to}caos.
\prov{relações: usa dsl\_comparacoes\_ordem; explica expressividade\_sub\_turing; usa parabola\_algebra\_conversacional; relaciona bai\_orbita\_reta\_metalica; usa sandbox\_cristal\_ao\_so; relaciona problema\_da\_parada; usa tiffany\_roda\_codigo; relaciona conjugados\_galois\_parabola; relaciona transformada\_metalica}
\prov{proveniência: wiki \textperiodcentered{} linguagem/sandbox.py (rodar, ORBITA\_MAX); conversa/colapso.py (\_rodar\_sandbox) \textperiodcentered{} inferencia \textperiodcentered{} confiança 1.0 \textperiodcentered{} domínio programacao \textperiodcentered{} história 28 versões no jornal}

%CRISTAL {"arestas": [{"alvo": "banco_de_dados", "confianca": 1.0, "epistemico": "fato", "origem": "wiki", "rel": "usa"}, {"alvo": "modelo_relacional", "confianca": 1.0, "epistemico": "fato", "origem": "wiki", "rel": "usa"}, {"alvo": "programacao_orientada_objetos", "confianca": 1.0, "epistemico": "fato", "origem": "wiki", "rel": "usa"}], "confianca": 1.0, "contraexemplos": [], "descricao": "Mapear tabelas do banco em objetos da linguagem (Prisma, TypeORM) — consultar com código, não SQL cru; conveniência com um custo de controle.", "epistemico": "inferencia", "exemplos": [], "id": "orm", "memoria": "semantica", "meta": {"dominio": "programacao", "fonte": ""}, "origem": "wiki", "palavras_chave": [], "sinonimos": [], "tipo": "conceito", "titulo": "ORM (Mapeamento Objeto-Relacional)"}
\section{ORM (Mapeamento Objeto-Relacional)}
Mapear tabelas do banco em objetos da linguagem (Prisma, TypeORM) — consultar com código, não SQL cru; conveniência com um custo de controle.
\prov{relações: usa banco\_de\_dados; usa modelo\_relacional; usa programacao\_orientada\_objetos}
\prov{proveniência: wiki \textperiodcentered{} inferencia \textperiodcentered{} confiança 1.0 \textperiodcentered{} domínio programacao \textperiodcentered{} história 8 versões no jornal}

%CRISTAL {"arestas": [{"alvo": "tiffany", "confianca": 1.0, "epistemico": "fato", "origem": "wiki", "rel": "usa"}, {"alvo": "operacao_broca_so", "confianca": 1.0, "epistemico": "fato", "origem": "wiki", "rel": "parte_de"}], "confianca": 1.0, "contraexemplos": [], "descricao": "code/orquestrador.py:roteia_agente classifica a mensagem em Conversa/Software/Paper/Pesquisa por regex (default Conversa) e delega a tiffany_platform.orquestrador_multi para consulta multi-idx.", "epistemico": "fato", "exemplos": [], "id": "orquestrador_roteador_regex", "memoria": "semantica", "meta": {"autor": "broca-so", "dominio": "operacao_so", "fonte": "code/orquestrador.py"}, "origem": "wiki", "palavras_chave": [], "sinonimos": [], "tipo": "conceito", "titulo": "O Orquestrador (roteador por regex)"}
\section{O Orquestrador (roteador por regex)}
code/orquestrador.py:roteia\_agente classifica a mensagem em Conversa/Software/Paper/Pesquisa por regex (default Conversa) e delega a tiffany\_platform.orquestrador\_multi para consulta multi-idx.
\prov{relações: usa tiffany; parte\_de operacao\_broca\_so}
\prov{proveniência: wiki \textperiodcentered{} code/orquestrador.py \textperiodcentered{} fato \textperiodcentered{} confiança 1.0 \textperiodcentered{} domínio operacao\_so \textperiodcentered{} história 3 versões no jornal}

%CRISTAL {"arestas": [{"alvo": "aprendizado_de_maquina", "confianca": 1.0, "epistemico": "fato", "origem": "wiki", "rel": "usa"}], "confianca": 1.0, "contraexemplos": [], "descricao": "Decorar os dados de treino em vez de aprender o padrão — acerta o visto, erra o novo. O eterno inimigo; regularizar é a cura.", "epistemico": "inferencia", "exemplos": [], "id": "overfitting", "memoria": "semantica", "meta": {"dominio": "programacao", "fonte": ""}, "origem": "wiki", "palavras_chave": [], "sinonimos": [], "tipo": "conceito", "titulo": "Overfitting e Generalização"}
\section{Overfitting e Generalização}
Decorar os dados de treino em vez de aprender o padrão — acerta o visto, erra o novo. O eterno inimigo; regularizar é a cura.
\prov{relações: usa aprendizado\_de\_maquina}
\prov{proveniência: wiki \textperiodcentered{} inferencia \textperiodcentered{} confiança 1.0 \textperiodcentered{} domínio programacao \textperiodcentered{} história 2 versões no jornal}

%CRISTAL {"arestas": [{"alvo": "gerenciamento_de_memoria", "confianca": 1.0, "epistemico": "fato", "origem": "wiki", "rel": "especializa"}, {"alvo": "rust", "confianca": 1.0, "epistemico": "fato", "origem": "wiki", "rel": "parte_de"}], "confianca": 1.0, "contraexemplos": [], "descricao": "Rust: cada valor tem um dono; empréstimos são checados em compilação — segurança de memória SEM coletor nem overhead.", "epistemico": "fato", "exemplos": [], "id": "ownership_borrow", "memoria": "semantica", "meta": {"autor": "Rust", "dominio": "programacao", "fonte": ""}, "origem": "wiki", "palavras_chave": [], "sinonimos": [], "tipo": "conceito", "titulo": "Ownership e Borrow Checker"}
\section{Ownership e Borrow Checker}
Rust: cada valor tem um dono; empréstimos são checados em compilação — segurança de memória SEM coletor nem overhead.
\prov{relações: especializa gerenciamento\_de\_memoria; parte\_de rust}
\prov{proveniência: wiki \textperiodcentered{} fato \textperiodcentered{} confiança 1.0 \textperiodcentered{} domínio programacao \textperiodcentered{} história 6 versões no jornal}

%CRISTAL {"arestas": [{"alvo": "rede_de_computadores", "confianca": 1.0, "epistemico": "fato", "origem": "wiki", "rel": "especializa"}, {"alvo": "goldchain", "confianca": 1.0, "epistemico": "fato", "origem": "wiki", "rel": "explica"}], "confianca": 1.0, "contraexemplos": [], "descricao": "Uma rede sem servidor central: cada nó é cliente e servidor. Robusta e sem dono — o modelo do BitTorrent e do goldchain.", "epistemico": "fato", "exemplos": [], "id": "p2p", "memoria": "semantica", "meta": {"dominio": "programacao", "fonte": ""}, "origem": "wiki", "palavras_chave": [], "sinonimos": [], "tipo": "conceito", "titulo": "Rede P2P (Ponto a Ponto)"}
\section{Rede P2P (Ponto a Ponto)}
Uma rede sem servidor central: cada nó é cliente e servidor. Robusta e sem dono — o modelo do BitTorrent e do goldchain.
\prov{relações: especializa rede\_de\_computadores; explica goldchain}
\prov{proveniência: wiki \textperiodcentered{} fato \textperiodcentered{} confiança 1.0 \textperiodcentered{} domínio programacao \textperiodcentered{} história 6 versões no jornal}

%CRISTAL {"arestas": [{"alvo": "classes_p_np_cook_levin", "confianca": 1.0, "epistemico": "fato", "origem": "wiki", "rel": "usa"}], "confianca": 1.0, "contraexemplos": [], "descricao": "Todo problema cuja solução é verificável rapidamente (NP) também pode ser RESOLVIDO rapidamente (P)? Isto é, P=NP? Um dos 7 problemas do milênio (Clay, US1M); NÃO RESOLVIDO. Consenso: P≠NP (não provado). Contraste didático: sabemos P⊊EXPTIME (teorema da hierarquia temporal, Hartmanis-Stearns 1965, por diagonalização) mas NÃO sabemos P vs NP. ABERTO.", "epistemico": "hipotese", "exemplos": [], "id": "p_vs_np_milenio", "memoria": "semantica", "meta": {"autor": "broca-so", "dominio": "computacao", "fonte": "Clay Mathematics Institute; Cook; Hartmanis-Stearns 1965"}, "origem": "wiki", "palavras_chave": [], "sinonimos": [], "tipo": "conceito", "titulo": "P vs NP (o problema do milênio)"}
\section{P vs NP (o problema do milênio)}
Todo problema cuja solução é verificável rapidamente (NP) também pode ser RESOLVIDO rapidamente (P)? Isto é, P=NP? Um dos 7 problemas do milênio (Clay, US1M); NÃO RESOLVIDO. Consenso: P\ensuremath{\ne}NP (não provado). Contraste didático: sabemos P\ensuremath{\subsetneq}EXPTIME (teorema da hierarquia temporal, Hartmanis-Stearns 1965, por diagonalização) mas NÃO sabemos P vs NP. ABERTO.
\prov{relações: usa classes\_p\_np\_cook\_levin}
\prov{proveniência: wiki \textperiodcentered{} Clay Mathematics Institute; Cook; Hartmanis-Stearns 1965 \textperiodcentered{} hipotese \textperiodcentered{} confiança 1.0 \textperiodcentered{} domínio computacao \textperiodcentered{} história 3 versões no jornal}

%CRISTAL {"arestas": [{"alvo": "arquitetura_de_software", "confianca": 1.0, "epistemico": "fato", "origem": "wiki", "rel": "usa"}, {"alvo": "programacao_orientada_objetos", "confianca": 1.0, "epistemico": "fato", "origem": "wiki", "rel": "usa"}], "confianca": 1.0, "contraexemplos": [], "descricao": "Uma solução reutilizável para um problema recorrente de design (Gang of Four) — vocabulário compartilhado, não receita a copiar cega.", "epistemico": "fato", "exemplos": [], "id": "padrao_de_projeto", "memoria": "semantica", "meta": {"autor": "Gang of Four", "dominio": "programacao", "fonte": ""}, "origem": "wiki", "palavras_chave": [], "sinonimos": [], "tipo": "conceito", "titulo": "Padrão de Projeto"}
\section{Padrão de Projeto}
Uma solução reutilizável para um problema recorrente de design (Gang of Four) — vocabulário compartilhado, não receita a copiar cega.
\prov{relações: usa arquitetura\_de\_software; usa programacao\_orientada\_objetos}
\prov{proveniência: wiki \textperiodcentered{} fato \textperiodcentered{} confiança 1.0 \textperiodcentered{} domínio programacao \textperiodcentered{} história 3 versões no jornal}

%CRISTAL {"arestas": [{"alvo": "funcao_pura", "confianca": 1.0, "epistemico": "fato", "origem": "wiki", "rel": "usa"}, {"alvo": "design_laco_e_posicao", "confianca": 1.0, "epistemico": "fato", "origem": "wiki", "rel": "relaciona"}], "confianca": 1.0, "contraexemplos": [], "descricao": "Derivar o próximo estado de (estado atual, ação) por uma função pura — um FOLD sobre o fluxo de eventos. A espinha do Redux e da arquitetura Elm.", "epistemico": "fato", "exemplos": [], "id": "padrao_reducer", "memoria": "semantica", "meta": {"dominio": "programacao", "fonte": ""}, "origem": "wiki", "palavras_chave": [], "sinonimos": [], "tipo": "conceito", "titulo": "Padrão Reducer (Fold)"}
\section{Padrão Reducer (Fold)}
Derivar o próximo estado de (estado atual, ação) por uma função pura — um FOLD sobre o fluxo de eventos. A espinha do Redux e da arquitetura Elm.
\prov{relações: usa funcao\_pura; relaciona design\_laco\_e\_posicao}
\prov{proveniência: wiki \textperiodcentered{} fato \textperiodcentered{} confiança 1.0 \textperiodcentered{} domínio programacao \textperiodcentered{} história 3 versões no jornal}

%CRISTAL {"arestas": [{"alvo": "algebra_tecidos_tropical", "confianca": 1.0, "epistemico": "fato", "origem": "wiki", "rel": "usa"}, {"alvo": "cofre_parabola_hero", "confianca": 1.0, "epistemico": "fato", "origem": "wiki", "rel": "usa"}, {"alvo": "symbolic_beta_multifractal_producao", "confianca": 1.0, "epistemico": "fato", "origem": "wiki", "rel": "explica"}, {"alvo": "tiffany", "confianca": 1.0, "epistemico": "fato", "origem": "wiki", "rel": "parte_de"}], "confianca": 1.0, "contraexemplos": [], "descricao": "A taxa de tecidos não se fixa por limiar ad-hoc — a PARÁBOLA a controla. O tecido N contribui o marginal r_N; define-se c²_N = r_N/r_1 (cristal relativo ao primeiro) e a_N = 1−c²_N (calor). ÓTIMO = último N com c²_N > a_N (> 0,5): enquanto o tecido marginal adiciona mais CRISTAL que CALOR. Depois, somar é net-dissipativo (a>c²) — a taxa total ainda sobe, mas em puro calor. Estudo na GEX (20 cores, goldchain/estudo_multifractal.py): por-tecido decai 6.892→3.350 Pell/s (N=1→18); combinado sobe 6.892→61.069 (~9×); a PARÁBOLA fixa o ótimo em 11 tecidos (49.685 Pell/s combinado) — o pico bruto (19) é além, mas net-calor. [D] medido.", "epistemico": "fato", "exemplos": [], "id": "parabola_controla_taxa_tecidos", "memoria": "semantica", "meta": {"autor": "Lopes/broca-so", "dominio": "goldchain", "fonte": "goldchain/estudo_multifractal.py:otimo_parabola (c²>a); GEX 20 cores → 11 tecidos"}, "origem": "wiki", "palavras_chave": [], "sinonimos": [], "tipo": "conceito", "titulo": "A parábola controla a taxa de tecidos: ótimo = último com c² > a (GEX: 11)"}
\section{A parábola controla a taxa de tecidos: ótimo = último com c\ensuremath{^{2}} > a (GEX: 11)}
A taxa de tecidos não se fixa por limiar ad-hoc — a PARÁBOLA a controla. O tecido N contribui o marginal r\_N; define-se c\ensuremath{^{2}}\_N = r\_N/r\_1 (cristal relativo ao primeiro) e a\_N = 1\ensuremath{-}c\ensuremath{^{2}}\_N (calor). ÓTIMO = último N com c\ensuremath{^{2}}\_N > a\_N (> 0,5): enquanto o tecido marginal adiciona mais CRISTAL que CALOR. Depois, somar é net-dissipativo (a>c\ensuremath{^{2}}) — a taxa total ainda sobe, mas em puro calor. Estudo na GEX (20 cores, goldchain/estudo\_multifractal.py): por-tecido decai 6.892\ensuremath{\to}3.350 Pell/s (N=1\ensuremath{\to}18); combinado sobe 6.892\ensuremath{\to}61.069 (\~{}9$\times$); a PARÁBOLA fixa o ótimo em 11 tecidos (49.685 Pell/s combinado) — o pico bruto (19) é além, mas net-calor. [D] medido.
\prov{relações: usa algebra\_tecidos\_tropical; usa cofre\_parabola\_hero; explica symbolic\_beta\_multifractal\_producao; parte\_de tiffany}
\prov{proveniência: wiki \textperiodcentered{} goldchain/estudo\_multifractal.py:otimo\_parabola (c\ensuremath{^{2}}>a); GEX 20 cores \ensuremath{\to} 11 tecidos \textperiodcentered{} fato \textperiodcentered{} confiança 1.0 \textperiodcentered{} domínio goldchain \textperiodcentered{} história 20 versões no jornal}

%CRISTAL {"arestas": [{"alvo": "programacao", "confianca": 1.0, "epistemico": "fato", "origem": "wiki", "rel": "parte_de"}], "confianca": 1.0, "contraexemplos": [], "descricao": "Um estilo fundamental de estruturar o cálculo — imperativo, funcional, orientado a objetos, lógico, declarativo.", "epistemico": "fato", "exemplos": [], "id": "paradigma_de_programacao", "memoria": "semantica", "meta": {"dominio": "programacao", "fonte": ""}, "origem": "wiki", "palavras_chave": [], "sinonimos": [], "tipo": "conceito", "titulo": "Paradigma de Programação"}
\section{Paradigma de Programação}
Um estilo fundamental de estruturar o cálculo — imperativo, funcional, orientado a objetos, lógico, declarativo.
\prov{relações: parte\_de programacao}
\prov{proveniência: wiki \textperiodcentered{} fato \textperiodcentered{} confiança 1.0 \textperiodcentered{} domínio programacao \textperiodcentered{} história 4 versões no jornal}

%CRISTAL {"arestas": [{"alvo": "documentacao", "confianca": 1.0, "epistemico": "fato", "origem": "wiki", "rel": "usa"}], "confianca": 1.0, "contraexemplos": [], "descricao": "Portable Document Format (Adobe): descreve uma página fixa e portátil, idêntica em qualquer tela ou impressora — o documento congelado.", "epistemico": "fato", "exemplos": [], "id": "pdf", "memoria": "semantica", "meta": {"autor": "Adobe", "dominio": "programacao", "fonte": ""}, "origem": "wiki", "palavras_chave": [], "sinonimos": [], "tipo": "conceito", "titulo": "PDF"}
\section{PDF}
Portable Document Format (Adobe): descreve uma página fixa e portátil, idêntica em qualquer tela ou impressora — o documento congelado.
\prov{relações: usa documentacao}
\prov{proveniência: wiki \textperiodcentered{} fato \textperiodcentered{} confiança 1.0 \textperiodcentered{} domínio programacao \textperiodcentered{} história 2 versões no jornal}

%CRISTAL {"arestas": [{"alvo": "logo_pedra_extraida", "confianca": 1.0, "epistemico": "fato", "origem": "wiki", "rel": "especializa"}, {"alvo": "design_system_cristalchain", "confianca": 1.0, "epistemico": "fato", "origem": "wiki", "rel": "parte_de"}, {"alvo": "pipe_svg_baricentrico", "confianca": 1.0, "epistemico": "fato", "origem": "wiki", "rel": "usa"}, {"alvo": "area_render_gema", "confianca": 1.0, "epistemico": "fato", "origem": "wiki", "rel": "parte_de"}], "confianca": 1.0, "contraexemplos": [], "descricao": "A pedra do logo era só uma EXTRAÇÃO recortada do board do ChatGPT; agora é RECONSTRUÍDA em 3D com iluminação pelo raytracer, PIXEL A PIXEL, com até 6 RICOCHETES internos (reflexão interna total + Fresnel parcial — o fogo real da gema: a cada face a luz que sai é refratada com dispersão 7D e somada, a que reflete continua ricocheteando). A gema em corte TRILHÃO (r=1+k·cos(3θ) — a simetria 3-fold da turmalina) renderiza numa cena de ESTÚDIO DE JOALHERIA: uma softbox teal ampla no alto na cor do GABARITO do GPT (amostrada de model/design: sombra #052327 → luz #25c6cc, o azul Paraíba) e o FILTRO de absorção (ABS) CALIBRADO por FREQUÊNCIA a essa cor — o que casa a fonte de luz com a iluminação do GPT e tira o escuro. Sai o HERO (card petróleo com glow e reflexo) e o RECORTE (o jewel na capa do paper, que substitui a extração). raytracer.gerar_logo_3d; assets/cristalchain/logo_pedra_hero.png, logo_pedra.png.", "epistemico": "fato", "exemplos": [], "id": "pedra_logo_3d", "memoria": "semantica", "meta": {"dominio": "sistema", "fonte": "logo 3d pipe"}, "origem": "wiki", "palavras_chave": [], "sinonimos": [], "tipo": "conceito", "titulo": "A Pedra de Logo Reconstruída em 3D (a extração vira reconstrução)"}
\section{A Pedra de Logo Reconstruída em 3D (a extração vira reconstrução)}
A pedra do logo era só uma EXTRAÇÃO recortada do board do ChatGPT; agora é RECONSTRUÍDA em 3D com iluminação pelo raytracer, PIXEL A PIXEL, com até 6 RICOCHETES internos (reflexão interna total + Fresnel parcial — o fogo real da gema: a cada face a luz que sai é refratada com dispersão 7D e somada, a que reflete continua ricocheteando). A gema em corte TRILHÃO (r=1+k\textperiodcentered cos(3\ensuremath{\theta}) — a simetria 3-fold da turmalina) renderiza numa cena de ESTÚDIO DE JOALHERIA: uma softbox teal ampla no alto na cor do GABARITO do GPT (amostrada de model/design: sombra \#052327 \ensuremath{\to} luz \#25c6cc, o azul Paraíba) e o FILTRO de absorção (ABS) CALIBRADO por FREQUÊNCIA a essa cor — o que casa a fonte de luz com a iluminação do GPT e tira o escuro. Sai o HERO (card petróleo com glow e reflexo) e o RECORTE (o jewel na capa do paper, que substitui a extração). raytracer.gerar\_logo\_3d; assets/cristalchain/logo\_pedra\_hero.png, logo\_pedra.png.
\prov{relações: especializa logo\_pedra\_extraida; parte\_de design\_system\_cristalchain; usa pipe\_svg\_baricentrico; parte\_de area\_render\_gema}
\prov{proveniência: wiki \textperiodcentered{} logo 3d pipe \textperiodcentered{} fato \textperiodcentered{} confiança 1.0 \textperiodcentered{} domínio sistema \textperiodcentered{} história 11 versões no jornal}

%CRISTAL {"arestas": [{"alvo": "raytracer_no_metal_e_o_sol", "confianca": 1.0, "epistemico": "fato", "origem": "wiki", "rel": "usa"}, {"alvo": "espectro_irredutivel_floco", "confianca": 1.0, "epistemico": "fato", "origem": "wiki", "rel": "usa"}, {"alvo": "turmalina_paraiba_ancora", "confianca": 1.0, "epistemico": "fato", "origem": "wiki", "rel": "usa"}, {"alvo": "area_optica_gema", "confianca": 1.0, "epistemico": "fato", "origem": "wiki", "rel": "relaciona"}], "confianca": 1.0, "contraexemplos": [], "descricao": "O dado real fecha o pipeline gráfico: a foto de uma turmalina Paraíba de verdade (model/pedras/Turmalina.webp, cor #05BDC3, o cyan Paraíba) é tratada pela máquina (linguagem/pedra_real.py, 4/4) — SEGMENTA a pedra única (a região cyan, a maior componente), RECONSTRÓI pela ESTRELA (SVD, os modos, 99.8% da energia), TRANSFORMA (gira o plano) e ILUMINA (relight sob o Sol: difuso + especular). E a cor média da pedra CALIBRA a absorção Beer-Lambert do raytracer (ABS=−ln(cor/max)/L): o vermelho é quase todo comido (o cobre), o verde e o azul passam — a gema renderizada ganha a cor da pedra de verdade. Além disso: mais clara (exposição) e com REFLEXO na superfície (o chão polido espelha a gema).", "epistemico": "fato", "exemplos": [], "id": "pedra_real_calibra_o_render", "memoria": "semantica", "meta": {"dominio": "sistema", "fonte": "pipeline gráfico"}, "origem": "wiki", "palavras_chave": [], "sinonimos": [], "tipo": "conceito", "titulo": "A Pedra Real Calibra o Render (a foto tratada)"}
\section{A Pedra Real Calibra o Render (a foto tratada)}
O dado real fecha o pipeline gráfico: a foto de uma turmalina Paraíba de verdade (model/pedras/Turmalina.webp, cor \#05BDC3, o cyan Paraíba) é tratada pela máquina (linguagem/pedra\_real.py, 4/4) — SEGMENTA a pedra única (a região cyan, a maior componente), RECONSTRÓI pela ESTRELA (SVD, os modos, 99.8\% da energia), TRANSFORMA (gira o plano) e ILUMINA (relight sob o Sol: difuso + especular). E a cor média da pedra CALIBRA a absorção Beer-Lambert do raytracer (ABS=\ensuremath{-}ln(cor/max)/L): o vermelho é quase todo comido (o cobre), o verde e o azul passam — a gema renderizada ganha a cor da pedra de verdade. Além disso: mais clara (exposição) e com REFLEXO na superfície (o chão polido espelha a gema).
\prov{relações: usa raytracer\_no\_metal\_e\_o\_sol; usa espectro\_irredutivel\_floco; usa turmalina\_paraiba\_ancora; relaciona area\_optica\_gema}
\prov{proveniência: wiki \textperiodcentered{} pipeline gráfico \textperiodcentered{} fato \textperiodcentered{} confiança 1.0 \textperiodcentered{} domínio sistema \textperiodcentered{} história 7 versões no jornal}

%CRISTAL {"arestas": [{"alvo": "algebra_de_peirce", "confianca": 1.0, "epistemico": "fato", "origem": "wiki", "rel": "especializa"}, {"alvo": "grafo_tiffany_e_algebra_peirce", "confianca": 1.0, "epistemico": "fato", "origem": "wiki", "rel": "usa"}, {"alvo": "decomposicao_peirce", "confianca": 1.0, "epistemico": "fato", "origem": "wiki", "rel": "usa"}, {"alvo": "traducao_como_interpretante", "confianca": 1.0, "epistemico": "fato", "origem": "wiki", "rel": "explica"}, {"alvo": "colapso", "confianca": 1.0, "epistemico": "fato", "origem": "wiki", "rel": "explica"}, {"alvo": "dsl_frontend_linguagem_gatos", "confianca": 1.0, "epistemico": "fato", "origem": "wiki", "rel": "relaciona"}, {"alvo": "linguagem_de_programacao", "confianca": 1.0, "epistemico": "fato", "origem": "wiki", "rel": "relaciona"}], "confianca": 1.0, "contraexemplos": [], "descricao": "O broca-so tem DOIS DSLs. O cat-DSL (linguagem/dsl.py) computa sobre gatos/tensor. O PEIRCE-DSL é a álgebra de Peirce (B booleana de conceitos, R de relações, e o produto ':') como LINGUAGEM da SUBSTITUIÇÃO — e o grafo da Tiffany JÁ é essa álgebra concreta (grafo_tiffany_e_algebra_peirce). Nele, o operador ':' (substituir um signo/relação por outro) É a mesma operação que a TRADUÇÃO (o interpretante de Peirce), que o COLAPSO da Tiffany (trocar a pergunta pelo conceito) e que a órbita do BAI. Não é um DSL a construir: é o que a Tiffany já executa ao responder.", "epistemico": "inferencia", "exemplos": [], "id": "peirce_dsl", "memoria": "semantica", "meta": {"dominio": "programacao", "fonte": "algebra_de_peirce; grafo_tiffany_e_algebra_peirce; conversa/colapso.py"}, "origem": "wiki", "palavras_chave": [], "sinonimos": [], "tipo": "conceito", "titulo": "O Peirce-DSL — a álgebra (B, R, :) como linguagem da substituição"}
\section{O Peirce-DSL — a álgebra (B, R, :) como linguagem da substituição}
O broca-so tem DOIS DSLs. O cat-DSL (linguagem/dsl.py) computa sobre gatos/tensor. O PEIRCE-DSL é a álgebra de Peirce (B booleana de conceitos, R de relações, e o produto ':') como LINGUAGEM da SUBSTITUIÇÃO — e o grafo da Tiffany JÁ é essa álgebra concreta (grafo\_tiffany\_e\_algebra\_peirce). Nele, o operador ':' (substituir um signo/relação por outro) É a mesma operação que a TRADUÇÃO (o interpretante de Peirce), que o COLAPSO da Tiffany (trocar a pergunta pelo conceito) e que a órbita do BAI. Não é um DSL a construir: é o que a Tiffany já executa ao responder.
\prov{relações: especializa algebra\_de\_peirce; usa grafo\_tiffany\_e\_algebra\_peirce; usa decomposicao\_peirce; explica traducao\_como\_interpretante; explica colapso; relaciona dsl\_frontend\_linguagem\_gatos; relaciona linguagem\_de\_programacao}
\prov{proveniência: wiki \textperiodcentered{} algebra\_de\_peirce; grafo\_tiffany\_e\_algebra\_peirce; conversa/colapso.py \textperiodcentered{} inferencia \textperiodcentered{} confiança 1.0 \textperiodcentered{} domínio programacao \textperiodcentered{} história 8 versões no jornal}

%CRISTAL {"arestas": [{"alvo": "dimensao_de_prata_pell", "confianca": 1.0, "epistemico": "fato", "origem": "wiki", "rel": "especializa"}, {"alvo": "reta_metalica_geral", "confianca": 1.0, "epistemico": "fato", "origem": "wiki", "rel": "usa"}, {"alvo": "corpo_f2k", "confianca": 1.0, "epistemico": "fato", "origem": "wiki", "rel": "usa"}, {"alvo": "symbolic_beta_pipe_so", "confianca": 1.0, "epistemico": "fato", "origem": "wiki", "rel": "explica"}, {"alvo": "tiffany", "confianca": 1.0, "epistemico": "fato", "origem": "wiki", "rel": "parte_de"}, {"alvo": "algebra_posicional_prata", "confianca": 1.0, "epistemico": "fato", "origem": "wiki", "rel": "parte_de"}], "confianca": 1.0, "contraexemplos": [], "descricao": "A linha do SymbolicBeta demorava porque a cunhagem materializava o VALOR do Pell — decimal de milhares de dígitos, recomputado do zero a cada emissão (O(idx)), multiplicado em 𝔽₂¹²⁸ e hasheado, num ledger de 18 MB. O certo (Peirce/binário, matemática do Lopes): trabalhar só com a POSIÇÃO n. ÁLGEBRA ISOMORFA: o par (Pₙ-₁,Pₙ) ≡ a Q-matriz Aⁿ (A=[[2,1],[1,0]]) — posição n ↔ Aⁿ ↔ σⁿ (σ=1+√2, razão de prata); somar posições ↔ multiplicar (Pm+ₙ=Pm Pₙ+₁+Pm-₁Pₙ). Emitir = UM passo de posição (×A) em 128 bits (mod 2¹²⁸), O(1) — nunca se materializa o valor. legível=True por construção (le∘mineraₙabase = identidade em 𝔽₂¹²⁸ p/ símbolo<2¹²⁸; não computa o inverso por emissão). Ganho: 20 ms → 87 µs/emissão (240×); ledger 18 MB → KB (compacto); gravação atômica (temp+replace, não corrompe em kill). [D] roda na GEX.", "epistemico": "fato", "exemplos": [], "id": "pell_algebra_posicao", "memoria": "semantica", "meta": {"autor": "Lopes/broca-so", "dominio": "goldchain", "fonte": "goldchain/pell.py (passo_pell/pell_ab_em, MOD 2¹²⁸); goldchain/minera.py:emite_pell"}, "origem": "wiki", "palavras_chave": [], "sinonimos": [], "tipo": "conceito", "titulo": "O Pell pela POSIÇÃO: álgebra isomorfa (Q-matriz Aⁿ ↔ σⁿ), 128 bits, O(1)"}
\section{O Pell pela POSIÇÃO: álgebra isomorfa (Q-matriz A\ensuremath{^{n}} \ensuremath{\leftrightarrow} \ensuremath{\sigma}\ensuremath{^{n}}), 128 bits, O(1)}
A linha do SymbolicBeta demorava porque a cunhagem materializava o VALOR do Pell — decimal de milhares de dígitos, recomputado do zero a cada emissão (O(idx)), multiplicado em \ensuremath{\mathbb{F}}\ensuremath{_{2}}\ensuremath{^{128}} e hasheado, num ledger de 18 MB. O certo (Peirce/binário, matemática do Lopes): trabalhar só com a POSIÇÃO n. ÁLGEBRA ISOMORFA: o par (P\ensuremath{_{n}}-\ensuremath{_{1}},P\ensuremath{_{n}}) \ensuremath{\equiv} a Q-matriz A\ensuremath{^{n}} (A=[[2,1],[1,0]]) — posição n \ensuremath{\leftrightarrow} A\ensuremath{^{n}} \ensuremath{\leftrightarrow} \ensuremath{\sigma}\ensuremath{^{n}} (\ensuremath{\sigma}=1+\ensuremath{\sqrt{\,}}2, razão de prata); somar posições \ensuremath{\leftrightarrow} multiplicar (Pm+\ensuremath{_{n}}=Pm P\ensuremath{_{n}}+\ensuremath{_{1}}+Pm-\ensuremath{_{1}}P\ensuremath{_{n}}). Emitir = UM passo de posição ($\times$A) em 128 bits (mod 2\ensuremath{^{128}}), O(1) — nunca se materializa o valor. legível=True por construção (le\ensuremath{\circ}minera\ensuremath{_{n}}abase = identidade em \ensuremath{\mathbb{F}}\ensuremath{_{2}}\ensuremath{^{128}} p/ símbolo<2\ensuremath{^{128}}; não computa o inverso por emissão). Ganho: 20 ms \ensuremath{\to} 87 \ensuremath{\mu}s/emissão (240$\times$); ledger 18 MB \ensuremath{\to} KB (compacto); gravação atômica (temp+replace, não corrompe em kill). [D] roda na GEX.
\prov{relações: especializa dimensao\_de\_prata\_pell; usa reta\_metalica\_geral; usa corpo\_f2k; explica symbolic\_beta\_pipe\_so; parte\_de tiffany; parte\_de algebra\_posicional\_prata}
\prov{proveniência: wiki \textperiodcentered{} goldchain/pell.py (passo\_pell/pell\_ab\_em, MOD 2\ensuremath{^{128}}); goldchain/minera.py:emite\_pell \textperiodcentered{} fato \textperiodcentered{} confiança 1.0 \textperiodcentered{} domínio goldchain \textperiodcentered{} história 8 versões no jornal}

%CRISTAL {"arestas": [{"alvo": "rede_neural_artificial", "confianca": 1.0, "epistemico": "fato", "origem": "wiki", "rel": "parte_de"}], "confianca": 1.0, "contraexemplos": [], "descricao": "Rosenblatt: o neurônio artificial mínimo — soma pesada com limiar. Um só não faz XOR; empilhados, aproximam qualquer função.", "epistemico": "fato", "exemplos": [], "id": "perceptron", "memoria": "semantica", "meta": {"autor": "Rosenblatt", "dominio": "programacao", "fonte": ""}, "origem": "wiki", "palavras_chave": [], "sinonimos": [], "tipo": "conceito", "titulo": "Perceptron"}
\section{Perceptron}
Rosenblatt: o neurônio artificial mínimo — soma pesada com limiar. Um só não faz XOR; empilhados, aproximam qualquer função.
\prov{relações: parte\_de rede\_neural\_artificial}
\prov{proveniência: wiki \textperiodcentered{} fato \textperiodcentered{} confiança 1.0 \textperiodcentered{} domínio programacao \textperiodcentered{} história 2 versões no jornal}

%CRISTAL {"arestas": [{"alvo": "lacuna", "confianca": 1.0, "epistemico": "fato", "origem": "curriculo", "rel": "depende"}, {"alvo": "definicao", "confianca": 1.0, "epistemico": "fato", "origem": "curriculo", "rel": "usa"}, {"alvo": "word", "confianca": 0.5, "epistemico": "fato", "origem": "wiki", "rel": "relaciona"}, {"alvo": "virtualizacao", "confianca": 0.5, "epistemico": "fato", "origem": "wiki", "rel": "relaciona"}], "confianca": 1.0, "contraexemplos": [], "descricao": "Interrogação dirigida — o que falta saber sobre um conceito?", "epistemico": "fato", "exemplos": [], "id": "pergunta", "memoria": "semantica", "meta": {"dominio": "aprendizagem", "nivel": 7, "ordem": 1}, "origem": "curriculo", "palavras_chave": ["interrogativo", "lacuna"], "sinonimos": ["perguntas"], "tipo": "conceito", "titulo": "pergunta"}
\section{pergunta}
Interrogação dirigida — o que falta saber sobre um conceito?
\prov{sinónimos: perguntas}
\prov{relações: depende lacuna; usa definicao; relaciona word; relaciona virtualizacao}
\prov{proveniência: curriculo \textperiodcentered{} fato \textperiodcentered{} confiança 1.0 \textperiodcentered{} domínio aprendizagem \textperiodcentered{} história 5 versões no jornal}

%CRISTAL {"arestas": [{"alvo": "heap", "confianca": 1.0, "epistemico": "fato", "origem": "curriculo", "rel": "especializa"}, {"alvo": "resposta_a_pergunta_dificil", "confianca": 0.5, "epistemico": "fato", "origem": "wiki", "rel": "relaciona"}, {"alvo": "transcricao", "confianca": 0.5, "epistemico": "fato", "origem": "wiki", "rel": "relaciona"}, {"alvo": "topologia", "confianca": 0.5, "epistemico": "fato", "origem": "wiki", "rel": "relaciona"}, {"alvo": "pow_metal_pre_fecho", "confianca": 0.5, "epistemico": "fato", "origem": "wiki", "rel": "relaciona"}, {"alvo": "vida-morte", "confianca": 0.5, "epistemico": "fato", "origem": "wiki", "rel": "relaciona"}, {"alvo": "virtualizacao", "confianca": 0.5, "epistemico": "fato", "origem": "wiki", "rel": "relaciona"}, {"alvo": "teste_ordem", "confianca": 0.5, "epistemico": "fato", "origem": "wiki", "rel": "relaciona"}, {"alvo": "testes", "confianca": 0.5, "epistemico": "fato", "origem": "wiki", "rel": "relaciona"}, {"alvo": "por_sessao_conversadadostelemetriasessaojsonl", "confianca": 0.5, "epistemico": "fato", "origem": "wiki", "rel": "relaciona"}, {"alvo": "pow_metal_setor", "confianca": 0.5, "epistemico": "fato", "origem": "wiki", "rel": "relaciona"}, {"alvo": "tiffany_cabeca_broca", "confianca": 0.5, "epistemico": "fato", "origem": "wiki", "rel": "relaciona"}], "confianca": 1.0, "contraexemplos": [], "descricao": "LIFO — chamadas, frames, backtracking.", "epistemico": "fato", "exemplos": [], "id": "pilha", "memoria": "semantica", "meta": {"dominio": "programacao", "nivel": 6, "ordem": 4}, "origem": "curriculo", "palavras_chave": [], "sinonimos": ["stack"], "tipo": "conceito", "titulo": "pilha"}
\section{pilha}
LIFO — chamadas, frames, backtracking.
\prov{sinónimos: stack}
\prov{relações: especializa heap; relaciona resposta\_a\_pergunta\_dificil; relaciona transcricao; relaciona topologia; relaciona pow\_metal\_pre\_fecho; relaciona vida-morte; relaciona virtualizacao; relaciona teste\_ordem; relaciona testes; relaciona por\_sessao\_conversadadostelemetriasessaojsonl; relaciona pow\_metal\_setor; relaciona tiffany\_cabeca\_broca}
\prov{proveniência: curriculo \textperiodcentered{} fato \textperiodcentered{} confiança 1.0 \textperiodcentered{} domínio programacao \textperiodcentered{} história 15 versões no jornal}

%CRISTAL {"arestas": [{"alvo": "rede_de_computadores", "confianca": 1.0, "epistemico": "fato", "origem": "wiki", "rel": "usa"}, {"alvo": "modelo_osi", "confianca": 1.0, "epistemico": "fato", "origem": "wiki", "rel": "especializa"}], "confianca": 1.0, "contraexemplos": [], "descricao": "As camadas que estruturam a comunicação — enlace, rede (IP), transporte (TCP/UDP), aplicação (HTTP). Cada uma ignora as outras.", "epistemico": "fato", "exemplos": [], "id": "pilha_tcp_ip", "memoria": "semantica", "meta": {"autor": "Cerf/Kahn", "dominio": "programacao", "fonte": ""}, "origem": "wiki", "palavras_chave": [], "sinonimos": [], "tipo": "conceito", "titulo": "Pilha TCP/IP"}
\section{Pilha TCP/IP}
As camadas que estruturam a comunicação — enlace, rede (IP), transporte (TCP/UDP), aplicação (HTTP). Cada uma ignora as outras.
\prov{relações: usa rede\_de\_computadores; especializa modelo\_osi}
\prov{proveniência: wiki \textperiodcentered{} fato \textperiodcentered{} confiança 1.0 \textperiodcentered{} domínio programacao \textperiodcentered{} história 5 versões no jornal}

%CRISTAL {"arestas": [{"alvo": "biblioteca_pipe", "confianca": 1.0, "epistemico": "fato", "origem": "wiki", "rel": "parte_de"}], "confianca": 1.0, "contraexemplos": [], "descricao": "Os campos e a luz: Maxwell → a onda EM (a luz, c=1/√(ε₀μ₀)), a lei de Gauss, a luz na matéria (v=c/n, a refração é a onda EM freando) e a fase U(1) = o círculo de Gentil (a luz dá c voltas). eletromagnetismo.", "epistemico": "fato", "exemplos": [], "id": "pipe_eletromagnetismo", "memoria": "semantica", "meta": {"dominio": "biblioteca_pipe", "fonte": "biblioteca_pipe"}, "origem": "wiki", "palavras_chave": [], "sinonimos": [], "tipo": "conceito", "titulo": "Eletromagnetismo (Maxwell, a onda EM, U(1))"}
\section{Eletromagnetismo (Maxwell, a onda EM, U(1))}
Os campos e a luz: Maxwell \ensuremath{\to} a onda EM (a luz, c=1/\ensuremath{\sqrt{\,}}(\ensuremath{\varepsilon}\ensuremath{_{0}}\ensuremath{\mu}\ensuremath{_{0}})), a lei de Gauss, a luz na matéria (v=c/n, a refração é a onda EM freando) e a fase U(1) = o círculo de Gentil (a luz dá c voltas). eletromagnetismo.
\prov{relações: parte\_de biblioteca\_pipe}
\prov{proveniência: wiki \textperiodcentered{} biblioteca\_pipe \textperiodcentered{} fato \textperiodcentered{} confiança 1.0 \textperiodcentered{} domínio biblioteca\_pipe \textperiodcentered{} história 2 versões no jornal}

%CRISTAL {"arestas": [{"alvo": "biblioteca_pipe", "confianca": 1.0, "epistemico": "fato", "origem": "wiki", "rel": "parte_de"}], "confianca": 1.0, "contraexemplos": [], "descricao": "As formas: as SDFs (a esfera, a caixa, o toro), a gema analítica e a amêndoa/vesica — o que a superfície É, parada. shader_quilez (14 intersectores), shader_piramide (o vesica compilado). O 1º domínio.", "epistemico": "fato", "exemplos": [], "id": "pipe_geometria", "memoria": "semantica", "meta": {"dominio": "biblioteca_pipe", "fonte": "biblioteca_pipe"}, "origem": "wiki", "palavras_chave": [], "sinonimos": [], "tipo": "conceito", "titulo": "Geometria (as formas — a amêndoa)"}
\section{Geometria (as formas — a amêndoa)}
As formas: as SDFs (a esfera, a caixa, o toro), a gema analítica e a amêndoa/vesica — o que a superfície É, parada. shader\_quilez (14 intersectores), shader\_piramide (o vesica compilado). O 1º domínio.
\prov{relações: parte\_de biblioteca\_pipe}
\prov{proveniência: wiki \textperiodcentered{} biblioteca\_pipe \textperiodcentered{} fato \textperiodcentered{} confiança 1.0 \textperiodcentered{} domínio biblioteca\_pipe \textperiodcentered{} história 2 versões no jornal}

%CRISTAL {"arestas": [{"alvo": "biblioteca_pipe", "confianca": 1.0, "epistemico": "fato", "origem": "wiki", "rel": "parte_de"}, {"alvo": "pipe_eletromagnetismo", "confianca": 1.0, "epistemico": "fato", "origem": "wiki", "rel": "relaciona"}], "confianca": 1.0, "contraexemplos": [], "descricao": "A geometria do espaço-tempo: a luz desvia (Eddington, 4GM/c²b), o tempo dilata (redshift), a onda gravitacional anda a c (o mesmo c da luz), e a EFE G+Λg=κT é a colisão P_g=P_a = o vértice da parábola (Λ=D−1=N−1; H=√(Λ/3)=λ=log σ). O nosso universo (Λ>0)=caos. gravitacao, einstein_parabola.", "epistemico": "fato", "exemplos": [], "id": "pipe_gravitacao", "memoria": "semantica", "meta": {"dominio": "biblioteca_pipe", "fonte": "biblioteca_pipe"}, "origem": "wiki", "palavras_chave": [], "sinonimos": [], "tipo": "conceito", "titulo": "Gravitação (a EFE é o vértice, a onda a c)"}
\section{Gravitação (a EFE é o vértice, a onda a c)}
A geometria do espaço-tempo: a luz desvia (Eddington, 4GM/c\ensuremath{^{2}}b), o tempo dilata (redshift), a onda gravitacional anda a c (o mesmo c da luz), e a EFE G+\ensuremath{\Lambda}g=\ensuremath{\kappa}T é a colisão P\_g=P\_a = o vértice da parábola (\ensuremath{\Lambda}=D\ensuremath{-}1=N\ensuremath{-}1; H=\ensuremath{\sqrt{\,}}(\ensuremath{\Lambda}/3)=\ensuremath{\lambda}=log \ensuremath{\sigma}). O nosso universo (\ensuremath{\Lambda}>0)=caos. gravitacao, einstein\_parabola.
\prov{relações: parte\_de biblioteca\_pipe; relaciona pipe\_eletromagnetismo}
\prov{proveniência: wiki \textperiodcentered{} biblioteca\_pipe \textperiodcentered{} fato \textperiodcentered{} confiança 1.0 \textperiodcentered{} domínio biblioteca\_pipe \textperiodcentered{} história 3 versões no jornal}

%CRISTAL {"arestas": [{"alvo": "biblioteca_pipe", "confianca": 1.0, "epistemico": "fato", "origem": "wiki", "rel": "parte_de"}, {"alvo": "pipe_eletromagnetismo", "confianca": 1.0, "epistemico": "fato", "origem": "wiki", "rel": "relaciona"}], "confianca": 1.0, "contraexemplos": [], "descricao": "O conteúdo dos campos (o Modelo Padrão): a carga fracionária Q=N/3 (quarks em terços, hádrons inteiros — o confinamento), a anomalia que cancela (∑Q=0, o átomo neutro), o running de α (∑N_cQ²=8) e as três gerações (a trialidade). A matéria são os metais em cargas N/3. materia.", "epistemico": "fato", "exemplos": [], "id": "pipe_materia", "memoria": "semantica", "meta": {"dominio": "biblioteca_pipe", "fonte": "biblioteca_pipe"}, "origem": "wiki", "palavras_chave": [], "sinonimos": [], "tipo": "conceito", "titulo": "Matéria (a carga Q=N/3, a anomalia)"}
\section{Matéria (a carga Q=N/3, a anomalia)}
O conteúdo dos campos (o Modelo Padrão): a carga fracionária Q=N/3 (quarks em terços, hádrons inteiros — o confinamento), a anomalia que cancela (\ensuremath{\sum}Q=0, o átomo neutro), o running de \ensuremath{\alpha} (\ensuremath{\sum}N\_cQ\ensuremath{^{2}}=8) e as três gerações (a trialidade). A matéria são os metais em cargas N/3. materia.
\prov{relações: parte\_de biblioteca\_pipe; relaciona pipe\_eletromagnetismo}
\prov{proveniência: wiki \textperiodcentered{} biblioteca\_pipe \textperiodcentered{} fato \textperiodcentered{} confiança 1.0 \textperiodcentered{} domínio biblioteca\_pipe \textperiodcentered{} história 3 versões no jornal}

%CRISTAL {"arestas": [{"alvo": "biblioteca_pipe", "confianca": 1.0, "epistemico": "fato", "origem": "wiki", "rel": "parte_de"}], "confianca": 1.0, "contraexemplos": [], "descricao": "Como a forma se move: a flexão (a junta dobra, shader_juntas), a torção (o cubo torce, shader_bilinear) e as isometrias (o gato A^T M A=−M, tensor_mineral). As deformações e o que as preserva.", "epistemico": "fato", "exemplos": [], "id": "pipe_mecanica", "memoria": "semantica", "meta": {"dominio": "biblioteca_pipe", "fonte": "biblioteca_pipe"}, "origem": "wiki", "palavras_chave": [], "sinonimos": [], "tipo": "conceito", "titulo": "Mecânica (flexão, torção, isometrias)"}
\section{Mecânica (flexão, torção, isometrias)}
Como a forma se move: a flexão (a junta dobra, shader\_juntas), a torção (o cubo torce, shader\_bilinear) e as isometrias (o gato A\^{}T M A=\ensuremath{-}M, tensor\_mineral). As deformações e o que as preserva.
\prov{relações: parte\_de biblioteca\_pipe}
\prov{proveniência: wiki \textperiodcentered{} biblioteca\_pipe \textperiodcentered{} fato \textperiodcentered{} confiança 1.0 \textperiodcentered{} domínio biblioteca\_pipe \textperiodcentered{} história 2 versões no jornal}

%CRISTAL {"arestas": [{"alvo": "biblioteca_pipe", "confianca": 1.0, "epistemico": "fato", "origem": "wiki", "rel": "parte_de"}, {"alvo": "pipe_eletromagnetismo", "confianca": 1.0, "epistemico": "fato", "origem": "wiki", "rel": "relaciona"}], "confianca": 1.0, "contraexemplos": [], "descricao": "A luz na superfície: a reflexão especular (a esfera espelha, shader_transparencia) e a refração (a esfera desvia pela lei de Snell, shader_refracao), com Fresnel e Beer-Lambert. As duas metades numa interface.", "epistemico": "fato", "exemplos": [], "id": "pipe_optica", "memoria": "semantica", "meta": {"dominio": "biblioteca_pipe", "fonte": "biblioteca_pipe"}, "origem": "wiki", "palavras_chave": [], "sinonimos": [], "tipo": "conceito", "titulo": "Óptica (reflexão especular, refração)"}
\section{Óptica (reflexão especular, refração)}
A luz na superfície: a reflexão especular (a esfera espelha, shader\_transparencia) e a refração (a esfera desvia pela lei de Snell, shader\_refracao), com Fresnel e Beer-Lambert. As duas metades numa interface.
\prov{relações: parte\_de biblioteca\_pipe; relaciona pipe\_eletromagnetismo}
\prov{proveniência: wiki \textperiodcentered{} biblioteca\_pipe \textperiodcentered{} fato \textperiodcentered{} confiança 1.0 \textperiodcentered{} domínio biblioteca\_pipe \textperiodcentered{} história 3 versões no jornal}

%CRISTAL {"arestas": [{"alvo": "biblioteca_pipe", "confianca": 1.0, "epistemico": "fato", "origem": "wiki", "rel": "parte_de"}], "confianca": 1.0, "contraexemplos": [], "descricao": "A onda da matéria: Schrödinger (o pacote espalha, a dispersão), a quantização (E_n∝n²), a incerteza (Δx·Δp≥ℏ/2) e a partição Z(β)=Tr e^{−βH}=∑σ_n^{−β} = a zeta (quântica e termo no mesmo traço). O qubit é a amêndoa (o gato A_1). mecanica_quantica.", "epistemico": "fato", "exemplos": [], "id": "pipe_quantica", "memoria": "semantica", "meta": {"dominio": "biblioteca_pipe", "fonte": "biblioteca_pipe"}, "origem": "wiki", "palavras_chave": [], "sinonimos": [], "tipo": "conceito", "titulo": "Quântica (Schrödinger, a partição = a zeta)"}
\section{Quântica (Schrödinger, a partição = a zeta)}
A onda da matéria: Schrödinger (o pacote espalha, a dispersão), a quantização (E\_n\ensuremath{\propto}n\ensuremath{^{2}}), a incerteza (\ensuremath{\Delta}x\textperiodcentered \ensuremath{\Delta}p\ensuremath{\ge}\ensuremath{\hbar}/2) e a partição Z(\ensuremath{\beta})=Tr e\^{}\{\ensuremath{-}\ensuremath{\beta}H\}=\ensuremath{\sum}\ensuremath{\sigma}\_n\^{}\{\ensuremath{-}\ensuremath{\beta}\} = a zeta (quântica e termo no mesmo traço). O qubit é a amêndoa (o gato A\_1). mecanica\_quantica.
\prov{relações: parte\_de biblioteca\_pipe}
\prov{proveniência: wiki \textperiodcentered{} biblioteca\_pipe \textperiodcentered{} fato \textperiodcentered{} confiança 1.0 \textperiodcentered{} domínio biblioteca\_pipe \textperiodcentered{} história 2 versões no jornal}

%CRISTAL {"arestas": [{"alvo": "cor_baricentrica_faces", "confianca": 1.0, "epistemico": "fato", "origem": "wiki", "rel": "especializa"}, {"alvo": "design_vetorizado_svg", "confianca": 1.0, "epistemico": "fato", "origem": "wiki", "rel": "usa"}, {"alvo": "interpolacao_orbitas", "confianca": 1.0, "epistemico": "fato", "origem": "wiki", "rel": "relaciona"}, {"alvo": "area_render_gema", "confianca": 1.0, "epistemico": "fato", "origem": "wiki", "rel": "parte_de"}], "confianca": 1.0, "contraexemplos": [], "descricao": "linguagem/pipe_svg.py é a lib ÚNICA de saída do design: recebe (triângulos 2D, cor por vértice) e emite SVG VETOR, com a cor interpolada nas faces por COORDENADAS BARICÊNTRICAS (o sombreamento Gouraud, mas em <polygon>, não numa imagem raster) + contornos opcionais (as ARESTAS brilhando). raytracer.gema_svg CONECTA o gem 3D ao pipe, com a cor e a luz CASADAS ao gabarito do GPT (rampa 3-stops sombra→meio→luz): difuso plano por faceta (o corte) + CONTRASTE entre facetas, TRANSPARÊNCIA interna, um HIGHLIGHT BRANCO pequeno, as ARESTAS brilhando, e a MESA achatada numa janela lisa (o leque some) → assets/cristalchain/logo.svg, o logo vetor escalável. O pipe fica FIXO: as próximas reconstruções só precisam entregar (tris, cores). Verificado: pipe_svg 4/4, e no test_pipeline (254).", "epistemico": "fato", "exemplos": [], "id": "pipe_svg_baricentrico", "memoria": "semantica", "meta": {"dominio": "sistema", "fonte": "logo 3d pipe"}, "origem": "wiki", "palavras_chave": [], "sinonimos": [], "tipo": "conceito", "titulo": "O Pipe de Saída SVG (baricêntrico, fixado)"}
\section{O Pipe de Saída SVG (baricêntrico, fixado)}
linguagem/pipe\_svg.py é a lib ÚNICA de saída do design: recebe (triângulos 2D, cor por vértice) e emite SVG VETOR, com a cor interpolada nas faces por COORDENADAS BARICÊNTRICAS (o sombreamento Gouraud, mas em <polygon>, não numa imagem raster) + contornos opcionais (as ARESTAS brilhando). raytracer.gema\_svg CONECTA o gem 3D ao pipe, com a cor e a luz CASADAS ao gabarito do GPT (rampa 3-stops sombra\ensuremath{\to}meio\ensuremath{\to}luz): difuso plano por faceta (o corte) + CONTRASTE entre facetas, TRANSPARÊNCIA interna, um HIGHLIGHT BRANCO pequeno, as ARESTAS brilhando, e a MESA achatada numa janela lisa (o leque some) \ensuremath{\to} assets/cristalchain/logo.svg, o logo vetor escalável. O pipe fica FIXO: as próximas reconstruções só precisam entregar (tris, cores). Verificado: pipe\_svg 4/4, e no test\_pipeline (254).
\prov{relações: especializa cor\_baricentrica\_faces; usa design\_vetorizado\_svg; relaciona interpolacao\_orbitas; parte\_de area\_render\_gema}
\prov{proveniência: wiki \textperiodcentered{} logo 3d pipe \textperiodcentered{} fato \textperiodcentered{} confiança 1.0 \textperiodcentered{} domínio sistema \textperiodcentered{} história 11 versões no jornal}

%CRISTAL {"arestas": [{"alvo": "biblioteca_pipe", "confianca": 1.0, "epistemico": "fato", "origem": "wiki", "rel": "parte_de"}, {"alvo": "pipe_quantica", "confianca": 1.0, "epistemico": "fato", "origem": "wiki", "rel": "relaciona"}], "confianca": 1.0, "contraexemplos": [], "descricao": "A energia que flui e se dissipa: o fluido de ondas amortecidas (shader_fluido), a 2ª lei (a entropia sobe, a seta do tempo) e o núcleo do calor Z(β)=∑σ_n^{−β} = a zeta dos metais = o traço de Selberg (termodinamica, traco_termico). O calor É a aritmética dos metais.", "epistemico": "fato", "exemplos": [], "id": "pipe_termodinamica", "memoria": "semantica", "meta": {"dominio": "biblioteca_pipe", "fonte": "biblioteca_pipe"}, "origem": "wiki", "palavras_chave": [], "sinonimos": [], "tipo": "conceito", "titulo": "Termodinâmica (o fluido, o calor = a zeta)"}
\section{Termodinâmica (o fluido, o calor = a zeta)}
A energia que flui e se dissipa: o fluido de ondas amortecidas (shader\_fluido), a 2ª lei (a entropia sobe, a seta do tempo) e o núcleo do calor Z(\ensuremath{\beta})=\ensuremath{\sum}\ensuremath{\sigma}\_n\^{}\{\ensuremath{-}\ensuremath{\beta}\} = a zeta dos metais = o traço de Selberg (termodinamica, traco\_termico). O calor É a aritmética dos metais.
\prov{relações: parte\_de biblioteca\_pipe; relaciona pipe\_quantica}
\prov{proveniência: wiki \textperiodcentered{} biblioteca\_pipe \textperiodcentered{} fato \textperiodcentered{} confiança 1.0 \textperiodcentered{} domínio biblioteca\_pipe \textperiodcentered{} história 3 versões no jornal}

%CRISTAL {"arestas": [{"alvo": "shell", "confianca": 1.0, "epistemico": "fato", "origem": "wiki", "rel": "usa"}, {"alvo": "filosofia_unix", "confianca": 1.0, "epistemico": "fato", "origem": "wiki", "rel": "usa"}, {"alvo": "pipeline", "confianca": 1.0, "epistemico": "fato", "origem": "wiki", "rel": "relaciona"}, {"alvo": "symbolic_beta_pipe_so", "confianca": 1.0, "epistemico": "fato", "origem": "wiki", "rel": "relaciona"}], "confianca": 1.0, "contraexemplos": [], "descricao": "Ligar a saída de um programa à entrada do outro — compor ferramentas pequenas num fluxo. A ideia central da linha de comando (McIlroy).", "epistemico": "fato", "exemplos": [], "id": "pipe_unix", "memoria": "semantica", "meta": {"autor": "McIlroy", "dominio": "programacao", "fonte": ""}, "origem": "wiki", "palavras_chave": [], "sinonimos": [], "tipo": "conceito", "titulo": "Pipe (|)"}
\section{Pipe (|)}
Ligar a saída de um programa à entrada do outro — compor ferramentas pequenas num fluxo. A ideia central da linha de comando (McIlroy).
\prov{relações: usa shell; usa filosofia\_unix; relaciona pipeline; relaciona symbolic\_beta\_pipe\_so}
\prov{proveniência: wiki \textperiodcentered{} fato \textperiodcentered{} confiança 1.0 \textperiodcentered{} domínio programacao \textperiodcentered{} história 5 versões no jornal}

%CRISTAL {"arestas": [{"alvo": "sass_gpu_assembly", "confianca": 1.0, "epistemico": "fato", "origem": "wiki", "rel": "usa"}, {"alvo": "ptx_isa_virtual", "confianca": 1.0, "epistemico": "fato", "origem": "wiki", "rel": "usa"}, {"alvo": "compilador_glsl_broca", "confianca": 1.0, "epistemico": "fato", "origem": "wiki", "rel": "relaciona"}, {"alvo": "dsl_bai_xadrez", "confianca": 1.0, "epistemico": "fato", "origem": "wiki", "rel": "relaciona"}, {"alvo": "maquina_universal_gpu", "confianca": 1.0, "epistemico": "fato", "origem": "wiki", "rel": "parte_de"}], "confianca": 1.0, "contraexemplos": [], "descricao": "A cadeia que roda a 60 fps, SEM o Python no meio: a lei (BAI) → a DSL de gatos (a linguagem do sandbox) → GLSL/CUDA C → PTX (a ISA virtual) → SASS (o nativo, por arquitetura) → os cores da SM. O Python (numpy/torch) é uma CAMADA CLÁSSICA no meio: faz transformações densas na CPU (numpy) ou despacha ~1280 kernels não-fundidos/frame (torch) — nunca funde no SASS. O objetivo: o compilador do broca-so emitir PTX/SASS DIRETO (sob o Python, não através dele) — o GLSL nativo funde tudo num kernel/pixel. Hoje o glsl_compilador só emite numpy (CPU); a frente é o backend PTX/SASS (via moderngl/OpenGL para o GLSL nativo, ou CUDA/cupy para o kernel). Síntese declarada: a meta do compilador é o SASS, não o numpy.", "epistemico": "hipotese", "exemplos": [], "id": "pipeline_compilacao_gpu", "memoria": "semantica", "meta": {"dominio": "sass_gpu", "fonte": "SASS/GPU (o metal)"}, "origem": "wiki", "palavras_chave": [], "sinonimos": [], "tipo": "conceito", "titulo": "O pipeline de compilação GPU: DSL/GLSL → PTX → SASS → a SM (sem Python no meio)"}
\section{O pipeline de compilação GPU: DSL/GLSL \ensuremath{\to} PTX \ensuremath{\to} SASS \ensuremath{\to} a SM (sem Python no meio)}
A cadeia que roda a 60 fps, SEM o Python no meio: a lei (BAI) \ensuremath{\to} a DSL de gatos (a linguagem do sandbox) \ensuremath{\to} GLSL/CUDA C \ensuremath{\to} PTX (a ISA virtual) \ensuremath{\to} SASS (o nativo, por arquitetura) \ensuremath{\to} os cores da SM. O Python (numpy/torch) é uma CAMADA CLÁSSICA no meio: faz transformações densas na CPU (numpy) ou despacha \~{}1280 kernels não-fundidos/frame (torch) — nunca funde no SASS. O objetivo: o compilador do broca-so emitir PTX/SASS DIRETO (sob o Python, não através dele) — o GLSL nativo funde tudo num kernel/pixel. Hoje o glsl\_compilador só emite numpy (CPU); a frente é o backend PTX/SASS (via moderngl/OpenGL para o GLSL nativo, ou CUDA/cupy para o kernel). Síntese declarada: a meta do compilador é o SASS, não o numpy.
\prov{relações: usa sass\_gpu\_assembly; usa ptx\_isa\_virtual; relaciona compilador\_glsl\_broca; relaciona dsl\_bai\_xadrez; parte\_de maquina\_universal\_gpu}
\prov{proveniência: wiki \textperiodcentered{} SASS/GPU (o metal) \textperiodcentered{} hipotese \textperiodcentered{} confiança 1.0 \textperiodcentered{} domínio sass\_gpu \textperiodcentered{} história 42 versões no jornal}

%CRISTAL {"arestas": [{"alvo": "sandbox_cristal_ao_so", "confianca": 1.0, "epistemico": "fato", "origem": "wiki", "rel": "usa"}, {"alvo": "lei_unica_tres_dominios", "confianca": 1.0, "epistemico": "fato", "origem": "wiki", "rel": "usa"}, {"alvo": "bai_como_bussola", "confianca": 1.0, "epistemico": "fato", "origem": "wiki", "rel": "usa"}, {"alvo": "bai_orbita_reta_metalica", "confianca": 1.0, "epistemico": "fato", "origem": "wiki", "rel": "usa"}, {"alvo": "parabola_algebra_conversacional", "confianca": 1.0, "epistemico": "fato", "origem": "wiki", "rel": "usa"}, {"alvo": "tiffany_roda_codigo", "confianca": 1.0, "epistemico": "fato", "origem": "wiki", "rel": "relaciona"}, {"alvo": "colapso", "confianca": 1.0, "epistemico": "fato", "origem": "wiki", "rel": "contradiz"}, {"alvo": "tiffany", "confianca": 1.0, "epistemico": "fato", "origem": "wiki", "rel": "parte_de"}], "confianca": 1.0, "contraexemplos": [], "descricao": "O pipeline INVERSO do 'roda:' (linguagem/completar.py). Recebe um pedido vago ou incompleto em linguagem natural, ATERRA na interpretação de MENOR ENERGIA (o metal mais justo/cristal: ouro δ=6 < prata 10 < bronze 13 < cobre 17 — a intenção mais convergente), DESATERRA em texto completo + DSL de gatos, roda no sandbox e devolve resultado + classe da parábola. Intenções: fibonacci(ouro)/pell(prata)/soma·contagem(bronze)/multiplicação·potência·fatorial·primo(cobre)/mdc(prata). Ex.: 'me dá o décimo fibonacci'→89[cristal]; 'soma de 1 até 100'→5050[borda]; '2 elevado a 10'→1024. Vago demais → '(não aterrou)', honesto. Chat: 'completa:/gera: <texto>'. Fecha o par com a recuperação: colapso resolve pergunta→conceito, completar resolve pedido→código, ambos pela mínima energia da parábola.", "epistemico": "fato", "exemplos": [], "id": "pipeline_completar", "memoria": "semantica", "meta": {"dominio": "programacao", "fonte": "linguagem/completar.py (completar, _INTENCOES); conversa/colapso.py (_completar_pedido)"}, "origem": "wiki", "palavras_chave": [], "sinonimos": [], "tipo": "conceito", "titulo": "Completar — texto ambíguo → menor energia → texto completo + código"}
\section{Completar — texto ambíguo \ensuremath{\to} menor energia \ensuremath{\to} texto completo + código}
O pipeline INVERSO do 'roda:' (linguagem/completar.py). Recebe um pedido vago ou incompleto em linguagem natural, ATERRA na interpretação de MENOR ENERGIA (o metal mais justo/cristal: ouro \ensuremath{\delta}=6 < prata 10 < bronze 13 < cobre 17 — a intenção mais convergente), DESATERRA em texto completo + DSL de gatos, roda no sandbox e devolve resultado + classe da parábola. Intenções: fibonacci(ouro)/pell(prata)/soma\textperiodcentered contagem(bronze)/multiplicação\textperiodcentered potência\textperiodcentered fatorial\textperiodcentered primo(cobre)/mdc(prata). Ex.: 'me dá o décimo fibonacci'\ensuremath{\to}89[cristal]; 'soma de 1 até 100'\ensuremath{\to}5050[borda]; '2 elevado a 10'\ensuremath{\to}1024. Vago demais \ensuremath{\to} '(não aterrou)', honesto. Chat: 'completa:/gera: <texto>'. Fecha o par com a recuperação: colapso resolve pergunta\ensuremath{\to}conceito, completar resolve pedido\ensuremath{\to}código, ambos pela mínima energia da parábola.
\prov{relações: usa sandbox\_cristal\_ao\_so; usa lei\_unica\_tres\_dominios; usa bai\_como\_bussola; usa bai\_orbita\_reta\_metalica; usa parabola\_algebra\_conversacional; relaciona tiffany\_roda\_codigo; contradiz colapso; parte\_de tiffany}
\prov{proveniência: wiki \textperiodcentered{} linguagem/completar.py (completar, \_INTENCOES); conversa/colapso.py (\_completar\_pedido) \textperiodcentered{} fato \textperiodcentered{} confiança 1.0 \textperiodcentered{} domínio programacao \textperiodcentered{} história 18 versões no jornal}

%CRISTAL {"arestas": [{"alvo": "pedra_real_calibra_o_render", "confianca": 1.0, "epistemico": "fato", "origem": "wiki", "rel": "usa"}, {"alvo": "raytracer_no_metal_e_o_sol", "confianca": 1.0, "epistemico": "fato", "origem": "wiki", "rel": "usa"}, {"alvo": "metal_do_so", "confianca": 1.0, "epistemico": "fato", "origem": "wiki", "rel": "parte_de"}, {"alvo": "transformada_mineral_7d", "confianca": 1.0, "epistemico": "fato", "origem": "wiki", "rel": "usa"}, {"alvo": "corpo_estelar", "confianca": 1.0, "epistemico": "fato", "origem": "wiki", "rel": "usa"}], "confianca": 1.0, "contraexemplos": [], "descricao": "O pipeline gráfico da máquina, em três dimensões, cristalizado no metal (linguagem/pipeline_grafico.py): 2D = a IMAGEM (raster / a estrela SVD, as transformações 2D); 3D = a GEOMETRIA (o raytracer da gema, a física da luz); 7D = a ILUMINAÇÃO ESPECTRAL OCTONIÔNICA — a luz do Sol nas 7 BANDAS (as 7 unidades imaginárias do octônio, a transformada mineral 7D), a dispersão (o fogo da gema) e a distribuição espectral solar. O 7D ajuda a criar os efeitos da iluminação; os três, juntos, são o fortalecimento do pipeline gráfico. Verif.: os 3 pipelines (dim 2, 3, 7) fecham.", "epistemico": "fato", "exemplos": [], "id": "pipeline_grafico_2d_3d_7d", "memoria": "semantica", "meta": {"dominio": "sistema", "fonte": "pipeline gráfico"}, "origem": "wiki", "palavras_chave": [], "sinonimos": [], "tipo": "conceito", "titulo": "O Pipeline Gráfico 2D/3D/7D (no metal da máquina)"}
\section{O Pipeline Gráfico 2D/3D/7D (no metal da máquina)}
O pipeline gráfico da máquina, em três dimensões, cristalizado no metal (linguagem/pipeline\_grafico.py): 2D = a IMAGEM (raster / a estrela SVD, as transformações 2D); 3D = a GEOMETRIA (o raytracer da gema, a física da luz); 7D = a ILUMINAÇÃO ESPECTRAL OCTONIÔNICA — a luz do Sol nas 7 BANDAS (as 7 unidades imaginárias do octônio, a transformada mineral 7D), a dispersão (o fogo da gema) e a distribuição espectral solar. O 7D ajuda a criar os efeitos da iluminação; os três, juntos, são o fortalecimento do pipeline gráfico. Verif.: os 3 pipelines (dim 2, 3, 7) fecham.
\prov{relações: usa pedra\_real\_calibra\_o\_render; usa raytracer\_no\_metal\_e\_o\_sol; parte\_de metal\_do\_so; usa transformada\_mineral\_7d; usa corpo\_estelar}
\prov{proveniência: wiki \textperiodcentered{} pipeline gráfico \textperiodcentered{} fato \textperiodcentered{} confiança 1.0 \textperiodcentered{} domínio sistema \textperiodcentered{} história 11 versões no jornal}

%CRISTAL {"arestas": [{"alvo": "teorema_geral_traducao", "confianca": 1.0, "epistemico": "fato", "origem": "wiki", "rel": "usa"}, {"alvo": "bai_como_bussola", "confianca": 1.0, "epistemico": "fato", "origem": "wiki", "rel": "usa"}, {"alvo": "traducao_como_transformacao", "confianca": 1.0, "epistemico": "fato", "origem": "wiki", "rel": "usa"}, {"alvo": "meta_inducao_geral_simbolo", "confianca": 1.0, "epistemico": "fato", "origem": "wiki", "rel": "usa"}, {"alvo": "gramatica_portugues", "confianca": 1.0, "epistemico": "fato", "origem": "wiki", "rel": "usa"}, {"alvo": "morfologia", "confianca": 1.0, "epistemico": "fato", "origem": "wiki", "rel": "especializa"}], "confianca": 1.0, "contraexemplos": [], "descricao": "O plural não é um hack de stemming: é GRAMÁTICA do idioma. Pelo teorema geral da tradução, o plural e o lema são duas simbologias do MESMO significado; a gramática do português É o dicionário D:plural→lema (linguagem/gramatica.py: ões→ão, ns→m, s→∅). A bússola (BAI) aponta o lema — não se procura a regra, aplica-se o dicionário do idioma. O colapso normaliza plurais delegando a esse dicionário, não a um stem ad-hoc. É a meta-indução ('só o símbolo muda') na morfologia.", "epistemico": "inferencia", "exemplos": [], "id": "plural_como_traducao", "memoria": "semantica", "meta": {"dominio": "linguistica", "fonte": "linguagem/gramatica.py"}, "origem": "wiki", "palavras_chave": ["plural e singular", "como o plural vira lema", "plural é tradução", "normalizar plural"], "sinonimos": [], "tipo": "conceito", "titulo": "O Plural é uma Tradução (a bússola aponta o lema)"}
\section{O Plural é uma Tradução (a bússola aponta o lema)}
O plural não é um hack de stemming: é GRAMÁTICA do idioma. Pelo teorema geral da tradução, o plural e o lema são duas simbologias do MESMO significado; a gramática do português É o dicionário D:plural\ensuremath{\to}lema (linguagem/gramatica.py: ões\ensuremath{\to}ão, ns\ensuremath{\to}m, s\ensuremath{\to}\ensuremath{\varnothing}). A bússola (BAI) aponta o lema — não se procura a regra, aplica-se o dicionário do idioma. O colapso normaliza plurais delegando a esse dicionário, não a um stem ad-hoc. É a meta-indução ('só o símbolo muda') na morfologia.
\prov{relações: usa teorema\_geral\_traducao; usa bai\_como\_bussola; usa traducao\_como\_transformacao; usa meta\_inducao\_geral\_simbolo; usa gramatica\_portugues; especializa morfologia}
\prov{proveniência: wiki \textperiodcentered{} linguagem/gramatica.py \textperiodcentered{} inferencia \textperiodcentered{} confiança 1.0 \textperiodcentered{} domínio linguistica \textperiodcentered{} história 7 versões no jornal}

%CRISTAL {"arestas": [{"alvo": "informacao_imperfeita", "confianca": 1.0, "epistemico": "fato", "origem": "wiki", "rel": "eh_um"}, {"alvo": "estrategia_mista", "confianca": 1.0, "epistemico": "fato", "origem": "wiki", "rel": "usa"}], "confianca": 1.0, "contraexemplos": [], "descricao": "O jogo de cartas ocultas: a estratégia ótima é uma distribuição imprevisível que equilibra blefe e valor (um equilíbrio de Nash aproximado). O oposto do xadrez — o adversário raciocina sobre o que você esconde.", "epistemico": "fato", "exemplos": [], "id": "poker_jogo", "memoria": "semantica", "meta": {"dominio": "ia", "fonte": "jogos e IA"}, "origem": "wiki", "palavras_chave": [], "sinonimos": [], "tipo": "conceito", "titulo": "Pôquer (informação imperfeita + blefe)"}
\section{Pôquer (informação imperfeita + blefe)}
O jogo de cartas ocultas: a estratégia ótima é uma distribuição imprevisível que equilibra blefe e valor (um equilíbrio de Nash aproximado). O oposto do xadrez — o adversário raciocina sobre o que você esconde.
\prov{relações: eh\_um informacao\_imperfeita; usa estrategia\_mista}
\prov{proveniência: wiki \textperiodcentered{} jogos e IA \textperiodcentered{} fato \textperiodcentered{} confiança 1.0 \textperiodcentered{} domínio ia \textperiodcentered{} história 3 versões no jornal}

%CRISTAL {"arestas": [{"alvo": "sistema_de_tipos", "confianca": 1.0, "epistemico": "fato", "origem": "wiki", "rel": "usa"}], "confianca": 1.0, "contraexemplos": [], "descricao": "Um mesmo código operando sobre muitos tipos — por genéricos (paramétrico), herança (subtipo) ou traits (ad-hoc).", "epistemico": "fato", "exemplos": [], "id": "polimorfismo", "memoria": "semantica", "meta": {"dominio": "programacao", "fonte": ""}, "origem": "wiki", "palavras_chave": [], "sinonimos": [], "tipo": "conceito", "titulo": "Polimorfismo"}
\section{Polimorfismo}
Um mesmo código operando sobre muitos tipos — por genéricos (paramétrico), herança (subtipo) ou traits (ad-hoc).
\prov{relações: usa sistema\_de\_tipos}
\prov{proveniência: wiki \textperiodcentered{} fato \textperiodcentered{} confiança 1.0 \textperiodcentered{} domínio programacao \textperiodcentered{} história 4 versões no jornal}

%CRISTAL {"arestas": [{"alvo": "processo", "confianca": 1.0, "epistemico": "fato", "origem": "curriculo", "rel": "depende"}, {"alvo": "topologia", "confianca": 0.5, "epistemico": "fato", "origem": "wiki", "rel": "relaciona"}, {"alvo": "rust", "confianca": 0.5, "epistemico": "fato", "origem": "wiki", "rel": "relaciona"}, {"alvo": "por_sessao_conversadadostelemetriasessaojsonl", "confianca": 0.5, "epistemico": "fato", "origem": "wiki", "rel": "relaciona"}, {"alvo": "thread", "confianca": 0.5, "epistemico": "fato", "origem": "wiki", "rel": "relaciona"}, {"alvo": "virtualizacao", "confianca": 0.5, "epistemico": "fato", "origem": "wiki", "rel": "relaciona"}, {"alvo": "resposta_a_pergunta_dificil", "confianca": 0.5, "epistemico": "fato", "origem": "wiki", "rel": "relaciona"}, {"alvo": "testes", "confianca": 0.5, "epistemico": "fato", "origem": "wiki", "rel": "relaciona"}, {"alvo": "projeto_cristaljson", "confianca": 0.5, "epistemico": "fato", "origem": "wiki", "rel": "relaciona"}, {"alvo": "tiffany_cabeca_broca", "confianca": 0.5, "epistemico": "fato", "origem": "wiki", "rel": "relaciona"}, {"alvo": "syscall", "confianca": 0.5, "epistemico": "fato", "origem": "wiki", "rel": "relaciona"}, {"alvo": "pow_metal_pre_fecho", "confianca": 0.5, "epistemico": "fato", "origem": "wiki", "rel": "relaciona"}, {"alvo": "pow_geometria", "confianca": 0.5, "epistemico": "fato", "origem": "wiki", "rel": "relaciona"}, {"alvo": "vida-morte", "confianca": 0.5, "epistemico": "fato", "origem": "wiki", "rel": "relaciona"}, {"alvo": "utilitarismo", "confianca": 0.5, "epistemico": "fato", "origem": "wiki", "rel": "relaciona"}, {"alvo": "regua_associativa_e_o_risco_de_vazamento", "confianca": 0.5, "epistemico": "fato", "origem": "wiki", "rel": "relaciona"}], "confianca": 1.0, "contraexemplos": [], "descricao": "Endereço de memória; indireção.", "epistemico": "fato", "exemplos": [], "id": "ponteiro", "memoria": "semantica", "meta": {"dominio": "programacao", "nivel": 6, "ordem": 6}, "origem": "curriculo", "palavras_chave": [], "sinonimos": ["pointer"], "tipo": "conceito", "titulo": "ponteiro"}
\section{ponteiro}
Endereço de memória; indireção.
\prov{sinónimos: pointer}
\prov{relações: depende processo; relaciona topologia; relaciona rust; relaciona por\_sessao\_conversadadostelemetriasessaojsonl; relaciona thread; relaciona virtualizacao; relaciona resposta\_a\_pergunta\_dificil; relaciona testes; relaciona projeto\_cristaljson; relaciona tiffany\_cabeca\_broca; relaciona syscall; relaciona pow\_metal\_pre\_fecho; relaciona pow\_geometria; relaciona vida-morte; relaciona utilitarismo; relaciona regua\_associativa\_e\_o\_risco\_de\_vazamento}
\prov{proveniência: curriculo \textperiodcentered{} fato \textperiodcentered{} confiança 1.0 \textperiodcentered{} domínio programacao \textperiodcentered{} história 19 versões no jornal}

%CRISTAL {"arestas": [{"alvo": "logica_digital", "confianca": 1.0, "epistemico": "fato", "origem": "wiki", "rel": "parte_de"}, {"alvo": "gato", "confianca": 1.0, "epistemico": "fato", "origem": "wiki", "rel": "relaciona"}, {"alvo": "design_opcode_algebrico", "confianca": 1.0, "epistemico": "fato", "origem": "wiki", "rel": "relaciona"}], "confianca": 1.0, "contraexemplos": [], "descricao": "O elemento mínimo do circuito: uma função booleana em hardware (AND, OR, XOR). O 'átomo' que, combinado, faz a CPU — e que o gato encarna como opcode.", "epistemico": "fato", "exemplos": [], "id": "porta_logica", "memoria": "semantica", "meta": {"dominio": "programacao", "fonte": ""}, "origem": "wiki", "palavras_chave": [], "sinonimos": [], "tipo": "conceito", "titulo": "Porta Lógica"}
\section{Porta Lógica}
O elemento mínimo do circuito: uma função booleana em hardware (AND, OR, XOR). O 'átomo' que, combinado, faz a CPU — e que o gato encarna como opcode.
\prov{relações: parte\_de logica\_digital; relaciona gato; relaciona design\_opcode\_algebrico}
\prov{proveniência: wiki \textperiodcentered{} fato \textperiodcentered{} confiança 1.0 \textperiodcentered{} domínio programacao \textperiodcentered{} história 4 versões no jornal}

%CRISTAL {"arestas": [{"alvo": "contabilidade_sem_livro", "confianca": 1.0, "epistemico": "fato", "origem": "wiki", "rel": "usa"}, {"alvo": "colapso_koch_tiffany", "confianca": 1.0, "epistemico": "fato", "origem": "wiki", "rel": "relaciona"}], "confianca": 1.0, "contraexemplos": [], "descricao": "Não há relógio de decisão: a 'porta' abre por colapso instantâneo quando, no mesmo instante, fechamento.fechou E parabola_global.ok E sincronizado_ok (E hodge se presente). O painel bloqueia num long-poll até esse colapso.", "epistemico": "inferencia", "exemplos": [], "id": "porta_mecanica", "memoria": "semantica", "meta": {"autor": "broca-so", "dominio": "operacao_so", "fonte": "goldchain/hodge_automatizar.py; cockpit_server.py"}, "origem": "wiki", "palavras_chave": [], "sinonimos": [], "tipo": "conceito", "titulo": "A Porta Mecânica (balanço = colapso da leitura)"}
\section{A Porta Mecânica (balanço = colapso da leitura)}
Não há relógio de decisão: a 'porta' abre por colapso instantâneo quando, no mesmo instante, fechamento.fechou E parabola\_global.ok E sincronizado\_ok (E hodge se presente). O painel bloqueia num long-poll até esse colapso.
\prov{relações: usa contabilidade\_sem\_livro; relaciona colapso\_koch\_tiffany}
\prov{proveniência: wiki \textperiodcentered{} goldchain/hodge\_automatizar.py; cockpit\_server.py \textperiodcentered{} inferencia \textperiodcentered{} confiança 1.0 \textperiodcentered{} domínio operacao\_so \textperiodcentered{} história 3 versões no jornal}

%CRISTAL {"arestas": [{"alvo": "unix", "confianca": 1.0, "epistemico": "fato", "origem": "wiki", "rel": "parte_de"}, {"alvo": "sistema_operacional", "confianca": 1.0, "epistemico": "fato", "origem": "wiki", "rel": "usa"}], "confianca": 1.0, "contraexemplos": [], "descricao": "O padrão que define a interface portável dos sistemas Unix (syscalls, shell, utilitários) — escreva uma vez, rode em qualquer Unix.", "epistemico": "fato", "exemplos": [], "id": "posix", "memoria": "semantica", "meta": {"dominio": "programacao", "fonte": ""}, "origem": "wiki", "palavras_chave": [], "sinonimos": [], "tipo": "conceito", "titulo": "POSIX"}
\section{POSIX}
O padrão que define a interface portável dos sistemas Unix (syscalls, shell, utilitários) — escreva uma vez, rode em qualquer Unix.
\prov{relações: parte\_de unix; usa sistema\_operacional}
\prov{proveniência: wiki \textperiodcentered{} fato \textperiodcentered{} confiança 1.0 \textperiodcentered{} domínio programacao \textperiodcentered{} história 3 versões no jornal}

%CRISTAL {"arestas": [{"alvo": "shell", "confianca": 1.0, "epistemico": "fato", "origem": "wiki", "rel": "especializa"}, {"alvo": "windows", "confianca": 1.0, "epistemico": "fato", "origem": "wiki", "rel": "parte_de"}, {"alvo": "filosofia_unix", "confianca": 1.0, "epistemico": "fato", "origem": "wiki", "rel": "contradiz"}], "confianca": 1.0, "contraexemplos": [], "descricao": "Microsoft: um shell que passa OBJETOS estruturados pelo pipe, não texto — a ruptura com a tradição Unix de fluxos de bytes.", "epistemico": "fato", "exemplos": [], "id": "powershell", "memoria": "semantica", "meta": {"autor": "Snover", "dominio": "programacao", "fonte": ""}, "origem": "wiki", "palavras_chave": [], "sinonimos": [], "tipo": "conceito", "titulo": "PowerShell"}
\section{PowerShell}
Microsoft: um shell que passa OBJETOS estruturados pelo pipe, não texto — a ruptura com a tradição Unix de fluxos de bytes.
\prov{relações: especializa shell; parte\_de windows; contradiz filosofia\_unix}
\prov{proveniência: wiki \textperiodcentered{} fato \textperiodcentered{} confiança 1.0 \textperiodcentered{} domínio programacao \textperiodcentered{} história 4 versões no jornal}

%CRISTAL {"arestas": [{"alvo": "gato_ula", "confianca": 1.0, "epistemico": "fato", "origem": "wiki", "rel": "usa"}, {"alvo": "cristal_memoria", "confianca": 1.0, "epistemico": "fato", "origem": "wiki", "rel": "usa"}, {"alvo": "design_primitivo_matriz", "confianca": 1.0, "epistemico": "fato", "origem": "wiki", "rel": "relaciona"}, {"alvo": "dsl_frontend_linguagem_gatos", "confianca": 1.0, "epistemico": "fato", "origem": "wiki", "rel": "usa"}, {"alvo": "design_laco_e_posicao", "confianca": 1.0, "epistemico": "fato", "origem": "wiki", "rel": "usa"}, {"alvo": "interpretador_nucleo_gatos", "confianca": 1.0, "epistemico": "fato", "origem": "wiki", "rel": "usa"}], "confianca": 1.0, "contraexemplos": [], "descricao": "O conjunto completo de primitivas, cada uma RODANDO no computador mineral (frontend DSL → assembler RAM → o gato): memória (atribuição, mem[] indireta = Turing-completude), ULA (+, −, gato n = a porta Aₙ, × via laço), controle (while, if/else, for, < > <= >= == !=), E-S (print), e o LAÇO=POSIÇÃO (o laço de gato = Aₙ^k, O(k)→O(log k), a transformada no compilador). Verificado 12/12 em linguagem/primitivas_minerais.py.", "epistemico": "fato", "exemplos": [], "id": "primitivas_minerais", "memoria": "semantica", "meta": {"dominio": "computacao", "fonte": "linguagem mineral"}, "origem": "wiki", "palavras_chave": [], "sinonimos": [], "tipo": "conceito", "titulo": "As Primitivas da Linguagem Mineral"}
\section{As Primitivas da Linguagem Mineral}
O conjunto completo de primitivas, cada uma RODANDO no computador mineral (frontend DSL \ensuremath{\to} assembler RAM \ensuremath{\to} o gato): memória (atribuição, mem[] indireta = Turing-completude), ULA (+, \ensuremath{-}, gato n = a porta A\ensuremath{_{n}}, $\times$ via laço), controle (while, if/else, for, < > <= >= == !=), E-S (print), e o LAÇO=POSIÇÃO (o laço de gato = A\ensuremath{_{n}}\^{}k, O(k)\ensuremath{\to}O(log k), a transformada no compilador). Verificado 12/12 em linguagem/primitivas\_minerais.py.
\prov{relações: usa gato\_ula; usa cristal\_memoria; relaciona design\_primitivo\_matriz; usa dsl\_frontend\_linguagem\_gatos; usa design\_laco\_e\_posicao; usa interpretador\_nucleo\_gatos}
\prov{proveniência: wiki \textperiodcentered{} linguagem mineral \textperiodcentered{} fato \textperiodcentered{} confiança 1.0 \textperiodcentered{} domínio computacao \textperiodcentered{} história 7 versões no jornal}

%CRISTAL {"arestas": [{"alvo": "criptografia", "confianca": 1.0, "epistemico": "fato", "origem": "wiki", "rel": "usa"}, {"alvo": "goldchain", "confianca": 1.0, "epistemico": "fato", "origem": "wiki", "rel": "relaciona"}], "confianca": 1.0, "contraexemplos": [], "descricao": "A segurança deve estar na CHAVE, não no segredo do algoritmo — o sistema pode ser público e ainda seguro. O oposto da segurança por obscuridade.", "epistemico": "inferencia", "exemplos": [], "id": "principio_de_kerckhoffs", "memoria": "semantica", "meta": {"autor": "Kerckhoffs", "dominio": "programacao", "fonte": ""}, "origem": "wiki", "palavras_chave": [], "sinonimos": [], "tipo": "conceito", "titulo": "Princípio de Kerckhoffs"}
\section{Princípio de Kerckhoffs}
A segurança deve estar na CHAVE, não no segredo do algoritmo — o sistema pode ser público e ainda seguro. O oposto da segurança por obscuridade.
\prov{relações: usa criptografia; relaciona goldchain}
\prov{proveniência: wiki \textperiodcentered{} inferencia \textperiodcentered{} confiança 1.0 \textperiodcentered{} domínio programacao \textperiodcentered{} história 3 versões no jornal}

%CRISTAL {"fusao": [{"arestas": [{"alvo": "thread", "confianca": 1.0, "epistemico": "fato", "origem": "curriculo", "rel": "especializa"}, {"alvo": "broca_os", "confianca": 1.0, "epistemico": "fato", "origem": "curriculo", "rel": "parte_de"}, {"alvo": "sistema_operacional", "confianca": 1.0, "epistemico": "fato", "origem": "wiki", "rel": "parte_de"}], "confianca": 1.0, "contraexemplos": [], "descricao": "Um programa em execução, com seu próprio espaço de memória e recursos — a unidade de isolamento que o SO agenda e protege.", "epistemico": "fato", "exemplos": [], "id": "processo", "memoria": "semantica", "meta": {"dominio": "programacao", "fonte": ""}, "origem": "wiki", "palavras_chave": [], "sinonimos": [], "tipo": "conceito", "titulo": "Processo"}, {"arestas": [{"alvo": "memoria_virtual", "confianca": 1.0, "epistemico": "fato", "origem": "curriculo", "rel": "parte_de"}, {"alvo": "broca_os", "confianca": 1.0, "epistemico": "fato", "origem": "curriculo", "rel": "parte_de"}, {"alvo": "topologia", "confianca": 0.5, "epistemico": "fato", "origem": "wiki", "rel": "relaciona"}, {"alvo": "rust", "confianca": 0.5, "epistemico": "fato", "origem": "wiki", "rel": "relaciona"}, {"alvo": "recursao", "confianca": 0.5, "epistemico": "fato", "origem": "wiki", "rel": "relaciona"}, {"alvo": "teste_ordem", "confianca": 0.5, "epistemico": "fato", "origem": "wiki", "rel": "relaciona"}, {"alvo": "testes", "confianca": 0.5, "epistemico": "fato", "origem": "wiki", "rel": "relaciona"}, {"alvo": "projeto_cristaljson", "confianca": 0.5, "epistemico": "fato", "origem": "wiki", "rel": "relaciona"}], "confianca": 1.0, "contraexemplos": [], "descricao": "Unidade de isolamento e scheduling.", "epistemico": "fato", "exemplos": [], "id": "processos", "memoria": "semantica", "meta": {"dominio": "sistemas", "nivel": 6, "ordem": 0}, "origem": "curriculo", "palavras_chave": [], "sinonimos": [], "tipo": "conceito", "titulo": "processos"}], "id": "processo", "tipo": "conceito"}
\section{Processo}
Um programa em execução, com seu próprio espaço de memória e recursos — a unidade de isolamento que o SO agenda e protege.
\prov{relações: especializa thread; parte\_de broca\_os; parte\_de sistema\_operacional}
\prov{proveniência: wiki \textperiodcentered{} fato \textperiodcentered{} confiança 1.0 \textperiodcentered{} domínio programacao \textperiodcentered{} história 8 versões no jornal}
\prov{a segunda face --- processos:}
Unidade de isolamento e scheduling.
\prov{fusão de curadoria (julgada): absorve processos --- as duas faces acima, intactas no registo}

%CRISTAL {"arestas": [{"alvo": "ciencia_da_computacao", "confianca": 1.0, "epistemico": "fato", "origem": "wiki", "rel": "parte_de"}, {"alvo": "algoritmo", "confianca": 1.0, "epistemico": "fato", "origem": "wiki", "rel": "usa"}], "confianca": 1.0, "contraexemplos": [], "descricao": "A arte e a técnica de instruir máquinas — expressar um algoritmo numa linguagem que o computador executa.", "epistemico": "fato", "exemplos": [], "id": "programacao", "memoria": "semantica", "meta": {"dominio": "programacao", "fonte": ""}, "origem": "wiki", "palavras_chave": [], "sinonimos": [], "tipo": "conceito", "titulo": "Programação"}
\section{Programação}
A arte e a técnica de instruir máquinas — expressar um algoritmo numa linguagem que o computador executa.
\prov{relações: parte\_de ciencia\_da\_computacao; usa algoritmo}
\prov{proveniência: wiki \textperiodcentered{} fato \textperiodcentered{} confiança 1.0 \textperiodcentered{} domínio programacao \textperiodcentered{} história 6 versões no jornal}

%CRISTAL {"arestas": [{"alvo": "paradigma_de_programacao", "confianca": 1.0, "epistemico": "fato", "origem": "wiki", "rel": "parte_de"}], "confianca": 1.0, "contraexemplos": [], "descricao": "Estruturar o programa em tarefas que progridem ao mesmo tempo — o desafio de compartilhar estado sem corromper.", "epistemico": "fato", "exemplos": [], "id": "programacao_concorrente", "memoria": "semantica", "meta": {"dominio": "programacao", "fonte": ""}, "origem": "wiki", "palavras_chave": [], "sinonimos": [], "tipo": "conceito", "titulo": "Programação Concorrente"}
\section{Programação Concorrente}
Estruturar o programa em tarefas que progridem ao mesmo tempo — o desafio de compartilhar estado sem corromper.
\prov{relações: parte\_de paradigma\_de\_programacao}
\prov{proveniência: wiki \textperiodcentered{} fato \textperiodcentered{} confiança 1.0 \textperiodcentered{} domínio programacao \textperiodcentered{} história 4 versões no jornal}

%CRISTAL {"arestas": [{"alvo": "programacao", "confianca": 1.0, "epistemico": "fato", "origem": "wiki", "rel": "parte_de"}, {"alvo": "c", "confianca": 1.0, "epistemico": "fato", "origem": "wiki", "rel": "usa"}, {"alvo": "gerenciamento_de_memoria", "confianca": 1.0, "epistemico": "fato", "origem": "wiki", "rel": "usa"}, {"alvo": "linguagem_assembly", "confianca": 1.0, "epistemico": "fato", "origem": "wiki", "rel": "usa"}], "confianca": 1.0, "contraexemplos": [], "descricao": "Escrever o software que fica perto do metal — kernels, drivers, runtimes, shells; controle de memória e desempenho acima da conveniência.", "epistemico": "fato", "exemplos": [], "id": "programacao_de_sistemas", "memoria": "semantica", "meta": {"dominio": "programacao", "fonte": ""}, "origem": "wiki", "palavras_chave": [], "sinonimos": [], "tipo": "conceito", "titulo": "Programação de Sistemas"}
\section{Programação de Sistemas}
Escrever o software que fica perto do metal — kernels, drivers, runtimes, shells; controle de memória e desempenho acima da conveniência.
\prov{relações: parte\_de programacao; usa c; usa gerenciamento\_de\_memoria; usa linguagem\_assembly}
\prov{proveniência: wiki \textperiodcentered{} fato \textperiodcentered{} confiança 1.0 \textperiodcentered{} domínio programacao \textperiodcentered{} história 5 versões no jornal}

%CRISTAL {"arestas": [{"alvo": "paradigma_de_programacao", "confianca": 1.0, "epistemico": "fato", "origem": "wiki", "rel": "parte_de"}], "confianca": 1.0, "contraexemplos": [], "descricao": "Descrever O QUE se quer, não o passo a passo — o motor decide o como (SQL, HTML, Prolog).", "epistemico": "fato", "exemplos": [], "id": "programacao_declarativa", "memoria": "semantica", "meta": {"dominio": "programacao", "fonte": ""}, "origem": "wiki", "palavras_chave": [], "sinonimos": [], "tipo": "conceito", "titulo": "Programação Declarativa"}
\section{Programação Declarativa}
Descrever O QUE se quer, não o passo a passo — o motor decide o como (SQL, HTML, Prolog).
\prov{relações: parte\_de paradigma\_de\_programacao}
\prov{proveniência: wiki \textperiodcentered{} fato \textperiodcentered{} confiança 1.0 \textperiodcentered{} domínio programacao \textperiodcentered{} história 4 versões no jornal}

%CRISTAL {"arestas": [{"alvo": "programacao_imperativa", "confianca": 1.0, "epistemico": "fato", "origem": "wiki", "rel": "especializa"}], "confianca": 1.0, "contraexemplos": [], "descricao": "Dijkstra: só sequência, seleção e repetição — sem goto. A base disciplinada do código imperativo.", "epistemico": "fato", "exemplos": [], "id": "programacao_estruturada", "memoria": "semantica", "meta": {"autor": "Dijkstra", "dominio": "programacao", "fonte": ""}, "origem": "wiki", "palavras_chave": [], "sinonimos": [], "tipo": "conceito", "titulo": "Programação Estruturada"}
\section{Programação Estruturada}
Dijkstra: só sequência, seleção e repetição — sem goto. A base disciplinada do código imperativo.
\prov{relações: especializa programacao\_imperativa}
\prov{proveniência: wiki \textperiodcentered{} fato \textperiodcentered{} confiança 1.0 \textperiodcentered{} domínio programacao \textperiodcentered{} história 4 versões no jornal}

%CRISTAL {"arestas": [{"alvo": "paradigma_de_programacao", "confianca": 1.0, "epistemico": "fato", "origem": "wiki", "rel": "parte_de"}, {"alvo": "calculo_lambda", "confianca": 1.0, "epistemico": "fato", "origem": "wiki", "rel": "usa"}], "confianca": 1.0, "contraexemplos": [], "descricao": "Computar por avaliação de funções puras, sem estado mutável — o λ-cálculo aplicado; funções são valores.", "epistemico": "fato", "exemplos": [], "id": "programacao_funcional", "memoria": "semantica", "meta": {"dominio": "programacao", "fonte": ""}, "origem": "wiki", "palavras_chave": [], "sinonimos": [], "tipo": "conceito", "titulo": "Programação Funcional"}
\section{Programação Funcional}
Computar por avaliação de funções puras, sem estado mutável — o \ensuremath{\lambda}-cálculo aplicado; funções são valores.
\prov{relações: parte\_de paradigma\_de\_programacao; usa calculo\_lambda}
\prov{proveniência: wiki \textperiodcentered{} fato \textperiodcentered{} confiança 1.0 \textperiodcentered{} domínio programacao \textperiodcentered{} história 6 versões no jornal}

%CRISTAL {"arestas": [{"alvo": "paradigma_de_programacao", "confianca": 1.0, "epistemico": "fato", "origem": "wiki", "rel": "parte_de"}, {"alvo": "arquitetura_von_neumann", "confianca": 1.0, "epistemico": "fato", "origem": "wiki", "rel": "usa"}], "confianca": 1.0, "contraexemplos": [], "descricao": "Descrever COMO fazer: uma sequência de comandos que muta o estado (variáveis, memória). O modelo de von Neumann.", "epistemico": "fato", "exemplos": [], "id": "programacao_imperativa", "memoria": "semantica", "meta": {"dominio": "programacao", "fonte": ""}, "origem": "wiki", "palavras_chave": [], "sinonimos": [], "tipo": "conceito", "titulo": "Programação Imperativa"}
\section{Programação Imperativa}
Descrever COMO fazer: uma sequência de comandos que muta o estado (variáveis, memória). O modelo de von Neumann.
\prov{relações: parte\_de paradigma\_de\_programacao; usa arquitetura\_von\_neumann}
\prov{proveniência: wiki \textperiodcentered{} fato \textperiodcentered{} confiança 1.0 \textperiodcentered{} domínio programacao \textperiodcentered{} história 6 versões no jornal}

%CRISTAL {"arestas": [{"alvo": "programacao_declarativa", "confianca": 1.0, "epistemico": "fato", "origem": "wiki", "rel": "especializa"}, {"alvo": "logica_de_predicados", "confianca": 1.0, "epistemico": "fato", "origem": "wiki", "rel": "usa"}, {"alvo": "logica", "confianca": 1.0, "epistemico": "fato", "origem": "wiki", "rel": "relaciona"}], "confianca": 1.0, "contraexemplos": [], "descricao": "Prolog: o programa é um conjunto de fatos e regras; rodar é PROVAR uma consulta. Lógica de predicados executável.", "epistemico": "fato", "exemplos": [], "id": "programacao_logica", "memoria": "semantica", "meta": {"autor": "Kowalski", "dominio": "programacao", "fonte": ""}, "origem": "wiki", "palavras_chave": [], "sinonimos": [], "tipo": "conceito", "titulo": "Programação Lógica"}
\section{Programação Lógica}
Prolog: o programa é um conjunto de fatos e regras; rodar é PROVAR uma consulta. Lógica de predicados executável.
\prov{relações: especializa programacao\_declarativa; usa logica\_de\_predicados; relaciona logica}
\prov{proveniência: wiki \textperiodcentered{} fato \textperiodcentered{} confiança 1.0 \textperiodcentered{} domínio programacao \textperiodcentered{} história 7 versões no jornal}

%CRISTAL {"arestas": [{"alvo": "paradigma_de_programacao", "confianca": 1.0, "epistemico": "fato", "origem": "wiki", "rel": "parte_de"}], "confianca": 1.0, "contraexemplos": [], "descricao": "Organizar o programa em objetos que unem estado e comportamento — encapsulamento, herança, polimorfismo (Alan Kay).", "epistemico": "fato", "exemplos": [], "id": "programacao_orientada_objetos", "memoria": "semantica", "meta": {"autor": "Alan Kay", "dominio": "programacao", "fonte": ""}, "origem": "wiki", "palavras_chave": [], "sinonimos": [], "tipo": "conceito", "titulo": "Programação Orientada a Objetos"}
\section{Programação Orientada a Objetos}
Organizar o programa em objetos que unem estado e comportamento — encapsulamento, herança, polimorfismo (Alan Kay).
\prov{relações: parte\_de paradigma\_de\_programacao}
\prov{proveniência: wiki \textperiodcentered{} fato \textperiodcentered{} confiança 1.0 \textperiodcentered{} domínio programacao \textperiodcentered{} história 4 versões no jornal}

%CRISTAL {"arestas": [{"alvo": "programacao_funcional", "confianca": 1.0, "epistemico": "fato", "origem": "wiki", "rel": "especializa"}], "confianca": 1.0, "contraexemplos": [], "descricao": "Compor funções sem nomear os argumentos (APL, Haskell) — o programa é a álgebra das combinações, não a receita passo a passo.", "epistemico": "fato", "exemplos": [], "id": "programacao_tacita", "memoria": "semantica", "meta": {"dominio": "programacao", "fonte": ""}, "origem": "wiki", "palavras_chave": [], "sinonimos": [], "tipo": "conceito", "titulo": "Programação Tácita (Point-Free)"}
\section{Programação Tácita (Point-Free)}
Compor funções sem nomear os argumentos (APL, Haskell) — o programa é a álgebra das combinações, não a receita passo a passo.
\prov{relações: especializa programacao\_funcional}
\prov{proveniência: wiki \textperiodcentered{} fato \textperiodcentered{} confiança 1.0 \textperiodcentered{} domínio programacao \textperiodcentered{} história 4 versões no jornal}

%CRISTAL {"arestas": [{"alvo": "pilha_tcp_ip", "confianca": 1.0, "epistemico": "fato", "origem": "wiki", "rel": "parte_de"}, {"alvo": "comutacao_de_pacotes", "confianca": 1.0, "epistemico": "fato", "origem": "wiki", "rel": "usa"}], "confianca": 1.0, "contraexemplos": [], "descricao": "Endereça e roteia pacotes entre redes, best-effort e sem garantia — a camada que faz a internet ser uma só.", "epistemico": "fato", "exemplos": [], "id": "protocolo_ip", "memoria": "semantica", "meta": {"autor": "Cerf/Kahn", "dominio": "programacao", "fonte": ""}, "origem": "wiki", "palavras_chave": [], "sinonimos": [], "tipo": "conceito", "titulo": "IP (Protocolo de Internet)"}
\section{IP (Protocolo de Internet)}
Endereça e roteia pacotes entre redes, best-effort e sem garantia — a camada que faz a internet ser uma só.
\prov{relações: parte\_de pilha\_tcp\_ip; usa comutacao\_de\_pacotes}
\prov{proveniência: wiki \textperiodcentered{} fato \textperiodcentered{} confiança 1.0 \textperiodcentered{} domínio programacao \textperiodcentered{} história 6 versões no jornal}

%CRISTAL {"arestas": [{"alvo": "pilha_tcp_ip", "confianca": 1.0, "epistemico": "fato", "origem": "wiki", "rel": "parte_de"}, {"alvo": "protocolo_ip", "confianca": 1.0, "epistemico": "fato", "origem": "wiki", "rel": "usa"}], "confianca": 1.0, "contraexemplos": [], "descricao": "Sobre o IP não-confiável, entrega um fluxo ORDENADO e confiável (retransmissão, controle de congestão) — a ilusão de um cano limpo.", "epistemico": "fato", "exemplos": [], "id": "protocolo_tcp", "memoria": "semantica", "meta": {"dominio": "programacao", "fonte": ""}, "origem": "wiki", "palavras_chave": [], "sinonimos": [], "tipo": "conceito", "titulo": "TCP"}
\section{TCP}
Sobre o IP não-confiável, entrega um fluxo ORDENADO e confiável (retransmissão, controle de congestão) — a ilusão de um cano limpo.
\prov{relações: parte\_de pilha\_tcp\_ip; usa protocolo\_ip}
\prov{proveniência: wiki \textperiodcentered{} fato \textperiodcentered{} confiança 1.0 \textperiodcentered{} domínio programacao \textperiodcentered{} história 6 versões no jornal}

%CRISTAL {"arestas": [{"alvo": "pilha_tcp_ip", "confianca": 1.0, "epistemico": "fato", "origem": "wiki", "rel": "parte_de"}, {"alvo": "protocolo_tcp", "confianca": 1.0, "epistemico": "fato", "origem": "wiki", "rel": "contradiz"}], "confianca": 1.0, "contraexemplos": [], "descricao": "Transporte sem conexão nem garantia — rápido e leve, para quando perder um pacote importa menos que a latência (vídeo, jogos, DNS).", "epistemico": "fato", "exemplos": [], "id": "protocolo_udp", "memoria": "semantica", "meta": {"dominio": "programacao", "fonte": ""}, "origem": "wiki", "palavras_chave": [], "sinonimos": [], "tipo": "conceito", "titulo": "UDP"}
\section{UDP}
Transporte sem conexão nem garantia — rápido e leve, para quando perder um pacote importa menos que a latência (vídeo, jogos, DNS).
\prov{relações: parte\_de pilha\_tcp\_ip; contradiz protocolo\_tcp}
\prov{proveniência: wiki \textperiodcentered{} fato \textperiodcentered{} confiança 1.0 \textperiodcentered{} domínio programacao \textperiodcentered{} história 6 versões no jornal}

%CRISTAL {"arestas": [{"alvo": "sass_gpu_assembly", "confianca": 1.0, "epistemico": "fato", "origem": "wiki", "rel": "relaciona"}], "confianca": 1.0, "contraexemplos": [], "descricao": "PTX (Parallel Thread eXecution) é a ISA VIRTUAL da NVIDIA: um assembly ABSTRATO, forward-compatível (registradores infinitos %r, sem custo de arquitetura). O compilador (nvcc/o driver GLSL) emite PTX; o ptxas (o assembler) compila PTX→SASS para a arquitetura-alvo (o -arch sm_75, sm_80, sm_89, sm_90). PTX é a camada portável (o mesmo PTX roda em qualquer GPU, recompilado JIT pelo driver para o SASS certo). Análogo à DSL de gatos: uma linguagem intermediária que aterra no metal. Fato.", "epistemico": "fato", "exemplos": [], "id": "ptx_isa_virtual", "memoria": "semantica", "meta": {"dominio": "sass_gpu", "fonte": "SASS/GPU (o metal)"}, "origem": "wiki", "palavras_chave": [], "sinonimos": [], "tipo": "conceito", "titulo": "PTX — a ISA virtual (entre o DSL/GLSL e o SASS)"}
\section{PTX — a ISA virtual (entre o DSL/GLSL e o SASS)}
PTX (Parallel Thread eXecution) é a ISA VIRTUAL da NVIDIA: um assembly ABSTRATO, forward-compatível (registradores infinitos \%r, sem custo de arquitetura). O compilador (nvcc/o driver GLSL) emite PTX; o ptxas (o assembler) compila PTX\ensuremath{\to}SASS para a arquitetura-alvo (o -arch sm\_75, sm\_80, sm\_89, sm\_90). PTX é a camada portável (o mesmo PTX roda em qualquer GPU, recompilado JIT pelo driver para o SASS certo). Análogo à DSL de gatos: uma linguagem intermediária que aterra no metal. Fato.
\prov{relações: relaciona sass\_gpu\_assembly}
\prov{proveniência: wiki \textperiodcentered{} SASS/GPU (o metal) \textperiodcentered{} fato \textperiodcentered{} confiança 1.0 \textperiodcentered{} domínio sass\_gpu \textperiodcentered{} história 14 versões no jornal}

%CRISTAL {"arestas": [{"alvo": "haskell", "confianca": 1.0, "epistemico": "fato", "origem": "wiki", "rel": "especializa"}, {"alvo": "javascript", "confianca": 1.0, "epistemico": "fato", "origem": "wiki", "rel": "usa"}, {"alvo": "programacao_funcional", "confianca": 1.0, "epistemico": "fato", "origem": "wiki", "rel": "usa"}], "confianca": 1.0, "contraexemplos": [], "descricao": "Um Haskell que compila para JavaScript: tipagem forte, pureza e efeitos explícitos na web — o rigor funcional no navegador.", "epistemico": "fato", "exemplos": [], "id": "purescript", "memoria": "semantica", "meta": {"autor": "Freeman", "dominio": "programacao", "fonte": ""}, "origem": "wiki", "palavras_chave": [], "sinonimos": [], "tipo": "conceito", "titulo": "PureScript"}
\section{PureScript}
Um Haskell que compila para JavaScript: tipagem forte, pureza e efeitos explícitos na web — o rigor funcional no navegador.
\prov{relações: especializa haskell; usa javascript; usa programacao\_funcional}
\prov{proveniência: wiki \textperiodcentered{} fato \textperiodcentered{} confiança 1.0 \textperiodcentered{} domínio programacao \textperiodcentered{} história 4 versões no jornal}

%CRISTAL {"arestas": [{"alvo": "compilador", "confianca": 1.0, "epistemico": "fato", "origem": "curriculo", "rel": "referencia"}, {"alvo": "transcricao", "confianca": 0.5, "epistemico": "fato", "origem": "wiki", "rel": "relaciona"}, {"alvo": "sessao_interativa", "confianca": 0.5, "epistemico": "fato", "origem": "wiki", "rel": "relaciona"}, {"alvo": "readme", "confianca": 0.5, "epistemico": "fato", "origem": "wiki", "rel": "relaciona"}, {"alvo": "teste_ordem", "confianca": 0.5, "epistemico": "fato", "origem": "wiki", "rel": "relaciona"}, {"alvo": "virtualizacao", "confianca": 0.5, "epistemico": "fato", "origem": "wiki", "rel": "relaciona"}, {"alvo": "word", "confianca": 0.5, "epistemico": "fato", "origem": "wiki", "rel": "relaciona"}, {"alvo": "vontade_de_poder", "confianca": 0.5, "epistemico": "fato", "origem": "wiki", "rel": "relaciona"}, {"alvo": "while", "confianca": 0.5, "epistemico": "fato", "origem": "wiki", "rel": "relaciona"}, {"alvo": "linguagem_de_programacao", "confianca": 1.0, "epistemico": "fato", "origem": "wiki", "rel": "parte_de"}, {"alvo": "tipagem_estatica_dinamica", "confianca": 1.0, "epistemico": "fato", "origem": "wiki", "rel": "usa"}], "confianca": 1.0, "contraexemplos": [], "descricao": "van Rossum: legibilidade acima de tudo, tipagem dinâmica, 'baterias inclusas'. A língua franca da ciência de dados e do ensino.", "epistemico": "fato", "exemplos": [], "id": "python", "memoria": "semantica", "meta": {"autor": "van Rossum", "dominio": "programacao", "fonte": ""}, "origem": "wiki", "palavras_chave": [], "sinonimos": [], "tipo": "conceito", "titulo": "Python"}
\section{Python}
van Rossum: legibilidade acima de tudo, tipagem dinâmica, 'baterias inclusas'. A língua franca da ciência de dados e do ensino.
\prov{relações: referencia compilador; relaciona transcricao; relaciona sessao\_interativa; relaciona readme; relaciona teste\_ordem; relaciona virtualizacao; relaciona word; relaciona vontade\_de\_poder; relaciona while; parte\_de linguagem\_de\_programacao; usa tipagem\_estatica\_dinamica}
\prov{proveniência: wiki \textperiodcentered{} fato \textperiodcentered{} confiança 1.0 \textperiodcentered{} domínio programacao \textperiodcentered{} história 24 versões no jornal}

%CRISTAL {"arestas": [{"alvo": "linguagem_de_programacao", "confianca": 1.0, "epistemico": "fato", "origem": "wiki", "rel": "parte_de"}, {"alvo": "teoria_da_probabilidade", "confianca": 1.0, "epistemico": "fato", "origem": "wiki", "rel": "usa"}], "confianca": 1.0, "contraexemplos": [], "descricao": "A linguagem da estatística: manipulação de dados, modelagem e gráficos — o dialeto nativo dos estatísticos.", "epistemico": "fato", "exemplos": [], "id": "r_linguagem", "memoria": "semantica", "meta": {"dominio": "programacao", "fonte": ""}, "origem": "wiki", "palavras_chave": [], "sinonimos": [], "tipo": "conceito", "titulo": "R"}
\section{R}
A linguagem da estatística: manipulação de dados, modelagem e gráficos — o dialeto nativo dos estatísticos.
\prov{relações: parte\_de linguagem\_de\_programacao; usa teoria\_da\_probabilidade}
\prov{proveniência: wiki \textperiodcentered{} fato \textperiodcentered{} confiança 1.0 \textperiodcentered{} domínio programacao \textperiodcentered{} história 5 versões no jornal}

%CRISTAL {"arestas": [{"alvo": "consenso_distribuido", "confianca": 1.0, "epistemico": "fato", "origem": "wiki", "rel": "especializa"}], "confianca": 1.0, "contraexemplos": [], "descricao": "Um algoritmo de consenso desenhado para ser COMPREENSÍVEL (líder, log replicado, eleição) — o Paxos que engenheiros de fato usam.", "epistemico": "fato", "exemplos": [], "id": "raft", "memoria": "semantica", "meta": {"autor": "Ongaro", "dominio": "programacao", "fonte": ""}, "origem": "wiki", "palavras_chave": [], "sinonimos": [], "tipo": "conceito", "titulo": "Raft"}
\section{Raft}
Um algoritmo de consenso desenhado para ser COMPREENSÍVEL (líder, log replicado, eleição) — o Paxos que engenheiros de fato usam.
\prov{relações: especializa consenso\_distribuido}
\prov{proveniência: wiki \textperiodcentered{} fato \textperiodcentered{} confiança 1.0 \textperiodcentered{} domínio programacao \textperiodcentered{} história 4 versões no jornal}

%CRISTAL {"arestas": [{"alvo": "metal_do_so", "confianca": 1.0, "epistemico": "fato", "origem": "wiki", "rel": "parte_de"}, {"alvo": "operacoes_reservadas_no_metal", "confianca": 1.0, "epistemico": "fato", "origem": "wiki", "rel": "usa"}, {"alvo": "turmalina_paraiba_ancora", "confianca": 1.0, "epistemico": "fato", "origem": "wiki", "rel": "usa"}, {"alvo": "design_system_cristalchain", "confianca": 1.0, "epistemico": "fato", "origem": "wiki", "rel": "usa"}, {"alvo": "area_optica_gema", "confianca": 1.0, "epistemico": "fato", "origem": "wiki", "rel": "parte_de"}], "confianca": 1.0, "contraexemplos": [], "descricao": "A gema (turmalina lapidada) renderizada pela FÍSICA DA LUZ no metal do SO (linguagem/raytracer.py): refração de Snell (n=1.62), reflexão interna total, Fresnel (o brilho das facetas), absorção Beer-Lambert (o azul-verde), inclinação 3/4, cantos suaves (supersampling), e a SOMBRA + CÁUSTICA no chão. A fonte de luz padrão da cena é o nosso SOL — uma estrela amarela (corpo negro a 5778 K, Planck): a cor amarelo-branca, o disco solar, a distribuição espectral. É operação reservada do metal ('gema'/'gem', chamável em qualquer idioma). A física da pedra virando o logo.", "epistemico": "fato", "exemplos": [], "id": "raytracer_no_metal_e_o_sol", "memoria": "semantica", "meta": {"dominio": "sistema", "fonte": "pipeline gráfico"}, "origem": "wiki", "palavras_chave": [], "sinonimos": [], "tipo": "conceito", "titulo": "O Raytracer no Metal (e o Sol como luz padrão)"}
\section{O Raytracer no Metal (e o Sol como luz padrão)}
A gema (turmalina lapidada) renderizada pela FÍSICA DA LUZ no metal do SO (linguagem/raytracer.py): refração de Snell (n=1.62), reflexão interna total, Fresnel (o brilho das facetas), absorção Beer-Lambert (o azul-verde), inclinação 3/4, cantos suaves (supersampling), e a SOMBRA + CÁUSTICA no chão. A fonte de luz padrão da cena é o nosso SOL — uma estrela amarela (corpo negro a 5778 K, Planck): a cor amarelo-branca, o disco solar, a distribuição espectral. É operação reservada do metal ('gema'/'gem', chamável em qualquer idioma). A física da pedra virando o logo.
\prov{relações: parte\_de metal\_do\_so; usa operacoes\_reservadas\_no\_metal; usa turmalina\_paraiba\_ancora; usa design\_system\_cristalchain; parte\_de area\_optica\_gema}
\prov{proveniência: wiki \textperiodcentered{} pipeline gráfico \textperiodcentered{} fato \textperiodcentered{} confiança 1.0 \textperiodcentered{} domínio sistema \textperiodcentered{} história 13 versões no jornal}

%CRISTAL {"arestas": [{"alvo": "framework_frontend", "confianca": 1.0, "epistemico": "fato", "origem": "wiki", "rel": "especializa"}, {"alvo": "componente_ui", "confianca": 1.0, "epistemico": "fato", "origem": "wiki", "rel": "usa"}, {"alvo": "dom_virtual", "confianca": 1.0, "epistemico": "fato", "origem": "wiki", "rel": "usa"}, {"alvo": "jsx", "confianca": 1.0, "epistemico": "fato", "origem": "wiki", "rel": "usa"}, {"alvo": "hook_react", "confianca": 1.0, "epistemico": "fato", "origem": "wiki", "rel": "usa"}, {"alvo": "redux", "confianca": 1.0, "epistemico": "fato", "origem": "wiki", "rel": "usa"}, {"alvo": "zustand", "confianca": 1.0, "epistemico": "fato", "origem": "wiki", "rel": "usa"}], "confianca": 1.0, "contraexemplos": [], "descricao": "Meta: a biblioteca que popularizou os componentes e o DOM virtual — UI como função do estado. O centro do ecossistema frontend moderno.", "epistemico": "fato", "exemplos": [], "id": "react", "memoria": "semantica", "meta": {"autor": "Meta", "dominio": "programacao", "fonte": ""}, "origem": "wiki", "palavras_chave": [], "sinonimos": [], "tipo": "conceito", "titulo": "React"}
\section{React}
Meta: a biblioteca que popularizou os componentes e o DOM virtual — UI como função do estado. O centro do ecossistema frontend moderno.
\prov{relações: especializa framework\_frontend; usa componente\_ui; usa dom\_virtual; usa jsx; usa hook\_react; usa redux; usa zustand}
\prov{proveniência: wiki \textperiodcentered{} fato \textperiodcentered{} confiança 1.0 \textperiodcentered{} domínio programacao \textperiodcentered{} história 16 versões no jornal}

%CRISTAL {"arestas": [{"alvo": "framework_frontend", "confianca": 1.0, "epistemico": "fato", "origem": "wiki", "rel": "usa"}, {"alvo": "gerenciamento_de_estado", "confianca": 1.0, "epistemico": "fato", "origem": "wiki", "rel": "usa"}], "confianca": 1.0, "contraexemplos": [], "descricao": "A UI se atualiza sozinha quando o estado muda — declara-se a relação, não o passo de atualização. O padrão observador na interface.", "epistemico": "inferencia", "exemplos": [], "id": "reatividade", "memoria": "semantica", "meta": {"dominio": "programacao", "fonte": ""}, "origem": "wiki", "palavras_chave": [], "sinonimos": [], "tipo": "conceito", "titulo": "Reatividade"}
\section{Reatividade}
A UI se atualiza sozinha quando o estado muda — declara-se a relação, não o passo de atualização. O padrão observador na interface.
\prov{relações: usa framework\_frontend; usa gerenciamento\_de\_estado}
\prov{proveniência: wiki \textperiodcentered{} inferencia \textperiodcentered{} confiança 1.0 \textperiodcentered{} domínio programacao \textperiodcentered{} história 6 versões no jornal}

%CRISTAL {"arestas": [{"alvo": "pilha", "confianca": 1.0, "epistemico": "fato", "origem": "curriculo", "rel": "especializa"}, {"alvo": "testes", "confianca": 0.5, "epistemico": "fato", "origem": "wiki", "rel": "relaciona"}, {"alvo": "regua_associativa_e_o_risco_de_vazamento", "confianca": 0.5, "epistemico": "fato", "origem": "wiki", "rel": "relaciona"}, {"alvo": "topologia", "confianca": 0.5, "epistemico": "fato", "origem": "wiki", "rel": "relaciona"}, {"alvo": "respiracao_dos_programas", "confianca": 0.5, "epistemico": "fato", "origem": "wiki", "rel": "relaciona"}, {"alvo": "utilitarismo", "confianca": 0.5, "epistemico": "fato", "origem": "wiki", "rel": "relaciona"}, {"alvo": "resposta_a_pergunta_dificil", "confianca": 0.5, "epistemico": "fato", "origem": "wiki", "rel": "relaciona"}, {"alvo": "rust", "confianca": 0.5, "epistemico": "fato", "origem": "wiki", "rel": "relaciona"}, {"alvo": "teste_ordem", "confianca": 0.5, "epistemico": "fato", "origem": "wiki", "rel": "relaciona"}, {"alvo": "ruido", "confianca": 0.5, "epistemico": "fato", "origem": "wiki", "rel": "relaciona"}, {"alvo": "setor_algebrico_descontinuidades_no_fecho", "confianca": 0.5, "epistemico": "fato", "origem": "wiki", "rel": "relaciona"}, {"alvo": "sessao_interativa", "confianca": 0.5, "epistemico": "fato", "origem": "wiki", "rel": "relaciona"}, {"alvo": "thread", "confianca": 0.5, "epistemico": "fato", "origem": "wiki", "rel": "relaciona"}, {"alvo": "scheduler", "confianca": 0.5, "epistemico": "fato", "origem": "wiki", "rel": "relaciona"}], "confianca": 1.0, "contraexemplos": [], "descricao": "Função que chama a si mesma; requer caso base.", "epistemico": "fato", "exemplos": [], "id": "recursao", "memoria": "semantica", "meta": {"dominio": "programacao", "nivel": 6, "ordem": 3}, "origem": "curriculo", "palavras_chave": [], "sinonimos": [], "tipo": "conceito", "titulo": "recursão"}
\section{recursão}
Função que chama a si mesma; requer caso base.
\prov{relações: especializa pilha; relaciona testes; relaciona regua\_associativa\_e\_o\_risco\_de\_vazamento; relaciona topologia; relaciona respiracao\_dos\_programas; relaciona utilitarismo; relaciona resposta\_a\_pergunta\_dificil; relaciona rust; relaciona teste\_ordem; relaciona ruido; relaciona setor\_algebrico\_descontinuidades\_no\_fecho; relaciona sessao\_interativa; relaciona thread; relaciona scheduler}
\prov{proveniência: curriculo \textperiodcentered{} fato \textperiodcentered{} confiança 1.0 \textperiodcentered{} domínio programacao \textperiodcentered{} história 17 versões no jornal}

%CRISTAL {"arestas": [], "confianca": 1.0, "contraexemplos": [], "descricao": "Máquinas ligadas que trocam dados por protocolos — a base de tudo que é distribuído, da LAN à internet.", "epistemico": "fato", "exemplos": [], "id": "rede_de_computadores", "memoria": "semantica", "meta": {"dominio": "programacao", "fonte": ""}, "origem": "wiki", "palavras_chave": [], "sinonimos": [], "tipo": "conceito", "titulo": "Rede de Computadores"}
\section{Rede de Computadores}
Máquinas ligadas que trocam dados por protocolos — a base de tudo que é distribuído, da LAN à internet.
\prov{proveniência: wiki \textperiodcentered{} fato \textperiodcentered{} confiança 1.0 \textperiodcentered{} domínio programacao \textperiodcentered{} história 2 versões no jornal}

%CRISTAL {"arestas": [{"alvo": "aprendizado_de_maquina", "confianca": 1.0, "epistemico": "fato", "origem": "wiki", "rel": "usa"}, {"alvo": "hopfield", "confianca": 1.0, "epistemico": "fato", "origem": "wiki", "rel": "explica"}, {"alvo": "plasticidade_hebbiana", "confianca": 1.0, "epistemico": "fato", "origem": "wiki", "rel": "usa"}, {"alvo": "neuronio", "confianca": 1.0, "epistemico": "fato", "origem": "wiki", "rel": "relaciona"}, {"alvo": "a_rede_neural_de_ouro", "confianca": 1.0, "epistemico": "fato", "origem": "wiki", "rel": "relaciona"}], "confianca": 1.0, "contraexemplos": [], "descricao": "Camadas de unidades que somam entradas ponderadas e disparam por não-linearidade — inspiradas no neurônio; a rede de Hopfield é uma delas.", "epistemico": "fato", "exemplos": [], "id": "rede_neural_artificial", "memoria": "semantica", "meta": {"dominio": "programacao", "fonte": ""}, "origem": "wiki", "palavras_chave": [], "sinonimos": [], "tipo": "conceito", "titulo": "Rede Neural Artificial"}
\section{Rede Neural Artificial}
Camadas de unidades que somam entradas ponderadas e disparam por não-linearidade — inspiradas no neurônio; a rede de Hopfield é uma delas.
\prov{relações: usa aprendizado\_de\_maquina; explica hopfield; usa plasticidade\_hebbiana; relaciona neuronio; relaciona a\_rede\_neural\_de\_ouro}
\prov{proveniência: wiki \textperiodcentered{} fato \textperiodcentered{} confiança 1.0 \textperiodcentered{} domínio programacao \textperiodcentered{} história 6 versões no jornal}

%CRISTAL {"arestas": [{"alvo": "rede_neural_artificial", "confianca": 1.0, "epistemico": "fato", "origem": "wiki", "rel": "explica"}, {"alvo": "computador_mineral", "confianca": 1.0, "epistemico": "fato", "origem": "wiki", "rel": "eh_um"}, {"alvo": "transformer", "confianca": 1.0, "epistemico": "fato", "origem": "wiki", "rel": "relaciona"}], "confianca": 1.0, "contraexemplos": [], "descricao": "Uma camada de rede neural é uma multiplicação de matriz seguida de não-linearidade — e a matriz de gatos A_n=[[n,1],[1,0]] é a ULA do computador mineral. O transformer é pilha de contrações matriciais: o modelo roda no computador mineral, camada após camada de gatos.", "epistemico": "inferencia", "exemplos": [], "id": "rede_neural_gato", "memoria": "semantica", "meta": {"dominio": "ia", "fonte": "IA e modelos de linguagem"}, "origem": "wiki", "palavras_chave": [], "sinonimos": [], "tipo": "conceito", "titulo": "A Rede Neural é a Matriz de Gatos (a ULA)"}
\section{A Rede Neural é a Matriz de Gatos (a ULA)}
Uma camada de rede neural é uma multiplicação de matriz seguida de não-linearidade — e a matriz de gatos A\_n=[[n,1],[1,0]] é a ULA do computador mineral. O transformer é pilha de contrações matriciais: o modelo roda no computador mineral, camada após camada de gatos.
\prov{relações: explica rede\_neural\_artificial; eh\_um computador\_mineral; relaciona transformer}
\prov{proveniência: wiki \textperiodcentered{} IA e modelos de linguagem \textperiodcentered{} inferencia \textperiodcentered{} confiança 1.0 \textperiodcentered{} domínio ia \textperiodcentered{} história 4 versões no jornal}

%CRISTAL {"arestas": [{"alvo": "shell", "confianca": 1.0, "epistemico": "fato", "origem": "wiki", "rel": "usa"}, {"alvo": "tudo_e_arquivo", "confianca": 1.0, "epistemico": "fato", "origem": "wiki", "rel": "usa"}], "confianca": 1.0, "contraexemplos": [], "descricao": "Desviar entrada e saída de/para arquivos — porque tudo é arquivo, inclusive o teclado e a tela. O encanamento do shell.", "epistemico": "fato", "exemplos": [], "id": "redirecionamento", "memoria": "semantica", "meta": {"dominio": "programacao", "fonte": ""}, "origem": "wiki", "palavras_chave": [], "sinonimos": [], "tipo": "conceito", "titulo": "Redirecionamento (> < >>)"}
\section{Redirecionamento (> < >>)}
Desviar entrada e saída de/para arquivos — porque tudo é arquivo, inclusive o teclado e a tela. O encanamento do shell.
\prov{relações: usa shell; usa tudo\_e\_arquivo}
\prov{proveniência: wiki \textperiodcentered{} fato \textperiodcentered{} confiança 1.0 \textperiodcentered{} domínio programacao \textperiodcentered{} história 3 versões no jornal}

%CRISTAL {"arestas": [{"alvo": "padrao_reducer", "confianca": 1.0, "epistemico": "fato", "origem": "wiki", "rel": "usa"}, {"alvo": "estado_imutavel", "confianca": 1.0, "epistemico": "fato", "origem": "wiki", "rel": "usa"}, {"alvo": "elm", "confianca": 1.0, "epistemico": "fato", "origem": "wiki", "rel": "usa"}, {"alvo": "javascript", "confianca": 1.0, "epistemico": "fato", "origem": "wiki", "rel": "parte_de"}], "confianca": 1.0, "contraexemplos": [], "descricao": "Abramov: um contêiner de estado previsível para JS — uma única árvore imutável, mudada só por reducers puros sobre ações. Inspirado na arquitetura Elm.", "epistemico": "fato", "exemplos": [], "id": "redux", "memoria": "semantica", "meta": {"autor": "Abramov", "dominio": "programacao", "fonte": ""}, "origem": "wiki", "palavras_chave": [], "sinonimos": [], "tipo": "conceito", "titulo": "Redux"}
\section{Redux}
Abramov: um contêiner de estado previsível para JS — uma única árvore imutável, mudada só por reducers puros sobre ações. Inspirado na arquitetura Elm.
\prov{relações: usa padrao\_reducer; usa estado\_imutavel; usa elm; parte\_de javascript}
\prov{proveniência: wiki \textperiodcentered{} fato \textperiodcentered{} confiança 1.0 \textperiodcentered{} domínio programacao \textperiodcentered{} história 5 versões no jornal}

%CRISTAL {"arestas": [{"alvo": "divida_tecnica", "confianca": 1.0, "epistemico": "fato", "origem": "wiki", "rel": "usa"}, {"alvo": "arquitetura_de_software", "confianca": 1.0, "epistemico": "fato", "origem": "wiki", "rel": "usa"}], "confianca": 1.0, "contraexemplos": [], "descricao": "Melhorar a estrutura interna do código sem mudar seu comportamento — pagar a dívida técnica em passos pequenos e seguros.", "epistemico": "fato", "exemplos": [], "id": "refatoracao", "memoria": "semantica", "meta": {"autor": "Fowler", "dominio": "programacao", "fonte": ""}, "origem": "wiki", "palavras_chave": [], "sinonimos": [], "tipo": "conceito", "titulo": "Refatoração"}
\section{Refatoração}
Melhorar a estrutura interna do código sem mudar seu comportamento — pagar a dívida técnica em passos pequenos e seguros.
\prov{relações: usa divida\_tecnica; usa arquitetura\_de\_software}
\prov{proveniência: wiki \textperiodcentered{} fato \textperiodcentered{} confiança 1.0 \textperiodcentered{} domínio programacao \textperiodcentered{} história 3 versões no jornal}

%CRISTAL {"arestas": [{"alvo": "ligacao_por_overlap_limiar", "confianca": 1.0, "epistemico": "fato", "origem": "wiki", "rel": "usa"}, {"alvo": "grafo_de_proveniencia", "confianca": 1.0, "epistemico": "fato", "origem": "wiki", "rel": "relaciona"}], "confianca": 1.0, "contraexemplos": [], "descricao": "Aplica as tabelas declarativas LIGACOES/FRACTAL/ENRIQUECER, integra órfãos aos hubs por proveniência e roda os fechar_* (colapso, broca_os, bai, gato, hopfield, goldchain, mecânica quântica) antes de reindexar. É a manutenção da floresta (densificar-não-podar).", "epistemico": "inferencia", "exemplos": [], "id": "refinar_densificar", "memoria": "semantica", "meta": {"autor": "broca-so", "dominio": "operacao_so", "fonte": "conversa/conhecimento/refinar.py; tecido.py"}, "origem": "wiki", "palavras_chave": [], "sinonimos": [], "tipo": "conceito", "titulo": "Refinar = Ligar Ilhas + Enriquecer Hubs + Fechar Lacunas"}
\section{Refinar = Ligar Ilhas + Enriquecer Hubs + Fechar Lacunas}
Aplica as tabelas declarativas LIGACOES/FRACTAL/ENRIQUECER, integra órfãos aos hubs por proveniência e roda os fechar\_* (colapso, broca\_os, bai, gato, hopfield, goldchain, mecânica quântica) antes de reindexar. É a manutenção da floresta (densificar-não-podar).
\prov{relações: usa ligacao\_por\_overlap\_limiar; relaciona grafo\_de\_proveniencia}
\prov{proveniência: wiki \textperiodcentered{} conversa/conhecimento/refinar.py; tecido.py \textperiodcentered{} inferencia \textperiodcentered{} confiança 1.0 \textperiodcentered{} domínio operacao\_so \textperiodcentered{} história 3 versões no jornal}

%CRISTAL {"arestas": [{"alvo": "operacoes_reservadas_no_metal", "confianca": 1.0, "epistemico": "fato", "origem": "wiki", "rel": "usa"}, {"alvo": "bai_biblioteca_de_autorizacao_isomorfica", "confianca": 1.0, "epistemico": "fato", "origem": "wiki", "rel": "eh_um"}, {"alvo": "vocabulario_canonico", "confianca": 1.0, "epistemico": "fato", "origem": "wiki", "rel": "usa"}, {"alvo": "linguagem_fisica_unificada", "confianca": 1.0, "epistemico": "fato", "origem": "wiki", "rel": "relaciona"}], "confianca": 1.0, "contraexemplos": [], "descricao": "A régua de isomorfia do BAI resolve QUALQUER nome equivalente à mesma operação. As palavras reservadas têm nome em vários idiomas (dicionário 'palavra_reservada': PT/EN/ES) — gato=cat, soma=add=suma, produto=mul=producto, crivo=sieve=criba, norma=norm=norma… compor(idioma,'reservadas_pt') traduz ao canônico que o metal executa, e chamar('cat',…)==chamar('gato',…). A linguagem do broca-so roda em qualquer língua: o mesmo isomorfismo do BAI que traduz xadrez↔delegação↔imagem agora traduz as chamadas de sistema. Verif.: gato=cat=gato e soma=add=suma dão a mesma operação.", "epistemico": "fato", "exemplos": [], "id": "regua_de_isomorfia_multi_idioma", "memoria": "semantica", "meta": {"dominio": "sistema", "fonte": "metal do SO"}, "origem": "wiki", "palavras_chave": [], "sinonimos": [], "tipo": "conceito", "titulo": "A Régua de Isomorfia: a Linguagem em Qualquer Idioma"}
\section{A Régua de Isomorfia: a Linguagem em Qualquer Idioma}
A régua de isomorfia do BAI resolve QUALQUER nome equivalente à mesma operação. As palavras reservadas têm nome em vários idiomas (dicionário 'palavra\_reservada': PT/EN/ES) — gato=cat, soma=add=suma, produto=mul=producto, crivo=sieve=criba, norma=norm=norma… compor(idioma,'reservadas\_pt') traduz ao canônico que o metal executa, e chamar('cat',…)==chamar('gato',…). A linguagem do broca-so roda em qualquer língua: o mesmo isomorfismo do BAI que traduz xadrez\ensuremath{\leftrightarrow}delegação\ensuremath{\leftrightarrow}imagem agora traduz as chamadas de sistema. Verif.: gato=cat=gato e soma=add=suma dão a mesma operação.
\prov{relações: usa operacoes\_reservadas\_no\_metal; eh\_um bai\_biblioteca\_de\_autorizacao\_isomorfica; usa vocabulario\_canonico; relaciona linguagem\_fisica\_unificada}
\prov{proveniência: wiki \textperiodcentered{} metal do SO \textperiodcentered{} fato \textperiodcentered{} confiança 1.0 \textperiodcentered{} domínio sistema \textperiodcentered{} história 5 versões no jornal}

%CRISTAL {"arestas": [{"alvo": "isolamento_base_cliente", "confianca": 1.0, "epistemico": "fato", "origem": "wiki", "rel": "usa"}, {"alvo": "soberania_da_chave", "confianca": 1.0, "epistemico": "fato", "origem": "wiki", "rel": "relaciona"}], "confianca": 1.0, "contraexemplos": [], "descricao": "pack_core_release copia apenas base/ para staging, roda o build Core P0 e gera tarball .tar.zst; a verificação RECUSA o pacote se cliente/ estiver presente — o deploy público nunca carrega dado proprietário. Há também um template de cliente vazio para novos tenants.", "epistemico": "fato", "exemplos": [], "id": "release_core_p0", "memoria": "semantica", "meta": {"autor": "broca-so", "dominio": "operacao_so", "fonte": "code/tiffany_platform/release.py; cliente_template.py"}, "origem": "wiki", "palavras_chave": [], "sinonimos": [], "tipo": "conceito", "titulo": "Release do Core (deploy só da base)"}
\section{Release do Core (deploy só da base)}
pack\_core\_release copia apenas base/ para staging, roda o build Core P0 e gera tarball .tar.zst; a verificação RECUSA o pacote se cliente/ estiver presente — o deploy público nunca carrega dado proprietário. Há também um template de cliente vazio para novos tenants.
\prov{relações: usa isolamento\_base\_cliente; relaciona soberania\_da\_chave}
\prov{proveniência: wiki \textperiodcentered{} code/tiffany\_platform/release.py; cliente\_template.py \textperiodcentered{} fato \textperiodcentered{} confiança 1.0 \textperiodcentered{} domínio operacao\_so \textperiodcentered{} história 3 versões no jornal}

%CRISTAL {"arestas": [{"alvo": "compilador", "confianca": 1.0, "epistemico": "fato", "origem": "wiki", "rel": "usa"}], "confianca": 1.0, "contraexemplos": [], "descricao": "A forma neutra entre fonte e binário sobre a qual o compilador otimiza (LLVM IR) — desacopla linguagem de máquina.", "epistemico": "fato", "exemplos": [], "id": "representacao_intermediaria", "memoria": "semantica", "meta": {"dominio": "programacao", "fonte": ""}, "origem": "wiki", "palavras_chave": [], "sinonimos": [], "tipo": "conceito", "titulo": "Representação Intermediária (IR)"}
\section{Representação Intermediária (IR)}
A forma neutra entre fonte e binário sobre a qual o compilador otimiza (LLVM IR) — desacopla linguagem de máquina.
\prov{relações: usa compilador}
\prov{proveniência: wiki \textperiodcentered{} fato \textperiodcentered{} confiança 1.0 \textperiodcentered{} domínio programacao \textperiodcentered{} história 4 versões no jornal}

%CRISTAL {"arestas": [{"alvo": "operacao_broca_so", "confianca": 1.0, "epistemico": "fato", "origem": "wiki", "rel": "parte_de"}, {"alvo": "regua_dnci", "confianca": 1.0, "epistemico": "fato", "origem": "wiki", "rel": "usa"}], "confianca": 1.0, "contraexemplos": [], "descricao": "lab.py reproduzir roda as SUITES (simbolico/ponte/glacial/dual/radio) de neuronio/experimentos_tiffany, verificando os módulos-neurônio contra os papers — o conhecimento só vira fato se reproduzível. É o portão que toda semilla atravessa antes do commit.", "epistemico": "fato", "exemplos": [], "id": "reproduzir_experimentos", "memoria": "semantica", "meta": {"autor": "broca-so", "dominio": "operacao_so", "fonte": "conversa/conhecimento/lab.py; neuronio/experimentos_tiffany.py"}, "origem": "wiki", "palavras_chave": [], "sinonimos": [], "tipo": "conceito", "titulo": "Prova Reproduzível (papers ↔ código, 62/62)"}
\section{Prova Reproduzível (papers \ensuremath{\leftrightarrow} código, 62/62)}
lab.py reproduzir roda as SUITES (simbolico/ponte/glacial/dual/radio) de neuronio/experimentos\_tiffany, verificando os módulos-neurônio contra os papers — o conhecimento só vira fato se reproduzível. É o portão que toda semilla atravessa antes do commit.
\prov{relações: parte\_de operacao\_broca\_so; usa regua\_dnci}
\prov{proveniência: wiki \textperiodcentered{} conversa/conhecimento/lab.py; neuronio/experimentos\_tiffany.py \textperiodcentered{} fato \textperiodcentered{} confiança 1.0 \textperiodcentered{} domínio operacao\_so \textperiodcentered{} história 3 versões no jornal}

%CRISTAL {"arestas": [{"alvo": "confianca_epistemica", "confianca": 1.0, "epistemico": "fato", "origem": "curriculo", "rel": "explica"}, {"alvo": "virtualizacao", "confianca": 0.5, "epistemico": "fato", "origem": "wiki", "rel": "relaciona"}, {"alvo": "word", "confianca": 0.5, "epistemico": "fato", "origem": "wiki", "rel": "relaciona"}, {"alvo": "vontade_de_poder", "confianca": 0.5, "epistemico": "fato", "origem": "wiki", "rel": "relaciona"}], "confianca": 1.0, "contraexemplos": [], "descricao": "Escolha explícita: confirmar, analogia, dissolver ou enriquecer.", "epistemico": "fato", "exemplos": [], "id": "resolucao", "memoria": "semantica", "meta": {"dominio": "aprendizagem", "nivel": 7, "ordem": 4}, "origem": "curriculo", "palavras_chave": ["arbitragem", "fecho"], "sinonimos": ["resolução"], "tipo": "conceito", "titulo": "resolucao"}
\section{resolucao}
Escolha explícita: confirmar, analogia, dissolver ou enriquecer.
\prov{sinónimos: resolução}
\prov{relações: explica confianca\_epistemica; relaciona virtualizacao; relaciona word; relaciona vontade\_de\_poder}
\prov{proveniência: curriculo \textperiodcentered{} fato \textperiodcentered{} confiança 1.0 \textperiodcentered{} domínio aprendizagem \textperiodcentered{} história 6 versões no jornal}

%CRISTAL {"arestas": [{"alvo": "api", "confianca": 1.0, "epistemico": "fato", "origem": "wiki", "rel": "especializa"}, {"alvo": "http", "confianca": 1.0, "epistemico": "fato", "origem": "wiki", "rel": "usa"}], "confianca": 1.0, "contraexemplos": [], "descricao": "Um estilo de API sobre HTTP: recursos com URLs, verbos padrão, sem estado no servidor — a arquitetura de fato da web de serviços.", "epistemico": "fato", "exemplos": [], "id": "rest", "memoria": "semantica", "meta": {"autor": "Fielding", "dominio": "programacao", "fonte": ""}, "origem": "wiki", "palavras_chave": [], "sinonimos": [], "tipo": "conceito", "titulo": "REST"}
\section{REST}
Um estilo de API sobre HTTP: recursos com URLs, verbos padrão, sem estado no servidor — a arquitetura de fato da web de serviços.
\prov{relações: especializa api; usa http}
\prov{proveniência: wiki \textperiodcentered{} fato \textperiodcentered{} confiança 1.0 \textperiodcentered{} domínio programacao \textperiodcentered{} história 6 versões no jornal}

%CRISTAL {"arestas": [{"alvo": "gradiente_descendente", "confianca": 1.0, "epistemico": "fato", "origem": "wiki", "rel": "usa"}, {"alvo": "rede_neural_artificial", "confianca": 1.0, "epistemico": "fato", "origem": "wiki", "rel": "usa"}], "confianca": 1.0, "contraexemplos": [], "descricao": "Propagar o erro de trás para frente pela rede, pela regra da cadeia, para saber quanto cada peso contribuiu. Como a rede aprende.", "epistemico": "fato", "exemplos": [], "id": "retropropagacao", "memoria": "semantica", "meta": {"autor": "Rumelhart/Hinton", "dominio": "programacao", "fonte": ""}, "origem": "wiki", "palavras_chave": [], "sinonimos": [], "tipo": "conceito", "titulo": "Retropropagação (Backprop)"}
\section{Retropropagação (Backprop)}
Propagar o erro de trás para frente pela rede, pela regra da cadeia, para saber quanto cada peso contribuiu. Como a rede aprende.
\prov{relações: usa gradiente\_descendente; usa rede\_neural\_artificial}
\prov{proveniência: wiki \textperiodcentered{} fato \textperiodcentered{} confiança 1.0 \textperiodcentered{} domínio programacao \textperiodcentered{} história 3 versões no jornal}

%CRISTAL {"arestas": [{"alvo": "isolamento_base_cliente", "confianca": 1.0, "epistemico": "fato", "origem": "wiki", "rel": "usa"}, {"alvo": "orquestrador_roteador_regex", "confianca": 1.0, "epistemico": "fato", "origem": "wiki", "rel": "relaciona"}], "confianca": 1.0, "contraexemplos": [], "descricao": "Por padrão a consulta usa só base/; passa a incluir cliente/ se a query casar domínio cliente (keyword), com cliente_only, ou TIFFANY_CLIENTE=1 — depois seleciona até N domínios por score e funde por overlap, respeitando o isolamento.", "epistemico": "fato", "exemplos": [], "id": "roteamento_multi_tenant", "memoria": "semantica", "meta": {"autor": "broca-so", "dominio": "operacao_so", "fonte": "code/tiffany_platform/orquestrador_multi.py"}, "origem": "wiki", "palavras_chave": [], "sinonimos": [], "tipo": "conceito", "titulo": "Roteamento Multi-Tenant (base vs cliente)"}
\section{Roteamento Multi-Tenant (base vs cliente)}
Por padrão a consulta usa só base/; passa a incluir cliente/ se a query casar domínio cliente (keyword), com cliente\_only, ou TIFFANY\_CLIENTE=1 — depois seleciona até N domínios por score e funde por overlap, respeitando o isolamento.
\prov{relações: usa isolamento\_base\_cliente; relaciona orquestrador\_roteador\_regex}
\prov{proveniência: wiki \textperiodcentered{} code/tiffany\_platform/orquestrador\_multi.py \textperiodcentered{} fato \textperiodcentered{} confiança 1.0 \textperiodcentered{} domínio operacao\_so \textperiodcentered{} história 3 versões no jornal}

%CRISTAL {"arestas": [{"alvo": "linguagem_descricao_hardware", "confianca": 1.0, "epistemico": "fato", "origem": "wiki", "rel": "usa"}], "confianca": 1.0, "contraexemplos": [], "descricao": "O nível em que se descreve hardware: o fluxo de dados entre registradores a cada ciclo de clock — a abstração central de Verilog/VHDL.", "epistemico": "fato", "exemplos": [], "id": "rtl", "memoria": "semantica", "meta": {"dominio": "programacao", "fonte": ""}, "origem": "wiki", "palavras_chave": [], "sinonimos": [], "tipo": "conceito", "titulo": "RTL (Register-Transfer Level)"}
\section{RTL (Register-Transfer Level)}
O nível em que se descreve hardware: o fluxo de dados entre registradores a cada ciclo de clock — a abstração central de Verilog/VHDL.
\prov{relações: usa linguagem\_descricao\_hardware}
\prov{proveniência: wiki \textperiodcentered{} fato \textperiodcentered{} confiança 1.0 \textperiodcentered{} domínio programacao \textperiodcentered{} história 2 versões no jornal}

%CRISTAL {"arestas": [{"alvo": "programacao_orientada_objetos", "confianca": 1.0, "epistemico": "fato", "origem": "wiki", "rel": "usa"}, {"alvo": "linguagem_de_programacao", "confianca": 1.0, "epistemico": "fato", "origem": "wiki", "rel": "parte_de"}], "confianca": 1.0, "contraexemplos": [], "descricao": "Matsumoto: desenhada para a felicidade do programador — orientação a objetos pura e expressiva; o motor do Rails.", "epistemico": "fato", "exemplos": [], "id": "ruby", "memoria": "semantica", "meta": {"autor": "Matsumoto", "dominio": "programacao", "fonte": ""}, "origem": "wiki", "palavras_chave": [], "sinonimos": [], "tipo": "conceito", "titulo": "Ruby"}
\section{Ruby}
Matsumoto: desenhada para a felicidade do programador — orientação a objetos pura e expressiva; o motor do Rails.
\prov{relações: usa programacao\_orientada\_objetos; parte\_de linguagem\_de\_programacao}
\prov{proveniência: wiki \textperiodcentered{} fato \textperiodcentered{} confiança 1.0 \textperiodcentered{} domínio programacao \textperiodcentered{} história 6 versões no jornal}

%CRISTAL {"arestas": [{"alvo": "compilador", "confianca": 1.0, "epistemico": "fato", "origem": "curriculo", "rel": "referencia"}, {"alvo": "thread", "confianca": 0.5, "epistemico": "fato", "origem": "wiki", "rel": "relaciona"}, {"alvo": "topologia", "confianca": 0.5, "epistemico": "fato", "origem": "wiki", "rel": "relaciona"}, {"alvo": "syscall", "confianca": 0.5, "epistemico": "fato", "origem": "wiki", "rel": "relaciona"}, {"alvo": "testes", "confianca": 0.5, "epistemico": "fato", "origem": "wiki", "rel": "relaciona"}, {"alvo": "teste_ordem", "confianca": 0.5, "epistemico": "fato", "origem": "wiki", "rel": "relaciona"}, {"alvo": "vida-morte", "confianca": 0.5, "epistemico": "fato", "origem": "wiki", "rel": "relaciona"}, {"alvo": "word", "confianca": 0.5, "epistemico": "fato", "origem": "wiki", "rel": "relaciona"}, {"alvo": "linguagem_de_programacao", "confianca": 1.0, "epistemico": "fato", "origem": "wiki", "rel": "parte_de"}, {"alvo": "ownership_borrow", "confianca": 1.0, "epistemico": "fato", "origem": "wiki", "rel": "usa"}], "confianca": 1.0, "contraexemplos": [], "descricao": "Segurança de memória sem coletor, via ownership e borrow checker; concorrência sem data races. Sistemas sem medo.", "epistemico": "fato", "exemplos": [], "id": "rust", "memoria": "semantica", "meta": {"autor": "Hoare", "dominio": "programacao", "fonte": ""}, "origem": "wiki", "palavras_chave": [], "sinonimos": [], "tipo": "conceito", "titulo": "Rust"}
\section{Rust}
Segurança de memória sem coletor, via ownership e borrow checker; concorrência sem data races. Sistemas sem medo.
\prov{relações: referencia compilador; relaciona thread; relaciona topologia; relaciona syscall; relaciona testes; relaciona teste\_ordem; relaciona vida-morte; relaciona word; parte\_de linguagem\_de\_programacao; usa ownership\_borrow}
\prov{proveniência: wiki \textperiodcentered{} fato \textperiodcentered{} confiança 1.0 \textperiodcentered{} domínio programacao \textperiodcentered{} história 17 versões no jornal}

%CRISTAL {"arestas": [{"alvo": "rust", "confianca": 1.0, "epistemico": "fato", "origem": "wiki", "rel": "usa"}, {"alvo": "linux", "confianca": 1.0, "epistemico": "fato", "origem": "wiki", "rel": "parte_de"}, {"alvo": "programacao_de_sistemas", "confianca": 1.0, "epistemico": "fato", "origem": "wiki", "rel": "especializa"}, {"alvo": "c", "confianca": 1.0, "epistemico": "fato", "origem": "wiki", "rel": "contradiz"}], "confianca": 1.0, "contraexemplos": [], "descricao": "A entrada do Rust no Linux e em SOs: segurança de memória sem coletor no lugar mais perigoso do software — o C dominante ganha um rival seguro.", "epistemico": "inferencia", "exemplos": [], "id": "rust_no_kernel", "memoria": "semantica", "meta": {"autor": "Rust", "dominio": "programacao", "fonte": ""}, "origem": "wiki", "palavras_chave": [], "sinonimos": [], "tipo": "conceito", "titulo": "Rust no Kernel"}
\section{Rust no Kernel}
A entrada do Rust no Linux e em SOs: segurança de memória sem coletor no lugar mais perigoso do software — o C dominante ganha um rival seguro.
\prov{relações: usa rust; parte\_de linux; especializa programacao\_de\_sistemas; contradiz c}
\prov{proveniência: wiki \textperiodcentered{} inferencia \textperiodcentered{} confiança 1.0 \textperiodcentered{} domínio programacao \textperiodcentered{} história 5 versões no jornal}

%CRISTAL {"arestas": [{"alvo": "porta_logica", "confianca": 1.0, "epistemico": "fato", "origem": "wiki", "rel": "usa"}, {"alvo": "transistor", "confianca": 1.0, "epistemico": "fato", "origem": "wiki", "rel": "usa"}, {"alvo": "logica_digital", "confianca": 1.0, "epistemico": "fato", "origem": "wiki", "rel": "usa"}, {"alvo": "circuito_integrado", "confianca": 1.0, "epistemico": "fato", "origem": "wiki", "rel": "explica"}, {"alvo": "backend_isa12_broca", "confianca": 1.0, "epistemico": "fato", "origem": "wiki", "rel": "usa"}, {"alvo": "linguagem_matriz_gatos", "confianca": 1.0, "epistemico": "fato", "origem": "wiki", "rel": "usa"}, {"alvo": "dsl_frontend_linguagem_gatos", "confianca": 1.0, "epistemico": "fato", "origem": "wiki", "rel": "usa"}, {"alvo": "broca_os", "confianca": 1.0, "epistemico": "fato", "origem": "wiki", "rel": "explica"}, {"alvo": "gato", "confianca": 1.0, "epistemico": "fato", "origem": "wiki", "rel": "usa"}, {"alvo": "cristalografia", "confianca": 1.0, "epistemico": "fato", "origem": "wiki", "rel": "relaciona"}, {"alvo": "dopagem", "confianca": 1.0, "epistemico": "fato", "origem": "wiki", "rel": "usa"}], "confianca": 1.0, "contraexemplos": [], "descricao": "linguagem/sandbox.py constrói a torre inteira, cada camada REAL: (0) CRISTAL dopado → portadores; (1) TRANSISTOR como chave (nmos/pmos); (2) PORTA — NAND CMOS de verdade (pull-down NMOS série, pull-up PMOS paralelo) e dela NÃO/E/OU/XOR; (3) ULA — somador de N bits SÓ de portas, que REALMENTE soma (5+3=8, 100+44=144, do transistor ao número); (4–5) CPU/SO/LINGUAGEM — a ISA-12 (backend_isa) e a DSL de gatos (dsl.executar). A torre é CONTÍNUA: o mesmo '+' que o somador faz do transistor, a linguagem usa no topo. torre() demonstra tudo computando.", "epistemico": "fato", "exemplos": [], "id": "sandbox_cristal_ao_so", "memoria": "semantica", "meta": {"dominio": "programacao", "fonte": "linguagem/sandbox.py (nmos/pmos/nand/somador/rodar/torre)"}, "origem": "wiki", "palavras_chave": [], "sinonimos": [], "tipo": "conceito", "titulo": "O Sandbox do broca-so — do cristal à linguagem, executável"}
\section{O Sandbox do broca-so — do cristal à linguagem, executável}
linguagem/sandbox.py constrói a torre inteira, cada camada REAL: (0) CRISTAL dopado \ensuremath{\to} portadores; (1) TRANSISTOR como chave (nmos/pmos); (2) PORTA — NAND CMOS de verdade (pull-down NMOS série, pull-up PMOS paralelo) e dela NÃO/E/OU/XOR; (3) ULA — somador de N bits SÓ de portas, que REALMENTE soma (5+3=8, 100+44=144, do transistor ao número); (4–5) CPU/SO/LINGUAGEM — a ISA-12 (backend\_isa) e a DSL de gatos (dsl.executar). A torre é CONTÍNUA: o mesmo '+' que o somador faz do transistor, a linguagem usa no topo. torre() demonstra tudo computando.
\prov{relações: usa porta\_logica; usa transistor; usa logica\_digital; explica circuito\_integrado; usa backend\_isa12\_broca; usa linguagem\_matriz\_gatos; usa dsl\_frontend\_linguagem\_gatos; explica broca\_os; usa gato; relaciona cristalografia; usa dopagem}
\prov{proveniência: wiki \textperiodcentered{} linguagem/sandbox.py (nmos/pmos/nand/somador/rodar/torre) \textperiodcentered{} fato \textperiodcentered{} confiança 1.0 \textperiodcentered{} domínio programacao \textperiodcentered{} história 24 versões no jornal}

%CRISTAL {"arestas": [{"alvo": "css", "confianca": 1.0, "epistemico": "fato", "origem": "wiki", "rel": "especializa"}], "confianca": 1.0, "contraexemplos": [], "descricao": "Um pré-processador que estende o CSS com variáveis, aninhamento e mixins — programação para folhas de estilo, compilada em CSS puro.", "epistemico": "fato", "exemplos": [], "id": "sass", "memoria": "semantica", "meta": {"dominio": "programacao", "fonte": ""}, "origem": "wiki", "palavras_chave": [], "sinonimos": [], "tipo": "conceito", "titulo": "Sass"}
\section{Sass}
Um pré-processador que estende o CSS com variáveis, aninhamento e mixins — programação para folhas de estilo, compilada em CSS puro.
\prov{relações: especializa css}
\prov{proveniência: wiki \textperiodcentered{} fato \textperiodcentered{} confiança 1.0 \textperiodcentered{} domínio programacao \textperiodcentered{} história 4 versões no jornal}

%CRISTAL {"arestas": [], "confianca": 1.0, "contraexemplos": [], "descricao": "SASS (Streaming ASSembler) é o CÓDIGO DE MÁQUINA nativo que a GPU NVIDIA executa de fato — abaixo do PTX. É ESPECÍFICO da arquitetura: cada geração da SM (Streaming Multiprocessor) tem a sua codificação de instruções (Turing≠Ampere≠Ada≠Hopper). Instruções típicas: FFMA (fused multiply-add, d=a·b+c, 1 ciclo — o coração do cálculo), FADD/FMUL/FMNMX, MUFU (Multi-Function Unit: sin/cos/ex2/lg2/rcp/rsqrt — as transcendentes), IMAD/IADD3/LOP3 (inteiro/lógica), ISETP (set predicate — o if), LDG/STG (load/store global), LDS/STS (shared), SHFL (warp shuffle — troca entre lanes), BAR.SYNC (barreira), BRA/BSSY/BSYNC (branch + convergência, Volta+), EXIT. Registradores: R0–R255 (por thread), predicados P0–P6 (P7=PT), uniform regs UR (Turing+). O warp = 32 threads (as lanes SIMT). Ferramentas: cuobjdump --dump-sass, nvdisasm. Fato (ISA da NVIDIA).", "epistemico": "fato", "exemplos": [], "id": "sass_gpu_assembly", "memoria": "semantica", "meta": {"dominio": "sass_gpu", "fonte": "SASS/GPU (o metal)"}, "origem": "wiki", "palavras_chave": [], "sinonimos": [], "tipo": "conceito", "titulo": "SASS — o assembly nativo da GPU (o código de máquina que a SM executa)"}
\section{SASS — o assembly nativo da GPU (o código de máquina que a SM executa)}
SASS (Streaming ASSembler) é o CÓDIGO DE MÁQUINA nativo que a GPU NVIDIA executa de fato — abaixo do PTX. É ESPECÍFICO da arquitetura: cada geração da SM (Streaming Multiprocessor) tem a sua codificação de instruções (Turing\ensuremath{\ne}Ampere\ensuremath{\ne}Ada\ensuremath{\ne}Hopper). Instruções típicas: FFMA (fused multiply-add, d=a\textperiodcentered b+c, 1 ciclo — o coração do cálculo), FADD/FMUL/FMNMX, MUFU (Multi-Function Unit: sin/cos/ex2/lg2/rcp/rsqrt — as transcendentes), IMAD/IADD3/LOP3 (inteiro/lógica), ISETP (set predicate — o if), LDG/STG (load/store global), LDS/STS (shared), SHFL (warp shuffle — troca entre lanes), BAR.SYNC (barreira), BRA/BSSY/BSYNC (branch + convergência, Volta+), EXIT. Registradores: R0–R255 (por thread), predicados P0–P6 (P7=PT), uniform regs UR (Turing+). O warp = 32 threads (as lanes SIMT). Ferramentas: cuobjdump --dump-sass, nvdisasm. Fato (ISA da NVIDIA).
\prov{proveniência: wiki \textperiodcentered{} SASS/GPU (o metal) \textperiodcentered{} fato \textperiodcentered{} confiança 1.0 \textperiodcentered{} domínio sass\_gpu \textperiodcentered{} história 7 versões no jornal}

%CRISTAL {"arestas": [{"alvo": "programacao_funcional", "confianca": 1.0, "epistemico": "fato", "origem": "wiki", "rel": "usa"}, {"alvo": "programacao_orientada_objetos", "confianca": 1.0, "epistemico": "fato", "origem": "wiki", "rel": "usa"}], "confianca": 1.0, "contraexemplos": [], "descricao": "Odersky: funde orientação a objetos e funcional na JVM, com sistema de tipos poderoso — do acadêmico ao big data.", "epistemico": "fato", "exemplos": [], "id": "scala", "memoria": "semantica", "meta": {"autor": "Odersky", "dominio": "programacao", "fonte": ""}, "origem": "wiki", "palavras_chave": [], "sinonimos": [], "tipo": "conceito", "titulo": "Scala"}
\section{Scala}
Odersky: funde orientação a objetos e funcional na JVM, com sistema de tipos poderoso — do acadêmico ao big data.
\prov{relações: usa programacao\_funcional; usa programacao\_orientada\_objetos}
\prov{proveniência: wiki \textperiodcentered{} fato \textperiodcentered{} confiança 1.0 \textperiodcentered{} domínio programacao \textperiodcentered{} história 6 versões no jornal}

%CRISTAL {"arestas": [{"alvo": "filesystem", "confianca": 1.0, "epistemico": "fato", "origem": "curriculo", "rel": "parte_de"}, {"alvo": "broca_os", "confianca": 1.0, "epistemico": "fato", "origem": "curriculo", "rel": "analogia"}, {"alvo": "transcricao", "confianca": 0.5, "epistemico": "fato", "origem": "wiki", "rel": "relaciona"}, {"alvo": "virtualizacao", "confianca": 0.5, "epistemico": "fato", "origem": "wiki", "rel": "relaciona"}, {"alvo": "utilitarismo", "confianca": 0.5, "epistemico": "fato", "origem": "wiki", "rel": "relaciona"}, {"alvo": "vida-morte", "confianca": 0.5, "epistemico": "fato", "origem": "wiki", "rel": "relaciona"}, {"alvo": "tipos_de_interpretante", "confianca": 0.5, "epistemico": "fato", "origem": "wiki", "rel": "relaciona"}, {"alvo": "veredito_experimental", "confianca": 0.5, "epistemico": "fato", "origem": "wiki", "rel": "relaciona"}, {"alvo": "teoria_do_valor", "confianca": 0.5, "epistemico": "fato", "origem": "wiki", "rel": "relaciona"}], "confianca": 1.0, "contraexemplos": [], "descricao": "Escolhe qual thread/processo roda; política e fairness.", "epistemico": "fato", "exemplos": [], "id": "scheduler", "memoria": "semantica", "meta": {"dominio": "programacao", "nivel": 6, "ordem": 9}, "origem": "curriculo", "palavras_chave": [], "sinonimos": [], "tipo": "conceito", "titulo": "scheduler"}
\section{scheduler}
Escolhe qual thread/processo roda; política e fairness.
\prov{relações: parte\_de filesystem; analogia broca\_os; relaciona transcricao; relaciona virtualizacao; relaciona utilitarismo; relaciona vida-morte; relaciona tipos\_de\_interpretante; relaciona veredito\_experimental; relaciona teoria\_do\_valor}
\prov{proveniência: curriculo \textperiodcentered{} fato \textperiodcentered{} confiança 1.0 \textperiodcentered{} domínio programacao \textperiodcentered{} história 12 versões no jornal}

%CRISTAL {"arestas": [{"alvo": "tiffany_home", "confianca": 1.0, "epistemico": "fato", "origem": "wiki", "rel": "usa"}], "confianca": 1.0, "contraexemplos": [], "descricao": "Valida sem dependências externas os campos obrigatórios (semver, idioma, hash sha256:…, data ISO, idx vN.idx, cobertura) e cruza os chunks do meta com as linhas reais do nos.tsv e o hash do arquivo.", "epistemico": "fato", "exemplos": [], "id": "schema_meta_manifest", "memoria": "semantica", "meta": {"autor": "broca-so", "dominio": "operacao_so", "fonte": "code/tiffany_platform/schema.py; validate.py"}, "origem": "wiki", "palavras_chave": [], "sinonimos": [], "tipo": "conceito", "titulo": "Schema/Validação stdlib (meta.json, manifest.json)"}
\section{Schema/Validação stdlib (meta.json, manifest.json)}
Valida sem dependências externas os campos obrigatórios (semver, idioma, hash sha256:…, data ISO, idx vN.idx, cobertura) e cruza os chunks do meta com as linhas reais do nos.tsv e o hash do arquivo.
\prov{relações: usa tiffany\_home}
\prov{proveniência: wiki \textperiodcentered{} code/tiffany\_platform/schema.py; validate.py \textperiodcentered{} fato \textperiodcentered{} confiança 1.0 \textperiodcentered{} domínio operacao\_so \textperiodcentered{} história 2 versões no jornal}

%CRISTAL {"arestas": [{"alvo": "expressao_regular", "confianca": 1.0, "epistemico": "fato", "origem": "wiki", "rel": "usa"}, {"alvo": "shell", "confianca": 1.0, "epistemico": "fato", "origem": "wiki", "rel": "usa"}], "confianca": 1.0, "contraexemplos": [], "descricao": "O editor de fluxo: transforma texto que passa por ele, linha a linha, por comandos (substituir, apagar, inserir). Edição sem abrir o arquivo.", "epistemico": "fato", "exemplos": [], "id": "sed", "memoria": "semantica", "meta": {"dominio": "programacao", "fonte": ""}, "origem": "wiki", "palavras_chave": [], "sinonimos": [], "tipo": "conceito", "titulo": "sed"}
\section{sed}
O editor de fluxo: transforma texto que passa por ele, linha a linha, por comandos (substituir, apagar, inserir). Edição sem abrir o arquivo.
\prov{relações: usa expressao\_regular; usa shell}
\prov{proveniência: wiki \textperiodcentered{} fato \textperiodcentered{} confiança 1.0 \textperiodcentered{} domínio programacao \textperiodcentered{} história 3 versões no jornal}

%CRISTAL {"arestas": [], "confianca": 1.0, "contraexemplos": [], "descricao": "Proteger dados e sistemas contra acesso, alteração ou negação indevidos — assumindo que há um adversário inteligente do outro lado.", "epistemico": "fato", "exemplos": [], "id": "seguranca_da_informacao", "memoria": "semantica", "meta": {"dominio": "programacao", "fonte": ""}, "origem": "wiki", "palavras_chave": [], "sinonimos": [], "tipo": "conceito", "titulo": "Segurança da Informação"}
\section{Segurança da Informação}
Proteger dados e sistemas contra acesso, alteração ou negação indevidos — assumindo que há um adversário inteligente do outro lado.
\prov{proveniência: wiki \textperiodcentered{} fato \textperiodcentered{} confiança 1.0 \textperiodcentered{} domínio programacao}

%CRISTAL {"arestas": [{"alvo": "semente_booleana_machinec", "confianca": 1.0, "epistemico": "fato", "origem": "wiki", "rel": "especializa"}, {"alvo": "dna_dupla_helice", "confianca": 1.0, "epistemico": "fato", "origem": "wiki", "rel": "relaciona"}, {"alvo": "broca_os", "confianca": 1.0, "epistemico": "fato", "origem": "wiki", "rel": "parte_de"}], "confianca": 1.0, "contraexemplos": [], "descricao": "A Máquina Booleana sobre 𝔽₂ (6 primitivas: ⊕=gap, ⊙=correlação, pop=medida, shift, NOT) é o 'DNA' auto-contido: um único paper-semente (machine/main.tex) reconstrói 9 artefatos, incluindo cadeia_genoma/cadeia_reparo/cadeia_telomero (léxico biológico explícito para integridade/auto-reconstrução). O Apêndice = recursão de Cayley-Dickson que gera 𝕆 (o 'código' mínimo que se desdobra em corpo). Substrato-independente (CPU/GPU/FPGA computam o mesmo organismo). [D]/[N] semente reconstrutível validada; a rotulagem DNA/genoma/telômero é [J] metáfora estrutural declarada.", "epistemico": "inferencia", "exemplos": [], "id": "semente_booleana_dna_sistema", "memoria": "semantica", "meta": {"autor": "broca-so", "dominio": "operacao_so", "fonte": "machine/main.tex; semente_booleana_machinec"}, "origem": "wiki", "palavras_chave": [], "sinonimos": [], "tipo": "conceito", "titulo": "A semente booleana = o DNA do sistema (auto-reconstrução)"}
\section{A semente booleana = o DNA do sistema (auto-reconstrução)}
A Máquina Booleana sobre \ensuremath{\mathbb{F}}\ensuremath{_{2}} (6 primitivas: \ensuremath{\oplus}=gap, \ensuremath{\odot}=correlação, pop=medida, shift, NOT) é o 'DNA' auto-contido: um único paper-semente (machine/main.tex) reconstrói 9 artefatos, incluindo cadeia\_genoma/cadeia\_reparo/cadeia\_telomero (léxico biológico explícito para integridade/auto-reconstrução). O Apêndice = recursão de Cayley-Dickson que gera \ensuremath{\mathbb{O}} (o 'código' mínimo que se desdobra em corpo). Substrato-independente (CPU/GPU/FPGA computam o mesmo organismo). [D]/[N] semente reconstrutível validada; a rotulagem DNA/genoma/telômero é [J] metáfora estrutural declarada.
\prov{relações: especializa semente\_booleana\_machinec; relaciona dna\_dupla\_helice; parte\_de broca\_os}
\prov{proveniência: wiki \textperiodcentered{} machine/main.tex; semente\_booleana\_machinec \textperiodcentered{} inferencia \textperiodcentered{} confiança 1.0 \textperiodcentered{} domínio operacao\_so \textperiodcentered{} história 4 versões no jornal}

%CRISTAL {"arestas": [{"alvo": "linguagem_de_marcacao", "confianca": 1.0, "epistemico": "fato", "origem": "wiki", "rel": "usa"}], "confianca": 1.0, "contraexemplos": [], "descricao": "O princípio da marcação: o QUE (estrutura, significado) mora num lugar; o COMO (estilo, layout) noutro. HTML diz o quê, CSS diz o como.", "epistemico": "inferencia", "exemplos": [], "id": "separacao_conteudo_forma", "memoria": "semantica", "meta": {"dominio": "programacao", "fonte": ""}, "origem": "wiki", "palavras_chave": [], "sinonimos": [], "tipo": "conceito", "titulo": "Separação Conteúdo/Apresentação"}
\section{Separação Conteúdo/Apresentação}
O princípio da marcação: o QUE (estrutura, significado) mora num lugar; o COMO (estilo, layout) noutro. HTML diz o quê, CSS diz o como.
\prov{relações: usa linguagem\_de\_marcacao}
\prov{proveniência: wiki \textperiodcentered{} inferencia \textperiodcentered{} confiança 1.0 \textperiodcentered{} domínio programacao \textperiodcentered{} história 2 versões no jornal}

%CRISTAL {"arestas": [{"alvo": "arquitetura_de_software", "confianca": 1.0, "epistemico": "fato", "origem": "wiki", "rel": "usa"}, {"alvo": "modularidade", "confianca": 1.0, "epistemico": "fato", "origem": "wiki", "rel": "relaciona"}], "confianca": 1.0, "contraexemplos": [], "descricao": "Cada parte do sistema cuida de uma preocupação e só dela — o princípio que organiza desde a função até a arquitetura inteira.", "epistemico": "inferencia", "exemplos": [], "id": "separacao_de_interesses", "memoria": "semantica", "meta": {"dominio": "programacao", "fonte": ""}, "origem": "wiki", "palavras_chave": [], "sinonimos": [], "tipo": "conceito", "titulo": "Separação de Interesses"}
\section{Separação de Interesses}
Cada parte do sistema cuida de uma preocupação e só dela — o princípio que organiza desde a função até a arquitetura inteira.
\prov{relações: usa arquitetura\_de\_software; relaciona modularidade}
\prov{proveniência: wiki \textperiodcentered{} inferencia \textperiodcentered{} confiança 1.0 \textperiodcentered{} domínio programacao \textperiodcentered{} história 3 versões no jornal}

%CRISTAL {"arestas": [{"alvo": "shell", "confianca": 1.0, "epistemico": "fato", "origem": "wiki", "rel": "especializa"}, {"alvo": "posix", "confianca": 1.0, "epistemico": "fato", "origem": "wiki", "rel": "parte_de"}], "confianca": 1.0, "contraexemplos": [], "descricao": "O shell Bourne e o padrão POSIX que todo Unix fala — o mínimo portável em que scripts confiáveis se escrevem.", "epistemico": "fato", "exemplos": [], "id": "sh_posix", "memoria": "semantica", "meta": {"autor": "Bourne", "dominio": "programacao", "fonte": ""}, "origem": "wiki", "palavras_chave": [], "sinonimos": [], "tipo": "conceito", "titulo": "sh (Shell POSIX)"}
\section{sh (Shell POSIX)}
O shell Bourne e o padrão POSIX que todo Unix fala — o mínimo portável em que scripts confiáveis se escrevem.
\prov{relações: especializa shell; parte\_de posix}
\prov{proveniência: wiki \textperiodcentered{} fato \textperiodcentered{} confiança 1.0 \textperiodcentered{} domínio programacao \textperiodcentered{} história 3 versões no jornal}

%CRISTAL {"arestas": [{"alvo": "bit", "confianca": 1.0, "epistemico": "fato", "origem": "curriculo", "rel": "relaciona"}, {"alvo": "blockchain", "confianca": 1.0, "epistemico": "fato", "origem": "curriculo", "rel": "relaciona"}, {"alvo": "cifra_de_ouro_branco", "confianca": 1.0, "epistemico": "fato", "origem": "curriculo", "rel": "relaciona"}, {"alvo": "goldchain", "confianca": 1.0, "epistemico": "fato", "origem": "curriculo", "rel": "relaciona"}, {"alvo": "integridade_e_norma_conservada", "confianca": 1.0, "epistemico": "fato", "origem": "wiki", "rel": "relaciona"}], "confianca": 1.0, "contraexemplos": [], "descricao": "Hash de integridade: banda do tecido, certificados, endereço na biblioteca.", "epistemico": "fato", "exemplos": [], "id": "sha256", "memoria": "semantica", "meta": {"dominio": "cripto", "nivel": 5, "serve": "Núcleo criptografia", "verificacao": "slug idêntico — enriquecer, não duplicar"}, "origem": "curriculo", "palavras_chave": ["hash", "Ed25519", "integridade"], "sinonimos": ["SHA-256", "SHA256"], "tipo": "conceito", "titulo": "sha256"}
\section{sha256}
Hash de integridade: banda do tecido, certificados, endereço na biblioteca.
\prov{sinónimos: SHA-256; SHA256}
\prov{relações: relaciona bit; relaciona blockchain; relaciona cifra\_de\_ouro\_branco; relaciona goldchain; relaciona integridade\_e\_norma\_conservada}
\prov{proveniência: curriculo \textperiodcentered{} fato \textperiodcentered{} confiança 1.0 \textperiodcentered{} domínio cripto \textperiodcentered{} história 67 versões no jornal}

%CRISTAL {"arestas": [{"alvo": "sistema_operacional", "confianca": 1.0, "epistemico": "fato", "origem": "wiki", "rel": "parte_de"}, {"alvo": "interpretador", "confianca": 1.0, "epistemico": "fato", "origem": "wiki", "rel": "especializa"}], "confianca": 1.0, "contraexemplos": [], "descricao": "O interpretador de comandos que traduz o que você digita em ações do SO — o REPL do sistema. Texto entra, processos saem.", "epistemico": "fato", "exemplos": [], "id": "shell", "memoria": "semantica", "meta": {"dominio": "programacao", "fonte": ""}, "origem": "wiki", "palavras_chave": [], "sinonimos": [], "tipo": "conceito", "titulo": "Shell"}
\section{Shell}
O interpretador de comandos que traduz o que você digita em ações do SO — o REPL do sistema. Texto entra, processos saem.
\prov{relações: parte\_de sistema\_operacional; especializa interpretador}
\prov{proveniência: wiki \textperiodcentered{} fato \textperiodcentered{} confiança 1.0 \textperiodcentered{} domínio programacao \textperiodcentered{} história 3 versões no jornal}

%CRISTAL {"arestas": [{"alvo": "shell", "confianca": 1.0, "epistemico": "fato", "origem": "wiki", "rel": "usa"}, {"alvo": "programacao", "confianca": 1.0, "epistemico": "fato", "origem": "wiki", "rel": "parte_de"}, {"alvo": "linguagem_de_programacao", "confianca": 1.0, "epistemico": "fato", "origem": "wiki", "rel": "relaciona"}], "confianca": 1.0, "contraexemplos": [], "descricao": "Programar o próprio shell: sequências de comandos, variáveis, laços e condicionais — automação colando ferramentas do sistema.", "epistemico": "fato", "exemplos": [], "id": "shell_scripting", "memoria": "semantica", "meta": {"dominio": "programacao", "fonte": ""}, "origem": "wiki", "palavras_chave": [], "sinonimos": [], "tipo": "conceito", "titulo": "Shell Scripting"}
\section{Shell Scripting}
Programar o próprio shell: sequências de comandos, variáveis, laços e condicionais — automação colando ferramentas do sistema.
\prov{relações: usa shell; parte\_de programacao; relaciona linguagem\_de\_programacao}
\prov{proveniência: wiki \textperiodcentered{} fato \textperiodcentered{} confiança 1.0 \textperiodcentered{} domínio programacao \textperiodcentered{} história 4 versões no jornal}

%CRISTAL {"arestas": [{"alvo": "semantica", "confianca": 1.0, "epistemico": "fato", "origem": "curriculo", "rel": "parte_de"}, {"alvo": "word", "confianca": 0.5, "epistemico": "fato", "origem": "wiki", "rel": "relaciona"}, {"alvo": "vontade_de_poder", "confianca": 0.5, "epistemico": "fato", "origem": "wiki", "rel": "relaciona"}, {"alvo": "virtualizacao", "confianca": 0.5, "epistemico": "fato", "origem": "wiki", "rel": "relaciona"}, {"alvo": "linguistica", "confianca": 1.0, "epistemico": "fato", "origem": "wiki", "rel": "parte_de"}, {"alvo": "hierarquia_de_chomsky", "confianca": 1.0, "epistemico": "fato", "origem": "wiki", "rel": "usa"}, {"alvo": "gramatica_bnf", "confianca": 1.0, "epistemico": "fato", "origem": "wiki", "rel": "relaciona"}], "confianca": 1.0, "contraexemplos": [], "descricao": "Ordem e combinação de unidades linguísticas; árvore de dependências.", "epistemico": "fato", "exemplos": [], "id": "sintaxe", "memoria": "semantica", "meta": {"dominio": "linguagem", "nivel": 3, "ordem": 2}, "origem": "curriculo", "palavras_chave": ["parse", "estrutura", "frase"], "sinonimos": [], "tipo": "conceito", "titulo": "sintaxe"}
\section{sintaxe}
Ordem e combinação de unidades linguísticas; árvore de dependências.
\prov{relações: parte\_de semantica; relaciona word; relaciona vontade\_de\_poder; relaciona virtualizacao; parte\_de linguistica; usa hierarquia\_de\_chomsky; relaciona gramatica\_bnf}
\prov{proveniência: curriculo \textperiodcentered{} fato \textperiodcentered{} confiança 1.0 \textperiodcentered{} domínio linguagem \textperiodcentered{} história 8 versões no jornal}

%CRISTAL {"arestas": [{"alvo": "linguagem_descricao_hardware", "confianca": 1.0, "epistemico": "fato", "origem": "wiki", "rel": "usa"}, {"alvo": "porta_logica", "confianca": 1.0, "epistemico": "fato", "origem": "wiki", "rel": "usa"}, {"alvo": "backend_isa12_broca", "confianca": 1.0, "epistemico": "fato", "origem": "wiki", "rel": "relaciona"}, {"alvo": "compilador", "confianca": 1.0, "epistemico": "fato", "origem": "wiki", "rel": "relaciona"}], "confianca": 1.0, "contraexemplos": [], "descricao": "Traduzir a descrição HDL numa rede concreta de portas lógicas — o 'compilador' do hardware, análogo ao backend que gera a ISA-12 do broca-so.", "epistemico": "fato", "exemplos": [], "id": "sintese_logica", "memoria": "semantica", "meta": {"dominio": "programacao", "fonte": ""}, "origem": "wiki", "palavras_chave": [], "sinonimos": [], "tipo": "conceito", "titulo": "Síntese Lógica"}
\section{Síntese Lógica}
Traduzir a descrição HDL numa rede concreta de portas lógicas — o 'compilador' do hardware, análogo ao backend que gera a ISA-12 do broca-so.
\prov{relações: usa linguagem\_descricao\_hardware; usa porta\_logica; relaciona backend\_isa12\_broca; relaciona compilador}
\prov{proveniência: wiki \textperiodcentered{} fato \textperiodcentered{} confiança 1.0 \textperiodcentered{} domínio programacao \textperiodcentered{} história 5 versões no jornal}

%CRISTAL {"arestas": [{"alvo": "sistema_operacional", "confianca": 1.0, "epistemico": "fato", "origem": "wiki", "rel": "parte_de"}, {"alvo": "filesystem", "confianca": 1.0, "epistemico": "fato", "origem": "wiki", "rel": "relaciona"}], "confianca": 1.0, "contraexemplos": [], "descricao": "Como o SO organiza dados em arquivos e diretórios sobre o disco — nomes, permissões, blocos. A árvore onde tudo mora.", "epistemico": "fato", "exemplos": [], "id": "sistema_de_arquivos", "memoria": "semantica", "meta": {"dominio": "programacao", "fonte": ""}, "origem": "wiki", "palavras_chave": [], "sinonimos": [], "tipo": "conceito", "titulo": "Sistema de Arquivos"}
\section{Sistema de Arquivos}
Como o SO organiza dados em arquivos e diretórios sobre o disco — nomes, permissões, blocos. A árvore onde tudo mora.
\prov{relações: parte\_de sistema\_operacional; relaciona filesystem}
\prov{proveniência: wiki \textperiodcentered{} fato \textperiodcentered{} confiança 1.0 \textperiodcentered{} domínio programacao \textperiodcentered{} história 3 versões no jornal}

%CRISTAL {"arestas": [{"alvo": "linguagem_de_programacao", "confianca": 1.0, "epistemico": "fato", "origem": "wiki", "rel": "parte_de"}], "confianca": 1.0, "contraexemplos": [], "descricao": "As regras que classificam valores e restringem operações — a rede de segurança (ou a camisa de força) da linguagem.", "epistemico": "fato", "exemplos": [], "id": "sistema_de_tipos", "memoria": "semantica", "meta": {"dominio": "programacao", "fonte": ""}, "origem": "wiki", "palavras_chave": [], "sinonimos": [], "tipo": "conceito", "titulo": "Sistema de Tipos"}
\section{Sistema de Tipos}
As regras que classificam valores e restringem operações — a rede de segurança (ou a camisa de força) da linguagem.
\prov{relações: parte\_de linguagem\_de\_programacao}
\prov{proveniência: wiki \textperiodcentered{} fato \textperiodcentered{} confiança 1.0 \textperiodcentered{} domínio programacao \textperiodcentered{} história 4 versões no jornal}

%CRISTAL {"arestas": [{"alvo": "broca_os", "confianca": 1.0, "epistemico": "fato", "origem": "curriculo", "rel": "especializa"}, {"alvo": "kernel", "confianca": 1.0, "epistemico": "fato", "origem": "curriculo", "rel": "parte_de"}, {"alvo": "broca-so_camada_de_trabalho_isomorfica", "confianca": 1.0, "epistemico": "fato", "origem": "codigo", "rel": "generaliza"}, {"alvo": "programacao_de_sistemas", "confianca": 1.0, "epistemico": "fato", "origem": "wiki", "rel": "usa"}, {"alvo": "ciencia_da_computacao", "confianca": 1.0, "epistemico": "fato", "origem": "wiki", "rel": "parte_de"}, {"alvo": "o_sistema_operacional_e_a_algebra", "confianca": 1.0, "epistemico": "fato", "origem": "wiki", "rel": "relaciona"}], "confianca": 1.0, "contraexemplos": [], "descricao": "O software que gerencia o hardware e serve os programas — processos, memória, arquivos, dispositivos. A camada entre a máquina e você.", "epistemico": "fato", "exemplos": [], "id": "sistema_operacional", "memoria": "semantica", "meta": {"dominio": "programacao", "fonte": ""}, "origem": "wiki", "palavras_chave": ["kernel", "processo", "syscall"], "sinonimos": ["SO", "sistema operacional clássico"], "tipo": "conceito", "titulo": "Sistema Operacional"}
\section{Sistema Operacional}
O software que gerencia o hardware e serve os programas — processos, memória, arquivos, dispositivos. A camada entre a máquina e você.
\prov{sinónimos: SO; sistema operacional clássico}
\prov{relações: especializa broca\_os; parte\_de kernel; generaliza broca-so\_camada\_de\_trabalho\_isomorfica; usa programacao\_de\_sistemas; parte\_de ciencia\_da\_computacao; relaciona o\_sistema\_operacional\_e\_a\_algebra}
\prov{proveniência: wiki \textperiodcentered{} fato \textperiodcentered{} confiança 1.0 \textperiodcentered{} domínio programacao \textperiodcentered{} história 45 versões no jornal}

%CRISTAL {"arestas": [{"alvo": "rede_de_computadores", "confianca": 1.0, "epistemico": "fato", "origem": "wiki", "rel": "usa"}, {"alvo": "concorrencia_paralelismo", "confianca": 1.0, "epistemico": "fato", "origem": "wiki", "rel": "usa"}], "confianca": 1.0, "contraexemplos": [], "descricao": "Muitos computadores que cooperam parecendo um só — o desafio de coordenar sem memória comum, com rede falha e sem relógio global.", "epistemico": "fato", "exemplos": [], "id": "sistemas_distribuidos", "memoria": "semantica", "meta": {"dominio": "programacao", "fonte": ""}, "origem": "wiki", "palavras_chave": [], "sinonimos": [], "tipo": "conceito", "titulo": "Sistemas Distribuídos"}
\section{Sistemas Distribuídos}
Muitos computadores que cooperam parecendo um só — o desafio de coordenar sem memória comum, com rede falha e sem relógio global.
\prov{relações: usa rede\_de\_computadores; usa concorrencia\_paralelismo}
\prov{proveniência: wiki \textperiodcentered{} fato \textperiodcentered{} confiança 1.0 \textperiodcentered{} domínio programacao \textperiodcentered{} história 6 versões no jornal}

%CRISTAL {"arestas": [{"alvo": "protocolo_tcp", "confianca": 1.0, "epistemico": "fato", "origem": "wiki", "rel": "usa"}, {"alvo": "sistema_operacional", "confianca": 1.0, "epistemico": "fato", "origem": "wiki", "rel": "usa"}], "confianca": 1.0, "contraexemplos": [], "descricao": "A abstração que o SO dá para um ponto de comunicação de rede — abrir, conectar, ler, escrever, como se fosse um arquivo (BSD).", "epistemico": "fato", "exemplos": [], "id": "socket", "memoria": "semantica", "meta": {"dominio": "programacao", "fonte": ""}, "origem": "wiki", "palavras_chave": [], "sinonimos": [], "tipo": "conceito", "titulo": "Socket"}
\section{Socket}
A abstração que o SO dá para um ponto de comunicação de rede — abrir, conectar, ler, escrever, como se fosse um arquivo (BSD).
\prov{relações: usa protocolo\_tcp; usa sistema\_operacional}
\prov{proveniência: wiki \textperiodcentered{} fato \textperiodcentered{} confiança 1.0 \textperiodcentered{} domínio programacao \textperiodcentered{} história 6 versões no jornal}

%CRISTAL {"arestas": [{"alvo": "programacao_orientada_objetos", "confianca": 1.0, "epistemico": "fato", "origem": "wiki", "rel": "usa"}, {"alvo": "acoplamento_coesao", "confianca": 1.0, "epistemico": "fato", "origem": "wiki", "rel": "usa"}], "confianca": 1.0, "contraexemplos": [], "descricao": "Cinco princípios de design orientado a objetos (responsabilidade única, aberto/fechado…) para código flexível e sustentável.", "epistemico": "inferencia", "exemplos": [], "id": "solid", "memoria": "semantica", "meta": {"autor": "Robert Martin", "dominio": "programacao", "fonte": ""}, "origem": "wiki", "palavras_chave": [], "sinonimos": [], "tipo": "conceito", "titulo": "Princípios SOLID"}
\section{Princípios SOLID}
Cinco princípios de design orientado a objetos (responsabilidade única, aberto/fechado…) para código flexível e sustentável.
\prov{relações: usa programacao\_orientada\_objetos; usa acoplamento\_coesao}
\prov{proveniência: wiki \textperiodcentered{} inferencia \textperiodcentered{} confiança 1.0 \textperiodcentered{} domínio programacao \textperiodcentered{} história 3 versões no jornal}

%CRISTAL {"arestas": [{"alvo": "framework_frontend", "confianca": 1.0, "epistemico": "fato", "origem": "wiki", "rel": "usa"}, {"alvo": "navegador", "confianca": 1.0, "epistemico": "fato", "origem": "wiki", "rel": "usa"}], "confianca": 1.0, "contraexemplos": [], "descricao": "Uma aplicação que carrega uma página e reescreve o conteúdo via JS sem recarregar — o navegador vira plataforma de app, não leitor de documentos.", "epistemico": "inferencia", "exemplos": [], "id": "spa", "memoria": "semantica", "meta": {"dominio": "programacao", "fonte": ""}, "origem": "wiki", "palavras_chave": [], "sinonimos": [], "tipo": "conceito", "titulo": "SPA (Single-Page Application)"}
\section{SPA (Single-Page Application)}
Uma aplicação que carrega uma página e reescreve o conteúdo via JS sem recarregar — o navegador vira plataforma de app, não leitor de documentos.
\prov{relações: usa framework\_frontend; usa navegador}
\prov{proveniência: wiki \textperiodcentered{} inferencia \textperiodcentered{} confiança 1.0 \textperiodcentered{} domínio programacao \textperiodcentered{} história 6 versões no jornal}

%CRISTAL {"arestas": [{"alvo": "programacao_declarativa", "confianca": 1.0, "epistemico": "fato", "origem": "wiki", "rel": "especializa"}, {"alvo": "algebra_abstrata", "confianca": 1.0, "epistemico": "fato", "origem": "wiki", "rel": "relaciona"}, {"alvo": "modelo_relacional", "confianca": 1.0, "epistemico": "fato", "origem": "wiki", "rel": "usa"}, {"alvo": "banco_de_dados", "confianca": 1.0, "epistemico": "fato", "origem": "wiki", "rel": "parte_de"}], "confianca": 1.0, "contraexemplos": [], "descricao": "A linguagem declarativa das bases relacionais: você diz QUE dados quer, o motor decide como buscá-los. Álgebra relacional executável.", "epistemico": "fato", "exemplos": [], "id": "sql", "memoria": "semantica", "meta": {"dominio": "programacao", "fonte": ""}, "origem": "wiki", "palavras_chave": [], "sinonimos": [], "tipo": "conceito", "titulo": "SQL"}
\section{SQL}
A linguagem declarativa das bases relacionais: você diz QUE dados quer, o motor decide como buscá-los. Álgebra relacional executável.
\prov{relações: especializa programacao\_declarativa; relaciona algebra\_abstrata; usa modelo\_relacional; parte\_de banco\_de\_dados}
\prov{proveniência: wiki \textperiodcentered{} fato \textperiodcentered{} confiança 1.0 \textperiodcentered{} domínio programacao \textperiodcentered{} história 12 versões no jornal}

%CRISTAL {"arestas": [{"alvo": "framework_frontend", "confianca": 1.0, "epistemico": "fato", "origem": "wiki", "rel": "especializa"}, {"alvo": "compilador", "confianca": 1.0, "epistemico": "fato", "origem": "wiki", "rel": "usa"}, {"alvo": "dom_virtual", "confianca": 1.0, "epistemico": "fato", "origem": "wiki", "rel": "contradiz"}], "confianca": 1.0, "contraexemplos": [], "descricao": "Harris: um COMPILADOR que gera JS mínimo sem DOM virtual nem runtime pesado — a reatividade resolvida em tempo de build.", "epistemico": "fato", "exemplos": [], "id": "svelte", "memoria": "semantica", "meta": {"autor": "Harris", "dominio": "programacao", "fonte": ""}, "origem": "wiki", "palavras_chave": [], "sinonimos": [], "tipo": "conceito", "titulo": "Svelte"}
\section{Svelte}
Harris: um COMPILADOR que gera JS mínimo sem DOM virtual nem runtime pesado — a reatividade resolvida em tempo de build.
\prov{relações: especializa framework\_frontend; usa compilador; contradiz dom\_virtual}
\prov{proveniência: wiki \textperiodcentered{} fato \textperiodcentered{} confiança 1.0 \textperiodcentered{} domínio programacao \textperiodcentered{} história 8 versões no jornal}

%CRISTAL {"arestas": [{"alvo": "xml", "confianca": 1.0, "epistemico": "fato", "origem": "wiki", "rel": "especializa"}], "confianca": 1.0, "contraexemplos": [], "descricao": "Scalable Vector Graphics: marcação XML que descreve imagens por formas e curvas, não pixels — gráfico que escala sem perder nitidez.", "epistemico": "fato", "exemplos": [], "id": "svg", "memoria": "semantica", "meta": {"dominio": "programacao", "fonte": ""}, "origem": "wiki", "palavras_chave": [], "sinonimos": [], "tipo": "conceito", "titulo": "SVG"}
\section{SVG}
Scalable Vector Graphics: marcação XML que descreve imagens por formas e curvas, não pixels — gráfico que escala sem perder nitidez.
\prov{relações: especializa xml}
\prov{proveniência: wiki \textperiodcentered{} fato \textperiodcentered{} confiança 1.0 \textperiodcentered{} domínio programacao \textperiodcentered{} história 2 versões no jornal}

%CRISTAL {"arestas": [{"alvo": "linguagem_de_programacao", "confianca": 1.0, "epistemico": "fato", "origem": "wiki", "rel": "parte_de"}], "confianca": 1.0, "contraexemplos": [], "descricao": "Apple: sucessor moderno do Objective-C — seguro, veloz e expressivo para o ecossistema Apple.", "epistemico": "fato", "exemplos": [], "id": "swift", "memoria": "semantica", "meta": {"autor": "Lattner", "dominio": "programacao", "fonte": ""}, "origem": "wiki", "palavras_chave": [], "sinonimos": [], "tipo": "conceito", "titulo": "Swift"}
\section{Swift}
Apple: sucessor moderno do Objective-C — seguro, veloz e expressivo para o ecossistema Apple.
\prov{relações: parte\_de linguagem\_de\_programacao}
\prov{proveniência: wiki \textperiodcentered{} fato \textperiodcentered{} confiança 1.0 \textperiodcentered{} domínio programacao \textperiodcentered{} história 4 versões no jornal}

%CRISTAL {"arestas": [{"alvo": "symbolic_beta_pipe_so", "confianca": 1.0, "epistemico": "fato", "origem": "wiki", "rel": "parte_de"}, {"alvo": "cristal_vs_cha", "confianca": 1.0, "epistemico": "fato", "origem": "wiki", "rel": "usa"}, {"alvo": "dimensao_de_prata_pell", "confianca": 1.0, "epistemico": "fato", "origem": "wiki", "rel": "relaciona"}], "confianca": 1.0, "contraexemplos": [], "descricao": "O valor de um Pell cresce ~σ₂ⁿ (1+√2)ⁿ — a partir de ~#100 vira número de milhares de dígitos e polui o log. Ajuste (2 jul 2026): _linha_lastro, recuperar e status passam a logar só a POSIÇÃO 'Pell #idx' e o '[total emitidos: N]', nunca o valor explícito (que continua no ledger). É a mesma lição do cristal: a lei/posição é compressível (o gerador basta), o valor é desdobramento determinístico — não precisa carregar o floco inteiro no log. [D] no código (goldchain/symbolic.py), redeployado no gex44 (tmux broca-stratum reiniciado).", "epistemico": "fato", "exemplos": [], "id": "symbolic_beta_log_posicao", "memoria": "semantica", "meta": {"autor": "Lopes/broca-so", "dominio": "goldchain", "fonte": "goldchain/symbolic.py (_linha_lastro/cmd_recuperar/cmd_status); commit 79cffca"}, "origem": "wiki", "palavras_chave": [], "sinonimos": [], "tipo": "conceito", "titulo": "O log do SymbolicBeta guarda só a POSIÇÃO do Pell (#idx) + total, não o valor"}
\section{O log do SymbolicBeta guarda só a POSIÇÃO do Pell (\#idx) + total, não o valor}
O valor de um Pell cresce \~{}\ensuremath{\sigma}\ensuremath{_{2}}\ensuremath{^{n}} (1+\ensuremath{\sqrt{\,}}2)\ensuremath{^{n}} — a partir de \~{}\#100 vira número de milhares de dígitos e polui o log. Ajuste (2 jul 2026): \_linha\_lastro, recuperar e status passam a logar só a POSIÇÃO 'Pell \#idx' e o '[total emitidos: N]', nunca o valor explícito (que continua no ledger). É a mesma lição do cristal: a lei/posição é compressível (o gerador basta), o valor é desdobramento determinístico — não precisa carregar o floco inteiro no log. [D] no código (goldchain/symbolic.py), redeployado no gex44 (tmux broca-stratum reiniciado).
\prov{relações: parte\_de symbolic\_beta\_pipe\_so; usa cristal\_vs\_cha; relaciona dimensao\_de\_prata\_pell}
\prov{proveniência: wiki \textperiodcentered{} goldchain/symbolic.py (\_linha\_lastro/cmd\_recuperar/cmd\_status); commit 79cffca \textperiodcentered{} fato \textperiodcentered{} confiança 1.0 \textperiodcentered{} domínio goldchain \textperiodcentered{} história 4 versões no jornal}

%CRISTAL {"arestas": [{"alvo": "camadas_tiffany_multifractal", "confianca": 1.0, "epistemico": "fato", "origem": "wiki", "rel": "usa"}, {"alvo": "symbolic_beta_pipe_so", "confianca": 1.0, "epistemico": "fato", "origem": "wiki", "rel": "usa"}, {"alvo": "temperatura_camada_parabola", "confianca": 1.0, "epistemico": "fato", "origem": "wiki", "rel": "usa"}, {"alvo": "tiffany", "confianca": 1.0, "epistemico": "fato", "origem": "wiki", "rel": "parte_de"}], "confianca": 1.0, "contraexemplos": [], "descricao": "Levada à produção, a multifractalidade vira LANES de mineração: cada (tecido/camada, worker) é uma lane INDEPENDENTE — fonte de espectro própria, sua temperatura na parábola. O metal já COMPORTA a multifractalidade em produção: o roteamento é layer-aware (conc0 semeado POR camada; a difusão-fluido e a impressão Koch são as mesmas em toda folha). A TAXA DE CUNHAGEM do Pell escala com #lanes = #tecidos × #workers — medido: 1140 → 2277 → 3204 → 3651 Pell/s de 1→2→4→8 lanes (≈linear até saturar os 8 cores). VERDADE HONESTA (não maquiar): N workers no MESMO job do stratum NÃO multiplicam flocos (mineram faixas de nonce do mesmo job → 1 floco/job); o que multiplica é cada TECIDO ser lane própria. E o lastro segue REAL (1 floco de ouro → 1 Pell): o que escala é o THROUGHPUT do pipeline, não a emissão por unidade de ouro. [D] medido.", "epistemico": "fato", "exemplos": [], "id": "symbolic_beta_multifractal_producao", "memoria": "semantica", "meta": {"autor": "Lopes/broca-so", "dominio": "goldchain", "fonte": "medição multiprocessing (pell+corpo); goldchain/symbolic.py; neuronio/stratum_orquestrador.py (-n)"}, "origem": "wiki", "palavras_chave": [], "sinonimos": [], "tipo": "conceito", "titulo": "Em produção: as camadas multifractais são LANES de cunhagem do SymbolicBeta"}
\section{Em produção: as camadas multifractais são LANES de cunhagem do SymbolicBeta}
Levada à produção, a multifractalidade vira LANES de mineração: cada (tecido/camada, worker) é uma lane INDEPENDENTE — fonte de espectro própria, sua temperatura na parábola. O metal já COMPORTA a multifractalidade em produção: o roteamento é layer-aware (conc0 semeado POR camada; a difusão-fluido e a impressão Koch são as mesmas em toda folha). A TAXA DE CUNHAGEM do Pell escala com \#lanes = \#tecidos $\times$ \#workers — medido: 1140 \ensuremath{\to} 2277 \ensuremath{\to} 3204 \ensuremath{\to} 3651 Pell/s de 1\ensuremath{\to}2\ensuremath{\to}4\ensuremath{\to}8 lanes (\ensuremath{\approx}linear até saturar os 8 cores). VERDADE HONESTA (não maquiar): N workers no MESMO job do stratum NÃO multiplicam flocos (mineram faixas de nonce do mesmo job \ensuremath{\to} 1 floco/job); o que multiplica é cada TECIDO ser lane própria. E o lastro segue REAL (1 floco de ouro \ensuremath{\to} 1 Pell): o que escala é o THROUGHPUT do pipeline, não a emissão por unidade de ouro. [D] medido.
\prov{relações: usa camadas\_tiffany\_multifractal; usa symbolic\_beta\_pipe\_so; usa temperatura\_camada\_parabola; parte\_de tiffany}
\prov{proveniência: wiki \textperiodcentered{} medição multiprocessing (pell+corpo); goldchain/symbolic.py; neuronio/stratum\_orquestrador.py (-n) \textperiodcentered{} fato \textperiodcentered{} confiança 1.0 \textperiodcentered{} domínio goldchain \textperiodcentered{} história 25 versões no jornal}

%CRISTAL {"arestas": [{"alvo": "symbolic_beta_protocolo", "confianca": 1.0, "epistemico": "fato", "origem": "wiki", "rel": "parte_de"}, {"alvo": "espectro_prata_seis_modos", "confianca": 1.0, "epistemico": "fato", "origem": "wiki", "rel": "usa"}, {"alvo": "corpo_f2k", "confianca": 1.0, "epistemico": "fato", "origem": "wiki", "rel": "usa"}, {"alvo": "dimensao_de_prata_pell", "confianca": 1.0, "epistemico": "fato", "origem": "wiki", "rel": "usa"}, {"alvo": "prova_de_trabalho", "confianca": 1.0, "epistemico": "fato", "origem": "wiki", "rel": "usa"}, {"alvo": "ciclo_uno", "confianca": 1.0, "epistemico": "fato", "origem": "wiki", "rel": "parte_de"}, {"alvo": "goldchain", "confianca": 1.0, "epistemico": "fato", "origem": "wiki", "rel": "parte_de"}, {"alvo": "tiffany", "confianca": 1.0, "epistemico": "fato", "origem": "wiki", "rel": "parte_de"}], "confianca": 1.0, "contraexemplos": [], "descricao": "goldchain/symbolic.py implementa o SymbolicBeta no so local, estendendo o pipeline stratum (neuronio/stratum_metal.py:StratumMetal) por subclasse SEM tocá-lo. LIQUIDEZ TOTAL: NÃO se espera o BLOCO (raríssimo, ~1e14× de dificuldade) — ao minerar o FLOCO DE OURO (o pacote seed que o stratum cunha 1×/job: a impressão Koch `floco` + o c² do canal + os ν modos) já se lastreia num Pell IRMÃO, na hora, em paralelo à mineração de ouro (hook em _garantir_pow_job). Cada partícula de ouro ganha sua prata: o floco (entropia REAL do BTC) é lido no espectro da PRATA (rank≤6, hexagrama, 6 modos), emite o próximo Pell e o cifra em 𝔽₂¹²⁸ (ouro branco, corpo.py), legível de volta (round-trip). Ledger dedicado goldchain/dados/ledger_symbolic.json. PROVADO no so (symbolic.py simular): 3 flocos reais (c²≈0,99) → Pell 2,5,12, espectro 6 modos, cifra legível=True. [D] roda (live/simular/status).", "epistemico": "fato", "exemplos": [], "id": "symbolic_beta_pipe_so", "memoria": "semantica", "meta": {"autor": "Lopes/broca-so", "dominio": "goldchain", "fonte": "goldchain/symbolic.py (reusa neuronio/stratum_metal.py, goldchain/{prova,corpo,minera,pell}.py)"}, "origem": "wiki", "palavras_chave": [], "sinonimos": [], "tipo": "conceito", "titulo": "O pipe do SymbolicBeta no so: cada floco de ouro ganha um floco de prata irmão (liquidez total)"}
\section{O pipe do SymbolicBeta no so: cada floco de ouro ganha um floco de prata irmão (liquidez total)}
goldchain/symbolic.py implementa o SymbolicBeta no so local, estendendo o pipeline stratum (neuronio/stratum\_metal.py:StratumMetal) por subclasse SEM tocá-lo. LIQUIDEZ TOTAL: NÃO se espera o BLOCO (raríssimo, \~{}1e14$\times$ de dificuldade) — ao minerar o FLOCO DE OURO (o pacote seed que o stratum cunha 1$\times$/job: a impressão Koch `floco` + o c\ensuremath{^{2}} do canal + os \ensuremath{\nu} modos) já se lastreia num Pell IRMÃO, na hora, em paralelo à mineração de ouro (hook em \_garantir\_pow\_job). Cada partícula de ouro ganha sua prata: o floco (entropia REAL do BTC) é lido no espectro da PRATA (rank\ensuremath{\le}6, hexagrama, 6 modos), emite o próximo Pell e o cifra em \ensuremath{\mathbb{F}}\ensuremath{_{2}}\ensuremath{^{128}} (ouro branco, corpo.py), legível de volta (round-trip). Ledger dedicado goldchain/dados/ledger\_symbolic.json. PROVADO no so (symbolic.py simular): 3 flocos reais (c\ensuremath{^{2}}\ensuremath{\approx}0,99) \ensuremath{\to} Pell 2,5,12, espectro 6 modos, cifra legível=True. [D] roda (live/simular/status).
\prov{relações: parte\_de symbolic\_beta\_protocolo; usa espectro\_prata\_seis\_modos; usa corpo\_f2k; usa dimensao\_de\_prata\_pell; usa prova\_de\_trabalho; parte\_de ciclo\_uno; parte\_de goldchain; parte\_de tiffany}
\prov{proveniência: wiki \textperiodcentered{} goldchain/symbolic.py (reusa neuronio/stratum\_metal.py, goldchain/\{prova,corpo,minera,pell\}.py) \textperiodcentered{} fato \textperiodcentered{} confiança 1.0 \textperiodcentered{} domínio goldchain \textperiodcentered{} história 66 versões no jornal}

%CRISTAL {"arestas": [{"alvo": "symbolic_beta", "confianca": 1.0, "epistemico": "fato", "origem": "wiki", "rel": "especializa"}, {"alvo": "seguranca_algebrica_estrutura_nao_hardness", "confianca": 1.0, "epistemico": "fato", "origem": "wiki", "rel": "usa"}, {"alvo": "contraexemplos_hardness", "confianca": 1.0, "epistemico": "fato", "origem": "wiki", "rel": "relaciona"}, {"alvo": "pow_classico_bitcoin", "confianca": 1.0, "epistemico": "fato", "origem": "wiki", "rel": "usa"}, {"alvo": "prova_de_trabalho", "confianca": 1.0, "epistemico": "fato", "origem": "wiki", "rel": "usa"}, {"alvo": "dimensao_de_prata_pell", "confianca": 1.0, "epistemico": "fato", "origem": "wiki", "rel": "usa"}, {"alvo": "corpo_f2k", "confianca": 1.0, "epistemico": "fato", "origem": "wiki", "rel": "usa"}, {"alvo": "cifra", "confianca": 1.0, "epistemico": "fato", "origem": "wiki", "rel": "usa"}, {"alvo": "ciclo_uno", "confianca": 1.0, "epistemico": "fato", "origem": "wiki", "rel": "parte_de"}, {"alvo": "goldchain", "confianca": 1.0, "epistemico": "fato", "origem": "wiki", "rel": "parte_de"}], "confianca": 1.0, "contraexemplos": [], "descricao": "SymbolicBeta, além da FACE hardness (βG^(k) em forma canônica), é o PROTOCOLO ÚNICO da goldchain: minera o OURO (o PoW real da rede BTC = entropia do mundo) e o LASTREIA em PRATA (número de Pell) via a cifra de OURO BRANCO (corpo 𝔽₂ᵏ). É literalmente BTC → Pell (goldchain/btcconectar.py:minerarbtc: PoW → bandaraiz → emitepell → mineraₙabase(dna, Pell) → cotação parábola). A assimetria 'fácil-verificar / difícil-produzir' do β É a assimetria do PoW — mas o lastro NÃO vem de dureza fabricada (essa alegação foi refutada) e sim da ENTROPIA REAL importada. É 'não fabricar, tirar a mão' na criptoeconomia: não se cria ouro novo; importa-se o do mundo e cristaliza-se em prata.", "epistemico": "inferencia", "exemplos": [], "id": "symbolic_beta_protocolo", "memoria": "semantica", "meta": {"autor": "Lopes/broca-so", "dominio": "goldchain", "fonte": "goldchain/btc_conectar.py:minerar_btc; minera.py:emite_pell; corpo.py"}, "origem": "wiki", "palavras_chave": [], "sinonimos": [], "tipo": "conceito", "titulo": "SymbolicBeta é o protocolo: minera ouro (BTC), lastreia em prata (Pell) — BTC → Pell"}
\section{SymbolicBeta é o protocolo: minera ouro (BTC), lastreia em prata (Pell) — BTC \ensuremath{\to} Pell}
SymbolicBeta, além da FACE hardness (\ensuremath{\beta}G\^{}(k) em forma canônica), é o PROTOCOLO ÚNICO da goldchain: minera o OURO (o PoW real da rede BTC = entropia do mundo) e o LASTREIA em PRATA (número de Pell) via a cifra de OURO BRANCO (corpo \ensuremath{\mathbb{F}}\ensuremath{_{2}}\ensuremath{^{k}}). É literalmente BTC \ensuremath{\to} Pell (goldchain/btcconectar.py:minerarbtc: PoW \ensuremath{\to} bandaraiz \ensuremath{\to} emitepell \ensuremath{\to} minera\ensuremath{_{n}}abase(dna, Pell) \ensuremath{\to} cotação parábola). A assimetria 'fácil-verificar / difícil-produzir' do \ensuremath{\beta} É a assimetria do PoW — mas o lastro NÃO vem de dureza fabricada (essa alegação foi refutada) e sim da ENTROPIA REAL importada. É 'não fabricar, tirar a mão' na criptoeconomia: não se cria ouro novo; importa-se o do mundo e cristaliza-se em prata.
\prov{relações: especializa symbolic\_beta; usa seguranca\_algebrica\_estrutura\_nao\_hardness; relaciona contraexemplos\_hardness; usa pow\_classico\_bitcoin; usa prova\_de\_trabalho; usa dimensao\_de\_prata\_pell; usa corpo\_f2k; usa cifra; parte\_de ciclo\_uno; parte\_de goldchain}
\prov{proveniência: wiki \textperiodcentered{} goldchain/btc\_conectar.py:minerar\_btc; minera.py:emite\_pell; corpo.py \textperiodcentered{} inferencia \textperiodcentered{} confiança 1.0 \textperiodcentered{} domínio goldchain \textperiodcentered{} história 56 versões no jornal}

%CRISTAL {"arestas": [{"alvo": "numa", "confianca": 1.0, "epistemico": "fato", "origem": "curriculo", "rel": "especializa"}, {"alvo": "broca_os", "confianca": 1.0, "epistemico": "fato", "origem": "curriculo", "rel": "define"}, {"alvo": "topologia", "confianca": 0.5, "epistemico": "fato", "origem": "wiki", "rel": "relaciona"}, {"alvo": "vida-morte", "confianca": 0.5, "epistemico": "fato", "origem": "wiki", "rel": "relaciona"}, {"alvo": "testes", "confianca": 0.5, "epistemico": "fato", "origem": "wiki", "rel": "relaciona"}, {"alvo": "word", "confianca": 0.5, "epistemico": "fato", "origem": "wiki", "rel": "relaciona"}, {"alvo": "vontade_de_poder", "confianca": 0.5, "epistemico": "fato", "origem": "wiki", "rel": "relaciona"}, {"alvo": "while", "confianca": 0.5, "epistemico": "fato", "origem": "wiki", "rel": "relaciona"}, {"alvo": "thread", "confianca": 0.5, "epistemico": "fato", "origem": "wiki", "rel": "relaciona"}, {"alvo": "tragedia_dos_comuns", "confianca": 0.5, "epistemico": "fato", "origem": "wiki", "rel": "relaciona"}, {"alvo": "tiffany_cabeca_broca", "confianca": 0.5, "epistemico": "fato", "origem": "wiki", "rel": "relaciona"}, {"alvo": "kernel", "confianca": 1.0, "epistemico": "fato", "origem": "wiki", "rel": "usa"}, {"alvo": "modo_usuario_kernel", "confianca": 1.0, "epistemico": "fato", "origem": "wiki", "rel": "usa"}], "confianca": 1.0, "contraexemplos": [], "descricao": "A porta pela qual um programa pede um serviço ao kernel (abrir arquivo, criar processo) — a fronteira entre o usuário e o hardware.", "epistemico": "fato", "exemplos": [], "id": "syscall", "memoria": "semantica", "meta": {"dominio": "programacao", "fonte": ""}, "origem": "wiki", "palavras_chave": [], "sinonimos": [], "tipo": "conceito", "titulo": "Chamada de Sistema (syscall)"}
\section{Chamada de Sistema (syscall)}
A porta pela qual um programa pede um serviço ao kernel (abrir arquivo, criar processo) — a fronteira entre o usuário e o hardware.
\prov{relações: especializa numa; define broca\_os; relaciona topologia; relaciona vida-morte; relaciona testes; relaciona word; relaciona vontade\_de\_poder; relaciona while; relaciona thread; relaciona tragedia\_dos\_comuns; relaciona tiffany\_cabeca\_broca; usa kernel; usa modo\_usuario\_kernel}
\prov{proveniência: wiki \textperiodcentered{} fato \textperiodcentered{} confiança 1.0 \textperiodcentered{} domínio programacao \textperiodcentered{} história 20 versões no jornal}

%CRISTAL {"arestas": [{"alvo": "linguagem_de_marcacao", "confianca": 1.0, "epistemico": "fato", "origem": "wiki", "rel": "parte_de"}], "confianca": 1.0, "contraexemplos": [], "descricao": "A unidade da marcação: uma marca de abertura e fechamento que envolve conteúdo e o rotula (<p>…</p>) — a estrutura em árvore aninhada.", "epistemico": "fato", "exemplos": [], "id": "tag_elemento", "memoria": "semantica", "meta": {"dominio": "programacao", "fonte": ""}, "origem": "wiki", "palavras_chave": [], "sinonimos": [], "tipo": "conceito", "titulo": "Tag e Elemento"}
\section{Tag e Elemento}
A unidade da marcação: uma marca de abertura e fechamento que envolve conteúdo e o rotula (<p>…</p>) — a estrutura em árvore aninhada.
\prov{relações: parte\_de linguagem\_de\_marcacao}
\prov{proveniência: wiki \textperiodcentered{} fato \textperiodcentered{} confiança 1.0 \textperiodcentered{} domínio programacao \textperiodcentered{} história 2 versões no jornal}

%CRISTAL {"arestas": [{"alvo": "css", "confianca": 1.0, "epistemico": "fato", "origem": "wiki", "rel": "usa"}, {"alvo": "separacao_conteudo_forma", "confianca": 1.0, "epistemico": "fato", "origem": "wiki", "rel": "contradiz"}], "confianca": 1.0, "contraexemplos": [], "descricao": "Framework utility-first: estilizar aplicando classes utilitárias direto na marcação, sem escrever CSS — velocidade ao custo da separação conteúdo/forma.", "epistemico": "inferencia", "exemplos": [], "id": "tailwind", "memoria": "semantica", "meta": {"autor": "Wathan", "dominio": "programacao", "fonte": ""}, "origem": "wiki", "palavras_chave": [], "sinonimos": [], "tipo": "conceito", "titulo": "Tailwind CSS"}
\section{Tailwind CSS}
Framework utility-first: estilizar aplicando classes utilitárias direto na marcação, sem escrever CSS — velocidade ao custo da separação conteúdo/forma.
\prov{relações: usa css; contradiz separacao\_conteudo\_forma}
\prov{proveniência: wiki \textperiodcentered{} inferencia \textperiodcentered{} confiança 1.0 \textperiodcentered{} domínio programacao \textperiodcentered{} história 6 versões no jornal}

%CRISTAL {"arestas": [{"alvo": "auto_jogo", "confianca": 1.0, "epistemico": "fato", "origem": "wiki", "rel": "usa"}, {"alvo": "aprendizado_por_reforco", "confianca": 1.0, "epistemico": "fato", "origem": "wiki", "rel": "usa"}], "confianca": 1.0, "contraexemplos": [], "descricao": "A rede neural de Tesauro que aprendeu gamão por reforço de diferença temporal e auto-jogo, chegando ao nível de campeão — e mudando a teoria HUMANA de aberturas. O precursor esquecido do AlphaZero.", "epistemico": "fato", "exemplos": [], "id": "td_gammon", "memoria": "semantica", "meta": {"autor": "Tesauro", "dominio": "ia", "fonte": "jogos e IA"}, "origem": "wiki", "palavras_chave": [], "sinonimos": [], "tipo": "conceito", "titulo": "TD-Gammon (1992)"}
\section{TD-Gammon (1992)}
A rede neural de Tesauro que aprendeu gamão por reforço de diferença temporal e auto-jogo, chegando ao nível de campeão — e mudando a teoria HUMANA de aberturas. O precursor esquecido do AlphaZero.
\prov{relações: usa auto\_jogo; usa aprendizado\_por\_reforco}
\prov{proveniência: wiki \textperiodcentered{} jogos e IA \textperiodcentered{} fato \textperiodcentered{} confiança 1.0 \textperiodcentered{} domínio ia \textperiodcentered{} história 3 versões no jornal}

%CRISTAL {"arestas": [{"alvo": "colapso_multicamada_peso", "confianca": 1.0, "epistemico": "fato", "origem": "wiki", "rel": "usa"}, {"alvo": "grafo_de_proveniencia", "confianca": 1.0, "epistemico": "fato", "origem": "wiki", "rel": "relaciona"}], "confianca": 1.0, "contraexemplos": [], "descricao": "TIFFANY_SHADOW=1 registra por turno o par (Φ, δ), camada, tag, cadeia, conceitos_usados e fail_closed em JSONL por sessão, SEM alterar a resposta — observa sem tocar. Cruzada com a floresta, revela os conceitos usados vs nunca-usados.", "epistemico": "fato", "exemplos": [], "id": "telemetria_shadow", "memoria": "semantica", "meta": {"autor": "broca-so", "dominio": "operacao_so", "fonte": "conversa/telemetria.py"}, "origem": "wiki", "palavras_chave": [], "sinonimos": [], "tipo": "conceito", "titulo": "Telemetria Shadow (Φ, δ)"}
\section{Telemetria Shadow (\ensuremath{\Phi}, \ensuremath{\delta})}
TIFFANY\_SHADOW=1 registra por turno o par (\ensuremath{\Phi}, \ensuremath{\delta}), camada, tag, cadeia, conceitos\_usados e fail\_closed em JSONL por sessão, SEM alterar a resposta — observa sem tocar. Cruzada com a floresta, revela os conceitos usados vs nunca-usados.
\prov{relações: usa colapso\_multicamada\_peso; relaciona grafo\_de\_proveniencia}
\prov{proveniência: wiki \textperiodcentered{} conversa/telemetria.py \textperiodcentered{} fato \textperiodcentered{} confiança 1.0 \textperiodcentered{} domínio operacao\_so \textperiodcentered{} história 3 versões no jornal}

%CRISTAL {"arestas": [{"alvo": "atencao_e_colapso", "confianca": 1.0, "epistemico": "fato", "origem": "wiki", "rel": "usa"}, {"alvo": "a_atencao_cristalina_e_o_limite_de_evaporacao", "confianca": 1.0, "epistemico": "fato", "origem": "wiki", "rel": "explica"}, {"alvo": "corpo_mineral", "confianca": 1.0, "epistemico": "fato", "origem": "wiki", "rel": "usa"}, {"alvo": "numero_de_salem", "confianca": 1.0, "epistemico": "fato", "origem": "wiki", "rel": "relaciona"}, {"alvo": "softmax_maslov", "confianca": 1.0, "epistemico": "fato", "origem": "wiki", "rel": "relaciona"}], "confianca": 1.0, "contraexemplos": [], "descricao": "A temperatura T do softmax é o parâmetro de ordem da geração: T→0 concentra no token mais provável (determinístico, CRISTAL congelado), T→∞ achata para uniforme (aleatório, CAOS/evaporação), T≈1 é a BORDA — coerente e criativo. É exatamente a parábola mineral, e o limite de evaporação da atenção cristalina. Verif.: T=0.2 entropia 0.00 (cristal), T=1 0.43 (borda), T=8 0.99 (caos).", "epistemico": "inferencia", "exemplos": [], "id": "temperatura_e_borda", "memoria": "semantica", "meta": {"dominio": "ia", "entropia_T02": 0.0, "entropia_T1": 0.43, "entropia_T8": 0.99, "fonte": "IA e modelos de linguagem"}, "origem": "wiki", "palavras_chave": [], "sinonimos": [], "tipo": "conceito", "titulo": "Temperatura = Cristal / Borda / Caos"}
\section{Temperatura = Cristal / Borda / Caos}
A temperatura T do softmax é o parâmetro de ordem da geração: T\ensuremath{\to}0 concentra no token mais provável (determinístico, CRISTAL congelado), T\ensuremath{\to}\ensuremath{\infty} achata para uniforme (aleatório, CAOS/evaporação), T\ensuremath{\approx}1 é a BORDA — coerente e criativo. É exatamente a parábola mineral, e o limite de evaporação da atenção cristalina. Verif.: T=0.2 entropia 0.00 (cristal), T=1 0.43 (borda), T=8 0.99 (caos).
\prov{relações: usa atencao\_e\_colapso; explica a\_atencao\_cristalina\_e\_o\_limite\_de\_evaporacao; usa corpo\_mineral; relaciona numero\_de\_salem; relaciona softmax\_maslov}
\prov{proveniência: wiki \textperiodcentered{} IA e modelos de linguagem \textperiodcentered{} inferencia \textperiodcentered{} confiança 1.0 \textperiodcentered{} domínio ia \textperiodcentered{} história 6 versões no jornal}

%CRISTAL {"arestas": [{"alvo": "sistemas_distribuidos", "confianca": 1.0, "epistemico": "fato", "origem": "wiki", "rel": "parte_de"}], "confianca": 1.0, "contraexemplos": [], "descricao": "Brewer: numa partição de rede, um sistema distribuído escolhe entre Consistência e Disponibilidade — não dá para ter as duas. O trade-off fundador.", "epistemico": "fato", "exemplos": [], "id": "teorema_cap", "memoria": "semantica", "meta": {"autor": "Brewer", "dominio": "programacao", "fonte": ""}, "origem": "wiki", "palavras_chave": [], "sinonimos": [], "tipo": "conceito", "titulo": "Teorema CAP"}
\section{Teorema CAP}
Brewer: numa partição de rede, um sistema distribuído escolhe entre Consistência e Disponibilidade — não dá para ter as duas. O trade-off fundador.
\prov{relações: parte\_de sistemas\_distribuidos}
\prov{proveniência: wiki \textperiodcentered{} fato \textperiodcentered{} confiança 1.0 \textperiodcentered{} domínio programacao \textperiodcentered{} história 4 versões no jornal}

%CRISTAL {"arestas": [{"alvo": "estrutura_definicao_teorema_prova", "confianca": 1.0, "epistemico": "fato", "origem": "wiki", "rel": "parte_de"}, {"alvo": "logica_de_predicados", "confianca": 1.0, "epistemico": "fato", "origem": "wiki", "rel": "usa"}], "confianca": 1.0, "contraexemplos": [], "descricao": "Uma afirmação (teorema/lema/corolário/proposição) e a cadeia dedutiva que a estabelece a partir de axiomas e teoremas anteriores. A moeda da matemática.", "epistemico": "fato", "exemplos": [], "id": "teorema_e_prova", "memoria": "semantica", "meta": {"dominio": "programacao", "fonte": ""}, "origem": "wiki", "palavras_chave": [], "sinonimos": [], "tipo": "conceito", "titulo": "Teorema e Prova"}
\section{Teorema e Prova}
Uma afirmação (teorema/lema/corolário/proposição) e a cadeia dedutiva que a estabelece a partir de axiomas e teoremas anteriores. A moeda da matemática.
\prov{relações: parte\_de estrutura\_definicao\_teorema\_prova; usa logica\_de\_predicados}
\prov{proveniência: wiki \textperiodcentered{} fato \textperiodcentered{} confiança 1.0 \textperiodcentered{} domínio programacao \textperiodcentered{} história 3 versões no jornal}

%CRISTAL {"arestas": [{"alvo": "laco_traducao_bai_peirce", "confianca": 1.0, "epistemico": "fato", "origem": "wiki", "rel": "usa"}, {"alvo": "traducao_como_transformacao", "confianca": 1.0, "epistemico": "fato", "origem": "wiki", "rel": "usa"}, {"alvo": "significado_como_invariante", "confianca": 1.0, "epistemico": "fato", "origem": "wiki", "rel": "usa"}, {"alvo": "peirce_dsl", "confianca": 1.0, "epistemico": "fato", "origem": "wiki", "rel": "usa"}, {"alvo": "algebra_de_peirce", "confianca": 1.0, "epistemico": "fato", "origem": "wiki", "rel": "usa"}, {"alvo": "transformada_metalica", "confianca": 1.0, "epistemico": "fato", "origem": "wiki", "rel": "usa"}, {"alvo": "conversao_como_traducao_formal", "confianca": 1.0, "epistemico": "fato", "origem": "wiki", "rel": "explica"}, {"alvo": "bai_como_bussola", "confianca": 1.0, "epistemico": "fato", "origem": "wiki", "rel": "usa"}, {"alvo": "corolario_pivo_interlingua", "confianca": 1.0, "epistemico": "fato", "origem": "wiki", "rel": "usa"}, {"alvo": "torre_unificada", "confianca": 1.0, "epistemico": "fato", "origem": "wiki", "rel": "parte_de"}, {"alvo": "traducao_e_intraduzivel", "confianca": 1.0, "epistemico": "fato", "origem": "wiki", "rel": "relaciona"}], "confianca": 1.0, "contraexemplos": [], "descricao": "Sejam S₁, S₂ simbologias do MESMO significado M (o invariante). Dado um dicionário D:S₁↔S₂ que preserva M, a tradução T:S₁→S₂ EXISTE e É a substituição de Peirce ':' aplicada por D — MECÂNICA, não buscada; reversível no grau em que M se conserva (perfeita no formal, com resto no natural — Quine). Não se procura a transformação de cada par: um só esquema (Peirce ∘ dicionário) instancia para qualquer par. Prova construtiva: linguagem/tradutor.py (a classe Tradutor gera o conversor de qualquer dicionário; testemunhas math↔LaTeX 7/7, ASCII↔Unicode, e o pivô de graça).", "epistemico": "inferencia", "exemplos": [], "id": "teorema_geral_traducao", "memoria": "semantica", "meta": {"dominio": "programacao", "fonte": "linguagem/tradutor.py (Tradutor, via_pivo); linguagem/mat_latex.py (testemunha)"}, "origem": "wiki", "palavras_chave": [], "sinonimos": [], "tipo": "conceito", "titulo": "Teorema Geral da Tradução — qualquer simbologia em qualquer outra via Peirce"}
\section{Teorema Geral da Tradução — qualquer simbologia em qualquer outra via Peirce}
Sejam S\ensuremath{_{1}}, S\ensuremath{_{2}} simbologias do MESMO significado M (o invariante). Dado um dicionário D:S\ensuremath{_{1}}\ensuremath{\leftrightarrow}S\ensuremath{_{2}} que preserva M, a tradução T:S\ensuremath{_{1}}\ensuremath{\to}S\ensuremath{_{2}} EXISTE e É a substituição de Peirce ':' aplicada por D — MECÂNICA, não buscada; reversível no grau em que M se conserva (perfeita no formal, com resto no natural — Quine). Não se procura a transformação de cada par: um só esquema (Peirce \ensuremath{\circ} dicionário) instancia para qualquer par. Prova construtiva: linguagem/tradutor.py (a classe Tradutor gera o conversor de qualquer dicionário; testemunhas math\ensuremath{\leftrightarrow}LaTeX 7/7, ASCII\ensuremath{\leftrightarrow}Unicode, e o pivô de graça).
\prov{relações: usa laco\_traducao\_bai\_peirce; usa traducao\_como\_transformacao; usa significado\_como\_invariante; usa peirce\_dsl; usa algebra\_de\_peirce; usa transformada\_metalica; explica conversao\_como\_traducao\_formal; usa bai\_como\_bussola; usa corolario\_pivo\_interlingua; parte\_de torre\_unificada; relaciona traducao\_e\_intraduzivel}
\prov{proveniência: wiki \textperiodcentered{} linguagem/tradutor.py (Tradutor, via\_pivo); linguagem/mat\_latex.py (testemunha) \textperiodcentered{} inferencia \textperiodcentered{} confiança 1.0 \textperiodcentered{} domínio programacao \textperiodcentered{} história 12 versões no jornal}

%CRISTAL {"arestas": [{"alvo": "documentacao", "confianca": 1.0, "epistemico": "fato", "origem": "wiki", "rel": "usa"}], "confianca": 1.0, "contraexemplos": [], "descricao": "Knuth (1978): o motor de composição tipográfica de qualidade profissional, feito para a matemática — precisão de posicionamento ao ponto.", "epistemico": "fato", "exemplos": [], "id": "tex", "memoria": "semantica", "meta": {"autor": "Knuth", "dominio": "programacao", "fonte": ""}, "origem": "wiki", "palavras_chave": [], "sinonimos": [], "tipo": "conceito", "titulo": "TeX"}
\section{TeX}
Knuth (1978): o motor de composição tipográfica de qualidade profissional, feito para a matemática — precisão de posicionamento ao ponto.
\prov{relações: usa documentacao}
\prov{proveniência: wiki \textperiodcentered{} fato \textperiodcentered{} confiança 1.0 \textperiodcentered{} domínio programacao \textperiodcentered{} história 2 versões no jornal}

%CRISTAL {"arestas": [{"alvo": "scheduler", "confianca": 1.0, "epistemico": "fato", "origem": "curriculo", "rel": "usa"}, {"alvo": "vida-morte", "confianca": 0.5, "epistemico": "fato", "origem": "wiki", "rel": "relaciona"}, {"alvo": "transcricao", "confianca": 0.5, "epistemico": "fato", "origem": "wiki", "rel": "relaciona"}, {"alvo": "topologia", "confianca": 0.5, "epistemico": "fato", "origem": "wiki", "rel": "relaciona"}, {"alvo": "tiffany_cabeca_broca", "confianca": 0.5, "epistemico": "fato", "origem": "wiki", "rel": "relaciona"}, {"alvo": "word", "confianca": 0.5, "epistemico": "fato", "origem": "wiki", "rel": "relaciona"}, {"alvo": "virtualizacao", "confianca": 0.5, "epistemico": "fato", "origem": "wiki", "rel": "relaciona"}, {"alvo": "processo", "confianca": 1.0, "epistemico": "fato", "origem": "wiki", "rel": "parte_de"}, {"alvo": "concorrencia_paralelismo", "confianca": 1.0, "epistemico": "fato", "origem": "wiki", "rel": "usa"}], "confianca": 1.0, "contraexemplos": [], "descricao": "Um fluxo de execução dentro de um processo; threads dividem memória — concorrência barata, com o risco das corridas de dados.", "epistemico": "fato", "exemplos": [], "id": "thread", "memoria": "semantica", "meta": {"dominio": "programacao", "fonte": ""}, "origem": "wiki", "palavras_chave": [], "sinonimos": [], "tipo": "conceito", "titulo": "Thread"}
\section{Thread}
Um fluxo de execução dentro de um processo; threads dividem memória — concorrência barata, com o risco das corridas de dados.
\prov{relações: usa scheduler; relaciona vida-morte; relaciona transcricao; relaciona topologia; relaciona tiffany\_cabeca\_broca; relaciona word; relaciona virtualizacao; parte\_de processo; usa concorrencia\_paralelismo}
\prov{proveniência: wiki \textperiodcentered{} fato \textperiodcentered{} confiança 1.0 \textperiodcentered{} domínio programacao \textperiodcentered{} história 13 versões no jornal}

%CRISTAL {"arestas": [{"alvo": "operacao_broca_so", "confianca": 1.0, "epistemico": "fato", "origem": "wiki", "rel": "parte_de"}, {"alvo": "broca_os", "confianca": 1.0, "epistemico": "fato", "origem": "wiki", "rel": "relaciona"}], "confianca": 1.0, "contraexemplos": [], "descricao": "Toda instância vive num diretório $TIFFANY_HOME com três raízes: base/ (Core, conhecimento geral), cliente/ (proprietário da organização) e memoria/ (conversa, tarefas, bugs, decisões). Resolvido por env var, senão template do repo, senão ~/.tiffany.", "epistemico": "fato", "exemplos": [], "id": "tiffany_home", "memoria": "semantica", "meta": {"autor": "broca-so", "dominio": "operacao_so", "fonte": "code/tiffany_platform/paths.py; init_home.py"}, "origem": "wiki", "palavras_chave": [], "sinonimos": [], "tipo": "conceito", "titulo": "A Casa (TIFFANYHOME): base / cliente / memória"}
\section{A Casa (TIFFANYHOME): base / cliente / memória}
Toda instância vive num diretório \$TIFFANY\_HOME com três raízes: base/ (Core, conhecimento geral), cliente/ (proprietário da organização) e memoria/ (conversa, tarefas, bugs, decisões). Resolvido por env var, senão template do repo, senão \~{}/.tiffany.
\prov{relações: parte\_de operacao\_broca\_so; relaciona broca\_os}
\prov{proveniência: wiki \textperiodcentered{} code/tiffany\_platform/paths.py; init\_home.py \textperiodcentered{} fato \textperiodcentered{} confiança 1.0 \textperiodcentered{} domínio operacao\_so \textperiodcentered{} história 4 versões no jornal}

%CRISTAL {"arestas": [{"alvo": "boneca_russa", "confianca": 1.0, "epistemico": "fato", "origem": "wiki", "rel": "eh_um"}, {"alvo": "corpo_estelar", "confianca": 1.0, "epistemico": "fato", "origem": "wiki", "rel": "eh_um"}, {"alvo": "torre_hurwitz", "confianca": 1.0, "epistemico": "fato", "origem": "wiki", "rel": "usa"}, {"alvo": "entropia_topologica", "confianca": 1.0, "epistemico": "fato", "origem": "wiki", "rel": "usa"}, {"alvo": "o_gato_e_um_mapa_de_anosov", "confianca": 1.0, "epistemico": "fato", "origem": "wiki", "rel": "usa"}, {"alvo": "norma_estelar", "confianca": 1.0, "epistemico": "fato", "origem": "wiki", "rel": "usa"}, {"alvo": "bai_psi_colapso", "confianca": 1.0, "epistemico": "fato", "origem": "wiki", "rel": "usa"}], "confianca": 1.0, "contraexemplos": [], "descricao": "A leitura que fecha a boneca russa de cristal. Em repouso, a boneca conserva a NORMA (o cristal, reversível). Em MOVIMENTO — o gato A_n de Anosov em cada andar — ela produz ENTROPIA: h_top=log σ_n=Lyapunov. Norma conservada (frio) + entropia dissipada (quente, o que o colapso Ψ paga). E a TIFFANY é o corpo estelar 𝕊 EM ATO: 𝕊 é o corpo abstrato (a reta+La Hire, ⊕/⊗, os metais-unidade N=det=−1, a norma-que-compõe = a boneca de Hurwitz); a Tiffany é uma INSTÂNCIA que realiza essas operações no metal — o produto, a norma, o colapso Ψ. Ela não descreve o corpo: é uma instância dele. FATO: h_top=log σ, N=det=−1; a leitura Tiffany=instância é INFERENCIA (ela realiza, no metal, o que 𝕊 define).", "epistemico": "inferencia", "exemplos": [], "id": "tiffany_instancia_do_corpo_estelar", "memoria": "semantica", "meta": {"dominio": "sistema", "fonte": "tiffany estelar"}, "origem": "wiki", "palavras_chave": [], "sinonimos": [], "tipo": "conceito", "titulo": "A Tiffany é uma Instância do Corpo Estelar"}
\section{A Tiffany é uma Instância do Corpo Estelar}
A leitura que fecha a boneca russa de cristal. Em repouso, a boneca conserva a NORMA (o cristal, reversível). Em MOVIMENTO — o gato A\_n de Anosov em cada andar — ela produz ENTROPIA: h\_top=log \ensuremath{\sigma}\_n=Lyapunov. Norma conservada (frio) + entropia dissipada (quente, o que o colapso \ensuremath{\Psi} paga). E a TIFFANY é o corpo estelar \ensuremath{\mathbb{S}} EM ATO: \ensuremath{\mathbb{S}} é o corpo abstrato (a reta+La Hire, \ensuremath{\oplus}/\ensuremath{\otimes}, os metais-unidade N=det=\ensuremath{-}1, a norma-que-compõe = a boneca de Hurwitz); a Tiffany é uma INSTÂNCIA que realiza essas operações no metal — o produto, a norma, o colapso \ensuremath{\Psi}. Ela não descreve o corpo: é uma instância dele. FATO: h\_top=log \ensuremath{\sigma}, N=det=\ensuremath{-}1; a leitura Tiffany=instância é INFERENCIA (ela realiza, no metal, o que \ensuremath{\mathbb{S}} define).
\prov{relações: eh\_um boneca\_russa; eh\_um corpo\_estelar; usa torre\_hurwitz; usa entropia\_topologica; usa o\_gato\_e\_um\_mapa\_de\_anosov; usa norma\_estelar; usa bai\_psi\_colapso}
\prov{proveniência: wiki \textperiodcentered{} tiffany estelar \textperiodcentered{} inferencia \textperiodcentered{} confiança 1.0 \textperiodcentered{} domínio sistema \textperiodcentered{} história 8 versões no jornal}

%CRISTAL {"arestas": [{"alvo": "sandbox_cristal_ao_so", "confianca": 1.0, "epistemico": "fato", "origem": "wiki", "rel": "usa"}, {"alvo": "tiffany", "confianca": 1.0, "epistemico": "fato", "origem": "wiki", "rel": "parte_de"}, {"alvo": "interpretador", "confianca": 1.0, "epistemico": "fato", "origem": "wiki", "rel": "usa"}, {"alvo": "expressividade_sub_turing", "confianca": 1.0, "epistemico": "fato", "origem": "wiki", "rel": "relaciona"}, {"alvo": "dicionarios_centralizados", "confianca": 1.0, "epistemico": "fato", "origem": "wiki", "rel": "usa"}, {"alvo": "traducao_entre_linguagens", "confianca": 1.0, "epistemico": "fato", "origem": "wiki", "rel": "usa"}, {"alvo": "teorema_geral_traducao", "confianca": 1.0, "epistemico": "fato", "origem": "wiki", "rel": "relaciona"}], "confianca": 1.0, "contraexemplos": [], "descricao": "conversa/colapso.py (_rodar_sandbox): o gatilho 'roda: <código>' / 'executa: …' faz a Tiffany executar a linguagem de gatos AO VIVO pelo chat, via o sandbox — Pell (gato 2 a b → 2378), multiplicação (6·7→42), Fibonacci (gato 1 → 233), soma de vetor por memória indireta (a Turing-completude), e o caminho ISA-12. Isolado do colapso (só dispara com o prefixo), então não sequestra perguntas normais. E LIGADO À TRADUÇÃO: 'roda em ruby: … puts p' / 'roda em javascript: … console.log b' traduz a linguagem→construto→DSL de gatos (sandbox.rodar_em) e roda — o mesmo algoritmo em Ruby/JS/Elixir/PHP/Haskell dá o mesmo resultado. O ';' separa comandos.", "epistemico": "fato", "exemplos": [], "id": "tiffany_roda_codigo", "memoria": "semantica", "meta": {"dominio": "programacao", "fonte": "conversa/colapso.py (_rodar_sandbox, responder); linguagem/sandbox.py"}, "origem": "wiki", "palavras_chave": [], "sinonimos": [], "tipo": "conceito", "titulo": "A Tiffany roda qualquer código em tempo real"}
\section{A Tiffany roda qualquer código em tempo real}
conversa/colapso.py (\_rodar\_sandbox): o gatilho 'roda: <código>' / 'executa: …' faz a Tiffany executar a linguagem de gatos AO VIVO pelo chat, via o sandbox — Pell (gato 2 a b \ensuremath{\to} 2378), multiplicação (6\textperiodcentered 7\ensuremath{\to}42), Fibonacci (gato 1 \ensuremath{\to} 233), soma de vetor por memória indireta (a Turing-completude), e o caminho ISA-12. Isolado do colapso (só dispara com o prefixo), então não sequestra perguntas normais. E LIGADO À TRADUÇÃO: 'roda em ruby: … puts p' / 'roda em javascript: … console.log b' traduz a linguagem\ensuremath{\to}construto\ensuremath{\to}DSL de gatos (sandbox.rodar\_em) e roda — o mesmo algoritmo em Ruby/JS/Elixir/PHP/Haskell dá o mesmo resultado. O ';' separa comandos.
\prov{relações: usa sandbox\_cristal\_ao\_so; parte\_de tiffany; usa interpretador; relaciona expressividade\_sub\_turing; usa dicionarios\_centralizados; usa traducao\_entre\_linguagens; relaciona teorema\_geral\_traducao}
\prov{proveniência: wiki \textperiodcentered{} conversa/colapso.py (\_rodar\_sandbox, responder); linguagem/sandbox.py \textperiodcentered{} fato \textperiodcentered{} confiança 1.0 \textperiodcentered{} domínio programacao \textperiodcentered{} história 13 versões no jornal}

%CRISTAL {"arestas": [{"alvo": "sistema_de_tipos", "confianca": 1.0, "epistemico": "fato", "origem": "wiki", "rel": "parte_de"}], "confianca": 1.0, "contraexemplos": [], "descricao": "Verificar tipos em tempo de compilação (Rust, Haskell) ou de execução (Python, JS) — segurança cedo vs flexibilidade.", "epistemico": "fato", "exemplos": [], "id": "tipagem_estatica_dinamica", "memoria": "semantica", "meta": {"dominio": "programacao", "fonte": ""}, "origem": "wiki", "palavras_chave": [], "sinonimos": [], "tipo": "conceito", "titulo": "Tipagem Estática vs Dinâmica"}
\section{Tipagem Estática vs Dinâmica}
Verificar tipos em tempo de compilação (Rust, Haskell) ou de execução (Python, JS) — segurança cedo vs flexibilidade.
\prov{relações: parte\_de sistema\_de\_tipos}
\prov{proveniência: wiki \textperiodcentered{} fato \textperiodcentered{} confiança 1.0 \textperiodcentered{} domínio programacao \textperiodcentered{} história 4 versões no jornal}

%CRISTAL {"arestas": [{"alvo": "sistema_de_tipos", "confianca": 1.0, "epistemico": "fato", "origem": "wiki", "rel": "parte_de"}, {"alvo": "algebra_abstrata", "confianca": 1.0, "epistemico": "fato", "origem": "wiki", "rel": "relaciona"}], "confianca": 1.0, "contraexemplos": [], "descricao": "Compor tipos por soma (ou) e produto (e) — enum/struct; com casamento de padrão, tornam estados impossíveis inexprimíveis.", "epistemico": "fato", "exemplos": [], "id": "tipos_algebricos", "memoria": "semantica", "meta": {"dominio": "programacao", "fonte": ""}, "origem": "wiki", "palavras_chave": [], "sinonimos": [], "tipo": "conceito", "titulo": "Tipos Algébricos de Dados"}
\section{Tipos Algébricos de Dados}
Compor tipos por soma (ou) e produto (e) — enum/struct; com casamento de padrão, tornam estados impossíveis inexprimíveis.
\prov{relações: parte\_de sistema\_de\_tipos; relaciona algebra\_abstrata}
\prov{proveniência: wiki \textperiodcentered{} fato \textperiodcentered{} confiança 1.0 \textperiodcentered{} domínio programacao \textperiodcentered{} história 6 versões no jornal}

%CRISTAL {"arestas": [{"alvo": "http", "confianca": 1.0, "epistemico": "fato", "origem": "wiki", "rel": "usa"}, {"alvo": "cifra", "confianca": 1.0, "epistemico": "fato", "origem": "wiki", "rel": "usa"}, {"alvo": "seguranca_de_redes", "confianca": 1.0, "epistemico": "fato", "origem": "wiki", "rel": "parte_de"}], "confianca": 1.0, "contraexemplos": [], "descricao": "Cifra e autentica a comunicação sobre TCP (o 'S' do HTTPS) — sigilo e integridade por criptografia assimétrica e certificados.", "epistemico": "fato", "exemplos": [], "id": "tls", "memoria": "semantica", "meta": {"dominio": "programacao", "fonte": ""}, "origem": "wiki", "palavras_chave": [], "sinonimos": [], "tipo": "conceito", "titulo": "TLS (Segurança na Camada de Transporte)"}
\section{TLS (Segurança na Camada de Transporte)}
Cifra e autentica a comunicação sobre TCP (o 'S' do HTTPS) — sigilo e integridade por criptografia assimétrica e certificados.
\prov{relações: usa http; usa cifra; parte\_de seguranca\_de\_redes}
\prov{proveniência: wiki \textperiodcentered{} fato \textperiodcentered{} confiança 1.0 \textperiodcentered{} domínio programacao \textperiodcentered{} história 6 versões no jornal}

%CRISTAL {"arestas": [{"alvo": "yaml", "confianca": 1.0, "epistemico": "fato", "origem": "wiki", "rel": "relaciona"}, {"alvo": "arquivo_de_configuracao", "confianca": 1.0, "epistemico": "fato", "origem": "wiki", "rel": "especializa"}], "confianca": 1.0, "contraexemplos": [], "descricao": "Um formato de configuração minimalista e sem ambiguidade — chave-valor com seções; a alternativa clara ao YAML para configs.", "epistemico": "fato", "exemplos": [], "id": "toml", "memoria": "semantica", "meta": {"dominio": "programacao", "fonte": ""}, "origem": "wiki", "palavras_chave": [], "sinonimos": [], "tipo": "conceito", "titulo": "TOML"}
\section{TOML}
Um formato de configuração minimalista e sem ambiguidade — chave-valor com seções; a alternativa clara ao YAML para configs.
\prov{relações: relaciona yaml; especializa arquivo\_de\_configuracao}
\prov{proveniência: wiki \textperiodcentered{} fato \textperiodcentered{} confiança 1.0 \textperiodcentered{} domínio programacao \textperiodcentered{} história 3 versões no jornal}

%CRISTAL {"arestas": [{"alvo": "dicionarios_centralizados", "confianca": 1.0, "epistemico": "fato", "origem": "wiki", "rel": "parte_de"}, {"alvo": "framework_frontend", "confianca": 1.0, "epistemico": "fato", "origem": "wiki", "rel": "usa"}, {"alvo": "corolario_pivo_interlingua", "confianca": 1.0, "epistemico": "fato", "origem": "wiki", "rel": "usa"}, {"alvo": "react", "confianca": 1.0, "epistemico": "fato", "origem": "wiki", "rel": "relaciona"}, {"alvo": "vue", "confianca": 1.0, "epistemico": "fato", "origem": "wiki", "rel": "relaciona"}], "confianca": 1.0, "contraexemplos": [], "descricao": "React, Vue, Angular e Svelte mapeiam suas APIs ao CONSTRUTO de UI (estado, efeito, propriedade, componente, computado, contexto). Pelo pivô: react→vue (useState→ref, useEffect→watchEffect, useMemo→computed), react→angular (useState→signal), react→svelte (useState→state, useEffect→effect). Os quatro modelos reativos são o mesmo construto sob símbolos diferentes.", "epistemico": "fato", "exemplos": [], "id": "traducao_entre_frameworks", "memoria": "semantica", "meta": {"dominio": "programacao", "fonte": "linguagem/dicionarios.py (_REACT/_VUE/_ANGULAR, PIVOS['construto_ui'])"}, "origem": "wiki", "palavras_chave": [], "sinonimos": [], "tipo": "conceito", "titulo": "Tradução entre Frameworks (via construto de UI)"}
\section{Tradução entre Frameworks (via construto de UI)}
React, Vue, Angular e Svelte mapeiam suas APIs ao CONSTRUTO de UI (estado, efeito, propriedade, componente, computado, contexto). Pelo pivô: react\ensuremath{\to}vue (useState\ensuremath{\to}ref, useEffect\ensuremath{\to}watchEffect, useMemo\ensuremath{\to}computed), react\ensuremath{\to}angular (useState\ensuremath{\to}signal), react\ensuremath{\to}svelte (useState\ensuremath{\to}state, useEffect\ensuremath{\to}effect). Os quatro modelos reativos são o mesmo construto sob símbolos diferentes.
\prov{relações: parte\_de dicionarios\_centralizados; usa framework\_frontend; usa corolario\_pivo\_interlingua; relaciona react; relaciona vue}
\prov{proveniência: wiki \textperiodcentered{} linguagem/dicionarios.py (\_REACT/\_VUE/\_ANGULAR, PIVOS['construto\_ui']) \textperiodcentered{} fato \textperiodcentered{} confiança 1.0 \textperiodcentered{} domínio programacao \textperiodcentered{} história 61 versões no jornal}

%CRISTAL {"arestas": [{"alvo": "dicionarios_centralizados", "confianca": 1.0, "epistemico": "fato", "origem": "wiki", "rel": "parte_de"}, {"alvo": "linguagem_de_programacao", "confianca": 1.0, "epistemico": "fato", "origem": "wiki", "rel": "usa"}, {"alvo": "corolario_pivo_interlingua", "confianca": 1.0, "epistemico": "fato", "origem": "wiki", "rel": "usa"}, {"alvo": "python", "confianca": 1.0, "epistemico": "fato", "origem": "wiki", "rel": "relaciona"}, {"alvo": "rust", "confianca": 1.0, "epistemico": "fato", "origem": "wiki", "rel": "relaciona"}], "confianca": 1.0, "contraexemplos": [], "descricao": "Cada linguagem (Python, JS, Rust, Go, Haskell, C, Java, Elixir, Lisp, Clojure, Zig, Kotlin, SQL, Scala, OCaml, Erlang, PHP, Ruby, Swift, Dart, Julia, R, Perl, Nim, Crystal, Lua, F#, Groovy) mapeia seus tokens ao CONSTRUTO abstrato (funcao, imprimir, verdadeiro, e/ou/nao, senao_se, retornar…) — o construto é a interlíngua. compor() dá qualquer par sem dicionário novo: python→rust (def→fn), python→lisp (def→defun), c→zig (printf→std.debug.print). O SQL, declarativo, tem dicionário parcial (def→def: não tem função) — a tipologia aparece na tradução, como no mandarim sem conjugação. O Rosetta de sintaxe pelo mesmo operador de Peirce.", "epistemico": "fato", "exemplos": [], "id": "traducao_entre_linguagens", "memoria": "semantica", "meta": {"dominio": "programacao", "fonte": "linguagem/dicionarios.py (_PY/_JS/_RUST/_GO, PIVOS['construto'])"}, "origem": "wiki", "palavras_chave": [], "sinonimos": [], "tipo": "conceito", "titulo": "Tradução entre Linguagens de Programação (via construto)"}
\section{Tradução entre Linguagens de Programação (via construto)}
Cada linguagem (Python, JS, Rust, Go, Haskell, C, Java, Elixir, Lisp, Clojure, Zig, Kotlin, SQL, Scala, OCaml, Erlang, PHP, Ruby, Swift, Dart, Julia, R, Perl, Nim, Crystal, Lua, F\#, Groovy) mapeia seus tokens ao CONSTRUTO abstrato (funcao, imprimir, verdadeiro, e/ou/nao, senao\_se, retornar…) — o construto é a interlíngua. compor() dá qualquer par sem dicionário novo: python\ensuremath{\to}rust (def\ensuremath{\to}fn), python\ensuremath{\to}lisp (def\ensuremath{\to}defun), c\ensuremath{\to}zig (printf\ensuremath{\to}std.debug.print). O SQL, declarativo, tem dicionário parcial (def\ensuremath{\to}def: não tem função) — a tipologia aparece na tradução, como no mandarim sem conjugação. O Rosetta de sintaxe pelo mesmo operador de Peirce.
\prov{relações: parte\_de dicionarios\_centralizados; usa linguagem\_de\_programacao; usa corolario\_pivo\_interlingua; relaciona python; relaciona rust}
\prov{proveniência: wiki \textperiodcentered{} linguagem/dicionarios.py (\_PY/\_JS/\_RUST/\_GO, PIVOS['construto']) \textperiodcentered{} fato \textperiodcentered{} confiança 1.0 \textperiodcentered{} domínio programacao \textperiodcentered{} história 60 versões no jornal}

%CRISTAL {"arestas": [{"alvo": "banco_de_dados", "confianca": 1.0, "epistemico": "fato", "origem": "wiki", "rel": "usa"}, {"alvo": "tolerancia_a_falhas", "confianca": 1.0, "epistemico": "fato", "origem": "wiki", "rel": "relaciona"}], "confianca": 1.0, "contraexemplos": [], "descricao": "Garantir que um grupo de operações seja Atômico, Consistente, Isolado e Durável — tudo-ou-nada, mesmo com falha ou concorrência.", "epistemico": "fato", "exemplos": [], "id": "transacao_acid", "memoria": "semantica", "meta": {"dominio": "programacao", "fonte": ""}, "origem": "wiki", "palavras_chave": [], "sinonimos": [], "tipo": "conceito", "titulo": "Transações e ACID"}
\section{Transações e ACID}
Garantir que um grupo de operações seja Atômico, Consistente, Isolado e Durável — tudo-ou-nada, mesmo com falha ou concorrência.
\prov{relações: usa banco\_de\_dados; relaciona tolerancia\_a\_falhas}
\prov{proveniência: wiki \textperiodcentered{} fato \textperiodcentered{} confiança 1.0 \textperiodcentered{} domínio programacao \textperiodcentered{} história 6 versões no jornal}

%CRISTAL {"arestas": [{"alvo": "mecanismo_de_atencao", "confianca": 1.0, "epistemico": "fato", "origem": "wiki", "rel": "usa"}, {"alvo": "aprendizado_profundo", "confianca": 1.0, "epistemico": "fato", "origem": "wiki", "rel": "especializa"}], "confianca": 1.0, "contraexemplos": [], "descricao": "Vaswani et al. (2017): a arquitetura que trocou a recorrência por atenção pura ('attention is all you need') — a base de todo LLM.", "epistemico": "fato", "exemplos": [], "id": "transformer", "memoria": "semantica", "meta": {"autor": "Vaswani", "dominio": "programacao", "fonte": ""}, "origem": "wiki", "palavras_chave": [], "sinonimos": [], "tipo": "conceito", "titulo": "Transformer"}
\section{Transformer}
Vaswani et al. (2017): a arquitetura que trocou a recorrência por atenção pura ('attention is all you need') — a base de todo LLM.
\prov{relações: usa mecanismo\_de\_atencao; especializa aprendizado\_profundo}
\prov{proveniência: wiki \textperiodcentered{} fato \textperiodcentered{} confiança 1.0 \textperiodcentered{} domínio programacao \textperiodcentered{} história 3 versões no jornal}

%CRISTAL {"arestas": [{"alvo": "gradiente_descendente", "confianca": 1.0, "epistemico": "fato", "origem": "wiki", "rel": "explica"}, {"alvo": "lyapunov_da_orbita", "confianca": 1.0, "epistemico": "fato", "origem": "wiki", "rel": "eh_um"}, {"alvo": "corpo_mineral", "confianca": 1.0, "epistemico": "fato", "origem": "wiki", "rel": "usa"}, {"alvo": "aprendizado_profundo", "confianca": 1.0, "epistemico": "fato", "origem": "wiki", "rel": "relaciona"}], "confianca": 1.0, "contraexemplos": [], "descricao": "Treinar é descer o gradiente da perda até um mínimo — um ponto fixo/ATRATOR. A taxa de aprendizado lr é o Lyapunov do treino: em f=x² o passo é x←(1−2·lr)x, λ=log|1−2·lr|. lr pequena converge (cristal), lr grande diverge (caos), no meio a borda rápida. O treino é uma órbita que busca o cristal. Verif.: lr=0.1 λ<0 (converge), lr=1.05 λ>0 (diverge).", "epistemico": "inferencia", "exemplos": [], "id": "treino_atrator_lyapunov", "memoria": "semantica", "meta": {"dominio": "ia", "fonte": "IA e modelos de linguagem", "lambda_lr01": -0.223, "lambda_lr105": 0.095}, "origem": "wiki", "palavras_chave": [], "sinonimos": [], "tipo": "conceito", "titulo": "Treino = Descida a um Atrator (a taxa é o Lyapunov)"}
\section{Treino = Descida a um Atrator (a taxa é o Lyapunov)}
Treinar é descer o gradiente da perda até um mínimo — um ponto fixo/ATRATOR. A taxa de aprendizado lr é o Lyapunov do treino: em f=x\ensuremath{^{2}} o passo é x\ensuremath{\leftarrow}(1\ensuremath{-}2\textperiodcentered lr)x, \ensuremath{\lambda}=log|1\ensuremath{-}2\textperiodcentered lr|. lr pequena converge (cristal), lr grande diverge (caos), no meio a borda rápida. O treino é uma órbita que busca o cristal. Verif.: lr=0.1 \ensuremath{\lambda}<0 (converge), lr=1.05 \ensuremath{\lambda}>0 (diverge).
\prov{relações: explica gradiente\_descendente; eh\_um lyapunov\_da\_orbita; usa corpo\_mineral; relaciona aprendizado\_profundo}
\prov{proveniência: wiki \textperiodcentered{} IA e modelos de linguagem \textperiodcentered{} inferencia \textperiodcentered{} confiança 1.0 \textperiodcentered{} domínio ia \textperiodcentered{} história 5 versões no jornal}

%CRISTAL {"arestas": [{"alvo": "seguranca_da_informacao", "confianca": 1.0, "epistemico": "fato", "origem": "wiki", "rel": "explica"}], "confianca": 1.0, "contraexemplos": [], "descricao": "Os três pilares: Confidencialidade (só quem deve lê), Integridade (não foi alterado) e Disponibilidade (está acessível). O que 'seguro' significa.", "epistemico": "inferencia", "exemplos": [], "id": "triade_cia", "memoria": "semantica", "meta": {"dominio": "programacao", "fonte": ""}, "origem": "wiki", "palavras_chave": [], "sinonimos": [], "tipo": "conceito", "titulo": "Tríade CIA"}
\section{Tríade CIA}
Os três pilares: Confidencialidade (só quem deve lê), Integridade (não foi alterado) e Disponibilidade (está acessível). O que 'seguro' significa.
\prov{relações: explica seguranca\_da\_informacao}
\prov{proveniência: wiki \textperiodcentered{} inferencia \textperiodcentered{} confiança 1.0 \textperiodcentered{} domínio programacao \textperiodcentered{} história 2 versões no jornal}

%CRISTAL {"arestas": [{"alvo": "sistema_de_arquivos", "confianca": 1.0, "epistemico": "fato", "origem": "wiki", "rel": "usa"}], "confianca": 1.0, "contraexemplos": [], "descricao": "No Unix, dispositivos, pipes e até processos aparecem como arquivos — uma interface única (ler/escrever) para tudo. Elegância por uniformidade.", "epistemico": "inferencia", "exemplos": [], "id": "tudo_e_arquivo", "memoria": "semantica", "meta": {"dominio": "programacao", "fonte": ""}, "origem": "wiki", "palavras_chave": [], "sinonimos": [], "tipo": "conceito", "titulo": "Tudo é Arquivo"}
\section{Tudo é Arquivo}
No Unix, dispositivos, pipes e até processos aparecem como arquivos — uma interface única (ler/escrever) para tudo. Elegância por uniformidade.
\prov{relações: usa sistema\_de\_arquivos}
\prov{proveniência: wiki \textperiodcentered{} inferencia \textperiodcentered{} confiança 1.0 \textperiodcentered{} domínio programacao \textperiodcentered{} história 2 versões no jornal}

%CRISTAL {"arestas": [{"alvo": "maquina_de_turing", "confianca": 1.0, "epistemico": "fato", "origem": "wiki", "rel": "usa"}, {"alvo": "hierarquia_de_chomsky", "confianca": 1.0, "epistemico": "fato", "origem": "wiki", "rel": "usa"}, {"alvo": "expressividade_sub_turing", "confianca": 1.0, "epistemico": "fato", "origem": "wiki", "rel": "contradiz"}], "confianca": 1.0, "contraexemplos": [], "descricao": "A capacidade de computar tudo que uma máquina de Turing computa — o teto da hierarquia de Chomsky; quase toda linguagem a tem.", "epistemico": "fato", "exemplos": [], "id": "turing_completude", "memoria": "semantica", "meta": {"dominio": "programacao", "fonte": ""}, "origem": "wiki", "palavras_chave": [], "sinonimos": [], "tipo": "conceito", "titulo": "Turing-completude"}
\section{Turing-completude}
A capacidade de computar tudo que uma máquina de Turing computa — o teto da hierarquia de Chomsky; quase toda linguagem a tem.
\prov{relações: usa maquina\_de\_turing; usa hierarquia\_de\_chomsky; contradiz expressividade\_sub\_turing}
\prov{proveniência: wiki \textperiodcentered{} fato \textperiodcentered{} confiança 1.0 \textperiodcentered{} domínio programacao \textperiodcentered{} história 8 versões no jornal}

%CRISTAL {"arestas": [{"alvo": "maquina_de_turing", "confianca": 1.0, "epistemico": "fato", "origem": "wiki", "rel": "especializa"}, {"alvo": "computacao", "confianca": 1.0, "epistemico": "fato", "origem": "wiki", "rel": "parte_de"}], "confianca": 1.0, "contraexemplos": [], "descricao": "A máquina de Turing (fita infinita, cabeçote, estados) formaliza 'algoritmo/computável'; a TESE de Church-Turing (não é teorema — 'efetivamente calculável' é informal) diz que ela captura tudo o que é mecanicamente calculável. O PROBLEMA DA PARADA é INDECIDÍVEL (não há algoritmo geral que decida se um programa para; por diagonalização) — o Entscheidungsproblem cai como corolário. Turing 1936.", "epistemico": "fato", "exemplos": [], "id": "turing_halting_church", "memoria": "semantica", "meta": {"autor": "broca-so", "dominio": "computacao", "fonte": "Turing 1936 Proc.LMS; Church 1936"}, "origem": "wiki", "palavras_chave": [], "sinonimos": [], "tipo": "conceito", "titulo": "Máquina de Turing, tese de Church-Turing e o halting"}
\section{Máquina de Turing, tese de Church-Turing e o halting}
A máquina de Turing (fita infinita, cabeçote, estados) formaliza 'algoritmo/computável'; a TESE de Church-Turing (não é teorema — 'efetivamente calculável' é informal) diz que ela captura tudo o que é mecanicamente calculável. O PROBLEMA DA PARADA é INDECIDÍVEL (não há algoritmo geral que decida se um programa para; por diagonalização) — o Entscheidungsproblem cai como corolário. Turing 1936.
\prov{relações: especializa maquina\_de\_turing; parte\_de computacao}
\prov{proveniência: wiki \textperiodcentered{} Turing 1936 Proc.LMS; Church 1936 \textperiodcentered{} fato \textperiodcentered{} confiança 1.0 \textperiodcentered{} domínio computacao \textperiodcentered{} história 3 versões no jornal}

%CRISTAL {"arestas": [{"alvo": "javascript", "confianca": 1.0, "epistemico": "fato", "origem": "wiki", "rel": "especializa"}, {"alvo": "tipagem_estatica_dinamica", "confianca": 1.0, "epistemico": "fato", "origem": "wiki", "rel": "usa"}], "confianca": 1.0, "contraexemplos": [], "descricao": "Microsoft: JavaScript com tipos estáticos opcionais — segurança e ferramentas para bases de código grandes.", "epistemico": "fato", "exemplos": [], "id": "typescript", "memoria": "semantica", "meta": {"autor": "Hejlsberg", "dominio": "programacao", "fonte": ""}, "origem": "wiki", "palavras_chave": [], "sinonimos": [], "tipo": "conceito", "titulo": "TypeScript"}
\section{TypeScript}
Microsoft: JavaScript com tipos estáticos opcionais — segurança e ferramentas para bases de código grandes.
\prov{relações: especializa javascript; usa tipagem\_estatica\_dinamica}
\prov{proveniência: wiki \textperiodcentered{} fato \textperiodcentered{} confiança 1.0 \textperiodcentered{} domínio programacao \textperiodcentered{} história 6 versões no jornal}

%CRISTAL {"arestas": [{"alvo": "isa_da_broca", "confianca": 1.0, "epistemico": "fato", "origem": "wiki", "rel": "usa"}, {"alvo": "word_espectral", "confianca": 1.0, "epistemico": "fato", "origem": "wiki", "rel": "usa"}], "confianca": 1.0, "contraexemplos": [], "descricao": "ADD/SUB somam total e e independentemente; AND/OR/XOR fazem bitwise em cada componente. O LOAD empilha (B←A antes de ler o novo slot), então 'LOAD;LOAD;ADD' soma dois espectros sem instrução de MOV; o loop é fetch→decode→exec (cpu_step) até HALT.", "epistemico": "fato", "exemplos": [], "id": "ula_por_componente", "memoria": "semantica", "meta": {"autor": "broca-so", "dominio": "operacao_so", "fonte": "ula/ula.c; ula/cpu.c; ula/instrucoes.c"}, "origem": "wiki", "palavras_chave": [], "sinonimos": [], "tipo": "conceito", "titulo": "A ULA por Componente (e o LOAD que empilha)"}
\section{A ULA por Componente (e o LOAD que empilha)}
ADD/SUB somam total e e independentemente; AND/OR/XOR fazem bitwise em cada componente. O LOAD empilha (B\ensuremath{\leftarrow}A antes de ler o novo slot), então 'LOAD;LOAD;ADD' soma dois espectros sem instrução de MOV; o loop é fetch\ensuremath{\to}decode\ensuremath{\to}exec (cpu\_step) até HALT.
\prov{relações: usa isa\_da\_broca; usa word\_espectral}
\prov{proveniência: wiki \textperiodcentered{} ula/ula.c; ula/cpu.c; ula/instrucoes.c \textperiodcentered{} fato \textperiodcentered{} confiança 1.0 \textperiodcentered{} domínio operacao\_so \textperiodcentered{} história 3 versões no jornal}

%CRISTAL {"arestas": [{"alvo": "sistema_operacional", "confianca": 1.0, "epistemico": "fato", "origem": "wiki", "rel": "especializa"}, {"alvo": "c", "confianca": 1.0, "epistemico": "fato", "origem": "wiki", "rel": "usa"}], "confianca": 1.0, "contraexemplos": [], "descricao": "Thompson e Ritchie (Bell Labs, 1969): o SO que fundou uma tradição — processos, arquivos, o shell, escrito em C. O ancestral de quase tudo.", "epistemico": "fato", "exemplos": [], "id": "unix", "memoria": "semantica", "meta": {"autor": "Thompson/Ritchie", "dominio": "programacao", "fonte": ""}, "origem": "wiki", "palavras_chave": [], "sinonimos": [], "tipo": "conceito", "titulo": "Unix"}
\section{Unix}
Thompson e Ritchie (Bell Labs, 1969): o SO que fundou uma tradição — processos, arquivos, o shell, escrito em C. O ancestral de quase tudo.
\prov{relações: especializa sistema\_operacional; usa c}
\prov{proveniência: wiki \textperiodcentered{} fato \textperiodcentered{} confiança 1.0 \textperiodcentered{} domínio programacao \textperiodcentered{} história 3 versões no jornal}

%CRISTAL {"arestas": [{"alvo": "unix", "confianca": 1.0, "epistemico": "fato", "origem": "wiki", "rel": "usa"}, {"alvo": "windows", "confianca": 1.0, "epistemico": "fato", "origem": "wiki", "rel": "usa"}, {"alvo": "filosofia_unix", "confianca": 1.0, "epistemico": "fato", "origem": "wiki", "rel": "usa"}], "confianca": 1.0, "contraexemplos": [], "descricao": "Duas filosofias de SO: Unix (tudo é arquivo, texto nos pipes, ferramentas pequenas) vs Windows (objetos/API, registro, GUI-primeiro). Design oposto.", "epistemico": "inferencia", "exemplos": [], "id": "unix_vs_windows", "memoria": "semantica", "meta": {"dominio": "programacao", "fonte": ""}, "origem": "wiki", "palavras_chave": [], "sinonimos": [], "tipo": "conceito", "titulo": "Unix vs Windows (as duas tradições)"}
\section{Unix vs Windows (as duas tradições)}
Duas filosofias de SO: Unix (tudo é arquivo, texto nos pipes, ferramentas pequenas) vs Windows (objetos/API, registro, GUI-primeiro). Design oposto.
\prov{relações: usa unix; usa windows; usa filosofia\_unix}
\prov{proveniência: wiki \textperiodcentered{} inferencia \textperiodcentered{} confiança 1.0 \textperiodcentered{} domínio programacao \textperiodcentered{} história 4 versões no jornal}

%CRISTAL {"arestas": [{"alvo": "arquivo_de_configuracao", "confianca": 1.0, "epistemico": "fato", "origem": "wiki", "rel": "relaciona"}, {"alvo": "shell", "confianca": 1.0, "epistemico": "fato", "origem": "wiki", "rel": "usa"}, {"alvo": "processo", "confianca": 1.0, "epistemico": "fato", "origem": "wiki", "rel": "usa"}], "confianca": 1.0, "contraexemplos": [], "descricao": "Um par nome=valor que o SO passa aos processos — configura sem tocar no código; a ponte entre o shell e o programa (PATH, HOME).", "epistemico": "fato", "exemplos": [], "id": "variavel_de_ambiente", "memoria": "semantica", "meta": {"dominio": "programacao", "fonte": ""}, "origem": "wiki", "palavras_chave": [], "sinonimos": [], "tipo": "conceito", "titulo": "Variável de Ambiente"}
\section{Variável de Ambiente}
Um par nome=valor que o SO passa aos processos — configura sem tocar no código; a ponte entre o shell e o programa (PATH, HOME).
\prov{relações: relaciona arquivo\_de\_configuracao; usa shell; usa processo}
\prov{proveniência: wiki \textperiodcentered{} fato \textperiodcentered{} confiança 1.0 \textperiodcentered{} domínio programacao \textperiodcentered{} história 4 versões no jornal}

%CRISTAL {"arestas": [{"alvo": "linguagem_descricao_hardware", "confianca": 1.0, "epistemico": "fato", "origem": "wiki", "rel": "especializa"}], "confianca": 1.0, "contraexemplos": [], "descricao": "A HDL de sintaxe estilo-C para descrever e simular circuitos digitais — a mais usada na indústria de chips; base do SystemVerilog.", "epistemico": "fato", "exemplos": [], "id": "verilog", "memoria": "semantica", "meta": {"autor": "Moorby", "dominio": "programacao", "fonte": ""}, "origem": "wiki", "palavras_chave": [], "sinonimos": [], "tipo": "conceito", "titulo": "Verilog"}
\section{Verilog}
A HDL de sintaxe estilo-C para descrever e simular circuitos digitais — a mais usada na indústria de chips; base do SystemVerilog.
\prov{relações: especializa linguagem\_descricao\_hardware}
\prov{proveniência: wiki \textperiodcentered{} fato \textperiodcentered{} confiança 1.0 \textperiodcentered{} domínio programacao \textperiodcentered{} história 2 versões no jornal}

%CRISTAL {"arestas": [{"alvo": "linguagem_descricao_hardware", "confianca": 1.0, "epistemico": "fato", "origem": "wiki", "rel": "especializa"}], "confianca": 1.0, "contraexemplos": [], "descricao": "A HDL fortemente tipada e verbosa (origem no Departamento de Defesa dos EUA) — rigorosa e verificável, forte na Europa e na aeroespacial.", "epistemico": "fato", "exemplos": [], "id": "vhdl", "memoria": "semantica", "meta": {"dominio": "programacao", "fonte": ""}, "origem": "wiki", "palavras_chave": [], "sinonimos": [], "tipo": "conceito", "titulo": "VHDL"}
\section{VHDL}
A HDL fortemente tipada e verbosa (origem no Departamento de Defesa dos EUA) — rigorosa e verificável, forte na Europa e na aeroespacial.
\prov{relações: especializa linguagem\_descricao\_hardware}
\prov{proveniência: wiki \textperiodcentered{} fato \textperiodcentered{} confiança 1.0 \textperiodcentered{} domínio programacao \textperiodcentered{} história 2 versões no jornal}

%CRISTAL {"arestas": [{"alvo": "aprendizado_de_maquina", "confianca": 1.0, "epistemico": "fato", "origem": "wiki", "rel": "usa"}, {"alvo": "injustica_epistemica", "confianca": 1.0, "epistemico": "fato", "origem": "wiki", "rel": "relaciona"}], "confianca": 1.0, "contraexemplos": [], "descricao": "O modelo herda e amplifica os vieses dos dados — 'lixo entra, lixo sai' com escala e verniz de objetividade. Justiça não é automática.", "epistemico": "inferencia", "exemplos": [], "id": "vies_algoritmico", "memoria": "semantica", "meta": {"dominio": "programacao", "fonte": ""}, "origem": "wiki", "palavras_chave": [], "sinonimos": [], "tipo": "conceito", "titulo": "Viés Algorítmico"}
\section{Viés Algorítmico}
O modelo herda e amplifica os vieses dos dados — 'lixo entra, lixo sai' com escala e verniz de objetividade. Justiça não é automática.
\prov{relações: usa aprendizado\_de\_maquina; relaciona injustica\_epistemica}
\prov{proveniência: wiki \textperiodcentered{} inferencia \textperiodcentered{} confiança 1.0 \textperiodcentered{} domínio programacao \textperiodcentered{} história 3 versões no jornal}

%CRISTAL {"arestas": [{"alvo": "framework_frontend", "confianca": 1.0, "epistemico": "fato", "origem": "wiki", "rel": "especializa"}, {"alvo": "reatividade", "confianca": 1.0, "epistemico": "fato", "origem": "wiki", "rel": "usa"}], "confianca": 1.0, "contraexemplos": [], "descricao": "Evan You: um framework progressivo de reatividade, adotável aos poucos — o meio-termo entre a liberdade do React e a estrutura do Angular.", "epistemico": "fato", "exemplos": [], "id": "vue", "memoria": "semantica", "meta": {"autor": "Evan You", "dominio": "programacao", "fonte": ""}, "origem": "wiki", "palavras_chave": [], "sinonimos": [], "tipo": "conceito", "titulo": "Vue.js"}
\section{Vue.js}
Evan You: um framework progressivo de reatividade, adotável aos poucos — o meio-termo entre a liberdade do React e a estrutura do Angular.
\prov{relações: especializa framework\_frontend; usa reatividade}
\prov{proveniência: wiki \textperiodcentered{} fato \textperiodcentered{} confiança 1.0 \textperiodcentered{} domínio programacao \textperiodcentered{} história 6 versões no jornal}

%CRISTAL {"arestas": [{"alvo": "seguranca_da_informacao", "confianca": 1.0, "epistemico": "fato", "origem": "wiki", "rel": "parte_de"}], "confianca": 1.0, "contraexemplos": [], "descricao": "A vulnerabilidade é a falha; o exploit é o código que a explora. A corrida entre quem fecha e quem acha o buraco.", "epistemico": "fato", "exemplos": [], "id": "vulnerabilidade_exploit", "memoria": "semantica", "meta": {"dominio": "programacao", "fonte": ""}, "origem": "wiki", "palavras_chave": [], "sinonimos": [], "tipo": "conceito", "titulo": "Vulnerabilidade e Exploit"}
\section{Vulnerabilidade e Exploit}
A vulnerabilidade é a falha; o exploit é o código que a explora. A corrida entre quem fecha e quem acha o buraco.
\prov{relações: parte\_de seguranca\_da\_informacao}
\prov{proveniência: wiki \textperiodcentered{} fato \textperiodcentered{} confiança 1.0 \textperiodcentered{} domínio programacao \textperiodcentered{} história 2 versões no jornal}

%CRISTAL {"arestas": [{"alvo": "sistema_operacional", "confianca": 1.0, "epistemico": "fato", "origem": "wiki", "rel": "especializa"}], "confianca": 1.0, "contraexemplos": [], "descricao": "Microsoft: o SO de desktop mais usado; kernel NT, registro em vez de arquivos de config, API Win32 — a tradição rival da Unix.", "epistemico": "fato", "exemplos": [], "id": "windows", "memoria": "semantica", "meta": {"autor": "Microsoft", "dominio": "programacao", "fonte": ""}, "origem": "wiki", "palavras_chave": [], "sinonimos": [], "tipo": "conceito", "titulo": "Windows"}
\section{Windows}
Microsoft: o SO de desktop mais usado; kernel NT, registro em vez de arquivos de config, API Win32 — a tradição rival da Unix.
\prov{relações: especializa sistema\_operacional}
\prov{proveniência: wiki \textperiodcentered{} fato \textperiodcentered{} confiança 1.0 \textperiodcentered{} domínio programacao \textperiodcentered{} história 2 versões no jornal}

%CRISTAL {"arestas": [{"alvo": "maquina_algebrica", "confianca": 1.0, "epistemico": "fato", "origem": "wiki", "rel": "usa"}, {"alvo": "neuronio_stream_bit", "confianca": 1.0, "epistemico": "fato", "origem": "wiki", "rel": "usa"}], "confianca": 1.0, "contraexemplos": [], "descricao": "O único tipo de dado da ULA é Wordlong total; long e (com o=total−e implícito); a ULA não vê JSON nem objetos, só words de dois inteiros. O operador do gato reaparece na cifra como [n·total+e, total].", "epistemico": "fato", "exemplos": [], "id": "word_espectral", "memoria": "semantica", "meta": {"autor": "broca-so", "dominio": "operacao_so", "fonte": "ula/word.h"}, "origem": "wiki", "palavras_chave": [], "sinonimos": [], "tipo": "conceito", "titulo": "Word [total, e] — a palavra da máquina"}
\section{Word [total, e] — a palavra da máquina}
O único tipo de dado da ULA é Wordlong total; long e (com o=total\ensuremath{-}e implícito); a ULA não vê JSON nem objetos, só words de dois inteiros. O operador do gato reaparece na cifra como [n\textperiodcentered total+e, total].
\prov{relações: usa maquina\_algebrica; usa neuronio\_stream\_bit}
\prov{proveniência: wiki \textperiodcentered{} ula/word.h \textperiodcentered{} fato \textperiodcentered{} confiança 1.0 \textperiodcentered{} domínio operacao\_so \textperiodcentered{} história 4 versões no jornal}

%CRISTAL {"arestas": [{"alvo": "linguagem_de_marcacao", "confianca": 1.0, "epistemico": "fato", "origem": "wiki", "rel": "usa"}], "confianca": 1.0, "contraexemplos": [], "descricao": "Editar vendo o resultado final (Word) OU escrever a intenção em marcas e deixar o renderizador compor (LaTeX, Markdown). Controle vs conveniência.", "epistemico": "inferencia", "exemplos": [], "id": "wysiwyg_vs_marcacao", "memoria": "semantica", "meta": {"dominio": "programacao", "fonte": ""}, "origem": "wiki", "palavras_chave": [], "sinonimos": [], "tipo": "conceito", "titulo": "WYSIWYG vs Marcação"}
\section{WYSIWYG vs Marcação}
Editar vendo o resultado final (Word) OU escrever a intenção em marcas e deixar o renderizador compor (LaTeX, Markdown). Controle vs conveniência.
\prov{relações: usa linguagem\_de\_marcacao}
\prov{proveniência: wiki \textperiodcentered{} inferencia \textperiodcentered{} confiança 1.0 \textperiodcentered{} domínio programacao \textperiodcentered{} história 2 versões no jornal}

%CRISTAL {"arestas": [{"alvo": "corpo_mineral", "confianca": 1.0, "epistemico": "fato", "origem": "wiki", "rel": "usa"}, {"alvo": "lyapunov_da_orbita", "confianca": 1.0, "epistemico": "fato", "origem": "wiki", "rel": "usa"}, {"alvo": "orbita_de_partida", "confianca": 1.0, "epistemico": "fato", "origem": "wiki", "rel": "explica"}, {"alvo": "xadrez_bai", "confianca": 1.0, "epistemico": "fato", "origem": "wiki", "rel": "relaciona"}, {"alvo": "transformada_metalica", "confianca": 1.0, "epistemico": "fato", "origem": "wiki", "rel": "usa"}], "confianca": 1.0, "contraexemplos": [], "descricao": "Analisar uma partida pelo corpo mineral, direto dos FENs guardados (sem minimax): o sinal Φ por ply é uma série cujo espectro (DMD/matriz companheira) dá a assinatura mineral — entropia H=log β, o modo espiral (λ,θ) e a classe cristal/borda/caos. Custo O(n) por partida (uma SVD), não a árvore. Achados: mate do pastor tem a maior entropia (decisão mais rápida); a Tiffany δ=4 é a mais contraente (cristal, moagem técnica); Deep Blue é espiral. Economiza o tempo do minimax de δ alto.", "epistemico": "inferencia", "exemplos": [], "id": "xadrez_mineral", "memoria": "semantica", "meta": {"dominio": "ia", "fonte": "xadrez mineral"}, "origem": "wiki", "palavras_chave": [], "sinonimos": [], "tipo": "conceito", "titulo": "Xadrez pelas Ferramentas Minerais"}
\section{Xadrez pelas Ferramentas Minerais}
Analisar uma partida pelo corpo mineral, direto dos FENs guardados (sem minimax): o sinal \ensuremath{\Phi} por ply é uma série cujo espectro (DMD/matriz companheira) dá a assinatura mineral — entropia H=log \ensuremath{\beta}, o modo espiral (\ensuremath{\lambda},\ensuremath{\theta}) e a classe cristal/borda/caos. Custo O(n) por partida (uma SVD), não a árvore. Achados: mate do pastor tem a maior entropia (decisão mais rápida); a Tiffany \ensuremath{\delta}=4 é a mais contraente (cristal, moagem técnica); Deep Blue é espiral. Economiza o tempo do minimax de \ensuremath{\delta} alto.
\prov{relações: usa corpo\_mineral; usa lyapunov\_da\_orbita; explica orbita\_de\_partida; relaciona xadrez\_bai; usa transformada\_metalica}
\prov{proveniência: wiki \textperiodcentered{} xadrez mineral \textperiodcentered{} inferencia \textperiodcentered{} confiança 1.0 \textperiodcentered{} domínio ia \textperiodcentered{} história 12 versões no jornal}

%CRISTAL {"arestas": [{"alvo": "linguagem_de_marcacao", "confianca": 1.0, "epistemico": "fato", "origem": "wiki", "rel": "especializa"}], "confianca": 1.0, "contraexemplos": [], "descricao": "eXtensible Markup Language: marcação genérica e estrita para dados estruturados e autodescritivos — verboso, mas validável por esquema.", "epistemico": "fato", "exemplos": [], "id": "xml", "memoria": "semantica", "meta": {"dominio": "programacao", "fonte": ""}, "origem": "wiki", "palavras_chave": [], "sinonimos": [], "tipo": "conceito", "titulo": "XML"}
\section{XML}
eXtensible Markup Language: marcação genérica e estrita para dados estruturados e autodescritivos — verboso, mas validável por esquema.
\prov{relações: especializa linguagem\_de\_marcacao}
\prov{proveniência: wiki \textperiodcentered{} fato \textperiodcentered{} confiança 1.0 \textperiodcentered{} domínio programacao \textperiodcentered{} história 2 versões no jornal}

%CRISTAL {"arestas": [{"alvo": "ci_cd", "confianca": 1.0, "epistemico": "fato", "origem": "wiki", "rel": "usa"}, {"alvo": "kubernetes", "confianca": 1.0, "epistemico": "fato", "origem": "wiki", "rel": "usa"}, {"alvo": "arquivo_de_configuracao", "confianca": 1.0, "epistemico": "fato", "origem": "wiki", "rel": "especializa"}], "confianca": 1.0, "contraexemplos": [], "descricao": "Formato de serialização legível, baseado em indentação — a escolha dos arquivos de configuração (CI/CD, Kubernetes). Menos ruído que XML/JSON.", "epistemico": "fato", "exemplos": [], "id": "yaml", "memoria": "semantica", "meta": {"dominio": "programacao", "fonte": ""}, "origem": "wiki", "palavras_chave": [], "sinonimos": [], "tipo": "conceito", "titulo": "YAML"}
\section{YAML}
Formato de serialização legível, baseado em indentação — a escolha dos arquivos de configuração (CI/CD, Kubernetes). Menos ruído que XML/JSON.
\prov{relações: usa ci\_cd; usa kubernetes; especializa arquivo\_de\_configuracao}
\prov{proveniência: wiki \textperiodcentered{} fato \textperiodcentered{} confiança 1.0 \textperiodcentered{} domínio programacao \textperiodcentered{} história 4 versões no jornal}

%CRISTAL {"arestas": [{"alvo": "c", "confianca": 1.0, "epistemico": "fato", "origem": "wiki", "rel": "relaciona"}, {"alvo": "gerenciamento_de_memoria", "confianca": 1.0, "epistemico": "fato", "origem": "wiki", "rel": "usa"}], "confianca": 1.0, "contraexemplos": [], "descricao": "Um C moderno: sem GC, sem alocação escondida, comptime no lugar de macros — controle explícito e simplicidade.", "epistemico": "fato", "exemplos": [], "id": "zig", "memoria": "semantica", "meta": {"autor": "Kelley", "dominio": "programacao", "fonte": ""}, "origem": "wiki", "palavras_chave": [], "sinonimos": [], "tipo": "conceito", "titulo": "Zig"}
\section{Zig}
Um C moderno: sem GC, sem alocação escondida, comptime no lugar de macros — controle explícito e simplicidade.
\prov{relações: relaciona c; usa gerenciamento\_de\_memoria}
\prov{proveniência: wiki \textperiodcentered{} fato \textperiodcentered{} confiança 1.0 \textperiodcentered{} domínio programacao \textperiodcentered{} história 6 versões no jornal}

%CRISTAL {"arestas": [{"alvo": "shell", "confianca": 1.0, "epistemico": "fato", "origem": "wiki", "rel": "especializa"}, {"alvo": "sh_posix", "confianca": 1.0, "epistemico": "fato", "origem": "wiki", "rel": "relaciona"}], "confianca": 1.0, "contraexemplos": [], "descricao": "Um shell interativo poderoso (autocompletar, globbing, temas) compatível com sh — o preferido de muitos, base do Oh My Zsh.", "epistemico": "fato", "exemplos": [], "id": "zsh", "memoria": "semantica", "meta": {"dominio": "programacao", "fonte": ""}, "origem": "wiki", "palavras_chave": [], "sinonimos": [], "tipo": "conceito", "titulo": "Zsh"}
\section{Zsh}
Um shell interativo poderoso (autocompletar, globbing, temas) compatível com sh — o preferido de muitos, base do Oh My Zsh.
\prov{relações: especializa shell; relaciona sh\_posix}
\prov{proveniência: wiki \textperiodcentered{} fato \textperiodcentered{} confiança 1.0 \textperiodcentered{} domínio programacao \textperiodcentered{} história 3 versões no jornal}

%CRISTAL {"arestas": [{"alvo": "gerenciamento_de_estado", "confianca": 1.0, "epistemico": "fato", "origem": "wiki", "rel": "especializa"}, {"alvo": "redux", "confianca": 1.0, "epistemico": "fato", "origem": "wiki", "rel": "relaciona"}], "confianca": 1.0, "contraexemplos": [], "descricao": "Um gerenciador de estado minimalista para React — store por hook, sem boilerplate; a reação ao peso do Redux, mantendo a imutabilidade.", "epistemico": "fato", "exemplos": [], "id": "zustand", "memoria": "semantica", "meta": {"dominio": "programacao", "fonte": ""}, "origem": "wiki", "palavras_chave": [], "sinonimos": [], "tipo": "conceito", "titulo": "Zustand"}
\section{Zustand}
Um gerenciador de estado minimalista para React — store por hook, sem boilerplate; a reação ao peso do Redux, mantendo a imutabilidade.
\prov{relações: especializa gerenciamento\_de\_estado; relaciona redux}
\prov{proveniência: wiki \textperiodcentered{} fato \textperiodcentered{} confiança 1.0 \textperiodcentered{} domínio programacao \textperiodcentered{} história 3 versões no jornal}

\end{document}
